\documentclass[11pt,a4paper]{article}
\usepackage[utf8]{inputenc}
\usepackage[T5]{fontenc}
\usepackage[vietnamese]{babel}
\usepackage[top=1.5cm,bottom=1.5cm,left=1.5cm,right=1.5cm]{geometry}
\usepackage{amsmath,amssymb,amsfonts}
\usepackage{tikz,tkz-tab}
\usetikzlibrary{calc,angles,quotes,arrows,arrows.meta,patterns,intersections,shapes.geometric,decorations.markings,decorations.pathmorphing,backgrounds,positioning}
\usepackage[most]{tcolorbox}
\usepackage{longtable,booktabs,array,multirow}
\usepackage{hyperref}
\usepackage{fancyhdr}

\hypersetup{
    colorlinks=true,
    linkcolor=blue!70!black,
    urlcolor=blue!70!black
}

\pagestyle{fancy}
\fancyhf{}
\fancyhead[L]{\small\textbf{HỆ THỐNG QMATH LMS -- BẢNG THỐNG KÊ CÂU HỎI ĐỔI SỐ CÓ HÌNH TIKZ}}
\fancyhead[R]{\small\textbf{Trang \thepage}}
\setlength{\headheight}{14pt}
\setlength{\parindent}{0pt}
\setlength{\parskip}{4pt}

% Macro hỗ trợ một số lệnh từ gói ex_test
\newcommand{\True}{\textbf{[ĐÚNG]} }
\newcommand{\heva}[1]{\begin{cases}#1\end{cases}}
\newcommand{\hoac}[1]{\left[\begin{aligned}#1\end{aligned}\right.}
\providecommand{\wideparen}[1]{\widehat{#1}}

\begin{document}

\begin{center}
{\LARGE\bfseries BẢNG THỐNG KÊ \& MÃ ĐỊNH DANH CÁC CÂU HỎI ĐỔI SỐ CÓ HÌNH TIKZ}\\[6pt]
{\large Hệ thống QMath LMS -- Tổng số: \textbf{456 câu có hình} (\textbf{350 mẫu hình duy nhất})}
\end{center}
\begin{tcolorbox}[colback=yellow!12!white,colframe=orange!80!black,title=\textbf{HƯỚNG DẪN SỬ DỤNG BẢNG MÃ HÌNH TIKZ NÀY}]
\small
1. Mỗi câu hỏi có hình dưới đây đều được gắn một \textbf{MÃ ĐỊNH DANH DUY NHẤT} ở góc trái tiêu đề, ví dụ: \texttt{[L09\_C4\_B2\_57]} (kèm theo mã \texttt{MapID} của QMath và mã nhóm hình \texttt{HINH\_xxx}).\\
2. Khi muốn thay thế hoặc chỉnh sửa hình vẽ của bất kỳ câu nào, bạn chỉ cần gửi \textbf{Mã câu} (ví dụ: \texttt{L09\_C4\_B2\_12}) hoặc \textbf{Mã nhóm hình} (ví dụ: \texttt{HINH\_085}) kèm mã TikZ mới.\\
3. Những câu có cùng \textbf{Mã nhóm hình (\texttt{HINH\_xxx})} dùng chung 1 khung hình TikZ: chỉ cần thay 1 lần là toàn bộ các câu cùng nhóm sẽ tự động được cập nhật hình mới!
\end{tcolorbox}
\tableofcontents
\newpage


\section{Toán 6}

\subsection{Chương 4 - Một số hình trong thực tiễn}

\begin{tcolorbox}[
    breakable,
    colback=white,
    colframe=blue!60!black,
    boxrule=0.8pt,
    arc=3pt,
    title={\textbf{Câu \#001} \quad \textbf{MÃ CÂU: \texttt{[L06\_C4\_B1\_05]}} \quad | \quad \textbf{Nhóm: \texttt{HINH\_001}} \quad | \quad \textbf{MapID: \texttt{6H4N1}} \hfill \textcolor{yellow!90!white}{[Hình gốc .pyc]}}
]
\textbf{Dạng toán:} Tính diện tích tam giác \hfill \textit{(Trắc nghiệm -- Nhận biết)}
\vspace{4pt}\hrule\vspace{6pt}
\begin{minipage}[t]{0.58\textwidth}
\vspace{0pt}
Tính diện tích tam giác sau:

\vspace{4pt}\noindent
\textbf{\textcolor{blue!75!black}{A. (Đ)}} $65cm^{2}$ \quad \textbf{B.} $260cm^{2}$ \quad \textbf{C.} $\dfrac{65}{2}cm^{2}$ \quad \textbf{D.} $130cm^{2}$
\end{minipage}\hfill
\begin{minipage}[t]{0.39\textwidth}
\vspace{0pt}
\centering

\resizebox{\linewidth}{!}{%
\begin{tikzpicture}
    % Vẽ đoạn thẳng BC
    \draw[thick, blue] (0,0) -- (6,0) node[below right] {$K$};
    \draw[thick, blue] (0,0) -- (-3,0) node[below left] {$F$};
    
    % Vẽ đoạn thẳng AB
    \draw[thick, blue] (-3,0) -- (0,3) node[above] {$E$};
    
    % Vẽ đoạn thẳng AC
    \draw[thick, blue] (0,3) -- (6,0);
    
    % Vẽ đường cao AH
    \draw[thick, blue] (0,3) -- (0,0) node[below] {$H$};
    
    % Ghi nhãn chiều dài đoạn AH
    \node at (0.6,1.5) {$10cm$};
    
    % Ghi nhãn chiều dài đoạn BC
    \node at (3,-0.7) {$13cm$};
    
    % Ký hiệu góc vuông tại H
    \draw (0,0) rectangle (0.3,0.3);
    
    % Đánh dấu các điểm bằng chấm tròn nhỏ
    \filldraw[blue] (0,0) circle (2pt);   % Điểm H
    \filldraw[blue] (-3,0) circle (2pt);  % Điểm B
    \filldraw[blue] (6,0) circle (2pt);   % Điểm C
    \filldraw[blue] (0,3) circle (2pt);   % Điểm A
    \draw[<->]   (-3,-0.5)--(6,-0.5);
\end{tikzpicture}
}
\end{minipage}

\vspace{3pt}\hrule\vspace{3pt}\footnotesize\textcolor{gray!80!black}{\textbf{Các mã dùng chung hình (HINH\_001):} \texttt{L06\_C4\_B1\_04}}
\end{tcolorbox}
\vspace{4pt}


\begin{tcolorbox}[
    breakable,
    colback=white,
    colframe=blue!60!black,
    boxrule=0.8pt,
    arc=3pt,
    title={\textbf{Câu \#002} \quad \textbf{MÃ CÂU: \texttt{[L06\_C4\_B1\_07]}} \quad | \quad \textbf{Nhóm: \texttt{HINH\_002}} \quad | \quad \textbf{MapID: \texttt{6H4N1}} \hfill \textcolor{yellow!90!white}{[Hình gốc .pyc]}}
]
\textbf{Dạng toán:} Tính diện tích tam giác vuông \hfill \textit{(Trắc nghiệm -- Nhận biết)}
\vspace{4pt}\hrule\vspace{6pt}
\begin{minipage}[t]{0.58\textwidth}
\vspace{0pt}
Tính diện tích tam giác sau:

\vspace{4pt}\noindent
\textbf{\textcolor{blue!75!black}{A. (Đ)}} $63cm^{2}$ \quad \textbf{B.} $\dfrac{63}{2}cm^{2}$ \quad \textbf{C.} $126cm^{2}$ \quad \textbf{D.} $252cm^{2}$
\end{minipage}\hfill
\begin{minipage}[t]{0.39\textwidth}
\vspace{0pt}
\centering

\resizebox{\linewidth}{!}{%
\begin{tikzpicture}
    % Vẽ cạnh DE (vuông góc)
    \draw[blue, thick] (0,0) -- (0,4) node[midway, left] {$18cm$};
    
    % Vẽ cạnh DF (nằm ngang)
    \draw[blue, thick] (0,0) -- (5,0) node[midway, below] {$7cm$};
    
    % Vẽ cạnh EF (chéo)
    \draw[blue, thick] (0,4) -- (5,0);
    
    % Đánh dấu góc vuông tại D
    \draw[blue] (0,0) rectangle (0.4,0.4);
    
    % Đặt các điểm D, E, F với chấm tròn nhỏ
    \filldraw[blue] (0,0) circle (2pt) node[below left] {$M$};
    \filldraw[blue] (0,4) circle (2pt) node[above left] {$N$};
    \filldraw[blue] (5,0) circle (2pt) node[below right] {$P$};
\end{tikzpicture}
}
\end{minipage}

\vspace{3pt}\hrule\vspace{3pt}\footnotesize\textcolor{gray!80!black}{\textbf{Các mã dùng chung hình (HINH\_002):} \texttt{L06\_C4\_B1\_06}}
\end{tcolorbox}
\vspace{4pt}


\begin{tcolorbox}[
    breakable,
    colback=white,
    colframe=blue!60!black,
    boxrule=0.8pt,
    arc=3pt,
    title={\textbf{Câu \#003} \quad \textbf{MÃ CÂU: \texttt{[L06\_C4\_B1\_09]}} \quad | \quad \textbf{Nhóm: \texttt{HINH\_003}} \quad | \quad \textbf{MapID: \texttt{6H4N1}} \hfill \textcolor{yellow!90!white}{[Hình gốc .pyc]}}
]
\textbf{Dạng toán:} Tính diện tích tứ giác có 2 góc vuông đối diện. \hfill \textit{(Trắc nghiệm -- Nhận biết)}
\vspace{4pt}\hrule\vspace{6pt}
\begin{minipage}[t]{0.58\textwidth}
\vspace{0pt}
Tính diện tích hình sau biết ${QK}=8cm$, ${EK}=18cm$, ${QF}=17cm$, ${EF}=10cm$

\vspace{4pt}\noindent
\textbf{A.} $\dfrac{157}{2}cm^{2}$ \quad \textbf{B.} $628cm^{2}$ \quad \textbf{\textcolor{blue!75!black}{C. (Đ)}} $157cm^{2}$ \quad \textbf{D.} $314cm^{2}$
\end{minipage}\hfill
\begin{minipage}[t]{0.39\textwidth}
\vspace{0pt}
\centering

\resizebox{\linewidth}{!}{%
\begin{tikzpicture}
    % Định nghĩa tọa độ của các điểm A và C (đường kính)
    \coordinate (A) at (0,0);
    \coordinate (C) at (6,0);
    \coordinate (B) at (1,2.236);  % Điểm B nằm trên đường tròn
    \coordinate (D) at (4,-2.828); % Điểm D nằm trên đường tròn
    
    \draw[blue, thick] (A) -- (B) -- (C) -- (D) -- cycle;
    
    \def\khvuong[size=#1](#2,#3,#4){%
        \draw ($(#3)!#1!(#2)$) --
        ($($(#3)!#1!(#2)$)+($(#3)!#1!(#4)$)-(#3)$)--
        ($(#3)!#1!(#4)$);
    }
    \khvuong[size=7pt](C,B,A);
        \khvuong[size=7pt](C,D,A);
    \filldraw[black] (A) circle (2pt) node[below left] {$Q$};
    \filldraw[black] (B) circle (2pt) node[above left] {$K$};
    \filldraw[black] (C) circle (2pt) node[below right] {$E$};
    \filldraw[black] (D) circle (2pt) node[below right] {$F$};
\end{tikzpicture}
}
\end{minipage}

\vspace{3pt}\hrule\vspace{3pt}\footnotesize\textcolor{gray!80!black}{\textbf{Các mã dùng chung hình (HINH\_003):} \texttt{L06\_C4\_B1\_08}}
\end{tcolorbox}
\vspace{4pt}


\begin{tcolorbox}[
    breakable,
    colback=white,
    colframe=blue!60!black,
    boxrule=0.8pt,
    arc=3pt,
    title={\textbf{Câu \#004} \quad \textbf{MÃ CÂU: \texttt{[L06\_C4\_B1\_01]}} \quad | \quad \textbf{Nhóm: \texttt{HINH\_004}} \quad | \quad \textbf{MapID: \texttt{6H4N1}} \hfill \textcolor{yellow!90!white}{[Hình gốc .pyc]}}
]
\textbf{Dạng toán:} Nhận diện tam giác đều \hfill \textit{(Tự luận -- Nhận biết)}
\vspace{4pt}\hrule\vspace{6pt}
\begin{minipage}[t]{0.58\textwidth}
\vspace{0pt}
Hình nào là tam giác đều?
\end{minipage}\hfill
\begin{minipage}[t]{0.39\textwidth}
\vspace{0pt}
\centering

\resizebox{\linewidth}{!}{%
\begin{tikzpicture}
    % Draw the equilateral triangle ACB
    \draw[thick] (0,0) -- (5,0) -- (2.5,4.33) -- cycle;
    
    % Label the vertices of the equilateral triangle
    \node[below left] at (0,0) {$Q$};
    \node[below right] at (5,0) {$H$};
    \node[above] at (2.5,4.33) {$K$};
    
    % Mark equal angles with arcs in equilateral triangle
    \draw (0.5,0) arc[start angle=0,end angle=60,radius=0.5];     % Arc for angle at A
    \draw (4.5,0) arc[start angle=180,end angle=120,radius=0.5];  % Arc for angle at C
    \draw (2.25,3.9) arc[start angle=240,end angle=300,radius=0.5];  % Arc for angle at B
    
    % Draw the obtuse triangle DEF at specific coordinates
    \draw[thick] (7,1) -- (12,1) -- (9,4) -- cycle;
    
    % Label the vertices of the obtuse triangle
    \node[below left] at (7,1) {$M$};
    \node[below right] at (12,1) {$N$};
    \node[above] at (9,4) {$I$};
    
    % Mark equal side notations (parallel lines and single line slashes)
    \node at ($(7,1)!0.5!(9,4)$) [rotate=40] {$//$};   % Symbol for side DF
    \node at ($(9,4)!0.5!(12,1)$) [rotate=20] {$/$};   % Symbol for side EF
    
\end{tikzpicture}
}
\end{minipage}

\end{tcolorbox}
\vspace{4pt}


\begin{tcolorbox}[
    breakable,
    colback=white,
    colframe=blue!60!black,
    boxrule=0.8pt,
    arc=3pt,
    title={\textbf{Câu \#005} \quad \textbf{MÃ CÂU: \texttt{[L06\_C4\_B1\_04]}} \quad | \quad \textbf{Nhóm: \texttt{HINH\_001}} \quad | \quad \textbf{MapID: \texttt{6H4N1}} \hfill \textcolor{yellow!90!white}{[Hình gốc .pyc]}}
]
\textbf{Dạng toán:} Tính diện tích tam giác \hfill \textit{(Tự luận -- Nhận biết)}
\vspace{4pt}\hrule\vspace{6pt}
\begin{minipage}[t]{0.58\textwidth}
\vspace{0pt}
Tính diện tích tam giác sau:
\end{minipage}\hfill
\begin{minipage}[t]{0.39\textwidth}
\vspace{0pt}
\centering

\resizebox{\linewidth}{!}{%
\begin{tikzpicture}
    % Vẽ đoạn thẳng BC
    \draw[thick, blue] (0,0) -- (6,0) node[below right] {$C$};
    \draw[thick, blue] (0,0) -- (-3,0) node[below left] {$B$};
    
    % Vẽ đoạn thẳng AB
    \draw[thick, blue] (-3,0) -- (0,3) node[above] {$A$};
    
    % Vẽ đoạn thẳng AC
    \draw[thick, blue] (0,3) -- (6,0);
    
    % Vẽ đường cao AH
    \draw[thick, blue] (0,3) -- (0,0) node[below] {$H$};
    
    % Ghi nhãn chiều dài đoạn AH
    \node at (0.6,1.5) {$11cm$};
    
    % Ghi nhãn chiều dài đoạn BC
    \node at (3,-0.7) {$6cm$};
    
    % Ký hiệu góc vuông tại H
    \draw (0,0) rectangle (0.3,0.3);
    
    % Đánh dấu các điểm bằng chấm tròn nhỏ
    \filldraw[blue] (0,0) circle (2pt);   % Điểm H
    \filldraw[blue] (-3,0) circle (2pt);  % Điểm B
    \filldraw[blue] (6,0) circle (2pt);   % Điểm C
    \filldraw[blue] (0,3) circle (2pt);   % Điểm A
    \draw[<->]   (-3,-0.5)--(6,-0.5);
\end{tikzpicture}
}
\end{minipage}

\vspace{3pt}\hrule\vspace{3pt}\footnotesize\textcolor{gray!80!black}{\textbf{Các mã dùng chung hình (HINH\_001):} \texttt{L06\_C4\_B1\_05}}
\end{tcolorbox}
\vspace{4pt}


\begin{tcolorbox}[
    breakable,
    colback=white,
    colframe=blue!60!black,
    boxrule=0.8pt,
    arc=3pt,
    title={\textbf{Câu \#006} \quad \textbf{MÃ CÂU: \texttt{[L06\_C4\_B1\_06]}} \quad | \quad \textbf{Nhóm: \texttt{HINH\_002}} \quad | \quad \textbf{MapID: \texttt{6H4N1}} \hfill \textcolor{yellow!90!white}{[Hình gốc .pyc]}}
]
\textbf{Dạng toán:} Tính diện tích tam giác vuông \hfill \textit{(Tự luận -- Nhận biết)}
\vspace{4pt}\hrule\vspace{6pt}
\begin{minipage}[t]{0.58\textwidth}
\vspace{0pt}
Tính diện tích tam giác sau:
\end{minipage}\hfill
\begin{minipage}[t]{0.39\textwidth}
\vspace{0pt}
\centering

\resizebox{\linewidth}{!}{%
\begin{tikzpicture}
    % Vẽ cạnh DE (vuông góc)
    \draw[blue, thick] (0,0) -- (0,4) node[midway, left] {$11cm$};
    
    % Vẽ cạnh DF (nằm ngang)
    \draw[blue, thick] (0,0) -- (5,0) node[midway, below] {$18cm$};
    
    % Vẽ cạnh EF (chéo)
    \draw[blue, thick] (0,4) -- (5,0);
    
    % Đánh dấu góc vuông tại D
    \draw[blue] (0,0) rectangle (0.4,0.4);
    
    % Đặt các điểm D, E, F với chấm tròn nhỏ
    \filldraw[blue] (0,0) circle (2pt) node[below left] {$E$};
    \filldraw[blue] (0,4) circle (2pt) node[above left] {$F$};
    \filldraw[blue] (5,0) circle (2pt) node[below right] {$G$};
\end{tikzpicture}
}
\end{minipage}

\vspace{3pt}\hrule\vspace{3pt}\footnotesize\textcolor{gray!80!black}{\textbf{Các mã dùng chung hình (HINH\_002):} \texttt{L06\_C4\_B1\_07}}
\end{tcolorbox}
\vspace{4pt}


\begin{tcolorbox}[
    breakable,
    colback=white,
    colframe=blue!60!black,
    boxrule=0.8pt,
    arc=3pt,
    title={\textbf{Câu \#007} \quad \textbf{MÃ CÂU: \texttt{[L06\_C4\_B1\_08]}} \quad | \quad \textbf{Nhóm: \texttt{HINH\_003}} \quad | \quad \textbf{MapID: \texttt{6H4N1}} \hfill \textcolor{yellow!90!white}{[Hình gốc .pyc]}}
]
\textbf{Dạng toán:} Tính diện tích tứ giác có 2 góc vuông đối diện \hfill \textit{(Tự luận -- Nhận biết)}
\vspace{4pt}\hrule\vspace{6pt}
\begin{minipage}[t]{0.58\textwidth}
\vspace{0pt}
Tính diện tích hình sau biết ${EF}=2cm$, ${GF}=6cm$, ${EH}=5cm$, ${GH}=3cm$
\end{minipage}\hfill
\begin{minipage}[t]{0.39\textwidth}
\vspace{0pt}
\centering

\resizebox{\linewidth}{!}{%
\begin{tikzpicture}
    % Định nghĩa tọa độ của các điểm A và C (đường kính)
    \coordinate (A) at (0,0);
    \coordinate (C) at (6,0);
    \coordinate (B) at (1,2.236);  % Điểm B nằm trên đường tròn
    \coordinate (D) at (4,-2.828); % Điểm D nằm trên đường tròn
    
    \draw[blue, thick] (A) -- (B) -- (C) -- (D) -- cycle;
    
    \def\khvuong[size=#1](#2,#3,#4){%
        \draw ($(#3)!#1!(#2)$) --
        ($($(#3)!#1!(#2)$)+($(#3)!#1!(#4)$)-(#3)$)--
        ($(#3)!#1!(#4)$);
    }
    \khvuong[size=7pt](C,B,A);
        \khvuong[size=7pt](C,D,A);
    \filldraw[black] (A) circle (2pt) node[below left] {$E$};
    \filldraw[black] (B) circle (2pt) node[above left] {$F$};
    \filldraw[black] (C) circle (2pt) node[below right] {$G$};
    \filldraw[black] (D) circle (2pt) node[below right] {$H$};
\end{tikzpicture}
}
\end{minipage}

\vspace{3pt}\hrule\vspace{3pt}\footnotesize\textcolor{gray!80!black}{\textbf{Các mã dùng chung hình (HINH\_003):} \texttt{L06\_C4\_B1\_09}}
\end{tcolorbox}
\vspace{4pt}


\begin{tcolorbox}[
    breakable,
    colback=white,
    colframe=blue!60!black,
    boxrule=0.8pt,
    arc=3pt,
    title={\textbf{Câu \#008} \quad \textbf{MÃ CÂU: \texttt{[L06\_C4\_B1\_10]}} \quad | \quad \textbf{Nhóm: \texttt{HINH\_005}} \quad | \quad \textbf{MapID: \texttt{6H4N1}} \hfill \textcolor{yellow!90!white}{[Hình gốc .pyc]}}
]
\textbf{Dạng toán:} Vẽ tam giác đều. \hfill \textit{(Tự luận -- Nhận biết)}
\vspace{4pt}\hrule\vspace{6pt}
\begin{minipage}[t]{0.58\textwidth}
\vspace{0pt}
Vẽ tam giác đều có cạnh bằng ${4}cm$.
\end{minipage}\hfill
\begin{minipage}[t]{0.39\textwidth}
\vspace{0pt}
\centering

\resizebox{\linewidth}{!}{%
\begin{tikzpicture}
    % Đặt các tọa độ cho các điểm A, B, C theo tỷ lệ nhỏ hơn
    \coordinate (B) at (0,0);   % Điểm B tại gốc tọa độ
    \coordinate (C) at (3,0);   % Điểm C cách B 3 đơn vị (thu nhỏ lại)
    \coordinate (A) at ($(B)!0.5!(C) + (0,2.6)$); % Điểm A tạo tam giác đều với BC
    
    % Vẽ các cạnh của tam giác ABC
    \draw[blue, thick] (A) -- (B) -- (C) -- cycle;
    
    % Vẽ 2 đường tròn đi qua A và B, và đi qua A và C
    \draw[green, thick] (B) circle (3cm);
    \draw[green, thick] (C) circle (3cm);
    
    % Đánh dấu các điểm A, B, C với chấm tròn nhỏ và tên
    \filldraw[red] (A) circle (2pt) node[above] {\textbf{A}};
    \filldraw[red] (B) circle (2pt) node[below left] {\textbf{B}};
    \filldraw[red] (C) circle (2pt) node[below right] {\textbf{C}};
    
    % Ghi nhãn độ dài cạnh BC
    \node at ($(B)!0.5!(C)$) [below] {$4cm$};
    
    \end{tikzpicture}
}
\end{minipage}

\end{tcolorbox}
\vspace{4pt}


\begin{tcolorbox}[
    breakable,
    colback=white,
    colframe=blue!60!black,
    boxrule=0.8pt,
    arc=3pt,
    title={\textbf{Câu \#009} \quad \textbf{MÃ CÂU: \texttt{[L06\_C4\_B1\_13]}} \quad | \quad \textbf{Nhóm: \texttt{HINH\_006}} \quad | \quad \textbf{MapID: \texttt{6H4N1}} \hfill \textcolor{yellow!90!white}{[Hình gốc .pyc]}}
]
\textbf{Dạng toán:} Nhận diện hình vuông \hfill \textit{(Tự luận -- Nhận biết)}
\vspace{4pt}\hrule\vspace{6pt}
\begin{minipage}[t]{0.58\textwidth}
\vspace{0pt}
Trong các hình sau hình nào là hình vuông?
\end{minipage}\hfill
\begin{minipage}[t]{0.39\textwidth}
\vspace{0pt}
\centering

\resizebox{\linewidth}{!}{%
\begin{tikzpicture}
    % Define points
    \coordinate (M) at (0, 3);
    \coordinate (N) at (5, 3);
    \coordinate (P) at (5, 0);
    \coordinate (Q) at (0, 0);
    
    % Draw the rectangle
    
    \draw[thick, blue] (M) -- (N) -- (P) -- (Q) -- cycle;
    
    % Right angle symbols
    \draw[thick] (M) rectangle +(0.3,-0.3);
    \draw[thick] (N) rectangle +(-0.3,-0.3);
    \draw[thick] (P) rectangle +(-0.3,0.3);
    \draw[thick] (Q) rectangle +(0.3,0.3);
    
    % Labels for points
    \fill[red] (M) circle (2pt);
    \fill[red] (N) circle (2pt);
    \fill[red] (P) circle (2pt);
    \fill[red] (Q) circle (2pt);
    
    \node[above left] at (M) {M};
    \node[above right] at (N) {N};
    \node[below right] at (P) {P};
    \node[below left] at (Q) {Q};
    
        \draw
    ($(M)!0.5!(N)$)node[rotate=20]{$/$} ($(P)!0.5!(Q)$)node[rotate=20]{$/$};
    \draw   ($(P)!0.5!(N)$)node[rotate=60]{$//$} ($(M)!0.5!(Q)$)node[rotate=60]{$//$};
    % Draw the square

    
    % Right angle symbols for square
        % Define points
    \coordinate (A) at (8, 0);
    \coordinate (B) at (8, 3);
    \coordinate (C) at (11, 3);
    \coordinate (D) at (11, 0);
    \draw[thick] (A) rectangle +(0.3,0.3);
    \draw[thick] (B) rectangle +(0.3,-0.3);
    \draw[thick] (C) rectangle +(-0.3,-0.3);
    \draw[thick] (D) rectangle +(-0.3,0.3);
    
    % Labels for points of square
    \fill[red] (A) circle (2pt);
    \fill[red] (B) circle (2pt);
    \fill[red] (C) circle (2pt);
    \fill[red] (D) circle (2pt);
    
    \node[below left] at (A) {A};
    \node[above left] at (B) {B};
    \node[above right] at (C) {C};
    \node[below right] at (D) {D};
    \draw[thick, blue] (A) -- (B) -- (C) -- (D) -- cycle;
        \draw
    ($(A)!0.5!(B)$)node[rotate=60]{$/$} ($(B)!0.5!(C)$)node[rotate=20]{$/$};
    \draw   ($(C)!0.5!(D)$)node[rotate=60]{$/$} ($(A)!0.5!(D)$)node[rotate=60]{$/$};
        % Right angle symbols for square
    % Define points
    \coordinate (E) at (15, 0);
    \coordinate (F) at (16, 3);
    \coordinate (G) at (14, 3);
    \coordinate (H) at (15, 6);
    \coordinate (O) at (15, 3);
    % Labels for points of square
    \fill[red] (E) circle (2pt);
    \fill[red] (F) circle (2pt);
    \fill[red] (G) circle (2pt);
    \fill[red] (H) circle (2pt);
        \fill[red] (O) circle (2pt);
    \node[below] at (E) {E};
    \node[above right] at (F) {F};
    \node[above left] at (G) {G};
    \node[above] at (H) {H};
    \draw[thick, blue] (E) -- (F) -- (H) -- (G) -- cycle;
        \draw[thick, blue] (E) -- (H);
    \draw[thick, blue] (F) -- (G);
        \node[above right] at (O) {O};
    \draw   ($(E)!0.5!(F)$)node[rotate=60]{$/$} ($(F)!0.5!(H)$)node[rotate=80]{$/$};
    \draw   ($(H)!0.5!(G)$)node[rotate=60]{$/$} ($(G)!0.5!(E)$)node[rotate=80]{$/$};
    
\end{tikzpicture}
}
\end{minipage}

\end{tcolorbox}
\vspace{4pt}


\begin{tcolorbox}[
    breakable,
    colback=white,
    colframe=blue!60!black,
    boxrule=0.8pt,
    arc=3pt,
    title={\textbf{Câu \#010} \quad \textbf{MÃ CÂU: \texttt{[L06\_C4\_B1\_15]}} \quad | \quad \textbf{Nhóm: \texttt{HINH\_007}} \quad | \quad \textbf{MapID: \texttt{6H4C1}} \hfill \textcolor{yellow!90!white}{[Hình gốc .pyc]}}
]
\textbf{Dạng toán:} Vẽ hình theo yêu cầu \hfill \textit{(Tự luận -- Vận dụng cao)}
\vspace{4pt}\hrule\vspace{6pt}
\begin{minipage}[t]{0.58\textwidth}
\vspace{0pt}
Vẽ hình vuông ${ABCD}$ có cạnh bằng ${8}cm$, vẽ về phía ngoài hình vuông các tam giác đều ${ABE}$, ${BCF}$. Vẽ bằng thước và compa.
\end{minipage}\hfill
\begin{minipage}[t]{0.39\textwidth}
\vspace{0pt}
\centering

\resizebox{\linewidth}{!}{%
\begin{tikzpicture}
    % Vẽ hình vuông ABCD
    \coordinate (A) at (0, 0);
    \coordinate (B) at (3, 0);
    \coordinate (C) at (3, -3);
    \coordinate (D) at (0, -3);
    
    % Vẽ các cạnh của hình vuông
    \draw[thick] (A) -- (B) -- (C) -- (D) -- cycle;
    
    % Vẽ nhãn các điểm của hình vuông
    \node[above left] at (A) {A};
    \node[above right] at (B) {B};
    \node[below right] at (C) {C};
    \node[below left] at (D) {D};
    
    % Vẽ tam giác đều ABE ra phía trên
    \coordinate (E) at (1.5, 2.598); % Đỉnh của tam giác đều ABE
    \draw[thick] (A) -- (E) -- (B) -- cycle;
    \node[above] at (E) {E};
    
    % Vẽ tam giác đều BCF ra phía bên phải
    \coordinate (F) at (6, -1.5); % Đỉnh của tam giác đều BCF
    \draw[thick] (B) -- (F) -- (C) -- cycle;
    \node[right] at (F) {F};
     \fill (A) circle (2pt);
    \fill (B) circle (2pt);
    \fill (C) circle (2pt);
    \fill (D) circle (2pt);
        \fill (E) circle (2pt);
    \fill (F) circle (2pt);
\end{tikzpicture}
}
\end{minipage}

\end{tcolorbox}
\vspace{4pt}


\begin{tcolorbox}[
    breakable,
    colback=white,
    colframe=blue!60!black,
    boxrule=0.8pt,
    arc=3pt,
    title={\textbf{Câu \#011} \quad \textbf{MÃ CÂU: \texttt{[L06\_C4\_B1\_16]}} \quad | \quad \textbf{Nhóm: \texttt{HINH\_008}} \quad | \quad \textbf{MapID: \texttt{6H4C1}} \hfill \textcolor{yellow!90!white}{[Hình gốc .pyc]}}
]
\textbf{Dạng toán:} Vẽ hình theo yêu cầu 2 \hfill \textit{(Tự luận -- Vận dụng cao)}
\vspace{4pt}\hrule\vspace{6pt}
\begin{minipage}[t]{0.58\textwidth}
\vspace{0pt}
Vẽ hình vuông ${MNPQ}$ có cạnh bằng ${2}cm$, vẽ về phía ngoài hình vuông ${MNPQ}$ các hình vuông đều ${MQAB}$, ${NPOC}$. Vẽ bằng thước và compa.
\end{minipage}\hfill
\begin{minipage}[t]{0.39\textwidth}
\vspace{0pt}
\centering

\resizebox{\linewidth}{!}{%
\begin{tikzpicture}
    % Vẽ các điểm của hình
    \coordinate (A) at (0, 0);  % Điểm A
    \coordinate (B) at (0, 3);  % Điểm B
    \coordinate (Q) at (3, 0);  % Điểm Q
    \coordinate (M) at (3, 3);  % Điểm M
    \coordinate (P) at (6, 0);  % Điểm P
    \coordinate (N) at (6, 3);  % Điểm N
    \coordinate (O) at (9, 0);  % Điểm O
    \coordinate (C) at (9, 3);  % Điểm C
    
    % Vẽ hình chữ nhật (hình vuông ghép đôi)
    \draw[thick] (A) -- (B) -- (M) -- (N) -- (C) -- (O) -- (P) -- (Q) -- cycle;
    \draw[thick] (M) -- (Q);
        \draw[thick] (N) -- (P);
    % Đánh dấu các điểm bằng chấm tròn nhỏ
    \fill (A) circle (2pt);
    \fill (B) circle (2pt);
    \fill (Q) circle (2pt);
    \fill (M) circle (2pt);
    \fill (P) circle (2pt);
    \fill (N) circle (2pt);
    \fill (O) circle (2pt);
    \fill (C) circle (2pt);
    
    % Đặt tên các điểm
    \node[below left] at (A) {A};
    \node[above left] at (B) {B};
    \node[below] at (Q) {Q};
    \node[above] at (M) {M};
    \node[below] at (P) {P};
    \node[above] at (N) {N};
    \node[below right] at (O) {O};
    \node[above right] at (C) {C};
    
    % Vẽ ký hiệu góc vuông
    \draw (M) +(0.3,0) -- +(0.3,-0.3) -- +(0,-0.3);
    \draw (N) +(-0.3,0) -- +(-0.3,-0.3) -- +(0,-0.3);
    \draw (Q) +(0,0.3) -- +(0.3,0.3) -- +(0.3,0);
    \draw (P) +(0,0.3) -- +(-0.3,0.3) -- +(-0.3,0);
\end{tikzpicture}
}
\end{minipage}

\end{tcolorbox}
\vspace{4pt}


\begin{tcolorbox}[
    breakable,
    colback=white,
    colframe=blue!60!black,
    boxrule=0.8pt,
    arc=3pt,
    title={\textbf{Câu \#012} \quad \textbf{MÃ CÂU: \texttt{[L06\_C4\_B1\_17]}} \quad | \quad \textbf{Nhóm: \texttt{HINH\_009}} \quad | \quad \textbf{MapID: \texttt{6H4C1}} \hfill \textcolor{yellow!90!white}{[Hình gốc .pyc]}}
]
\textbf{Dạng toán:} Vẽ hình theo yêu cầu 3 \hfill \textit{(Tự luận -- Vận dụng cao)}
\vspace{4pt}\hrule\vspace{6pt}
\begin{minipage}[t]{0.58\textwidth}
\vspace{0pt}
Vẽ tam giác đều ${ABC}$ có cạnh bằng ${4}cm$, vẽ về phía ngoài tam giác các hình vuông đều ${ABDE}$ và ${BCMN}$. Vẽ bằng thước và compa.
\end{minipage}\hfill
\begin{minipage}[t]{0.39\textwidth}
\vspace{0pt}
\centering

\resizebox{\linewidth}{!}{%
\begin{tikzpicture}
    
    % Định nghĩa các điểm của tam giác đều ABC
    \coordinate (B) at (0, 0);    % Điểm B
    \coordinate (C) at (3, 0);    % Điểm C
    \coordinate (A) at (1.5, 2.598);  % Điểm A (đỉnh tam giác đều)
    
    % Vẽ tam giác đều ABC
    \draw[thick] (A) -- (B) -- (C) -- cycle;
    
    % Đánh dấu các điểm của tam giác đều ABC
    \fill (A) circle (2pt);
    \fill (B) circle (2pt);
    \fill (C) circle (2pt);
    
    % Nhãn các điểm
    \node[above] at (A) {A};
    \node[below left] at (B) {B};
    \node[below right] at (C) {C};
    
    % Vẽ hình vuông ABDE ra phía ngoài cạnh AB
    \coordinate (D) at (-2.598, 1.299);  % Điểm D
    \coordinate (E) at (-1.299, 3.897);  % Điểm E
    \draw[thick] (A) -- (E) -- (D) -- (B) -- cycle;
    
    % Đánh dấu các điểm của hình vuông ABDE
    \fill (D) circle (2pt);
    \fill (E) circle (2pt);
    
    % Nhãn các điểm
    \node[left] at (D) {D};
    \node[above] at (E) {E};
    
    % Vẽ hình vuông BCMN ra phía ngoài cạnh BC
    \coordinate (M) at (3, -3);   % Điểm M
    \coordinate (N) at (0, -3);   % Điểm N
    \draw[thick] (B) -- (N) -- (M) -- (C) -- cycle;
    
    % Đánh dấu các điểm của hình vuông BCMN
    \fill (M) circle (2pt);
    \fill (N) circle (2pt);
    
    % Nhãn các điểm
    \node[below right] at (M) {M};
    \node[below left] at (N) {N};
\end{tikzpicture}
}
\end{minipage}

\end{tcolorbox}
\vspace{4pt}


\begin{tcolorbox}[
    breakable,
    colback=white,
    colframe=blue!60!black,
    boxrule=0.8pt,
    arc=3pt,
    title={\textbf{Câu \#013} \quad \textbf{MÃ CÂU: \texttt{[L06\_C4\_B1\_24]}} \quad | \quad \textbf{Nhóm: \texttt{HINH\_010}} \quad | \quad \textbf{MapID: \texttt{6H4N1}} \hfill \textcolor{yellow!90!white}{[Hình gốc .pyc]}}
]
\textbf{Dạng toán:} Đếm số hình vuông \hfill \textit{(Tự luận -- Nhận biết)}
\vspace{4pt}\hrule\vspace{6pt}
\begin{minipage}[t]{0.58\textwidth}
\vspace{0pt}
Đếm số hình vuông trong hình sau
\end{minipage}\hfill
\begin{minipage}[t]{0.39\textwidth}
\vspace{0pt}
\centering

\resizebox{\linewidth}{!}{%
\begin{tikzpicture}
    % Vẽ lưới 2x2
    \draw[step=1cm] (0, 0) grid (2, 2);
    
    % Vẽ viền ngoài hình vuông 2x2
    \draw[thick] (0, 0) rectangle (2, 2);
    
    % Điền các vòng tròn đặc tại các đỉnh
    \fill (0, 0) circle (2pt); % Góc dưới bên trái
    \fill (0, 2) circle (2pt); % Góc trên bên trái
    \fill (2, 0) circle (2pt); % Góc dưới bên phải
    \fill (2, 2) circle (2pt); % Góc trên bên phải
    
    % Điền các vòng tròn đặc tại các giao điểm bên trong
    \fill (0, 1) circle (2pt);
    \fill (1, 0) circle (2pt);
    \fill (1, 1) circle (2pt);
    \fill (1, 2) circle (2pt);
    \fill (2, 1) circle (2pt);
\end{tikzpicture}
}
\end{minipage}

\end{tcolorbox}
\vspace{4pt}


\begin{tcolorbox}[
    breakable,
    colback=white,
    colframe=blue!60!black,
    boxrule=0.8pt,
    arc=3pt,
    title={\textbf{Câu \#014} \quad \textbf{MÃ CÂU: \texttt{[L06\_C4\_B1\_25]}} \quad | \quad \textbf{Nhóm: \texttt{HINH\_011}} \quad | \quad \textbf{MapID: \texttt{6H4V1}} \hfill \textcolor{yellow!90!white}{[Hình gốc .pyc]}}
]
\textbf{Dạng toán:} Đếm số hình vuông \hfill \textit{(Tự luận -- Vận dụng)}
\vspace{4pt}\hrule\vspace{6pt}
\begin{minipage}[t]{0.58\textwidth}
\vspace{0pt}
Đếm số hình vuông trong hình sau
\end{minipage}\hfill
\begin{minipage}[t]{0.39\textwidth}
\vspace{0pt}
\centering

\resizebox{\linewidth}{!}{%
\begin{tikzpicture}
    % Vẽ lưới 3x3
    \draw[step=1cm] (0, 0) grid (3, 3);
    
    % Vẽ viền ngoài của hình vuông 3x3
    \draw[thick] (0, 0) rectangle (3, 3);
    
    % Vẽ các vòng tròn đặc tại các đỉnh và giao điểm
    \fill (0, 0) circle (2pt);
    \fill (0, 1) circle (2pt);
    \fill (0, 2) circle (2pt);
    \fill (0, 3) circle (2pt);
    
    \fill (1, 0) circle (2pt);
    \fill (1, 1) circle (2pt);
    \fill (1, 2) circle (2pt);
    \fill (1, 3) circle (2pt);
    
    \fill (2, 0) circle (2pt);
    \fill (2, 1) circle (2pt);
    \fill (2, 2) circle (2pt);
    \fill (2, 3) circle (2pt);
    
    \fill (3, 0) circle (2pt);
    \fill (3, 1) circle (2pt);
    \fill (3, 2) circle (2pt);
    \fill (3, 3) circle (2pt);

\end{tikzpicture}
}
\end{minipage}

\end{tcolorbox}
\vspace{4pt}


\begin{tcolorbox}[
    breakable,
    colback=white,
    colframe=blue!60!black,
    boxrule=0.8pt,
    arc=3pt,
    title={\textbf{Câu \#015} \quad \textbf{MÃ CÂU: \texttt{[L06\_C4\_B1\_26]}} \quad | \quad \textbf{Nhóm: \texttt{HINH\_012}} \quad | \quad \textbf{MapID: \texttt{6H4C1}} \hfill \textcolor{yellow!90!white}{[Hình gốc .pyc]}}
]
\textbf{Dạng toán:} Đếm số hình vuông \hfill \textit{(Tự luận -- Vận dụng cao)}
\vspace{4pt}\hrule\vspace{6pt}
\begin{minipage}[t]{0.58\textwidth}
\vspace{0pt}
Đếm số hình vuông trong hình sau
\end{minipage}\hfill
\begin{minipage}[t]{0.39\textwidth}
\vspace{0pt}
\centering

\resizebox{\linewidth}{!}{%
\begin{tikzpicture}
    % Vẽ lưới
    \draw[step=1cm] (0, 0) grid (4, 4);
    
    % Vẽ hình vuông bên ngoài
    \draw[thick] (0, 0) rectangle (4, 4);
    
    % Vẽ các hình tròn đặc tại các điểm lưới
    \fill (0, 0) circle (2pt);
    \fill (0, 1) circle (2pt);
    \fill (0, 2) circle (2pt);
    \fill (0, 3) circle (2pt);
    \fill (0, 4) circle (2pt);

    \fill (1, 0) circle (2pt);
    \fill (1, 1) circle (2pt);
    \fill (1, 2) circle (2pt);
    \fill (1, 3) circle (2pt);
    \fill (1, 4) circle (2pt);

    \fill (2, 0) circle (2pt);
    \fill (2, 1) circle (2pt);
    \fill (2, 2) circle (2pt);
    \fill (2, 3) circle (2pt);
    \fill (2, 4) circle (2pt);

    \fill (3, 0) circle (2pt);
    \fill (3, 1) circle (2pt);
    \fill (3, 2) circle (2pt);
    \fill (3, 3) circle (2pt);
    \fill (3, 4) circle (2pt);

    \fill (4, 0) circle (2pt);
    \fill (4, 1) circle (2pt);
    \fill (4, 2) circle (2pt);
    \fill (4, 3) circle (2pt);
    \fill (4, 4) circle (2pt);

    % Vẽ các hình tròn rỗng tại các góc dưới cùng
    \draw[thick] (0, 0) circle (2pt);
    \draw[thick] (4, 0) circle (2pt);
    
\end{tikzpicture}
}
\end{minipage}

\end{tcolorbox}
\vspace{4pt}


\begin{tcolorbox}[
    breakable,
    colback=white,
    colframe=blue!60!black,
    boxrule=0.8pt,
    arc=3pt,
    title={\textbf{Câu \#016} \quad \textbf{MÃ CÂU: \texttt{[L06\_C4\_B2\_02]}} \quad | \quad \textbf{Nhóm: \texttt{HINH\_013}} \quad | \quad \textbf{MapID: \texttt{6H4N2}} \hfill \textcolor{yellow!90!white}{[Hình gốc .pyc]}}
]
\textbf{Dạng toán:} Nhận diện các hình \hfill \textit{(Trắc nghiệm -- Nhận biết)}
\vspace{4pt}\hrule\vspace{6pt}
\begin{minipage}[t]{0.58\textwidth}
\vspace{0pt}
Trong các hình sau hình nào là hình thang cân.

\vspace{4pt}\noindent
\textbf{A.} c) d) \quad \textbf{B.} a) b) \quad \textbf{\textcolor{blue!75!black}{C. (Đ)}} a) c) \quad \textbf{D.} b) d)
\end{minipage}\hfill
\begin{minipage}[t]{0.39\textwidth}
\vspace{0pt}
\centering

\resizebox{\linewidth}{!}{%
\begin{tikzpicture}

% Hình a) - Trapezoid
\draw[thick] (0,0) -- (2,0) -- (2.5,1.5) -- (-0.5,1.5) -- cycle;
\node at (1, -0.5) {\textbf{a)}};

% Hình b) - Parallelogram
\draw[thick] (3,0) -- (5,0) -- (6,1.5) -- (4,1.5) -- cycle;
\node at (4.5, -0.5) {\textbf{b)}};

% Hình c) - Rectangle
\draw[thick] (7,0) -- (9,0) -- (9,1.5) -- (7,1.5) -- cycle;
\node at (8, -0.5) {\textbf{c)}};

% Hình d) - Rhombus
\draw[thick] (10,0.75) -- (11.5,0) -- (13,0.75) -- (11.5,1.5) -- cycle;
\node at (11.5, -0.5) {\textbf{d)}};

\end{tikzpicture}
}
\end{minipage}

\end{tcolorbox}
\vspace{4pt}


\begin{tcolorbox}[
    breakable,
    colback=white,
    colframe=blue!60!black,
    boxrule=0.8pt,
    arc=3pt,
    title={\textbf{Câu \#017} \quad \textbf{MÃ CÂU: \texttt{[L06\_C4\_B2\_01]}} \quad | \quad \textbf{Nhóm: \texttt{HINH\_014}} \quad | \quad \textbf{MapID: \texttt{6H4N2}} \hfill \textcolor{yellow!90!white}{[Hình gốc .pyc]}}
]
\textbf{Dạng toán:} Nhận diện các hình \hfill \textit{(Trả lời ngắn -- Nhận biết)}
\vspace{4pt}\hrule\vspace{6pt}
\begin{minipage}[t]{0.58\textwidth}
\vspace{0pt}
Trong các hình sau hình nào là hình chữ nhật, hình nào là hình bình hành.


\textbf{Đáp án điền số:} 	exttt{}
\end{minipage}\hfill
\begin{minipage}[t]{0.39\textwidth}
\vspace{0pt}
\centering

\resizebox{\linewidth}{!}{%
\begin{tikzpicture}

    % Trapezoid - a)
    \draw[thick] (0,0) -- (2,0) -- (2.5,1.5) -- (-0.5,1.5) -- cycle;
    \node at (1, -0.5) {a)};

    % Parallelogram - b)
    \draw[thick] (3,0) -- (5,0) -- (6,1.5) -- (4,1.5) -- cycle;
    \node at (4.5, -0.5) {b)};

    % Rectangle - c)
    \draw[thick] (7,0) -- (9,0) -- (9,1.5) -- (7,1.5) -- cycle;
    \node at (8, -0.5) {c)};

    % Rhombus - d)
    \draw[thick] (10,0.75) -- (11.5,0) -- (13,0.75) -- (11.5,1.5) -- cycle;
    \node at (11.5, -0.5) {d)};

    \end{tikzpicture}
}
\end{minipage}

\end{tcolorbox}
\vspace{4pt}


\begin{tcolorbox}[
    breakable,
    colback=white,
    colframe=blue!60!black,
    boxrule=0.8pt,
    arc=3pt,
    title={\textbf{Câu \#018} \quad \textbf{MÃ CÂU: \texttt{[L06\_C4\_B2\_08]}} \quad | \quad \textbf{Nhóm: \texttt{HINH\_015}} \quad | \quad \textbf{MapID: \texttt{6H4N2}} \hfill \textcolor{yellow!90!white}{[Hình gốc .pyc]}}
]
\textbf{Dạng toán:} Mô tả được các yếu tố cơ bản của chữ nhật \hfill \textit{(Trả lời ngắn -- Nhận biết)}
\vspace{4pt}\hrule\vspace{6pt}
\begin{minipage}[t]{0.58\textwidth}
\vspace{0pt}
Nêu tên đỉnh, cạnh, đường chéo, hai cạnh đối của hình chữ nhật  ${GHFE}$


\textbf{Đáp án điền số:} 	exttt{}
\end{minipage}\hfill
\begin{minipage}[t]{0.39\textwidth}
\vspace{0pt}
\centering

\resizebox{\linewidth}{!}{%
\begin{tikzpicture}
    % Các đỉnh của hình chữ nhật
    % Các đỉnh của hình chữ nhật
    \coordinate [label=below left:{$G$}] (E) at (0, 0);
    \coordinate [label=above left:{$H$}] (F) at (0, 3);
    \coordinate [label=above right:{$F$}] (H) at (5, 3);
    \coordinate [label=below right:{$E$}] (G) at (5, 0);
    % Vẽ hình chữ nhật
    \draw[thick] (E) -- (F) -- (G) -- (H) -- cycle;
    
    % Vẽ các đường chéo
    \draw[thick] (E) -- (G);
    \draw[thick] (F) -- (H);

    % Các điểm
    \fill (E) circle (2pt);
    \fill (F) circle (2pt);
    \fill (H) circle (2pt);
    \fill (G) circle (2pt);
\end{tikzpicture}
}
\end{minipage}

\vspace{3pt}\hrule\vspace{3pt}\footnotesize\textcolor{gray!80!black}{\textbf{Các mã dùng chung hình (HINH\_015):} \texttt{L06\_C4\_B2\_11}}
\end{tcolorbox}
\vspace{4pt}


\begin{tcolorbox}[
    breakable,
    colback=white,
    colframe=blue!60!black,
    boxrule=0.8pt,
    arc=3pt,
    title={\textbf{Câu \#019} \quad \textbf{MÃ CÂU: \texttt{[L06\_C4\_B2\_09]}} \quad | \quad \textbf{Nhóm: \texttt{HINH\_016}} \quad | \quad \textbf{MapID: \texttt{6H4N2}} \hfill \textcolor{yellow!90!white}{[Hình gốc .pyc]}}
]
\textbf{Dạng toán:} Mô tả được yếu tố cơ bản của hình thoi \hfill \textit{(Trả lời ngắn -- Nhận biết)}
\vspace{4pt}\hrule\vspace{6pt}
\begin{minipage}[t]{0.58\textwidth}
\vspace{0pt}
Nêu tên đỉnh, cạnh, đường chéo, hai cạnh đối của hình thoi  ${PQRT}$


\textbf{Đáp án điền số:} 	exttt{}
\end{minipage}\hfill
\begin{minipage}[t]{0.39\textwidth}
\vspace{0pt}
\centering

\resizebox{\linewidth}{!}{%
\begin{tikzpicture}
    % Các đỉnh của hình thoi
    \coordinate [label=left:{$P$}] (G) at (0, 0);
    \coordinate [label=above:{$Q$}] (H) at (2, 3);
    \coordinate [label=right:{$R$}] (I) at (4, 0);
    \coordinate [label=below:{$T$}] (K) at (2, -3);

    % Vẽ các cạnh của hình thoi
    \draw[thick] (G) -- (H) -- (I) -- (K) -- cycle;
    
    % Vẽ các đường chéo
    \draw[thick] (H) -- (K);
    \draw[thick] (G) -- (I);

    % Vẽ các điểm
    \fill (G) circle (2pt);
    \fill (H) circle (2pt);
    \fill (I) circle (2pt);
    \fill (K) circle (2pt);
\end{tikzpicture}
}
\end{minipage}

\end{tcolorbox}
\vspace{4pt}


\begin{tcolorbox}[
    breakable,
    colback=white,
    colframe=blue!60!black,
    boxrule=0.8pt,
    arc=3pt,
    title={\textbf{Câu \#020} \quad \textbf{MÃ CÂU: \texttt{[L06\_C4\_B2\_10]}} \quad | \quad \textbf{Nhóm: \texttt{HINH\_017}} \quad | \quad \textbf{MapID: \texttt{6H4N2}} \hfill \textcolor{yellow!90!white}{[Hình gốc .pyc]}}
]
\textbf{Dạng toán:} Mô tả được yếu tố cơ bản của hình bình hành \hfill \textit{(Trả lời ngắn -- Nhận biết)}
\vspace{4pt}\hrule\vspace{6pt}
\begin{minipage}[t]{0.58\textwidth}
\vspace{0pt}
Nêu tên đỉnh, cạnh, đường chéo, hai cạnh đối của hình bình hành  ${PQRT}$


\textbf{Đáp án điền số:} 	exttt{}
\end{minipage}\hfill
\begin{minipage}[t]{0.39\textwidth}
\vspace{0pt}
\centering

\resizebox{\linewidth}{!}{%
\begin{tikzpicture}

    % Định nghĩa các đỉnh của hình bình hành
    \coordinate [label=above left:{$P$}] (A) at (0, 0);
    \coordinate [label=above right:{$Q$}] (B) at (4, 0);
    \coordinate [label=below right:{$R$}] (C) at (5, -3);
    \coordinate [label=below left:{$T$}] (D) at (1, -3);

    % Vẽ các cạnh của hình bình hành
    \draw[thick] (A) -- (B) -- (C) -- (D) -- cycle;

    % Vẽ các điểm
    \fill (A) circle (2pt);
    \fill (B) circle (2pt);
    \fill (C) circle (2pt);
    \fill (D) circle (2pt);

\end{tikzpicture}
}
\end{minipage}

\end{tcolorbox}
\vspace{4pt}


\begin{tcolorbox}[
    breakable,
    colback=white,
    colframe=blue!60!black,
    boxrule=0.8pt,
    arc=3pt,
    title={\textbf{Câu \#021} \quad \textbf{MÃ CÂU: \texttt{[L06\_C4\_B2\_11]}} \quad | \quad \textbf{Nhóm: \texttt{HINH\_015}} \quad | \quad \textbf{MapID: \texttt{6H4N2}} \hfill \textcolor{yellow!90!white}{[Hình gốc .pyc]}}
]
\textbf{Dạng toán:} Tính được độ dài các cạnh của hình chữ nhật \hfill \textit{(Trả lời ngắn -- Nhận biết)}
\vspace{4pt}\hrule\vspace{6pt}
\begin{minipage}[t]{0.58\textwidth}
\vspace{0pt}
Cho hình chữ nhật  ${GHFE}$ có ${GH}=15cm$, ${GE}=20cm$, ${GF}=25cm$. Tính ${FH}$, ${FE}$, ${EH}$


\textbf{Đáp án điền số:} 	exttt{}
\end{minipage}\hfill
\begin{minipage}[t]{0.39\textwidth}
\vspace{0pt}
\centering

\resizebox{\linewidth}{!}{%
\begin{tikzpicture}
    % Các đỉnh của hình chữ nhật
    \coordinate [label=below left:{$G$}] (E) at (0, 0);
    \coordinate [label=above left:{$H$}] (F) at (0, 3);
    \coordinate [label=above right:{$F$}] (H) at (5, 3);
    \coordinate [label=below right:{$E$}] (G) at (5, 0);

    % Vẽ hình chữ nhật
    \draw[thick] (E) -- (F) -- (G) -- (H) -- cycle;
    
    % Vẽ các đường chéo
    \draw[thick] (E) -- (G);
    \draw[thick] (F) -- (H);

    % Các điểm
    \fill (E) circle (2pt);
    \fill (F) circle (2pt);
    \fill (H) circle (2pt);
    \fill (G) circle (2pt);
\end{tikzpicture}
}
\end{minipage}

\vspace{3pt}\hrule\vspace{3pt}\footnotesize\textcolor{gray!80!black}{\textbf{Các mã dùng chung hình (HINH\_015):} \texttt{L06\_C4\_B2\_08}}
\end{tcolorbox}
\vspace{4pt}


\begin{tcolorbox}[
    breakable,
    colback=white,
    colframe=blue!60!black,
    boxrule=0.8pt,
    arc=3pt,
    title={\textbf{Câu \#022} \quad \textbf{MÃ CÂU: \texttt{[L06\_C4\_B2\_12]}} \quad | \quad \textbf{Nhóm: \texttt{HINH\_018}} \quad | \quad \textbf{MapID: \texttt{6H4N2}} \hfill \textcolor{yellow!90!white}{[Hình gốc .pyc]}}
]
\textbf{Dạng toán:} Tính các cạnh còn lại của hình thoi \hfill \textit{(Trả lời ngắn -- Nhận biết)}
\vspace{4pt}\hrule\vspace{6pt}
\begin{minipage}[t]{0.58\textwidth}
\vspace{0pt}
Cho hình thoi  ${GHFE}$ biết $GH=29cm$. Tính độ dài ${FH}$, ${FE}$, ${GE}$.


\textbf{Đáp án điền số:} 	exttt{}
\end{minipage}\hfill
\begin{minipage}[t]{0.39\textwidth}
\vspace{0pt}
\centering

\resizebox{\linewidth}{!}{%
\begin{tikzpicture}
    % Các đỉnh của hình thoi
    \coordinate [label=left:{$G$}] (G) at (0, 0);
    \coordinate [label=above:{$H$}] (H) at (2, 3);
    \coordinate [label=right:{$F$}] (I) at (4, 0);
    \coordinate [label=below:{$E$}] (K) at (2, -3);

    % Vẽ các cạnh của hình thoi
    \draw[thick] (G) -- (H) -- (I) -- (K) -- cycle;
    
    % Vẽ các đường chéo
    \draw[thick] (H) -- (K);
    \draw[thick] (G) -- (I);

    % Vẽ các điểm
    \fill (G) circle (2pt);
    \fill (H) circle (2pt);
    \fill (I) circle (2pt);
    \fill (K) circle (2pt);
\end{tikzpicture}
}
\end{minipage}

\end{tcolorbox}
\vspace{4pt}


\begin{tcolorbox}[
    breakable,
    colback=white,
    colframe=blue!60!black,
    boxrule=0.8pt,
    arc=3pt,
    title={\textbf{Câu \#023} \quad \textbf{MÃ CÂU: \texttt{[L06\_C4\_B2\_13]}} \quad | \quad \textbf{Nhóm: \texttt{HINH\_019}} \quad | \quad \textbf{MapID: \texttt{6H4N2}} \hfill \textcolor{yellow!90!white}{[Hình gốc .pyc]}}
]
\textbf{Dạng toán:} Tính các yếu tố còn lại của hình thoi \hfill \textit{(Trả lời ngắn -- Nhận biết)}
\vspace{4pt}\hrule\vspace{6pt}
\begin{minipage}[t]{0.58\textwidth}
\vspace{0pt}
Cho hình thoi  ${IJQK}$, gọi ${O}$ là giao điểm của hai đường chéo ${IK}$ và ${JQ}$ biết $IJ=25cm$, $OI=40cm$, $OJ=20cm$. Tính độ dài ${QJ}$, ${QK}$, ${IK}$, ${IQ}$, ${JK}$.


\textbf{Đáp án điền số:} 	exttt{}
\end{minipage}\hfill
\begin{minipage}[t]{0.39\textwidth}
\vspace{0pt}
\centering

\resizebox{\linewidth}{!}{%
\begin{tikzpicture}
    % Các đỉnh của hình thoi
    \coordinate [label=left:{$I$}] (G) at (0, 0);
    \coordinate [label=above:{$J$}] (H) at (2, 3);
    \coordinate [label=right:{$Q$}] (I) at (4, 0);
    \coordinate [label=below:{$K$}] (K) at (2, -3);
 \coordinate [label=below right:{$O$}] (O) at (2, 0);
    % Vẽ các cạnh của hình thoi
    \draw[thick] (G) -- (H) -- (I) -- (K) -- cycle;
    
    % Vẽ các đường chéo
    \draw[thick] (H) -- (K);
    \draw[thick] (G) -- (I);

    % Vẽ các điểm
    \fill (G) circle (2pt);
    \fill (H) circle (2pt);
    \fill (I) circle (2pt);
    \fill (K) circle (2pt);
\end{tikzpicture}
}
\end{minipage}

\end{tcolorbox}
\vspace{4pt}


\begin{tcolorbox}[
    breakable,
    colback=white,
    colframe=blue!60!black,
    boxrule=0.8pt,
    arc=3pt,
    title={\textbf{Câu \#024} \quad \textbf{MÃ CÂU: \texttt{[L06\_C4\_B2\_15]}} \quad | \quad \textbf{Nhóm: \texttt{HINH\_020}} \quad | \quad \textbf{MapID: \texttt{6H4N2}} \hfill \textcolor{yellow!90!white}{[Hình gốc .pyc]}}
]
\textbf{Dạng toán:} Mô tả được yếu tố cơ bản của hình thang cân \hfill \textit{(Trả lời ngắn -- Nhận biết)}
\vspace{4pt}\hrule\vspace{6pt}
\begin{minipage}[t]{0.58\textwidth}
\vspace{0pt}
Nêu tên đỉnh, cạnh, đường chéo, đáy lớn đáy nhỏ của hình thang cân ${IJQK}$


\textbf{Đáp án điền số:} 	exttt{}
\end{minipage}\hfill
\begin{minipage}[t]{0.39\textwidth}
\vspace{0pt}
\centering

\resizebox{\linewidth}{!}{%
\begin{tikzpicture}
    % Tọa độ các đỉnh của hình thang cân ABCD
    \coordinate (A) at (0,0);   % Đỉnh A
    \coordinate (B) at (6,0);   % Đỉnh B
    \coordinate (D) at (2,3);   % Đỉnh D
    \coordinate (C) at (4,3);   % Đỉnh C

    % Vẽ hình thang cân ABCD
    \draw[thick] (A) -- (B) -- (C) -- (D) -- cycle;

    % Ghi nhãn các đỉnh
    \node[below left] at (A) {$I$};
    \node[below right] at (B) {$J$};
    \node[above right] at (C) {$Q$};
    \node[above left] at (D) {$K$};
\end{tikzpicture}
}
\end{minipage}

\end{tcolorbox}
\vspace{4pt}


\begin{tcolorbox}[
    breakable,
    colback=white,
    colframe=blue!60!black,
    boxrule=0.8pt,
    arc=3pt,
    title={\textbf{Câu \#025} \quad \textbf{MÃ CÂU: \texttt{[L06\_C4\_B3\_32]}} \quad | \quad \textbf{Nhóm: \texttt{HINH\_021}} \quad | \quad \textbf{MapID: \texttt{6H4V3}} \hfill \textcolor{yellow!90!white}{[Hình gốc .pyc]}}
]
\textbf{Dạng toán:} Bài toán thực tế về hình bình hành và hình vuông \hfill \textit{(Trắc nghiệm -- Vận dụng)}
\vspace{4pt}\hrule\vspace{6pt}
\begin{minipage}[t]{0.58\textwidth}
\vspace{0pt}
Một mảnh vườn hình bình hành, giữa vườn người ta đào một cái ao hình vuông cạnh ${4}m$ như hình vẽ minh hoạ. Biết chiều cao h dài gấp $3$ cạnh hình vuông. Tính diện tích còn lại của vườn sau khi đào ao?

\vspace{4pt}\noindent
\textbf{A.} ${60}m^{2}$ \quad \textbf{B.} ${228}m^{2}$ \quad \textbf{\textcolor{blue!75!black}{C. (Đ)}} ${212}m^{2}$ \quad \textbf{D.} ${224}m^{2}$
\end{minipage}\hfill
\begin{minipage}[t]{0.39\textwidth}
\vspace{0pt}
\centering

\resizebox{\linewidth}{!}{%
\begin{tikzpicture}
    % Tọa độ các điểm của hình bình hành
    \coordinate (A) at (0,0);     % Góc dưới bên trái
    \coordinate (B) at (10,0);    % Góc dưới bên phải
    \coordinate (C) at (8,4);     % Góc trên bên phải
    \coordinate (D) at (-2,4);    % Góc trên bên trái
    
    \coordinate (N) at (0,4);    % Góc trên bên trái
    % Tọa độ các điểm của hình vuông nhỏ bên trong (cạnh = 3)
    \coordinate (E) at (3,2);     % Điểm trung tâm của hình vuông
    \coordinate (F) at (3,3.5);   % Điểm trên
    \coordinate (G) at (1.5,2);   % Điểm trái
    \coordinate (H) at (3,0.5);   % Điểm dưới
    \coordinate (I) at (4.5,2);   % Điểm phải
    
    % Vẽ hình bình hành
    \draw[thick] (A) -- (B) -- (C) -- (D) -- cycle;
        \draw[thick] (A) -- (N);
    % Vẽ hình vuông nhỏ
    \fill[cyan!30] (F) -- (G) -- (H) -- (I) -- cycle; % Tô màu cho hình vuông
    \draw[thick] (F) -- (G) -- (H) -- (I) -- cycle;   % Vẽ các cạnh của hình vuông
    
    % Vẽ chiều cao của hình bình hành từ cạnh dưới lên cạnh trên

    \node[right] at ($(A)!0.5!(N)$) {$h$}; % Chú thích chiều cao
    
    % Chú thích kích thước các cạnh
    \node[above] at ($(D)!0.5!(C)$) {$19$};   % Chú thích cạnh trên của hình bình hành
    \node[left] at ($(G)!0.5!(F)$) {$4$};     % Chú thích cạnh hình vuông nhỏ

    
    % Vẽ góc vuông tại chân đường cao
    \draw (A) -- ++(0.3,0) -- ++(0,0.3) -- ++(-0.3,0);
\end{tikzpicture}
}
\end{minipage}

\vspace{3pt}\hrule\vspace{3pt}\footnotesize\textcolor{gray!80!black}{\textbf{Các mã dùng chung hình (HINH\_021):} \texttt{L06\_C4\_B3\_31}}
\end{tcolorbox}
\vspace{4pt}


\begin{tcolorbox}[
    breakable,
    colback=white,
    colframe=blue!60!black,
    boxrule=0.8pt,
    arc=3pt,
    title={\textbf{Câu \#026} \quad \textbf{MÃ CÂU: \texttt{[L06\_C4\_B3\_48]}} \quad | \quad \textbf{Nhóm: \texttt{HINH\_022}} \quad | \quad \textbf{MapID: \texttt{6H4V3}} \hfill \textcolor{yellow!90!white}{[Hình gốc .pyc]}}
]
\textbf{Dạng toán:} TF-M3. Tính diện tích và tiền trồng hoa và cỏ của mảnh đất gồm HCN và hình bình hành \hfill \textit{(Đúng-Sai -- Vận dụng)}
\vspace{4pt}\hrule\vspace{6pt}
\begin{minipage}[t]{0.58\textwidth}
\vspace{0pt}
Bác Lan có một mảnh đất hình chữ nhật có diện tích là $540$ m$^{2}$, chiều rộng $15$ m. Ở giữa đất bác làm lối đi như hình vẽ, phần đất còn lại để trồng cây. Xét tính đúng-sai của các khẳng định sau.

\vspace{4pt}\noindent
\textbf{\textcolor{blue!75!black}{a) (Đ)}} Chiều dài mảnh đất là $36$ (m)\\[2pt]
\textbf{b)} Chu vi mảnh đất là $51$ (m)\\[2pt]
\textbf{c)} Diện tích của phần đất trồng cây là $492$ m$^{2}$\\[2pt]
\textbf{\textcolor{blue!75!black}{d) (Đ)}} Số tiền bác Lan phải trả khi xây tường rào là $5940000$ (đồng)
\end{minipage}\hfill
\begin{minipage}[t]{0.39\textwidth}
\vspace{0pt}
\centering

\resizebox{\linewidth}{!}{%
\begin{tikzpicture}
    \def\dai{5}
    \def\rong{4}
    \def\loidi{0.5}
    \draw[latex-latex] (\dai+0.5,0)--++(90:\rong); 
    \draw (0,0)--++(90:\rong)--++(0:\dai) node[midway,above]{$3$ m}--++(-90:\rong)node[fill=white,inner sep=4pt,rectangle,midway,right]{\color{black} $15$ m}--++(180:\dai) node[midway,below]{$3$ m}
    (\dai/2-\loidi/2,0) rectangle (\dai/2+\loidi/2,\rong)
    ;
    \end{tikzpicture}
}
\end{minipage}

\end{tcolorbox}
\vspace{4pt}


\begin{tcolorbox}[
    breakable,
    colback=white,
    colframe=blue!60!black,
    boxrule=0.8pt,
    arc=3pt,
    title={\textbf{Câu \#027} \quad \textbf{MÃ CÂU: \texttt{[L06\_C4\_B3\_31]}} \quad | \quad \textbf{Nhóm: \texttt{HINH\_021}} \quad | \quad \textbf{MapID: \texttt{6H4V3}} \hfill \textcolor{yellow!90!white}{[Hình gốc .pyc]}}
]
\textbf{Dạng toán:} Bài toán thực tế về hình bình hành và hình vuông \hfill \textit{(Tự luận -- Vận dụng)}
\vspace{4pt}\hrule\vspace{6pt}
\begin{minipage}[t]{0.58\textwidth}
\vspace{0pt}
Một mảnh vườn hình bình hành, giữa vườn người ta đào một cái ao hình vuông cạnh ${9}m$ như hình vẽ minh hoạ. Biết chiều cao h dài gấp $4$ cạnh hình vuông. Tính diện tích còn lại của vườn sau khi đào ao?
\end{minipage}\hfill
\begin{minipage}[t]{0.39\textwidth}
\vspace{0pt}
\centering

\resizebox{\linewidth}{!}{%
\begin{tikzpicture}
    % Tọa độ các điểm của hình bình hành
    \coordinate (A) at (0,0);     % Góc dưới bên trái
    \coordinate (B) at (10,0);    % Góc dưới bên phải
    \coordinate (C) at (8,4);     % Góc trên bên phải
    \coordinate (D) at (-2,4);    % Góc trên bên trái
    
    \coordinate (N) at (0,4);    % Góc trên bên trái
    % Tọa độ các điểm của hình vuông nhỏ bên trong (cạnh = 3)
    \coordinate (E) at (3,2);     % Điểm trung tâm của hình vuông
    \coordinate (F) at (3,3.5);   % Điểm trên
    \coordinate (G) at (1.5,2);   % Điểm trái
    \coordinate (H) at (3,0.5);   % Điểm dưới
    \coordinate (I) at (4.5,2);   % Điểm phải
    
    % Vẽ hình bình hành
    \draw[thick] (A) -- (B) -- (C) -- (D) -- cycle;
        \draw[thick] (A) -- (N);
    % Vẽ hình vuông nhỏ
    \fill[cyan!30] (F) -- (G) -- (H) -- (I) -- cycle; % Tô màu cho hình vuông
    \draw[thick] (F) -- (G) -- (H) -- (I) -- cycle;   % Vẽ các cạnh của hình vuông
    
    % Vẽ chiều cao của hình bình hành từ cạnh dưới lên cạnh trên

    \node[right] at ($(A)!0.5!(N)$) {$h$}; % Chú thích chiều cao
    
    % Chú thích kích thước các cạnh
    \node[above] at ($(D)!0.5!(C)$) {$19$};   % Chú thích cạnh trên của hình bình hành
    \node[left] at ($(G)!0.5!(F)$) {$9$};     % Chú thích cạnh hình vuông nhỏ

    
    % Vẽ góc vuông tại chân đường cao
    \draw (A) -- ++(0.3,0) -- ++(0,0.3) -- ++(-0.3,0);
\end{tikzpicture}
}
\end{minipage}

\vspace{3pt}\hrule\vspace{3pt}\footnotesize\textcolor{gray!80!black}{\textbf{Các mã dùng chung hình (HINH\_021):} \texttt{L06\_C4\_B3\_32}}
\end{tcolorbox}
\vspace{4pt}


\begin{tcolorbox}[
    breakable,
    colback=white,
    colframe=blue!60!black,
    boxrule=0.8pt,
    arc=3pt,
    title={\textbf{Câu \#028} \quad \textbf{MÃ CÂU: \texttt{[L06\_C4\_B3\_39]}} \quad | \quad \textbf{Nhóm: \texttt{HINH\_023}} \quad | \quad \textbf{MapID: \texttt{6H4C3}} \hfill \textcolor{yellow!90!white}{[Hình gốc .pyc]}}
]
\textbf{Dạng toán:} Bài toán thực tế liên quan đến diện tích hình vuông. \hfill \textit{(Tự luận -- Vận dụng cao)}
\vspace{4pt}\hrule\vspace{6pt}
\begin{minipage}[t]{0.58\textwidth}
\vspace{0pt}
Người ta mở rộng một cái ao hình vuông về bốn phía như hình vẽ. Sau khi mở rộng, diện tích ao tăng thêm ${720}m^{2}$ . Tìm diện tích ao cũ?
\end{minipage}\hfill
\begin{minipage}[t]{0.39\textwidth}
\vspace{0pt}
\centering

\resizebox{\linewidth}{!}{%
\begin{tikzpicture}
    % Tạo hình vuông lớn (kích thước nhỏ hơn)
    \draw[thick] (-3,-3) rectangle (3,3); % Hình vuông lớn với kích thước 6x6 (tọa độ từ -3 đến 3)
    
    % Tạo hình vuông nhỏ ở giữa (thu nhỏ kích thước)
    \fill[blue!70] (-1,-1) rectangle (1,1); % Hình vuông nhỏ với cạnh bằng 2 (tọa độ từ -1 đến 1)
    
    % Vẽ các đường mũi tên và chú thích khoảng cách
    \draw[<->, thick, blue] (-3,0) -- (-1,0) node[midway, above] {$10m$}; % Mũi tên trái
    \draw[<->, thick, blue] (1,0) -- (3,0) node[midway, above] {$10m$};  % Mũi tên phải
    \draw[<->, thick, blue] (0,-3) -- (0,-1) node[midway, right] {$10m$}; % Mũi tên dưới
    \draw[<->, thick, blue] (0,1) -- (0,3) node[midway, right] {$10m$};  % Mũi tên trên
    
\end{tikzpicture}
}
\end{minipage}

\end{tcolorbox}
\vspace{4pt}


\begin{tcolorbox}[
    breakable,
    colback=white,
    colframe=blue!60!black,
    boxrule=0.8pt,
    arc=3pt,
    title={\textbf{Câu \#029} \quad \textbf{MÃ CÂU: \texttt{[L06\_C4\_B3\_43]}} \quad | \quad \textbf{Nhóm: \texttt{HINH\_024}} \quad | \quad \textbf{MapID: \texttt{6H4V3}} \hfill \textcolor{yellow!90!white}{[Hình gốc .pyc]}}
]
\textbf{Dạng toán:} Toán thực tế tính tiền làm sân khấu hình thang cân \hfill \textit{(Tự luận -- Vận dụng)}
\vspace{4pt}\hrule\vspace{6pt}
\begin{minipage}[t]{0.58\textwidth}
\vspace{0pt}
Trên một mảnh đất hình chữ nhật có chiều dài $26$ m, chiều rộng $23$ m. Người ta làm một sân khấu hình thang cân có kích thước như hình vẽ.

a) Tính diện tích mảnh đất và sân khấu.

b) Để trang trí sân khấu, mỗi mét vuông tốn $700000$ đồng. Khán đài phía dưới trang trí với mỗi mét vuông $600000$ đồng (không tính phần sân khấu). Tính tổng chi phí trang trí.\\
\end{minipage}\hfill
\begin{minipage}[t]{0.39\textwidth}
\vspace{0pt}
\centering

\resizebox{\linewidth}{!}{%
\begin{tikzpicture}[scale=1, font=\footnotesize, line join=round, line cap=round, >=stealth]
        \draw (0,0)--(6,0)--(6,-4)--(0,-4)--cycle
              (1.5,0)coordinate (A)--(2,-1.5)--(4,-1.5)coordinate (C)--(4.5,0) (4,-1.5)--(4,0)coordinate (B);
        \draw [dashed] (1.5,0) ..controls+(40:1) and+(140:1).. (4.5,0);
        \draw pic[draw,angle radius=2mm] {right angle = A--B--C};
        \fill (3,-1.5)+(-90:2mm) node{$9$ m}
              (4,-.75)+(180:2mm) node[rotate=90]{$8$ m}
              (3,.5)+(90:2mm) node{$11$ m};
    \end{tikzpicture}
}
\end{minipage}

\end{tcolorbox}
\vspace{4pt}


\begin{tcolorbox}[
    breakable,
    colback=white,
    colframe=blue!60!black,
    boxrule=0.8pt,
    arc=3pt,
    title={\textbf{Câu \#030} \quad \textbf{MÃ CÂU: \texttt{[L06\_C4\_B3\_42]}} \quad | \quad \textbf{Nhóm: \texttt{HINH\_025}} \quad | \quad \textbf{MapID: \texttt{6H4V3}} \hfill \textcolor{yellow!90!white}{[Hình gốc .pyc]}}
]
\textbf{Dạng toán:} Toán thực tế tính tiền trồng vườn hoa hình thoi \hfill \textit{(Tự luận -- Vận dụng)}
\vspace{4pt}\hrule\vspace{6pt}
\begin{minipage}[t]{0.58\textwidth}
\vspace{0pt}
Trường của bạn An có khu vườn hình chữ nhật với chiều rộng $8$ m và chiều dài lớn hơn chiều rộng $5$ m. Nhà trường muốn làm một bồn hoa hình thoi trong khu vườn (như hình vẽ). Giả sử mỗi mét vuông trồng được $5$ cây hoa và mỗi cây hoa có giá $16000$ đồng (giá cây hoa và giá trồng hoa). Tính số tiền cần để mua đủ hoa trồng kín bồn hoa đó.\\
\end{minipage}\hfill
\begin{minipage}[t]{0.39\textwidth}
\vspace{0pt}
\centering

\resizebox{\linewidth}{!}{%
\begin{tikzpicture}[scale=.75, font=\footnotesize, line join=round, line cap=round, >=stealth]
    \path 
    (0,0) coordinate (A)
    (0,3) coordinate (B)
    (6,3) coordinate (C)
    (6,0) coordinate (D)
    ($(A)!0.5!(B)$) coordinate (E)
    ($(C)!0.5!(B)$) coordinate (F)
    ($(C)!0.5!(D)$) coordinate (G)
    ($(A)!0.5!(D)$) coordinate (H);
    \draw (A)--(B)--(C)--(D)--(A) (E)--(F)--(G)--(H)--(E) ;
    \fill[violet!20] (E)--(F)--(G)--(H)--(E); 
    \end{tikzpicture}
}
\end{minipage}

\end{tcolorbox}
\vspace{4pt}


\begin{tcolorbox}[
    breakable,
    colback=white,
    colframe=blue!60!black,
    boxrule=0.8pt,
    arc=3pt,
    title={\textbf{Câu \#031} \quad \textbf{MÃ CÂU: \texttt{[L06\_C4\_B3\_44]}} \quad | \quad \textbf{Nhóm: \texttt{HINH\_026}} \quad | \quad \textbf{MapID: \texttt{6H4V3}} \hfill \textcolor{yellow!90!white}{[Hình gốc .pyc]}}
]
\textbf{Dạng toán:} Tính diện tích trồng cỏ của mảnh đất HCN cắt 4 góc \hfill \textit{(Tự luận -- Vận dụng)}
\vspace{4pt}\hrule\vspace{6pt}
\begin{minipage}[t]{0.58\textwidth}
\vspace{0pt}
Cho mảnh vườn hình chữ nhật có độ dài hai cạnh là $9$ m và $20$ m. Người ta trồng $4$ bồn cỏ có dạng hình tam giác xung quanh như hình vẽ, phần còn lại trồng hoa. 

Tính diện tích phần đất dùng để trồng cỏ.

Mỗi mét vuông trồng cỏ chi phí mua hạt giống hết $10000$ đồng, công trồng hết $50000$ đồng. Mỗi mét vuông trồng hoa chi phí mua hoa giống hết $80000$ đồng, công trồng hết $70000$ đồng. Hỏi tổng chi phí làm mảnh vườn hết bao nhiêu tiền?\\
\end{minipage}\hfill
\begin{minipage}[t]{0.39\textwidth}
\vspace{0pt}
\centering

\resizebox{\linewidth}{!}{%
\begin{tikzpicture}[scale=0.6, font=\footnotesize, line join=round, line cap=round,>=stealth]
    \path
    (0,0) coordinate (A) 
    (0,6) coordinate (B) 
    (10,0) coordinate (D) 
    ($(B)+(D)-(A)$) coordinate (C)
    (0,2) coordinate (E)
    (0,4) coordinate (F)
    (3,6) coordinate (G)
    (7,6) coordinate (H)
    (10,4) coordinate (I)
    (10,2) coordinate (J)
    (7,0) coordinate (K)
    (3,0) coordinate (L)
    ;
    \draw (A)--(B)--(C)--(D)--cycle;
    \foreach \i/\j/\goc in {A/E/180,F/B/180,I/C/0,D/J/0}{\draw($(\i)!0.5!(\j)$) node[shift={(\goc:14pt)}]{$2$ cm};}
    \foreach \m/\n/\g in {A/L/-90,B/G/90,H/C/90,K/D/-90}{\draw($(\m)!0.5!(\n)$) node[shift={(\g:9pt)}]{$5$ cm};}
        \foreach \x/\y/\z in {A/E/L,F/B/G,H/C/I,J/D/K}{\draw [fill=cyan](\x)--(\y)--(\z)--cycle;}
    \draw[dashed] ($(A)-(0,1)$)--(A)
                  (D)--($(D)-(0,1)$);
    \draw[<->] ($(A)-(0,1)$)--($(D)-(0,1)$) node[sloped,pos=0.5,midway,below]{$20$ m};  
    %%%%%%%
    \draw[dashed] ($(C)+(2,0)$)--(C)
(D)--($(D)+(2,0)$);
\draw[<->] ($(D)+(2,0)$)--($(C)+(2,0)$) node[pos=0.5,right]{$9$ m}; 
%%%%%% 
    %\foreach \x/\g in {A/180,B/90,C/90,D/-40}  \fill[black] (\x)  +(\g:3mm) node {$\x$};
\end{tikzpicture}
}
\end{minipage}

\end{tcolorbox}
\vspace{4pt}


\begin{tcolorbox}[
    breakable,
    colback=white,
    colframe=blue!60!black,
    boxrule=0.8pt,
    arc=3pt,
    title={\textbf{Câu \#032} \quad \textbf{MÃ CÂU: \texttt{[L06\_C4\_B3\_45]}} \quad | \quad \textbf{Nhóm: \texttt{HINH\_027}} \quad | \quad \textbf{MapID: \texttt{6H4V3}} \hfill \textcolor{yellow!90!white}{[Hình gốc .pyc]}}
]
\textbf{Dạng toán:} Tính diện tích và tiền trồng hoa và cỏ của mảnh đất gồm HCN và hthang \hfill \textit{(Tự luận -- Vận dụng)}
\vspace{4pt}\hrule\vspace{6pt}
\begin{minipage}[t]{0.58\textwidth}
\vspace{0pt}
Một mảnh đất có hình dạng ghép bởi một hình chữ nhật và một hình thang. Kích thước của mảnh đất được mô tả bằng bản vẽ dưới đây. 

Trên mảnh đất đó, người ta trồng ngô thì thu được $10000$ đồng/$m^{2}$, trồng mía thì thu được $20000$ đồng/$m^{2}$, trồng thanh long thì thu được $70000$ đồng/$m^{2}$

Em hãy chọn phương án có doanh thu tốt hơn trong hai phương án sau 

Phương án 1: Trồng mía trên cả mảnh đất.

Phương án 2: Trồng ngô trên phần mảnh đất hình thang và trồng thanh long trên mảnh đất hình chữ nhật.\\
\end{minipage}\hfill
\begin{minipage}[t]{0.39\textwidth}
\vspace{0pt}
\centering

\resizebox{\linewidth}{!}{%
\begin{tikzpicture}[line join = round,line cap = round, thick, font = \small, scale = 0.7]
        \path 
        (0:0) coordinate (D)
        +(0:9) coordinate (C)
        +(90:7.2) coordinate (A)
        ($(A)+(0:6)$) coordinate (B)
        ($(A)+(90:3.6)$) coordinate (E)
        ($(B)+(90:3.6)$) coordinate (F)  
        ;
        \draw 
        (A)--(B)--(C)--(D)--cycle
        pic[draw, angle radius = 8pt]{right angle = A--D--C}
        pic[draw, angle radius = 8pt]{right angle = D--A--B}
        (A)--(E)--(F)--(B)
        ;
        \path (D)--(E)--++(-180:0.35) coordinate (xt)
        ($(D)-(E)+(xt)$) coordinate (yt);
        \draw[>=stealth,|<->|] (xt)--(yt) node[font=\scriptsize,midway,left]{$30\ \mathrm{m}$};
    \path (E)--(F)--++(90:0.35) coordinate (xt)
    ($(E)-(F)+(xt)$) coordinate (yt);
    \draw[>=stealth,|<->|] (xt)--(yt) node[font=\scriptsize,midway,above]{$9\ \mathrm{m}$};
    \path (F)--(B)--++(0:0.35) coordinate (xt)
    ($(F)-(B)+(xt)$) coordinate (yt);
    \draw[>=stealth,<->] (xt)--(yt) node[font=\scriptsize,midway,right]{$6\ \mathrm{m}$};
    \path ($(C)!0.06!-90:(B)$) coordinate (xt_1);
    \path ($(B)-(C)+(xt_1)$) coordinate (yt_1);
    \draw[>=stealth,<->] (xt_1)--(yt_1) node[font=\scriptsize,midway,sloped,above]{$26\ \mathrm{m}$};
    \path (D)--(C)--++(-90:0.35) coordinate (xt)
    ($(D)-(C)+(xt)$) coordinate (yt);
    \draw[>=stealth,|<->|] (xt)--(yt) node[font=\scriptsize,midway,sloped,below]{$19\ \mathrm{m}$};
    \end{tikzpicture}
}
\end{minipage}

\end{tcolorbox}
\vspace{4pt}


\begin{tcolorbox}[
    breakable,
    colback=white,
    colframe=blue!60!black,
    boxrule=0.8pt,
    arc=3pt,
    title={\textbf{Câu \#033} \quad \textbf{MÃ CÂU: \texttt{[L06\_C4\_B3\_46]}} \quad | \quad \textbf{Nhóm: \texttt{HINH\_028}} \quad | \quad \textbf{MapID: \texttt{6H4V3}} \hfill \textcolor{yellow!90!white}{[Hình gốc .pyc]}}
]
\textbf{Dạng toán:} Tính d.tích và tiền trồng cỏ của mảnh đất gồm HCN và hình bình hành \hfill \textit{(Tự luận -- Vận dụng)}
\vspace{4pt}\hrule\vspace{6pt}
\begin{minipage}[t]{0.58\textwidth}
\vspace{0pt}
Một mảnh đất hình chữ nhật có chiều dài là $40$ m, chiều rộng $30$ m. Người ta chia mảnh đất thành hai khu, một khu trồng hoa và một khu trồng cỏ. Hoa sẽ trồng trong khu vực hình bình hành và cỏ trồng ở phần đất còn lại (như hình vẽ).

a) Tính diện tích mảnh đất hình chữ nhật.

b) Tính diện tích phần trồng hoa.

c) Tính số tiền cần trả để trồng cỏ, biết giá mỗi mét vuông trồng cỏ là $40000$ đồng?\\
\end{minipage}\hfill
\begin{minipage}[t]{0.39\textwidth}
\vspace{0pt}
\centering

\resizebox{\linewidth}{!}{%
\begin{tikzpicture}[scale=.9, line cap=round, line join=round,font=\footnotesize,>=stealth]
            \def\dai{6} \def\rong{5} 
            \pgfmathsetmacro\goc{atan((\dai)/(2*\rong))}
            \path (0,0) coordinate (A)++(0:\dai) coordinate (B)++(90:\rong) coordinate (C)++(180:\dai) coordinate(D)
            ($(A)!.5!(B)$) coordinate (M) ($(C)!.5!(D)$) coordinate (N);
            \draw (A)--(B)--(C)--(D)--cycle
            (N)--(B) (D)--(M)
            ($(A)!.5!(C)$)node[rotate=-90+\goc]{Trồng hoa};
            \foreach \i in {1.5,4.5}{
                \node[rotate=-90+\goc] at ($(A)!\i/6!(C)$) {Trồng cỏ};
            }
            \draw[<->,dashed] ([xshift=-.4cm]A)--node[midway,rotate=90,fill=white]{$40$ m}([xshift=-.4cm]D);
            \draw[<->,dashed] ([yshift=.3cm]D)--node[midway,fill=white]{$15$ m}([yshift=.3cm]N);
            \draw[<->,dashed] ([yshift=.3cm]C)--node[midway,fill=white]{$15$ m}([yshift=.3cm]N);
        \end{tikzpicture}
}
\end{minipage}

\end{tcolorbox}
\vspace{4pt}


\begin{tcolorbox}[
    breakable,
    colback=white,
    colframe=blue!60!black,
    boxrule=0.8pt,
    arc=3pt,
    title={\textbf{Câu \#034} \quad \textbf{MÃ CÂU: \texttt{[L06\_C4\_B3\_47]}} \quad | \quad \textbf{Nhóm: \texttt{HINH\_029}} \quad | \quad \textbf{MapID: \texttt{6H4V3}} \hfill \textcolor{yellow!90!white}{[Hình gốc .pyc]}}
]
\textbf{Dạng toán:} Tính diện tích và tiền trồng hoa và cỏ của mảnh đất gồm HCN và hình bình hành \hfill \textit{(Tự luận -- Vận dụng)}
\vspace{4pt}\hrule\vspace{6pt}
\begin{minipage}[t]{0.58\textwidth}
\vspace{0pt}
Bác Lan có một mảnh đất hình chữ nhật có diện tích là $559$ m$^{2}$, chiều rộng $13$ m. Ở giữa đất bác làm lối đi như hình vẽ, phần đất còn lại để trồng cây.

a) Tính chu vi mảnh đất hình chữ nhật.

b) Tính diện tích của phần đất trồng cây.

c) Bác Lan định thuê thợ xây tường rào xung quanh khu đất chỉ để $1$ cửa đi vào rộng $3$ m. Giá làm $1$ m tường rào là $80000$ đồng. Số tiền bác Lan phải trả khi xây tường rào là bao nhiêu? )\\
\end{minipage}\hfill
\begin{minipage}[t]{0.39\textwidth}
\vspace{0pt}
\centering

\resizebox{\linewidth}{!}{%
\begin{tikzpicture}
    \def\dai{5}
    \def\rong{4}
    \def\loidi{0.5}
    \draw[latex-latex] (\dai+0.5,0)--++(90:\rong); 
    \draw (0,0)--++(90:\rong)--++(0:\dai) node[midway,above]{$3$ m}--++(-90:\rong)node[fill=white,inner sep=4pt,rectangle,midway,right]{\color{black} $13$ m}--++(180:\dai) node[midway,below]{$3$ m}
    (\dai/2-\loidi/2,0) rectangle (\dai/2+\loidi/2,\rong)
    ;
    \end{tikzpicture}
}
\end{minipage}

\end{tcolorbox}
\vspace{4pt}


\subsection{Chương 5 - Tính đối xứng của hình phẳng trong tự nhiên}

\begin{tcolorbox}[
    breakable,
    colback=white,
    colframe=blue!60!black,
    boxrule=0.8pt,
    arc=3pt,
    title={\textbf{Câu \#035} \quad \textbf{MÃ CÂU: \texttt{[L06\_C5\_B1\_01]}} \quad | \quad \textbf{Nhóm: \texttt{HINH\_030}} \quad | \quad \textbf{MapID: \texttt{6H5N1}} \hfill \textcolor{yellow!90!white}{[Hình gốc .pyc]}}
]
\textbf{Dạng toán:} Tìm trục đối xứng \hfill \textit{(Tự luận -- Nhận biết)}
\vspace{4pt}\hrule\vspace{6pt}
\begin{minipage}[t]{0.58\textwidth}
\vspace{0pt}
Tìm trục đối xứng của hình ảnh sau
\end{minipage}\hfill
\begin{minipage}[t]{0.39\textwidth}
\vspace{0pt}
\centering

\resizebox{\linewidth}{!}{%
\begin{tikzpicture}
    
    % Vẽ trục đối xứng (trục tung)
    %\draw[thick, dashed] (0,-4) -- (0,4) node[above] {Trục đối xứng};
    
    % Vẽ các hình đối xứng qua trục
    % Vẽ hình tròn đối xứng
    \draw[fill=blue!20] (-1,2) circle(0.5cm);
    \draw[fill=blue!20] (1,2) circle(0.5cm);
    
    % Vẽ các hình tam giác đối xứng
    \draw[fill=red!30] (-2,0) -- (-1,1) -- (-3,1) -- cycle;
    \draw[fill=red!30] (2,0) -- (3,1) -- (1,1) -- cycle;
    
    % Vẽ các đường thẳng đối xứng
    \draw[thick, green] (-1,-2) -- (-1,-3);
    \draw[thick, green] (1,-2) -- (1,-3);
    
    % Vẽ các hình chữ nhật đối xứng
    \draw[fill=yellow!30] (-3,-1) rectangle (-2,-2);
    \draw[fill=yellow!30] (3,-1) rectangle (2,-2);
    
    % Đánh dấu các điểm đối xứng
    \filldraw[red] (-1,2) circle(2pt);
    \filldraw[red] (1,2) circle(2pt);
    \filldraw[red] (-1,-3) circle(2pt);
    \filldraw[red] (1,-3) circle(2pt);
    
\end{tikzpicture}
}
\end{minipage}

\end{tcolorbox}
\vspace{4pt}


\begin{tcolorbox}[
    breakable,
    colback=white,
    colframe=blue!60!black,
    boxrule=0.8pt,
    arc=3pt,
    title={\textbf{Câu \#036} \quad \textbf{MÃ CÂU: \texttt{[L06\_C5\_B1\_02]}} \quad | \quad \textbf{Nhóm: \texttt{HINH\_031}} \quad | \quad \textbf{MapID: \texttt{6H5N1}} \hfill \textcolor{yellow!90!white}{[Hình gốc .pyc]}}
]
\textbf{Dạng toán:} Tìm trục đối xứng ở hình cờ Tổ Quốc \hfill \textit{(Tự luận -- Nhận biết)}
\vspace{4pt}\hrule\vspace{6pt}
\begin{minipage}[t]{0.58\textwidth}
\vspace{0pt}
Tìm trục đối xứng của hình ảnh sau
\end{minipage}\hfill
\begin{minipage}[t]{0.39\textwidth}
\vspace{0pt}
\centering

\resizebox{\linewidth}{!}{%
\begin{tikzpicture}
    % Vẽ nền cờ đỏ
    \fill[red] (0,0) rectangle (6,4);
    
    % Tâm của ngôi sao
    \coordinate (O) at (3,2);
    
    % Bán kính và các điểm cho ngôi sao vàng
    \def\R{1} % Bán kính lớn của ngôi sao
    \def\r{0.4} % Bán kính nhỏ của ngôi sao
    
    % Các điểm của ngôi sao (theo chiều kim đồng hồ)
    \foreach \i in {18,90,162,234,306} {
        \coordinate (\i) at ($(O)+(\i:\R)$); % Điểm trên bán kính lớn
    }
    
    % Các điểm trong của ngôi sao (theo chiều kim đồng hồ)
    \foreach \i in {54,126,198,270,342} {
        \coordinate (\i) at ($(O)+(\i:\r)$); % Điểm trên bán kính nhỏ
    }
    
    % Vẽ ngôi sao vàng
    \fill[yellow] (18) -- (54) -- (90) -- (126) -- (162) -- (198) -- (234) -- (270) -- (306) -- (342) -- cycle;
    
\end{tikzpicture}
}
\end{minipage}

\end{tcolorbox}
\vspace{4pt}


\begin{tcolorbox}[
    breakable,
    colback=white,
    colframe=blue!60!black,
    boxrule=0.8pt,
    arc=3pt,
    title={\textbf{Câu \#037} \quad \textbf{MÃ CÂU: \texttt{[L06\_C5\_B1\_03]}} \quad | \quad \textbf{Nhóm: \texttt{HINH\_032}} \quad | \quad \textbf{MapID: \texttt{6H5N1}} \hfill \textcolor{yellow!90!white}{[Hình gốc .pyc]}}
]
\textbf{Dạng toán:} Tìm trục đối xứng ở hình lá cây \hfill \textit{(Tự luận -- Nhận biết)}
\vspace{4pt}\hrule\vspace{6pt}
\begin{minipage}[t]{0.58\textwidth}
\vspace{0pt}
Tìm trục đối xứng của hình ảnh sau
\end{minipage}\hfill
\begin{minipage}[t]{0.39\textwidth}
\vspace{0pt}
\centering

\resizebox{\linewidth}{!}{%
\begin{tikzpicture}
    
    % Vẽ hình chiếc lá đối xứng với màu xanh lá cây
    \fill[green!40!yellow] 
    (0,2) .. controls (0.6,1.5) and (1.2,0.5) .. (0,-2) % Nửa phải của lá
    (0,-2) .. controls (-1.2,0.5) and (-0.6,1.5) .. (0,2);  % Nửa trái của lá
    
    % Viền của lá
    \draw[thick, green!70!black] 
    (0,2) .. controls (0.6,1.5) and (1.2,0.5) .. (0,-2)  % Nửa phải của lá
    (0,-2) .. controls (-1.2,0.5) and (-0.6,1.5) .. (0,2);  % Nửa trái của lá
    
    % Vẽ gân lá
    \draw[thick, green!50!black]
    (0,2) -- (0,-2);  % Gân chính giữa
    
    % Vẽ các gân lá đối xứng (bên phải)
    \draw[thick, green!50!black] 
    (0,0) .. controls (0.3,0.5) .. (0.6,0.8);  % Gân phải trên
    \draw[thick, green!50!black] 
    (0,-0.5) .. controls (0.3,-0.8) .. (0.6,-0.5);  % Gân phải dưới
    
    % Vẽ các gân lá đối xứng (bên trái)
    \draw[thick, green!50!black] 
    (0,0) .. controls (-0.3,0.5) .. (-0.6,0.8);  % Gân trái trên
    \draw[thick, green!50!black] 
    (0,-0.5) .. controls (-0.3,-0.8) .. (-0.6,-0.5);  % Gân trái dưới
    
    % Thêm các chi tiết hoạt hình (những giọt sương)
    \fill[blue!70!cyan] (0.4,1.2) circle(1.5pt);  % Giọt sương bên phải
        \fill[blue!70!cyan] (-0.4,1.2) circle(1.5pt);  % Giọt sương bên phải
    \fill[blue!70!cyan] (-0.5,0.3) circle(1.5pt);  % Giọt sương bên trái
    \fill[blue!70!cyan] (0.5,0.3) circle(1.5pt);  % Giọt sương phía dưới
        \fill[blue!70!cyan] (-0.2,-0.3) circle(1.5pt);  % Giọt sương bên trái
    \fill[blue!70!cyan] (0.2,-0.3) circle(1.5pt);  % Giọt sương phía dưới
\end{tikzpicture}
}
\end{minipage}

\end{tcolorbox}
\vspace{4pt}


\begin{tcolorbox}[
    breakable,
    colback=white,
    colframe=blue!60!black,
    boxrule=0.8pt,
    arc=3pt,
    title={\textbf{Câu \#038} \quad \textbf{MÃ CÂU: \texttt{[L06\_C5\_B1\_04]}} \quad | \quad \textbf{Nhóm: \texttt{HINH\_033}} \quad | \quad \textbf{MapID: \texttt{6H5N1}} \hfill \textcolor{yellow!90!white}{[Hình gốc .pyc]}}
]
\textbf{Dạng toán:} Tìm trục đối xứng của đa giác \hfill \textit{(Tự luận -- Nhận biết)}
\vspace{4pt}\hrule\vspace{6pt}
\begin{minipage}[t]{0.58\textwidth}
\vspace{0pt}
Tìm trục đối xứng của hình ảnh sau
\end{minipage}\hfill
\begin{minipage}[t]{0.39\textwidth}
\vspace{0pt}
\centering

\resizebox{\linewidth}{!}{%
\begin{tikzpicture}
    % Vẽ lục giác đều với bán kính 2
    \draw[thick] (0,0) -- (30:2) -- (90:2) -- (150:2) -- (210:2) -- (270:2) -- (330:2) -- cycle;
\end{tikzpicture}
}
\end{minipage}

\end{tcolorbox}
\vspace{4pt}


\begin{tcolorbox}[
    breakable,
    colback=white,
    colframe=blue!60!black,
    boxrule=0.8pt,
    arc=3pt,
    title={\textbf{Câu \#039} \quad \textbf{MÃ CÂU: \texttt{[L06\_C5\_B1\_05]}} \quad | \quad \textbf{Nhóm: \texttt{HINH\_034}} \quad | \quad \textbf{MapID: \texttt{6H5N1}} \hfill \textcolor{yellow!90!white}{[Hình gốc .pyc]}}
]
\textbf{Dạng toán:} Tìm trục đối xứng của chữ thập \hfill \textit{(Tự luận -- Nhận biết)}
\vspace{4pt}\hrule\vspace{6pt}
\begin{minipage}[t]{0.58\textwidth}
\vspace{0pt}
Tìm trục đối xứng của hình ảnh sau
\end{minipage}\hfill
\begin{minipage}[t]{0.39\textwidth}
\vspace{0pt}
\centering

\resizebox{\linewidth}{!}{%
\begin{tikzpicture}
    
    % Kích thước của lá cờ Anh
    \def\flagWidth{8}
    \def\flagHeight{5}
    
    % Vẽ nền cờ màu trắng
    \fill[white] (0,0) rectangle (\flagWidth,\flagHeight);
    
    % Vẽ đường thẳng đứng màu đỏ ở giữa
    \fill[red] ({\flagWidth/2-0.3}, 0) rectangle ({\flagWidth/2+0.3}, \flagHeight);
    
    % Vẽ đường ngang màu đỏ ở giữa
    \fill[red] (0, {\flagHeight/2-0.3}) rectangle (\flagWidth, {\flagHeight/2+0.3});
    
\end{tikzpicture}
}
\end{minipage}

\end{tcolorbox}
\vspace{4pt}


\begin{tcolorbox}[
    breakable,
    colback=white,
    colframe=blue!60!black,
    boxrule=0.8pt,
    arc=3pt,
    title={\textbf{Câu \#040} \quad \textbf{MÃ CÂU: \texttt{[L06\_C5\_B1\_06]}} \quad | \quad \textbf{Nhóm: \texttt{HINH\_035}} \quad | \quad \textbf{MapID: \texttt{6H5N1}} \hfill \textcolor{yellow!90!white}{[Hình gốc .pyc]}}
]
\textbf{Dạng toán:} Tìm trục đối xứng của ngôi nhà \hfill \textit{(Tự luận -- Nhận biết)}
\vspace{4pt}\hrule\vspace{6pt}
\begin{minipage}[t]{0.58\textwidth}
\vspace{0pt}
Tìm trục đối xứng của hình ảnh sau
\end{minipage}\hfill
\begin{minipage}[t]{0.39\textwidth}
\vspace{0pt}
\centering

\resizebox{\linewidth}{!}{%
\begin{tikzpicture}
    % Vẽ thân nhà (hình chữ nhật)
    \draw[fill=blue!20] (-2,0) rectangle (2,-2);
    
    % Vẽ mái nhà (hình tam giác) với màu đỏ
    \draw[fill=red] (-2,0) -- (0,1.5) -- (2,0) -- cycle;
    
    % Vẽ cửa chính (hình chữ nhật nhỏ)
    \draw[fill=brown!50] (-0.5,-2) rectangle (0.5,-1);
    
    % Vẽ hai cửa sổ (hình vuông đối xứng)
    \draw[fill=white] (-1.5,-0.5) rectangle (-0.5,-1.5);
    \draw[fill=white] (0.5,-0.5) rectangle (1.5,-1.5);
    
    % Vẽ các đường chia cửa sổ
    \draw (-1,-0.5) -- (-1,-1.5);
    \draw (-1.5,-1) -- (-0.5,-1);
    \draw (1,-0.5) -- (1,-1.5);
    \draw (0.5,-1) -- (1.5,-1);
\end{tikzpicture}
}
\end{minipage}

\end{tcolorbox}
\vspace{4pt}


\begin{tcolorbox}[
    breakable,
    colback=white,
    colframe=blue!60!black,
    boxrule=0.8pt,
    arc=3pt,
    title={\textbf{Câu \#041} \quad \textbf{MÃ CÂU: \texttt{[L06\_C5\_B1\_07]}} \quad | \quad \textbf{Nhóm: \texttt{HINH\_036}} \quad | \quad \textbf{MapID: \texttt{6H5N1}} \hfill \textcolor{yellow!90!white}{[Hình gốc .pyc]}}
]
\textbf{Dạng toán:} Tìm trục đối xứng của cờ ISRAEL \hfill \textit{(Tự luận -- Nhận biết)}
\vspace{4pt}\hrule\vspace{6pt}
\begin{minipage}[t]{0.58\textwidth}
\vspace{0pt}
Tìm trục đối xứng của hình ảnh sau
\end{minipage}\hfill
\begin{minipage}[t]{0.39\textwidth}
\vspace{0pt}
\centering

\resizebox{\linewidth}{!}{%
\begin{tikzpicture}
    
    % Kích thước cờ nhỏ (6:4 thay vì 11:8)
    \draw[white, fill=white] (0,0) rectangle (6,4);
    
    % Vẽ hai dải ngang màu xanh lam
    \draw[fill=blue] (0,3.6) rectangle (6,4);
    \draw[fill=blue] (0,0.4) rectangle (6,0);
    
    % Tọa độ ngôi sao David lớn hơn
    \begin{scope}[shift={(3,2)}, scale=1.0] % Phóng to ngôi sao David
        % Tam giác hướng lên
        \draw[blue, thick] (0,1.5) -- (1.3,-0.75) -- (-1.3,-0.75) -- cycle;
        % Tam giác hướng xuống
        \draw[blue, thick] (0,-1.5) -- (1.3,0.75) -- (-1.3,0.75) -- cycle;
    \end{scope}
    
\end{tikzpicture}
}
\end{minipage}

\end{tcolorbox}
\vspace{4pt}


\begin{tcolorbox}[
    breakable,
    colback=white,
    colframe=blue!60!black,
    boxrule=0.8pt,
    arc=3pt,
    title={\textbf{Câu \#042} \quad \textbf{MÃ CÂU: \texttt{[L06\_C5\_B2\_01]}} \quad | \quad \textbf{Nhóm: \texttt{HINH\_037}} \quad | \quad \textbf{MapID: \texttt{6H5N2}} \hfill \textcolor{yellow!90!white}{[Hình gốc .pyc]}}
]
\textbf{Dạng toán:} Tìm trục đối xứng tâm đối xứng của quạt \hfill \textit{(Tự luận -- Nhận biết)}
\vspace{4pt}\hrule\vspace{6pt}
\begin{minipage}[t]{0.58\textwidth}
\vspace{0pt}
Tìm trục đối xứng và tâm đối xứng của hình ảnh sau
\end{minipage}\hfill
\begin{minipage}[t]{0.39\textwidth}
\vspace{0pt}
\centering

\resizebox{\linewidth}{!}{%
\begin{tikzpicture}
    
    % Tâm đối xứng của hình
    \coordinate (O) at (0,0);
    
    % Vẽ các tam giác đối xứng qua tâm O
    \fill[blue!40] (0,0) -- (2,1) -- (1,2) -- cycle;  % Tam giác trên bên phải
    \fill[blue!40] (0,0) -- (2,-1) -- (1,-2) -- cycle;  % Tam giác dưới bên phải
    \fill[blue!40] (0,0) -- (-2,-1) -- (-1,-2) -- cycle;  % Tam giác dưới bên trái
    \fill[blue!40] (0,0) -- (-2,1) -- (-1,2) -- cycle;  % Tam giác trên bên trái
    
    % Vẽ đường viền cho các tam giác
    \draw[thick, blue!70] (0,0) -- (2,1) -- (1,2) -- cycle;
    \draw[thick, blue!70] (0,0) -- (2,-1) -- (1,-2) -- cycle;
    \draw[thick, blue!70] (0,0) -- (-2,-1) -- (-1,-2) -- cycle;
    \draw[thick, blue!70] (0,0) -- (-2,1) -- (-1,2) -- cycle;
    
    % Vẽ tâm đối xứng
    \fill[red] (O) circle(2pt);  % Tâm đối xứng màu đỏ
    
\end{tikzpicture}
}
\end{minipage}

\end{tcolorbox}
\vspace{4pt}


\begin{tcolorbox}[
    breakable,
    colback=white,
    colframe=blue!60!black,
    boxrule=0.8pt,
    arc=3pt,
    title={\textbf{Câu \#043} \quad \textbf{MÃ CÂU: \texttt{[L06\_C5\_B2\_02]}} \quad | \quad \textbf{Nhóm: \texttt{HINH\_038}} \quad | \quad \textbf{MapID: \texttt{6H5N2}} \hfill \textcolor{yellow!90!white}{[Hình gốc .pyc]}}
]
\textbf{Dạng toán:} Tìm trục đối xứng tâm đối xứng của hình3sep \hfill \textit{(Tự luận -- Nhận biết)}
\vspace{4pt}\hrule\vspace{6pt}
\begin{minipage}[t]{0.58\textwidth}
\vspace{0pt}
Tìm trục đối xứng và tâm đối xứng của hình ảnh sau
\end{minipage}\hfill
\begin{minipage}[t]{0.39\textwidth}
\vspace{0pt}
\centering

\resizebox{\linewidth}{!}{%
\begin{tikzpicture}[scale =0.5]
    
    % Vẽ nền tròn đối xứng
    \draw[blue, thick] (0,0) circle (5);
    
    % Vẽ các đường xoắn ốc
    \foreach \angle in {0, 45, 90, 135, 180, 225, 270, 315} {
        \draw[red, thick, domain=0:4, samples=100] plot ({\x*cos(\angle) - \x*sin(\angle)}, {\x*sin(\angle) + \x*cos(\angle)});
    }
    
    % Vẽ các ngôi sao đối xứng qua tâm
    \foreach \a in {0, 90, 180, 270} {
        \draw[fill=yellow, rotate=\a] (3,0) -- (3.5,0.5) -- (4,0) -- (3.5,-0.5) -- cycle;
    }
    
    % Vẽ các đường thẳng chia hình tròn thành các phần đối xứng
    \foreach \ang in {0, 30, 60, 90, 120, 150, 180, 210, 240, 270, 300, 330} {
        \draw[green, thick] (0,0) -- ({5*cos(\ang)}, {5*sin(\ang)});
    }
    
    % Thêm các vòng tròn nhỏ đối xứng qua tâm
    \foreach \radius in {1, 2, 3, 4} {
        \draw[dashed, purple] (0,0) circle (\radius);
    }
    
\end{tikzpicture}
}
\end{minipage}

\end{tcolorbox}
\vspace{4pt}


\subsection{Chương 8 -  Những hình hình học cơ bản}

\begin{tcolorbox}[
    breakable,
    colback=white,
    colframe=blue!60!black,
    boxrule=0.8pt,
    arc=3pt,
    title={\textbf{Câu \#044} \quad \textbf{MÃ CÂU: \texttt{[L06\_C8\_B1\_17]}} \quad | \quad \textbf{Nhóm: \texttt{HINH\_039}} \quad | \quad \textbf{MapID: \texttt{6H8N1}} \hfill \textcolor{yellow!90!white}{[Hình gốc .pyc]}}
]
\textbf{Dạng toán:} Nhận biết điểm thuộc đường thẳng \hfill \textit{(Trắc nghiệm -- Nhận biết)}
\vspace{4pt}\hrule\vspace{6pt}
\begin{minipage}[t]{0.58\textwidth}
\vspace{0pt}
Quan sát hình vẽ và chọn khẳng định đúng.

\vspace{4pt}\noindent
\textbf{A.} ${D\in i}$ \quad \textbf{B.} ${N\in a}$ \quad \textbf{\textcolor{blue!75!black}{C. (Đ)}} ${D\not\in m}$ \quad \textbf{D.} ${J\not\in m}$
\end{minipage}\hfill
\begin{minipage}[t]{0.39\textwidth}
\vspace{0pt}
\centering

\resizebox{\linewidth}{!}{%
\begin{tikzpicture}[>=stealth,line join=round,line cap=round,font=\footnotesize,scale=0.6]
        \draw
        (-1,0) coordinate[label=above:$z$](A1)
        (0,0) coordinate[label=below:$D$](A)
        (6,0) coordinate[label=above right:$N$](C)
        (6.5,0) coordinate(C2)
        (3.5,0) coordinate(D)
        (4,5) coordinate[label=left:$J$](B);
        \tkzDefPointBy[homothety=center B ratio 1.2](A)\tkzGetPoint{A2}
        \tkzDefPointBy[homothety=center B ratio 1.2](C)\tkzGetPoint{C1}
        \tkzDefPointBy[homothety=center B ratio 1.2](D)\tkzGetPoint{D1}
        \tkzDefPointBy[homothety=center A ratio 1.2](B)\tkzGetPoint{B1}
        \tkzDefPointBy[homothety=center D ratio 1.2](B)\tkzGetPoint{B2}
        \tkzDefPointBy[homothety=center C ratio 1.2](B)\tkzGetPoint{B3}
        \draw (A1)--(C2) (A2)--(B1) (D1)--(B2) (C1)--(B3);
        \foreach \x in {A,B,C}
        {\fill (\x) circle (2pt);}
        \node[below] at (current bounding box.south){Hình a};
        \path
        (D1) node[left]{$m$}
        (B1) node[right]{$a$}
        (B3) node[left]{$i$};
        \end{tikzpicture}
}
\end{minipage}

\end{tcolorbox}
\vspace{4pt}


\begin{tcolorbox}[
    breakable,
    colback=white,
    colframe=blue!60!black,
    boxrule=0.8pt,
    arc=3pt,
    title={\textbf{Câu \#045} \quad \textbf{MÃ CÂU: \texttt{[L06\_C8\_B1\_18]}} \quad | \quad \textbf{Nhóm: \texttt{HINH\_040}} \quad | \quad \textbf{MapID: \texttt{6H8N1}} \hfill \textcolor{yellow!90!white}{[Hình gốc .pyc]}}
]
\textbf{Dạng toán:} Chỉ ra ba điểm thẳng hàng \hfill \textit{(Trắc nghiệm -- Nhận biết)}
\vspace{4pt}\hrule\vspace{6pt}
\begin{minipage}[t]{0.58\textwidth}
\vspace{0pt}
Trong hình dưới, cho bốn điểm ${B}$, ${G}$, ${H}$, ${J}$ thuộc đường thẳng ${z}$ và điểm ${Q}$ không thuộc đường thẳng ${z}$. Hãy chọn bộ ba điểm thẳng hàng.

\vspace{4pt}\noindent
\textbf{A.} $\left( B, J, Q\right)$ \quad \textbf{B.} $\left(G, H, Q\right)$ \quad \textbf{\textcolor{blue!75!black}{C. (Đ)}} $\left(G, H, J\right)$ \quad \textbf{D.} $\left(B, H, Q\right)$
\end{minipage}\hfill
\begin{minipage}[t]{0.39\textwidth}
\vspace{0pt}
\centering

\resizebox{\linewidth}{!}{%
\begin{tikzpicture}[>=stealth,line join=round,line cap=round,font=\footnotesize,scale=1]
    \draw
    (0,0) coordinate(O)
    (1,0) coordinate[label=below:$B$](A)
    (1.5,0) coordinate[label=below:$G$](B)
    (2.5,0) coordinate[label=below:$H$](C)
    (4,0) coordinate[label=below:$J$](D)
    (3,0.5) coordinate[label=above:$Q$](E)
    (4.5,0) coordinate[label=above:$z$](F);
    \draw (O)--(F); 
    \foreach \x in {A,B,C,D,E}
    {\fill (\x) circle (2pt);}
    
    \end{tikzpicture}
}
\end{minipage}

\end{tcolorbox}
\vspace{4pt}


\begin{tcolorbox}[
    breakable,
    colback=white,
    colframe=blue!60!black,
    boxrule=0.8pt,
    arc=3pt,
    title={\textbf{Câu \#046} \quad \textbf{MÃ CÂU: \texttt{[L06\_C8\_B1\_19]}} \quad | \quad \textbf{Nhóm: \texttt{HINH\_041}} \quad | \quad \textbf{MapID: \texttt{6H8N1}} \hfill \textcolor{yellow!90!white}{[Hình gốc .pyc]}}
]
\textbf{Dạng toán:} Chỉ ra điểm nằm giữa \hfill \textit{(Trắc nghiệm -- Nhận biết)}
\vspace{4pt}\hrule\vspace{6pt}
\begin{minipage}[t]{0.58\textwidth}
\vspace{0pt}
Trong hình dưới, hãy chỉ ra tất cả các điểm nằm giữa hai điểm ${F}$ và ${Q}$.

\vspace{4pt}\noindent
\textbf{A.} ${R}$, ${T}$ \quad \textbf{\textcolor{blue!75!black}{B. (Đ)}} ${X}$, ${R}$, ${T}$ \quad \textbf{C.} ${T}$ \quad \textbf{D.} ${X}$
\end{minipage}\hfill
\begin{minipage}[t]{0.39\textwidth}
\vspace{0pt}
\centering

\resizebox{\linewidth}{!}{%
\begin{tikzpicture}[>=stealth,line join=round,line cap=round,font=\footnotesize,scale=1]
    \draw
    (0,0) coordinate(O)
    (1,0) coordinate[label=below:$F$](A)
    (2,0) coordinate[label=below:$X$](B)
    (3.6,0) coordinate[label=below:$R$](C)
    (5,0) coordinate[label=below:$T$](D)
    (7,0) coordinate[label=below:$Q$](N)
    (7.5,0) coordinate(H);
    \draw (O)--(H); 
    \foreach \x in {A,B,C,D,N}
    {\fill (\x) circle (2pt);}
    \node[below] at (current bounding box.south){Hình a};
    \end{tikzpicture}
}
\end{minipage}

\end{tcolorbox}
\vspace{4pt}


\begin{tcolorbox}[
    breakable,
    colback=white,
    colframe=blue!60!black,
    boxrule=0.8pt,
    arc=3pt,
    title={\textbf{Câu \#047} \quad \textbf{MÃ CÂU: \texttt{[L06\_C8\_B1\_20]}} \quad | \quad \textbf{Nhóm: \texttt{HINH\_042}} \quad | \quad \textbf{MapID: \texttt{6H8N1}} \hfill \textcolor{yellow!90!white}{[Hình gốc .pyc]}}
]
\textbf{Dạng toán:} Chỉ ra điểm không nằm giữa hai điểm \hfill \textit{(Trắc nghiệm -- Nhận biết)}
\vspace{4pt}\hrule\vspace{6pt}
\begin{minipage}[t]{0.58\textwidth}
\vspace{0pt}
Trong hình dưới, hãy chỉ ra tất cả các điểm không nằm giữa hai điểm ${J}$ và ${P}$.

\vspace{4pt}\noindent
\textbf{A.} ${O}$ \quad \textbf{B.} ${S}$, ${P}$ \quad \textbf{C.} ${O}$, ${P}$ \quad \textbf{\textcolor{blue!75!black}{D. (Đ)}} ${N}$
\end{minipage}\hfill
\begin{minipage}[t]{0.39\textwidth}
\vspace{0pt}
\centering

\resizebox{\linewidth}{!}{%
\begin{tikzpicture}[>=stealth,line join=round,line cap=round,font=\footnotesize,scale=1]
    \draw
    (0,0) coordinate(O)
    (1,0) coordinate[label=below:$J$](A)
    (2,0) coordinate[label=below:$S$](B)
    (3.6,0) coordinate[label=below:$O$](C)
    (5,0) coordinate[label=below:$P$](D)
    (7,0) coordinate[label=below:$N$](N)
    (7.5,0) coordinate(H);
    \draw (O)--(H); 
    \foreach \x in {A,B,C,D,N}
    {\fill (\x) circle (2pt);}
    \node[below] at (current bounding box.south){Hình a};
    \end{tikzpicture}
}
\end{minipage}

\end{tcolorbox}
\vspace{4pt}


\begin{tcolorbox}[
    breakable,
    colback=white,
    colframe=blue!60!black,
    boxrule=0.8pt,
    arc=3pt,
    title={\textbf{Câu \#048} \quad \textbf{MÃ CÂU: \texttt{[L06\_C8\_B1\_21]}} \quad | \quad \textbf{Nhóm: \texttt{HINH\_043}} \quad | \quad \textbf{MapID: \texttt{6H8N1}} \hfill \textcolor{yellow!90!white}{[Hình gốc .pyc]}}
]
\textbf{Dạng toán:} Đếm số giao điểm \hfill \textit{(Trắc nghiệm -- Nhận biết)}
\vspace{4pt}\hrule\vspace{6pt}
\begin{minipage}[t]{0.58\textwidth}
\vspace{0pt}
Đếm số giao điểm tạo bởi ba đường thẳng trong hình.

\vspace{4pt}\noindent
\textbf{A.} $1$ \quad \textbf{\textcolor{blue!75!black}{B. (Đ)}} $2$ \quad \textbf{C.} $0$ \quad \textbf{D.} $3$
\end{minipage}\hfill
\begin{minipage}[t]{0.39\textwidth}
\vspace{0pt}
\centering

\resizebox{\linewidth}{!}{%
\begin{tikzpicture}[>=stealth,line join=round,line cap=round,font=\footnotesize,scale=1]
        \draw (0,0)--(4,0) (0,1)--(4,1) (1.5,-0.5)--(3,2);
        
        \end{tikzpicture}
}
\end{minipage}

\end{tcolorbox}
\vspace{4pt}


\begin{tcolorbox}[
    breakable,
    colback=white,
    colframe=blue!60!black,
    boxrule=0.8pt,
    arc=3pt,
    title={\textbf{Câu \#049} \quad \textbf{MÃ CÂU: \texttt{[L06\_C8\_B1\_22]}} \quad | \quad \textbf{Nhóm: \texttt{HINH\_044}} \quad | \quad \textbf{MapID: \texttt{6H8H1}} \hfill \textcolor{yellow!90!white}{[Hình gốc .pyc]}}
]
\textbf{Dạng toán:} Chỉ ra các đường thẳng song song \hfill \textit{(Trắc nghiệm -- Thông hiểu)}
\vspace{4pt}\hrule\vspace{6pt}
\begin{minipage}[t]{0.58\textwidth}
\vspace{0pt}
Chọn khẳng định đúng.

\vspace{4pt}\noindent
\textbf{A.} ${FJ//VT}$ \quad \textbf{\textcolor{blue!75!black}{B. (Đ)}} ${JG//TV}$ \quad \textbf{C.} ${FG//NT}$ \quad \textbf{D.} ${JG//FG}$
\end{minipage}\hfill
\begin{minipage}[t]{0.39\textwidth}
\vspace{0pt}
\centering

\resizebox{\linewidth}{!}{%
\begin{tikzpicture}[>=stealth,line join=round,line cap=round,font=\footnotesize,scale=0.6]
    \draw
    (-1,0) coordinate(A1)
    (0,0) coordinate[label=above:$G$](A)
    (6,0) coordinate[label=above right:$J$](C)
    (6.5,0) coordinate(C2)
    (3,4) coordinate[label=above:$F$](B);
    \tkzDefMidPoint(A,B)\tkzGetPoint{E}
    \tkzDefMidPoint(A,C)\tkzGetPoint{D}
    \tkzDefMidPoint(C,B)\tkzGetPoint{F}
    \tkzDefPointBy[homothety=center B ratio 1.2](A)\tkzGetPoint{A2}
    \tkzDefPointBy[homothety=center B ratio 1.2](C)\tkzGetPoint{C1}
    \tkzDefPointBy[homothety=center A ratio 1.2](B)\tkzGetPoint{B1}
    \tkzDefPointBy[homothety=center C ratio 1.2](B)\tkzGetPoint{B3}
    \tkzDefPointBy[homothety=center C ratio 1.2](B)\tkzGetPoint{B3}
    \tkzDefPointBy[homothety=center D ratio 1.2](E)\tkzGetPoint{E1}
    \tkzDefPointBy[homothety=center E ratio 1.2](D)\tkzGetPoint{D2}
    \tkzDefPointBy[homothety=center F ratio 1.2](E)\tkzGetPoint{E2}
    \tkzDefPointBy[homothety=center E ratio 1.2](F)\tkzGetPoint{F1}
    \tkzDefPointBy[homothety=center D ratio 1.2](F)\tkzGetPoint{F2}
    \tkzDefPointBy[homothety=center F ratio 1.2](D)\tkzGetPoint{D1}
    \draw (A1)--(C2) (A2)--(B1) (C1)--(B3) (E1)--(D2) (F1)--(E2) (F2)--(D1);
    \foreach \x in {A,C,D,E,F}
    {\fill (\x) circle (2pt);}
    \node[below] at (current bounding box.south){Hình a};
    \path
    (E) node[above]{$T$}
    (F) node[above]{$V$}
    (D) node[above]{$N$};
    \end{tikzpicture}
}
\end{minipage}

\end{tcolorbox}
\vspace{4pt}


\begin{tcolorbox}[
    breakable,
    colback=white,
    colframe=blue!60!black,
    boxrule=0.8pt,
    arc=3pt,
    title={\textbf{Câu \#050} \quad \textbf{MÃ CÂU: \texttt{[L06\_C8\_B1\_01]}} \quad | \quad \textbf{Nhóm: \texttt{HINH\_045}} \quad | \quad \textbf{MapID: \texttt{6H8N1}} \hfill \textcolor{yellow!90!white}{[Hình gốc .pyc]}}
]
\textbf{Dạng toán:} Kể tên và vẽ các đường thẳng \hfill \textit{(Tự luận -- Nhận biết)}
\vspace{4pt}\hrule\vspace{6pt}
\begin{minipage}[t]{0.58\textwidth}
\vspace{0pt}
a) Kể tên các đường thẳng có trong hình a)?

b) Vẽ vào vở $3$ điểm như Hình b). Vẽ các đường thẳng đi qua hai trong $3$ điểm đó.
\begin{center}

    \begin{minipage}{0.5\textwidth}

\end{minipage}
\begin{minipage}{0.5\textwidth}
\begin{tikzpicture}[>=stealth,line join=round,line cap=round,font=\footnotesize,scale=.75]
\draw
(0,0) coordinate[label=above:$C$](N)
(3,1.5) coordinate[label=left:$B$](M)
(5,-.5) coordinate[label=below:$J$](P);
\foreach \x in {M,N,P}
{\fill (\x) circle (2pt);}

\path
(2,-1) node[below]{Hình b};
\end{tikzpicture}
\end{minipage}
    
\end{center}
\end{minipage}\hfill
\begin{minipage}[t]{0.39\textwidth}
\vspace{0pt}
\centering

\resizebox{\linewidth}{!}{%
\begin{tikzpicture}[>=stealth,line join=round,line cap=round,font=\footnotesize,scale=0.75]
\path
(5,2.5) coordinate (A)
(-1,-1.5) coordinate (B)
(-.6,.3) coordinate (C)
(6,-3) coordinate (D)
(6,-1.5) coordinate (E)
(-0.6,1.8) coordinate (F);
\draw (A)--(B);
\tkzLabelSegments[color=black,pos=.9](A,B){$d$}
\draw (C)--(D);
\tkzLabelSegments[color=black,pos=.9](C,D){$z$}
\draw (E)--(F);
\tkzLabelSegments[color=black,pos=.9](E,F){$a$}
\node[below] at (current bounding box.south){Hình a};
\end{tikzpicture}
}
\end{minipage}

\end{tcolorbox}
\vspace{4pt}


\begin{tcolorbox}[
    breakable,
    colback=white,
    colframe=blue!60!black,
    boxrule=0.8pt,
    arc=3pt,
    title={\textbf{Câu \#051} \quad \textbf{MÃ CÂU: \texttt{[L06\_C8\_B1\_03]}} \quad | \quad \textbf{Nhóm: \texttt{HINH\_046}} \quad | \quad \textbf{MapID: \texttt{6H8N1}} \hfill \textcolor{yellow!90!white}{[Hình gốc .pyc]}}
]
\textbf{Dạng toán:} Đặt tên, gọi tên điểm đường. \hfill \textit{(Tự luận -- Nhận biết)}
\vspace{4pt}\hrule\vspace{6pt}
\begin{minipage}[t]{0.58\textwidth}
\vspace{0pt}
a) Hãy đặt tên cho các điểm và đường thẳng trong Hình a.

 b) Hãy nêu 3 cách gọi tên đường thẳng trong Hình b
\begin{center}

    \begin{tikzpicture}[>=stealth,line join=round,line cap=round,font=\footnotesize,scale=1]
    \draw
    (0,0) coordinate(O)
    (1,0) coordinate[label=above:{$C$}](A)
    (5,0) coordinate[label=above:$M$](B)
    (7,0) coordinate[label=above:$N$](C)
    (8,0) coordinate[label=above:$B$](D)
    (9,0) coordinate(E);
    \draw (O)--(E); 
    \foreach \x in {A,B,C,D}
    {\fill (\x) circle (2pt);}
    \node[below] at (current bounding box.south){Hình b};
    \end{tikzpicture}

    
\end{center}
\end{minipage}\hfill
\begin{minipage}[t]{0.39\textwidth}
\vspace{0pt}
\centering

\resizebox{\linewidth}{!}{%
\begin{tikzpicture}[>=stealth,line join=round,line cap=round,font=\footnotesize,scale=1]
    \draw
    (-1,0) coordinate(A)
    (0,0) coordinate(O)
    (0,2) coordinate(B)
    (0,-2) coordinate(C)
    (5,0) coordinate(D)
    (1.5,-2) coordinate(E)
    (4,2) coordinate(F);
    \tkzInterLL(A,D)(E,F)\tkzGetPoint{G}
    \draw (B)--(C) (A)--(D) (E)--(F);
        \foreach \x in {O,G}
    {\fill (\x) circle (2pt);}
    \node[below] at (current bounding box.south){Hình a};
        \end{tikzpicture}
}
\end{minipage}

\end{tcolorbox}
\vspace{4pt}


\begin{tcolorbox}[
    breakable,
    colback=white,
    colframe=blue!60!black,
    boxrule=0.8pt,
    arc=3pt,
    title={\textbf{Câu \#052} \quad \textbf{MÃ CÂU: \texttt{[L06\_C8\_B1\_05]}} \quad | \quad \textbf{Nhóm: \texttt{HINH\_047}} \quad | \quad \textbf{MapID: \texttt{6H8N1}} \hfill \textcolor{yellow!90!white}{[Hình gốc .pyc]}}
]
\textbf{Dạng toán:} Kể tên và vẽ các đường thẳng \hfill \textit{(Tự luận -- Nhận biết)}
\vspace{4pt}\hrule\vspace{6pt}
\begin{minipage}[t]{0.58\textwidth}
\vspace{0pt}
Trong Hình vẽ:

Điểm ${G}$ thuộc những đường thẳng nào? 

Điểm ${F}$ không thuộc những đường thẳng nào?

Đường thẳng nào không chứa điểm ${O}$?

Sử dụng kí hiệu để mô tả các mối quan hệ trên
\end{minipage}\hfill
\begin{minipage}[t]{0.39\textwidth}
\vspace{0pt}
\centering

\resizebox{\linewidth}{!}{%
\begin{tikzpicture}[>=stealth,line join=round,line cap=round,font=\footnotesize,scale=0.6]
        \draw
        (-1,0) coordinate[label=above:$x$](A1)
        (0,0) coordinate[label=below:$F$](A)
        (6,0) coordinate[label=above right:$O$](C)
        (6.5,0) coordinate(C2)
        (3.5,0) coordinate(D)
        (4,5) coordinate[label=left:$G$](B);
        \tkzDefPointBy[homothety=center B ratio 1.2](A)\tkzGetPoint{A2}
        \tkzDefPointBy[homothety=center B ratio 1.2](C)\tkzGetPoint{C1}
        \tkzDefPointBy[homothety=center B ratio 1.2](D)\tkzGetPoint{D1}
        \tkzDefPointBy[homothety=center A ratio 1.2](B)\tkzGetPoint{B1}
        \tkzDefPointBy[homothety=center D ratio 1.2](B)\tkzGetPoint{B2}
        \tkzDefPointBy[homothety=center C ratio 1.2](B)\tkzGetPoint{B3}
        \draw (A1)--(C2) (A2)--(B1) (D1)--(B2) (C1)--(B3);
        \foreach \x in {A,B,C}
        {\fill (\x) circle (2pt);}
        \node[below] at (current bounding box.south){Hình a};
        \path
        (D1) node[left]{$z$}
        (B1) node[right]{$d$}
        (B3) node[left]{$c$};
        \end{tikzpicture}
}
\end{minipage}

\end{tcolorbox}
\vspace{4pt}


\begin{tcolorbox}[
    breakable,
    colback=white,
    colframe=blue!60!black,
    boxrule=0.8pt,
    arc=3pt,
    title={\textbf{Câu \#053} \quad \textbf{MÃ CÂU: \texttt{[L06\_C8\_B1\_07]}} \quad | \quad \textbf{Nhóm: \texttt{HINH\_048}} \quad | \quad \textbf{MapID: \texttt{6H8N1}} \hfill \textcolor{yellow!90!white}{[Hình gốc .pyc]}}
]
\textbf{Dạng toán:} Chỉ ra ba điểm thẳng hàng \hfill \textit{(Tự luận -- Nhận biết)}
\vspace{4pt}\hrule\vspace{6pt}
\begin{minipage}[t]{0.58\textwidth}
\vspace{0pt}
Trong Hình a, cho bốn điểm ${A}$, ${D}$, ${F}$, ${G}$ thuộc đường thẳng ${b}$ và điểm ${I}$ không thuộc đường thẳng ${b}$. Hãy nêu các bộ ba điểm thẳng hàng và bộ ba điểm không thẳng hàng
\end{minipage}\hfill
\begin{minipage}[t]{0.39\textwidth}
\vspace{0pt}
\centering

\resizebox{\linewidth}{!}{%
\begin{tikzpicture}[>=stealth,line join=round,line cap=round,font=\footnotesize,scale=1]
    \draw
    (0,0) coordinate(O)
    (1,0) coordinate[label=below:$A$](A)
    (1.5,0) coordinate[label=below:$D$](B)
    (2.5,0) coordinate[label=below:$F$](C)
    (4,0) coordinate[label=below:$G$](D)
    (3,0.5) coordinate[label=above:$I$](E)
    (4.5,0) coordinate[label=above:$b$](F);
    \draw (O)--(F); 
    \foreach \x in {A,B,C,D,E}
    {\fill (\x) circle (2pt);}
    \node[below] at (current bounding box.south){Hình a};
    \end{tikzpicture}
}
\end{minipage}

\end{tcolorbox}
\vspace{4pt}


\begin{tcolorbox}[
    breakable,
    colback=white,
    colframe=blue!60!black,
    boxrule=0.8pt,
    arc=3pt,
    title={\textbf{Câu \#054} \quad \textbf{MÃ CÂU: \texttt{[L06\_C8\_B1\_08]}} \quad | \quad \textbf{Nhóm: \texttt{HINH\_049}} \quad | \quad \textbf{MapID: \texttt{6H8N1}} \hfill \textcolor{yellow!90!white}{[Hình gốc .pyc]}}
]
\textbf{Dạng toán:} Chỉ ra điểm nằm giữa, không nằm giữa \hfill \textit{(Tự luận -- Nhận biết)}
\vspace{4pt}\hrule\vspace{6pt}
\begin{minipage}[t]{0.58\textwidth}
\vspace{0pt}
Trong Hình a, hãy chỉ ra các điểm:

a) Nằm giữa hai điểm ${A}$ và ${I}$.

b) Không nằm giữa hai điểm ${X}$ và ${V}$.
\end{minipage}\hfill
\begin{minipage}[t]{0.39\textwidth}
\vspace{0pt}
\centering

\resizebox{\linewidth}{!}{%
\begin{tikzpicture}[>=stealth,line join=round,line cap=round,font=\footnotesize,scale=1]
    \draw
    (0,0) coordinate(O)
    (1,0) coordinate[label=below:$A$](A)
    (2,0) coordinate[label=below:$X$](B)
    (3.6,0) coordinate[label=below:$L$](C)
    (5,0) coordinate[label=below:$V$](D)
    (7,0) coordinate[label=below:$I$](N)
    (7.5,0) coordinate(H);
    \draw (O)--(H); 
    \foreach \x in {A,B,C,D,N}
    {\fill (\x) circle (2pt);}
    \node[below] at (current bounding box.south){Hình a};
    \end{tikzpicture}
}
\end{minipage}

\end{tcolorbox}
\vspace{4pt}


\begin{tcolorbox}[
    breakable,
    colback=white,
    colframe=blue!60!black,
    boxrule=0.8pt,
    arc=3pt,
    title={\textbf{Câu \#055} \quad \textbf{MÃ CÂU: \texttt{[L06\_C8\_B1\_09]}} \quad | \quad \textbf{Nhóm: \texttt{HINH\_050}} \quad | \quad \textbf{MapID: \texttt{6H8N1}} \hfill \textcolor{yellow!90!white}{[Hình gốc .pyc]}}
]
\textbf{Dạng toán:} Nhận biết các điểm thẳng hang \hfill \textit{(Tự luận -- Nhận biết)}
\vspace{4pt}\hrule\vspace{6pt}
\begin{minipage}[t]{0.58\textwidth}
\vspace{0pt}
Xem Hình a và trả lời:

 Có bao nhiêu bộ ba điểm thẳng hàng?

a) Hãy nêu ít nhất hai bộ ba điểm không thẳng hàng?

b) Bốn điểm ${F}$, ${M}$, ${N}$, ${R}$ có thẳng hàng hay không?.
\end{minipage}\hfill
\begin{minipage}[t]{0.39\textwidth}
\vspace{0pt}
\centering

\resizebox{\linewidth}{!}{%
\begin{tikzpicture}[>=stealth,line join=round,line cap=round,font=\footnotesize,scale=0.75]
    \draw
    (0,2) coordinate[label=above:$y$](D1)
    (0,0) coordinate[label=below:$m$](D2)
    (0,-1.8) coordinate[label=below:$d$](D3)
    (3,0) coordinate[label=below:$R$](S)
    (4.5,2) coordinate(E1)
    (5.5,-2) coordinate(E2);
    \tkzDefPointBy[homothety=center D1 ratio 2](S)\tkzGetPoint{D4}
    \tkzDefPointBy[homothety=center D2 ratio 2](S)\tkzGetPoint{D5}
    \tkzDefPointBy[homothety=center D3 ratio 2](S)\tkzGetPoint{D6}
    \draw (E1)--(E2) (D1)--(D4) (D2)--(D5) (D3)--(D6);
    \tkzInterLL(E1,E2)(S,D4)\tkzGetPoint{C}
    \tkzInterLL(E1,E2)(S,D5)\tkzGetPoint{B}
    \tkzInterLL(E1,E2)(S,D6)\tkzGetPoint{A}
    \foreach \x in {A,B,C,S}
    {\fill (\x) circle (2pt);}
    \node[below] at (current bounding box.south){Hình a};
    \path
    (A) node[left]{$F$}
    (B) node[above right]{$M$}
    (C) node[right]{$N$};
    \end{tikzpicture}
}
\end{minipage}

\end{tcolorbox}
\vspace{4pt}


\begin{tcolorbox}[
    breakable,
    colback=white,
    colframe=blue!60!black,
    boxrule=0.8pt,
    arc=3pt,
    title={\textbf{Câu \#056} \quad \textbf{MÃ CÂU: \texttt{[L06\_C8\_B1\_10]}} \quad | \quad \textbf{Nhóm: \texttt{HINH\_051}} \quad | \quad \textbf{MapID: \texttt{6H8V1}} \hfill \textcolor{yellow!90!white}{[Hình gốc .pyc]}}
]
\textbf{Dạng toán:} Vận dụng tìm điểm đường chưa biết \hfill \textit{(Tự luận -- Vận dụng)}
\vspace{4pt}\hrule\vspace{6pt}
\begin{minipage}[t]{0.58\textwidth}
\vspace{0pt}
Hình a mô tả $4$ đường thẳng và $5$ điểm có tên là ${C}$, ${F}$, ${K}$, ${P}$ và ${W}$, trong đó ta chỉ biết vị trí của điểm ${C}$. Hãy điền tên của các điểm còn lại, biết rằng:

a) Điểm ${P}$ nằm trên $3$ trong $4$ đường thẳng.

b) Ba điểm ${C}$, ${F}$, ${K}$ thẳng hàng.

c) Ba điểm ${F}$, ${P}$, ${W}$ thẳng hàng.
\end{minipage}\hfill
\begin{minipage}[t]{0.39\textwidth}
\vspace{0pt}
\centering

\resizebox{\linewidth}{!}{%
\begin{tikzpicture}[>=stealth,line join=round,line cap=round,font=\footnotesize,scale=0.6]
        \draw
        (-1,0) coordinate(A1)
        (0,0) coordinate[label=below:$C$](A)
        (6,0) coordinate(C)
        (6.5,0) coordinate(C2)
        (3.5,0) coordinate(D)
        (4,5) coordinate(B);
        \tkzDefPointBy[homothety=center B ratio 1.2](A)\tkzGetPoint{A2}
        \tkzDefPointBy[homothety=center B ratio 1.2](C)\tkzGetPoint{C1}
        \tkzDefPointBy[homothety=center B ratio 1.2](D)\tkzGetPoint{D1}
        \tkzDefPointBy[homothety=center B ratio 0.6](D)\tkzGetPoint{D2}
        \tkzDefPointBy[homothety=center A ratio 1.2](B)\tkzGetPoint{B1}
        \tkzDefPointBy[homothety=center D ratio 1.2](B)\tkzGetPoint{B2}
        \tkzDefPointBy[homothety=center C ratio 1.2](B)\tkzGetPoint{B3}
        \draw (A1)--(C2) (A2)--(B1) (D1)--(B2) (C1)--(B3);
        \foreach \x in {A,B,C,D,D2}
        {\fill (\x) circle (2pt);}
        \node[below] at (current bounding box.south){Hình a};
        \end{tikzpicture}
}
\end{minipage}

\end{tcolorbox}
\vspace{4pt}


\begin{tcolorbox}[
    breakable,
    colback=white,
    colframe=blue!60!black,
    boxrule=0.8pt,
    arc=3pt,
    title={\textbf{Câu \#057} \quad \textbf{MÃ CÂU: \texttt{[L06\_C8\_B1\_11]}} \quad | \quad \textbf{Nhóm: \texttt{HINH\_052}} \quad | \quad \textbf{MapID: \texttt{6H8V1}} \hfill \textcolor{yellow!90!white}{[Hình gốc .pyc]}}
]
\textbf{Dạng toán:} Vận dụng vẽ thêm điểm thoả mãn yêu cầu \hfill \textit{(Tự luận -- Vận dụng)}
\vspace{4pt}\hrule\vspace{6pt}
\begin{minipage}[t]{0.58\textwidth}
\vspace{0pt}
a) Cho hai điểm ${U}$ và ${N}$. Hãy vẽ thêm điểm ${V}$ sao cho ba điểm ${U}$, ${V}$, ${N}$ thẳng hàng và điểm ${U}$, ${N}$ nằm cùng phía đối với điểm ${V}$.

b) Vẽ hai điểm trên $1$ tờ giấy trắng. Không dùng thước, em hãy tìm cách vẽ thêm một điểm thứ ba sao cho 3 điểm đó thẳng hàng.
\end{minipage}\hfill
\begin{minipage}[t]{0.39\textwidth}
\vspace{0pt}
\centering

\resizebox{\linewidth}{!}{%
\begin{tikzpicture}[>=stealth,line join=round,line cap=round,font=\footnotesize,scale=0.6]
\draw
(0,0) coordinate(O)
(1,0) coordinate[label=above:$U$](A)
(2.5,0) coordinate[label=above:$N$](M)
(4,0) coordinate[label=above:$V$](N)
(4.5,0) coordinate(D);
\draw (O)--(D);
\foreach \x in {A,M,N}
{\fill (\x) circle (2pt);}
    \end{tikzpicture}
}
\end{minipage}

\end{tcolorbox}
\vspace{4pt}


\begin{tcolorbox}[
    breakable,
    colback=white,
    colframe=blue!60!black,
    boxrule=0.8pt,
    arc=3pt,
    title={\textbf{Câu \#058} \quad \textbf{MÃ CÂU: \texttt{[L06\_C8\_B1\_12]}} \quad | \quad \textbf{Nhóm: \texttt{HINH\_053}} \quad | \quad \textbf{MapID: \texttt{6H8H1}} \hfill \textcolor{yellow!90!white}{[Hình gốc .pyc]}}
]
\textbf{Dạng toán:} Quan hệ giữa 3 điểm thẳng hàng \hfill \textit{(Tự luận -- Thông hiểu)}
\vspace{4pt}\hrule\vspace{6pt}
\begin{minipage}[t]{0.58\textwidth}
\vspace{0pt}
Nhìn Hình a, em hãy trả lời các câu hỏi sau:

a) Những điểm nào thẳng hàng?

b) Điểm nào nằm giữa hai điểm?

c) Hai điểm nào nằm cùng phía đối với điểm thứ ba?

d) Điểm nào không thẳng hàng với hai điểm ${P}$ và ${R}$?
\end{minipage}\hfill
\begin{minipage}[t]{0.39\textwidth}
\vspace{0pt}
\centering

\resizebox{\linewidth}{!}{%
\begin{tikzpicture}[>=stealth,line join=round,line cap=round,font=\footnotesize,scale=0.75]
   
    \draw
    (0,0) coordinate[label=below right:$P$](E)
    (-0.6,0) coordinate[label=above left:$H$](B)
    (3,0) coordinate[label=above:$R$](H)
    (0,1) coordinate[label=right:$K$](C)
    (0,3) coordinate[label=left:$F$](A);
    \tkzInterLL(B,C)(A,H)\tkzGetPoint{D}
    \tkzDefPointBy[homothety=center B ratio 1.2](D)\tkzGetPoint{D1}
    \tkzDefPointBy[homothety=center D ratio 1.2](B)\tkzGetPoint{B1}
    \tkzDefPointBy[homothety=center B ratio 1.2](A)\tkzGetPoint{A1}
    \tkzDefPointBy[homothety=center A ratio 1.2](B)\tkzGetPoint{B2}
    \draw (B)+(-0.4,0)--(H)--++(0.4,0) (A)+(0,0.4)--(E)--++(0,-0.4) (H)+(0.4,-0.4)--(A)--++(-0.4,0.4) (B1)--(D1) (A1)--(B2);    
    \foreach \x in {E,B,H,C,A,D}
    {\fill (\x) circle (2pt);}
    \path
    (D) node[right]{$N$};
    \node[below] at (current bounding box.south){Hình a};
    
    \end{tikzpicture}
}
\end{minipage}

\end{tcolorbox}
\vspace{4pt}


\begin{tcolorbox}[
    breakable,
    colback=white,
    colframe=blue!60!black,
    boxrule=0.8pt,
    arc=3pt,
    title={\textbf{Câu \#059} \quad \textbf{MÃ CÂU: \texttt{[L06\_C8\_B1\_13]}} \quad | \quad \textbf{Nhóm: \texttt{HINH\_054}} \quad | \quad \textbf{MapID: \texttt{6H8N1}} \hfill \textcolor{yellow!90!white}{[Hình gốc .pyc]}}
]
\textbf{Dạng toán:} Vẽ đường thẳng đi qua 2 điểm \hfill \textit{(Tự luận -- Nhận biết)}
\vspace{4pt}\hrule\vspace{6pt}
\begin{minipage}[t]{0.58\textwidth}
\vspace{0pt}
Đánh dấu ba điểm phân biệt ${H}$, ${I}$ và ${Y}$ trên một tờ giấy trắng sao cho chúng không thẳng hàng.

a) Hãy vẽ các đường thẳng đi qua hai trong ba điểm ấy. Đó là những đường thẳng nào?

b) Hãy chỉ ra hai đường thẳng cắt nhau và giao điểm của chúng.
\end{minipage}\hfill
\begin{minipage}[t]{0.39\textwidth}
\vspace{0pt}
\centering

\resizebox{\linewidth}{!}{%
\begin{tikzpicture}[>=stealth,line join=round,line cap=round,font=\footnotesize,scale=0.5]
        \draw
        (-1,0) coordinate(A1)
        (0,0) coordinate[label=below:$H$](A)
        (6,0) coordinate[label=above right:$Y$](C)
        (6.5,0) coordinate(C2)
        (4,5) coordinate[label=right:$I$](B);
        \tkzDefPointBy[homothety=center B ratio 1.2](A)\tkzGetPoint{A2}
        \tkzDefPointBy[homothety=center B ratio 1.2](C)\tkzGetPoint{C1}
        \tkzDefPointBy[homothety=center A ratio 1.2](B)\tkzGetPoint{B1}
        \tkzDefPointBy[homothety=center C ratio 1.2](B)\tkzGetPoint{B2}
        \tkzDefPointBy[homothety=center C ratio 1.2](B)\tkzGetPoint{B3}
        \draw (A1)--(C2) (A2)--(B1) (C1)--(B3);
        \foreach \x in {A,B,C}
        {\fill (\x) circle (2pt);}
        \end{tikzpicture}
}
\end{minipage}

\end{tcolorbox}
\vspace{4pt}


\begin{tcolorbox}[
    breakable,
    colback=white,
    colframe=blue!60!black,
    boxrule=0.8pt,
    arc=3pt,
    title={\textbf{Câu \#060} \quad \textbf{MÃ CÂU: \texttt{[L06\_C8\_B1\_15]}} \quad | \quad \textbf{Nhóm: \texttt{HINH\_055}} \quad | \quad \textbf{MapID: \texttt{6H8N1}} \hfill \textcolor{yellow!90!white}{[Hình gốc .pyc]}}
]
\textbf{Dạng toán:} Đếm số giao điểm \hfill \textit{(Tự luận -- Nhận biết)}
\vspace{4pt}\hrule\vspace{6pt}
\begin{minipage}[t]{0.58\textwidth}
\vspace{0pt}
Đếm số giao điểm tạo bởi ba đường thẳng trong mỗi Hình
\begin{center}

    \begin{minipage}{0.25\textwidth}

\end{minipage}
\begin{minipage}{0.25\textwidth}
    \begin{tikzpicture}[>=stealth,line join=round,line cap=round,font=\footnotesize,scale=1]
    \draw (0,0)--(4,0) (0,1)--(4,1) (1.5,-0.5)--(3,2);
    \path 
    (2,-0.6) node[below]{b)};
    \end{tikzpicture}
\end{minipage}
\begin{minipage}{0.25\textwidth}
    \begin{tikzpicture}[>=stealth,line join=round,line cap=round,font=\footnotesize,scale=1]
    \draw (0,-1.2)--(0,1.2) (1.2,1.2)--(-1.2,-1.2) (1.2,-1.2)--(-1.2,1.2);
    \path 
    (0,-1.4) node[below]{c)};
    \end{tikzpicture}
\end{minipage}
\begin{minipage}{0.25\textwidth}
    \begin{tikzpicture}[>=stealth,line join=round,line cap=round,font=\footnotesize,scale=1]
    \draw (0,0)--(4,0) (0,-0.6)--(2,1.6) (1.6,1.6)--(3.6,-0.6);
    \path 
    (2,-0.6) node[below]{d)};
    \end{tikzpicture}
\end{minipage}

    
\end{center}
\end{minipage}\hfill
\begin{minipage}[t]{0.39\textwidth}
\vspace{0pt}
\centering

\resizebox{\linewidth}{!}{%
\begin{tikzpicture}[>=stealth,line join=round,line cap=round,font=\footnotesize,scale=1]
\draw (0,0)--(4,0) (0,1)--(4,1) (0,2)--(4,2);
\path 
(2,-0.6) node[below]{a)};
\end{tikzpicture}
}
\end{minipage}

\end{tcolorbox}
\vspace{4pt}


\begin{tcolorbox}[
    breakable,
    colback=white,
    colframe=blue!60!black,
    boxrule=0.8pt,
    arc=3pt,
    title={\textbf{Câu \#061} \quad \textbf{MÃ CÂU: \texttt{[L06\_C8\_B1\_16]}} \quad | \quad \textbf{Nhóm: \texttt{HINH\_056}} \quad | \quad \textbf{MapID: \texttt{6H8N1}} \hfill \textcolor{yellow!90!white}{[Hình gốc .pyc]}}
]
\textbf{Dạng toán:} Liệt kê các đường thẳng song song \hfill \textit{(Tự luận -- Nhận biết)}
\vspace{4pt}\hrule\vspace{6pt}
\begin{minipage}[t]{0.58\textwidth}
\vspace{0pt}
Hãy liệt kê các đường thẳng song song trong Hình
\end{minipage}\hfill
\begin{minipage}[t]{0.39\textwidth}
\vspace{0pt}
\centering

\resizebox{\linewidth}{!}{%
\begin{tikzpicture}[>=stealth,line join=round,line cap=round,font=\footnotesize,scale=0.6]
    \draw
    (-1,0) coordinate(A1)
    (0,0) coordinate[label=above:$E$](A)
    (6,0) coordinate[label=above right:$J$](C)
    (6.5,0) coordinate(C2)
    (3,4) coordinate[label=above:$D$](B);
    \tkzDefMidPoint(A,B)\tkzGetPoint{E}
    \tkzDefMidPoint(A,C)\tkzGetPoint{D}
    \tkzDefMidPoint(C,B)\tkzGetPoint{F}
    \tkzDefPointBy[homothety=center B ratio 1.2](A)\tkzGetPoint{A2}
    \tkzDefPointBy[homothety=center B ratio 1.2](C)\tkzGetPoint{C1}
    \tkzDefPointBy[homothety=center A ratio 1.2](B)\tkzGetPoint{B1}
    \tkzDefPointBy[homothety=center C ratio 1.2](B)\tkzGetPoint{B3}
    \tkzDefPointBy[homothety=center C ratio 1.2](B)\tkzGetPoint{B3}
    \tkzDefPointBy[homothety=center D ratio 1.2](E)\tkzGetPoint{E1}
    \tkzDefPointBy[homothety=center E ratio 1.2](D)\tkzGetPoint{D2}
    \tkzDefPointBy[homothety=center F ratio 1.2](E)\tkzGetPoint{E2}
    \tkzDefPointBy[homothety=center E ratio 1.2](F)\tkzGetPoint{F1}
    \tkzDefPointBy[homothety=center D ratio 1.2](F)\tkzGetPoint{F2}
    \tkzDefPointBy[homothety=center F ratio 1.2](D)\tkzGetPoint{D1}
    \draw (A1)--(C2) (A2)--(B1) (C1)--(B3) (E1)--(D2) (F1)--(E2) (F2)--(D1);
    \foreach \x in {A,C,D,E,F}
    {\fill (\x) circle (2pt);}
    \node[below] at (current bounding box.south){Hình a};
    \path
    (E) node[above]{$Q$}
    (F) node[above]{$U$}
    (D) node[above]{$N$};
    \end{tikzpicture}
}
\end{minipage}

\end{tcolorbox}
\vspace{4pt}


\begin{tcolorbox}[
    breakable,
    colback=white,
    colframe=blue!60!black,
    boxrule=0.8pt,
    arc=3pt,
    title={\textbf{Câu \#062} \quad \textbf{MÃ CÂU: \texttt{[L06\_C8\_B2\_11]}} \quad | \quad \textbf{Nhóm: \texttt{HINH\_057}} \quad | \quad \textbf{MapID: \texttt{6H8N2}} \hfill \textcolor{yellow!90!white}{[Hình gốc .pyc]}}
]
\textbf{Dạng toán:} Nhận biết hai tia đối nhau, trùng nhau \hfill \textit{(Trắc nghiệm -- Nhận biết)}
\vspace{4pt}\hrule\vspace{6pt}
\begin{minipage}[t]{0.58\textwidth}
\vspace{0pt}
Quan sát hình vẽ chọn khẳng định đúng.

\vspace{4pt}\noindent
\textbf{A.} Tia ${Lx}$ và tia ${Ly}$ là hai tia trùng nhau. \quad \textbf{\textcolor{blue!75!black}{B. (Đ)}} Tia ${Ex}$ và tia ${Ey}$ là hai tia đối nhau \quad \textbf{C.} Tia ${Ex}$ và tia ${Ey}$ là hai tia trùng nhau \quad \textbf{D.} Tia ${IE}$ và tia ${Ix}$ là hai tia đối nhau
\end{minipage}\hfill
\begin{minipage}[t]{0.39\textwidth}
\vspace{0pt}
\centering

\resizebox{\linewidth}{!}{%
\begin{tikzpicture}[>=stealth,line join=round,line cap=round,font=\footnotesize,scale=1]
                \coordinate[label=above:{$N$}] (O) at (3,2);
                \coordinate[label=above:{$x$}] (x) at (0,0);
                \coordinate[label=above:{$y$}] (y) at (6,0);
                \coordinate[label=below:{$E$}] (A) at (1,0);
                \coordinate[label=below right:{$I$}] (B) at (2,0);
                \coordinate[label=below left:{$L$}] (C) at (5,0);

                \draw[thick] (x)--(y);
                \draw[shorten >=-1cm] (O)--(A);
                \draw[shorten >=-1cm] (O)--(B);
                \draw[shorten >=-1cm] (O)--(C);
                \foreach \p in {A,B,C,O} 
                    \fill (\p) circle (1.3pt);
            \end{tikzpicture}
}
\end{minipage}

\end{tcolorbox}
\vspace{4pt}


\begin{tcolorbox}[
    breakable,
    colback=white,
    colframe=blue!60!black,
    boxrule=0.8pt,
    arc=3pt,
    title={\textbf{Câu \#063} \quad \textbf{MÃ CÂU: \texttt{[L06\_C8\_B2\_12]}} \quad | \quad \textbf{Nhóm: \texttt{HINH\_058}} \quad | \quad \textbf{MapID: \texttt{6H8N2}} \hfill \textcolor{yellow!90!white}{[Hình gốc .pyc]}}
]
\textbf{Dạng toán:} Tìm tia đối \hfill \textit{(Trắc nghiệm -- Nhận biết)}
\vspace{4pt}\hrule\vspace{6pt}
\begin{minipage}[t]{0.58\textwidth}
\vspace{0pt}
Quan sát hình vẽ. Tìm tia đối của tia ${Qx}$.

\vspace{4pt}\noindent
\textbf{A.} ${By}$ \quad \textbf{B.} ${QB}$ \quad \textbf{\textcolor{blue!75!black}{C. (Đ)}} ${Qy}$ \quad \textbf{D.} ${My}$
\end{minipage}\hfill
\begin{minipage}[t]{0.39\textwidth}
\vspace{0pt}
\centering

\resizebox{\linewidth}{!}{%
\begin{tikzpicture}[>=stealth,line join=round,line cap=round,font=\footnotesize,scale=1]
                \coordinate[label=above:{$x$}] (x) at (-3,0);
                \coordinate[label=above:{$y$}] (y) at (3,0);
                \coordinate[label=above:{$Q$}] (O) at (0,0);
                \coordinate[label=above:{$B$}] (A) at (-1,0);
                \coordinate[label=above:{$M$}] (B) at (2,0);

                \draw[thick] (x)--(y);
                \foreach \p in {A,B,O}
                    \fill (\p) circle (1.3pt);  
            \end{tikzpicture}
}
\end{minipage}

\end{tcolorbox}
\vspace{4pt}


\begin{tcolorbox}[
    breakable,
    colback=white,
    colframe=blue!60!black,
    boxrule=0.8pt,
    arc=3pt,
    title={\textbf{Câu \#064} \quad \textbf{MÃ CÂU: \texttt{[L06\_C8\_B2\_13]}} \quad | \quad \textbf{Nhóm: \texttt{HINH\_059}} \quad | \quad \textbf{MapID: \texttt{6H8N2}} \hfill \textcolor{yellow!90!white}{[Hình gốc .pyc]}}
]
\textbf{Dạng toán:} Nhận biết vị trí của điểm \hfill \textit{(Trắc nghiệm -- Nhận biết)}
\vspace{4pt}\hrule\vspace{6pt}
\begin{minipage}[t]{0.58\textwidth}
\vspace{0pt}
Quan sát hình vẽ, em hãy chọn khẳng định đúng.

\vspace{4pt}\noindent
\textbf{A.} ${G}$ nằm giữa ${Q}$ và ${U}$ \quad \textbf{\textcolor{blue!75!black}{B. (Đ)}} ${Q}$ nằm giữa ${G}$ và ${L}$ \quad \textbf{C.} ${G}$ nằm giữa ${Q}$ và ${L}$ \quad \textbf{D.} ${L}$ nằm giữa ${Q}$ và ${U}$
\end{minipage}\hfill
\begin{minipage}[t]{0.39\textwidth}
\vspace{0pt}
\centering

\resizebox{\linewidth}{!}{%
\begin{tikzpicture}[>=stealth,line join=round,line cap=round,font=\footnotesize,scale=1]
                \coordinate[label=above:{$G$}] (A) at (0,0);
                \coordinate[label=above:{$L$}] (B) at (8,0);
                \coordinate[label=above:{$Q$}] (C) at (3,0);
                \coordinate[label=above:{$U$}] (D) at (5.5,0);

                \draw[thick] (-1,0)--(8.5,0);
                \foreach \p in {A,B,C,D}
                    \fill (\p) circle (1.3pt);
            \end{tikzpicture}
}
\end{minipage}

\end{tcolorbox}
\vspace{4pt}


\begin{tcolorbox}[
    breakable,
    colback=white,
    colframe=blue!60!black,
    boxrule=0.8pt,
    arc=3pt,
    title={\textbf{Câu \#065} \quad \textbf{MÃ CÂU: \texttt{[L06\_C8\_B2\_14]}} \quad | \quad \textbf{Nhóm: \texttt{HINH\_060}} \quad | \quad \textbf{MapID: \texttt{6H8N2}} \hfill \textcolor{yellow!90!white}{[Hình gốc .pyc]}}
]
\textbf{Dạng toán:} Nhận biết điểm thuộc tia \hfill \textit{(Trắc nghiệm -- Nhận biết)}
\vspace{4pt}\hrule\vspace{6pt}
\begin{minipage}[t]{0.58\textwidth}
\vspace{0pt}
Quan sát hình vẽ, em hãy chọn khẳng định đúng.

\vspace{4pt}\noindent
\textbf{\textcolor{blue!75!black}{A. (Đ)}} ${M}$ thuộc tia ${Wx}$ \quad \textbf{B.} ${M}$ thuộc tia ${Ny}$ \quad \textbf{C.} ${M}$ không thuộc tia ${Wx}$ \quad \textbf{D.} ${M}$ thuộc tia ${Wy}$
\end{minipage}\hfill
\begin{minipage}[t]{0.39\textwidth}
\vspace{0pt}
\centering

\resizebox{\linewidth}{!}{%
\begin{tikzpicture}[>=stealth,line join=round,line cap=round,font=\footnotesize,scale=1]
                \coordinate[label=above:{$x$}] (x) at (-3,0);
                \coordinate[label=above:{$y$}] (y) at (3,0);
                \coordinate[label=above:{$W$}] (O) at (0,0);
                \coordinate[label=above:{$M$}] (A) at (-1,0);
                \coordinate[label=above:{$N$}] (B) at (2,0);

                \draw[thick] (x)--(y);
                \foreach \p in {A,B,O}
                    \fill (\p) circle (1.3pt);  
            \end{tikzpicture}
}
\end{minipage}

\end{tcolorbox}
\vspace{4pt}


\begin{tcolorbox}[
    breakable,
    colback=white,
    colframe=blue!60!black,
    boxrule=0.8pt,
    arc=3pt,
    title={\textbf{Câu \#066} \quad \textbf{MÃ CÂU: \texttt{[L06\_C8\_B2\_02]}} \quad | \quad \textbf{Nhóm: \texttt{HINH\_061}} \quad | \quad \textbf{MapID: \texttt{6H8N2}} \hfill \textcolor{yellow!90!white}{[Hình gốc .pyc]}}
]
\textbf{Dạng toán:} Vẽ, kể tên các tia \hfill \textit{(Tự luận -- Nhận biết)}
\vspace{4pt}\hrule\vspace{6pt}
\begin{minipage}[t]{0.58\textwidth}
\vspace{0pt}
Cho đường thẳng ${xy}$. Lấy điểm ${C}$, ${K}$, ${T}$ thuộc ${xy}$ theo thứ tự đó. Điểm ${Y}$ không thuộc đường thẳng ${xy}$.

a) Vẽ các tia ${YC}$, ${YK}$, ${YT}$.

b) Kể tên các tia đối nhau trong hình vẽ.

c) Kể tên các tia trùng nhau trong hình vẽ.

d) Tia ${Cx}$ và ${Ky}$ có phải là hai tia đối nhau hai không?
\end{minipage}\hfill
\begin{minipage}[t]{0.39\textwidth}
\vspace{0pt}
\centering

\resizebox{\linewidth}{!}{%
\begin{tikzpicture}[>=stealth,line join=round,line cap=round,font=\footnotesize,scale=1]
                \coordinate[label=above:{$Y$}] (O) at (3,2);
                \coordinate[label=above:{$x$}] (x) at (0,0);
                \coordinate[label=above:{$y$}] (y) at (6,0);
                \coordinate[label=below:{$C$}] (A) at (1,0);
                \coordinate[label=below right:{$K$}] (B) at (2,0);
                \coordinate[label=below left:{$T$}] (C) at (5,0);

                \draw[thick] (x)--(y);
                \draw[shorten >=-1cm] (O)--(A);
                \draw[shorten >=-1cm] (O)--(B);
                \draw[shorten >=-1cm] (O)--(C);
                \foreach \p in {A,B,C,O} 
                    \fill (\p) circle (1.3pt);
            \end{tikzpicture}
}
\end{minipage}

\end{tcolorbox}
\vspace{4pt}


\begin{tcolorbox}[
    breakable,
    colback=white,
    colframe=blue!60!black,
    boxrule=0.8pt,
    arc=3pt,
    title={\textbf{Câu \#067} \quad \textbf{MÃ CÂU: \texttt{[L06\_C8\_B2\_03]}} \quad | \quad \textbf{Nhóm: \texttt{HINH\_062}} \quad | \quad \textbf{MapID: \texttt{6H8N2}} \hfill \textcolor{yellow!90!white}{[Hình gốc .pyc]}}
]
\textbf{Dạng toán:} Vẽ, kể tên các tia đối nhau \hfill \textit{(Tự luận -- Nhận biết)}
\vspace{4pt}\hrule\vspace{6pt}
\begin{minipage}[t]{0.58\textwidth}
\vspace{0pt}
Vẽ hai tia đối nhau ${Ox}$ và ${Oy}$. Lấy ${C\in Ox}$ và ${D\in Oy}$.

a) Viết tên các tia trùng với tia ${Cy}$.

b) Hai tia ${CD}$ và ${Oy}$ có đối nhau không? Vì sao?

c) Viết tên các tia đối của tia ${Cy}$.
\end{minipage}\hfill
\begin{minipage}[t]{0.39\textwidth}
\vspace{0pt}
\centering

\resizebox{\linewidth}{!}{%
\begin{tikzpicture}[>=stealth,line join=round,line cap=round,font=\footnotesize,scale=1]
                \coordinate[label=above:{$x$}] (x) at (-3,0);
                \coordinate[label=above:{$y$}] (y) at (3,0);
                \coordinate[label=above:{$O$}] (O) at (0,0);
                \coordinate[label=above:{$C$}] (A) at (-1,0);
                \coordinate[label=above:{$D$}] (B) at (2,0);

                \draw[thick] (x)--(y);
                \foreach \p in {A,B,O}
                    \fill (\p) circle (1.3pt);  
            \end{tikzpicture}
}
\end{minipage}

\end{tcolorbox}
\vspace{4pt}


\begin{tcolorbox}[
    breakable,
    colback=white,
    colframe=blue!60!black,
    boxrule=0.8pt,
    arc=3pt,
    title={\textbf{Câu \#068} \quad \textbf{MÃ CÂU: \texttt{[L06\_C8\_B2\_04]}} \quad | \quad \textbf{Nhóm: \texttt{HINH\_063}} \quad | \quad \textbf{MapID: \texttt{6H8N2}} \hfill \textcolor{yellow!90!white}{[Hình gốc .pyc]}}
]
\textbf{Dạng toán:} Vẽ tia theo yêu cầu \hfill \textit{(Tự luận -- Nhận biết)}
\vspace{4pt}\hrule\vspace{6pt}
\begin{minipage}[t]{0.58\textwidth}
\vspace{0pt}
a) Lấy $3$ điểm ${O}$, ${P}$, ${Q}$ không thẳng hàng.

a) Vẽ hai tia ${OP}$ và ${OQ}$.

b) Vẽ tia ${Ox}$ cắt đường thẳng ${PQ}$ tại điểm ${W}$ nằm giữa hai điểm ${P}$ và ${Q}$.
\end{minipage}\hfill
\begin{minipage}[t]{0.39\textwidth}
\vspace{0pt}
\centering

\resizebox{\linewidth}{!}{%
\begin{tikzpicture}[>=stealth,line join=round,line cap=round,font=\footnotesize,scale=1]
                \coordinate[label=below left:{$O$}] (A) at (0,0);
                \coordinate[label=below:{$P$}] (B) at (7,0);
                \coordinate[label=above:{$Q$}] (C) at (3,3.5);
                \coordinate[label=above right:{$x$}] (x) at (6,2);
                \coordinate[label=above:{$W$}] (I) at (intersection of A--x and B--C);
                \foreach \p in {A,B,C,I}
                    \fill (\p) circle (1.4pt);
                \draw[shorten >=-1cm] (A)--(B);
                \draw[shorten >=-1cm] (A)--(C);
                \draw[shorten >=-1cm,shorten <=-1cm] (B)--(C);
                \draw (A)--(x);
            \end{tikzpicture}
}
\end{minipage}

\end{tcolorbox}
\vspace{4pt}


\begin{tcolorbox}[
    breakable,
    colback=white,
    colframe=blue!60!black,
    boxrule=0.8pt,
    arc=3pt,
    title={\textbf{Câu \#069} \quad \textbf{MÃ CÂU: \texttt{[L06\_C8\_B2\_05]}} \quad | \quad \textbf{Nhóm: \texttt{HINH\_064}} \quad | \quad \textbf{MapID: \texttt{6H8H2}} \hfill \textcolor{yellow!90!white}{[Hình gốc .pyc]}}
]
\textbf{Dạng toán:} Vẽ tia theo yêu cầu \hfill \textit{(Tự luận -- Thông hiểu)}
\vspace{4pt}\hrule\vspace{6pt}
\begin{minipage}[t]{0.58\textwidth}
\vspace{0pt}
Học sinh vẽ lại hình bên dưới.

a) Hãy viết tên ba điểm thẳng hàng và nói rõ điểm nào nằm giữa hai điểm còn lại.

b) Hãy viết tên hai tia đối nhau.

c) Hãy vẽ tia ${TO}$ thêm vào hình đó.

d) Vẽ tia gốc ${U}$ sao cho tia đó cắt tia ${TO}$.
\end{minipage}\hfill
\begin{minipage}[t]{0.39\textwidth}
\vspace{0pt}
\centering

\resizebox{\linewidth}{!}{%
\begin{tikzpicture}[>=stealth,line join=round,line cap=round,font=\footnotesize,scale=0.6]
            \coordinate[label=below:{$K$}] (K) at (0,0);
        \coordinate[label=below:{$O$}] (M) at (2,0);
        \coordinate[label=below:{$R$}] (A) at (5,0);
        \coordinate[label=above:{$T$}] (S) at (-1,1);
        \coordinate[label=above:{$U$}] (T) at (4,2);
        \draw[thick] (-1.5,0)--(6,0);
        \foreach \p in {S,K,M,T,A}
            \fill (\p) circle (1.3pt);
            \end{tikzpicture}
}
\end{minipage}

\end{tcolorbox}
\vspace{4pt}


\begin{tcolorbox}[
    breakable,
    colback=white,
    colframe=blue!60!black,
    boxrule=0.8pt,
    arc=3pt,
    title={\textbf{Câu \#070} \quad \textbf{MÃ CÂU: \texttt{[L06\_C8\_B2\_06]}} \quad | \quad \textbf{Nhóm: \texttt{HINH\_065}} \quad | \quad \textbf{MapID: \texttt{6H8V2}} \hfill \textcolor{yellow!90!white}{[Hình gốc .pyc]}}
]
\textbf{Dạng toán:} Vẽ tia theo yêu cầu \hfill \textit{(Tự luận -- Vận dụng)}
\vspace{4pt}\hrule\vspace{6pt}
\begin{minipage}[t]{0.58\textwidth}
\vspace{0pt}
Vẽ hai tia ${Fx}$, ${Fy}$ đối nhau. Lấy điểm ${I}$ thuộc tia ${Fx}$, điểm ${S}$ thuộc tia ${Fy}$.

a) Kể tên hai tia đối nhau gốc ${I}$, hai tia trùng nhau gốc ${S}$.

b) Hai tia ${Fx}$ và ${Sy}$ có trùng nhau không? Vì sao?

c) Hai tia ${Ix}$ và ${Fy}$ có đối nhau không? Vì sao 

d) Vẽ đường thẳng ${d}$ cắt đường thẳng ${xy}$ tại một điểm nằm giữa hai điểm ${F}$ và ${S}$.
\end{minipage}\hfill
\begin{minipage}[t]{0.39\textwidth}
\vspace{0pt}
\centering

\resizebox{\linewidth}{!}{%
\begin{tikzpicture}[>=stealth,line join=round,line cap=round,font=\footnotesize,scale=1]
                \coordinate[label=above:{$F$}] (A) at (0,0);
                \coordinate[label=above:{$x$}] (x) at (-3,0);
                \coordinate[label=above:{$y$}] (y) at (3,0);
                \coordinate[label=above:{$I$}] (B) at (-1,0);
                \coordinate[label=above:{$S$}] (C) at (2,0);
                \draw[thick] (x)--(y);
                \foreach \p in {A,B,C}
                    \fill (\p) circle (1.3pt);
            \end{tikzpicture}
}
\end{minipage}

\end{tcolorbox}
\vspace{4pt}


\begin{tcolorbox}[
    breakable,
    colback=white,
    colframe=blue!60!black,
    boxrule=0.8pt,
    arc=3pt,
    title={\textbf{Câu \#071} \quad \textbf{MÃ CÂU: \texttt{[L06\_C8\_B2\_07]}} \quad | \quad \textbf{Nhóm: \texttt{HINH\_066}} \quad | \quad \textbf{MapID: \texttt{6H8V2}} \hfill \textcolor{yellow!90!white}{[Hình gốc .pyc]}}
]
\textbf{Dạng toán:} Viết tên các tia đối, tia trùng nhau. \hfill \textit{(Tự luận -- Vận dụng)}
\vspace{4pt}\hrule\vspace{6pt}
\begin{minipage}[t]{0.58\textwidth}
\vspace{0pt}
Vẽ đường thẳng ${BC}$. Lấy hai điểm ${G}$, ${S}$ trên đường thẳng $AB$ theo thứ tự ${B}$, ${G}$, ${S}$, ${C}$.

a) Viết tên các tia đối nhau.

b) Viết tên các tia trùng nhau.

c) Lấy điểm ${Y}$ nằm ngoài đường thẳng ${BC}$. Vẽ tia ${YB}$, đường thẳng ${YC}$.
\end{minipage}\hfill
\begin{minipage}[t]{0.39\textwidth}
\vspace{0pt}
\centering

\resizebox{\linewidth}{!}{%
\begin{tikzpicture}[>=stealth,line join=round,line cap=round,font=\footnotesize,scale=1]
                \coordinate[label=above:{$B$}] (A) at (0,0);
                \coordinate[label=above:{$C$}] (B) at (8,0);
                \coordinate[label=above:{$G$}] (C) at (3,0);
                \coordinate[label=above:{$S$}] (D) at (5.5,0);

                \draw[thick] (-1,0)--(8.5,0);
                \foreach \p in {A,B,C,D}
                    \fill (\p) circle (1.3pt);
            \end{tikzpicture}
}
\end{minipage}

\end{tcolorbox}
\vspace{4pt}


\begin{tcolorbox}[
    breakable,
    colback=white,
    colframe=blue!60!black,
    boxrule=0.8pt,
    arc=3pt,
    title={\textbf{Câu \#072} \quad \textbf{MÃ CÂU: \texttt{[L06\_C8\_B2\_08]}} \quad | \quad \textbf{Nhóm: \texttt{HINH\_067}} \quad | \quad \textbf{MapID: \texttt{6H8V2}} \hfill \textcolor{yellow!90!white}{[Hình gốc .pyc]}}
]
\textbf{Dạng toán:} Vẽ hình theo yêu cầu \hfill \textit{(Tự luận -- Vận dụng)}
\vspace{4pt}\hrule\vspace{6pt}
\begin{minipage}[t]{0.58\textwidth}
\vspace{0pt}
Lấy ba điểm ${J}$, ${P}$, ${Q}$ không thẳng hàng.

a) Vẽ tia ${Jx}$ cắt đường thẳng ${PQ}$ tại điểm ${R}$ nằm giữa ${P}$ và ${Q}$.

b) Vẽ tia ${Jy}$ cắt tia ${QP}$ tại điểm ${V}$ không nằm giữa ${Q}$ và ${P}$.

c) Trong hình vẽ trên, những điểm nào là điểm nằm giữa hai điểm khác?
\end{minipage}\hfill
\begin{minipage}[t]{0.39\textwidth}
\vspace{0pt}
\centering

\resizebox{\linewidth}{!}{%
\begin{tikzpicture}[>=stealth,line join=round,line cap=round,font=\footnotesize,scale=1]
              cap=round,font=\footnotesize,scale=1,thick]
                \coordinate[label=above:{$J$}] (A) at (3,3.5);
                \coordinate[label=below:{$P$}] (B) at (0,0);
                \coordinate[label=below:{$Q$}] (C) at (5,0);

                \coordinate[label=above right:{$x$}] (x) at (4,-1);
                \coordinate[label=below left:{$R$}] (M) at (intersection of A--x and B--C);
                \coordinate[label=left:{$y$}] (y) at (-3,-1);
                \coordinate[label=below:{$V$}] (N) at (intersection of A--y and B--C);

                % \draw[shorten >=-1cm] (A)--(B);
                % \draw[shorten >=-1cm] (A)--(C);
                \draw[shorten >=-1.2cm,shorten <=-1cm] (C)--(N);
                \draw[blue] (A)--(x);
                \draw[red,shorten >=-1.2cm] (A)--(N);
                \foreach \p in {A,B,C,M,N}
                    \fill (\p) circle (1.4pt);
            \end{tikzpicture}
}
\end{minipage}

\end{tcolorbox}
\vspace{4pt}


\begin{tcolorbox}[
    breakable,
    colback=white,
    colframe=blue!60!black,
    boxrule=0.8pt,
    arc=3pt,
    title={\textbf{Câu \#073} \quad \textbf{MÃ CÂU: \texttt{[L06\_C8\_B2\_09]}} \quad | \quad \textbf{Nhóm: \texttt{HINH\_068}} \quad | \quad \textbf{MapID: \texttt{6H8V2}} \hfill \textcolor{yellow!90!white}{[Hình gốc .pyc]}}
]
\textbf{Dạng toán:} Viết hình theo diễn đạt \hfill \textit{(Tự luận -- Vận dụng)}
\vspace{4pt}\hrule\vspace{6pt}
\begin{minipage}[t]{0.58\textwidth}
\vspace{0pt}
Trên đường thẳng ${xy}$, lấy điểm ${C}$ và ${D}$ (phân biệt). Qua ${D}$ vẽ đường thẳng ${pq}$ và qua ${C}$ vẽ đường thẳng ${mn}$ sao cho ${pq}$ cắt ${mn}$ tại ${J}$.

a) Vẽ hình theo diễn đạt trên.

b) Trong hình vẽ có tất cả bao nhiêu tia?

c) Qua ${D}$ vẽ đường thẳng ${uv}$ cắt ${CJ}$ tại điểm ${Y}$ nằm giữa ${C}$ và ${J}$. Kể tên các tia đối nhau, các tia trùng nhau trên đường thẳng ${mn}$.

d) Hai tia ${JY}$ và ${Cm}$ có trùng nhau không? Giải thích.
\end{minipage}\hfill
\begin{minipage}[t]{0.39\textwidth}
\vspace{0pt}
\centering

\resizebox{\linewidth}{!}{%
\begin{tikzpicture}[>=stealth,line join=round,line cap=round,font=\footnotesize,scale=0.6]
                
                    \coordinate[label=above:{$J$}] (C) at (3,3);
                    \coordinate[label=below:{$C$}] (A) at (0,0);
                    \coordinate[label=below left:{$D$}] (B) at (6,0);

                    \coordinate[label=above:{$x$}] (x) at (-1.5,0);
                    \coordinate[label=above:{$y$}] (y) at (7,0);

                    \coordinate[label=above left:{$n$}] (n) at ($(C)!1.3!(A)$);
                    \coordinate[label=above:{$m$}] (m) at ($(A)!1.3!(C)$);
                    \coordinate[label=above:{$p$}] (p) at ($(B)!1.3!(C)$);
                    \coordinate[label=below:{$q$}] (q) at ($(C)!1.3!(B)$);

                    \draw[thick] (m)--(n) (x)--(y) (p)--(q);
                    \foreach \p in {A,B,C}
                        \fill (\p) circle (1.4pt);
            \end{tikzpicture}
}
\end{minipage}

\end{tcolorbox}
\vspace{4pt}


\begin{tcolorbox}[
    breakable,
    colback=white,
    colframe=blue!60!black,
    boxrule=0.8pt,
    arc=3pt,
    title={\textbf{Câu \#074} \quad \textbf{MÃ CÂU: \texttt{[L06\_C8\_B3\_17]}} \quad | \quad \textbf{Nhóm: \texttt{HINH\_069}} \quad | \quad \textbf{MapID: \texttt{6H8N3}} \hfill \textcolor{yellow!90!white}{[Hình gốc .pyc]}}
]
\textbf{Dạng toán:} Kể tên các đoạn thẳng \hfill \textit{(Trắc nghiệm -- Nhận biết)}
\vspace{4pt}\hrule\vspace{6pt}
\begin{minipage}[t]{0.58\textwidth}
\vspace{0pt}
Em hãy kể tên các đoạn thẳng có đầu mút là hai trong ba điểm ${B}$, ${N}$, ${O}$, ${U}$.

\vspace{4pt}\noindent
\textbf{\textcolor{blue!75!black}{A. (Đ)}} ${BN}$, ${BO}$, ${BU}$, ${NO}$, ${NU}$ và ${OU}$ \quad \textbf{B.} ${BN}$, ${BO}$, ${BU}$, ${NO}$, ${NU}$ \quad \textbf{C.} ${NO}$, ${BO}$ \quad \textbf{D.} ${BN}$, ${NO}$
\end{minipage}\hfill
\begin{minipage}[t]{0.39\textwidth}
\vspace{0pt}
\centering

\resizebox{\linewidth}{!}{%
\begin{tikzpicture}[>=stealth,line join=round,line cap=round,font=\footnotesize,scale=1]
    \path
    (0,0) coordinate (A)
    (2,-0.5) coordinate (B)
    (4,0) coordinate (C)
    (1,-1.5) coordinate (D)
    ;
    \filldraw[black] (A) circle (1pt) node[above right] {$B$};
        \filldraw[black] (B) circle (1pt) node[above left] {$N$};
    \filldraw[black] (C) circle (1pt) node[above right] {$O$};
        \filldraw[black] (D) circle (1pt) node[above left] {$U$};

    \end{tikzpicture}
}
\end{minipage}

\end{tcolorbox}
\vspace{4pt}


\begin{tcolorbox}[
    breakable,
    colback=white,
    colframe=blue!60!black,
    boxrule=0.8pt,
    arc=3pt,
    title={\textbf{Câu \#075} \quad \textbf{MÃ CÂU: \texttt{[L06\_C8\_B3\_01]}} \quad | \quad \textbf{Nhóm: \texttt{HINH\_070}} \quad | \quad \textbf{MapID: \texttt{6H8N3}} \hfill \textcolor{yellow!90!white}{[Hình gốc .pyc]}}
]
\textbf{Dạng toán:} Kể tên các đoạn thẳng \hfill \textit{(Tự luận -- Nhận biết)}
\vspace{4pt}\hrule\vspace{6pt}
\begin{minipage}[t]{0.58\textwidth}
\vspace{0pt}
a) Hãy đọc tên các đoạn thẳng có trong hình 1.

a) Với 4 điểm ${B}$, ${F}$, ${J}$, ${K}$ như hình 2, em hãy kể tên các đoạn thẳng có đầu mút là 

b) Hai trong ba điểm ${B}$, ${F}$, ${J}$.

c) Hai trong ba điểm ${B}$, ${F}$, ${J}$, ${K}$.
\end{minipage}\hfill
\begin{minipage}[t]{0.39\textwidth}
\vspace{0pt}
\centering

\resizebox{\linewidth}{!}{%
\begin{tikzpicture}[>=stealth,line join=round,line cap=round,font=\footnotesize,scale=1]
          \path
(2,3) coordinate (A)
(0,0) coordinate (B)
(5,0) coordinate (C)
;
\draw (A)--(B)--(C)--cycle;
\filldraw[black] (A) circle (1pt) node[above right] {$B$};
    \filldraw[black] (B) circle (1pt) node[above left] {$F$};
\filldraw[black] (C) circle (1pt) node[above right] {$J$};


\end{tikzpicture}
\begin{tikzpicture}[>=stealth,line join=round,line cap=round,font=\footnotesize,scale=1]
\path
(0,0) coordinate (A)
(2,-0.5) coordinate (B)
(4,0) coordinate (C)
(1,-1.5) coordinate (D)
;

\filldraw[black] (A) circle (1pt) node[above right] {$B$};
    \filldraw[black] (B) circle (1pt) node[above left] {$F$};
\filldraw[black] (C) circle (1pt) node[above right] {$J$};
    \filldraw[black] (D) circle (1pt) node[above left] {$K$};

            \end{tikzpicture}
}
\end{minipage}

\end{tcolorbox}
\vspace{4pt}


\begin{tcolorbox}[
    breakable,
    colback=white,
    colframe=blue!60!black,
    boxrule=0.8pt,
    arc=3pt,
    title={\textbf{Câu \#076} \quad \textbf{MÃ CÂU: \texttt{[L06\_C8\_B3\_02]}} \quad | \quad \textbf{Nhóm: \texttt{HINH\_071}} \quad | \quad \textbf{MapID: \texttt{6H8N3}} \hfill \textcolor{yellow!90!white}{[Hình gốc .pyc]}}
]
\textbf{Dạng toán:} Kể tên các đoạn thẳng 2 \hfill \textit{(Tự luận -- Nhận biết)}
\vspace{4pt}\hrule\vspace{6pt}
\begin{minipage}[t]{0.58\textwidth}
\vspace{0pt}
Nhìn hình bên và cho biết

a) Hãy kể tên các đường thẳng đi qua 2 trong 4 điểm ${G}$, ${Q}$, ${J}$, ${K}$.

b) Hai tia ${QG}$ và ${QJ}$ có phải là 2 tia đối nhau không? Vì sao?

c) Cho biết có tất cả bao nhiêu đoạn thẳng?
\end{minipage}\hfill
\begin{minipage}[t]{0.39\textwidth}
\vspace{0pt}
\centering

\resizebox{\linewidth}{!}{%
\begin{tikzpicture}[>=stealth,line join=round,line cap=round,font=\footnotesize,scale=1]
\path
(0,0) coordinate (O)
(2,0) coordinate (B)
(1,2) coordinate (C)
(-2.5,0) coordinate (A)
;
\draw (-3,0)node[above]{$x$}--(3,0)node[above]{$y$}
(A)--(C)--(B) (O)--($(O)!1.3!(C)$)node[above left]{$z$};
\filldraw[black] (A) circle (1pt) node[above right] {$G$};
    \filldraw[black] (B) circle (1pt) node[above left] {$J$};
\filldraw[black] (C) circle (1pt) node[above right] {$K$};
    \filldraw[black] (O) circle (1pt) node[above left] {$Q$};
            \end{tikzpicture}
}
\end{minipage}

\end{tcolorbox}
\vspace{4pt}


\begin{tcolorbox}[
    breakable,
    colback=white,
    colframe=blue!60!black,
    boxrule=0.8pt,
    arc=3pt,
    title={\textbf{Câu \#077} \quad \textbf{MÃ CÂU: \texttt{[L06\_C8\_B3\_03]}} \quad | \quad \textbf{Nhóm: \texttt{HINH\_072}} \quad | \quad \textbf{MapID: \texttt{6H8N3}} \hfill \textcolor{yellow!90!white}{[Hình gốc .pyc]}}
]
\textbf{Dạng toán:} Kể tên các đoạn thẳng, tia chung gốc, tia đối nhau \hfill \textit{(Tự luận -- Nhận biết)}
\vspace{4pt}\hrule\vspace{6pt}
\begin{minipage}[t]{0.58\textwidth}
\vspace{0pt}
Cho ba điểm ${A}$, ${K}$, ${N}$ không thẳng hàng. Vẽ tia ${AK}$, đoạn thẳng ${AN}$, đường thẳng ${KN}$. Vẽ tia ${Ax}$ cắt đường thẳng ${KN}$ tại điểm ${W}$ nằm giữa hai điểm ${K}$ và ${N}$.

a) Kể tên các đoạn thẳng có trong hình.

b) Kể tên hai tia gốc ${A}$ trùng nhau.

c) Kể tên hai tia gốc ${W}$ đối nhau.
\end{minipage}\hfill
\begin{minipage}[t]{0.39\textwidth}
\vspace{0pt}
\centering

\resizebox{\linewidth}{!}{%
\begin{tikzpicture}[>=stealth,line join=round,line cap=round,font=\footnotesize,scale=1]
\path
(0,0) coordinate (A)
(2,0) coordinate (M)
(1,2) coordinate (P)
($(M)!0.3!(P)$) coordinate (K)
;
\draw (A)--($(A)!1.3!(M)$) (A)--(P) ($(M)!1.2!(P)$)--($(P)!1.2!(M)$)
(A)--($(A)!1.6!(K)$)node[above]{$x$};
    \filldraw[black] (A) circle (1pt) node[above right] {$A$};
    \filldraw[black] (M) circle (1pt) node[above left] {$K$};
    \filldraw[black] (P) circle (1pt) node[above] {$N$};
   \filldraw[black] (K) circle (1pt) node[above right]{$W$};
               \end{tikzpicture}
}
\end{minipage}

\end{tcolorbox}
\vspace{4pt}


\begin{tcolorbox}[
    breakable,
    colback=white,
    colframe=blue!60!black,
    boxrule=0.8pt,
    arc=3pt,
    title={\textbf{Câu \#078} \quad \textbf{MÃ CÂU: \texttt{[L06\_C8\_B3\_04]}} \quad | \quad \textbf{Nhóm: \texttt{HINH\_073}} \quad | \quad \textbf{MapID: \texttt{6H8H3}} \hfill \textcolor{yellow!90!white}{[Hình gốc .pyc]}}
]
\textbf{Dạng toán:} Vẽ tia đối \hfill \textit{(Tự luận -- Thông hiểu)}
\vspace{4pt}\hrule\vspace{6pt}
\begin{minipage}[t]{0.58\textwidth}
\vspace{0pt}
Vẽ hai tia ${Yx}$ và ${Yy}$ đối nhau.

 Lấy điểm ${F}$ thuộc tia ${Yx}$ và điểm ${H}$ thuộc tia ${Yy}$. Viết tên các tia trùng với tia ${Fy}$.

 Hai tia ${Fx}$ và ${Hy}$ có đối nhau không? Vì sao?

 Lấy điểm ${O}$ nằm ngoài đường thẳng ${xy}$. Vẽ đoạn thẳng ${FO}$ và tia ${HO}$.
\end{minipage}\hfill
\begin{minipage}[t]{0.39\textwidth}
\vspace{0pt}
\centering

\resizebox{\linewidth}{!}{%
\begin{tikzpicture}[>=stealth,line join=round,line cap=round,font=\footnotesize,scale=1]
\path
(0,0) coordinate (O)
(-2,0) coordinate (A)
(3,0) coordinate (B)
;
\draw (-4,0)node[above]{$x$}--(5,0)node[above]{$y$};
    \filldraw[black] (A) circle (1pt) node[above right] {$F$};
    \filldraw[black] (B) circle (1pt) node[above left] {$H$};
 
\end{tikzpicture}
\begin{tikzpicture}[>=stealth,line join=round,line cap=round,font=\footnotesize,scale=1]
\path
(0,0) coordinate (O)
(-2,0) coordinate (A)
(3,0) coordinate (B)
(1,2) coordinate (C)
;
\draw (-4,0)node[above]{$x$}--(5,0)node[above]{$y$} (A)--(C) (B)--($(B)!1.3!(C)$);
  \filldraw[black] (A) circle (1pt) node[above right] {$F$};
    \filldraw[black] (B) circle (1pt) node[above left] {$H$};
    \filldraw[black] (C) circle (1pt) node[above] {$O$};

               \end{tikzpicture}
}
\end{minipage}

\end{tcolorbox}
\vspace{4pt}


\begin{tcolorbox}[
    breakable,
    colback=white,
    colframe=blue!60!black,
    boxrule=0.8pt,
    arc=3pt,
    title={\textbf{Câu \#079} \quad \textbf{MÃ CÂU: \texttt{[L06\_C8\_B3\_05]}} \quad | \quad \textbf{Nhóm: \texttt{HINH\_074}} \quad | \quad \textbf{MapID: \texttt{6H8H3}} \hfill \textcolor{yellow!90!white}{[Hình gốc .pyc]}}
]
\textbf{Dạng toán:} Kể, vẽ các đoạn thẳng, tia chung gốc, tia đối nhau \hfill \textit{(Tự luận -- Thông hiểu)}
\vspace{4pt}\hrule\vspace{6pt}
\begin{minipage}[t]{0.58\textwidth}
\vspace{0pt}
Vẽ ba điểm ${A}$, ${K}$, ${P}$ thuộc đường thẳng ${xy}$ sao cho điểm ${K}$ nằm giữa hai điểm ${A}$ và ${P}$.

Kể tên hai tia đối nhau gốc ${K}$.

 Kể tên hai tia trùng với tia ${AP}$.

 Vẽ đường thẳng ${b}$ qua ${A}$ sao cho ${b}$ và ${xy}$ là hai đường thẳng không trùng nhau. Vẽ hai điểm ${T}$, ${U}$ thuộc đường thẳng ${b}$ sao cho điểm ${A}$ nằm giữa hai điểm ${T}$ và ${U}$.

Kể tên các đoạn thẳng có trong hình.
\end{minipage}\hfill
\begin{minipage}[t]{0.39\textwidth}
\vspace{0pt}
\centering

\resizebox{\linewidth}{!}{%
\begin{tikzpicture}[>=stealth,line join=round,line cap=round,font=\footnotesize,scale=1]
\path
(0,0) coordinate (A)
(2,0) coordinate (B)
(5,0) coordinate (C)
(3,1.5) coordinate[label=above:$b$] (m)
($(A)!0.6!(m)$) coordinate (D)
($(D)!2!(A)$) coordinate (I)
;
\draw (-2,0)node[above]{$x$}--(7,0)node[above]{$y$} ($(m)!1.2!(I)$)--(m);
  \filldraw[black] (A) circle (1pt) node[above right] {$A$};
    \filldraw[black] (B) circle (1pt) node[above left] {$K$};
    \filldraw[black] (C) circle (1pt) node[above] {$P$};
 \filldraw[black] (I) circle (1pt) node[above left] {$U$};
    \filldraw[black] (D) circle (1pt) node[above] {$T$};

               \end{tikzpicture}
}
\end{minipage}

\end{tcolorbox}
\vspace{4pt}


\begin{tcolorbox}[
    breakable,
    colback=white,
    colframe=blue!60!black,
    boxrule=0.8pt,
    arc=3pt,
    title={\textbf{Câu \#080} \quad \textbf{MÃ CÂU: \texttt{[L06\_C8\_B3\_06]}} \quad | \quad \textbf{Nhóm: \texttt{HINH\_075}} \quad | \quad \textbf{MapID: \texttt{6H8H3}} \hfill \textcolor{yellow!90!white}{[Hình gốc .pyc]}}
]
\textbf{Dạng toán:} Vẽ tia đối, đếm đoạn thẳng \hfill \textit{(Tự luận -- Thông hiểu)}
\vspace{4pt}\hrule\vspace{6pt}
\begin{minipage}[t]{0.58\textwidth}
\vspace{0pt}
Vẽ hai tia ${Ax}$ và ${Ay}$ đối nhau. Lấy điểm ${I}$ thuộc tia ${Ax}$ và điểm ${K}$ thuộc tia ${Ay}$.

Trên hình có bao nhiêu đoạn thẳng? Kể tên?

 Tia ${KI}$ trùng với những tia nào?

Tia ${Ix}$ đối với những tia nào?
\end{minipage}\hfill
\begin{minipage}[t]{0.39\textwidth}
\vspace{0pt}
\centering

\resizebox{\linewidth}{!}{%
\begin{tikzpicture}[>=stealth,line join=round,line cap=round,font=\footnotesize,scale=1]
\path
(0,0) coordinate (A)
(-2,0) coordinate (B)
(3,0) coordinate (C)
;
\draw (-4,0)node[above]{$x$}--(4,0)node[above]{$y$};

  \filldraw[black] (A) circle (1pt) node[above right] {$A$};
    \filldraw[black] (B) circle (1pt) node[above left] {$I$};
    \filldraw[black] (C) circle (1pt) node[above] {$K$};

               \end{tikzpicture}
}
\end{minipage}

\end{tcolorbox}
\vspace{4pt}


\begin{tcolorbox}[
    breakable,
    colback=white,
    colframe=blue!60!black,
    boxrule=0.8pt,
    arc=3pt,
    title={\textbf{Câu \#081} \quad \textbf{MÃ CÂU: \texttt{[L06\_C8\_B3\_07]}} \quad | \quad \textbf{Nhóm: \texttt{HINH\_076}} \quad | \quad \textbf{MapID: \texttt{6H8V3}} \hfill \textcolor{yellow!90!white}{[Hình gốc .pyc]}}
]
\textbf{Dạng toán:} Vẽ hình theo cách diễn đạt \hfill \textit{(Tự luận -- Vận dụng)}
\vspace{4pt}\hrule\vspace{6pt}
\begin{minipage}[t]{0.58\textwidth}
\vspace{0pt}
Vẽ đường thẳng ${JM}$, lấy điểm ${K}$ nằm giữa hai điểm ${J}$ và ${M}$, điểm ${Q}$ nằm ngoài đường thẳng ${JM}$.

Vẽ tia ${QJ}$, đoạn thẳng ${QK}$ và đường thẳng ${QM}$.

Vẽ tia ${Kx}$ cắt đoạn thẳng ${QJ}$ tại ${U}$ nằm giữa ${Q}$ và ${J}$.

Vẽ tia ${Ky}$ là tia đối của tia ${Kx}$.

Dựa vào hình hãy trả lời các câu hỏi sau

Kể tên hai tia trùng nhau gốc ${J}$.

Kể tên hai tia đối nhau gốc ${K}$.
\end{minipage}\hfill
\begin{minipage}[t]{0.39\textwidth}
\vspace{0pt}
\centering

\resizebox{\linewidth}{!}{%
\begin{tikzpicture}[>=stealth,line join=round,line cap=round,font=\footnotesize,scale=1]
\path
(0,0) coordinate (C)
(5,0) coordinate (D)
(3,0) coordinate (N)
(2.5,2) coordinate (O)
($(N)+(150:4)$) coordinate[label=above:$x$] (x)
(intersection of N--x and O--C) coordinate (K)
($(x)!1.5!(N)$) coordinate[label=below:$y$] (y)
;
\draw (N)--(O)--($(O)!1.3!(C)$) ($(D)!1.2!(O)$)--($(O)!1.2!(D)$) (x)--(y);
\draw (-2,0)--(7,0);
  \filldraw[black] (C) circle (1pt) node[above right] {$J$};
    \filldraw[black] (N) circle (1pt) node[above left] {$K$};
    \filldraw[black] (D) circle (1pt) node[above] {$M$};
\filldraw[black] (O) circle (1pt) node[above left] {$Q$};
    \filldraw[black] (K) circle (1pt) node[above] {$U$};
    \end{tikzpicture}
}
\end{minipage}

\end{tcolorbox}
\vspace{4pt}


\begin{tcolorbox}[
    breakable,
    colback=white,
    colframe=blue!60!black,
    boxrule=0.8pt,
    arc=3pt,
    title={\textbf{Câu \#082} \quad \textbf{MÃ CÂU: \texttt{[L06\_C8\_B3\_08]}} \quad | \quad \textbf{Nhóm: \texttt{HINH\_077}} \quad | \quad \textbf{MapID: \texttt{6H8V3}} \hfill \textcolor{yellow!90!white}{[Hình gốc .pyc]}}
]
\textbf{Dạng toán:} Vẽ hình theo cách diễn đạt 2 \hfill \textit{(Tự luận -- Vận dụng)}
\vspace{4pt}\hrule\vspace{6pt}
\begin{minipage}[t]{0.58\textwidth}
\vspace{0pt}
Vẽ hình theo cách diễn đạt bằng lời sau (Vẽ trên cùng một hình).

Cho ba điểm ${A}$, ${E}$, ${H}$ không thẳng hàng, vẽ đoạn thẳng ${AE}$.

Vẽ điểm ${R}$ thuộc đoạn thẳng ${AE}$.

Vẽ tia ${AH}$.

Vẽ đường thẳng ${HP}$ sao cho điểm ${R}$ nằm giữa hai điểm ${H}$ và ${P}$.
\end{minipage}\hfill
\begin{minipage}[t]{0.39\textwidth}
\vspace{0pt}
\centering

\resizebox{\linewidth}{!}{%
\begin{tikzpicture}[>=stealth,line join=round,line cap=round,font=\footnotesize,scale=1]
\path
(0,0) coordinate (A)
(5,0) coordinate (B)
(3,2) coordinate (C)
($(A)!0.4!(B)$) coordinate (M)
($(C)!1.6!(M)$) coordinate (D)
;
\draw (A)--(B);
\draw (A)--($(A)!1.3!(C)$);
\draw ($(C)!1.2!(D)$)--($(D)!1.2!(C)$);
  \filldraw[black] (A) circle (1pt) node[above right] {$A$};
    \filldraw[black] (B) circle (1pt) node[above left] {$E$};
    \filldraw[black] (C) circle (1pt) node[above] {$H$};
\filldraw[black] (M) circle (1pt) node[above left] {$R$};
    \filldraw[black] (D) circle (1pt) node[above] {$P$};
    \end{tikzpicture}
}
\end{minipage}

\end{tcolorbox}
\vspace{4pt}


\begin{tcolorbox}[
    breakable,
    colback=white,
    colframe=blue!60!black,
    boxrule=0.8pt,
    arc=3pt,
    title={\textbf{Câu \#083} \quad \textbf{MÃ CÂU: \texttt{[L06\_C8\_B3\_09]}} \quad | \quad \textbf{Nhóm: \texttt{HINH\_078}} \quad | \quad \textbf{MapID: \texttt{6H8V3}} \hfill \textcolor{yellow!90!white}{[Hình gốc .pyc]}}
]
\textbf{Dạng toán:} Vẽ hình, kể tên các tia chung gốc, tia đối \hfill \textit{(Tự luận -- Vận dụng)}
\vspace{4pt}\hrule\vspace{6pt}
\begin{minipage}[t]{0.58\textwidth}
\vspace{0pt}
Trên cùng một đường thẳng, lấy điểm ${L}$ nằm giữa $2$ điểm ${J}$ và ${O}$, điểm ${R}$ nằm giữa $2$ điểm ${J}$ và ${L}$, điểm ${W}$ nằm giữa $2$ điểm ${L}$ và ${O}$.

a) Kể tên các tia trùng nhau gốc ${W}$.

b) Kể tên hai tia đối nhau gốc ${W}$.
\end{minipage}\hfill
\begin{minipage}[t]{0.39\textwidth}
\vspace{0pt}
\centering

\resizebox{\linewidth}{!}{%
\begin{tikzpicture}[>=stealth,line join=round,line cap=round,font=\footnotesize,scale=1]
\foreach \x/\diem in {0/A,1/M,3/B,4.5/N,7/C}{\path (\x,0) coordinate (\diem);}
\draw (-2,0)--++(0:12);

  \filldraw[black] (A) circle (1pt) node[above right] {$J$};
    \filldraw[black] (B) circle (1pt) node[above left] {$L$};
    \filldraw[black] (C) circle (1pt) node[above] {$O$};
\filldraw[black] (M) circle (1pt) node[above left] {$R$};
    \filldraw[black] (N) circle (1pt) node[above] {$W$};
    \end{tikzpicture}
}
\end{minipage}

\end{tcolorbox}
\vspace{4pt}


\begin{tcolorbox}[
    breakable,
    colback=white,
    colframe=blue!60!black,
    boxrule=0.8pt,
    arc=3pt,
    title={\textbf{Câu \#084} \quad \textbf{MÃ CÂU: \texttt{[L06\_C8\_B3\_10]}} \quad | \quad \textbf{Nhóm: \texttt{HINH\_079}} \quad | \quad \textbf{MapID: \texttt{6H8V3}} \hfill \textcolor{yellow!90!white}{[Hình gốc .pyc]}}
]
\textbf{Dạng toán:} Vẽ hình, kể tên các đoạn thẳng \hfill \textit{(Tự luận -- Vận dụng)}
\vspace{4pt}\hrule\vspace{6pt}
\begin{minipage}[t]{0.58\textwidth}
\vspace{0pt}
Vẽ đường thẳng ${a}$. Lấy $3$ điểm ${F}$, ${H}$, ${Q}$ thuộc đường thẳng ${a}$, điểm ${R}$ không thuộc đường thẳng ${a}$. Kẻ các đường thẳng qua các cặp điểm.

Có tất cả bao nhiêu đường thẳng? Kể tên?

${R}$ là giao điểm của những đường thẳng nào?.
\end{minipage}\hfill
\begin{minipage}[t]{0.39\textwidth}
\vspace{0pt}
\centering

\resizebox{\linewidth}{!}{%
\begin{tikzpicture}[>=stealth,line join=round,line cap=round,font=\footnotesize,scale=1]

\foreach \x/\diem in {0/A,2/B,5/C}{\path (\x,0) coordinate (\diem);}
\path (4,2) coordinate (D);
\foreach \x/\y in {D/A,D/B,D/C}{\draw ($(\x)!1.2!(\y)$) -- ($(\y)!1.2!(\x)$);}
\draw (-1,0)--(6,0)node[above right]{$a$};
  \filldraw[black] (A) circle (1pt) node[above right] {$F$};
    \filldraw[black] (B) circle (1pt) node[above left] {$H$};
    \filldraw[black] (C) circle (1pt) node[above] {$Q$};
\filldraw[black] (D) circle (1pt) node[above left] {$R$};
   
    \end{tikzpicture}
}
\end{minipage}

\end{tcolorbox}
\vspace{4pt}


\begin{tcolorbox}[
    breakable,
    colback=white,
    colframe=blue!60!black,
    boxrule=0.8pt,
    arc=3pt,
    title={\textbf{Câu \#085} \quad \textbf{MÃ CÂU: \texttt{[L06\_C8\_B3\_11]}} \quad | \quad \textbf{Nhóm: \texttt{HINH\_080}} \quad | \quad \textbf{MapID: \texttt{6H8N3}} \hfill \textcolor{yellow!90!white}{[Hình gốc .pyc]}}
]
\textbf{Dạng toán:} Đo độ dài đoạn thẳng và sắp xếp thứ tự \hfill \textit{(Tự luận -- Nhận biết)}
\vspace{4pt}\hrule\vspace{6pt}
\begin{minipage}[t]{0.58\textwidth}
\vspace{0pt}
Cho hình vẽ bên. Hãy đo độ dài các đoạn thẳng ${AB}$, ${BD}$, ${DO}$, ${OS}$, ${SY}$, ${AY}$ và sắp xếp các đoạn thẳng đó theo thứ tự độ dài từ bé đến lớn..
\end{minipage}\hfill
\begin{minipage}[t]{0.39\textwidth}
\vspace{0pt}
\centering

\resizebox{\linewidth}{!}{%
\begin{tikzpicture}[>=stealth,line join=round,line cap=round,font=\footnotesize,scale=1]

\draw (0,0)node[below left]{$B$}--++(0:4)node[below right]{$D$}--++(90:2)node[above right]{$O$}--++(180:3)node[above right]{$S$}--++(90:1)node[above]{$Y$}--++(180:1)node[above]{$A$}--cycle;

  
    \end{tikzpicture}
}
\end{minipage}

\end{tcolorbox}
\vspace{4pt}


\begin{tcolorbox}[
    breakable,
    colback=white,
    colframe=blue!60!black,
    boxrule=0.8pt,
    arc=3pt,
    title={\textbf{Câu \#086} \quad \textbf{MÃ CÂU: \texttt{[L06\_C8\_B3\_12]}} \quad | \quad \textbf{Nhóm: \texttt{HINH\_081}} \quad | \quad \textbf{MapID: \texttt{6H8N3}} \hfill \textcolor{yellow!90!white}{[Hình gốc .pyc]}}
]
\textbf{Dạng toán:} Đo độ dài đoạn thẳng và vẽ \hfill \textit{(Tự luận -- Nhận biết)}
\vspace{4pt}\hrule\vspace{6pt}
\begin{minipage}[t]{0.58\textwidth}
\vspace{0pt}
Cho hai đoạn thẳng ${DF}$ và ${KM}$ như hình bên

Đo độ dài hai đoạn thẳng trên.

Vẽ đoạn thẳng ${RX}$ có độ dài bằng tổng độ dài hai đoạn thẳng trên.
\end{minipage}\hfill
\begin{minipage}[t]{0.39\textwidth}
\vspace{0pt}
\centering

\resizebox{\linewidth}{!}{%
\begin{tikzpicture}[>=stealth,line join=round,line cap=round,font=\footnotesize,scale=1]
\path
(0,0) coordinate (A)
(0:3) coordinate (B)
(0,-1) coordinate (C)
($(C)+(-20:4)$) coordinate (D)
;
\draw (A)--(B) (C)--(D);
    \filldraw[black] (A) circle (1pt) node[above right] {$D$};
    \filldraw[black] (B) circle (1pt) node[above left] {$F$};
    \filldraw[black] (C) circle (1pt) node[above] {$K$};
   \filldraw[black] (D) circle (1pt) node[above right]{$M$};
    \end{tikzpicture}
}
\end{minipage}

\end{tcolorbox}
\vspace{4pt}


\begin{tcolorbox}[
    breakable,
    colback=white,
    colframe=blue!60!black,
    boxrule=0.8pt,
    arc=3pt,
    title={\textbf{Câu \#087} \quad \textbf{MÃ CÂU: \texttt{[L06\_C8\_B3\_13]}} \quad | \quad \textbf{Nhóm: \texttt{HINH\_082}} \quad | \quad \textbf{MapID: \texttt{6H8H3}} \hfill \textcolor{yellow!90!white}{[Hình gốc .pyc]}}
]
\textbf{Dạng toán:} Vẽ đoạn thẳng, tia, điểm theo yêu cầu \hfill \textit{(Tự luận -- Thông hiểu)}
\vspace{4pt}\hrule\vspace{6pt}
\begin{minipage}[t]{0.58\textwidth}
\vspace{0pt}
Vẽ hình theo cách diễn đạt bằng lời sau 

Vẽ đoạn thẳng ${RS=6}$cm.

Vẽ điểm ${Y}$ trên đoạn thẳng ${RS}$ sao cho ${RY=YS}$.

Vẽ tia ${YC}$ sao cho ${C}$ không thẳng hàng với ${R}$ và ${S}$.

Vẽ tia ${YT}$ là tia đối của tia ${YC}$.
\end{minipage}\hfill
\begin{minipage}[t]{0.39\textwidth}
\vspace{0pt}
\centering

\resizebox{\linewidth}{!}{%
\begin{tikzpicture}[>=stealth,line join=round,line cap=round,font=\footnotesize,scale=1]
\path
(0,0) coordinate (B)
(0:4) coordinate (C)
;
\draw (B)--($(B)!0.5!(C)$)node[above]{$6$}--(C);
\filldraw[black] (B) circle (1pt) node[above left] {$R$};
    \filldraw[black] (C) circle (1pt) node[above] {$S$};
\end{tikzpicture}
\begin{tikzpicture}[>=stealth,line join=round,line cap=round,font=\footnotesize,scale=1]
\path
(0,0) coordinate (B)
(0:4) coordinate (C)
(0:2) coordinate (M)
;
\draw (B)--($(B)!0.5!(M)$)node[above]{$3$}--(M)--($(M)!0.5!(C)$)node[above]{$3$}--(C);
  \filldraw[black] (B) circle (1pt) node[above left] {$R$};
    \filldraw[black] (C) circle (1pt) node[above] {$S$};
   \filldraw[black] (M) circle (1pt) node[above right]{$Y$};
\end{tikzpicture}
\begin{tikzpicture}[>=stealth,line join=round,line cap=round,font=\footnotesize,scale=1]
\path
(0,0) coordinate (B)
(0:4) coordinate (C)
(0:2) coordinate (M)
($(M)!0.5!30:(C)$) coordinate (A)
;
\draw (B)--($(B)!0.5!(M)$)node[above]{2cm}--(M)--($(M)!0.5!(C)$)node[above]{$3$}--(C) (M)--($(M)!1.2!(A)$);
 \filldraw[black] (B) circle (1pt) node[above left] {$R$};
    \filldraw[black] (C) circle (1pt) node[above] {$S$};
   \filldraw[black] (M) circle (1pt) node[above right]{$Y$};
    \filldraw[black] (A) circle (1pt) node[below left] {$C$};
\end{tikzpicture}
\begin{tikzpicture}[>=stealth,line join=round,line cap=round,font=\footnotesize,scale=1]
\path
(0,0) coordinate (B)
(0:4) coordinate (C)
(0:2) coordinate (M)
($(M)!0.5!30:(C)$) coordinate (A)
($(A)!2!(M)$) coordinate (D)
;
\draw (B)--($(B)!0.5!(M)$)node[above]{2cm}--(M)--($(M)!0.5!(C)$)node[above]{$3$}--(C) ($(M)!1.2!(D)$)--($(M)!1.2!(A)$);
 \filldraw[black] (B) circle (1pt) node[above left] {$R$};
    \filldraw[black] (C) circle (1pt) node[above] {$S$};
   \filldraw[black] (M) circle (1pt) node[above right]{$Y$};
    \filldraw[black] (A) circle (1pt) node[below left] {$C$};
    \filldraw[black] (D) circle (1pt) node[below left] {$T$};
    \end{tikzpicture}
}
\end{minipage}

\end{tcolorbox}
\vspace{4pt}


\begin{tcolorbox}[
    breakable,
    colback=white,
    colframe=blue!60!black,
    boxrule=0.8pt,
    arc=3pt,
    title={\textbf{Câu \#088} \quad \textbf{MÃ CÂU: \texttt{[L06\_C8\_B4\_01]}} \quad | \quad \textbf{Nhóm: \texttt{HINH\_083}} \quad | \quad \textbf{MapID: \texttt{6H8N4}} \hfill \textcolor{yellow!90!white}{[Hình gốc .pyc]}}
]
\textbf{Dạng toán:} Vẽ hình theo yêu cầu \hfill \textit{(Tự luận -- Nhận biết)}
\vspace{4pt}\hrule\vspace{6pt}
\begin{minipage}[t]{0.58\textwidth}
\vspace{0pt}
Học sinh vẽ hình theo cách diễn đạt bằng lời sau: (Vẽ trên cùng một hình)

Vẽ ba điểm ${F}$, ${H}$, ${M}$ không thẳng hàng.

Vẽ điểm ${R}$ là trung điểm của đoạn thẳng ${HM}$.

Vẽ tia ${RF}$.

Vẽ tia ${RN}$ là tia đối của tia ${RF}$ sao cho ${R}$ cũng là trung điểm của đoạn thẳng ${FN}$.
\end{minipage}\hfill
\begin{minipage}[t]{0.39\textwidth}
\vspace{0pt}
\centering

\resizebox{\linewidth}{!}{%
\begin{tikzpicture}[line join=round, line cap=round,>=stealth,thick,scale=1]

    \coordinate[label=below:$F$] (A) at (0,0);
    \coordinate[label=below:$H$] (B) at (2,0);
    \coordinate[label=above:$M$] (C) at (4,2);
    \coordinate[label=below right:$R$] (M) at ($(B)!.5!(C)$);
    \coordinate (H) at ($(M)!5/4!(A)$);
    \coordinate[label=below right:$N$] (D) at ($(A)!2!(M)$);
    \coordinate (K) at ($(M)!5/4!(D)$);
    \draw[] (B)--(C) (K)--(M)--(H) ;
    \foreach \diem in {A,B,C,M,D}\fill (\diem)circle(1.5pt);
\end{tikzpicture}
}
\end{minipage}

\end{tcolorbox}
\vspace{4pt}


\begin{tcolorbox}[
    breakable,
    colback=white,
    colframe=blue!60!black,
    boxrule=0.8pt,
    arc=3pt,
    title={\textbf{Câu \#089} \quad \textbf{MÃ CÂU: \texttt{[L06\_C8\_B4\_02]}} \quad | \quad \textbf{Nhóm: \texttt{HINH\_084}} \quad | \quad \textbf{MapID: \texttt{6H8N4}} \hfill \textcolor{yellow!90!white}{[Hình gốc .pyc]}}
]
\textbf{Dạng toán:} Vẽ hình theo yêu cầu 2 \hfill \textit{(Tự luận -- Nhận biết)}
\vspace{4pt}\hrule\vspace{6pt}
\begin{minipage}[t]{0.58\textwidth}
\vspace{0pt}
Học sinh vẽ hình theo cách diễn đạt bằng lời sau: (Vẽ trên cùng một hình)

Vẽ đoạn thẳng ${AD}$.

Vẽ điểm ${O}$ là trung điểm của đoạn thẳng ${AD}$.

Lấy điểm ${F}$ sao cho ba điểm ${A}$, ${D}$, ${F}$ không thẳng hàng, vẽ đường thẳng ${FO}$.

 Vẽ tia ${FD}$.
\end{minipage}\hfill
\begin{minipage}[t]{0.39\textwidth}
\vspace{0pt}
\centering

\resizebox{\linewidth}{!}{%
\begin{tikzpicture}[line join=round, line cap=round,>=stealth,thick,scale=1]

    \coordinate[label=below:$A$] (A) at (0,0);
    \coordinate[label=below right:$D$] (B) at (4,0);
    \coordinate[label=below:$O$] (K) at ($(A)!.5!(B)$);
    \coordinate[label=right:$F$] (H) at (5,2);
    \coordinate[] (I) at ($(H)!5/4!(B)$);
    \draw[] (A)--(B) (I)--(H)--(K);
    \foreach \diem in {A,B,K,H}\fill (\diem)circle(1.5pt);
\end{tikzpicture}
}
\end{minipage}

\end{tcolorbox}
\vspace{4pt}


\section{Toán 7}

\subsection{Chương 3 - Góc và đường thẳng song song}

\begin{tcolorbox}[
    breakable,
    colback=white,
    colframe=blue!60!black,
    boxrule=0.8pt,
    arc=3pt,
    title={\textbf{Câu \#090} \quad \textbf{MÃ CÂU: \texttt{[L07\_C3\_B1\_02]}} \quad | \quad \textbf{Nhóm: \texttt{HINH\_085}} \quad | \quad \textbf{MapID: \texttt{7H3N1}} \hfill \textcolor{yellow!90!white}{[Hình gốc .pyc]}}
]
\textbf{Dạng toán:} Tính số đo góc ở vị trí đối đỉnh \hfill \textit{(Trắc nghiệm -- Nhận biết)}
\vspace{4pt}\hrule\vspace{6pt}
\begin{minipage}[t]{0.58\textwidth}
\vspace{0pt}
Cho hình vẽ sau. Biết $\widehat{O_{1}} =115^{\circ}$, tính số đo góc $\widehat{O_{2}}$.

\vspace{4pt}\noindent
\textbf{A.} $\widehat{O_{2} }=65^{\circ}$ \quad \textbf{\textcolor{blue!75!black}{B. (Đ)}} $\widehat{O_{2} }=115^{\circ} $ \quad \textbf{C.} $\widehat{O_{2} }=125^{\circ}$ \quad \textbf{D.} $\widehat{O_{2} }=105^{\circ}$
\end{minipage}\hfill
\begin{minipage}[t]{0.39\textwidth}
\vspace{0pt}
\centering

\resizebox{\linewidth}{!}{%
\begin{tikzpicture}
            % Vẽ các điểm A, B, C, D
            \coordinate (A) at (0, 0);
            \coordinate (B) at (4, 4);
            \coordinate (C) at (6, 0);
            \coordinate (D) at (2, -3);
            
            % Vẽ đường thẳng đi qua A và C
            \draw[thick] (A) -- (C);
            
            % Vẽ đường thẳng đi qua B và D
            \draw[thick] (B) -- (D) ;
            
            % Tìm giao điểm O của 2 đường thẳng
            \coordinate (O) at (intersection of A--C and B--D);
            
            % Đánh dấu các điểm A, B, C, D, O

            \fill (O) circle (2pt) node[above right] {O};
             \draw 
            (O) pic[draw,angle radius = 10] {{angle = B--O--A}} node[shift={{(-49:-15pt)}}]{{$1$}} 
            ;
             \draw 
            (O) pic[draw,angle radius = 10] {{angle = D--O--C}} node[shift={{(-45:14pt)}}]{{$2$}} 
            ;
        \end{tikzpicture}
}
\end{minipage}

\end{tcolorbox}
\vspace{4pt}


\begin{tcolorbox}[
    breakable,
    colback=white,
    colframe=blue!60!black,
    boxrule=0.8pt,
    arc=3pt,
    title={\textbf{Câu \#091} \quad \textbf{MÃ CÂU: \texttt{[L07\_C3\_B1\_03]}} \quad | \quad \textbf{Nhóm: \texttt{HINH\_086}} \quad | \quad \textbf{MapID: \texttt{7H3N1}} \hfill \textcolor{yellow!90!white}{[Hình gốc .pyc]}}
]
\textbf{Dạng toán:} Tính số đo góc ở vị trí kề bù ( cần test lại) \hfill \textit{(Trắc nghiệm -- Nhận biết)}
\vspace{4pt}\hrule\vspace{6pt}
\begin{minipage}[t]{0.58\textwidth}
\vspace{0pt}
Cho hình vẽ sau. Biết $\widehat{O_{1}} =65^{\circ}$, tính số đo góc $\widehat{O_{2}}$.

\vspace{4pt}\noindent
\textbf{A.} $\widehat{O_{2} }=105^{\circ}$ \quad \textbf{\textcolor{blue!75!black}{B. (Đ)}} $\widehat{O_{2} }=115^{\circ} $ \quad \textbf{C.} $\widehat{O_{2} }=75^{\circ}$ \quad \textbf{D.} $\widehat{O_{2} }=65^{\circ}$
\end{minipage}\hfill
\begin{minipage}[t]{0.39\textwidth}
\vspace{0pt}
\centering

\resizebox{\linewidth}{!}{%
\begin{tikzpicture}
            % Vẽ các điểm A, B, C, D
            \coordinate (A) at (0, 0);
            \coordinate (B) at (4, 4);
            \coordinate (C) at (6, 0);
            \coordinate (D) at (2, -3);
            
            % Vẽ đường thẳng đi qua A và C
            \draw[thick] (A) -- (C);
            
            % Vẽ đường thẳng đi qua B và D
            \draw[thick] (B) -- (D) ;
            
            % Tìm giao điểm O của 2 đường thẳng
            \coordinate (O) at (intersection of A--C and B--D);
            
            % Đánh dấu các điểm A, B, C, D, O

            \fill (O) circle (2pt) node[above right] {O};
            \draw 
            (O)  node[shift={{(49:-15pt)}}]{{$1$}} 
            ;
            \draw 
            (O)  node[shift={{(-45:14pt)}}]{{$2$}} 
            ;
        \end{tikzpicture}
}
\end{minipage}

\end{tcolorbox}
\vspace{4pt}


\begin{tcolorbox}[
    breakable,
    colback=white,
    colframe=blue!60!black,
    boxrule=0.8pt,
    arc=3pt,
    title={\textbf{Câu \#092} \quad \textbf{MÃ CÂU: \texttt{[L07\_C3\_B1\_04]}} \quad | \quad \textbf{Nhóm: \texttt{HINH\_087}} \quad | \quad \textbf{MapID: \texttt{7H3N1}} \hfill \textcolor{yellow!90!white}{[Hình gốc .pyc]}}
]
\textbf{Dạng toán:} Đếm số lượng các cặp góc ở ví trí đặc biệt \hfill \textit{(Trắc nghiệm -- Nhận biết)}
\vspace{4pt}\hrule\vspace{6pt}
\begin{minipage}[t]{0.58\textwidth}
\vspace{0pt}
Cho hình vẽ sau. Có bao nhiêu cặp góc bù nhau trong hình?.

\vspace{4pt}\noindent
\textbf{A.} $5$ \quad \textbf{B.} $3$ \quad \textbf{C.} $2$ \quad \textbf{\textcolor{blue!75!black}{D. (Đ)}} $4$
\end{minipage}\hfill
\begin{minipage}[t]{0.39\textwidth}
\vspace{0pt}
\centering

\resizebox{\linewidth}{!}{%
\begin{tikzpicture}
            % Vẽ các điểm A, B, C, D
            \coordinate (A) at (0, 0);
            \coordinate (B) at (4, 4);
            \coordinate (C) at (6, 0);
            \coordinate (D) at (2, -3);
            
            % Vẽ đường thẳng đi qua A và C
            \draw[thick] (A) -- (C);
            
            % Vẽ đường thẳng đi qua B và D
            \draw[thick] (B) -- (D) ;
            
            % Tìm giao điểm O của 2 đường thẳng
            \coordinate (O) at (intersection of A--C and B--D);
            
            % Đánh dấu các điểm A, B, C, D, O

            \fill (O) circle (2pt) node[above right] {O};
               \draw 
            (O) node[shift={{(-49:-15pt)}}]{{$1$}} 
            ;
            \draw 
            (O) node[shift={{(-45:14pt)}}]{{$2$}} 
            ;
            \draw 
            (O) node[shift={{(45:-14pt)}}]{{$3$}} 
            ;
            \draw 
            (O) node[shift={{(55:24pt)}}]{{$4$}} 
            ;
        \end{tikzpicture}
}
\end{minipage}

\end{tcolorbox}
\vspace{4pt}


\begin{tcolorbox}[
    breakable,
    colback=white,
    colframe=blue!60!black,
    boxrule=0.8pt,
    arc=3pt,
    title={\textbf{Câu \#093} \quad \textbf{MÃ CÂU: \texttt{[L07\_C3\_B1\_05]}} \quad | \quad \textbf{Nhóm: \texttt{HINH\_088}} \quad | \quad \textbf{MapID: \texttt{7H3N1}} \hfill \textcolor{yellow!90!white}{[Hình gốc .pyc]}}
]
\textbf{Dạng toán:} Tính so đo góc khi cho phân giác \hfill \textit{(Trắc nghiệm -- Nhận biết)}
\vspace{4pt}\hrule\vspace{6pt}
\begin{minipage}[t]{0.58\textwidth}
\vspace{0pt}
Cho hình vẽ, biết ${Oz}$ là tia phân giác của góc $\widehat{xOy}$ và $\widehat{xOy}=120^{\circ}$. Số đo $\widehat{xOz}$ và $\widehat{yOz}$ là

\vspace{4pt}\noindent
\textbf{A.} $\widehat{xOz}=60^{\circ}$ và $\widehat{yOz} =30^{\circ}$ \quad \textbf{\textcolor{blue!75!black}{B. (Đ)}} $\widehat{xOz}=\widehat{yOz} =60^{\circ}$ \quad \textbf{C.} $\widehat{xOz}=\widehat{yOz}=120^{\circ}$ \quad \textbf{D.} $\widehat{xOz} =30^{\circ}$ và $\widehat{yOz}=60^{\circ}$
\end{minipage}\hfill
\begin{minipage}[t]{0.39\textwidth}
\vspace{0pt}
\centering

\resizebox{\linewidth}{!}{%
\begin{tikzpicture}
            % Vẽ các điểm A, B, C, D
             \coordinate [label=below left:$O$] (A) at (0, 0);
            \coordinate (C) at (-3.5, 5.94);
            \coordinate (B) at (6, 0);
            \coordinate (D) at (2.57, 4.5);
            
            % Vẽ đường thẳng đi qua A và C
            \draw[thick] (A) -- (B);
            % Vẽ đường thẳng đi qua A và C
            \draw[thick] (A) -- (C);
            % Vẽ đường thẳng đi qua B và D
            \draw[thick] (A) -- (D) ;
            \draw 
            (A) pic[draw,angle radius = 16] {{angle = B--A--D}}   ;
            \draw 
            (A) pic[draw,angle radius = 13] {{angle = D--A--C}}   ;
             % Thêm nhãn độ dài cạnh
            \node at (6, -0.5) {{$x$ }};  % Nhãn độ dài cạnh AC với vị trí cố định
                \node at (2.57, 4) {{$z$ }};
                    \node at (-3.3, 5.2) {{$y$ }};
        \end{tikzpicture}
}
\end{minipage}

\end{tcolorbox}
\vspace{4pt}


\begin{tcolorbox}[
    breakable,
    colback=white,
    colframe=blue!60!black,
    boxrule=0.8pt,
    arc=3pt,
    title={\textbf{Câu \#094} \quad \textbf{MÃ CÂU: \texttt{[L07\_C3\_B1\_06]}} \quad | \quad \textbf{Nhóm: \texttt{HINH\_089}} \quad | \quad \textbf{MapID: \texttt{7H3N1}} \hfill \textcolor{yellow!90!white}{[Hình gốc .pyc]}}
]
\textbf{Dạng toán:} Tính số đo góc kề nhau. \hfill \textit{(Trắc nghiệm -- Nhận biết)}
\vspace{4pt}\hrule\vspace{6pt}
\begin{minipage}[t]{0.58\textwidth}
\vspace{0pt}
Cho hình vẽ sau: biết $\widehat{O_{1}} =65^{\circ}$; $\widehat{O_{3}} =81^{\circ}$. Tính số đo góc $\widehat{O_{2}}$

\vspace{4pt}\noindent
\textbf{A.} $\widehat{O_{2}}=44^{\circ}$ \quad \textbf{\textcolor{blue!75!black}{B. (Đ)}} $\widehat{O_{2}}=34^{\circ}$ \quad \textbf{C.} $\widehat{O_{2}}=24^{\circ}$ \quad \textbf{D.} $\widehat{O_{2}}=39^{\circ}$
\end{minipage}\hfill
\begin{minipage}[t]{0.39\textwidth}
\vspace{0pt}
\centering

\resizebox{\linewidth}{!}{%
\begin{tikzpicture}
            % Vẽ các điểm A, B, C, D
            \coordinate (A) at (0, 0);
            \coordinate (B) at (4, 4);
            \coordinate (C) at (6, 0);
            \coordinate (D) at (2, -3);
            \coordinate (E) at (4, -4);
                % Tìm giao điểm O của 2 đường thẳng
            \coordinate (O) at (intersection of A--C and B--D);
            
            % Vẽ đường thẳng đi qua A và C
            \draw[thick] (A) -- (C);
            
            % Vẽ đường thẳng đi qua B và D
            \draw[thick] (O) -- (D) ;
                \draw[thick] (O) -- (E) ;
        
            % Đánh dấu các điểm A, B, C, D, O
            
            \fill (O) circle (2pt) node[above right] {O};
            \draw 
            (O) pic[draw,angle radius = 8] {{angle = E--O--C}} node[shift={{(-25:15pt)}}]{{$1$}}   
            ;
            \draw 
            (O) pic[draw,angle radius = 14] {{angle = D--O--E}} pic[draw,angle radius = 16] {{angle = D--O--E}}  node[shift={{(85:-25pt)}}]{{$2$}}   
            ;
            \draw 
        (O)pic[draw,angle radius = 10] {{angle = A--O--D}} pic[draw,angle radius = 12] {{angle = A--O--D}} pic[draw,angle radius = 14] {{angle = A--O--D}}  node[shift={{(45:-25pt)}}]{{$3$}}   
        ;
         \node at (0.1, 0.1) {{$x$ }};  % Nhãn độ dài cạnh AC với vị trí cố định
        \node at (6.1, 0.1) {{$z$ }};
        \node at (2.1, -3.1) {{$y$ }};
            \node at (4.1, -3.8) {{$t$ }};
        \end{tikzpicture}
}
\end{minipage}

\vspace{3pt}\hrule\vspace{3pt}\footnotesize\textcolor{gray!80!black}{\textbf{Các mã dùng chung hình (HINH\_089):} \texttt{L07\_C3\_B1\_08}}
\end{tcolorbox}
\vspace{4pt}


\begin{tcolorbox}[
    breakable,
    colback=white,
    colframe=blue!60!black,
    boxrule=0.8pt,
    arc=3pt,
    title={\textbf{Câu \#095} \quad \textbf{MÃ CÂU: \texttt{[L07\_C3\_B1\_07]}} \quad | \quad \textbf{Nhóm: \texttt{HINH\_090}} \quad | \quad \textbf{MapID: \texttt{7H3N1}} \hfill \textcolor{yellow!90!white}{[Hình gốc .pyc]}}
]
\textbf{Dạng toán:} Liệt các các cặp góc ở vị trí đặc biệt. \hfill \textit{(Tự luận -- Nhận biết)}
\vspace{4pt}\hrule\vspace{6pt}
\begin{minipage}[t]{0.58\textwidth}
\vspace{0pt}
Liệt kê các cặp góc đối đỉnh, kề bù, kề nhau trong hình vẽ dưới đây
\end{minipage}\hfill
\begin{minipage}[t]{0.39\textwidth}
\vspace{0pt}
\centering

\resizebox{\linewidth}{!}{%
\begin{tikzpicture}
            \coordinate (O) at (3/2, 3/2);
            \coordinate (D) at (2, 0);
            \coordinate (C) at (5, 3);
            \coordinate (B) at (1, 3);
            \coordinate (A) at (-2, 0);
            \
            % Vẽ tứ giác
            \draw[thick] (O) -- (A) -- (C) ;
            \draw[thick] (A) -- (B)--(C)--(D)--(A);
        \draw[thick] (O) -- (B) -- (D) ;
            % Ghi nhãn các điểm
            \node[above right] at (O) {$O $};
            \node[above left] at (B) {$B $};
            \node[above right] at (C) {$C $};
            \node[below right] at (D) {$D $};
            \node[below left] at (A){$A $};
                \draw
            ($(B)!0.4!(O)$)node[rotate=80]{$/$} ($(O)!0.4!(D)$)node[rotate=80]{$/$};
    \end{tikzpicture}
}
\end{minipage}

\vspace{3pt}\hrule\vspace{3pt}\footnotesize\textcolor{gray!80!black}{\textbf{Các mã dùng chung hình (HINH\_090):} \texttt{L08\_C3\_B3\_01, L08\_C3\_B3\_09}}
\end{tcolorbox}
\vspace{4pt}


\begin{tcolorbox}[
    breakable,
    colback=white,
    colframe=blue!60!black,
    boxrule=0.8pt,
    arc=3pt,
    title={\textbf{Câu \#096} \quad \textbf{MÃ CÂU: \texttt{[L07\_C3\_B1\_08]}} \quad | \quad \textbf{Nhóm: \texttt{HINH\_089}} \quad | \quad \textbf{MapID: \texttt{7H3N1}} \hfill \textcolor{yellow!90!white}{[Hình gốc .pyc]}}
]
\textbf{Dạng toán:} Liệt các các cặp góc ở vị trí đặc biệt. \hfill \textit{(Tự luận -- Nhận biết)}
\vspace{4pt}\hrule\vspace{6pt}
\begin{minipage}[t]{0.58\textwidth}
\vspace{0pt}
Cho hình vẽ sau: biết $\widehat{O_{1}} =65^{\circ}$; $\widehat{O_{3}} =81^{\circ}$. Tính số đo góc $\widehat{O_{3}}$
\end{minipage}\hfill
\begin{minipage}[t]{0.39\textwidth}
\vspace{0pt}
\centering

\resizebox{\linewidth}{!}{%
\begin{tikzpicture}
            % Vẽ các điểm A, B, C, D
            \coordinate (A) at (0, 0);
            \coordinate (B) at (4, 4);
            \coordinate (C) at (6, 0);
            \coordinate (D) at (2, -3);
            \coordinate (E) at (4, -4);
                % Tìm giao điểm O của 2 đường thẳng
            \coordinate (O) at (intersection of A--C and B--D);
            
            % Vẽ đường thẳng đi qua A và C
            \draw[thick] (A) -- (C);
            
            % Vẽ đường thẳng đi qua B và D
            \draw[thick] (O) -- (D) ;
                \draw[thick] (O) -- (E) ;
        
            % Đánh dấu các điểm A, B, C, D, O
            
            \fill (O) circle (2pt) node[above right] {O};
            \draw 
            (O) pic[draw,angle radius = 8] {{angle = E--O--C}} node[shift={{(-25:15pt)}}]{{$1$}}   
            ;
            \draw 
            (O) pic[draw,angle radius = 14] {{angle = D--O--E}} pic[draw,angle radius = 16] {{angle = D--O--E}}  node[shift={{(85:-25pt)}}]{{$2$}}   
            ;
            \draw 
        (O)pic[draw,angle radius = 10] {{angle = A--O--D}} pic[draw,angle radius = 12] {{angle = A--O--D}} pic[draw,angle radius = 14] {{angle = A--O--D}}  node[shift={{(45:-25pt)}}]{{$3$}}   
        ;
         \node at (0.1, 0.1) {{$x$ }};  % Nhãn độ dài cạnh AC với vị trí cố định
        \node at (6.1, 0.1) {{$z$ }};
        \node at (2.1, -3.1) {{$y$ }};
            \node at (4.1, -3.8) {{$t$ }};
        \end{tikzpicture}
}
\end{minipage}

\vspace{3pt}\hrule\vspace{3pt}\footnotesize\textcolor{gray!80!black}{\textbf{Các mã dùng chung hình (HINH\_089):} \texttt{L07\_C3\_B1\_06}}
\end{tcolorbox}
\vspace{4pt}


\begin{tcolorbox}[
    breakable,
    colback=white,
    colframe=blue!60!black,
    boxrule=0.8pt,
    arc=3pt,
    title={\textbf{Câu \#097} \quad \textbf{MÃ CÂU: \texttt{[L07\_C3\_B2\_11]}} \quad | \quad \textbf{Nhóm: \texttt{HINH\_091}} \quad | \quad \textbf{MapID: \texttt{7H3N2}} \hfill \textcolor{yellow!90!white}{[Hình gốc .pyc]}}
]
\textbf{Dạng toán:} Nhận biết vị trí đặc biệt của các cặp góc \hfill \textit{(Trắc nghiệm -- Nhận biết)}
\vspace{4pt}\hrule\vspace{6pt}
\begin{minipage}[t]{0.58\textwidth}
\vspace{0pt}
Cho hình sau. Khẳng định nào trong các khẳng định sau đúng?

\vspace{4pt}\noindent
\textbf{A.} $\widehat {F_{2}}$ và $\widehat {E_{3}}$ là hai góc đồng vị \quad \textbf{B.} $\widehat {F_{3}}$ và $\widehat {E_{2}}$ là hai góc so le ngoài \quad \textbf{C.} $\widehat {F_{1}}$ và $\widehat {E_{4}}$ là hai góc so le trong \quad \textbf{\textcolor{blue!75!black}{D. (Đ)}} $\widehat {F_{4}}$ và $\widehat {E_{4}}$ là hai góc đồng vị
\end{minipage}\hfill
\begin{minipage}[t]{0.39\textwidth}
\vspace{0pt}
\centering

\resizebox{\linewidth}{!}{%
\begin{tikzpicture}
            % Vẽ các điểm A, B, C, D
            \coordinate (A) at (-2, 0);
            \coordinate (C) at (-3, -2);
            \coordinate (B) at (2, 0);
            \coordinate (D) at (-1, 2);
            \coordinate (E) at (-4, 0);
            \coordinate (F) at (3, -2);
            \coordinate (M) at (1, 2);
            \coordinate (N) at (4, 0);
            
            \draw[thick] (E) -- (N);
            \draw[thick] (C) -- (D) ;
            \draw[thick] (M) -- (F) ;
            
            % Đánh dấu các điểm A, B, 
            \fill (B) circle (2pt) node[above] {$E$};
            \fill (A) circle (2pt) node[above] {$F$};
            \draw 
            (B) node[shift={(-25:15pt)}]{$1$}   
            ;
            \draw 
            (B) node[shift={(45:-15pt)}]{$2$}   
            ;
            \draw 
            (B) node[shift={(-50:-15pt)}]{$3$}   
            ;
            \draw 
            (B) node[shift={(50:15pt)}]{$4$}   
            ;
            \draw 
            (A)  node[shift={(-50:12pt)}]{$1$}   
            ;
            \draw 
            (A) node[shift={(50:-13pt)}]{$2$}   
            ;
            \draw 
            (A)  node[shift={(-50:-15pt)}]{$3$}   
            ;
            \draw 
            (A) node[shift={(50:15pt)}]{$4$}   
            ;
            
        \end{tikzpicture}
}
\end{minipage}

\vspace{3pt}\hrule\vspace{3pt}\footnotesize\textcolor{gray!80!black}{\textbf{Các mã dùng chung hình (HINH\_091):} \texttt{L07\_C3\_B2\_06, L07\_C3\_B3\_04}}
\end{tcolorbox}
\vspace{4pt}


\begin{tcolorbox}[
    breakable,
    colback=white,
    colframe=blue!60!black,
    boxrule=0.8pt,
    arc=3pt,
    title={\textbf{Câu \#098} \quad \textbf{MÃ CÂU: \texttt{[L07\_C3\_B2\_12]}} \quad | \quad \textbf{Nhóm: \texttt{HINH\_092}} \quad | \quad \textbf{MapID: \texttt{7H3N2}} \hfill \textcolor{yellow!90!white}{[Hình gốc .pyc]}}
]
\textbf{Dạng toán:} Nhận biết hai đường thẳng song song \hfill \textit{(Trắc nghiệm -- Nhận biết)}
\vspace{4pt}\hrule\vspace{6pt}
\begin{minipage}[t]{0.58\textwidth}
\vspace{0pt}
Cho hình vẽ sau biết $\widehat { I_{4} }=\widehat{ J_{1} }=47^{\circ}$. Chọn khẳng sai trong các khẳng định sau?

\vspace{4pt}\noindent
\textbf{A.} $\widehat{ I_{1} }=47^{\circ}$ \quad \textbf{\textcolor{blue!75!black}{B. (Đ)}} $\widehat{ J_{4} } = 48^{\circ}$ \quad \textbf{C.} $y//x$ \quad \textbf{D.} $\widehat { I_{2} } =133^{\circ} $
\end{minipage}\hfill
\begin{minipage}[t]{0.39\textwidth}
\vspace{0pt}
\centering

\resizebox{\linewidth}{!}{%
\begin{tikzpicture}
            \coordinate (E) at (1.14, 3.06);
            \coordinate (D) at (4.46, 0);
            \coordinate (C) at (-0.6, 0);
            \coordinate (B) at (5, 2);
            \coordinate (A) at (0, 2);
            \coordinate (F) at (3.6, -1.44);
            \coordinate (H) at (2.81, 0);
            \coordinate (G) at (1.72, 2);
            % Vẽ tứ giác
            \draw[thick] (C) -- (H)--(D);
            \draw[thick] (A) -- (G) -- (B);
            \draw[thick] (E) -- (G) -- (H)--(F);
            % Ghi nhãn các điểm
            
            \node[above] at (G) {$I$};
            \node[above ] at (H){$J$};
                \node[above] at (A) {$y$};
                        \node[above] at (C) {$x$};
            \draw 
            (G) pic[draw,angle radius = 8] {angle = E--G--A} node[shift={(-25:15pt)}]{$1$}  ;
            \draw (G) node[shift={(30:11pt)}]{$2$};
            \draw  (G) node[shift={(50:-15pt)}]{$3$}  ;
            \draw (G)  node[shift={(-25:-15pt)}]{$4$};  
            \draw (H) pic[draw,angle radius = 10] {angle = F--H--D};
            \draw (H) node[shift={(-25:15pt)}]{$1$};
            \draw (H) node[shift={(40:15pt)}]{$2$}  ;
            \draw  (H) node[shift={(50:-13pt)}]{$3$} ;
            \draw (H) node[shift={(-25:-15pt)}]{$4$}  ;     
        \end{tikzpicture}
}
\end{minipage}

\vspace{3pt}\hrule\vspace{3pt}\footnotesize\textcolor{gray!80!black}{\textbf{Các mã dùng chung hình (HINH\_092):} \texttt{L07\_C3\_B2\_07, L07\_C3\_B3\_09, L07\_C3\_B3\_05}}
\end{tcolorbox}
\vspace{4pt}


\begin{tcolorbox}[
    breakable,
    colback=white,
    colframe=blue!60!black,
    boxrule=0.8pt,
    arc=3pt,
    title={\textbf{Câu \#099} \quad \textbf{MÃ CÂU: \texttt{[L07\_C3\_B2\_13]}} \quad | \quad \textbf{Nhóm: \texttt{HINH\_093}} \quad | \quad \textbf{MapID: \texttt{7H3H2}} \hfill \textcolor{yellow!90!white}{[Hình gốc .pyc]}}
]
\textbf{Dạng toán:} Thông hiểu hai đường thẳng song song dựa vào 2 góc đồng vị \hfill \textit{(Trắc nghiệm -- Thông hiểu)}
\vspace{4pt}\hrule\vspace{6pt}
\begin{minipage}[t]{0.58\textwidth}
\vspace{0pt}
Cho hình vẽ biết $\widehat { A_{4} }= \widehat{ B_{4} }=53^{\circ}$. Trong các khẳng định sau khẳng định nào sai.

\vspace{4pt}\noindent
\textbf{A.} $ \widehat{ B_{4} }=53^{\circ}$ \quad \textbf{B.} $n //m$ \quad \textbf{C.} $\widehat{  A_{2} }=127^{\circ}$ \quad \textbf{\textcolor{blue!75!black}{D. (Đ)}} $\widehat{ B_{3} } =107^{\circ}$
\end{minipage}\hfill
\begin{minipage}[t]{0.39\textwidth}
\vspace{0pt}
\centering

\resizebox{\linewidth}{!}{%
\begin{tikzpicture}
            \coordinate (E) at (1.14, 3.06);
            \coordinate (D) at (4.46, 0);
            \coordinate (C) at (-0.6, 0);
            \coordinate (B) at (5, 2);
            \coordinate (A) at (0, 2);
            \coordinate (F) at (3.6, -1.44);
            \coordinate (H) at (2.81, 0);
            \coordinate (G) at (1.72, 2);
            % Vẽ tứ giác
            \draw[thick] (C) -- (H)--(D);
            \draw[thick] (A) -- (G) -- (B);
            \draw[thick] (E) -- (G) -- (H)--(F);
            % Ghi nhãn các điểm
            
            \node[above] at (G) {$A$};
            \node[above ] at (H){$B$};
                \node[above] at (A) {$n$};
                        \node[above] at (C) {$m$};
            \draw 
            (G) pic[draw,angle radius = 8] {angle = E--G--A} node[shift={(-25:15pt)}]{$1$}  ;
            \draw (G) node[shift={(30:11pt)}]{$2$};
            \draw  (G) node[shift={(50:-15pt)}]{$3$}  ;
            \draw (G)  node[shift={(-25:-15pt)}]{$4$};  
            \draw (H) pic[draw,angle radius = 10] {angle = E--H--C};
            \draw (H) node[shift={(-25:15pt)}]{$1$};
            \draw (H) node[shift={(40:15pt)}]{$2$}  ;
            \draw  (H) node[shift={(50:-13pt)}]{$3$} ;
            \draw (H) node[shift={(-25:-15pt)}]{$4$}  ;     
        \end{tikzpicture}
}
\end{minipage}

\vspace{3pt}\hrule\vspace{3pt}\footnotesize\textcolor{gray!80!black}{\textbf{Các mã dùng chung hình (HINH\_093):} \texttt{L07\_C3\_B2\_08, L07\_C3\_B3\_10, L07\_C3\_B3\_06}}
\end{tcolorbox}
\vspace{4pt}


\begin{tcolorbox}[
    breakable,
    colback=white,
    colframe=blue!60!black,
    boxrule=0.8pt,
    arc=3pt,
    title={\textbf{Câu \#100} \quad \textbf{MÃ CÂU: \texttt{[L07\_C3\_B2\_14]}} \quad | \quad \textbf{Nhóm: \texttt{HINH\_094}} \quad | \quad \textbf{MapID: \texttt{7H3H2}} \hfill \textcolor{yellow!90!white}{[Hình gốc .pyc]}}
]
\textbf{Dạng toán:} Thông hiểu 2 đường thẳng song song \hfill \textit{(Trắc nghiệm -- Thông hiểu)}
\vspace{4pt}\hrule\vspace{6pt}
\begin{minipage}[t]{0.58\textwidth}
\vspace{0pt}
Cho hình vẽ biết $\widehat { A _{1} }= \widehat{ B _{4} }=56^{\circ}$. Chọn khẳng định sai trong các khẳng định sau.

\vspace{4pt}\noindent
\textbf{A.} $d //d'$ \quad \textbf{B.} $\widehat { A _{2} }= \widehat{ B _{3} }=124^{\circ}$ \quad \textbf{\textcolor{blue!75!black}{C. (Đ)}} $\widehat{ A _{3} }=\widehat{  B _{2} }=114^{\circ}$ \quad \textbf{D.} $\widehat { A _{4} }=\widehat{ B _{1} }= 56^{\circ}$
\end{minipage}\hfill
\begin{minipage}[t]{0.39\textwidth}
\vspace{0pt}
\centering

\resizebox{\linewidth}{!}{%
\begin{tikzpicture}
            \coordinate (E) at (1.14, 3.06);
            \coordinate (D) at (4.46, 0);
            \coordinate (C) at (-0.6, 0);
            \coordinate (B) at (5, 2);
            \coordinate (A) at (0, 2);
            \coordinate (F) at (3.6, -1.44);
            \coordinate (H) at (2.81, 0);
            \coordinate (G) at (1.72, 2);
            % Vẽ tứ giác
            \draw[thick] (C) -- (H)--(D);
            \draw[thick] (A) -- (G) -- (B);
            \draw[thick] (E) -- (G) -- (H)--(F);
            % Ghi nhãn các điểm
            
            \node[above] at (G) {$A $};
            \node[above ] at (H){$B $};
                      \node[above] at (A) {$d$};
                        \node[above] at (C) {$d'$};
            \draw 
            (G) pic[draw,angle radius = 8] {angle = H--G--B} node[shift={(-25:15pt)}]{$1$}  ;
            \draw (G) node[shift={(30:11pt)}]{$2$};
            \draw  (G) node[shift={(50:-15pt)}]{$3$}  ;
            \draw (G)  node[shift={(-25:-15pt)}]{$4$};  
            \draw (H) pic[draw,angle radius = 10] {angle = E--H--C};
            \draw (H) node[shift={(-25:15pt)}]{$1$};
            \draw (H) node[shift={(40:15pt)}]{$2$}  ;
            \draw  (H) node[shift={(50:-13pt)}]{$3$} ;
            \draw (H) node[shift={(-25:-15pt)}]{$4$}  ;     
        \end{tikzpicture}
}
\end{minipage}

\vspace{3pt}\hrule\vspace{3pt}\footnotesize\textcolor{gray!80!black}{\textbf{Các mã dùng chung hình (HINH\_094):} \texttt{L07\_C3\_B2\_09, L07\_C3\_B3\_11, L07\_C3\_B3\_07}}
\end{tcolorbox}
\vspace{4pt}


\begin{tcolorbox}[
    breakable,
    colback=white,
    colframe=blue!60!black,
    boxrule=0.8pt,
    arc=3pt,
    title={\textbf{Câu \#101} \quad \textbf{MÃ CÂU: \texttt{[L07\_C3\_B2\_15]}} \quad | \quad \textbf{Nhóm: \texttt{HINH\_095}} \quad | \quad \textbf{MapID: \texttt{7H3H2}} \hfill \textcolor{yellow!90!white}{[Hình gốc .pyc]}}
]
\textbf{Dạng toán:} Thông hiểu 2 đường thẳng song song 2 \hfill \textit{(Trắc nghiệm -- Thông hiểu)}
\vspace{4pt}\hrule\vspace{6pt}
\begin{minipage}[t]{0.58\textwidth}
\vspace{0pt}
Cho hình vẽ biết $\widehat { A_{1} }=53^{\circ}$, $\widehat { B_{3} }=127^{\circ}$. Trong các khẳng định sau khẳng định nào sai.

\vspace{4pt}\noindent
\textbf{A.} $\widehat { A_{2} } =127^{\circ} $ m \quad \textbf{\textcolor{blue!75!black}{B. (Đ)}} $\widehat{  B_{2} }=53^{\circ}$ \quad \textbf{C.} $m  //n $ \quad \textbf{D.} $\widehat{ B_{4} } = 53^{\circ}$
\end{minipage}\hfill
\begin{minipage}[t]{0.39\textwidth}
\vspace{0pt}
\centering

\resizebox{\linewidth}{!}{%
\begin{tikzpicture}
            \coordinate (E) at (1.14, 3.06);
            \coordinate (D) at (4.46, 0);
            \coordinate (C) at (-0.6, 0);
            \coordinate (B) at (5, 2);
            \coordinate (A) at (0, 2);
            \coordinate (F) at (3.6, -1.44);
            \coordinate (H) at (2.81, 0);
            \coordinate (G) at (1.72, 2);
            % Vẽ tứ giác
            \draw[thick] (C) -- (H)--(D);
            \draw[thick] (A) -- (G) -- (B);
            \draw[thick] (E) -- (G) -- (H)--(F);
            % Ghi nhãn các điểm
            
            \node[above] at (G) {$A$};
            \node[above ] at (H){$B$};
                      \node[above] at (A) {$m $};
                        \node[above] at (C) {$n $};
            \draw 
            (G) node[shift={(-25:15pt)}]{$1$}  ;
            \draw (G) node[shift={(30:11pt)}]{$2$};
            \draw  (G) node[shift={(50:-15pt)}]{$3$}  ;
            \draw (G)  node[shift={(-25:-15pt)}]{$4$};  
           
            \draw (H) node[shift={(-25:15pt)}]{$1$};
            \draw (H) node[shift={(40:15pt)}]{$2$}  ;
            \draw  (H) node[shift={(50:-13pt)}]{$3$} ;
            \draw (H) node[shift={(-25:-15pt)}]{$4$}  ;     
        \end{tikzpicture}
}
\end{minipage}

\vspace{3pt}\hrule\vspace{3pt}\footnotesize\textcolor{gray!80!black}{\textbf{Các mã dùng chung hình (HINH\_095):} \texttt{L07\_C3\_B2\_16, L07\_C3\_B2\_17, L07\_C3\_B2\_10, L07\_C3\_B3\_12, L07\_C3\_B3\_13, L07\_C3\_B3\_14, L07\_C3\_B3\_08}}
\end{tcolorbox}
\vspace{4pt}


\begin{tcolorbox}[
    breakable,
    colback=white,
    colframe=blue!60!black,
    boxrule=0.8pt,
    arc=3pt,
    title={\textbf{Câu \#102} \quad \textbf{MÃ CÂU: \texttt{[L07\_C3\_B2\_16]}} \quad | \quad \textbf{Nhóm: \texttt{HINH\_095}} \quad | \quad \textbf{MapID: \texttt{7H3H2}} \hfill \textcolor{yellow!90!white}{[Hình gốc .pyc]}}
]
\textbf{Dạng toán:} Tính số đo góc 1. \hfill \textit{(Trắc nghiệm -- Thông hiểu)}
\vspace{4pt}\hrule\vspace{6pt}
\begin{minipage}[t]{0.58\textwidth}
\vspace{0pt}
Cho hình vẽ biết $a //b$ và $\widehat { A_{1} }=47^{\circ}$. Trong các khẳng định sau khẳng định nào sai.

\vspace{4pt}\noindent
\textbf{\textcolor{blue!75!black}{A. (Đ)}} $\widehat{  B_{2} }=47^{\circ}$ \quad \textbf{B.} $\widehat { A_{4} }=47^{\circ}$ \quad \textbf{C.} $\widehat { A_{2} } =133^{\circ} $ \quad \textbf{D.} $\widehat{ B_{4} } = 47^{\circ}$
\end{minipage}\hfill
\begin{minipage}[t]{0.39\textwidth}
\vspace{0pt}
\centering

\resizebox{\linewidth}{!}{%
\begin{tikzpicture}
            \coordinate (E) at (1.14, 3.06);
            \coordinate (D) at (4.46, 0);
            \coordinate (C) at (-0.6, 0);
            \coordinate (B) at (5, 2);
            \coordinate (A) at (0, 2);
            \coordinate (F) at (3.6, -1.44);
            \coordinate (H) at (2.81, 0);
            \coordinate (G) at (1.72, 2);
            % Vẽ tứ giác
            \draw[thick] (C) -- (H)--(D);
            \draw[thick] (A) -- (G) -- (B);
            \draw[thick] (E) -- (G) -- (H)--(F);
            % Ghi nhãn các điểm
            
            \node[above] at (G) {$A$};
            \node[above ] at (H){$B$};
                      \node[above] at (A) {$a$};
                        \node[above] at (C) {$b$};
            \draw 
            (G) node[shift={(-25:15pt)}]{$1$}  ;
            \draw (G) node[shift={(30:11pt)}]{$2$};
            \draw  (G) node[shift={(50:-15pt)}]{$3$}  ;
            \draw (G)  node[shift={(-25:-15pt)}]{$4$};  
           
            \draw (H) node[shift={(-25:15pt)}]{$1$};
            \draw (H) node[shift={(40:15pt)}]{$2$}  ;
            \draw  (H) node[shift={(50:-13pt)}]{$3$} ;
            \draw (H) node[shift={(-25:-15pt)}]{$4$}  ;     
        \end{tikzpicture}
}
\end{minipage}

\vspace{3pt}\hrule\vspace{3pt}\footnotesize\textcolor{gray!80!black}{\textbf{Các mã dùng chung hình (HINH\_095):} \texttt{L07\_C3\_B2\_15, L07\_C3\_B2\_17, L07\_C3\_B2\_10, L07\_C3\_B3\_12, L07\_C3\_B3\_13, L07\_C3\_B3\_14, L07\_C3\_B3\_08}}
\end{tcolorbox}
\vspace{4pt}


\begin{tcolorbox}[
    breakable,
    colback=white,
    colframe=blue!60!black,
    boxrule=0.8pt,
    arc=3pt,
    title={\textbf{Câu \#103} \quad \textbf{MÃ CÂU: \texttt{[L07\_C3\_B2\_17]}} \quad | \quad \textbf{Nhóm: \texttt{HINH\_095}} \quad | \quad \textbf{MapID: \texttt{7H3H2}} \hfill \textcolor{yellow!90!white}{[Hình gốc .pyc]}}
]
\textbf{Dạng toán:} Tính số đo góc 2. \hfill \textit{(Trắc nghiệm -- Thông hiểu)}
\vspace{4pt}\hrule\vspace{6pt}
\begin{minipage}[t]{0.58\textwidth}
\vspace{0pt}
Cho hình vẽ biết $a //b$ và $\widehat { A _{4} }=63^{\circ}$. Trong các khẳng định sau khẳng định nào sai.

\vspace{4pt}\noindent
\textbf{A.} $\widehat { A _{1} }=63^{\circ}$ \quad \textbf{B.} $\widehat { B _{1} } =63^{\circ} $ \quad \textbf{C.} $\widehat{ B _{3} } = 117^{\circ}$ \quad \textbf{\textcolor{blue!75!black}{D. (Đ)}} $\widehat{  A _{3} }=63^{\circ}$
\end{minipage}\hfill
\begin{minipage}[t]{0.39\textwidth}
\vspace{0pt}
\centering

\resizebox{\linewidth}{!}{%
\begin{tikzpicture}
            \coordinate (E) at (1.14, 3.06);
            \coordinate (D) at (4.46, 0);
            \coordinate (C) at (-0.6, 0);
            \coordinate (B) at (5, 2);
            \coordinate (A) at (0, 2);
            \coordinate (F) at (3.6, -1.44);
            \coordinate (H) at (2.81, 0);
            \coordinate (G) at (1.72, 2);
            % Vẽ tứ giác
            \draw[thick] (C) -- (H)--(D);
            \draw[thick] (A) -- (G) -- (B);
            \draw[thick] (E) -- (G) -- (H)--(F);
            % Ghi nhãn các điểm
            
            \node[above] at (G) {$A $};
            \node[above ] at (H){$B $};
                      \node[above] at (A) {$a$};
                        \node[above] at (C) {$b$};
            \draw 
            (G) node[shift={(-25:15pt)}]{$1$}  ;
            \draw (G) node[shift={(30:11pt)}]{$2$};
            \draw  (G) node[shift={(50:-15pt)}]{$3$}  ;
            \draw (G)  node[shift={(-25:-15pt)}]{$4$};  
           
            \draw (H) node[shift={(-25:15pt)}]{$1$};
            \draw (H) node[shift={(40:15pt)}]{$2$}  ;
            \draw  (H) node[shift={(50:-13pt)}]{$3$} ;
            \draw (H) node[shift={(-25:-15pt)}]{$4$}  ;     
        \end{tikzpicture}
}
\end{minipage}

\vspace{3pt}\hrule\vspace{3pt}\footnotesize\textcolor{gray!80!black}{\textbf{Các mã dùng chung hình (HINH\_095):} \texttt{L07\_C3\_B2\_15, L07\_C3\_B2\_16, L07\_C3\_B2\_10, L07\_C3\_B3\_12, L07\_C3\_B3\_13, L07\_C3\_B3\_14, L07\_C3\_B3\_08}}
\end{tcolorbox}
\vspace{4pt}


\begin{tcolorbox}[
    breakable,
    colback=white,
    colframe=blue!60!black,
    boxrule=0.8pt,
    arc=3pt,
    title={\textbf{Câu \#104} \quad \textbf{MÃ CÂU: \texttt{[L07\_C3\_B2\_18]}} \quad | \quad \textbf{Nhóm: \texttt{HINH\_096}} \quad | \quad \textbf{MapID: \texttt{7H3V2}} \hfill \textcolor{yellow!90!white}{[Hình gốc .pyc]}}
]
\textbf{Dạng toán:} Chứng minh các đường thẳng song song 1. \hfill \textit{(Trắc nghiệm -- Vận dụng)}
\vspace{4pt}\hrule\vspace{6pt}
\begin{minipage}[t]{0.58\textwidth}
\vspace{0pt}
Cho hình vẽ sau. Trong các khẳng định sau khẳng định nào sai?

\vspace{4pt}\noindent
\textbf{A.} $\widehat{QKF }= 60^{\circ}$ \quad \textbf{B.} $QE//HF$ \quad \textbf{C.} $\widehat{FKI }= 23^{\circ}$ \quad \textbf{\textcolor{blue!75!black}{D. (Đ)}} $KF$ là tia phân giác của $\widehat{ QKI}$
\end{minipage}\hfill
\begin{minipage}[t]{0.39\textwidth}
\vspace{0pt}
\centering

\resizebox{\linewidth}{!}{%
\begin{tikzpicture}
            \coordinate (E) at (4, 3);
            \coordinate (D) at (3, 2);
            \coordinate (C) at (-6, 2);
            \coordinate (B) at (-2, 2);
            \coordinate (A) at (-1, 3);
            \coordinate (F) at (0, 0);
       
            % Vẽ tứ giác
            \draw[thick] (C) -- (B)--(D);
            \draw[thick] (A) -- (E);
            \draw[thick] (A) -- (B);
            \draw[thick] (B) --(F);
            % Ghi nhãn các điểm
            
            \node[above] at (A) {$Q$};
            \node[above ] at (B){$K$};
                      \node[above] at (C) {$H$};
                        \node[above] at (D) {$F$};
                         \node[above] at (E) {$E$};
                        \node[above] at (F) {$I$};
         
    \draw 
            (A) pic[draw,angle radius = 10] {angle = B--A--E} node[shift={(-20:20pt)}]{$ 120^\circ $}   
            ;    
            \draw 
            (B) pic[draw,angle radius = 10] {angle = F--B--A} pic[draw,angle radius = 12] {angle = F--B--A} node[shift={(-5:25pt)}]{$ 83^\circ $}   
            ;    
            \draw 
            (B) pic[draw,angle radius = 10] {angle = C--B--F}  pic[draw,angle radius = 8] {angle = C--B--F}  pic[draw,angle radius = 12] {angle = C--B--F} node[shift={(20:-20pt)}]{$ 157^\circ $}   
            ;  
        \end{tikzpicture}
}
\end{minipage}

\end{tcolorbox}
\vspace{4pt}


\begin{tcolorbox}[
    breakable,
    colback=white,
    colframe=blue!60!black,
    boxrule=0.8pt,
    arc=3pt,
    title={\textbf{Câu \#105} \quad \textbf{MÃ CÂU: \texttt{[L07\_C3\_B2\_19]}} \quad | \quad \textbf{Nhóm: \texttt{HINH\_097}} \quad | \quad \textbf{MapID: \texttt{7H3H2}} \hfill \textcolor{yellow!90!white}{[Hình gốc .pyc]}}
]
\textbf{Dạng toán:} Chứng minh hai đường thẳng song song 2 \hfill \textit{(Trắc nghiệm -- Thông hiểu)}
\vspace{4pt}\hrule\vspace{6pt}
\begin{minipage}[t]{0.58\textwidth}
\vspace{0pt}
Cho hình vẽ sau. Trong các khẳng định sau khẳng định nào sai?

\vspace{4pt}\noindent
\textbf{A.} $By// Cz $ \quad \textbf{\textcolor{blue!75!black}{B. (Đ)}} ${Ax}$ và ${Cz}$ không song song với nhau \quad \textbf{C.} $Ax //By $ \quad \textbf{D.} $Ax //Cz  $
\end{minipage}\hfill
\begin{minipage}[t]{0.39\textwidth}
\vspace{0pt}
\centering

\resizebox{\linewidth}{!}{%
\begin{tikzpicture}
            \coordinate (E) at (4, 1);
            \coordinate (D) at (3, 2);
            \coordinate (C) at (0, 0);
            \coordinate (B) at (-1, 1);
            \coordinate (A) at (-2, 2);
            \coordinate (F) at (5, 0);
             \coordinate (I) at (-3, 3);
                  \coordinate (M) at (-1, 0);
             \coordinate (N) at (1, -1);
            % Vẽ tứ giác
            \draw[thick] (C) -- (B);
            \draw[thick] (A) -- (D);
            \draw[thick] (N)--(C)--(B) -- (A)--(I);
            \draw[thick] (B) --(E);
             \draw[thick] (M)--(C) --(F);
            % Ghi nhãn các điểm
            
            \node[below left] at (A) {$A$};
            \node[below left] at (B){$B$};
                      \node[above right] at (C) {$C$};
                  \node[below] at (E) {$y$};
                         \node[below] at (D) {$x$};
                        \node[below] at (F) {$z$};
          
                    \node[below left] at (I) {$a$};
    \draw 
            (A) pic[draw,angle radius = 7] {angle = D--A--I} node[shift={(65:13pt)}]{$ 137^\circ $}   
            ;    
            \draw 
            (B) pic[draw,angle radius = 7] {angle = E--B--I} node[shift={(65:13pt)}]{$ 137^\circ $}   
            ;    
              \draw 
            (C) pic[draw,angle radius = 7] {angle = M--C--N} node[shift={(25: -20pt)}]{$ 137^\circ $}   
            ;  
        \end{tikzpicture}
}
\end{minipage}

\end{tcolorbox}
\vspace{4pt}


\begin{tcolorbox}[
    breakable,
    colback=white,
    colframe=blue!60!black,
    boxrule=0.8pt,
    arc=3pt,
    title={\textbf{Câu \#106} \quad \textbf{MÃ CÂU: \texttt{[L07\_C3\_B2\_20]}} \quad | \quad \textbf{Nhóm: \texttt{HINH\_099}} \quad | \quad \textbf{MapID: \texttt{7H3V2}} \hfill \textcolor{yellow!90!white}{[Hình gốc .pyc]}}
]
\textbf{Dạng toán:} Chứng minh các đường thẳng song song 4. \hfill \textit{(Trắc nghiệm -- Vận dụng)}
\vspace{4pt}\hrule\vspace{6pt}
\begin{minipage}[t]{0.58\textwidth}
\vspace{0pt}
Cho hình vẽ sau. Trong các khẳng định sau khẳng định nào sai?

\vspace{4pt}\noindent
\textbf{A.} $F K //N I $ \quad \textbf{\textcolor{blue!75!black}{B. (Đ)}} $\widehat {KFN}=124^{\circ} $ \quad \textbf{C.} $EM //N I  $ \quad \textbf{D.} $EM //F K $
\end{minipage}\hfill
\begin{minipage}[t]{0.39\textwidth}
\vspace{0pt}
\centering

\resizebox{\linewidth}{!}{%
\begin{tikzpicture}
            \coordinate (E) at (4, 3);
            \coordinate (D) at (3, 2);
            \coordinate (C) at (-6, 2);
            \coordinate (B) at (-2, 2);
            \coordinate (A) at (-1, 3);
            \coordinate (F) at (0, 0);
        \coordinate (M) at (4, 0);
            % Vẽ tứ giác
            \draw[thick] (C) -- (B);
            \draw[thick] (A) -- (E);
            \draw[thick] (A) -- (B);
            \draw[thick] (B) --(F);
             \draw[thick] (M) --(F);
            % Ghi nhãn các điểm
            
            \node[above] at (A) {$E$};
            \node[below ] at (B){$F$};
                      \node[above] at (C) {$K$};
                        \node[below] at (M) {$I$};
                         \node[above] at (E) {$M$};
                        \node[below] at (F) {$N$};
         
    \draw 
            (A) pic[draw,angle radius = 10] {angle = B--A--E} node[shift={(-20:20pt)}]{$ 124^\circ $}   
            ;    
            \draw 
            (B) pic[draw,angle radius = 10] {angle = F--B--A} pic[draw,angle radius = 12] {angle = F--B--A} node[shift={(-5:25pt)}]{$ 95^\circ $}   
            ;    
            \draw 
            (B) pic[draw,angle radius = 10] {angle = A--B--C}  node[shift={(-20:-25pt)}]{$ 124^\circ $}   
            ;  
              \draw 
            (F) pic[draw,angle radius = 10] {angle = M--F--B}  pic[draw,angle radius = 8] {angle = M--F--B}  pic[draw,angle radius = 12] {angle = M--F--B} node[shift={(30: 25pt)}]{$ 141^\circ $}   
            ;  
        \end{tikzpicture}
}
\end{minipage}

\end{tcolorbox}
\vspace{4pt}


\begin{tcolorbox}[
    breakable,
    colback=white,
    colframe=blue!60!black,
    boxrule=0.8pt,
    arc=3pt,
    title={\textbf{Câu \#107} \quad \textbf{MÃ CÂU: \texttt{[L07\_C3\_B2\_20]}} \quad | \quad \textbf{Nhóm: \texttt{HINH\_099}} \quad | \quad \textbf{MapID: \texttt{7H3V2}} \hfill \textcolor{yellow!90!white}{[Hình gốc .pyc]}}
]
\textbf{Dạng toán:} Chứng minh các đường thẳng song song 4. \hfill \textit{(Trắc nghiệm -- Vận dụng)}
\vspace{4pt}\hrule\vspace{6pt}
\begin{minipage}[t]{0.58\textwidth}
\vspace{0pt}
Cho hình vẽ sau. Trong các khẳng định sau khẳng định nào sai?

\vspace{4pt}\noindent
\textbf{A.} $AM //N I  $ \quad \textbf{\textcolor{blue!75!black}{B. (Đ)}} $\widehat {CBN}=136^{\circ} $ \quad \textbf{C.} $B C //N I $ \quad \textbf{D.} $AM //B C $
\end{minipage}\hfill
\begin{minipage}[t]{0.39\textwidth}
\vspace{0pt}
\centering

\resizebox{\linewidth}{!}{%
\begin{tikzpicture}
            \coordinate (E) at (4, 3);
            \coordinate (D) at (3, 2);
            \coordinate (C) at (-6, 2);
            \coordinate (B) at (-2, 2);
            \coordinate (A) at (-1, 3);
            \coordinate (F) at (0, 0);
        \coordinate (M) at (4, 0);
            % Vẽ tứ giác
            \draw[thick] (C) -- (B);
            \draw[thick] (A) -- (E);
            \draw[thick] (A) -- (B);
            \draw[thick] (B) --(F);
             \draw[thick] (M) --(F);
            % Ghi nhãn các điểm
            
            \node[above] at (A) {$A$};
            \node[below ] at (B){$B$};
                      \node[above] at (C) {$C$};
                        \node[below] at (M) {$I$};
                         \node[above] at (E) {$M$};
                        \node[below] at (F) {$N$};
         
    \draw 
            (A) pic[draw,angle radius = 10] {angle = B--A--E} node[shift={(-20:20pt)}]{$ 136^\circ $}   
            ;    
            \draw 
            (B) pic[draw,angle radius = 10] {angle = F--B--A} pic[draw,angle radius = 12] {angle = F--B--A} node[shift={(-5:25pt)}]{$ 89^\circ $}   
            ;    
            \draw 
            (B) pic[draw,angle radius = 10] {angle = A--B--C}  node[shift={(-20:-25pt)}]{$ 136^\circ $}   
            ;  
              \draw 
            (F) pic[draw,angle radius = 10] {angle = M--F--B}  pic[draw,angle radius = 8] {angle = M--F--B}  pic[draw,angle radius = 12] {angle = M--F--B} node[shift={(30: 25pt)}]{$ 135^\circ $}   
            ;  
        \end{tikzpicture}
}
\end{minipage}

\end{tcolorbox}
\vspace{4pt}


\begin{tcolorbox}[
    breakable,
    colback=white,
    colframe=blue!60!black,
    boxrule=0.8pt,
    arc=3pt,
    title={\textbf{Câu \#108} \quad \textbf{MÃ CÂU: \texttt{[L07\_C3\_B2\_22]}} \quad | \quad \textbf{Nhóm: \texttt{HINH\_100}} \quad | \quad \textbf{MapID: \texttt{7H3V2}} \hfill \textcolor{yellow!90!white}{[Hình gốc .pyc]}}
]
\textbf{Dạng toán:} Tính số đo góc 1 \hfill \textit{(Trắc nghiệm -- Vận dụng)}
\vspace{4pt}\hrule\vspace{6pt}
\begin{minipage}[t]{0.58\textwidth}
\vspace{0pt}
Cho hình vẽ sau biết ${Ja//Qb}$ tìm ${x}$ ?

\vspace{4pt}\noindent
\textbf{A.} $x=33^{\circ}  $ \quad \textbf{B.} $x=63^{\circ} $ \quad \textbf{\textcolor{blue!75!black}{C. (Đ)}} $x=53^{\circ} $ \quad \textbf{D.} $x=43^{\circ}$
\end{minipage}\hfill
\begin{minipage}[t]{0.39\textwidth}
\vspace{0pt}
\centering

\resizebox{\linewidth}{!}{%
\begin{tikzpicture}
            \coordinate (E) at (5, 3);
            \coordinate (C) at (0, 0);
            \coordinate (B) at (0.5, 3);
            \coordinate (A) at (-1, 2);
            \coordinate (F) at (-4, 0);
          
            % Vẽ tứ giác
            \draw[thick] (F)--(C) -- (A)--(B)--(E);
            % Ghi nhãn các điểm
            
            \node[above] at (A) {$I$};
            \node[above ] at (B){$J$};
                      \node[below] at (C) {$Q$};
                         \node[above] at (E) {$a$};
                        \node[below] at (F) {$b$};
    \draw 
            (A)  node[shift={(10:10pt)}]{$x^\circ $}   
            ;    
            \draw 
            (B) node[shift={(-50:10pt)}]{$ 164^\circ $}   
            ;    
            \draw 
            (C) node[shift={(-35:-15pt)}]{$ 37^\circ $}   
            ;  
            \fill[black] (A) circle (1.5pt) ; 
            \fill[black] (B) circle (1.5pt)  ;
             \fill[black] (C) circle (1.5pt)  ;
             
        \end{tikzpicture}
}
\end{minipage}

\end{tcolorbox}
\vspace{4pt}


\begin{tcolorbox}[
    breakable,
    colback=white,
    colframe=blue!60!black,
    boxrule=0.8pt,
    arc=3pt,
    title={\textbf{Câu \#109} \quad \textbf{MÃ CÂU: \texttt{[L07\_C3\_B2\_23]}} \quad | \quad \textbf{Nhóm: \texttt{HINH\_101}} \quad | \quad \textbf{MapID: \texttt{7H3V2}} \hfill \textcolor{yellow!90!white}{[Hình gốc .pyc]}}
]
\textbf{Dạng toán:} Tính số đo góc 2 \hfill \textit{(Trắc nghiệm -- Vận dụng)}
\vspace{4pt}\hrule\vspace{6pt}
\begin{minipage}[t]{0.58\textwidth}
\vspace{0pt}
Cho hình vẽ sau biết ${Ga//Eb}$ tìm ${x}$ ?

\vspace{4pt}\noindent
\textbf{A.} $x=74^{\circ} $ \quad \textbf{\textcolor{blue!75!black}{B. (Đ)}} $x=72^{\circ} $ \quad \textbf{C.} $x=68^{\circ}  $ \quad \textbf{D.} $x=76^{\circ}$
\end{minipage}\hfill
\begin{minipage}[t]{0.39\textwidth}
\vspace{0pt}
\centering

\resizebox{\linewidth}{!}{%
\begin{tikzpicture}
            \coordinate (E) at (3, 4);
            \coordinate (C) at (-1, 1);
            \coordinate (B) at (0.5, 2);
            \coordinate (A) at (0, 4);
            \coordinate (F) at (-3, 0);
           \coordinate (D) at (0, 0);
            % Vẽ tứ giác
            \draw[thick] (F)--(D)--(C) -- (B)--(A)--(E);
            % Ghi nhãn các điểm
              \node[above right] at (D) {$E$};
            \node[above left] at (A) {$G$};
            \node[above right ] at (B){$H$};
                      \node[above left] at (C) {$F$};
                         \node[above] at (E) {$a$};
                        \node[above] at (F) {$b$};
    \draw 
            (A)  node[shift={(-40:13pt)}]{$48^\circ $}   
            ;    
            \draw 
            (B) node[shift={(-50:-10pt)}]{$ x^\circ $}   
            ;    
            \draw 
            (C) node[shift={(-5:15pt)}]{$ 63^\circ $}   
            ;  
             \draw 
            (D) node[shift={(-21:-16pt)}]{$ 39^\circ $}   
            ;  
           \fill[black] (A) circle (1.5pt) ; 
            \fill[black] (B) circle (1.5pt)  ;
             \fill[black] (C) circle (1.5pt)  ;
              \fill[black] (D) circle (1.5pt)  ;
        \end{tikzpicture}
}
\end{minipage}

\end{tcolorbox}
\vspace{4pt}


\begin{tcolorbox}[
    breakable,
    colback=white,
    colframe=blue!60!black,
    boxrule=0.8pt,
    arc=3pt,
    title={\textbf{Câu \#110} \quad \textbf{MÃ CÂU: \texttt{[L07\_C3\_B2\_01]}} \quad | \quad \textbf{Nhóm: \texttt{HINH\_102}} \quad | \quad \textbf{MapID: \texttt{7H3N2}} \hfill \textcolor{yellow!90!white}{[Hình gốc .pyc]}}
]
\textbf{Dạng toán:} Nhận biết vị trí đặc biệt của các cặp góc 1 \hfill \textit{(Tự luận -- Nhận biết)}
\vspace{4pt}\hrule\vspace{6pt}
\begin{minipage}[t]{0.58\textwidth}
\vspace{0pt}
Cho hình sau:

a) Kể tên các góc so le trong.

b) Kể tên các góc đồng vị.

c) Kể tên các góc trong cùng phía.

d) Kể tên các góc so le ngoài.
\end{minipage}\hfill
\begin{minipage}[t]{0.39\textwidth}
\vspace{0pt}
\centering

\resizebox{\linewidth}{!}{%
\begin{tikzpicture}
            % Vẽ các điểm A, B, C, D
            \coordinate (A) at (-2, 0);
            \coordinate (C) at (-3, -2);
            \coordinate (B) at (2, 0);
            \coordinate (D) at (-1, 2);
            \coordinate (E) at (-4, 0);
            \coordinate (F) at (3, -2);
            \coordinate (M) at (1, 2);
            \coordinate (N) at (4, 0);
            
            \draw[thick] (E) -- (N);
            \draw[thick] (C) -- (D) ;
            \draw[thick] (M) -- (F) ;
            
            % Đánh dấu các điểm A, B, 
            \fill (B) circle (2pt) node[above] {$N$};
            \fill (A) circle (2pt) node[above] {$M$};
            \draw 
            (B) node[shift={(-25:15pt)}]{$4$}   
            ;
            \draw 
            (B) node[shift={(45:-15pt)}]{$3$}   
            ;
            \draw 
            (B) node[shift={(-50:-15pt)}]{$2$}   
            ;
            \draw 
            (B) node[shift={(50:15pt)}]{$1$}   
            ;
            \draw 
            (A)  node[shift={(-50:12pt)}]{$2$}   
            ;
            \draw 
            (A) node[shift={(50:-13pt)}]{$1$}   
            ;
            \draw 
            (A)  node[shift={(-50:-15pt)}]{$4$}   
            ;
            \draw 
            (A) node[shift={(50:15pt)}]{$3$}   
            ;
            
        \end{tikzpicture}
}
\end{minipage}

\end{tcolorbox}
\vspace{4pt}


\begin{tcolorbox}[
    breakable,
    colback=white,
    colframe=blue!60!black,
    boxrule=0.8pt,
    arc=3pt,
    title={\textbf{Câu \#111} \quad \textbf{MÃ CÂU: \texttt{[L07\_C3\_B2\_02]}} \quad | \quad \textbf{Nhóm: \texttt{HINH\_103}} \quad | \quad \textbf{MapID: \texttt{7H3N2}} \hfill \textcolor{yellow!90!white}{[Hình gốc .pyc]}}
]
\textbf{Dạng toán:} Nhận biết vị trí đặc biệt của các cặp góc 2 \hfill \textit{(Tự luận -- Nhận biết)}
\vspace{4pt}\hrule\vspace{6pt}
\begin{minipage}[t]{0.58\textwidth}
\vspace{0pt}
Xem hình vẽ bên và điền vào chỗ trống .....

a) $\widehat{ MGF}$ và $\widehat{ GFH}$ là ...

b) $\widehat{ FGM}$ và $\widehat{ IEF}$ là ...

c) $\widehat{ GMF}$ và $\widehat{ EIF}$ là ...

 d) $\widehat{ FGH}$ và $\widehat{ GFM}$ là ...
\end{minipage}\hfill
\begin{minipage}[t]{0.39\textwidth}
\vspace{0pt}
\centering

\resizebox{\linewidth}{!}{%
\begin{tikzpicture}
            \coordinate (R) at (1, 3);
            \coordinate (P) at (0, 0);
            \coordinate (O) at (3, 3);
         \coordinate (I) at (6, 0);
            \coordinate (N) at (2, 2);
              \coordinate (T) at (4, 2);
            % Vẽ tứ giác
             \draw[thick] (O) -- (P);
              \draw[thick] (N) -- (T);
            \draw[thick] (P) -- (R) -- (O) -- (I)--(P);
    % Ghi nhãn các điểm
            
            \node[above left] at (R) {$H$};
            \node[above] at (O) {$F$};
            \node[below right] at (I) {$M$};
            \node[below left] at (P) {$G$};
                 \node[above left] at (N) {$E$};
                  \node[above right] at (T) {$I$};  
         \fill [black] (R) circle (1pt);
        \fill [black] (O) circle (1pt);
        \fill [black] (I) circle (1pt);
        \fill [black] (P) circle (1pt);
         \fill [black] (N) circle (1pt);
        \fill [black] (T) circle (1pt);
    \end{tikzpicture}
}
\end{minipage}

\end{tcolorbox}
\vspace{4pt}


\begin{tcolorbox}[
    breakable,
    colback=white,
    colframe=blue!60!black,
    boxrule=0.8pt,
    arc=3pt,
    title={\textbf{Câu \#112} \quad \textbf{MÃ CÂU: \texttt{[L07\_C3\_B2\_03]}} \quad | \quad \textbf{Nhóm: \texttt{HINH\_104}} \quad | \quad \textbf{MapID: \texttt{7H3H2}} \hfill \textcolor{yellow!90!white}{[Hình gốc .pyc]}}
]
\textbf{Dạng toán:} Tính số đo góc 1 \hfill \textit{(Tự luận -- Thông hiểu)}
\vspace{4pt}\hrule\vspace{6pt}
\begin{minipage}[t]{0.58\textwidth}
\vspace{0pt}
Cho hình vẽ 

a) Kể tên các góc bằng nhau có trong hình.

b) Cho  $\widehat { P_{4} }= \widehat{ Q_{1} }=58^{\circ}$ tính số đo các góc còn lại.
\end{minipage}\hfill
\begin{minipage}[t]{0.39\textwidth}
\vspace{0pt}
\centering

\resizebox{\linewidth}{!}{%
\begin{tikzpicture}
            \coordinate (E) at (1.14, 3.06);
            \coordinate (D) at (4.46, 0);
            \coordinate (C) at (-0.6, 0);
            \coordinate (B) at (5, 2);
            \coordinate (A) at (0, 2);
            \coordinate (F) at (3.6, -1.44);
            \coordinate (H) at (2.81, 0);
            \coordinate (G) at (1.72, 2);
            % Vẽ tứ giác
            \draw[thick] (C) -- (H)--(D);
            \draw[thick] (A) -- (G) -- (B);
            \draw[thick] (E) -- (G) -- (H)--(F);
            % Ghi nhãn các điểm
            
            \node[above] at (G) {$P$};
            \node[above ] at (H){$Q$};
            \draw 
            (G) pic[draw,angle radius = 8] {angle = E--G--A} node[shift={(-25:15pt)}]{$1$}  ;
            \draw (G) node[shift={(30:11pt)}]{$2$};
            \draw  (G) node[shift={(50:-15pt)}]{$3$}  ;
            \draw (G)  node[shift={(-25:-15pt)}]{$4$};  
            \draw (H) pic[draw,angle radius = 10] {angle = F--H--D};
            \draw (H) node[shift={(-25:15pt)}]{$1$};
            \draw (H) node[shift={(40:15pt)}]{$2$}  ;
            \draw  (H) node[shift={(50:-13pt)}]{$3$} ;
            \draw (H) node[shift={(-25:-15pt)}]{$4$}  ;     
        \end{tikzpicture}
}
\end{minipage}

\vspace{3pt}\hrule\vspace{3pt}\footnotesize\textcolor{gray!80!black}{\textbf{Các mã dùng chung hình (HINH\_104):} \texttt{L07\_C3\_B3\_01}}
\end{tcolorbox}
\vspace{4pt}


\begin{tcolorbox}[
    breakable,
    colback=white,
    colframe=blue!60!black,
    boxrule=0.8pt,
    arc=3pt,
    title={\textbf{Câu \#113} \quad \textbf{MÃ CÂU: \texttt{[L07\_C3\_B2\_04]}} \quad | \quad \textbf{Nhóm: \texttt{HINH\_105}} \quad | \quad \textbf{MapID: \texttt{7H3H2}} \hfill \textcolor{yellow!90!white}{[Hình gốc .pyc]}}
]
\textbf{Dạng toán:} Tính số đo góc 2. \hfill \textit{(Tự luận -- Thông hiểu)}
\vspace{4pt}\hrule\vspace{6pt}
\begin{minipage}[t]{0.58\textwidth}
\vspace{0pt}
Cho hình vẽ 

a) Kể tên các góc bằng nhau có trong hình.

b) Cho  $\widehat { A _{4} }= \widehat{ B _{4} }=40^{\circ}$ tính số đo các góc còn lại.
\end{minipage}\hfill
\begin{minipage}[t]{0.39\textwidth}
\vspace{0pt}
\centering

\resizebox{\linewidth}{!}{%
\begin{tikzpicture}
            \coordinate (E) at (1.14, 3.06);
            \coordinate (D) at (4.46, 0);
            \coordinate (C) at (-0.6, 0);
            \coordinate (B) at (5, 2);
            \coordinate (A) at (0, 2);
            \coordinate (F) at (3.6, -1.44);
            \coordinate (H) at (2.81, 0);
            \coordinate (G) at (1.72, 2);
            % Vẽ tứ giác
            \draw[thick] (C) -- (H)--(D);
            \draw[thick] (A) -- (G) -- (B);
            \draw[thick] (E) -- (G) -- (H)--(F);
            % Ghi nhãn các điểm
            
            \node[above] at (G) {$A $};
            \node[above ] at (H){$B $};
            \draw 
            (G) pic[draw,angle radius = 8] {angle = E--G--A} node[shift={(-25:15pt)}]{$1$}  ;
            \draw (G) node[shift={(30:11pt)}]{$2$};
            \draw  (G) node[shift={(50:-15pt)}]{$3$}  ;
            \draw (G)  node[shift={(-25:-15pt)}]{$4$};  
            \draw (H) pic[draw,angle radius = 10] {angle = E--H--C};
            \draw (H) node[shift={(-25:15pt)}]{$1$};
            \draw (H) node[shift={(40:15pt)}]{$2$}  ;
            \draw  (H) node[shift={(50:-13pt)}]{$3$} ;
            \draw (H) node[shift={(-25:-15pt)}]{$4$}  ;     
        \end{tikzpicture}
}
\end{minipage}

\vspace{3pt}\hrule\vspace{3pt}\footnotesize\textcolor{gray!80!black}{\textbf{Các mã dùng chung hình (HINH\_105):} \texttt{L07\_C3\_B3\_02}}
\end{tcolorbox}
\vspace{4pt}


\begin{tcolorbox}[
    breakable,
    colback=white,
    colframe=blue!60!black,
    boxrule=0.8pt,
    arc=3pt,
    title={\textbf{Câu \#114} \quad \textbf{MÃ CÂU: \texttt{[L07\_C3\_B2\_05]}} \quad | \quad \textbf{Nhóm: \texttt{HINH\_106}} \quad | \quad \textbf{MapID: \texttt{7H3H2}} \hfill \textcolor{yellow!90!white}{[Hình gốc .pyc]}}
]
\textbf{Dạng toán:} Tính số đo góc 3. \hfill \textit{(Tự luận -- Thông hiểu)}
\vspace{4pt}\hrule\vspace{6pt}
\begin{minipage}[t]{0.58\textwidth}
\vspace{0pt}
Cho hình vẽ 

a) Kể tên các góc bằng nhau có trong hình.

b) Cho  $\widehat { E_{1} }= \widehat{ F_{4} }=41^{\circ}$ tính số đo các góc còn lại.
\end{minipage}\hfill
\begin{minipage}[t]{0.39\textwidth}
\vspace{0pt}
\centering

\resizebox{\linewidth}{!}{%
\begin{tikzpicture}
            \coordinate (E) at (1.14, 3.06);
            \coordinate (D) at (4.46, 0);
            \coordinate (C) at (-0.6, 0);
            \coordinate (B) at (5, 2);
            \coordinate (A) at (0, 2);
            \coordinate (F) at (3.6, -1.44);
            \coordinate (H) at (2.81, 0);
            \coordinate (G) at (1.72, 2);
            % Vẽ tứ giác
            \draw[thick] (C) -- (H)--(D);
            \draw[thick] (A) -- (G) -- (B);
            \draw[thick] (E) -- (G) -- (H)--(F);
            % Ghi nhãn các điểm
            
            \node[above] at (G) {$E$};
            \node[above ] at (H){$F$};
            \draw 
            (G) pic[draw,angle radius = 8] {angle = H--G--B} node[shift={(-25:15pt)}]{$1$}  ;
            \draw (G) node[shift={(30:11pt)}]{$2$};
            \draw  (G) node[shift={(50:-15pt)}]{$3$}  ;
            \draw (G)  node[shift={(-25:-15pt)}]{$4$};  
            \draw (H) pic[draw,angle radius = 10] {angle = E--H--C};
            \draw (H) node[shift={(-25:15pt)}]{$1$};
            \draw (H) node[shift={(40:15pt)}]{$2$}  ;
            \draw  (H) node[shift={(50:-13pt)}]{$3$} ;
            \draw (H) node[shift={(-25:-15pt)}]{$4$}  ;     
        \end{tikzpicture}
}
\end{minipage}

\vspace{3pt}\hrule\vspace{3pt}\footnotesize\textcolor{gray!80!black}{\textbf{Các mã dùng chung hình (HINH\_106):} \texttt{L07\_C3\_B3\_03}}
\end{tcolorbox}
\vspace{4pt}


\begin{tcolorbox}[
    breakable,
    colback=white,
    colframe=blue!60!black,
    boxrule=0.8pt,
    arc=3pt,
    title={\textbf{Câu \#115} \quad \textbf{MÃ CÂU: \texttt{[L07\_C3\_B2\_06]}} \quad | \quad \textbf{Nhóm: \texttt{HINH\_091}} \quad | \quad \textbf{MapID: \texttt{7H3H2}} \hfill \textcolor{yellow!90!white}{[Hình gốc .pyc]}}
]
\textbf{Dạng toán:} Tính số đo góc 4. \hfill \textit{(Tự luận -- Thông hiểu)}
\vspace{4pt}\hrule\vspace{6pt}
\begin{minipage}[t]{0.58\textwidth}
\vspace{0pt}
Cho hình vẽ sau biết $\widehat{ J_{4}} =160^{\circ}$, $\widehat{ I_{2}} =64^{\circ}$

 Tính số đo của các góc còn lại.
\end{minipage}\hfill
\begin{minipage}[t]{0.39\textwidth}
\vspace{0pt}
\centering

\resizebox{\linewidth}{!}{%
\begin{tikzpicture}
            % Vẽ các điểm A, B, C, D
            \coordinate (A) at (-2, 0);
            \coordinate (C) at (-3, -2);
            \coordinate (B) at (2, 0);
            \coordinate (D) at (-1, 2);
            \coordinate (E) at (-4, 0);
            \coordinate (F) at (3, -2);
            \coordinate (M) at (1, 2);
            \coordinate (N) at (4, 0);
            
            \draw[thick] (E) -- (N);
            \draw[thick] (C) -- (D) ;
            \draw[thick] (M) -- (F) ;
            
            % Đánh dấu các điểm A, B, 
            \fill (B) circle (2pt) node[above] {$J$};
            \fill (A) circle (2pt) node[above] {$I$};
            \draw 
            (B) node[shift={(-25:15pt)}]{$1$}   
            ;
            \draw 
            (B) node[shift={(45:-15pt)}]{$2$}   
            ;
            \draw 
            (B) node[shift={(-50:-15pt)}]{$3$}   
            ;
            \draw 
            (B) node[shift={(50:15pt)}]{$4$}   
            ;
            \draw 
            (A)  node[shift={(-50:12pt)}]{$1$}   
            ;
            \draw 
            (A) node[shift={(50:-13pt)}]{$2$}   
            ;
            \draw 
            (A)  node[shift={(-50:-15pt)}]{$3$}   
            ;
            \draw 
            (A) node[shift={(50:15pt)}]{$4$}   
            ;
            
        \end{tikzpicture}
}
\end{minipage}

\vspace{3pt}\hrule\vspace{3pt}\footnotesize\textcolor{gray!80!black}{\textbf{Các mã dùng chung hình (HINH\_091):} \texttt{L07\_C3\_B2\_11, L07\_C3\_B3\_04}}
\end{tcolorbox}
\vspace{4pt}


\begin{tcolorbox}[
    breakable,
    colback=white,
    colframe=blue!60!black,
    boxrule=0.8pt,
    arc=3pt,
    title={\textbf{Câu \#116} \quad \textbf{MÃ CÂU: \texttt{[L07\_C3\_B2\_07]}} \quad | \quad \textbf{Nhóm: \texttt{HINH\_092}} \quad | \quad \textbf{MapID: \texttt{7H3N2}} \hfill \textcolor{yellow!90!white}{[Hình gốc .pyc]}}
]
\textbf{Dạng toán:} Nhận biết hai đường thẳng song song. \hfill \textit{(Tự luận -- Nhận biết)}
\vspace{4pt}\hrule\vspace{6pt}
\begin{minipage}[t]{0.58\textwidth}
\vspace{0pt}
Cho hình vẽ biết $\widehat { A_{4} }=\widehat{ B_{1} }=45^{\circ}$ hai đường thẳng ${a}$ và ${b}$ có song song với nhau không? tính số đo các góc còn lại.
\end{minipage}\hfill
\begin{minipage}[t]{0.39\textwidth}
\vspace{0pt}
\centering

\resizebox{\linewidth}{!}{%
\begin{tikzpicture}
            \coordinate (E) at (1.14, 3.06);
            \coordinate (D) at (4.46, 0);
            \coordinate (C) at (-0.6, 0);
            \coordinate (B) at (5, 2);
            \coordinate (A) at (0, 2);
            \coordinate (F) at (3.6, -1.44);
            \coordinate (H) at (2.81, 0);
            \coordinate (G) at (1.72, 2);
            % Vẽ tứ giác
            \draw[thick] (C) -- (H)--(D);
            \draw[thick] (A) -- (G) -- (B);
            \draw[thick] (E) -- (G) -- (H)--(F);
            % Ghi nhãn các điểm
            
            \node[above] at (G) {$A$};
            \node[above ] at (H){$B$};
                \node[above] at (A) {$a$};
                        \node[above] at (C) {$b$};
            \draw 
            (G) pic[draw,angle radius = 8] {angle = E--G--A} node[shift={(-25:15pt)}]{$1$}  ;
            \draw (G) node[shift={(30:11pt)}]{$2$};
            \draw  (G) node[shift={(50:-15pt)}]{$3$}  ;
            \draw (G)  node[shift={(-25:-15pt)}]{$4$};  
            \draw (H) pic[draw,angle radius = 10] {angle = F--H--D};
            \draw (H) node[shift={(-25:15pt)}]{$1$};
            \draw (H) node[shift={(40:15pt)}]{$2$}  ;
            \draw  (H) node[shift={(50:-13pt)}]{$3$} ;
            \draw (H) node[shift={(-25:-15pt)}]{$4$}  ;     
        \end{tikzpicture}
}
\end{minipage}

\vspace{3pt}\hrule\vspace{3pt}\footnotesize\textcolor{gray!80!black}{\textbf{Các mã dùng chung hình (HINH\_092):} \texttt{L07\_C3\_B2\_12, L07\_C3\_B3\_09, L07\_C3\_B3\_05}}
\end{tcolorbox}
\vspace{4pt}


\begin{tcolorbox}[
    breakable,
    colback=white,
    colframe=blue!60!black,
    boxrule=0.8pt,
    arc=3pt,
    title={\textbf{Câu \#117} \quad \textbf{MÃ CÂU: \texttt{[L07\_C3\_B2\_08]}} \quad | \quad \textbf{Nhóm: \texttt{HINH\_093}} \quad | \quad \textbf{MapID: \texttt{7H3H2}} \hfill \textcolor{yellow!90!white}{[Hình gốc .pyc]}}
]
\textbf{Dạng toán:} Chứng minh hai đường thẳng song song dựa vào 2 góc đồng vị \hfill \textit{(Tự luận -- Thông hiểu)}
\vspace{4pt}\hrule\vspace{6pt}
\begin{minipage}[t]{0.58\textwidth}
\vspace{0pt}
Cho hình vẽ biết $\widehat { A _{4} }= \widehat{ B _{4} }=42^{\circ}$ hai đường thẳng ${y}$ và ${x}$ có song song với nhau không? tính số đo các góc còn lại.
\end{minipage}\hfill
\begin{minipage}[t]{0.39\textwidth}
\vspace{0pt}
\centering

\resizebox{\linewidth}{!}{%
\begin{tikzpicture}
            \coordinate (E) at (1.14, 3.06);
            \coordinate (D) at (4.46, 0);
            \coordinate (C) at (-0.6, 0);
            \coordinate (B) at (5, 2);
            \coordinate (A) at (0, 2);
            \coordinate (F) at (3.6, -1.44);
            \coordinate (H) at (2.81, 0);
            \coordinate (G) at (1.72, 2);
            % Vẽ tứ giác
            \draw[thick] (C) -- (H)--(D);
            \draw[thick] (A) -- (G) -- (B);
            \draw[thick] (E) -- (G) -- (H)--(F);
            % Ghi nhãn các điểm
            
            \node[above] at (G) {$A $};
            \node[above ] at (H){$B $};
                \node[above] at (A) {$y$};
                        \node[above] at (C) {$x$};
            \draw 
            (G) pic[draw,angle radius = 8] {angle = E--G--A} node[shift={(-25:15pt)}]{$1$}  ;
            \draw (G) node[shift={(30:11pt)}]{$2$};
            \draw  (G) node[shift={(50:-15pt)}]{$3$}  ;
            \draw (G)  node[shift={(-25:-15pt)}]{$4$};  
            \draw (H) pic[draw,angle radius = 10] {angle = E--H--C};
            \draw (H) node[shift={(-25:15pt)}]{$1$};
            \draw (H) node[shift={(40:15pt)}]{$2$}  ;
            \draw  (H) node[shift={(50:-13pt)}]{$3$} ;
            \draw (H) node[shift={(-25:-15pt)}]{$4$}  ;     
        \end{tikzpicture}
}
\end{minipage}

\vspace{3pt}\hrule\vspace{3pt}\footnotesize\textcolor{gray!80!black}{\textbf{Các mã dùng chung hình (HINH\_093):} \texttt{L07\_C3\_B2\_13, L07\_C3\_B3\_10, L07\_C3\_B3\_06}}
\end{tcolorbox}
\vspace{4pt}


\begin{tcolorbox}[
    breakable,
    colback=white,
    colframe=blue!60!black,
    boxrule=0.8pt,
    arc=3pt,
    title={\textbf{Câu \#118} \quad \textbf{MÃ CÂU: \texttt{[L07\_C3\_B2\_09]}} \quad | \quad \textbf{Nhóm: \texttt{HINH\_094}} \quad | \quad \textbf{MapID: \texttt{7H3N2}} \hfill \textcolor{yellow!90!white}{[Hình gốc .pyc]}}
]
\textbf{Dạng toán:} Nhận biết 2 đường thẳng song song dựa vào 2 góc so le trong. \hfill \textit{(Tự luận -- Nhận biết)}
\vspace{4pt}\hrule\vspace{6pt}
\begin{minipage}[t]{0.58\textwidth}
\vspace{0pt}
Cho hình vẽ biết $\widehat { I_{1} }= \widehat{ J_{4} }=59^{\circ}$, hỏi hai đường thẳng ${m }$ và ${n }$ có song song với nhau không? tính số đo các góc còn lại.
\end{minipage}\hfill
\begin{minipage}[t]{0.39\textwidth}
\vspace{0pt}
\centering

\resizebox{\linewidth}{!}{%
\begin{tikzpicture}
            \coordinate (E) at (1.14, 3.06);
            \coordinate (D) at (4.46, 0);
            \coordinate (C) at (-0.6, 0);
            \coordinate (B) at (5, 2);
            \coordinate (A) at (0, 2);
            \coordinate (F) at (3.6, -1.44);
            \coordinate (H) at (2.81, 0);
            \coordinate (G) at (1.72, 2);
            % Vẽ tứ giác
            \draw[thick] (C) -- (H)--(D);
            \draw[thick] (A) -- (G) -- (B);
            \draw[thick] (E) -- (G) -- (H)--(F);
            % Ghi nhãn các điểm
            
            \node[above] at (G) {$I$};
            \node[above ] at (H){$J$};
                      \node[above] at (A) {$m $};
                        \node[above] at (C) {$n $};
            \draw 
            (G) pic[draw,angle radius = 8] {angle = H--G--B} node[shift={(-25:15pt)}]{$1$}  ;
            \draw (G) node[shift={(30:11pt)}]{$2$};
            \draw  (G) node[shift={(50:-15pt)}]{$3$}  ;
            \draw (G)  node[shift={(-25:-15pt)}]{$4$};  
            \draw (H) pic[draw,angle radius = 10] {angle = E--H--C};
            \draw (H) node[shift={(-25:15pt)}]{$1$};
            \draw (H) node[shift={(40:15pt)}]{$2$}  ;
            \draw  (H) node[shift={(50:-13pt)}]{$3$} ;
            \draw (H) node[shift={(-25:-15pt)}]{$4$}  ;     
        \end{tikzpicture}
}
\end{minipage}

\vspace{3pt}\hrule\vspace{3pt}\footnotesize\textcolor{gray!80!black}{\textbf{Các mã dùng chung hình (HINH\_094):} \texttt{L07\_C3\_B2\_14, L07\_C3\_B3\_11, L07\_C3\_B3\_07}}
\end{tcolorbox}
\vspace{4pt}


\begin{tcolorbox}[
    breakable,
    colback=white,
    colframe=blue!60!black,
    boxrule=0.8pt,
    arc=3pt,
    title={\textbf{Câu \#119} \quad \textbf{MÃ CÂU: \texttt{[L07\_C3\_B2\_10]}} \quad | \quad \textbf{Nhóm: \texttt{HINH\_095}} \quad | \quad \textbf{MapID: \texttt{7H3H2}} \hfill \textcolor{yellow!90!white}{[Hình gốc .pyc]}}
]
\textbf{Dạng toán:} Thông hiểu 2 đường thẳng song song \hfill \textit{(Tự luận -- Thông hiểu)}
\vspace{4pt}\hrule\vspace{6pt}
\begin{minipage}[t]{0.58\textwidth}
\vspace{0pt}
Cho hình vẽ biết $\widehat { A_{1} }=59^{\circ}$, $\widehat { B_{3} }=121^{\circ}$ hỏi hai đường thẳng ${x}$ và ${y}$ có song song với nhau không? tính số đo các góc còn lại.
\end{minipage}\hfill
\begin{minipage}[t]{0.39\textwidth}
\vspace{0pt}
\centering

\resizebox{\linewidth}{!}{%
\begin{tikzpicture}
            \coordinate (E) at (1.14, 3.06);
            \coordinate (D) at (4.46, 0);
            \coordinate (C) at (-0.6, 0);
            \coordinate (B) at (5, 2);
            \coordinate (A) at (0, 2);
            \coordinate (F) at (3.6, -1.44);
            \coordinate (H) at (2.81, 0);
            \coordinate (G) at (1.72, 2);
            % Vẽ tứ giác
            \draw[thick] (C) -- (H)--(D);
            \draw[thick] (A) -- (G) -- (B);
            \draw[thick] (E) -- (G) -- (H)--(F);
            % Ghi nhãn các điểm
            
            \node[above] at (G) {$A$};
            \node[above ] at (H){$B$};
                      \node[above] at (A) {$x$};
                        \node[above] at (C) {$y$};
            \draw 
            (G) node[shift={(-25:15pt)}]{$1$}  ;
            \draw (G) node[shift={(30:11pt)}]{$2$};
            \draw  (G) node[shift={(50:-15pt)}]{$3$}  ;
            \draw (G)  node[shift={(-25:-15pt)}]{$4$};  
           
            \draw (H) node[shift={(-25:15pt)}]{$1$};
            \draw (H) node[shift={(40:15pt)}]{$2$}  ;
            \draw  (H) node[shift={(50:-13pt)}]{$3$} ;
            \draw (H) node[shift={(-25:-15pt)}]{$4$}  ;     
        \end{tikzpicture}
}
\end{minipage}

\vspace{3pt}\hrule\vspace{3pt}\footnotesize\textcolor{gray!80!black}{\textbf{Các mã dùng chung hình (HINH\_095):} \texttt{L07\_C3\_B2\_15, L07\_C3\_B2\_16, L07\_C3\_B2\_17, L07\_C3\_B3\_12, L07\_C3\_B3\_13, L07\_C3\_B3\_14, L07\_C3\_B3\_08}}
\end{tcolorbox}
\vspace{4pt}


\begin{tcolorbox}[
    breakable,
    colback=white,
    colframe=blue!60!black,
    boxrule=0.8pt,
    arc=3pt,
    title={\textbf{Câu \#120} \quad \textbf{MÃ CÂU: \texttt{[L07\_C3\_B3\_09]}} \quad | \quad \textbf{Nhóm: \texttt{HINH\_092}} \quad | \quad \textbf{MapID: \texttt{7H3N3}} \hfill \textcolor{yellow!90!white}{[Hình gốc .pyc]}}
]
\textbf{Dạng toán:} Nhận biết hai đường thẳng song song \hfill \textit{(Trắc nghiệm -- Nhận biết)}
\vspace{4pt}\hrule\vspace{6pt}
\begin{minipage}[t]{0.58\textwidth}
\vspace{0pt}
Cho hình vẽ sau biết $\widehat { E_{4} }=\widehat{ F_{1} }=50^{\circ}$. Chọn khẳng sai trong các khẳng định sau?

\vspace{4pt}\noindent
\textbf{A.} $\widehat{ E_{1} }=50^{\circ}$ \quad \textbf{B.} $x//y$ \quad \textbf{C.} $\widehat { E_{2} } =130^{\circ} $ \quad \textbf{\textcolor{blue!75!black}{D. (Đ)}} $\widehat{ F_{4} } = 51^{\circ}$
\end{minipage}\hfill
\begin{minipage}[t]{0.39\textwidth}
\vspace{0pt}
\centering

\resizebox{\linewidth}{!}{%
\begin{tikzpicture}
            \coordinate (E) at (1.14, 3.06);
            \coordinate (D) at (4.46, 0);
            \coordinate (C) at (-0.6, 0);
            \coordinate (B) at (5, 2);
            \coordinate (A) at (0, 2);
            \coordinate (F) at (3.6, -1.44);
            \coordinate (H) at (2.81, 0);
            \coordinate (G) at (1.72, 2);
            % Vẽ tứ giác
            \draw[thick] (C) -- (H)--(D);
            \draw[thick] (A) -- (G) -- (B);
            \draw[thick] (E) -- (G) -- (H)--(F);
            % Ghi nhãn các điểm
            
            \node[above] at (G) {$E$};
            \node[above ] at (H){$F$};
                \node[above] at (A) {$x$};
                        \node[above] at (C) {$y$};
            \draw 
            (G) pic[draw,angle radius = 8] {angle = E--G--A} node[shift={(-25:15pt)}]{$1$}  ;
            \draw (G) node[shift={(30:11pt)}]{$2$};
            \draw  (G) node[shift={(50:-15pt)}]{$3$}  ;
            \draw (G)  node[shift={(-25:-15pt)}]{$4$};  
            \draw (H) pic[draw,angle radius = 10] {angle = F--H--D};
            \draw (H) node[shift={(-25:15pt)}]{$1$};
            \draw (H) node[shift={(40:15pt)}]{$2$}  ;
            \draw  (H) node[shift={(50:-13pt)}]{$3$} ;
            \draw (H) node[shift={(-25:-15pt)}]{$4$}  ;     
        \end{tikzpicture}
}
\end{minipage}

\vspace{3pt}\hrule\vspace{3pt}\footnotesize\textcolor{gray!80!black}{\textbf{Các mã dùng chung hình (HINH\_092):} \texttt{L07\_C3\_B2\_12, L07\_C3\_B2\_07, L07\_C3\_B3\_05}}
\end{tcolorbox}
\vspace{4pt}


\begin{tcolorbox}[
    breakable,
    colback=white,
    colframe=blue!60!black,
    boxrule=0.8pt,
    arc=3pt,
    title={\textbf{Câu \#121} \quad \textbf{MÃ CÂU: \texttt{[L07\_C3\_B3\_10]}} \quad | \quad \textbf{Nhóm: \texttt{HINH\_093}} \quad | \quad \textbf{MapID: \texttt{7H3H3}} \hfill \textcolor{yellow!90!white}{[Hình gốc .pyc]}}
]
\textbf{Dạng toán:} Thông hiểu hai đường thẳng song song dựa vào 2 góc đồng vị \hfill \textit{(Trắc nghiệm -- Thông hiểu)}
\vspace{4pt}\hrule\vspace{6pt}
\begin{minipage}[t]{0.58\textwidth}
\vspace{0pt}
Cho hình vẽ biết $\widehat { A_{4} }= \widehat{ B_{4} }=58^{\circ}$. Trong các khẳng định sau khẳng định nào sai.

\vspace{4pt}\noindent
\textbf{\textcolor{blue!75!black}{A. (Đ)}} $\widehat{ B_{3} } =102^{\circ}$ \quad \textbf{B.} $\widehat{  A_{2} }=122^{\circ}$ \quad \textbf{C.} $ \widehat{ B_{4} }=58^{\circ}$ \quad \textbf{D.} $a //b$
\end{minipage}\hfill
\begin{minipage}[t]{0.39\textwidth}
\vspace{0pt}
\centering

\resizebox{\linewidth}{!}{%
\begin{tikzpicture}
            \coordinate (E) at (1.14, 3.06);
            \coordinate (D) at (4.46, 0);
            \coordinate (C) at (-0.6, 0);
            \coordinate (B) at (5, 2);
            \coordinate (A) at (0, 2);
            \coordinate (F) at (3.6, -1.44);
            \coordinate (H) at (2.81, 0);
            \coordinate (G) at (1.72, 2);
            % Vẽ tứ giác
            \draw[thick] (C) -- (H)--(D);
            \draw[thick] (A) -- (G) -- (B);
            \draw[thick] (E) -- (G) -- (H)--(F);
            % Ghi nhãn các điểm
            
            \node[above] at (G) {$A$};
            \node[above ] at (H){$B$};
                \node[above] at (A) {$a$};
                        \node[above] at (C) {$b$};
            \draw 
            (G) pic[draw,angle radius = 8] {angle = E--G--A} node[shift={(-25:15pt)}]{$1$}  ;
            \draw (G) node[shift={(30:11pt)}]{$2$};
            \draw  (G) node[shift={(50:-15pt)}]{$3$}  ;
            \draw (G)  node[shift={(-25:-15pt)}]{$4$};  
            \draw (H) pic[draw,angle radius = 10] {angle = E--H--C};
            \draw (H) node[shift={(-25:15pt)}]{$1$};
            \draw (H) node[shift={(40:15pt)}]{$2$}  ;
            \draw  (H) node[shift={(50:-13pt)}]{$3$} ;
            \draw (H) node[shift={(-25:-15pt)}]{$4$}  ;     
        \end{tikzpicture}
}
\end{minipage}

\vspace{3pt}\hrule\vspace{3pt}\footnotesize\textcolor{gray!80!black}{\textbf{Các mã dùng chung hình (HINH\_093):} \texttt{L07\_C3\_B2\_13, L07\_C3\_B2\_08, L07\_C3\_B3\_06}}
\end{tcolorbox}
\vspace{4pt}


\begin{tcolorbox}[
    breakable,
    colback=white,
    colframe=blue!60!black,
    boxrule=0.8pt,
    arc=3pt,
    title={\textbf{Câu \#122} \quad \textbf{MÃ CÂU: \texttt{[L07\_C3\_B3\_11]}} \quad | \quad \textbf{Nhóm: \texttt{HINH\_094}} \quad | \quad \textbf{MapID: \texttt{7H3H3}} \hfill \textcolor{yellow!90!white}{[Hình gốc .pyc]}}
]
\textbf{Dạng toán:} Thông hiểu 2 đường thẳng song song \hfill \textit{(Trắc nghiệm -- Thông hiểu)}
\vspace{4pt}\hrule\vspace{6pt}
\begin{minipage}[t]{0.58\textwidth}
\vspace{0pt}
Cho hình vẽ biết $\widehat { A _{1} }= \widehat{ B _{4} }=56^{\circ}$. Chọn khẳng định sai trong các khẳng định sau.

\vspace{4pt}\noindent
\textbf{\textcolor{blue!75!black}{A. (Đ)}} $\widehat{ A _{3} }=\widehat{  B _{2} }=114^{\circ}$ \quad \textbf{B.} $x //y$ \quad \textbf{C.} $\widehat { A _{4} }=\widehat{ B _{1} }= 56^{\circ}$ \quad \textbf{D.} $\widehat { A _{2} }= \widehat{ B _{3} }=124^{\circ}$
\end{minipage}\hfill
\begin{minipage}[t]{0.39\textwidth}
\vspace{0pt}
\centering

\resizebox{\linewidth}{!}{%
\begin{tikzpicture}
            \coordinate (E) at (1.14, 3.06);
            \coordinate (D) at (4.46, 0);
            \coordinate (C) at (-0.6, 0);
            \coordinate (B) at (5, 2);
            \coordinate (A) at (0, 2);
            \coordinate (F) at (3.6, -1.44);
            \coordinate (H) at (2.81, 0);
            \coordinate (G) at (1.72, 2);
            % Vẽ tứ giác
            \draw[thick] (C) -- (H)--(D);
            \draw[thick] (A) -- (G) -- (B);
            \draw[thick] (E) -- (G) -- (H)--(F);
            % Ghi nhãn các điểm
            
            \node[above] at (G) {$A $};
            \node[above ] at (H){$B $};
                      \node[above] at (A) {$x$};
                        \node[above] at (C) {$y$};
            \draw 
            (G) pic[draw,angle radius = 8] {angle = H--G--B} node[shift={(-25:15pt)}]{$1$}  ;
            \draw (G) node[shift={(30:11pt)}]{$2$};
            \draw  (G) node[shift={(50:-15pt)}]{$3$}  ;
            \draw (G)  node[shift={(-25:-15pt)}]{$4$};  
            \draw (H) pic[draw,angle radius = 10] {angle = E--H--C};
            \draw (H) node[shift={(-25:15pt)}]{$1$};
            \draw (H) node[shift={(40:15pt)}]{$2$}  ;
            \draw  (H) node[shift={(50:-13pt)}]{$3$} ;
            \draw (H) node[shift={(-25:-15pt)}]{$4$}  ;     
        \end{tikzpicture}
}
\end{minipage}

\vspace{3pt}\hrule\vspace{3pt}\footnotesize\textcolor{gray!80!black}{\textbf{Các mã dùng chung hình (HINH\_094):} \texttt{L07\_C3\_B2\_14, L07\_C3\_B2\_09, L07\_C3\_B3\_07}}
\end{tcolorbox}
\vspace{4pt}


\begin{tcolorbox}[
    breakable,
    colback=white,
    colframe=blue!60!black,
    boxrule=0.8pt,
    arc=3pt,
    title={\textbf{Câu \#123} \quad \textbf{MÃ CÂU: \texttt{[L07\_C3\_B3\_12]}} \quad | \quad \textbf{Nhóm: \texttt{HINH\_095}} \quad | \quad \textbf{MapID: \texttt{7H3H3}} \hfill \textcolor{yellow!90!white}{[Hình gốc .pyc]}}
]
\textbf{Dạng toán:} Thông hiểu 2 đường thẳng song song 2 \hfill \textit{(Trắc nghiệm -- Thông hiểu)}
\vspace{4pt}\hrule\vspace{6pt}
\begin{minipage}[t]{0.58\textwidth}
\vspace{0pt}
Cho hình vẽ biết $\widehat { E_{1} }=45^{\circ}$, $\widehat { F_{3} }=135^{\circ}$. Trong các khẳng định sau khẳng định nào sai.

\vspace{4pt}\noindent
\textbf{A.} $\widehat{ F_{4} } = 45^{\circ}$ \quad \textbf{\textcolor{blue!75!black}{B. (Đ)}} $\widehat{  F_{2} }=45^{\circ}$ \quad \textbf{C.} $\widehat { E_{2} } =135^{\circ} $ m \quad \textbf{D.} $x //y$
\end{minipage}\hfill
\begin{minipage}[t]{0.39\textwidth}
\vspace{0pt}
\centering

\resizebox{\linewidth}{!}{%
\begin{tikzpicture}
            \coordinate (E) at (1.14, 3.06);
            \coordinate (D) at (4.46, 0);
            \coordinate (C) at (-0.6, 0);
            \coordinate (B) at (5, 2);
            \coordinate (A) at (0, 2);
            \coordinate (F) at (3.6, -1.44);
            \coordinate (H) at (2.81, 0);
            \coordinate (G) at (1.72, 2);
            % Vẽ tứ giác
            \draw[thick] (C) -- (H)--(D);
            \draw[thick] (A) -- (G) -- (B);
            \draw[thick] (E) -- (G) -- (H)--(F);
            % Ghi nhãn các điểm
            
            \node[above] at (G) {$E$};
            \node[above ] at (H){$F$};
                      \node[above] at (A) {$x$};
                        \node[above] at (C) {$y$};
            \draw 
            (G) node[shift={(-25:15pt)}]{$1$}  ;
            \draw (G) node[shift={(30:11pt)}]{$2$};
            \draw  (G) node[shift={(50:-15pt)}]{$3$}  ;
            \draw (G)  node[shift={(-25:-15pt)}]{$4$};  
           
            \draw (H) node[shift={(-25:15pt)}]{$1$};
            \draw (H) node[shift={(40:15pt)}]{$2$}  ;
            \draw  (H) node[shift={(50:-13pt)}]{$3$} ;
            \draw (H) node[shift={(-25:-15pt)}]{$4$}  ;     
        \end{tikzpicture}
}
\end{minipage}

\vspace{3pt}\hrule\vspace{3pt}\footnotesize\textcolor{gray!80!black}{\textbf{Các mã dùng chung hình (HINH\_095):} \texttt{L07\_C3\_B2\_15, L07\_C3\_B2\_16, L07\_C3\_B2\_17, L07\_C3\_B2\_10, L07\_C3\_B3\_13, L07\_C3\_B3\_14, L07\_C3\_B3\_08}}
\end{tcolorbox}
\vspace{4pt}


\begin{tcolorbox}[
    breakable,
    colback=white,
    colframe=blue!60!black,
    boxrule=0.8pt,
    arc=3pt,
    title={\textbf{Câu \#124} \quad \textbf{MÃ CÂU: \texttt{[L07\_C3\_B3\_13]}} \quad | \quad \textbf{Nhóm: \texttt{HINH\_095}} \quad | \quad \textbf{MapID: \texttt{7H3H3}} \hfill \textcolor{yellow!90!white}{[Hình gốc .pyc]}}
]
\textbf{Dạng toán:} Tính số đo góc 1. \hfill \textit{(Trắc nghiệm -- Thông hiểu)}
\vspace{4pt}\hrule\vspace{6pt}
\begin{minipage}[t]{0.58\textwidth}
\vspace{0pt}
Cho hình vẽ biết $a //b$ và $\widehat { A _{1} }=50^{\circ}$. Trong các khẳng định sau khẳng định nào sai.

\vspace{4pt}\noindent
\textbf{A.} $\widehat{ B _{4} } = 50^{\circ}$ \quad \textbf{B.} $\widehat { A _{4} }=50^{\circ}$ \quad \textbf{\textcolor{blue!75!black}{C. (Đ)}} $\widehat{  B _{2} }=50^{\circ}$ \quad \textbf{D.} $\widehat { A _{2} } =130^{\circ} $
\end{minipage}\hfill
\begin{minipage}[t]{0.39\textwidth}
\vspace{0pt}
\centering

\resizebox{\linewidth}{!}{%
\begin{tikzpicture}
            \coordinate (E) at (1.14, 3.06);
            \coordinate (D) at (4.46, 0);
            \coordinate (C) at (-0.6, 0);
            \coordinate (B) at (5, 2);
            \coordinate (A) at (0, 2);
            \coordinate (F) at (3.6, -1.44);
            \coordinate (H) at (2.81, 0);
            \coordinate (G) at (1.72, 2);
            % Vẽ tứ giác
            \draw[thick] (C) -- (H)--(D);
            \draw[thick] (A) -- (G) -- (B);
            \draw[thick] (E) -- (G) -- (H)--(F);
            % Ghi nhãn các điểm
            
            \node[above] at (G) {$A $};
            \node[above ] at (H){$B $};
                      \node[above] at (A) {$a$};
                        \node[above] at (C) {$b$};
            \draw 
            (G) node[shift={(-25:15pt)}]{$1$}  ;
            \draw (G) node[shift={(30:11pt)}]{$2$};
            \draw  (G) node[shift={(50:-15pt)}]{$3$}  ;
            \draw (G)  node[shift={(-25:-15pt)}]{$4$};  
           
            \draw (H) node[shift={(-25:15pt)}]{$1$};
            \draw (H) node[shift={(40:15pt)}]{$2$}  ;
            \draw  (H) node[shift={(50:-13pt)}]{$3$} ;
            \draw (H) node[shift={(-25:-15pt)}]{$4$}  ;     
        \end{tikzpicture}
}
\end{minipage}

\vspace{3pt}\hrule\vspace{3pt}\footnotesize\textcolor{gray!80!black}{\textbf{Các mã dùng chung hình (HINH\_095):} \texttt{L07\_C3\_B2\_15, L07\_C3\_B2\_16, L07\_C3\_B2\_17, L07\_C3\_B2\_10, L07\_C3\_B3\_12, L07\_C3\_B3\_14, L07\_C3\_B3\_08}}
\end{tcolorbox}
\vspace{4pt}


\begin{tcolorbox}[
    breakable,
    colback=white,
    colframe=blue!60!black,
    boxrule=0.8pt,
    arc=3pt,
    title={\textbf{Câu \#125} \quad \textbf{MÃ CÂU: \texttt{[L07\_C3\_B3\_14]}} \quad | \quad \textbf{Nhóm: \texttt{HINH\_095}} \quad | \quad \textbf{MapID: \texttt{7H3H3}} \hfill \textcolor{yellow!90!white}{[Hình gốc .pyc]}}
]
\textbf{Dạng toán:} Tính số đo góc 2. \hfill \textit{(Trắc nghiệm -- Thông hiểu)}
\vspace{4pt}\hrule\vspace{6pt}
\begin{minipage}[t]{0.58\textwidth}
\vspace{0pt}
Cho hình vẽ biết $x //y$ và $\widehat { A_{4} }=46^{\circ}$. Trong các khẳng định sau khẳng định nào sai.

\vspace{4pt}\noindent
\textbf{A.} $\widehat{ B_{3} } = 134^{\circ}$ \quad \textbf{B.} $\widehat { A_{1} }=46^{\circ}$ \quad \textbf{\textcolor{blue!75!black}{C. (Đ)}} $\widehat{  A_{3} }=46^{\circ}$ \quad \textbf{D.} $\widehat { B_{1} } =46^{\circ} $
\end{minipage}\hfill
\begin{minipage}[t]{0.39\textwidth}
\vspace{0pt}
\centering

\resizebox{\linewidth}{!}{%
\begin{tikzpicture}
            \coordinate (E) at (1.14, 3.06);
            \coordinate (D) at (4.46, 0);
            \coordinate (C) at (-0.6, 0);
            \coordinate (B) at (5, 2);
            \coordinate (A) at (0, 2);
            \coordinate (F) at (3.6, -1.44);
            \coordinate (H) at (2.81, 0);
            \coordinate (G) at (1.72, 2);
            % Vẽ tứ giác
            \draw[thick] (C) -- (H)--(D);
            \draw[thick] (A) -- (G) -- (B);
            \draw[thick] (E) -- (G) -- (H)--(F);
            % Ghi nhãn các điểm
            
            \node[above] at (G) {$A$};
            \node[above ] at (H){$B$};
                      \node[above] at (A) {$x$};
                        \node[above] at (C) {$y$};
            \draw 
            (G) node[shift={(-25:15pt)}]{$1$}  ;
            \draw (G) node[shift={(30:11pt)}]{$2$};
            \draw  (G) node[shift={(50:-15pt)}]{$3$}  ;
            \draw (G)  node[shift={(-25:-15pt)}]{$4$};  
           
            \draw (H) node[shift={(-25:15pt)}]{$1$};
            \draw (H) node[shift={(40:15pt)}]{$2$}  ;
            \draw  (H) node[shift={(50:-13pt)}]{$3$} ;
            \draw (H) node[shift={(-25:-15pt)}]{$4$}  ;     
        \end{tikzpicture}
}
\end{minipage}

\vspace{3pt}\hrule\vspace{3pt}\footnotesize\textcolor{gray!80!black}{\textbf{Các mã dùng chung hình (HINH\_095):} \texttt{L07\_C3\_B2\_15, L07\_C3\_B2\_16, L07\_C3\_B2\_17, L07\_C3\_B2\_10, L07\_C3\_B3\_12, L07\_C3\_B3\_13, L07\_C3\_B3\_08}}
\end{tcolorbox}
\vspace{4pt}


\begin{tcolorbox}[
    breakable,
    colback=white,
    colframe=blue!60!black,
    boxrule=0.8pt,
    arc=3pt,
    title={\textbf{Câu \#126} \quad \textbf{MÃ CÂU: \texttt{[L07\_C3\_B3\_15]}} \quad | \quad \textbf{Nhóm: \texttt{HINH\_107}} \quad | \quad \textbf{MapID: \texttt{7H3V3}} \hfill \textcolor{yellow!90!white}{[Hình gốc .pyc]}}
]
\textbf{Dạng toán:} Chứng minh hai đường thẳng song song. \hfill \textit{(Trắc nghiệm -- Vận dụng)}
\vspace{4pt}\hrule\vspace{6pt}
\begin{minipage}[t]{0.58\textwidth}
\vspace{0pt}
Cho hình vẽ sau. Trong các khẳng định sau khẳng định nào sai?

\vspace{4pt}\noindent
\textbf{A.} $\widehat{DBN }= 54^{\circ}$ \quad \textbf{\textcolor{blue!75!black}{B. (Đ)}} $BD$ là tia phân giác của $\widehat{ ABN}$ \quad \textbf{C.} $AM//CD$ \quad \textbf{D.} $\widehat{DBN }= 23^{\circ}$
\end{minipage}\hfill
\begin{minipage}[t]{0.39\textwidth}
\vspace{0pt}
\centering

\resizebox{\linewidth}{!}{%
\begin{tikzpicture}
            \coordinate (E) at (4, 3);
            \coordinate (D) at (3, 2);
            \coordinate (C) at (-6, 2);
            \coordinate (B) at (-2, 2);
            \coordinate (A) at (-1, 3);
            \coordinate (F) at (0, 0);
       
            % Vẽ tứ giác
            \draw[thick] (C) -- (B)--(D);
            \draw[thick] (A) -- (E);
            \draw[thick] (A) -- (B);
            \draw[thick] (B) --(F);
            % Ghi nhãn các điểm
            
            \node[above] at (A) {$A$};
            \node[above ] at (B){$B$};
                      \node[above] at (C) {$C$};
                        \node[above] at (D) {$D$};
                         \node[above] at (E) {$M$};
                        \node[above] at (F) {$N$};
         
    \draw 
            (A) pic[draw,angle radius = 10] {angle = B--A--E} node[shift={(-20:20pt)}]{$ 126^\circ $}   
            ;    
            \draw 
            (B) pic[draw,angle radius = 10] {angle = F--B--A} pic[draw,angle radius = 12] {angle = F--B--A} node[shift={(-5:25pt)}]{$ 77^\circ $}   
            ;    
            \draw 
            (B) pic[draw,angle radius = 10] {angle = C--B--F}  pic[draw,angle radius = 8] {angle = C--B--F}  pic[draw,angle radius = 12] {angle = C--B--F} node[shift={(20:-20pt)}]{$ 157^\circ $}   
            ;  
        \end{tikzpicture}
}
\end{minipage}

\end{tcolorbox}
\vspace{4pt}


\begin{tcolorbox}[
    breakable,
    colback=white,
    colframe=blue!60!black,
    boxrule=0.8pt,
    arc=3pt,
    title={\textbf{Câu \#127} \quad \textbf{MÃ CÂU: \texttt{[L07\_C3\_B3\_16]}} \quad | \quad \textbf{Nhóm: \texttt{HINH\_108}} \quad | \quad \textbf{MapID: \texttt{7H3H3}} \hfill \textcolor{yellow!90!white}{[Hình gốc .pyc]}}
]
\textbf{Dạng toán:} Chứng minh hai đường thẳng song song 2 \hfill \textit{(Trắc nghiệm -- Thông hiểu)}
\vspace{4pt}\hrule\vspace{6pt}
\begin{minipage}[t]{0.58\textwidth}
\vspace{0pt}
Cho hình vẽ sau. Trong các khẳng định sau khẳng định nào sai?

\vspace{4pt}\noindent
\textbf{\textcolor{blue!75!black}{A. (Đ)}} ${Ix}$ và ${Qz}$ không song song với nhau \quad \textbf{B.} $Jy// Qz $ \quad \textbf{C.} $Ix //Jy $ \quad \textbf{D.} $Ix //Qz  $
\end{minipage}\hfill
\begin{minipage}[t]{0.39\textwidth}
\vspace{0pt}
\centering

\resizebox{\linewidth}{!}{%
\begin{tikzpicture}
            \coordinate (E) at (4, 1);
            \coordinate (D) at (3, 2);
            \coordinate (C) at (0, 0);
            \coordinate (B) at (-1, 1);
            \coordinate (A) at (-2, 2);
            \coordinate (F) at (5, 0);
             \coordinate (I) at (-3, 3);
                  \coordinate (M) at (-1, 0);
             \coordinate (N) at (1, -1);
            % Vẽ tứ giác
            \draw[thick] (C) -- (B);
            \draw[thick] (A) -- (D);
            \draw[thick] (N)--(C)--(B) -- (A)--(I);
            \draw[thick] (B) --(E);
             \draw[thick] (M)--(C) --(F);
            % Ghi nhãn các điểm
            
            \node[below left] at (A) {$I$};
            \node[below left] at (B){$J$};
                      \node[above right] at (C) {$Q$};
                  \node[below] at (E) {$y$};
                         \node[below] at (D) {$x$};
                        \node[below] at (F) {$z$};
          
                    \node[below left] at (I) {$a$};
    \draw 
            (A) pic[draw,angle radius = 7] {angle = D--A--I} node[shift={(65:13pt)}]{$ 134^\circ $}   
            ;    
            \draw 
            (B) pic[draw,angle radius = 7] {angle = E--B--I} node[shift={(65:13pt)}]{$ 134^\circ $}   
            ;    
              \draw 
            (C) pic[draw,angle radius = 7] {angle = M--C--N} node[shift={(25: -20pt)}]{$ 134^\circ $}   
            ;  
        \end{tikzpicture}
}
\end{minipage}

\end{tcolorbox}
\vspace{4pt}


\begin{tcolorbox}[
    breakable,
    colback=white,
    colframe=blue!60!black,
    boxrule=0.8pt,
    arc=3pt,
    title={\textbf{Câu \#128} \quad \textbf{MÃ CÂU: \texttt{[L07\_C3\_B3\_17]}} \quad | \quad \textbf{Nhóm: \texttt{HINH\_109}} \quad | \quad \textbf{MapID: \texttt{7H3V3}} \hfill \textcolor{yellow!90!white}{[Hình gốc .pyc]}}
]
\textbf{Dạng toán:} Chứng minh hai đường thẳng song song 3 \hfill \textit{(Trắc nghiệm -- Vận dụng)}
\vspace{4pt}\hrule\vspace{6pt}
\begin{minipage}[t]{0.58\textwidth}
\vspace{0pt}
Cho hình vẽ sau. Trong các khẳng định sau khẳng định nào sai?

\vspace{4pt}\noindent
\textbf{A.} $GI //H F $ \quad \textbf{\textcolor{blue!75!black}{B. (Đ)}} $\widehat {FHJ}=129^{\circ} $ \quad \textbf{C.} $H F //J O $ \quad \textbf{D.} $GI //J O  $
\end{minipage}\hfill
\begin{minipage}[t]{0.39\textwidth}
\vspace{0pt}
\centering

\resizebox{\linewidth}{!}{%
\begin{tikzpicture}
            \coordinate (E) at (4, 3);
            \coordinate (D) at (3, 2);
            \coordinate (C) at (-6, 2);
            \coordinate (B) at (-2, 2);
            \coordinate (A) at (-1, 3);
            \coordinate (F) at (0, 0);
        \coordinate (M) at (4, 0);
            % Vẽ tứ giác
            \draw[thick] (C) -- (B);
            \draw[thick] (A) -- (E);
            \draw[thick] (A) -- (B);
            \draw[thick] (B) --(F);
             \draw[thick] (M) --(F);
            % Ghi nhãn các điểm
            
            \node[above] at (A) {$G$};
            \node[below ] at (B){$H$};
                      \node[above] at (C) {$F$};
                        \node[below] at (M) {$O$};
                         \node[above] at (E) {$I$};
                        \node[below] at (F) {$J$};
         
    \draw 
            (A) pic[draw,angle radius = 10] {angle = B--A--E} node[shift={(-20:20pt)}]{$ 129^\circ $}   
            ;    
            \draw 
            (B) pic[draw,angle radius = 10] {angle = F--B--A} pic[draw,angle radius = 12] {angle = F--B--A} node[shift={(-5:25pt)}]{$ 93^\circ $}   
            ;    
            \draw 
            (B) pic[draw,angle radius = 10] {angle = A--B--C}  node[shift={(-20:-25pt)}]{$ 129^\circ $}   
            ;  
              \draw 
            (F) pic[draw,angle radius = 10] {angle = M--F--B}  pic[draw,angle radius = 8] {angle = M--F--B}  pic[draw,angle radius = 12] {angle = M--F--B} node[shift={(30: 25pt)}]{$ 138^\circ $}   
            ;  
        \end{tikzpicture}
}
\end{minipage}

\end{tcolorbox}
\vspace{4pt}


\begin{tcolorbox}[
    breakable,
    colback=white,
    colframe=blue!60!black,
    boxrule=0.8pt,
    arc=3pt,
    title={\textbf{Câu \#129} \quad \textbf{MÃ CÂU: \texttt{[L07\_C3\_B3\_18]}} \quad | \quad \textbf{Nhóm: \texttt{HINH\_110}} \quad | \quad \textbf{MapID: \texttt{7H3V3}} \hfill \textcolor{yellow!90!white}{[Hình gốc .pyc]}}
]
\textbf{Dạng toán:} Chứng minh hai đường thẳng song song 4 \hfill \textit{(Trắc nghiệm -- Vận dụng)}
\vspace{4pt}\hrule\vspace{6pt}
\begin{minipage}[t]{0.58\textwidth}
\vspace{0pt}
Cho hình vẽ sau. Trong các khẳng định sau khẳng định nào sai?

\vspace{4pt}\noindent
\textbf{\textcolor{blue!75!black}{A. (Đ)}} $\widehat {KFN}=124^{\circ} $ \quad \textbf{B.} $F K //N I $ \quad \textbf{C.} $EM //F K $ \quad \textbf{D.} $EM //N I  $
\end{minipage}\hfill
\begin{minipage}[t]{0.39\textwidth}
\vspace{0pt}
\centering

\resizebox{\linewidth}{!}{%
\begin{tikzpicture}
            \coordinate (E) at (4, 3);
            \coordinate (D) at (3, 2);
            \coordinate (C) at (-6, 2);
            \coordinate (B) at (-2, 2);
            \coordinate (A) at (-1, 3);
            \coordinate (F) at (0, 0);
             \coordinate (N) at (0, 4);
        \coordinate (M) at (4, 0);
            % Vẽ tứ giác
            \draw[thick] (C) -- (B);
            \draw[thick] (A) -- (E);
            \draw[thick] (N)--(A) -- (B);
            \draw[thick] (B) --(F);
             \draw[thick] (M) --(F);
            % Ghi nhãn các điểm
            
            \node[above] at (A) {$E$};
            \node[below ] at (B){$F$};
                      \node[above] at (C) {$K$};
                        \node[below] at (M) {$I$};
                         \node[above] at (E) {$M$};
                        \node[below] at (F) {$N$};
          \node[below] at (N) {$J$};
    \draw 
            (A) pic[draw,angle radius = 10] {angle = E--A--N} node[shift={(20:20pt)}]{$ 56^\circ $}   
            ;    
            \draw 
            (B) pic[draw,angle radius = 10] {angle = F--B--A} pic[draw,angle radius = 12] {angle = F--B--A} node[shift={(-5:25pt)}]{$ 97^\circ $}   
            ;    
            \draw 
            (B) node[shift={(-35:-10pt)}]{$ 124^\circ $}   
            ;  
              \draw 
            (F) pic[draw,angle radius = 10] {angle = M--F--B}  pic[draw,angle radius = 8] {angle = M--F--B}  pic[draw,angle radius = 12] {angle = M--F--B} node[shift={(30: 25pt)}]{$ 139^\circ $}   
            ;  
        \end{tikzpicture}
}
\end{minipage}

\end{tcolorbox}
\vspace{4pt}


\begin{tcolorbox}[
    breakable,
    colback=white,
    colframe=blue!60!black,
    boxrule=0.8pt,
    arc=3pt,
    title={\textbf{Câu \#130} \quad \textbf{MÃ CÂU: \texttt{[L07\_C3\_B3\_19]}} \quad | \quad \textbf{Nhóm: \texttt{HINH\_111}} \quad | \quad \textbf{MapID: \texttt{7H3V3}} \hfill \textcolor{yellow!90!white}{[Hình gốc .pyc]}}
]
\textbf{Dạng toán:} Tính số đo góc 1 \hfill \textit{(Trắc nghiệm -- Vận dụng)}
\vspace{4pt}\hrule\vspace{6pt}
\begin{minipage}[t]{0.58\textwidth}
\vspace{0pt}
Cho hình vẽ sau biết $Qa//Rb$ tìm ${x}$ ?

\vspace{4pt}\noindent
\textbf{A.} $x=29^{\circ}$ \quad \textbf{\textcolor{blue!75!black}{B. (Đ)}} $x=39^{\circ} $ \quad \textbf{C.} $x=19^{\circ}  $ \quad \textbf{D.} $x=49^{\circ} $
\end{minipage}\hfill
\begin{minipage}[t]{0.39\textwidth}
\vspace{0pt}
\centering

\resizebox{\linewidth}{!}{%
\begin{tikzpicture}
            \coordinate (E) at (5, 3);
            \coordinate (C) at (0, 0);
            \coordinate (B) at (0.5, 3);
            \coordinate (A) at (-1, 2);
            \coordinate (F) at (-4, 0);
          
            % Vẽ tứ giác
            \draw[thick] (F)--(C) -- (A)--(B)--(E);
            % Ghi nhãn các điểm
            
            \node[above] at (A) {$P$};
            \node[above ] at (B){$Q$};
                      \node[below] at (C) {$R$};
                         \node[above] at (E) {$a$};
                        \node[below] at (F) {$b$};
    \draw 
            (A)  node[shift={(10:10pt)}]{$x^\circ $}   
            ;    
            \draw 
            (B) node[shift={(-50:10pt)}]{$ 178^\circ $}   
            ;    
            \draw 
            (C) node[shift={(-35:-15pt)}]{$ 37^\circ $}   
            ;  
            \fill[black] (A) circle (1.5pt) ; 
            \fill[black] (B) circle (1.5pt)  ;
             \fill[black] (C) circle (1.5pt)  ;
             
        \end{tikzpicture}
}
\end{minipage}

\end{tcolorbox}
\vspace{4pt}


\begin{tcolorbox}[
    breakable,
    colback=white,
    colframe=blue!60!black,
    boxrule=0.8pt,
    arc=3pt,
    title={\textbf{Câu \#131} \quad \textbf{MÃ CÂU: \texttt{[L07\_C3\_B3\_20]}} \quad | \quad \textbf{Nhóm: \texttt{HINH\_112}} \quad | \quad \textbf{MapID: \texttt{7H3V3}} \hfill \textcolor{yellow!90!white}{[Hình gốc .pyc]}}
]
\textbf{Dạng toán:} Tính số đo góc 2 \hfill \textit{(Trắc nghiệm -- Vận dụng)}
\vspace{4pt}\hrule\vspace{6pt}
\begin{minipage}[t]{0.58\textwidth}
\vspace{0pt}
Cho hình vẽ sau biết $Aa//Db$ tìm ${x}$ ?

\vspace{4pt}\noindent
\textbf{A.} $x=185^{\circ} $ \quad \textbf{B.} $x=175^{\circ}  $ \quad \textbf{C.} $x=155^{\circ}$ \quad \textbf{\textcolor{blue!75!black}{D. (Đ)}} $x=165^{\circ} $
\end{minipage}\hfill
\begin{minipage}[t]{0.39\textwidth}
\vspace{0pt}
\centering

\resizebox{\linewidth}{!}{%
\begin{tikzpicture}
            \coordinate (E) at (3, 4);
            \coordinate (C) at (-1, 1);
            \coordinate (B) at (0.5, 2);
            \coordinate (A) at (0, 4);
            \coordinate (F) at (-3, 0);
           \coordinate (D) at (0, 0);
            % Vẽ tứ giác
            \draw[thick] (F)--(D)--(C) -- (B)--(A)--(E);
            % Ghi nhãn các điểm
              \node[above right] at (D) {$D$};
            \node[above left] at (A) {$A$};
            \node[above right ] at (B){$B$};
                      \node[above left] at (C) {$C$};
                         \node[above] at (E) {$a$};
                        \node[above] at (F) {$b$};
    \draw 
            (A)  node[shift={(-40:13pt)}]{$44^\circ $}   
            ;    
            \draw 
            (B) node[shift={(-50:-10pt)}]{$ x^\circ $}   
            ;    
            \draw 
            (C) node[shift={(-5:15pt)}]{$ 68^\circ $}   
            ;  
             \draw 
            (D) node[shift={(-21:-16pt)}]{$ 39^\circ $}   
            ;  
           \fill[black] (A) circle (1.5pt) ; 
            \fill[black] (B) circle (1.5pt)  ;
             \fill[black] (C) circle (1.5pt)  ;
              \fill[black] (D) circle (1.5pt)  ;
        \end{tikzpicture}
}
\end{minipage}

\end{tcolorbox}
\vspace{4pt}


\begin{tcolorbox}[
    breakable,
    colback=white,
    colframe=blue!60!black,
    boxrule=0.8pt,
    arc=3pt,
    title={\textbf{Câu \#132} \quad \textbf{MÃ CÂU: \texttt{[L07\_C3\_B3\_01]}} \quad | \quad \textbf{Nhóm: \texttt{HINH\_104}} \quad | \quad \textbf{MapID: \texttt{7H3N3}} \hfill \textcolor{yellow!90!white}{[Hình gốc .pyc]}}
]
\textbf{Dạng toán:} Tính số đo góc 1 \hfill \textit{(Tự luận -- Nhận biết)}
\vspace{4pt}\hrule\vspace{6pt}
\begin{minipage}[t]{0.58\textwidth}
\vspace{0pt}
Cho hình vẽ 

a) Kể tên các góc bằng nhau có trong hình.

b) Cho  $\widehat { I_{4} }= \widehat{ J_{1} }=48^{\circ}$ tính số đo các góc còn lại.
\end{minipage}\hfill
\begin{minipage}[t]{0.39\textwidth}
\vspace{0pt}
\centering

\resizebox{\linewidth}{!}{%
\begin{tikzpicture}
            \coordinate (E) at (1.14, 3.06);
            \coordinate (D) at (4.46, 0);
            \coordinate (C) at (-0.6, 0);
            \coordinate (B) at (5, 2);
            \coordinate (A) at (0, 2);
            \coordinate (F) at (3.6, -1.44);
            \coordinate (H) at (2.81, 0);
            \coordinate (G) at (1.72, 2);
            % Vẽ tứ giác
            \draw[thick] (C) -- (H)--(D);
            \draw[thick] (A) -- (G) -- (B);
            \draw[thick] (E) -- (G) -- (H)--(F);
            % Ghi nhãn các điểm
            
            \node[above] at (G) {$I$};
            \node[above ] at (H){$J$};
            \draw 
            (G) pic[draw,angle radius = 8] {angle = E--G--A} node[shift={(-25:15pt)}]{$1$}  ;
            \draw (G) node[shift={(30:11pt)}]{$2$};
            \draw  (G) node[shift={(50:-15pt)}]{$3$}  ;
            \draw (G)  node[shift={(-25:-15pt)}]{$4$};  
            \draw (H) pic[draw,angle radius = 10] {angle = F--H--D};
            \draw (H) node[shift={(-25:15pt)}]{$1$};
            \draw (H) node[shift={(40:15pt)}]{$2$}  ;
            \draw  (H) node[shift={(50:-13pt)}]{$3$} ;
            \draw (H) node[shift={(-25:-15pt)}]{$4$}  ;     
        \end{tikzpicture}
}
\end{minipage}

\vspace{3pt}\hrule\vspace{3pt}\footnotesize\textcolor{gray!80!black}{\textbf{Các mã dùng chung hình (HINH\_104):} \texttt{L07\_C3\_B2\_03}}
\end{tcolorbox}
\vspace{4pt}


\begin{tcolorbox}[
    breakable,
    colback=white,
    colframe=blue!60!black,
    boxrule=0.8pt,
    arc=3pt,
    title={\textbf{Câu \#133} \quad \textbf{MÃ CÂU: \texttt{[L07\_C3\_B3\_02]}} \quad | \quad \textbf{Nhóm: \texttt{HINH\_105}} \quad | \quad \textbf{MapID: \texttt{7H3H3}} \hfill \textcolor{yellow!90!white}{[Hình gốc .pyc]}}
]
\textbf{Dạng toán:} Tính số đo góc 2 \hfill \textit{(Tự luận -- Thông hiểu)}
\vspace{4pt}\hrule\vspace{6pt}
\begin{minipage}[t]{0.58\textwidth}
\vspace{0pt}
Cho hình vẽ 

a) Kể tên các góc bằng nhau có trong hình.

b) Cho  $\widehat { A _{4} }= \widehat{ B _{4} }=65^{\circ}$ tính số đo các góc còn lại.
\end{minipage}\hfill
\begin{minipage}[t]{0.39\textwidth}
\vspace{0pt}
\centering

\resizebox{\linewidth}{!}{%
\begin{tikzpicture}
            \coordinate (E) at (1.14, 3.06);
            \coordinate (D) at (4.46, 0);
            \coordinate (C) at (-0.6, 0);
            \coordinate (B) at (5, 2);
            \coordinate (A) at (0, 2);
            \coordinate (F) at (3.6, -1.44);
            \coordinate (H) at (2.81, 0);
            \coordinate (G) at (1.72, 2);
            % Vẽ tứ giác
            \draw[thick] (C) -- (H)--(D);
            \draw[thick] (A) -- (G) -- (B);
            \draw[thick] (E) -- (G) -- (H)--(F);
            % Ghi nhãn các điểm
            
            \node[above] at (G) {$A $};
            \node[above ] at (H){$B $};
            \draw 
            (G) pic[draw,angle radius = 8] {angle = E--G--A} node[shift={(-25:15pt)}]{$1$}  ;
            \draw (G) node[shift={(30:11pt)}]{$2$};
            \draw  (G) node[shift={(50:-15pt)}]{$3$}  ;
            \draw (G)  node[shift={(-25:-15pt)}]{$4$};  
            \draw (H) pic[draw,angle radius = 10] {angle = E--H--C};
            \draw (H) node[shift={(-25:15pt)}]{$1$};
            \draw (H) node[shift={(40:15pt)}]{$2$}  ;
            \draw  (H) node[shift={(50:-13pt)}]{$3$} ;
            \draw (H) node[shift={(-25:-15pt)}]{$4$}  ;     
        \end{tikzpicture}
}
\end{minipage}

\vspace{3pt}\hrule\vspace{3pt}\footnotesize\textcolor{gray!80!black}{\textbf{Các mã dùng chung hình (HINH\_105):} \texttt{L07\_C3\_B2\_04}}
\end{tcolorbox}
\vspace{4pt}


\begin{tcolorbox}[
    breakable,
    colback=white,
    colframe=blue!60!black,
    boxrule=0.8pt,
    arc=3pt,
    title={\textbf{Câu \#134} \quad \textbf{MÃ CÂU: \texttt{[L07\_C3\_B3\_03]}} \quad | \quad \textbf{Nhóm: \texttt{HINH\_106}} \quad | \quad \textbf{MapID: \texttt{7H3H3}} \hfill \textcolor{yellow!90!white}{[Hình gốc .pyc]}}
]
\textbf{Dạng toán:} Tính số đo góc 3 \hfill \textit{(Tự luận -- Thông hiểu)}
\vspace{4pt}\hrule\vspace{6pt}
\begin{minipage}[t]{0.58\textwidth}
\vspace{0pt}
Cho hình vẽ 

a) Kể tên các góc bằng nhau có trong hình.

b) Cho  $\widehat { A _{1} }= \widehat{ B _{4} }=47^{\circ}$ tính số đo các góc còn lại.
\end{minipage}\hfill
\begin{minipage}[t]{0.39\textwidth}
\vspace{0pt}
\centering

\resizebox{\linewidth}{!}{%
\begin{tikzpicture}
            \coordinate (E) at (1.14, 3.06);
            \coordinate (D) at (4.46, 0);
            \coordinate (C) at (-0.6, 0);
            \coordinate (B) at (5, 2);
            \coordinate (A) at (0, 2);
            \coordinate (F) at (3.6, -1.44);
            \coordinate (H) at (2.81, 0);
            \coordinate (G) at (1.72, 2);
            % Vẽ tứ giác
            \draw[thick] (C) -- (H)--(D);
            \draw[thick] (A) -- (G) -- (B);
            \draw[thick] (E) -- (G) -- (H)--(F);
            % Ghi nhãn các điểm
            
            \node[above] at (G) {$A $};
            \node[above ] at (H){$B $};
            \draw 
            (G) pic[draw,angle radius = 8] {angle = H--G--B} node[shift={(-25:15pt)}]{$1$}  ;
            \draw (G) node[shift={(30:11pt)}]{$2$};
            \draw  (G) node[shift={(50:-15pt)}]{$3$}  ;
            \draw (G)  node[shift={(-25:-15pt)}]{$4$};  
            \draw (H) pic[draw,angle radius = 10] {angle = E--H--C};
            \draw (H) node[shift={(-25:15pt)}]{$1$};
            \draw (H) node[shift={(40:15pt)}]{$2$}  ;
            \draw  (H) node[shift={(50:-13pt)}]{$3$} ;
            \draw (H) node[shift={(-25:-15pt)}]{$4$}  ;     
        \end{tikzpicture}
}
\end{minipage}

\vspace{3pt}\hrule\vspace{3pt}\footnotesize\textcolor{gray!80!black}{\textbf{Các mã dùng chung hình (HINH\_106):} \texttt{L07\_C3\_B2\_05}}
\end{tcolorbox}
\vspace{4pt}


\begin{tcolorbox}[
    breakable,
    colback=white,
    colframe=blue!60!black,
    boxrule=0.8pt,
    arc=3pt,
    title={\textbf{Câu \#135} \quad \textbf{MÃ CÂU: \texttt{[L07\_C3\_B3\_04]}} \quad | \quad \textbf{Nhóm: \texttt{HINH\_091}} \quad | \quad \textbf{MapID: \texttt{7H3H3}} \hfill \textcolor{yellow!90!white}{[Hình gốc .pyc]}}
]
\textbf{Dạng toán:} Tính số đo góc 4. \hfill \textit{(Tự luận -- Thông hiểu)}
\vspace{4pt}\hrule\vspace{6pt}
\begin{minipage}[t]{0.58\textwidth}
\vspace{0pt}
Cho hình vẽ sau biết $\widehat{ Q_{4}} =156^{\circ}$, $\widehat{ P_{2}} =63^{\circ}$

 Tính số đo của các góc còn lại.
\end{minipage}\hfill
\begin{minipage}[t]{0.39\textwidth}
\vspace{0pt}
\centering

\resizebox{\linewidth}{!}{%
\begin{tikzpicture}
            % Vẽ các điểm A, B, C, D
            \coordinate (A) at (-2, 0);
            \coordinate (C) at (-3, -2);
            \coordinate (B) at (2, 0);
            \coordinate (D) at (-1, 2);
            \coordinate (E) at (-4, 0);
            \coordinate (F) at (3, -2);
            \coordinate (M) at (1, 2);
            \coordinate (N) at (4, 0);
            
            \draw[thick] (E) -- (N);
            \draw[thick] (C) -- (D) ;
            \draw[thick] (M) -- (F) ;
            
            % Đánh dấu các điểm A, B, 
            \fill (B) circle (2pt) node[above] {$Q$};
            \fill (A) circle (2pt) node[above] {$P$};
            \draw 
            (B) node[shift={(-25:15pt)}]{$1$}   
            ;
            \draw 
            (B) node[shift={(45:-15pt)}]{$2$}   
            ;
            \draw 
            (B) node[shift={(-50:-15pt)}]{$3$}   
            ;
            \draw 
            (B) node[shift={(50:15pt)}]{$4$}   
            ;
            \draw 
            (A)  node[shift={(-50:12pt)}]{$1$}   
            ;
            \draw 
            (A) node[shift={(50:-13pt)}]{$2$}   
            ;
            \draw 
            (A)  node[shift={(-50:-15pt)}]{$3$}   
            ;
            \draw 
            (A) node[shift={(50:15pt)}]{$4$}   
            ;
            
        \end{tikzpicture}
}
\end{minipage}

\vspace{3pt}\hrule\vspace{3pt}\footnotesize\textcolor{gray!80!black}{\textbf{Các mã dùng chung hình (HINH\_091):} \texttt{L07\_C3\_B2\_11, L07\_C3\_B2\_06}}
\end{tcolorbox}
\vspace{4pt}


\begin{tcolorbox}[
    breakable,
    colback=white,
    colframe=blue!60!black,
    boxrule=0.8pt,
    arc=3pt,
    title={\textbf{Câu \#136} \quad \textbf{MÃ CÂU: \texttt{[L07\_C3\_B3\_05]}} \quad | \quad \textbf{Nhóm: \texttt{HINH\_092}} \quad | \quad \textbf{MapID: \texttt{7H3N3}} \hfill \textcolor{yellow!90!white}{[Hình gốc .pyc]}}
]
\textbf{Dạng toán:} Nhận biết hai đường thẳng song song \hfill \textit{(Tự luận -- Nhận biết)}
\vspace{4pt}\hrule\vspace{6pt}
\begin{minipage}[t]{0.58\textwidth}
\vspace{0pt}
Cho hình vẽ biết $\widehat { G_{4} }=\widehat{ H_{1} }=49^{\circ}$ hai đường thẳng ${m }$ và ${n }$ có song song với nhau không? tính số đo các góc còn lại.
\end{minipage}\hfill
\begin{minipage}[t]{0.39\textwidth}
\vspace{0pt}
\centering

\resizebox{\linewidth}{!}{%
\begin{tikzpicture}
            \coordinate (E) at (1.14, 3.06);
            \coordinate (D) at (4.46, 0);
            \coordinate (C) at (-0.6, 0);
            \coordinate (B) at (5, 2);
            \coordinate (A) at (0, 2);
            \coordinate (F) at (3.6, -1.44);
            \coordinate (H) at (2.81, 0);
            \coordinate (G) at (1.72, 2);
            % Vẽ tứ giác
            \draw[thick] (C) -- (H)--(D);
            \draw[thick] (A) -- (G) -- (B);
            \draw[thick] (E) -- (G) -- (H)--(F);
            % Ghi nhãn các điểm
            
            \node[above] at (G) {$G$};
            \node[above ] at (H){$H$};
                \node[above] at (A) {$m $};
                        \node[above] at (C) {$n $};
            \draw 
            (G) pic[draw,angle radius = 8] {angle = E--G--A} node[shift={(-25:15pt)}]{$1$}  ;
            \draw (G) node[shift={(30:11pt)}]{$2$};
            \draw  (G) node[shift={(50:-15pt)}]{$3$}  ;
            \draw (G)  node[shift={(-25:-15pt)}]{$4$};  
            \draw (H) pic[draw,angle radius = 10] {angle = F--H--D};
            \draw (H) node[shift={(-25:15pt)}]{$1$};
            \draw (H) node[shift={(40:15pt)}]{$2$}  ;
            \draw  (H) node[shift={(50:-13pt)}]{$3$} ;
            \draw (H) node[shift={(-25:-15pt)}]{$4$}  ;     
        \end{tikzpicture}
}
\end{minipage}

\vspace{3pt}\hrule\vspace{3pt}\footnotesize\textcolor{gray!80!black}{\textbf{Các mã dùng chung hình (HINH\_092):} \texttt{L07\_C3\_B2\_12, L07\_C3\_B2\_07, L07\_C3\_B3\_09}}
\end{tcolorbox}
\vspace{4pt}


\begin{tcolorbox}[
    breakable,
    colback=white,
    colframe=blue!60!black,
    boxrule=0.8pt,
    arc=3pt,
    title={\textbf{Câu \#137} \quad \textbf{MÃ CÂU: \texttt{[L07\_C3\_B3\_06]}} \quad | \quad \textbf{Nhóm: \texttt{HINH\_093}} \quad | \quad \textbf{MapID: \texttt{7H3H3}} \hfill \textcolor{yellow!90!white}{[Hình gốc .pyc]}}
]
\textbf{Dạng toán:} Chứng minh hai đường thẳng song song dựa vào 2 góc đồng vị \hfill \textit{(Tự luận -- Thông hiểu)}
\vspace{4pt}\hrule\vspace{6pt}
\begin{minipage}[t]{0.58\textwidth}
\vspace{0pt}
Cho hình vẽ biết $\widehat { M_{4} }= \widehat{ N_{4} }=45^{\circ}$ hai đường thẳng ${ b}$ và ${ a}$ có song song với nhau không? tính số đo các góc còn lại.
\end{minipage}\hfill
\begin{minipage}[t]{0.39\textwidth}
\vspace{0pt}
\centering

\resizebox{\linewidth}{!}{%
\begin{tikzpicture}
            \coordinate (E) at (1.14, 3.06);
            \coordinate (D) at (4.46, 0);
            \coordinate (C) at (-0.6, 0);
            \coordinate (B) at (5, 2);
            \coordinate (A) at (0, 2);
            \coordinate (F) at (3.6, -1.44);
            \coordinate (H) at (2.81, 0);
            \coordinate (G) at (1.72, 2);
            % Vẽ tứ giác
            \draw[thick] (C) -- (H)--(D);
            \draw[thick] (A) -- (G) -- (B);
            \draw[thick] (E) -- (G) -- (H)--(F);
            % Ghi nhãn các điểm
            
            \node[above] at (G) {$M$};
            \node[above ] at (H){$N$};
                \node[above] at (A) {$ b$};
                        \node[above] at (C) {$ a$};
            \draw 
            (G) pic[draw,angle radius = 8] {angle = E--G--A} node[shift={(-25:15pt)}]{$1$}  ;
            \draw (G) node[shift={(30:11pt)}]{$2$};
            \draw  (G) node[shift={(50:-15pt)}]{$3$}  ;
            \draw (G)  node[shift={(-25:-15pt)}]{$4$};  
            \draw (H) pic[draw,angle radius = 10] {angle = E--H--C};
            \draw (H) node[shift={(-25:15pt)}]{$1$};
            \draw (H) node[shift={(40:15pt)}]{$2$}  ;
            \draw  (H) node[shift={(50:-13pt)}]{$3$} ;
            \draw (H) node[shift={(-25:-15pt)}]{$4$}  ;     
        \end{tikzpicture}
}
\end{minipage}

\vspace{3pt}\hrule\vspace{3pt}\footnotesize\textcolor{gray!80!black}{\textbf{Các mã dùng chung hình (HINH\_093):} \texttt{L07\_C3\_B2\_13, L07\_C3\_B2\_08, L07\_C3\_B3\_10}}
\end{tcolorbox}
\vspace{4pt}


\begin{tcolorbox}[
    breakable,
    colback=white,
    colframe=blue!60!black,
    boxrule=0.8pt,
    arc=3pt,
    title={\textbf{Câu \#138} \quad \textbf{MÃ CÂU: \texttt{[L07\_C3\_B3\_07]}} \quad | \quad \textbf{Nhóm: \texttt{HINH\_094}} \quad | \quad \textbf{MapID: \texttt{7H3N3}} \hfill \textcolor{yellow!90!white}{[Hình gốc .pyc]}}
]
\textbf{Dạng toán:} Nhận biết 2 đường thẳng song song dựa vào 2 góc so le trong \hfill \textit{(Tự luận -- Nhận biết)}
\vspace{4pt}\hrule\vspace{6pt}
\begin{minipage}[t]{0.58\textwidth}
\vspace{0pt}
Cho hình vẽ biết $\widehat { P_{1} }= \widehat{ Q_{4} }=53^{\circ}$, hỏi hai đường thẳng ${d}$ và ${d'}$ có song song với nhau không? tính số đo các góc còn lại.
\end{minipage}\hfill
\begin{minipage}[t]{0.39\textwidth}
\vspace{0pt}
\centering

\resizebox{\linewidth}{!}{%
\begin{tikzpicture}
            \coordinate (E) at (1.14, 3.06);
            \coordinate (D) at (4.46, 0);
            \coordinate (C) at (-0.6, 0);
            \coordinate (B) at (5, 2);
            \coordinate (A) at (0, 2);
            \coordinate (F) at (3.6, -1.44);
            \coordinate (H) at (2.81, 0);
            \coordinate (G) at (1.72, 2);
            % Vẽ tứ giác
            \draw[thick] (C) -- (H)--(D);
            \draw[thick] (A) -- (G) -- (B);
            \draw[thick] (E) -- (G) -- (H)--(F);
            % Ghi nhãn các điểm
            
            \node[above] at (G) {$P$};
            \node[above ] at (H){$Q$};
                      \node[above] at (A) {$d$};
                        \node[above] at (C) {$d'$};
            \draw 
            (G) pic[draw,angle radius = 8] {angle = H--G--B} node[shift={(-25:15pt)}]{$1$}  ;
            \draw (G) node[shift={(30:11pt)}]{$2$};
            \draw  (G) node[shift={(50:-15pt)}]{$3$}  ;
            \draw (G)  node[shift={(-25:-15pt)}]{$4$};  
            \draw (H) pic[draw,angle radius = 10] {angle = E--H--C};
            \draw (H) node[shift={(-25:15pt)}]{$1$};
            \draw (H) node[shift={(40:15pt)}]{$2$}  ;
            \draw  (H) node[shift={(50:-13pt)}]{$3$} ;
            \draw (H) node[shift={(-25:-15pt)}]{$4$}  ;     
        \end{tikzpicture}
}
\end{minipage}

\vspace{3pt}\hrule\vspace{3pt}\footnotesize\textcolor{gray!80!black}{\textbf{Các mã dùng chung hình (HINH\_094):} \texttt{L07\_C3\_B2\_14, L07\_C3\_B2\_09, L07\_C3\_B3\_11}}
\end{tcolorbox}
\vspace{4pt}


\begin{tcolorbox}[
    breakable,
    colback=white,
    colframe=blue!60!black,
    boxrule=0.8pt,
    arc=3pt,
    title={\textbf{Câu \#139} \quad \textbf{MÃ CÂU: \texttt{[L07\_C3\_B3\_08]}} \quad | \quad \textbf{Nhóm: \texttt{HINH\_095}} \quad | \quad \textbf{MapID: \texttt{7H3H3}} \hfill \textcolor{yellow!90!white}{[Hình gốc .pyc]}}
]
\textbf{Dạng toán:} Thông hiểu 2 đường thẳng song song \hfill \textit{(Tự luận -- Thông hiểu)}
\vspace{4pt}\hrule\vspace{6pt}
\begin{minipage}[t]{0.58\textwidth}
\vspace{0pt}
Cho hình vẽ biết $\widehat { A_{1} }=40^{\circ}$, $\widehat { B_{3} }=140^{\circ}$ hỏi hai đường thẳng ${m }$ và ${n }$ có song song với nhau không? tính số đo các góc còn lại.
\end{minipage}\hfill
\begin{minipage}[t]{0.39\textwidth}
\vspace{0pt}
\centering

\resizebox{\linewidth}{!}{%
\begin{tikzpicture}
            \coordinate (E) at (1.14, 3.06);
            \coordinate (D) at (4.46, 0);
            \coordinate (C) at (-0.6, 0);
            \coordinate (B) at (5, 2);
            \coordinate (A) at (0, 2);
            \coordinate (F) at (3.6, -1.44);
            \coordinate (H) at (2.81, 0);
            \coordinate (G) at (1.72, 2);
            % Vẽ tứ giác
            \draw[thick] (C) -- (H)--(D);
            \draw[thick] (A) -- (G) -- (B);
            \draw[thick] (E) -- (G) -- (H)--(F);
            % Ghi nhãn các điểm
            
            \node[above] at (G) {$A$};
            \node[above ] at (H){$B$};
                      \node[above] at (A) {$m $};
                        \node[above] at (C) {$n $};
            \draw 
            (G) node[shift={(-25:15pt)}]{$1$}  ;
            \draw (G) node[shift={(30:11pt)}]{$2$};
            \draw  (G) node[shift={(50:-15pt)}]{$3$}  ;
            \draw (G)  node[shift={(-25:-15pt)}]{$4$};  
           
            \draw (H) node[shift={(-25:15pt)}]{$1$};
            \draw (H) node[shift={(40:15pt)}]{$2$}  ;
            \draw  (H) node[shift={(50:-13pt)}]{$3$} ;
            \draw (H) node[shift={(-25:-15pt)}]{$4$}  ;     
        \end{tikzpicture}
}
\end{minipage}

\vspace{3pt}\hrule\vspace{3pt}\footnotesize\textcolor{gray!80!black}{\textbf{Các mã dùng chung hình (HINH\_095):} \texttt{L07\_C3\_B2\_15, L07\_C3\_B2\_16, L07\_C3\_B2\_17, L07\_C3\_B2\_10, L07\_C3\_B3\_12, L07\_C3\_B3\_13, L07\_C3\_B3\_14}}
\end{tcolorbox}
\vspace{4pt}


\begin{tcolorbox}[
    breakable,
    colback=white,
    colframe=blue!60!black,
    boxrule=0.8pt,
    arc=3pt,
    title={\textbf{Câu \#140} \quad \textbf{MÃ CÂU: \texttt{[L07\_C3\_B4\_03]}} \quad | \quad \textbf{Nhóm: \texttt{HINH\_113}} \quad | \quad \textbf{MapID: \texttt{7H3N4}} \hfill \textcolor{yellow!90!white}{[Hình gốc .pyc]}}
]
\textbf{Dạng toán:} Diễn đạt bằng lời các định lí \hfill \textit{(Tự luận -- Nhận biết)}
\vspace{4pt}\hrule\vspace{6pt}
\begin{minipage}[t]{0.58\textwidth}
\vspace{0pt}
Diễn đạt bằng lời định lí sau
\end{minipage}\hfill
\begin{minipage}[t]{0.39\textwidth}
\vspace{0pt}
\centering

\resizebox{\linewidth}{!}{%
\begin{tikzpicture}
                \coordinate (E) at (2, 2);
                \coordinate (C) at (-2, 1);
                \coordinate (B) at (4, 2);
                \coordinate (A) at (-2, 2);
                \coordinate (F) at (1, 1);
               \coordinate (D) at (4, 1);
                \coordinate (H) at (0, 0);
                  \coordinate (G) at (3, 3);
                \draw[thick]  (C)--(F)--(D);
                \draw[thick]  (B)--(E)--(A);
                \draw[thick]  (H)--(F)--(E)--(G);
                % Ghi nhãn các điểm
                  \node[above left] at (E) {$H $};
                \node[above left] at (F) {$K $};
                 \draw 
                (E) pic[draw,angle radius = 10] {angle = B--E--G} 
                ;
                  \draw 
                (F) pic[draw,angle radius = 10] {angle = D--F--E} 
                ;
               \fill[black] (E) circle (1.5pt) ; 
                \fill[black] (F) circle (1.5pt)  ;
              
            \end{tikzpicture}
}
\end{minipage}

\end{tcolorbox}
\vspace{4pt}


\begin{tcolorbox}[
    breakable,
    colback=white,
    colframe=blue!60!black,
    boxrule=0.8pt,
    arc=3pt,
    title={\textbf{Câu \#141} \quad \textbf{MÃ CÂU: \texttt{[L07\_C3\_B4\_04]}} \quad | \quad \textbf{Nhóm: \texttt{HINH\_114}} \quad | \quad \textbf{MapID: \texttt{7H3V4}} \hfill \textcolor{yellow!90!white}{[Hình gốc .pyc]}}
]
\textbf{Dạng toán:} Chứng minh định lí. \hfill \textit{(Tự luận -- Vận dụng)}
\vspace{4pt}\hrule\vspace{6pt}
\begin{minipage}[t]{0.58\textwidth}
\vspace{0pt}
Chứng minh rằng nếu hai góc nhọn ${xGy}$ và ${x'G'y'}$ có $Gx // G'x'$, $Gy // G'y'$ thì $\widehat{xGy } = \widehat{x'G'y' } $
\end{minipage}\hfill
\begin{minipage}[t]{0.39\textwidth}
\vspace{0pt}
\centering

\resizebox{\linewidth}{!}{%
\begin{tikzpicture}
                \coordinate (E) at (1, 6);
                \coordinate (C) at (-1, 3);
                \coordinate (B) at (0, 0);
                \coordinate (A) at (-2, -1);
                \coordinate (F) at (2, 3);
               \coordinate (D) at (0, 2);
               \coordinate (G) at (2,5);
                \draw[thick]  (C)--(A)--(B);
                \draw[thick]  (E)--(D)--(F);
                \draw[thick]  (A)--(D)--(G);
                % Ghi nhãn các điểm
                  \node[below left] at (A) {$G$};
                \node[below left] at (D) {$G'$};
                \node[below left] at (C) {$x$};
                \node[above left] at (0.2, -0.5) {$y$};
                  \node[below left] at (E) {$x'$};
                \node[above left] at (2.5, 2.5) {$y'$};
                 \draw 
                (A) node[shift={(45:15pt)}]{$1$}   ;
                 \draw 
                (A) node[shift={(65:25pt)}]{$2$}   ;
                  \draw 
                  (D) node[shift={(45:15pt)}]{$1$}   ;
                 \draw 
                (D) node[shift={(65:25pt)}]{$2$}   ;
                  \fill[black] (A) circle (1.5pt); 
               \fill[black] (D) circle (1.5pt) ;
              
            \end{tikzpicture}
}
\end{minipage}

\end{tcolorbox}
\vspace{4pt}


\begin{tcolorbox}[
    breakable,
    colback=white,
    colframe=blue!60!black,
    boxrule=0.8pt,
    arc=3pt,
    title={\textbf{Câu \#142} \quad \textbf{MÃ CÂU: \texttt{[L07\_C3\_B6\_01]}} \quad | \quad \textbf{Nhóm: \texttt{HINH\_115}} \quad | \quad \textbf{MapID: \texttt{7H3H6}} \hfill \textcolor{yellow!90!white}{[Hình gốc .pyc]}}
]
\textbf{Dạng toán:} TF-M2. Tính số đo góc 1 \hfill \textit{(Đúng-Sai -- Thông hiểu)}
\vspace{4pt}\hrule\vspace{6pt}
\begin{minipage}[t]{0.58\textwidth}
\vspace{0pt}
Hình vẽ bên dưới thoả mãn các yêu cầu sau: góc $\widehat{ECK}$ bằng $68^{\circ}$; ${N}$ thuộc ${CK}$; ${LM}$ đi qua ${N}$; ${LM}$ song song với ${CE}$. Xét tính đúng-sai của các khẳng định sau.

\vspace{4pt}\noindent
\textbf{a)} $\widehat{KNM} = 73^{\circ}$\\[2pt]
\textbf{b)} $\widehat{LNC} = 79^{\circ}$\\[2pt]
\textbf{\textcolor{blue!75!black}{c) (Đ)}} $\widehat{CNM} = 112^{\circ}$\\[2pt]
\textbf{\textcolor{blue!75!black}{d) (Đ)}} Vẽ ${C}m$ là tia phân giác của góc ${KCE}$. Tia ${C}m$ cắt ${LM}$ tại ${O}$. Số đo $\widehat{NOC}= 34^{\circ}$
\end{minipage}\hfill
\begin{minipage}[t]{0.39\textwidth}
\vspace{0pt}
\centering

\resizebox{\linewidth}{!}{%
\begin{tikzpicture}[>=stealth,line join=round,line cap=round,font=\footnotesize,scale=1]
    \def\a{4}
    \def\b{3}
    \path
    (0,0) coordinate (A)
    (\a,0) coordinate (B)
    (60:\b) coordinate (F)
    ($(A)!1.3!(F)$) coordinate (C)
    ($(F)+(-0.5*\a,0)$) coordinate (D)
    ($(F)+(1.1*\a,0)$) coordinate (E);
    \draw (D)--(E) (C)--(A)--(B);
    \filldraw[black] (A) circle (2pt) node[above left ] {$C$};
    \filldraw[black] (B) circle (2pt) node[above right] {$E$};
    \filldraw[black] (C) circle (2pt) node[above ] {$K$};
   \filldraw[black] (D) circle (2pt) node[below]{$L$};
 \filldraw[black] (E) circle (2pt) node[above left ] {$M$};
    \filldraw[black] (F) circle (2pt) node[above right] {$N$};
    
    \end{tikzpicture}
}
\end{minipage}

\end{tcolorbox}
\vspace{4pt}


\begin{tcolorbox}[
    breakable,
    colback=white,
    colframe=blue!60!black,
    boxrule=0.8pt,
    arc=3pt,
    title={\textbf{Câu \#143} \quad \textbf{MÃ CÂU: \texttt{[L07\_C3\_B5\_01]}} \quad | \quad \textbf{Nhóm: \texttt{HINH\_116}} \quad | \quad \textbf{MapID: \texttt{7H3V5}} \hfill \textcolor{yellow!90!white}{[Hình gốc .pyc]}}
]
\textbf{Dạng toán:} Bài tổng hợp 1 \hfill \textit{(Tự luận -- Vận dụng)}
\vspace{4pt}\hrule\vspace{6pt}
\begin{minipage}[t]{0.58\textwidth}
\vspace{0pt}
1)  Vẽ hình, viết giả thiết và kết luận của định lí: Hai đường thẳng phân biệt cùng song song với một đường thẳng thứ ba thì chúng song song với nhau.

2) Cho hình vẽ, biết $A B  \parallel H I $, $\widehat{B A D }=60^\circ$, $\widehat{A D C }=60^\circ$, $\widehat{H K D }=114^\circ$ 

a) Giải thích vì sao $A B  \parallel C D $? 

b) Tính số đo góc $\widehat{A D K }$.
\end{minipage}\hfill
\begin{minipage}[t]{0.39\textwidth}
\vspace{0pt}
\centering

\resizebox{\linewidth}{!}{%
\begin{tikzpicture}[font=\footnotesize,line cap=round,line join=round,>=stealth,x=1.0cm,y=1.0cm,scale=0.7]
    \path 
    (0,0) coordinate (D)
    
    (-5,0) coordinate (C)   
    ($(D)!.3!60:(C)$) coordinate (F)
    ($(D)!.3!-56:(C)$) coordinate (A)   
    ($(A)+(D)-(C)$) coordinate (B)
    ($(F)+(D)-(C)$) coordinate (G)
    ($(F)+(C)-(D)$) coordinate (E)  
    ;
    \draw (D)--(C) (D)--(A) (D)--(F)  (A)--(B)   (G)--(E);
    
     \filldraw[black] (A) circle (2pt) node[above right] {$A $};
    \filldraw[black] (B) circle (2pt) node[above left] {$B $};
    \filldraw[black] (C) circle (2pt) node[above] {$C $};
   \filldraw[black] (D) circle (2pt) node[above right]{$D $};
    \filldraw[black] (E) circle (2pt) node[below left] {$H $};
    \filldraw[black] (F) circle (2pt) node[above] {$K $};
     \filldraw[black] (G) circle (2pt) node[above] {$I $};
    \end{tikzpicture}
}
\end{minipage}

\end{tcolorbox}
\vspace{4pt}


\begin{tcolorbox}[
    breakable,
    colback=white,
    colframe=blue!60!black,
    boxrule=0.8pt,
    arc=3pt,
    title={\textbf{Câu \#144} \quad \textbf{MÃ CÂU: \texttt{[L07\_C3\_B5\_02]}} \quad | \quad \textbf{Nhóm: \texttt{HINH\_117}} \quad | \quad \textbf{MapID: \texttt{7H3V5}} \hfill \textcolor{yellow!90!white}{[Hình gốc .pyc]}}
]
\textbf{Dạng toán:} Bài tổng hợp 2 \hfill \textit{(Tự luận -- Vận dụng)}
\vspace{4pt}\hrule\vspace{6pt}
\begin{minipage}[t]{0.58\textwidth}
\vspace{0pt}
Cho hình vẽ sau 

Biết rằng $\widehat{NMx }=31^\circ$; $\widehat{MNP}=57^\circ$; $\widehat{yNP}=26^\circ$; $\widehat{NPz}=26^\circ$.Chứng minh $Mx  \parallel Pz$.
\end{minipage}\hfill
\begin{minipage}[t]{0.39\textwidth}
\vspace{0pt}
\centering

\resizebox{\linewidth}{!}{%
\begin{tikzpicture}[line join = round, line cap = round,>=stealth,font=\footnotesize,scale=1.0]
    % Khai báo điểm, xác định các điểm trong hình vẽ
    \def \dodai{1.3};
    \path 
    (0:0) coordinate (S)
    (140:2*\dodai) coordinate (K)
    ($(K)+(3.6*\dodai,0)$) coordinate (x)
    (-150:2*\dodai) coordinate (Q)
    ($(Q)+(4*\dodai,0)$) coordinate (m)
    (180:3*\dodai) coordinate (k)
    ;
    %% Vẽ hình
    \pgfresetboundingbox
    % Vẽ đoạn thẳng
    \draw
    (K)--(x)
    (k)--(S)
    (Q)--(m)
    (K)--(S)--(Q)
    ;
    %\path 
    %pic[draw, angle radius=5.5mm,{$26^\circ$}] {angle = S--K--x} % Đánh dấu góc thường
    %pic[angle radius=4.5mm,{$31^\circ$}] {angle = k--S--Q} % Đánh dấu góc thường
    %pic[angle radius=5.5mm,{$31^\circ$}] {angle = m--Q--S} % Đánh dấu góc thường
    %;
        \filldraw[black] (S) circle (2pt) node[below right] {$N$};
    \filldraw[black] (K) circle (2pt) node[above left] {$M$};
    \filldraw[black] (Q) circle (2pt) node[below] {$P$};
   \filldraw[black] (k) node[above] {$y$};
    \filldraw[black] (x)  node[above] {$x $};
    \filldraw[black] (m) node[above] {$z$};
     
\end{tikzpicture}
}
\end{minipage}

\end{tcolorbox}
\vspace{4pt}


\begin{tcolorbox}[
    breakable,
    colback=white,
    colframe=blue!60!black,
    boxrule=0.8pt,
    arc=3pt,
    title={\textbf{Câu \#145} \quad \textbf{MÃ CÂU: \texttt{[L07\_C3\_B5\_03]}} \quad | \quad \textbf{Nhóm: \texttt{HINH\_118}} \quad | \quad \textbf{MapID: \texttt{7H3H5}} \hfill \textcolor{yellow!90!white}{[Hình gốc .pyc]}}
]
\textbf{Dạng toán:} Bài tổng hợp 3 \hfill \textit{(Tự luận -- Thông hiểu)}
\vspace{4pt}\hrule\vspace{6pt}
\begin{minipage}[t]{0.58\textwidth}
\vspace{0pt}
Cho hình vẽ, biết $x  \parallel y$; $GH \perp x $ và $\widehat{E_{1}}=131^{\circ}$.

a) Chứng minh $GH \perp y$.

b) Tính $\widehat{F_{1}}$.
\end{minipage}\hfill
\begin{minipage}[t]{0.39\textwidth}
\vspace{0pt}
\centering

\resizebox{\linewidth}{!}{%
\begin{tikzpicture}[scale=0.8, font=\footnotesize, line join=round, line cap=round, >=stealth]
    \def\a{8}
    \def\b{3}
    \path (0,0) coordinate (a) 
    (0:\a)  coordinate (a') node[right]{$x $}
    (-90:\b) coordinate (b) 
    ($(b) + (0:\a)$) coordinate (b') node[right]{$y$}
    ($(a)!0.8!(a')$) coordinate (A)
    ($(b)!0.8!(b')$) coordinate (B)
    ($(b)!0.2!(b')$) coordinate (C)
    ($(C)!1!60:(b)$) coordinate (m)
    (intersection of a--a' and C--m)coordinate (D);
    \pgfresetboundingbox
    \draw (a)--(a') (b)--(b')
    (C)--(D)--([turn]0:1.5) (D)--(C)--([turn]0:1.5) (A)--(B)--([turn]0:1.5) (B)--(A)--([turn]0:1.5)
    ($(C)+(100:0.4)$) node[scale=0.8]{$1$}
    ($(D)+(100:0.4)$) node[scale=0.8]{$1$};
  
    \foreach \x/\y/\z in {D/A/B}\pic[draw,angle radius=2.5mm]{right angle=\x--\y--\z};
        \filldraw[black] (A) circle (2pt) node[above right] {$G$};
    \filldraw[black] (B) circle (2pt) node[above right] {$H$};
    \filldraw[black] (C) circle (2pt) node[above right] {$F$};
   \filldraw[black] (D) circle (2pt) node[above right]{$E$};
\end{tikzpicture}
}
\end{minipage}

\end{tcolorbox}
\vspace{4pt}


\begin{tcolorbox}[
    breakable,
    colback=white,
    colframe=blue!60!black,
    boxrule=0.8pt,
    arc=3pt,
    title={\textbf{Câu \#146} \quad \textbf{MÃ CÂU: \texttt{[L07\_C3\_B5\_04]}} \quad | \quad \textbf{Nhóm: \texttt{HINH\_119}} \quad | \quad \textbf{MapID: \texttt{7H3V5}} \hfill \textcolor{yellow!90!white}{[Hình gốc .pyc]}}
]
\textbf{Dạng toán:} Bài tổng hợp 4 \hfill \textit{(Tự luận -- Vận dụng)}
\vspace{4pt}\hrule\vspace{6pt}
\begin{minipage}[t]{0.58\textwidth}
\vspace{0pt}
Cho hình vẽ biết $\widehat{G_{1}}=\widehat{H_{1}} = 62^{\circ}$, $\widehat{F_{1}} = 52^{\circ}$.

a) Chứng minh: $x  \parallel y$.

b) Cho $\widehat{F_{1}}=52^{\circ}$. Tính $\widehat{E_{1}}$, $\widehat{E_{2}}$.

c) Biết ${Fz}$ là tia phân giác của $\widehat{GFE}$. Tính $\widehat{F_{2}}$.
\end{minipage}\hfill
\begin{minipage}[t]{0.39\textwidth}
\vspace{0pt}
\centering

\resizebox{\linewidth}{!}{%
\begin{tikzpicture}[scale=1,>=stealth, font=\footnotesize, line join=round, line cap=round]
    \path 
    (0,0) coordinate (b)
    (6,0) coordinate (b')
    (0,2) coordinate (a)
    (6,2) coordinate (a')
    ;
    \coordinate (A) at ($(a)!0.25!(a')$);
    \coordinate (E) at ($(a)!0.6!(a')$);
    \coordinate (A') at ($(A)+({1*cos(70)},{1*sin(70)})$);
    \coordinate (M) at ($(E)+({1*cos(-50)},{1*sin(-50)})$);
    \coordinate (x) at ($(E)+({2.7*cos(-115)},{2.7*sin(-115)})$);
    \path 
    (intersection of b--b' and A--A') coordinate (B)
    (intersection of b--b' and E--M) coordinate (F)
    ;
    \draw ($(A)!-0.3!(B)$)--($(A)!1.2!(B)$);
    \draw (x)--(E)--([turn]0:0.7cm);
    \draw ($(E)!-0.3!(F)$)--($(E)!1.2!(F)$);
    \draw (a)node[above]{$x $}--(a')(b)node[above]{$y$}--(b');
    \node at (x)[right]{$z$};
    \foreach \x/\y/\z in {F/-25/1,A/35/1,B/35/1,F/155/2,E/-25/1,E/-83/2,E/-150/3}
    \fill (\x) circle (0pt) node[shift={(\y:3mm)}]{$\z$};
    \pic[draw,angle radius=0.5cm,double] {angle = x--E--F};
    \pic[draw,angle radius=0.6cm,double,mark=||] {angle = A--E--x};
        \filldraw[black] (A) circle (2pt) node[above left] {$G$};
    \filldraw[black] (B) circle (2pt) node[above left] {$H$};
    \filldraw[black] (E) circle (2pt) node[above] {$F$};
    \filldraw[black] (F) circle (2pt) node[above] {$E$};
\end{tikzpicture}
}
\end{minipage}

\end{tcolorbox}
\vspace{4pt}


\subsection{Chương 4 - Tam giác bằng nhau}

\begin{tcolorbox}[
    breakable,
    colback=white,
    colframe=blue!60!black,
    boxrule=0.8pt,
    arc=3pt,
    title={\textbf{Câu \#147} \quad \textbf{MÃ CÂU: \texttt{[L07\_C4\_B1\_06]}} \quad | \quad \textbf{Nhóm: \texttt{HINH\_120}} \quad | \quad \textbf{MapID: \texttt{7H4N1}} \hfill \textcolor{yellow!90!white}{[Hình gốc .pyc]}}
]
\textbf{Dạng toán:} Cho hình vẽ tìm số đo x, y. \hfill \textit{(Trắc nghiệm -- Nhận biết)}
\vspace{4pt}\hrule\vspace{6pt}
\begin{minipage}[t]{0.58\textwidth}
\vspace{0pt}
Tìm số đo ${x}$, ${y}$ trong hình vẽ.

\vspace{4pt}\noindent
\textbf{A.} $x= 112^{\circ}$ và $y= 64^{\circ}$ \quad \textbf{\textcolor{blue!75!black}{B. (Đ)}} $x= 116^{\circ}$ và $y= 66^{\circ}$ \quad \textbf{C.} $x= 116^{\circ}$ và $y= 69^{\circ}$ \quad \textbf{D.} $x= 119^{\circ}$ và $y= 66^{\circ}$
\end{minipage}\hfill
\begin{minipage}[t]{0.39\textwidth}
\vspace{0pt}
\centering

\resizebox{\linewidth}{!}{%
\begin{tikzpicture}
    % Define points
    \coordinate (B) at (0,0);
    \coordinate (C) at (4,0);
    \coordinate (A) at (3,4);
        \coordinate (D) at (-2,0);
    % Draw the triangle
    \draw (B) -- (A) -- (C) -- cycle;
        \draw (B) -- (D);
    % Add points with small circles
    \filldraw (B) circle (1pt); % Small circle at B
    \filldraw (C) circle (1pt); % Small circle at C
    \filldraw (A) circle (1pt); % Small circle at A
    
    % Add labels
    \node[above] at (A) {$E$};
    \node[below] at (B) {$F$};
    \node[right] at (C) {$G$};
    
    % Draw angles

    \node at (0.6,0.3) {$64^\circ$};

    \node at (3.7,0.2) {$y$};
    \node at (2.9,3.3) {$50^\circ$};
    
    % Mark angles with labels
    \node[below] at (-0.3,0.4) {$x$};
\end{tikzpicture}
}
\end{minipage}

\vspace{3pt}\hrule\vspace{3pt}\footnotesize\textcolor{gray!80!black}{\textbf{Các mã dùng chung hình (HINH\_120):} \texttt{L07\_C4\_B1\_05}}
\end{tcolorbox}
\vspace{4pt}


\begin{tcolorbox}[
    breakable,
    colback=white,
    colframe=blue!60!black,
    boxrule=0.8pt,
    arc=3pt,
    title={\textbf{Câu \#148} \quad \textbf{MÃ CÂU: \texttt{[L07\_C4\_B1\_19]}} \quad | \quad \textbf{Nhóm: \texttt{HINH\_121}} \quad | \quad \textbf{MapID: \texttt{7H4N1}} \hfill \textcolor{yellow!90!white}{[Hình gốc .pyc]}}
]
\textbf{Dạng toán:} Cho số đo góc ngoài tìm các góc trong tam giác \hfill \textit{(Trắc nghiệm -- Nhận biết)}
\vspace{4pt}\hrule\vspace{6pt}
\begin{minipage}[t]{0.58\textwidth}
\vspace{0pt}
Tìm số đo ${x}$, ${y}$ trong hình vẽ.

\vspace{4pt}\noindent
\textbf{\textcolor{blue!75!black}{A. (Đ)}} $x=42^{\circ}$ và $y=85^{\circ}$ \quad \textbf{B.} $x=42^{\circ}$ và $y=86^{\circ}$ \quad \textbf{C.} $x=46^{\circ}$ và $y=85^{\circ}$ \quad \textbf{D.} $x=45^{\circ}$ và $y=87^{\circ}$
\end{minipage}\hfill
\begin{minipage}[t]{0.39\textwidth}
\vspace{0pt}
\centering

\resizebox{\linewidth}{!}{%
\begin{tikzpicture}
    % Vẽ các điểm của tam giác
    \coordinate [label=left:A] (A) at (0, 0);
    \coordinate [label=above:B] (B) at (-2, 4);
    \coordinate [label=right:C] (C) at (6, 0);
        \coordinate [label=above right:D] (D) at (3, 3/2);
            \coordinate [label=above right:E] (E) at (-4, 10/2);
    % Tạo điểm D nằm giữa B và C

    
    % Vẽ các cạnh của tam giác
    \draw (A) -- (B) -- (C) -- cycle;
    \draw (A) -- (D);
        \draw (E) -- (B);
    % Đánh dấu điểm D
    \fill (D) circle (1.2pt);
        \fill (A) circle (1.2pt);
            \fill (B) circle (1.2pt);
                \fill (C) circle (1.2pt);
                    \fill (E) circle (1.2pt);
      \node at (0.1,0.3) {$y^{\circ}$};
    \node at (0.9,0.2) {$24^{\circ}$};
        \node at (2.3,1.5) {$66^{\circ}$};
            \node at (5,0.3) {$x^{\circ}$};
                \node at (-2.3,3.7) {$151^{\circ}$};
\end{tikzpicture}
}
\end{minipage}

\end{tcolorbox}
\vspace{4pt}


\begin{tcolorbox}[
    breakable,
    colback=white,
    colframe=blue!60!black,
    boxrule=0.8pt,
    arc=3pt,
    title={\textbf{Câu \#149} \quad \textbf{MÃ CÂU: \texttt{[L07\_C4\_B1\_05]}} \quad | \quad \textbf{Nhóm: \texttt{HINH\_120}} \quad | \quad \textbf{MapID: \texttt{7H4N1}} \hfill \textcolor{yellow!90!white}{[Hình gốc .pyc]}}
]
\textbf{Dạng toán:} Cho hình vẽ tìm số đo x, y \hfill \textit{(Tự luận -- Nhận biết)}
\vspace{4pt}\hrule\vspace{6pt}
\begin{minipage}[t]{0.58\textwidth}
\vspace{0pt}
Tìm số đo ${x}$, ${y}$ trong hình vẽ.
\end{minipage}\hfill
\begin{minipage}[t]{0.39\textwidth}
\vspace{0pt}
\centering

\resizebox{\linewidth}{!}{%
\begin{tikzpicture}
    % Define points
    \coordinate (B) at (0,0);
    \coordinate (C) at (4,0);
    \coordinate (A) at (3,4);
        \coordinate (D) at (-2,0);
    % Draw the triangle
    \draw (B) -- (A) -- (C) -- cycle;
        \draw (B) -- (D);
    % Add points with small circles
    \filldraw (B) circle (1pt); % Small circle at B
    \filldraw (C) circle (1pt); % Small circle at C
    \filldraw (A) circle (1pt); % Small circle at A
    
    % Add labels
    \node[above] at (A) {$P$};
    \node[below] at (B) {$Q$};
    \node[right] at (C) {$R$};
    
    % Draw angles

    \node at (0.6,0.3) {$60^\circ$};

    \node at (3.7,0.2) {$y$};
    \node at (2.9,3.3) {$48^\circ$};
    
    % Mark angles with labels
    \node[below] at (-0.3,0.4) {$x$};
\end{tikzpicture}
}
\end{minipage}

\vspace{3pt}\hrule\vspace{3pt}\footnotesize\textcolor{gray!80!black}{\textbf{Các mã dùng chung hình (HINH\_120):} \texttt{L07\_C4\_B1\_06}}
\end{tcolorbox}
\vspace{4pt}


\begin{tcolorbox}[
    breakable,
    colback=white,
    colframe=blue!60!black,
    boxrule=0.8pt,
    arc=3pt,
    title={\textbf{Câu \#150} \quad \textbf{MÃ CÂU: \texttt{[L07\_C4\_B1\_17]}} \quad | \quad \textbf{Nhóm: \texttt{HINH\_122}} \quad | \quad \textbf{MapID: \texttt{7H4V1}} \hfill \textcolor{yellow!90!white}{[Hình gốc .pyc]}}
]
\textbf{Dạng toán:} Cho phân giác AD của tam giác ABC, tính góc ADB biết B-C=a \hfill \textit{(Tự luận -- Vận dụng)}
\vspace{4pt}\hrule\vspace{6pt}
\begin{minipage}[t]{0.58\textwidth}
\vspace{0pt}
Cho tam giác ${ABC}$, tia phân giác ${AD}$ của góc $\widehat{A}$ cắt ${BC}$ tại ${D}$.

 Tính $\widehat{ADB}$ biết $\widehat{B}-\widehat{C}=46^{\circ}  $.
\end{minipage}\hfill
\begin{minipage}[t]{0.39\textwidth}
\vspace{0pt}
\centering

\resizebox{\linewidth}{!}{%
\begin{tikzpicture}
    % Points
    \coordinate [label=above:$A$] (A) at (0,3);
    \coordinate [label=below left:$B$] (B) at (-2,0);
    \coordinate [label=below:$D$] (D) at (0,0);
    \coordinate [label=below right:$C$] (C) at (3,0);
    
    % Triangle sides
    \draw [thick] (B) -- (A) -- (C) -- cycle;
    
    % Altitude
    \draw [thick] (A) -- (D);
    
    % Angles
    \draw [thick] (B) -- (A) -- (D);
    \draw [thick] (C) -- (A) -- (D);
    
    % Angle marks

    \node at (-0.1,2.6) {$1$};
    \node at (0.2,2.6) {$2$};
    
    % Points
    \filldraw [black] (A) circle (1.5pt);
    \filldraw [black] (B) circle (1.5pt);
    \filldraw [black] (C) circle (1.5pt);
    \filldraw [black] (D) circle (1.5pt);
\end{tikzpicture}
}
\end{minipage}

\end{tcolorbox}
\vspace{4pt}


\begin{tcolorbox}[
    breakable,
    colback=white,
    colframe=blue!60!black,
    boxrule=0.8pt,
    arc=3pt,
    title={\textbf{Câu \#151} \quad \textbf{MÃ CÂU: \texttt{[L07\_C4\_B3\_01]}} \quad | \quad \textbf{Nhóm: \texttt{HINH\_123}} \quad | \quad \textbf{MapID: \texttt{7H4N3}} \hfill \textcolor{yellow!90!white}{[Hình gốc .pyc]}}
]
\textbf{Dạng toán:} Viết các tam giác bằng nhau theo TH c.c.c \hfill \textit{(Tự luận -- Nhận biết)}
\vspace{4pt}\hrule\vspace{6pt}
\begin{minipage}[t]{0.58\textwidth}
\vspace{0pt}
Với mỗi hình hãy viết các cặp tam giác bằng nhau.
\end{minipage}\hfill
\begin{minipage}[t]{0.39\textwidth}
\vspace{0pt}
\centering

\resizebox{\linewidth}{!}{%
\begin{tikzpicture}
    
    % Coordinates of points
    \coordinate (A) at (-4,0);
    \coordinate (B) at (-2,0);
    \coordinate (C) at (-1,2);
    \coordinate (D) at (-1,-2);
    
    % Draw the triangles
    \draw (A) -- (B) -- (C) -- cycle; % Triangle ABC
    \draw (A) -- (B) -- (D) -- cycle; % Triangle ACD
    

    % Labels
    \node[below left] at (A) {$M$};
    \node[ right] at (B) {$N$};
    \node[right] at (C) {$P$};
    \node[below right] at (D) {$Q$};
       \filldraw [black] (A) circle (1.5pt);
    \filldraw [black] (B) circle (1.5pt);
    \filldraw [black] (C) circle (1.5pt);
    \filldraw [black] (D) circle (1.5pt);
    \draw
    ($(A)!0.5!(C)$)node[rotate=20]{$/$} ($(A)!0.5!(D)$)node[rotate=20]{$/$}
    ($(B)!0.5!(D)$)node[rotate=90]{$//$} ($(B)!0.5!(C)$)node[rotate=-90]{$//$};
    
    
    
        % Coordinates of points
    \coordinate (E) at (3,2);
    \coordinate (F) at (2,0);
    \coordinate (M) at (5,0);
    \coordinate (N) at (6,2);
    
    % Draw the triangles
    \draw (E) -- (F) -- (M) --(N)-- cycle; 
\draw (F)--(N);
    % Labels
    \node[above left] at (E) {$E$};
    \node[ left] at (F) {$F$};
    \node[right] at (M) {$I$};
    \node[above right] at (N){$J$};
    \filldraw [black] (E) circle (1.5pt);
    \filldraw [black] (F) circle (1.5pt);
    \filldraw [black] (M) circle (1.5pt);
    \filldraw [black] (N) circle (1.5pt);
    \draw
    ($(E)!0.5!(F)$)node[rotate=20]{$/$} ($(N)!0.5!(M)$)node[rotate=20]{$/$}
    ($(E)!0.5!(N)$)node[rotate=50]{$//$} ($(F)!0.5!(M)$)node[rotate=50]{$//$};
\end{tikzpicture}
}
\end{minipage}

\end{tcolorbox}
\vspace{4pt}


\begin{tcolorbox}[
    breakable,
    colback=white,
    colframe=blue!60!black,
    boxrule=0.8pt,
    arc=3pt,
    title={\textbf{Câu \#152} \quad \textbf{MÃ CÂU: \texttt{[L07\_C4\_B3\_02]}} \quad | \quad \textbf{Nhóm: \texttt{HINH\_124}} \quad | \quad \textbf{MapID: \texttt{7H4H3}} \hfill \textcolor{yellow!90!white}{[Hình gốc .pyc]}}
]
\textbf{Dạng toán:} Viết các tam giác bằng nhau \hfill \textit{(Tự luận -- Thông hiểu)}
\vspace{4pt}\hrule\vspace{6pt}
\begin{minipage}[t]{0.58\textwidth}
\vspace{0pt}
Hãy viết các tam giác bằng nhau.
\end{minipage}\hfill
\begin{minipage}[t]{0.39\textwidth}
\vspace{0pt}
\centering

\resizebox{\linewidth}{!}{%
\begin{tikzpicture}
            
            % Coordinates of points
            \coordinate (A) at (-2,2);
            \coordinate (B) at (-2,-2);
            \coordinate (C) at (2,-2);
            \coordinate (D) at (2,2);
            
            % Draw the triangles
            \draw (A) -- (B) -- (C) -- (D)--cycle; % Triangle ABC
            
            % Labels
            \node[above left] at (A) {$E$};
            \node[ left] at (B) {$F$};
            \node[right] at (C) {$K$};
            \node[above right] at (D) {$H$};
            \filldraw [black] (A) circle (1.5pt);
            \filldraw [black] (B) circle (1.5pt);
            \filldraw [black] (C) circle (1.5pt);
            \filldraw [black] (D) circle (1.5pt);
            % Coordinates of points
            \coordinate (E) at (-1,2);
            \coordinate (F) at (-2,-1);
            \coordinate (M) at (1,-2);
            \coordinate (N) at (2,1);
            
            % Draw the triangles
            \draw (E) -- (F) -- (M) --(N)-- cycle; 
            \draw (F)--(N);
            % Labels
            \node[above left] at (E) {$M$};
            \node[ left] at (F) {$N$};
            \node[below] at (M) {$I$};
            \node[above right] at (N) {$J$};
            \filldraw [black] (E) circle (1.5pt);
            \filldraw [black] (F) circle (1.5pt);
            \filldraw [black] (M) circle (1.5pt);
            \filldraw [black] (N) circle (1.5pt);
            \draw
            ($(D)!0.5!(E)$)node[rotate=20]{$/$} ($(A)!0.5!(F)$)node[rotate=50]{$/$}
            ($(B)!0.5!(M)$)node[rotate=50]{$/$} ($(C)!0.5!(N)$)node[rotate=50]{$/$};
            \draw
            ($(F)!0.5!(E)$)node[rotate=50]{$//$} ($(M)!0.5!(F)$)node[rotate=20]{$//$}
            ($(N)!0.5!(M)$)node[rotate=50]{$//$} ($(E)!0.5!(N)$)node[rotate=20]{$//$};
        \end{tikzpicture}
}
\end{minipage}

\end{tcolorbox}
\vspace{4pt}


\begin{tcolorbox}[
    breakable,
    colback=white,
    colframe=blue!60!black,
    boxrule=0.8pt,
    arc=3pt,
    title={\textbf{Câu \#153} \quad \textbf{MÃ CÂU: \texttt{[L07\_C4\_B3\_03]}} \quad | \quad \textbf{Nhóm: \texttt{HINH\_125}} \quad | \quad \textbf{MapID: \texttt{7H4N3}} \hfill \textcolor{yellow!90!white}{[Hình gốc .pyc]}}
]
\textbf{Dạng toán:} Chứng minh hai tam giác bằng nhau \hfill \textit{(Tự luận -- Nhận biết)}
\vspace{4pt}\hrule\vspace{6pt}
\begin{minipage}[t]{0.58\textwidth}
\vspace{0pt}
Cho hình vẽ bên. Chứng minh 

a) $\Delta{IJQ} = \Delta{IJK}$ 

b) ${IJ}$ là tia phân giác của $\widehat{ KIQ} $
\end{minipage}\hfill
\begin{minipage}[t]{0.39\textwidth}
\vspace{0pt}
\centering

\resizebox{\linewidth}{!}{%
\begin{tikzpicture}
    
    % Coordinates of points
    \coordinate (A) at (-4,0);
    \coordinate (B) at (-2,0);
    \coordinate (C) at (-1,2);
    \coordinate (D) at (-1,-2);
    
    % Draw the triangles
    \draw (A) -- (B) -- (C) -- cycle; % Triangle ABC
    \draw (A) -- (B) -- (D) -- cycle; % Triangle ACD
    

    % Labels
    \node[below left] at (A) {$I$};
    \node[ right] at (B) {$J$};
    \node[right] at (C) {$Q$};
    \node[below right] at (D) {$K$};
       \filldraw [black] (A) circle (1.5pt);
    \filldraw [black] (B) circle (1.5pt);
    \filldraw [black] (C) circle (1.5pt);
    \filldraw [black] (D) circle (1.5pt);
    \draw
    ($(A)!0.5!(C)$)node[rotate=20]{$/$} ($(A)!0.5!(D)$)node[rotate=20]{$/$}
    ($(B)!0.5!(D)$)node[rotate=90]{$//$} ($(B)!0.5!(C)$)node[rotate=-90]{$//$};
    \end{tikzpicture}
}
\end{minipage}

\vspace{3pt}\hrule\vspace{3pt}\footnotesize\textcolor{gray!80!black}{\textbf{Các mã dùng chung hình (HINH\_125):} \texttt{L07\_C4\_B3\_04}}
\end{tcolorbox}
\vspace{4pt}


\begin{tcolorbox}[
    breakable,
    colback=white,
    colframe=blue!60!black,
    boxrule=0.8pt,
    arc=3pt,
    title={\textbf{Câu \#154} \quad \textbf{MÃ CÂU: \texttt{[L07\_C4\_B3\_04]}} \quad | \quad \textbf{Nhóm: \texttt{HINH\_125}} \quad | \quad \textbf{MapID: \texttt{7H4H3}} \hfill \textcolor{yellow!90!white}{[Hình gốc .pyc]}}
]
\textbf{Dạng toán:} Chứng minh hai tam giác bằng nhau \hfill \textit{(Tự luận -- Thông hiểu)}
\vspace{4pt}\hrule\vspace{6pt}
\begin{minipage}[t]{0.58\textwidth}
\vspace{0pt}
Cho hình vẽ bên. Chứng minh 

a) $\Delta{IJQ} = \Delta{IJK}$ 

b) $\widehat{IQJ} = \Delta{IKJ}$c) ${IJ}$ là tia phân giác của $\widehat{ KIQ} $
\end{minipage}\hfill
\begin{minipage}[t]{0.39\textwidth}
\vspace{0pt}
\centering

\resizebox{\linewidth}{!}{%
\begin{tikzpicture}
    
    % Coordinates of points
    \coordinate (A) at (-4,0);
    \coordinate (B) at (-2,0);
    \coordinate (C) at (-1,2);
    \coordinate (D) at (-1,-2);
    
    % Draw the triangles
    \draw (A) -- (B) -- (C) -- cycle; % Triangle ABC
    \draw (A) -- (B) -- (D) -- cycle; % Triangle ACD
    

    % Labels
    \node[below left] at (A) {$I$};
    \node[ right] at (B) {$J$};
    \node[right] at (C) {$Q$};
    \node[below right] at (D) {$K$};
       \filldraw [black] (A) circle (1.5pt);
    \filldraw [black] (B) circle (1.5pt);
    \filldraw [black] (C) circle (1.5pt);
    \filldraw [black] (D) circle (1.5pt);
    \draw
    ($(A)!0.5!(C)$)node[rotate=20]{$/$} ($(A)!0.5!(D)$)node[rotate=20]{$/$}
    ($(B)!0.5!(D)$)node[rotate=90]{$//$} ($(B)!0.5!(C)$)node[rotate=-90]{$//$};
     \end{tikzpicture}
}
\end{minipage}

\vspace{3pt}\hrule\vspace{3pt}\footnotesize\textcolor{gray!80!black}{\textbf{Các mã dùng chung hình (HINH\_125):} \texttt{L07\_C4\_B3\_03}}
\end{tcolorbox}
\vspace{4pt}


\begin{tcolorbox}[
    breakable,
    colback=white,
    colframe=blue!60!black,
    boxrule=0.8pt,
    arc=3pt,
    title={\textbf{Câu \#155} \quad \textbf{MÃ CÂU: \texttt{[L07\_C4\_B3\_05]}} \quad | \quad \textbf{Nhóm: \texttt{HINH\_126}} \quad | \quad \textbf{MapID: \texttt{7H4H3}} \hfill \textcolor{yellow!90!white}{[Hình gốc .pyc]}}
]
\textbf{Dạng toán:} Bài toán liên quan đến tam giác cân. \hfill \textit{(Tự luận -- Thông hiểu)}
\vspace{4pt}\hrule\vspace{6pt}
\begin{minipage}[t]{0.58\textwidth}
\vspace{0pt}
Cho tam giác ${EFG}$ có $EF=EG$. Gọi ${H}$ là trung điểm của cạnh ${FG}$. Chứng minh 

a)  $\widehat{F} = \Delta{G}$ 

b) $EH$ là phân giác của góc $\widehat{FEG}$ c) $EH \bot FG$
\end{minipage}\hfill
\begin{minipage}[t]{0.39\textwidth}
\vspace{0pt}
\centering

\resizebox{\linewidth}{!}{%
\begin{tikzpicture}
    
    % Coordinates of points
    \coordinate (A) at (0,4);
    \coordinate (B) at (-2,0);
    \coordinate (C) at (2,0);
    \coordinate (M) at (0,0);
    
    % Draw the triangles
    \draw (A) -- (B) -- (C) -- cycle; % Triangle ABC
    \draw (A) -- (M); 
    

    % Labels
    \node[above] at (A) {$E$};
    \node[ left] at (B) {$F$};
    \node[right] at (C) {$G$};
    \node[below ] at (M) {$H$};
       \filldraw [black] (A) circle (1.5pt);
    \filldraw [black] (B) circle (1.5pt);
    \filldraw [black] (C) circle (1.5pt);
    \filldraw [black] (M) circle (1.5pt);
    \draw
    ($(A)!0.5!(C)$)node[rotate=20]{$/$} ($(A)!0.5!(B)$)node[rotate=20]{$/$};
    \end{tikzpicture}
}
\end{minipage}

\end{tcolorbox}
\vspace{4pt}


\begin{tcolorbox}[
    breakable,
    colback=white,
    colframe=blue!60!black,
    boxrule=0.8pt,
    arc=3pt,
    title={\textbf{Câu \#156} \quad \textbf{MÃ CÂU: \texttt{[L07\_C4\_B3\_08]}} \quad | \quad \textbf{Nhóm: \texttt{HINH\_127}} \quad | \quad \textbf{MapID: \texttt{7H4V3}} \hfill \textcolor{yellow!90!white}{[Hình gốc .pyc]}}
]
\textbf{Dạng toán:} Bài toán tổng hợp 2 (3 ý) \hfill \textit{(Tự luận -- Vận dụng)}
\vspace{4pt}\hrule\vspace{6pt}
\begin{minipage}[t]{0.58\textwidth}
\vspace{0pt}
Cho tam giác ${ABC}$ có ${AB=AC}$. Lấy hai điểm ${D}$, ${E}$ thuộc cạnh ${BC}$ sao cho $BD=DE=EC$. Biết $AD=AE$.

a) Chứng minh $\widehat{ EAB }= \widehat{DAC}$ 

b) Gọi ${M}$ là trung điểm của ${BC}$. Chứng minh ${AM}$ là phân giác của $\widehat{DAE}$

c) Giả sử $\widehat {DA E} =60^{\circ}$. Tính các góc còn lại của tam giác ${DAE }$
\end{minipage}\hfill
\begin{minipage}[t]{0.39\textwidth}
\vspace{0pt}
\centering

\resizebox{\linewidth}{!}{%
\begin{tikzpicture}
    % Vẽ tam giác ABC cân tại A
    \coordinate (A) at (0,5);  % Đỉnh A
    \coordinate (B) at (-3,0); % Đỉnh B
    \coordinate (C) at (3,0);  % Đỉnh C
    
    % Tính các điểm D và E trên BC sao cho BD = DE = EC
    \coordinate (D) at ($(B)!1/3!(C)$);  % Điểm D chia BC thành 1/3
    \coordinate (E) at ($(B)!2/3!(C)$);  % Điểm E chia BC thành 2/3
    
    % Tính trung điểm M của BC
    \coordinate (M) at ($(B)!0.5!(C)$);  % Điểm M là trung điểm của BC
    
    % Vẽ các cạnh của tam giác
    \draw (A) -- (B) -- (C) -- cycle;
        \draw (A) -- (D);
            \draw (A) -- (E);
    % Vẽ các điểm D, E và M
    \draw (B) -- (D) -- (E) -- (C);
    \draw[thick] (A) -- (M);  % Vẽ đường nối A với M
    
    % Gán nhãn cho các đỉnh và điểm
    \node[above] at (A) {$A$};
    \node[below left] at (B) {$B$};
    \node[below right] at (C) {$C$};
    \node[below] at (D) {$D$};
    \node[below] at (E) {$E$};
    \node[below] at (M) {$M$};  % Gán nhãn cho trung điểm M
     \draw
    ($(A)!0.5!(C)$)node[rotate=90]{$///$} ($(A)!0.5!(B)$)node[rotate=40]{$///$};
     \draw
    ($(B)!0.5!(D)$)node[rotate=20]{$/$} ($(D)!0.7!(E)$)node[rotate=20]{$/$} ($(C)!0.5!(E)$)node[rotate=20]{$/$};
     \draw
    ($(A)!0.5!(D)$)node[rotate=90]{$//$} ($(A)!0.5!(E)$)node[rotate=90]{$//$};
    
\end{tikzpicture}
}
\end{minipage}

\end{tcolorbox}
\vspace{4pt}


\begin{tcolorbox}[
    breakable,
    colback=white,
    colframe=blue!60!black,
    boxrule=0.8pt,
    arc=3pt,
    title={\textbf{Câu \#157} \quad \textbf{MÃ CÂU: \texttt{[L07\_C4\_B4\_05]}} \quad | \quad \textbf{Nhóm: \texttt{HINH\_128}} \quad | \quad \textbf{MapID: \texttt{7H4H4}} \hfill \textcolor{yellow!90!white}{[Hình gốc .pyc]}}
]
\textbf{Dạng toán:} CM hai đoạn thẳng, hai góc bằng nhau. \hfill \textit{(Tự luận -- Thông hiểu)}
\vspace{4pt}\hrule\vspace{6pt}
\begin{minipage}[t]{0.58\textwidth}
\vspace{0pt}
Cho tam giác ${O B C }$ có $O B <O C $. Phân giác của gốc $\widehat{O }$ cắt cạnh ${B C }$ tại điểm ${D }$. Trên cạnh ${O C }$ lấy điểm ${E }$ sao cho $O E =O B $. Chứng minh 

   a) $\Delta{O B D } =\Delta{O E D }$ 

b) ${D O }$ là tia phân giác của $\widehat{B D E }$. Từ đó suy ra $\widehat{O B C }> \widehat{O C B }$
\end{minipage}\hfill
\begin{minipage}[t]{0.39\textwidth}
\vspace{0pt}
\centering

\resizebox{\linewidth}{!}{%
\begin{tikzpicture}
    % Vẽ các điểm A, B, C
    \coordinate [label=above: {$O $}] (A) at (0, 3);
    \coordinate [label=below left:{$B $}] (B) at (-2, 0);
    \coordinate [label=below right:{$C $}] (C) at (4, 0);
    
    % Vẽ các cạnh của tam giác
    \draw (A) -- (B) -- (C) -- cycle;
    
    % Ghi chú các cạnh
    \draw (A) -- (B) ;
    \draw (A) -- (C) ;
    \draw (B) -- (C) ;
    
    % Tính tọa độ của điểm D (phân giác cắt BC tại D)
    \coordinate [label=below right:{$D $}] (D) at (0.51, 0);
    
    % Vẽ đường phân giác AD
    \draw (A) -- (D);
    
    % Tính tọa độ của điểm E (AE = AB)
    \coordinate [label=above right:{$E $}] (E) at (2.88, 0.84);
    
    % Vẽ đoạn thẳng AE
    \draw (D) -- (E) ;
    
\end{tikzpicture}
}
\end{minipage}

\end{tcolorbox}
\vspace{4pt}


\begin{tcolorbox}[
    breakable,
    colback=white,
    colframe=blue!60!black,
    boxrule=0.8pt,
    arc=3pt,
    title={\textbf{Câu \#158} \quad \textbf{MÃ CÂU: \texttt{[L07\_C4\_B4\_07]}} \quad | \quad \textbf{Nhóm: \texttt{HINH\_129}} \quad | \quad \textbf{MapID: \texttt{7H4V4}} \hfill \textcolor{yellow!90!white}{[Hình gốc .pyc]}}
]
\textbf{Dạng toán:} Bài tổng hợp 4 ý \hfill \textit{(Tự luận -- Vận dụng)}
\vspace{4pt}\hrule\vspace{6pt}
\begin{minipage}[t]{0.58\textwidth}
\vspace{0pt}
Cho tam giác ${PQR}$ có $\widehat{P} =90^{\circ}$, trên cạnh ${QR}$ lấy điểm ${H}$ sao cho $QP=QH$. Tia phân giác $\widehat{Q}$ cắt ${PR}$ tại ${T}$ 

a) Chứng minh $\Delta{PQT} =\Delta {HQT}$

b) Chứng minh $TP=TH$ 

c) Xác định số đo $\widehat{QHT}$ 

d) Xác định độ lớn $\widehat{Q}$ để $\widehat{HTQ}=\widehat{HTR}$
\end{minipage}\hfill
\begin{minipage}[t]{0.39\textwidth}
\vspace{0pt}
\centering

\resizebox{\linewidth}{!}{%
\begin{tikzpicture}
  % Vẽ các điểm A, B, C
\coordinate [label=below left: {$P$}] (A) at (0, 0);
\coordinate [label=below right:{$Q$}] (B) at (3, 0);
\coordinate [label=above:{$R$}] (C) at (0, 4);
\coordinate [label=above right:{$H$}] (E) at (1.2, 2.4);
\coordinate [label=below left:{$T$}] (D) at (0, 1.5);
% Vẽ các cạnh của tam giác
    \path   (A) pic[draw=black,fill=gray!25, angle radius = 5] {right angle = B--A--C};
\draw (A) -- (B) -- (C) -- cycle;
  \draw
 ($(B)!0.7!(E)$)node[rotate=20]{$/$} ($(B)!0.5!(A)$)node[rotate=20]{$/$};
\draw (B) -- (D);
% Vẽ đoạn thẳng AE
\draw (A) -- (E) ;
    \draw (D) -- (E) ;
\end{tikzpicture}
}
\end{minipage}

\end{tcolorbox}
\vspace{4pt}


\begin{tcolorbox}[
    breakable,
    colback=white,
    colframe=blue!60!black,
    boxrule=0.8pt,
    arc=3pt,
    title={\textbf{Câu \#159} \quad \textbf{MÃ CÂU: \texttt{[L07\_C4\_B4\_08]}} \quad | \quad \textbf{Nhóm: \texttt{HINH\_130}} \quad | \quad \textbf{MapID: \texttt{7H4H4}} \hfill \textcolor{yellow!90!white}{[Hình gốc .pyc]}}
]
\textbf{Dạng toán:} Bài tổng hợp 2 ý \hfill \textit{(Tự luận -- Thông hiểu)}
\vspace{4pt}\hrule\vspace{6pt}
\begin{minipage}[t]{0.58\textwidth}
\vspace{0pt}
Cho $\Delta{QKH}$, ${O}$ kà trung điểm của ${KH}$. Trên tia đối của tia ${OQ}$ lấy điểm ${I}$ sao cho $OI=OQ$. Chứng minh 

  a) $\Delta{QKO}=\Delta{IHO} $ 

b) $QK=HI$ và $QH//KI$
\end{minipage}\hfill
\begin{minipage}[t]{0.39\textwidth}
\vspace{0pt}
\centering

\resizebox{\linewidth}{!}{%
\begin{tikzpicture}
  % Vẽ các điểm A, B, C
\coordinate [label=above: {$Q$}] (A) at (0, 4);
\coordinate [label=left:{$K$}] (B) at (-1, 0);
\coordinate [label=right:{$H$}] (C) at (3, 0);
\coordinate [label=below left:{$O$}] (M) at (1, 0);
\coordinate [label=below left:{$I$}] (E) at (2, -4);
% Vẽ các cạnh của tam giác
\draw (A) -- (B) -- (C) -- cycle;
  \draw
 ($(M)!0.7!(E)$)node[rotate=90]{$/$} ($(M)!0.5!(A)$)node[rotate=90]{$/$};
 \draw
 ($(M)!0.7!(B)$)node[rotate=40]{$//$} ($(M)!0.5!(C)$)node[rotate=40]{$//$};
% Vẽ đoạn thẳng AE
\draw (A) -- (M)--(E) ;
   
\end{tikzpicture}
}
\end{minipage}

\end{tcolorbox}
\vspace{4pt}


\begin{tcolorbox}[
    breakable,
    colback=white,
    colframe=blue!60!black,
    boxrule=0.8pt,
    arc=3pt,
    title={\textbf{Câu \#160} \quad \textbf{MÃ CÂU: \texttt{[L07\_C4\_B5\_02]}} \quad | \quad \textbf{Nhóm: \texttt{HINH\_131}} \quad | \quad \textbf{MapID: \texttt{7H4N5}} \hfill \textcolor{yellow!90!white}{[Hình gốc .pyc]}}
]
\textbf{Dạng toán:} Nhận biết trường hợp 2 tam giác bằng nhau c.g.c \hfill \textit{(Tự luận -- Nhận biết)}
\vspace{4pt}\hrule\vspace{6pt}
\begin{minipage}[t]{0.58\textwidth}
\vspace{0pt}
Với mỗi hình hãy viết các cặp tam giác bằng nhau.
\end{minipage}\hfill
\begin{minipage}[t]{0.39\textwidth}
\vspace{0pt}
\centering

\resizebox{\linewidth}{!}{%
\begin{tikzpicture}
            
            % Coordinates of points
            \coordinate (A) at (-4,0);
            \coordinate (B) at (-2,0);
            \coordinate (C) at (-1,2);
            \coordinate (D) at (-1,-2);
            
            % Draw the triangles
            \draw (A) -- (B) -- (C) -- cycle; % Triangle ABC
            \draw (A) -- (B) -- (D) -- cycle; % Triangle ACD
            
            
            % Labels
            \node[below left] at (A) {$E$};
            \node[ right] at (B) {$F$};
            \node[right] at (C) {$G$};
            \node[below right] at (D) {$H$};
            \filldraw [black] (A) circle (1.5pt);
            \filldraw [black] (B) circle (1.5pt);
            \filldraw [black] (C) circle (1.5pt);
            \filldraw [black] (D) circle (1.5pt);
            % Coordinates of points
            \coordinate (E) at (3,2);
            \coordinate (F) at (2,0);
            \coordinate (M) at (5,0);
            \coordinate (N) at (6,2);
            
            % Draw the triangles
            \draw (E) -- (F) -- (M) --(N)-- cycle; 
            \draw (F)--(N);
            % Labels
            \node[above left] at (E) {$M$};
            \node[ left] at (F) {$N$};
            \node[right] at (M) {$I$};
            \node[above right] at (N){$J$};
            \filldraw [black] (E) circle (1.5pt);
            \filldraw [black] (F) circle (1.5pt);
            \filldraw [black] (M) circle (1.5pt);
            \filldraw [black] (N) circle (1.5pt);
        \draw 
        (A) pic[draw,angle radius = 15] {angle = B--A--C} ;
        \draw 
        (A) pic[draw,angle radius = 13] {angle = D--A--B} ;
            \draw 
        (B) pic[draw,angle radius = 8] {angle = C--B--A} pic[draw,angle radius = 10] {angle = C--B--A} ;
        \draw 
        (B) pic[draw,angle radius = 13] {angle = A--B--D}  pic[draw,angle radius = 11] {angle = A--B--D}    ;
        
            \draw 
        (N) pic[draw,angle radius = 15] {angle = E--N--F} ;
        \draw 
        (N) pic[draw,angle radius = 13] {angle = M--F--N} ;
        \draw 
        (F) pic[draw,angle radius = 8] {angle = N--F--E} pic[draw,angle radius = 10] {angle = N--F--E} ;
        \draw 
        (N) pic[draw,angle radius = 13] {angle = F--N--M}  pic[draw,angle radius = 11] {angle = F--N--M}    ;
        \end{tikzpicture}
}
\end{minipage}

\vspace{3pt}\hrule\vspace{3pt}\footnotesize\textcolor{gray!80!black}{\textbf{Các mã dùng chung hình (HINH\_131):} \texttt{L07\_C4\_B5\_03}}
\end{tcolorbox}
\vspace{4pt}


\begin{tcolorbox}[
    breakable,
    colback=white,
    colframe=blue!60!black,
    boxrule=0.8pt,
    arc=3pt,
    title={\textbf{Câu \#161} \quad \textbf{MÃ CÂU: \texttt{[L07\_C4\_B5\_03]}} \quad | \quad \textbf{Nhóm: \texttt{HINH\_131}} \quad | \quad \textbf{MapID: \texttt{7H4N5}} \hfill \textcolor{yellow!90!white}{[Hình gốc .pyc]}}
]
\textbf{Dạng toán:} Nhận dạng 2 tam giác bằng nhau \hfill \textit{(Tự luận -- Nhận biết)}
\vspace{4pt}\hrule\vspace{6pt}
\begin{minipage}[t]{0.58\textwidth}
\vspace{0pt}
Với mỗi hình hãy viết các cặp tam giác bằng nhau.
\end{minipage}\hfill
\begin{minipage}[t]{0.39\textwidth}
\vspace{0pt}
\centering

\resizebox{\linewidth}{!}{%
\begin{tikzpicture}
                
                % Coordinates of points
                \coordinate (A) at (-4,0);
                \coordinate (B) at (-2,0);
                \coordinate (C) at (-1,2);
                \coordinate (D) at (-1,-2);
                
                % Draw the triangles
                \draw (A) -- (B) -- (C) -- cycle; % Triangle ABC
                \draw (A) -- (B) -- (D) -- cycle; % Triangle ACD
                
                
                % Labels
                \node[below left] at (A) {$Q$};
                \node[ right] at (B) {$K$};
                \node[right] at (C) {$H$};
                \node[below right] at (D) {$F$};
                \filldraw [black] (A) circle (1.5pt);
                \filldraw [black] (B) circle (1.5pt);
                \filldraw [black] (C) circle (1.5pt);
                \filldraw [black] (D) circle (1.5pt);
                % Coordinates of points
                \coordinate (E) at (3,2);
                \coordinate (F) at (2,0);
                \coordinate (M) at (5,0);
                \coordinate (N) at (6,2);
                
                % Draw the triangles
                \draw (E) -- (F) -- (M) --(N)-- cycle; 
                \draw (F)--(N);
                % Labels
                \node[above left] at (E) {$E$};
                \node[ left] at (F) {$I$};
                \node[right] at (M) {$G$};
                \node[above right] at (N){$J$};
                \filldraw [black] (E) circle (1.5pt);
                \filldraw [black] (F) circle (1.5pt);
                \filldraw [black] (M) circle (1.5pt);
                \filldraw [black] (N) circle (1.5pt);
            \draw 
            (A) pic[draw,angle radius = 15] {angle = B--A--C} ;
            \draw 
            (A) pic[draw,angle radius = 13] {angle = D--A--B} ;
                \draw 
            (B) pic[draw,angle radius = 8] {angle = C--B--A} pic[draw,angle radius = 10] {angle = C--B--A} ;
            \draw 
            (B) pic[draw,angle radius = 13] {angle = A--B--D}  pic[draw,angle radius = 11] {angle = A--B--D}    ;
            
                \draw 
            (N) pic[draw,angle radius = 15] {angle = E--N--F} ;
            \draw 
            (N) pic[draw,angle radius = 13] {angle = M--F--N} ;
            \draw 
            (F) pic[draw,angle radius = 8] {angle = N--F--E} pic[draw,angle radius = 10] {angle = N--F--E} ;
            \draw 
            (N) pic[draw,angle radius = 13] {angle = F--N--M}  pic[draw,angle radius = 11] {angle = F--N--M}    ;
            \end{tikzpicture}
}
\end{minipage}

\vspace{3pt}\hrule\vspace{3pt}\footnotesize\textcolor{gray!80!black}{\textbf{Các mã dùng chung hình (HINH\_131):} \texttt{L07\_C4\_B5\_02}}
\end{tcolorbox}
\vspace{4pt}


\begin{tcolorbox}[
    breakable,
    colback=white,
    colframe=blue!60!black,
    boxrule=0.8pt,
    arc=3pt,
    title={\textbf{Câu \#162} \quad \textbf{MÃ CÂU: \texttt{[L07\_C4\_B5\_05]}} \quad | \quad \textbf{Nhóm: \texttt{HINH\_132}} \quad | \quad \textbf{MapID: \texttt{7H4H5}} \hfill \textcolor{yellow!90!white}{[Hình gốc .pyc]}}
]
\textbf{Dạng toán:} Bài tổng hợp CM 2 tam giác bắng nhau 2 dt song song cm trung điểm \hfill \textit{(Tự luận -- Thông hiểu)}
\vspace{4pt}\hrule\vspace{6pt}
\begin{minipage}[t]{0.58\textwidth}
\vspace{0pt}
Cho $\triangle {EFG}$, vẽ điểm ${P}$ là trung điểm của ${FG}$. Trên tia đối của tia ${PE}$ lấy điểm ${L}$ sao cho ${PE = PL}$.

 a) CHứng minh $\triangle {EFP} = \triangle {LGP}$ và ${EF \parallel LG}$.

b) Kẻ ${FM \perp EP}$ (${M \in EP}$), ${GO \perp LP}$ ($O \in LP$). Chứng minh rằng ${P}$ là trung điểm của ${MO}$ \\
\end{minipage}\hfill
\begin{minipage}[t]{0.39\textwidth}
\vspace{0pt}
\centering

\resizebox{\linewidth}{!}{%
\begin{tikzpicture}[scale=1, font=\footnotesize, line join=round, line cap=round, >=stealth]
    \def\r{2.5}

    \def\gocA{130}
    \path (0,0) coordinate (O)
    (\gocA:\r) coordinate (A)
    (210:\r) coordinate (B)
    (-30:\r) coordinate (C)
    ($(C)!0.5!(B)$) coordinate (M)
    ($(A)!2!(M)$) coordinate (D)
    ($(A)!(B)!(M)$) coordinate (E) 
    ($(A)!(C)!(M)$) coordinate (F) 
    ;
    \draw (A)--(B)--(C)--cycle (C)--(D)--(A)  (B)--(E)  (C)--(F)
    ;
    \foreach \x/\y in {M/B,M/C}{
        \draw ($(\x)!0.5!(\y)!4pt!90:(\x)$)--($(\x)!0.5!(\y)!4pt!-90:(\x)$);}
    \foreach \x/\y/\z in{B/E/A,C/F/D}
    \path pic[draw=black,angle radius = 6] {right angle = \x--\y--\z} ; 
    \filldraw[black] (A) circle (2pt) node[above left ] {$E$};
    \filldraw[black] (B) circle (2pt) node[below left] {$F$};
    \filldraw[black] (C) circle (2pt) node[above ] {$G$};
   \filldraw[black] (D) circle (2pt) node[below]{$L$};
     \filldraw[black] (E) circle (2pt) node[above right ] {$M$};
    \filldraw[black] (F) circle (2pt) node[below left] {$O$};

    \filldraw[black] (M) circle (2pt) node[above right] {$P$};

    \end{tikzpicture}
}
\end{minipage}

\vspace{3pt}\hrule\vspace{3pt}\footnotesize\textcolor{gray!80!black}{\textbf{Các mã dùng chung hình (HINH\_132):} \texttt{L07\_C4\_B8\_01}}
\end{tcolorbox}
\vspace{4pt}


\begin{tcolorbox}[
    breakable,
    colback=white,
    colframe=blue!60!black,
    boxrule=0.8pt,
    arc=3pt,
    title={\textbf{Câu \#163} \quad \textbf{MÃ CÂU: \texttt{[L07\_C4\_B8\_01]}} \quad | \quad \textbf{Nhóm: \texttt{HINH\_132}} \quad | \quad \textbf{MapID: \texttt{7H4H8}} \hfill \textcolor{yellow!90!white}{[Hình gốc .pyc]}}
]
\textbf{Dạng toán:} TF-M2. CM 2 tam giác bắng nhau, 2 dt song song, cm trung điểm \hfill \textit{(Đúng-Sai -- Thông hiểu)}
\vspace{4pt}\hrule\vspace{6pt}
\begin{minipage}[t]{0.58\textwidth}
\vspace{0pt}
Cho $\triangle {BDH}$, vẽ điểm ${N}$ là trung điểm của ${DH}$. Trên tia đối của tia ${NB}$ lấy điểm ${I}$ sao cho ${NB = NI}$. Xét tính đúng-sai của các khẳng định sau.

\vspace{4pt}\noindent
\textbf{\textcolor{blue!75!black}{a) (Đ)}} $\triangle {BDN} = \triangle {IHN}$\\[2pt]
\textbf{b)} $\widehat{B D N}=\widehat{H I N}$\\[2pt]
\textbf{\textcolor{blue!75!black}{c) (Đ)}} ${BD \parallel IH}$\\[2pt]
\textbf{\textcolor{blue!75!black}{d) (Đ)}} Kẻ ${DK \perp BN}$ (${K \in BN}$), ${HL \perp IN}$ ($L \in IN$). Khi đó ${N}$ là trung điểm của ${KL}$
\end{minipage}\hfill
\begin{minipage}[t]{0.39\textwidth}
\vspace{0pt}
\centering

\resizebox{\linewidth}{!}{%
\begin{tikzpicture}[scale=1, font=\footnotesize, line join=round, line cap=round, >=stealth]
    \def\r{2.5}

    \def\gocA{130}
    \path (0,0) coordinate (O)
    (\gocA:\r) coordinate (A)
    (210:\r) coordinate (B)
    (-30:\r) coordinate (C)
    ($(C)!0.5!(B)$) coordinate (M)
    ($(A)!2!(M)$) coordinate (D)
    ($(A)!(B)!(M)$) coordinate (E) 
    ($(A)!(C)!(M)$) coordinate (F) 
    ;
    \draw (A)--(B)--(C)--cycle (C)--(D)--(A)  (B)--(E)  (C)--(F)
    ;
    \foreach \x/\y in {M/B,M/C}{
        \draw ($(\x)!0.5!(\y)!4pt!90:(\x)$)--($(\x)!0.5!(\y)!4pt!-90:(\x)$);}
    \foreach \x/\y/\z in{B/E/A,C/F/D}
    \path pic[draw=black,angle radius = 6] {right angle = \x--\y--\z} ; 
    \filldraw[black] (A) circle (2pt) node[above left ] {$B$};
    \filldraw[black] (B) circle (2pt) node[below left] {$D$};
    \filldraw[black] (C) circle (2pt) node[above ] {$H$};
   \filldraw[black] (D) circle (2pt) node[below]{$I$};
     \filldraw[black] (E) circle (2pt) node[above right ] {$K$};
    \filldraw[black] (F) circle (2pt) node[below left] {$L$};

    \filldraw[black] (M) circle (2pt) node[above right] {$N$};

    \end{tikzpicture}
}
\end{minipage}

\vspace{3pt}\hrule\vspace{3pt}\footnotesize\textcolor{gray!80!black}{\textbf{Các mã dùng chung hình (HINH\_132):} \texttt{L07\_C4\_B5\_05}}
\end{tcolorbox}
\vspace{4pt}


\begin{tcolorbox}[
    breakable,
    colback=white,
    colframe=blue!60!black,
    boxrule=0.8pt,
    arc=3pt,
    title={\textbf{Câu \#164} \quad \textbf{MÃ CÂU: \texttt{[L07\_C4\_B8\_02]}} \quad | \quad \textbf{Nhóm: \texttt{HINH\_133}} \quad | \quad \textbf{MapID: \texttt{7H4H8}} \hfill \textcolor{yellow!90!white}{[Hình gốc .pyc]}}
]
\textbf{Dạng toán:} TF-M2. CM 2 tam giác bắng nhau 2 dt song song cm trung điểm \hfill \textit{(Đúng-Sai -- Thông hiểu)}
\vspace{4pt}\hrule\vspace{6pt}
\begin{minipage}[t]{0.58\textwidth}
\vspace{0pt}
Cho $\triangle {BCG}$ có ${BC=BG}$. Gọi ${Q}$ là trung điểm của ${CG}$. Xét tính đúng-sai của các khẳng định sau.

\vspace{4pt}\noindent
\textbf{a)} $\triangle {BQC}=\triangle {BGQ}$\\[2pt]
\textbf{b)} ${BQ\perp CG}$\\[2pt]
\textbf{\textcolor{blue!75!black}{c) (Đ)}} Trên tia đối của tia ${QB}$ lấy điểm ${I}$ sao cho ${QB=QI}$, khi đó ${BC=IG}$\\[2pt]
\textbf{d)} Kẻ ${QJ\perp BC}(J\in BC)$, ${QP\perp IG}$ ($P\in IG$). Khi đó $\widehat{JQP} < 180^{\circ}$
\end{minipage}\hfill
\begin{minipage}[t]{0.39\textwidth}
\vspace{0pt}
\centering

\resizebox{\linewidth}{!}{%
\begin{tikzpicture}[scale=1, font=\footnotesize, line join=round, line cap=round, >=stealth]  
    \def\a{5}
    \def\h{3}
    \path   (0:0) coordinate (B)
            ++(0:\a) coordinate (C)
            ($(B)!0.5!(C)$) coordinate (I)  
            ($(I)+(90:\h)$) coordinate (A)
            ($(A)!2!(I)$) coordinate (K)
            ($(A)!(I)!(B)$) coordinate (E)
            ($(K)!(I)!(C)$) coordinate (F)
            ;
    \draw 
    (A)--(I)--(K)--(C) (F)--(I)--(E)
    (A)--(B)--(C)--cycle;
    \draw 
    pic[draw,angle radius=3mm]{right angle=A--I--C}
    pic[draw,angle radius=3mm]{right angle=A--E--I}
    pic[draw,angle radius=3mm]{right angle=I--F--C}
    ;
   \filldraw[black] (A) circle (2pt) node[above left ] {$B$};
    \filldraw[black] (B) circle (2pt) node[below left] {$C$};
    \filldraw[black] (C) circle (2pt) node[above ] {$G$};
   \filldraw[black] (K) circle (2pt) node[below]{$I$};
     \filldraw[black] (E) circle (2pt) node[above left ] {$J$};
    \filldraw[black] (F) circle (2pt) node[below right] {$P$};
    \filldraw[black] (I) circle (2pt) node[below left] {$Q$};

        \end{tikzpicture}
}
\end{minipage}

\end{tcolorbox}
\vspace{4pt}


\section{Toán 8}

\subsection{Chương 1 - Đa thức}

\begin{tcolorbox}[
    breakable,
    colback=white,
    colframe=blue!60!black,
    boxrule=0.8pt,
    arc=3pt,
    title={\textbf{Câu \#165} \quad \textbf{MÃ CÂU: \texttt{[L08\_C1\_B5\_10]}} \quad | \quad \textbf{Nhóm: \texttt{HINH\_134}} \quad | \quad \textbf{MapID: \texttt{8D1V5}} \hfill \textcolor{yellow!90!white}{[Hình gốc .pyc]}}
]
\textbf{Dạng toán:} Cho hình vẽ mô phỏng viết biểu thức tính diện tích \hfill \textit{(Tự luận -- Vận dụng)}
\vspace{4pt}\hrule\vspace{6pt}
\begin{minipage}[t]{0.58\textwidth}
\vspace{0pt}
Tính diện tích phần tô màu trong hình dưới.
\end{minipage}\hfill
\begin{minipage}[t]{0.39\textwidth}
\vspace{0pt}
\centering

\resizebox{\linewidth}{!}{%
\begin{tikzpicture}
    % Tô màu phần không thuộc hình chữ nhật nhỏ
    \fill[blue!20] (0,0) rectangle (10,5); % Hình chữ nhật lớn tăng chiều dài
    \fill[white] (3,1) rectangle (7,3.5); % Hình chữ nhật nhỏ tăng chiều dài
    
    % Viền của hình chữ nhật lớn
    \draw[thick] (0,0) rectangle (10,5); % Điều chỉnh tọa độ cho hình lớn
    \node at (10.9,2.5) {$10 x + 4 y$};
    \node at (5,-0.5) {$5 x + 3 y$};
    
    % Viền của hình chữ nhật nhỏ
    \draw[thick] (3,1) rectangle (7,3.5); % Điều chỉnh tọa độ cho hình nhỏ
    \node at (7.9,2.5) {$4 x + y$};
    \node at (5,0.5)  {$3 x - 4 y$};
\end{tikzpicture}
}
\end{minipage}

\end{tcolorbox}
\vspace{4pt}


\subsection{Chương 3 -  Tứ giác}

\begin{tcolorbox}[
    breakable,
    colback=white,
    colframe=blue!60!black,
    boxrule=0.8pt,
    arc=3pt,
    title={\textbf{Câu \#166} \quad \textbf{MÃ CÂU: \texttt{[L08\_C3\_B1\_01]}} \quad | \quad \textbf{Nhóm: \texttt{HINH\_135}} \quad | \quad \textbf{MapID: \texttt{8H3N1}} \hfill \textcolor{yellow!90!white}{[Hình gốc .pyc]}}
]
\textbf{Dạng toán:} Tính số đo góc trong tứ giác \hfill \textit{(Trắc nghiệm -- Nhận biết)}
\vspace{4pt}\hrule\vspace{6pt}
\begin{minipage}[t]{0.58\textwidth}
\vspace{0pt}
Tìm số đo ${x}$ trong hình sau.

\vspace{4pt}\noindent
\textbf{A.} $215^{\circ}$ \quad \textbf{B.} $173^{\circ}$ \quad \textbf{C.} $172^{\circ}$ \quad \textbf{\textcolor{blue!75!black}{D. (Đ)}} $100^{\circ}$
\end{minipage}\hfill
\begin{minipage}[t]{0.39\textwidth}
\vspace{0pt}
\centering

\resizebox{\linewidth}{!}{%
\begin{tikzpicture}
        % Draw the quadrilateral
        \coordinate (A) at (0, 0);
        \coordinate (B) at (1, 3);
        \coordinate (C) at (5, 4);
        \coordinate (D) at (6, 0);
        
        \draw (A) -- (B) -- (C) -- (D) -- cycle;
        
        % Label the points
        \node[below left] at (A) {$A$};
        \node[above left] at (B) {$B$};
        \node[above] at (C) {$C$};
        \node[below right] at (D) {$D$};
        
        % Draw and label the angles
       \draw 
    (B) pic[draw,angle radius = 12] {angle =A--B--C} node[shift={(-35:25pt)}]{$ 115^\circ $} 
    (A) pic[draw,angle radius = 12] {angle =D--A--B} node[shift={(25:20pt)}]{$ 72^\circ $}  
    (D) pic[draw,angle radius = 12] {angle =C--D--A} node[shift={(-25:-20pt)}]{$ 73^\circ $} 
    (C) pic[draw,angle radius = 12] {angle =B--C--D} node[shift={(55:-20pt)}]{$ x^\circ $}   
    ;
        
    \end{tikzpicture}
}
\end{minipage}

\end{tcolorbox}
\vspace{4pt}


\begin{tcolorbox}[
    breakable,
    colback=white,
    colframe=blue!60!black,
    boxrule=0.8pt,
    arc=3pt,
    title={\textbf{Câu \#167} \quad \textbf{MÃ CÂU: \texttt{[L08\_C3\_B1\_02]}} \quad | \quad \textbf{Nhóm: \texttt{HINH\_136}} \quad | \quad \textbf{MapID: \texttt{8H3N1}} \hfill \textcolor{yellow!90!white}{[Hình gốc .pyc]}}
]
\textbf{Dạng toán:} Tính số đo góc trong tứ giác có 1 góc vuông \hfill \textit{(Trắc nghiệm -- Nhận biết)}
\vspace{4pt}\hrule\vspace{6pt}
\begin{minipage}[t]{0.58\textwidth}
\vspace{0pt}
Tìm số đo ${x}$ trong hình sau.

\vspace{4pt}\noindent
\textbf{A.} $117^{\circ}$ \quad \textbf{\textcolor{blue!75!black}{B. (Đ)}} $127^{\circ}$ \quad \textbf{C.} $137^{\circ}$ \quad \textbf{D.} $132^{\circ}$
\end{minipage}\hfill
\begin{minipage}[t]{0.39\textwidth}
\vspace{0pt}
\centering

\resizebox{\linewidth}{!}{%
\begin{tikzpicture}
        % Vẽ các đỉnh của hình tứ giác
        \coordinate (E) at (0, 0);
        \coordinate (F) at (0, 3);
        \coordinate (G) at (2, 4);
        \coordinate (H) at (5, 0);
        
        % Vẽ các cạnh
        \draw (E) -- (F) -- (G) -- (H) -- cycle;
        
        % Đánh dấu góc vuông tại điểm E
        \draw (E) rectangle +(0.3, 0.3);
        
        % Ghi các giá trị góc
        \node at (1.9, 3.5) {$82^\circ$};
        \node at (4.4, 0.3) {$61^\circ$};
        \node at (0.3, 2.8) {$x$};
        
        % Ghi các điểm E, F, G, H
        \node[below left] at (E) {E};
        \node[above left] at (F) {F};
        \node[above right] at (G) {G};
        \node[below right] at (H) {H};
        
    \end{tikzpicture}
}
\end{minipage}

\end{tcolorbox}
\vspace{4pt}


\begin{tcolorbox}[
    breakable,
    colback=white,
    colframe=blue!60!black,
    boxrule=0.8pt,
    arc=3pt,
    title={\textbf{Câu \#168} \quad \textbf{MÃ CÂU: \texttt{[L08\_C3\_B1\_03]}} \quad | \quad \textbf{Nhóm: \texttt{HINH\_137}} \quad | \quad \textbf{MapID: \texttt{8H3N1}} \hfill \textcolor{yellow!90!white}{[Hình gốc .pyc]}}
]
\textbf{Dạng toán:} Tính số đo góc trong hình thang cân \hfill \textit{(Trắc nghiệm -- Nhận biết)}
\vspace{4pt}\hrule\vspace{6pt}
\begin{minipage}[t]{0.58\textwidth}
\vspace{0pt}
Tìm số đo ${x}$ trong hình sau.

\vspace{4pt}\noindent
\textbf{A.} $x={72}^{\circ}$ \quad \textbf{B.} $x={67}^{\circ}$ \quad \textbf{C.} $x={61}^{\circ}$ \quad \textbf{\textcolor{blue!75!black}{D. (Đ)}} $x={62}^{\circ}$
\end{minipage}\hfill
\begin{minipage}[t]{0.39\textwidth}
\vspace{0pt}
\centering

\resizebox{\linewidth}{!}{%
\begin{tikzpicture}
        % Các điểm
        \coordinate (N) at (0, 4);
        \coordinate (P) at (6, 4);
        \coordinate (M) at (1, 0);
        \coordinate (Q) at (5, 0);
        
        % Các cạnh của hình thang
        \draw (N) -- (P) -- (Q) -- (M) -- cycle;
        
        % Ghi giá trị góc
        \node at (1.5, 0.4) {$118^\circ$};
        \node at (4.6, 0.4) {$118^\circ$};
        \node at (0.4, 3.7) {$x$};
        \node at (5.5, 3.7) {$x$};
        
        % Ghi tên các đỉnh
        \node[above left] at (N) {$N$};
        \node[above right] at (P) {$P$};
        \node[below left] at (M) {$M$};
        \node[below right] at (Q) {$Q$};
    \end{tikzpicture}
}
\end{minipage}

\end{tcolorbox}
\vspace{4pt}


\begin{tcolorbox}[
    breakable,
    colback=white,
    colframe=blue!60!black,
    boxrule=0.8pt,
    arc=3pt,
    title={\textbf{Câu \#169} \quad \textbf{MÃ CÂU: \texttt{[L08\_C3\_B1\_04]}} \quad | \quad \textbf{Nhóm: \texttt{HINH\_138}} \quad | \quad \textbf{MapID: \texttt{8H3H1}} \hfill \textcolor{yellow!90!white}{[Hình gốc .pyc]}}
]
\textbf{Dạng toán:} Tính số đo góc trong tứ giác có 2 góc vuông đối diện \hfill \textit{(Trắc nghiệm -- Thông hiểu)}
\vspace{4pt}\hrule\vspace{6pt}
\begin{minipage}[t]{0.58\textwidth}
\vspace{0pt}
Tìm số đo ${x}$ trong hình sau.

\vspace{4pt}\noindent
\textbf{A.} $x={98}^{\circ}$ \quad \textbf{\textcolor{blue!75!black}{B. (Đ)}} $x={97}^{\circ}$ \quad \textbf{C.} $x={102}^{\circ}$ \quad \textbf{D.} $x={99}^{\circ}$
\end{minipage}\hfill
\begin{minipage}[t]{0.39\textwidth}
\vspace{0pt}
\centering

\resizebox{\linewidth}{!}{%
\begin{tikzpicture}
    
    % Draw the points and the lines
    \coordinate (A) at (-2, 2);
    \coordinate (B) at (2.07, 3.36);
    \coordinate (C) at (2, 0);
    \coordinate (D) at (-1.36, 0.07);
    
    % Draw the quadrilateral
    \draw (A) -- (B) -- (C) -- (D) -- cycle;
    
    
    % Label the points
    \node at (A) [left] {$A$};
    \node at (B) [above] {$B$};
    \node at (C) [right] {$C$};
    \node at (D) [below] {$D$};
    
    % Label the angles and side x
    \node at (1.6,2.8) {$83^\circ$};
    \node at (-1,0.3) {$x^\circ$};
    
\def\khvuong[size=#1](#2,#3,#4){%
    \draw ($(#3)!#1!(#2)$) --
    ($($(#3)!#1!(#2)$)+($(#3)!#1!(#4)$)-(#3)$)--
    ($(#3)!#1!(#4)$);
}
\khvuong[size=5pt](B,A,D);%
\khvuong[size=5pt](B,C,D);%

    \end{tikzpicture}
}
\end{minipage}

\end{tcolorbox}
\vspace{4pt}


\begin{tcolorbox}[
    breakable,
    colback=white,
    colframe=blue!60!black,
    boxrule=0.8pt,
    arc=3pt,
    title={\textbf{Câu \#170} \quad \textbf{MÃ CÂU: \texttt{[L08\_C3\_B1\_05]}} \quad | \quad \textbf{Nhóm: \texttt{HINH\_139}} \quad | \quad \textbf{MapID: \texttt{8H3H1}} \hfill \textcolor{yellow!90!white}{[Hình gốc .pyc]}}
]
\textbf{Dạng toán:} Tính số đo góc kề bù trong tứ giác có 2 góc vuông đối diện \hfill \textit{(Trắc nghiệm -- Thông hiểu)}
\vspace{4pt}\hrule\vspace{6pt}
\begin{minipage}[t]{0.58\textwidth}
\vspace{0pt}
Tìm số đo $\widehat{B}_{1}$ trong hình sau.

\vspace{4pt}\noindent
\textbf{A.} $x={72}^{\circ}$ \quad \textbf{B.} $x={73}^{\circ}$ \quad \textbf{\textcolor{blue!75!black}{C. (Đ)}} $x={108}^{\circ}$ \quad \textbf{D.} $x={109}^{\circ}$
\end{minipage}\hfill
\begin{minipage}[t]{0.39\textwidth}
\vspace{0pt}
\centering

\resizebox{\linewidth}{!}{%
\begin{tikzpicture}
            
            % Draw the points and the lines
            \coordinate (A) at (-2, 2);
            \coordinate (B) at (2.07, 3.36);
            \coordinate (C) at (2, 0);
            \coordinate (D) at (-1.36, 0.07);
            \coordinate (E) at (2.1,5);
            % Draw the quadrilateral
            \draw (A) -- (B) -- (C) -- (D) -- cycle;
                \draw (C) -- (B) -- (E);
            
            % Label the points
            \node at (A) [left] {$A$};
            \node at (B) [above right] {$B$};
            \node at (C) [right] {$C$};
            \node at (D) [below] {$D$};
        
            % Label the angles and side x
            \node at (1.9,3.5) {$1$};
            \node at (-1,0.3) {$108^\circ$};
            
            \def\khvuong[size=#1](#2,#3,#4){%
                \draw ($(#3)!#1!(#2)$) --
                ($($(#3)!#1!(#2)$)+($(#3)!#1!(#4)$)-(#3)$)--
                ($(#3)!#1!(#4)$);
            }
            \khvuong[size=5pt](B,A,D);%
            \khvuong[size=5pt](B,C,D);%

    \end{tikzpicture}
}
\end{minipage}

\end{tcolorbox}
\vspace{4pt}


\begin{tcolorbox}[
    breakable,
    colback=white,
    colframe=blue!60!black,
    boxrule=0.8pt,
    arc=3pt,
    title={\textbf{Câu \#171} \quad \textbf{MÃ CÂU: \texttt{[L08\_C3\_B1\_07]}} \quad | \quad \textbf{Nhóm: \texttt{HINH\_140}} \quad | \quad \textbf{MapID: \texttt{8H3V1}} \hfill \textcolor{yellow!90!white}{[Hình gốc .pyc]}}
]
\textbf{Dạng toán:} Tính số đo góc kiểu 2 \hfill \textit{(Trắc nghiệm -- Vận dụng)}
\vspace{4pt}\hrule\vspace{6pt}
\begin{minipage}[t]{0.58\textwidth}
\vspace{0pt}
Tính số đo của $\widehat{A}_{1}$

\vspace{4pt}\noindent
\textbf{A.} $101^{\circ} $ \quad \textbf{B.} $102^{\circ} $ \quad \textbf{C.} $99^{\circ} $ \quad \textbf{\textcolor{blue!75!black}{D. (Đ)}} $100^{\circ} $
\end{minipage}\hfill
\begin{minipage}[t]{0.39\textwidth}
\vspace{0pt}
\centering

\resizebox{\linewidth}{!}{%
\begin{tikzpicture}
    
    % Define the points A, B, C, and D
    \coordinate (A) at (0.4,4);
    \coordinate (B) at (5,2.5);
    \coordinate (C) at (5,0);
    \coordinate (D) at (0,0);
        \coordinate (E) at (0.5,5);
    % Draw the lines
    \draw (A) -- (B) -- (C) -- (D) -- cycle;
    \draw (D) -- (A)--(E);
    \draw (B) -- (C)--(5,4);
    % Add the labels
    \node[above left] at (A) {A};
    \node[below right] at (B) {B};
    \node[below right] at (C) {C};
    \node[below left] at (D) {D};
    
    % Add angle labels
   % Ghi giá trị góc
    \node at (0.4, 0.3) {$83^\circ$};
    \node at (4.7, 3.0) {$73^\circ$};
    \node at (0.6, 4.2) {$1$};
     % Đánh dấu góc vuông tại điểm E
    \draw (C) rectangle +(-0.3, 0.3);
    \end{tikzpicture}
}
\end{minipage}

\end{tcolorbox}
\vspace{4pt}


\begin{tcolorbox}[
    breakable,
    colback=white,
    colframe=blue!60!black,
    boxrule=0.8pt,
    arc=3pt,
    title={\textbf{Câu \#172} \quad \textbf{MÃ CÂU: \texttt{[L08\_C3\_B1\_08]}} \quad | \quad \textbf{Nhóm: \texttt{HINH\_141}} \quad | \quad \textbf{MapID: \texttt{8H3V1}} \hfill \textcolor{yellow!90!white}{[Hình gốc .pyc]}}
]
\textbf{Dạng toán:} Tính số đo góc kiểu 2 \hfill \textit{(Tự luận -- Vận dụng)}
\vspace{4pt}\hrule\vspace{6pt}
\begin{minipage}[t]{0.58\textwidth}
\vspace{0pt}
Tìm số đo các góc trong hình sau:
\end{minipage}\hfill
\begin{minipage}[t]{0.39\textwidth}
\vspace{0pt}
\centering

\resizebox{\linewidth}{!}{%
\begin{tikzpicture}
    
    % Define points
    \coordinate (K) at (0, 0);
    \coordinate (I) at (6, 0);
    \coordinate (H) at (5, 2);
    \coordinate (G) at (3, 3);
        \coordinate (A) at (-3, 0);
    % Draw lines
    \draw[thick] (K) -- (I) -- (H) -- (G) -- cycle;
    \draw[thick] (H) -- (I);
        \draw[thick] (A) -- (I);
        \draw[thick] (G) -- (H)--(6, 1.5);
    % Labels for points
    \node[below left] at (K) {K};
    \node[below right] at (I) {I};
    \node[above right] at (H) {H};
    \node[above] at (G) {G};
    
    % Labels for angles
    \node at (-0.3, 0.3) {$142^\circ$};
    \node at (3.1, 2.6) {$104^\circ$};
        \node at (5.6, 1.5) {$11^\circ$};
    \node at (4.8, 1.8) {?};
        \node at (0.5, 0.2) {?};
            \node at (5.7, 0.2) {?};
    \end{tikzpicture}
}
\end{minipage}

\end{tcolorbox}
\vspace{4pt}


\begin{tcolorbox}[
    breakable,
    colback=white,
    colframe=blue!60!black,
    boxrule=0.8pt,
    arc=3pt,
    title={\textbf{Câu \#173} \quad \textbf{MÃ CÂU: \texttt{[L08\_C3\_B2\_10]}} \quad | \quad \textbf{Nhóm: \texttt{HINH\_142}} \quad | \quad \textbf{MapID: \texttt{8H3N2}} \hfill \textcolor{yellow!90!white}{[Hình gốc .pyc]}}
]
\textbf{Dạng toán:} Tính số đo góc trong hình thang. \hfill \textit{(Trắc nghiệm -- Nhận biết)}
\vspace{4pt}\hrule\vspace{6pt}
\begin{minipage}[t]{0.58\textwidth}
\vspace{0pt}
Cho hình vẽ sau. Chọn khẳng định sai trong các khẳng định sau?

\vspace{4pt}\noindent
\textbf{A.} $\widehat{A B C } = 118^{\circ}$ \quad \textbf{B.} $x= 62^{\circ}$ thì tứ giác ${A B C D }$ hình thang cân. \quad \textbf{C.} Tứ giác ${A B C D }$ là hình thang. \quad \textbf{\textcolor{blue!75!black}{D. (Đ)}} $\widehat{B A D } = 62^{\circ}$
\end{minipage}\hfill
\begin{minipage}[t]{0.39\textwidth}
\vspace{0pt}
\centering

\resizebox{\linewidth}{!}{%
\begin{tikzpicture}
    \coordinate (D) at (0, 0);
            \coordinate (C) at (4, 0);
            \coordinate (B) at (3, 2);
            \coordinate (A) at (1, 2);
            \coordinate (E) at (6, 2);
            % Vẽ các cạnh của hình thang
            \draw[thick] (A) -- (B) -- (C) -- (D) -- cycle;
        
            % Gán nhãn cho các điểm
            \node[above left] at (A) {$A $};
            \node[above right] at (B) {$B $};
            \node[below right] at (C) {$C $};
            \node[below left] at (D) {$D $};
            % Gán biến x tại góc D
            \node[left] at (0.5,0.2) {$x$};
            
            % Vẽ đoạn thẳng nối AB
            \draw[thick] (A) -- (5,2);
            \draw 
            (B) pic[draw,angle radius = 6] {angle = C--B--E} node[shift={(-25:16pt)}]{$ 62^\circ $} 
            (C) pic[draw,angle radius = 6] {angle = B--C--D} node[shift={(-32:-15pt)}]{$ 62^\circ $}   
            ;

    \end{tikzpicture}
}
\end{minipage}

\end{tcolorbox}
\vspace{4pt}


\begin{tcolorbox}[
    breakable,
    colback=white,
    colframe=blue!60!black,
    boxrule=0.8pt,
    arc=3pt,
    title={\textbf{Câu \#174} \quad \textbf{MÃ CÂU: \texttt{[L08\_C3\_B2\_11]}} \quad | \quad \textbf{Nhóm: \texttt{HINH\_143}} \quad | \quad \textbf{MapID: \texttt{8H3N2}} \hfill \textcolor{yellow!90!white}{[Hình gốc .pyc]}}
]
\textbf{Dạng toán:} Chứng minh hình thang cân và Tính số đo góc trong hình thang. \hfill \textit{(Trắc nghiệm -- Nhận biết)}
\vspace{4pt}\hrule\vspace{6pt}
\begin{minipage}[t]{0.58\textwidth}
\vspace{0pt}
Cho hình thang ${IJQK}$. Chọn khẳng định sai trong các khẳng định sau?

\vspace{4pt}\noindent
\textbf{\textcolor{blue!75!black}{A. (Đ)}} $\widehat{I} > \widehat{J} $ \quad \textbf{B.} Tứ giác ${IJQK}$ là hình thang cân. \quad \textbf{C.} $\widehat{J} =53^{\circ} $ \quad \textbf{D.} $\widehat{I} =53^{\circ} $
\end{minipage}\hfill
\begin{minipage}[t]{0.39\textwidth}
\vspace{0pt}
\centering

\resizebox{\linewidth}{!}{%
\begin{tikzpicture}
        % Define points
    \coordinate (D) at (0, 0);
    \coordinate (C) at (4, 0);
    \coordinate (B) at (5, 2);
    \coordinate (A) at (-1, 2);
        \coordinate (E) at (6, 2);
    % Vẽ các cạnh của hình thang
    \draw[thick] (0,0) -- (4,0) -- (5,2) -- (-1,2) -- cycle;
    % Gán nhãn cho các điểm
    \node[above left] at (-1,2) {$I$};
    \node[above right] at (5,2) {$J$};
    \node[below right] at (4,0) {$Q$};
    \node[below left] at (0,0) {$K$};
    \draw 
            (0,0) pic[draw,angle radius = 6] {angle = C--D--A} node[shift={(45:16pt)}]{$ 127^\circ $} 
            (4,0) pic[draw,angle radius = 6] {angle = B--C--D} node[shift={(-45:-15pt)}]{$ 127^\circ $}   
            ;
    \end{tikzpicture}
}
\end{minipage}

\end{tcolorbox}
\vspace{4pt}


\begin{tcolorbox}[
    breakable,
    colback=white,
    colframe=blue!60!black,
    boxrule=0.8pt,
    arc=3pt,
    title={\textbf{Câu \#175} \quad \textbf{MÃ CÂU: \texttt{[L08\_C3\_B2\_12]}} \quad | \quad \textbf{Nhóm: \texttt{HINH\_144}} \quad | \quad \textbf{MapID: \texttt{8H3H2}} \hfill \textcolor{yellow!90!white}{[Hình gốc .pyc]}}
]
\textbf{Dạng toán:} Chứng minh hình thang cân và góc bằng nhau \hfill \textit{(Trắc nghiệm -- Thông hiểu)}
\vspace{4pt}\hrule\vspace{6pt}
\begin{minipage}[t]{0.58\textwidth}
\vspace{0pt}
Cho hình thang ${IJQM}$ biết  $IQ=JM$. Trong các khẳng định sau khẳng định nào sai?

\vspace{4pt}\noindent
\textbf{A.} $\widehat{IJM}= \widehat{IQM}$ \quad \textbf{B.} ${IM} = {JQ}$ \quad \textbf{\textcolor{blue!75!black}{C. (Đ)}} $\widehat{IJQ}= \widehat{IMQ}$ \quad \textbf{D.} $\widehat{MQI}= \widehat{JMQ}$
\end{minipage}\hfill
\begin{minipage}[t]{0.39\textwidth}
\vspace{0pt}
\centering

\resizebox{\linewidth}{!}{%
\begin{tikzpicture}
            % Define points
            \coordinate (D) at (-2, 0);
            \coordinate (C) at (6, 0);
            \coordinate (B) at (4, 3);
            \coordinate (A) at (0, 3);
        
            % Vẽ các cạnh của hình thang
            \draw (-2,0) -- (6,0) -- (4,3) -- (0,3) -- cycle;
                \draw (-2,0) -- (4,3);
                \draw (6,0) -- (0,3);
    % Gán nhãn cho các điểm
    \node[above left] at (0,3) {$I$};
            \node[above right] at (4,3) {$J$};
            \node[below right] at (6,0) {$Q$};
            \node[below left] at (-2,0) {$M$};
   
    \end{tikzpicture}
}
\end{minipage}

\vspace{3pt}\hrule\vspace{3pt}\footnotesize\textcolor{gray!80!black}{\textbf{Các mã dùng chung hình (HINH\_144):} \texttt{L08\_C3\_B2\_14, L08\_C3\_B2\_03, L08\_C3\_B2\_05}}
\end{tcolorbox}
\vspace{4pt}


\begin{tcolorbox}[
    breakable,
    colback=white,
    colframe=blue!60!black,
    boxrule=0.8pt,
    arc=3pt,
    title={\textbf{Câu \#176} \quad \textbf{MÃ CÂU: \texttt{[L08\_C3\_B2\_13]}} \quad | \quad \textbf{Nhóm: \texttt{HINH\_145}} \quad | \quad \textbf{MapID: \texttt{8H3V2}} \hfill \textcolor{yellow!90!white}{[Hình gốc .pyc]}}
]
\textbf{Dạng toán:} Bài tổng hợp \hfill \textit{(Trắc nghiệm -- Vận dụng)}
\vspace{4pt}\hrule\vspace{6pt}
\begin{minipage}[t]{0.58\textwidth}
\vspace{0pt}
Cho tam giác ${A B C }$, hai đường phân giác $\widehat{B }$, $\widehat{C }$ cắt nhau tại ${O }$. Qua ${O }$ kẻ đường thẳng song song với $B C $  đường thẳng này cắt ${A B }$, ${A C }$ lần lượt tại ${M }, {N }$. Trong các khẳng định sau khẳng định nào sai?

\vspace{4pt}\noindent
\textbf{A.} Tứ giác ${B C O M }$ là hình thang. \quad \textbf{B.} Tứ giác ${B C N O }$ là hình thang \quad \textbf{C.} $M N =M B +N C $ \quad \textbf{\textcolor{blue!75!black}{D. (Đ)}} Tứ giác ${M N C B }$ là hình thang cân
\end{minipage}\hfill
\begin{minipage}[t]{0.39\textwidth}
\vspace{0pt}
\centering

\resizebox{\linewidth}{!}{%
\begin{tikzpicture}
    % Vẽ tam giác ACB (hoán đổi vị trí của A và C)
    \coordinate (A) at (-2, 4);
    \coordinate (B) at (-3, 0);
    \coordinate (C) at (2, 0);
        \coordinate (M) at (-2.66, 1.35);
    \coordinate (N) at (0.65, 1.35);
        \coordinate (O) at (-1.27, 1.35);
    % Vẽ các cạnh tam giác
    \draw [thick] (A) -- (B) -- (C) -- cycle;
        \draw [thick] (M) -- (O) -- (N);
            \draw [thick] (B) -- (O);
                \draw [thick] (C) -- (O);
    % Gán nhãn cho các điểm
    \node[above] at (A) {$A $};
    \node[below right] at (B) {$B $};
    \node[below left] at (C) {$C $};
    % Tính toán giao điểm của hai đường phân giác

    % Vẽ điểm O và gán nhãn
    \filldraw[black] (O) circle (1pt); % Vẽ điểm O
    \node[above right] at (O) {O };

    \filldraw[blue] (M) circle (1pt);
    \filldraw[blue] (N) circle (1pt);
    
    % Gán nhãn cho M và N
    \node[above left] at (M) {$M $};
    \node[above right] at (N) {$N $};
    \draw
        pic[draw,blue,angle radius=3mm] {angle = O--C--B} pic[draw,blue,angle radius=3.7mm] {angle = O--C--B};
    \draw
    pic[draw,blue,angle radius=4mm] {angle = N--C--O} pic[draw,blue,angle radius=4.7mm] {angle = N--C--O};
    \draw
     pic[draw,blue,angle radius=3.4mm] {angle = C--B--O};
    \draw
    pic[draw,blue,angle radius=4.4mm] {angle = O--B--M} ;
    \end{tikzpicture}
}
\end{minipage}

\end{tcolorbox}
\vspace{4pt}


\begin{tcolorbox}[
    breakable,
    colback=white,
    colframe=blue!60!black,
    boxrule=0.8pt,
    arc=3pt,
    title={\textbf{Câu \#177} \quad \textbf{MÃ CÂU: \texttt{[L08\_C3\_B2\_14]}} \quad | \quad \textbf{Nhóm: \texttt{HINH\_144}} \quad | \quad \textbf{MapID: \texttt{8H3V2}} \hfill \textcolor{yellow!90!white}{[Hình gốc .pyc]}}
]
\textbf{Dạng toán:} Bài tổng hợp 2 \hfill \textit{(Trắc nghiệm -- Vận dụng)}
\vspace{4pt}\hrule\vspace{6pt}
\begin{minipage}[t]{0.58\textwidth}
\vspace{0pt}
Cho hình thang cân ${A B C D }$ có $A B  // C D $ và $A B  < C D $, biết $A D = B A $. Trong các khẳng định sau khẳng định nào sai?

\vspace{4pt}\noindent
\textbf{\textcolor{blue!75!black}{A. (Đ)}} ${A C }$ là tia phân giác của $\widehat{B A D  }$ \quad \textbf{B.} $A B =C B $ \quad \textbf{C.} ${D B }$ là tia phân giác của $\widehat{A D C  }$ \quad \textbf{D.} ${C A }$ là tia phân giác của $\widehat{B C D  }$
\end{minipage}\hfill
\begin{minipage}[t]{0.39\textwidth}
\vspace{0pt}
\centering

\resizebox{\linewidth}{!}{%
\begin{tikzpicture}
            % Define points
            \coordinate (D) at (-2, 0);
            \coordinate (C) at (6, 0);
            \coordinate (B) at (4, 3);
            \coordinate (A) at (0, 3);
        
            % Vẽ các cạnh của hình thang
            \draw (-2,0) -- (6,0) -- (4,3) -- (0,3) -- cycle;
                \draw (-2,0) -- (4,3);
                \draw (6,0) -- (0,3);
    % Gán nhãn cho các điểm
    \node[above left] at (0,3) {$A $};
            \node[above right] at (4,3) {$B $};
            \node[below right] at (6,0) {$C $};
            \node[below left] at (-2,0) {$D $};
   
    \end{tikzpicture}
}
\end{minipage}

\vspace{3pt}\hrule\vspace{3pt}\footnotesize\textcolor{gray!80!black}{\textbf{Các mã dùng chung hình (HINH\_144):} \texttt{L08\_C3\_B2\_12, L08\_C3\_B2\_03, L08\_C3\_B2\_05}}
\end{tcolorbox}
\vspace{4pt}


\begin{tcolorbox}[
    breakable,
    colback=white,
    colframe=blue!60!black,
    boxrule=0.8pt,
    arc=3pt,
    title={\textbf{Câu \#178} \quad \textbf{MÃ CÂU: \texttt{[L08\_C3\_B2\_15]}} \quad | \quad \textbf{Nhóm: \texttt{HINH\_146}} \quad | \quad \textbf{MapID: \texttt{8H3V2}} \hfill \textcolor{yellow!90!white}{[Hình gốc .pyc]}}
]
\textbf{Dạng toán:} Bài tổng hợp 3 \hfill \textit{(Trắc nghiệm -- Vận dụng)}
\vspace{4pt}\hrule\vspace{6pt}
\begin{minipage}[t]{0.58\textwidth}
\vspace{0pt}
Cho hình thang cân ${EFKH}$ có $EF // KH$. Lấy ${J}, {I}$ lần lượt là trung điểm của ${EF}$, ${KH}$

 Trong các khẳng định sau khẳng định nào sai?

\vspace{4pt}\noindent
\textbf{A.} $\Delta {EIF}$ cân tại ${I}$ \quad \textbf{B.} $EI=FI$ \quad \textbf{\textcolor{blue!75!black}{C. (Đ)}} $\Delta {EIH}$ cân tại ${I}$ \quad \textbf{D.} ${ IJ}$ là đường cao của hình thang
\end{minipage}\hfill
\begin{minipage}[t]{0.39\textwidth}
\vspace{0pt}
\centering

\resizebox{\linewidth}{!}{%
\begin{tikzpicture}
            % Define points
            \coordinate (D) at (-2, 0);
            \coordinate (C) at (6, 0);
            \coordinate (B) at (4, 3);
            \coordinate (A) at (0, 3);
            \coordinate (M) at (2, 0);
            \coordinate (N) at (2, 3);
            % Vẽ các cạnh của hình thang
            \draw[thick] (-2,0) -- (6,0) -- (4,3) -- (0,3) -- cycle;
                \draw[thick] (-2,0) -- (4,3);
                \draw[thick] (6,0) -- (0,3);
                 \draw[thick] (M) -- (N);
                 \draw[thick] (M) -- (B)--(A)-- cycle;
    % Gán nhãn cho các điểm
    \node[above left] at (0,3) {$E$};
            \node[above right] at (4,3) {$F$};
            \node[below right] at (6,0) {$K$};
            \node[below left] at (-2,0) {$H$};
                \node[below] at (M) {$I$};
                  \node[above] at (N) {$J$};
    \end{tikzpicture}
}
\end{minipage}

\vspace{3pt}\hrule\vspace{3pt}\footnotesize\textcolor{gray!80!black}{\textbf{Các mã dùng chung hình (HINH\_146):} \texttt{L08\_C3\_B2\_06}}
\end{tcolorbox}
\vspace{4pt}


\begin{tcolorbox}[
    breakable,
    colback=white,
    colframe=blue!60!black,
    boxrule=0.8pt,
    arc=3pt,
    title={\textbf{Câu \#179} \quad \textbf{MÃ CÂU: \texttt{[L08\_C3\_B2\_16]}} \quad | \quad \textbf{Nhóm: \texttt{HINH\_147}} \quad | \quad \textbf{MapID: \texttt{8H3V2}} \hfill \textcolor{yellow!90!white}{[Hình gốc .pyc]}}
]
\textbf{Dạng toán:} Bài tổng hợp 4 \hfill \textit{(Trắc nghiệm -- Vận dụng)}
\vspace{4pt}\hrule\vspace{6pt}
\begin{minipage}[t]{0.58\textwidth}
\vspace{0pt}
Cho tam giác ${IJQ}$ cân tại ${I}$ hai đường trung tuyến ${JD}$ và ${QE}$. Chọn khẳng định sai trong các khẳng định sau:

\vspace{4pt}\noindent
\textbf{A.} $\Delta{IEQ} =\Delta{IDJ} $ \quad \textbf{B.} Tứ giác ${JQDE}$ là hình thang cân \quad \textbf{C.} Tam giác ${IED}$ là tam giác cân \quad \textbf{\textcolor{blue!75!black}{D. (Đ)}} $EQ$ và $DJ$ cắt nhau tại trung điểm của mỗi đường
\end{minipage}\hfill
\begin{minipage}[t]{0.39\textwidth}
\vspace{0pt}
\centering

\resizebox{\linewidth}{!}{%
\begin{tikzpicture}
    % Định nghĩa các điểm
    \coordinate (A) at (0, 4);
    \coordinate (B) at (-2, 0);
    \coordinate (C) at (2, 0);
    \coordinate (D) at (1, 2);
        \coordinate (E) at (-1, 2);
    % Vẽ tứ giác
    \draw[thick] (A) -- (B) -- (C) -- cycle;
        \draw[thick] (D) -- (B);
        \draw[thick] (C) -- (E);
         \draw[thick] (D) -- (E);
    % Ghi nhãn các điểm
    \node[above] at (A) {$I$};
    \node[above left] at (B) {$J$};
    \node[above right] at (C) {$Q$};
    \node[above right] at (D) {$D$};
        \node[above left] at (E){$E$};
    \draw
    ($(B)!0.5!(E)$)node[rotate=40]{$//$} ($(A)!0.5!(E)$)node[rotate=40]{$//$}
    ($(A)!0.5!(D)$)node[rotate=-30]{$//$} ($(D)!0.5!(C)$)node[rotate=-30]{$//$};   
    \end{tikzpicture}
}
\end{minipage}

\vspace{3pt}\hrule\vspace{3pt}\footnotesize\textcolor{gray!80!black}{\textbf{Các mã dùng chung hình (HINH\_147):} \texttt{L08\_C3\_B2\_07}}
\end{tcolorbox}
\vspace{4pt}


\begin{tcolorbox}[
    breakable,
    colback=white,
    colframe=blue!60!black,
    boxrule=0.8pt,
    arc=3pt,
    title={\textbf{Câu \#180} \quad \textbf{MÃ CÂU: \texttt{[L08\_C3\_B2\_17]}} \quad | \quad \textbf{Nhóm: \texttt{HINH\_148}} \quad | \quad \textbf{MapID: \texttt{8H3V2}} \hfill \textcolor{yellow!90!white}{[Hình gốc .pyc]}}
]
\textbf{Dạng toán:} Bài tổng hợp 5 \hfill \textit{(Trắc nghiệm -- Vận dụng)}
\vspace{4pt}\hrule\vspace{6pt}
\begin{minipage}[t]{0.58\textwidth}
\vspace{0pt}
Cho hình thang cân ${MNPQ}$ có $ MN // PQ$ và $MN < PQ$, hai đường cao ${ME, NF }$. Chọn khẳng định sai trong các khẳng định sau?

\vspace{4pt}\noindent
\textbf{\textcolor{blue!75!black}{A. (Đ)}} $MN=2.FP$ \quad \textbf{B.} $MN=EF$ \quad \textbf{C.} $FP = \dfrac { QP- MN}{2}$ \quad \textbf{D.} $\Delta {MEQ}=\Delta {NFP}$
\end{minipage}\hfill
\begin{minipage}[t]{0.39\textwidth}
\vspace{0pt}
\centering

\resizebox{\linewidth}{!}{%
\begin{tikzpicture}
            % Define points
            \coordinate (D) at (-2, 0);
            \coordinate (C) at (6, 0);
            \coordinate (B) at (4, 3);
            \coordinate (A) at (0, 3);
            \coordinate (H) at (0, 0);
            \coordinate (K) at (4, 0);
            % Vẽ các cạnh của hình thang
            \draw[thick] (-2,0) -- (6,0) -- (4,3) -- (0,3) -- cycle;
                 \draw[thick] (A) -- (H);
                  \draw[thick] (B) -- (K);

    % Gán nhãn cho các điểm
    \node[above left] at (0,3) {$M$};
            \node[above right] at (4,3) {$N$};
            \node[below right] at (6,0) {$P$};
            \node[below left] at (-2,0) {$Q$};
                \node[below] at (H) {$E$};
                  \node[below] at (K) {$F$};
    \def\khvuong[size=#1](#2,#3,#4){%
    \draw ($(#3)!#1!(#2)$) --
    ($($(#3)!#1!(#2)$)+($(#3)!#1!(#4)$)-(#3)$)--
    ($(#3)!#1!(#4)$);
    }
    \khvuong[size=5pt](A,H,D);
    \khvuong[size=5pt](B,K,C);
    \end{tikzpicture}
}
\end{minipage}

\vspace{3pt}\hrule\vspace{3pt}\footnotesize\textcolor{gray!80!black}{\textbf{Các mã dùng chung hình (HINH\_148):} \texttt{L08\_C3\_B2\_08}}
\end{tcolorbox}
\vspace{4pt}


\begin{tcolorbox}[
    breakable,
    colback=white,
    colframe=blue!60!black,
    boxrule=0.8pt,
    arc=3pt,
    title={\textbf{Câu \#181} \quad \textbf{MÃ CÂU: \texttt{[L08\_C3\_B2\_18]}} \quad | \quad \textbf{Nhóm: \texttt{HINH\_149}} \quad | \quad \textbf{MapID: \texttt{8H3C2}} \hfill \textcolor{yellow!90!white}{[Hình gốc .pyc]}}
]
\textbf{Dạng toán:} Bài tổng hợp 4 \hfill \textit{(Trắc nghiệm -- Vận dụng cao)}
\vspace{4pt}\hrule\vspace{6pt}
\begin{minipage}[t]{0.58\textwidth}
\vspace{0pt}
Cho hình thang cân ${EFGN}$ có ${EF} // {GN}$ và ${EF} < {GN}$. Gọi ${O}$ là giao điểm của ${EN}$ và ${FG}$, ${M}$ là giao điểm của ${EG}$ và ${FN}$.

Chọn khẳng định sai trong các khẳng định sau?

\vspace{4pt}\noindent
\textbf{A.} ${O}$, ${M}$ và trung điểm của ${NG}$ thẳng hàng \quad \textbf{B.} $\Delta{EFN}=\Delta{FEG}$ \quad \textbf{C.} ${MG}={MN}$ \quad \textbf{\textcolor{blue!75!black}{D. (Đ)}} Tam giác ${OEF}$ là tam giác đều
\end{minipage}\hfill
\begin{minipage}[t]{0.39\textwidth}
\vspace{0pt}
\centering

\resizebox{\linewidth}{!}{%
\begin{tikzpicture}
    % Định nghĩa các điểm
    \coordinate (O) at (0, 4);
    \coordinate (D) at (-2, 0);
    \coordinate (C) at (2, 0);
    \coordinate (B) at (1, 2);
        \coordinate (A) at (-1, 2);
            \coordinate (E) at (0, 4/3);
    % Vẽ tứ giác
    \draw[thick] (O) -- (D) -- (C) -- cycle;
        \draw[thick] (D) -- (B);
        \draw[thick] (C) -- (A);
         \draw[thick] (B) -- (A);
    % Ghi nhãn các điểm
    \node[above] at (O) {$O$};
    \node[above right] at (B) {$F$};
    \node[below right] at (C) {$G$};
    \node[below left] at (D) {$N$};
        \node[above left] at (A){$E$};
   \node[above] at (E) {$M$};
    \end{tikzpicture}
}
\end{minipage}

\vspace{3pt}\hrule\vspace{3pt}\footnotesize\textcolor{gray!80!black}{\textbf{Các mã dùng chung hình (HINH\_149):} \texttt{L08\_C3\_B2\_09}}
\end{tcolorbox}
\vspace{4pt}


\begin{tcolorbox}[
    breakable,
    colback=white,
    colframe=blue!60!black,
    boxrule=0.8pt,
    arc=3pt,
    title={\textbf{Câu \#182} \quad \textbf{MÃ CÂU: \texttt{[L08\_C3\_B2\_01]}} \quad | \quad \textbf{Nhóm: \texttt{HINH\_150}} \quad | \quad \textbf{MapID: \texttt{8H3N2}} \hfill \textcolor{yellow!90!white}{[Hình gốc .pyc]}}
]
\textbf{Dạng toán:} Tính số đo góc trong hình thang \hfill \textit{(Tự luận -- Nhận biết)}
\vspace{4pt}\hrule\vspace{6pt}
\begin{minipage}[t]{0.58\textwidth}
\vspace{0pt}
Cho hình vẽ 

a) Chứng minh ${MNPQ}$ là hình thang.

b) Số đo ${x}$ bằng bao nhiêu thì ${MNPQ}$ hình thang cân.
\end{minipage}\hfill
\begin{minipage}[t]{0.39\textwidth}
\vspace{0pt}
\centering

\resizebox{\linewidth}{!}{%
\begin{tikzpicture}
        % Define points
    \coordinate (D) at (0, 0);
    \coordinate (C) at (4, 0);
    \coordinate (B) at (5, 2);
    \coordinate (A) at (-1, 2);
        \coordinate (E) at (6, 2);
    % Vẽ các cạnh của hình thang
    \draw (0,0) -- (4,0) -- (5,2) -- (-1,2) -- cycle;
    \draw (5,2) -- (6,2);
    % Gán nhãn cho các điểm
    \node[above left] at (-1,2) {$M$};
    \node[above right] at (5,2) {$N$};
    \node[below right] at (4,0) {$P$};
    \node[below left] at (0,0) {$Q$};
    % Gán biến x tại góc D
    \node[left] at (0.4,0.2) {$x$};
    
    % Vẽ đoạn thẳng nối AB
    \draw (-1,2) -- (5,2);
    \draw 
    (5,2) pic[draw,angle radius = 6] {angle = C--B--E} node[shift={(-25:16pt)}]{$ 139^\circ $} 
    (4,0) pic[draw,angle radius = 6] {angle = B--C--D} node[shift={(-45:-15pt)}]{$ 139^\circ $}   
    ;

    \end{tikzpicture}
}
\end{minipage}

\end{tcolorbox}
\vspace{4pt}


\begin{tcolorbox}[
    breakable,
    colback=white,
    colframe=blue!60!black,
    boxrule=0.8pt,
    arc=3pt,
    title={\textbf{Câu \#183} \quad \textbf{MÃ CÂU: \texttt{[L08\_C3\_B2\_02]}} \quad | \quad \textbf{Nhóm: \texttt{HINH\_151}} \quad | \quad \textbf{MapID: \texttt{8H3N2}} \hfill \textcolor{yellow!90!white}{[Hình gốc .pyc]}}
]
\textbf{Dạng toán:} Chứng minh hình thang cân và Tính số đo góc trong hình thang \hfill \textit{(Tự luận -- Nhận biết)}
\vspace{4pt}\hrule\vspace{6pt}
\begin{minipage}[t]{0.58\textwidth}
\vspace{0pt}
Cho hình vẽ 

a) Cho biết hình thang ${PQRM}$ là hình thang gì?

b) Tính $\widehat{P}$, $\widehat{Q}$.
\end{minipage}\hfill
\begin{minipage}[t]{0.39\textwidth}
\vspace{0pt}
\centering

\resizebox{\linewidth}{!}{%
\begin{tikzpicture}
        % Define points
    \coordinate (D) at (0, 0);
    \coordinate (C) at (4, 0);
    \coordinate (B) at (5, 2);
    \coordinate (A) at (-1, 2);
        \coordinate (E) at (6, 2);
    % Vẽ các cạnh của hình thang
    \draw (0,0) -- (4,0) -- (5,2) -- (-1,2) -- cycle;
    % Gán nhãn cho các điểm
    \node[above left] at (-1,2) {$P$};
    \node[above right] at (5,2) {$Q$};
    \node[below right] at (4,0) {$R$};
    \node[below left] at (0,0) {$M$};
    \draw 
            (0,0) pic[draw,angle radius = 6] {angle = C--D--A} node[shift={(45:16pt)}]{$ 129^\circ $} 
            (4,0) pic[draw,angle radius = 6] {angle = B--C--D} node[shift={(-45:-15pt)}]{$ 129^\circ $}   
            ;
    \end{tikzpicture}
}
\end{minipage}

\end{tcolorbox}
\vspace{4pt}


\begin{tcolorbox}[
    breakable,
    colback=white,
    colframe=blue!60!black,
    boxrule=0.8pt,
    arc=3pt,
    title={\textbf{Câu \#184} \quad \textbf{MÃ CÂU: \texttt{[L08\_C3\_B2\_03]}} \quad | \quad \textbf{Nhóm: \texttt{HINH\_144}} \quad | \quad \textbf{MapID: \texttt{8H3H2}} \hfill \textcolor{yellow!90!white}{[Hình gốc .pyc]}}
]
\textbf{Dạng toán:} Chứng minh hình thang cân và góc bằng nhau \hfill \textit{(Tự luận -- Thông hiểu)}
\vspace{4pt}\hrule\vspace{6pt}
\begin{minipage}[t]{0.58\textwidth}
\vspace{0pt}
Cho hình thang ${MNPQ}$ biết  $MP=NQ$  

a) Hình thang ${MNPQ}$ là hình thang gì?

b) Chứng minh $\widehat{QPM}= \widehat{NQP}$.
\end{minipage}\hfill
\begin{minipage}[t]{0.39\textwidth}
\vspace{0pt}
\centering

\resizebox{\linewidth}{!}{%
\begin{tikzpicture}
            % Define points
            \coordinate (D) at (-2, 0);
            \coordinate (C) at (6, 0);
            \coordinate (B) at (4, 3);
            \coordinate (A) at (0, 3);
        
            % Vẽ các cạnh của hình thang
            \draw (-2,0) -- (6,0) -- (4,3) -- (0,3) -- cycle;
                \draw (-2,0) -- (4,3);
                \draw (6,0) -- (0,3);
    % Gán nhãn cho các điểm
    \node[above left] at (0,3) {$M$};
            \node[above right] at (4,3) {$N$};
            \node[below right] at (6,0) {$P$};
            \node[below left] at (-2,0) {$Q$};
   
    \end{tikzpicture}
}
\end{minipage}

\vspace{3pt}\hrule\vspace{3pt}\footnotesize\textcolor{gray!80!black}{\textbf{Các mã dùng chung hình (HINH\_144):} \texttt{L08\_C3\_B2\_12, L08\_C3\_B2\_14, L08\_C3\_B2\_05}}
\end{tcolorbox}
\vspace{4pt}


\begin{tcolorbox}[
    breakable,
    colback=white,
    colframe=blue!60!black,
    boxrule=0.8pt,
    arc=3pt,
    title={\textbf{Câu \#185} \quad \textbf{MÃ CÂU: \texttt{[L08\_C3\_B2\_04]}} \quad | \quad \textbf{Nhóm: \texttt{HINH\_152}} \quad | \quad \textbf{MapID: \texttt{8H3V2}} \hfill \textcolor{yellow!90!white}{[Hình gốc .pyc]}}
]
\textbf{Dạng toán:} Bài tổng hợp 1. \hfill \textit{(Tự luận -- Vận dụng)}
\vspace{4pt}\hrule\vspace{6pt}
\begin{minipage}[t]{0.58\textwidth}
\vspace{0pt}
Cho tam giác ${IJQ}$, hai đường phân giác $\widehat{J}$, $\widehat{Q}$ cắt nhau tại ${O}$. Qua ${O}$ kẻ đường thẳng song song với $JQ$  đường thẳng này cắt ${IJ}$, ${IQ}$ lần lượt tại ${M}, {N}$.

a) Tứ giác ${JQOM}$ và tứ giác ${JQNO}$ là hình gì?

b) Chứng minh rằng $MN=MJ+NQ$
\end{minipage}\hfill
\begin{minipage}[t]{0.39\textwidth}
\vspace{0pt}
\centering

\resizebox{\linewidth}{!}{%
\begin{tikzpicture}
    % Vẽ tam giác ACB (hoán đổi vị trí của A và C)
    \coordinate (A) at (-2, 4);
    \coordinate (B) at (-3, 0);
    \coordinate (C) at (2, 0);
        \coordinate (M) at (-2.66, 1.35);
    \coordinate (N) at (0.65, 1.35);
        \coordinate (O) at (-1.27, 1.35);
    % Vẽ các cạnh tam giác
    \draw [thick] (A) -- (B) -- (C) -- cycle;
        \draw [thick] (M) -- (O) -- (N);
            \draw [thick] (B) -- (O);
                \draw [thick] (C) -- (O);
    % Gán nhãn cho các điểm
    \node[above] at (A) {$I$};
    \node[below right] at (B) {$J$};
    \node[below left] at (C) {$Q$};
    % Tính toán giao điểm của hai đường phân giác

    % Vẽ điểm O và gán nhãn
    \filldraw[black] (O) circle (1pt); % Vẽ điểm O
    \node[above right] at (O) {O};

    \filldraw[blue] (M) circle (1pt);
    \filldraw[blue] (N) circle (1pt);
    
    % Gán nhãn cho M và N
    \node[above left] at (M) {$M$};
    \node[above right] at (N) {$N$};
    \draw
        pic[draw,blue,angle radius=3mm] {angle = O--C--B} pic[draw,blue,angle radius=3.7mm] {angle = O--C--B};
    \draw
    pic[draw,blue,angle radius=4mm] {angle = N--C--O} pic[draw,blue,angle radius=4.7mm] {angle = N--C--O};
    \draw
     pic[draw,blue,angle radius=3.4mm] {angle = C--B--O};
    \draw
    pic[draw,blue,angle radius=4.4mm] {angle = O--B--M} ;
    \end{tikzpicture}
}
\end{minipage}

\end{tcolorbox}
\vspace{4pt}


\begin{tcolorbox}[
    breakable,
    colback=white,
    colframe=blue!60!black,
    boxrule=0.8pt,
    arc=3pt,
    title={\textbf{Câu \#186} \quad \textbf{MÃ CÂU: \texttt{[L08\_C3\_B2\_05]}} \quad | \quad \textbf{Nhóm: \texttt{HINH\_144}} \quad | \quad \textbf{MapID: \texttt{8H3V2}} \hfill \textcolor{yellow!90!white}{[Hình gốc .pyc]}}
]
\textbf{Dạng toán:} Bài tổng hợp 2. \hfill \textit{(Tự luận -- Vận dụng)}
\vspace{4pt}\hrule\vspace{6pt}
\begin{minipage}[t]{0.58\textwidth}
\vspace{0pt}
Cho hình thang cân ${IJQM}$ có $IJ // QM$ và $IJ < QM$, biết $IM= JI$

a) Chứng minh rằng $IJ=QJ$. 

b) Chứng minh rằng ${MJ}$ là tia phân giác của $\widehat{IMQ }$.
\end{minipage}\hfill
\begin{minipage}[t]{0.39\textwidth}
\vspace{0pt}
\centering

\resizebox{\linewidth}{!}{%
\begin{tikzpicture}
            % Define points
            \coordinate (D) at (-2, 0);
            \coordinate (C) at (6, 0);
            \coordinate (B) at (4, 3);
            \coordinate (A) at (0, 3);
        
            % Vẽ các cạnh của hình thang
            \draw (-2,0) -- (6,0) -- (4,3) -- (0,3) -- cycle;
                \draw (-2,0) -- (4,3);
                \draw (6,0) -- (0,3);
    % Gán nhãn cho các điểm
    \node[above left] at (0,3) {$I$};
            \node[above right] at (4,3) {$J$};
            \node[below right] at (6,0) {$Q$};
            \node[below left] at (-2,0) {$M$};
   
    \end{tikzpicture}
}
\end{minipage}

\vspace{3pt}\hrule\vspace{3pt}\footnotesize\textcolor{gray!80!black}{\textbf{Các mã dùng chung hình (HINH\_144):} \texttt{L08\_C3\_B2\_12, L08\_C3\_B2\_14, L08\_C3\_B2\_03}}
\end{tcolorbox}
\vspace{4pt}


\begin{tcolorbox}[
    breakable,
    colback=white,
    colframe=blue!60!black,
    boxrule=0.8pt,
    arc=3pt,
    title={\textbf{Câu \#187} \quad \textbf{MÃ CÂU: \texttt{[L08\_C3\_B2\_06]}} \quad | \quad \textbf{Nhóm: \texttt{HINH\_146}} \quad | \quad \textbf{MapID: \texttt{8H3V2}} \hfill \textcolor{yellow!90!white}{[Hình gốc .pyc]}}
]
\textbf{Dạng toán:} Bài tổng hợp 3. \hfill \textit{(Tự luận -- Vận dụng)}
\vspace{4pt}\hrule\vspace{6pt}
\begin{minipage}[t]{0.58\textwidth}
\vspace{0pt}
Cho hình thang cân ${IJQK}$ có $IJ // QK$. Lấy ${N}, {M}$ lần lượt là trung điểm của ${IJ}$, ${QK}$

a) Chứng minh $IM=JM$.

b) Chứng minh ${ MN}$ là đường cao của hình thang.
\end{minipage}\hfill
\begin{minipage}[t]{0.39\textwidth}
\vspace{0pt}
\centering

\resizebox{\linewidth}{!}{%
\begin{tikzpicture}
            % Define points
            \coordinate (D) at (-2, 0);
            \coordinate (C) at (6, 0);
            \coordinate (B) at (4, 3);
            \coordinate (A) at (0, 3);
            \coordinate (M) at (2, 0);
            \coordinate (N) at (2, 3);
            % Vẽ các cạnh của hình thang
            \draw[thick] (-2,0) -- (6,0) -- (4,3) -- (0,3) -- cycle;
                \draw[thick] (-2,0) -- (4,3);
                \draw[thick] (6,0) -- (0,3);
                 \draw[thick] (M) -- (N);
                 \draw[thick] (M) -- (B)--(A)-- cycle;
    % Gán nhãn cho các điểm
    \node[above left] at (0,3) {$I$};
            \node[above right] at (4,3) {$J$};
            \node[below right] at (6,0) {$Q$};
            \node[below left] at (-2,0) {$K$};
                \node[below] at (M) {$M$};
                  \node[above] at (N) {$N$};
    \end{tikzpicture}
}
\end{minipage}

\vspace{3pt}\hrule\vspace{3pt}\footnotesize\textcolor{gray!80!black}{\textbf{Các mã dùng chung hình (HINH\_146):} \texttt{L08\_C3\_B2\_15}}
\end{tcolorbox}
\vspace{4pt}


\begin{tcolorbox}[
    breakable,
    colback=white,
    colframe=blue!60!black,
    boxrule=0.8pt,
    arc=3pt,
    title={\textbf{Câu \#188} \quad \textbf{MÃ CÂU: \texttt{[L08\_C3\_B2\_07]}} \quad | \quad \textbf{Nhóm: \texttt{HINH\_147}} \quad | \quad \textbf{MapID: \texttt{8H3V2}} \hfill \textcolor{yellow!90!white}{[Hình gốc .pyc]}}
]
\textbf{Dạng toán:} Bài tổng hợp 4. \hfill \textit{(Tự luận -- Vận dụng)}
\vspace{4pt}\hrule\vspace{6pt}
\begin{minipage}[t]{0.58\textwidth}
\vspace{0pt}
Cho tam giác ${ABC}$ cận tại ${A}$ hai đường trung tuyến ${BM}$ và ${CN}$.

a) Chứng minh rằng tam giác ${ANM}$ là tam giác cân. 

b) Chứng minh tứ giác ${BCMN}$ là hình thang cân.
\end{minipage}\hfill
\begin{minipage}[t]{0.39\textwidth}
\vspace{0pt}
\centering

\resizebox{\linewidth}{!}{%
\begin{tikzpicture}
    % Định nghĩa các điểm
    \coordinate (A) at (0, 4);
    \coordinate (B) at (-2, 0);
    \coordinate (C) at (2, 0);
    \coordinate (D) at (1, 2);
        \coordinate (E) at (-1, 2);
    % Vẽ tứ giác
    \draw[thick] (A) -- (B) -- (C) -- cycle;
        \draw[thick] (D) -- (B);
        \draw[thick] (C) -- (E);
         \draw[thick] (D) -- (E);
    % Ghi nhãn các điểm
    \node[above] at (A) {$A$};
    \node[above left] at (B) {$B$};
    \node[above right] at (C) {$C$};
    \node[above right] at (D) {$M$};
        \node[above left] at (E){$N$};
    \draw
    ($(B)!0.5!(E)$)node[rotate=40]{$//$} ($(A)!0.5!(E)$)node[rotate=40]{$//$}
    ($(A)!0.5!(D)$)node[rotate=-30]{$//$} ($(D)!0.5!(C)$)node[rotate=-30]{$//$};   
    \end{tikzpicture}
}
\end{minipage}

\vspace{3pt}\hrule\vspace{3pt}\footnotesize\textcolor{gray!80!black}{\textbf{Các mã dùng chung hình (HINH\_147):} \texttt{L08\_C3\_B2\_16}}
\end{tcolorbox}
\vspace{4pt}


\begin{tcolorbox}[
    breakable,
    colback=white,
    colframe=blue!60!black,
    boxrule=0.8pt,
    arc=3pt,
    title={\textbf{Câu \#189} \quad \textbf{MÃ CÂU: \texttt{[L08\_C3\_B2\_08]}} \quad | \quad \textbf{Nhóm: \texttt{HINH\_148}} \quad | \quad \textbf{MapID: \texttt{8H3V2}} \hfill \textcolor{yellow!90!white}{[Hình gốc .pyc]}}
]
\textbf{Dạng toán:} Bài tổng hợp 5 \hfill \textit{(Tự luận -- Vận dụng)}
\vspace{4pt}\hrule\vspace{6pt}
\begin{minipage}[t]{0.58\textwidth}
\vspace{0pt}
Cho hình thang cân ${EFKH}$ có $EF // KH$ và $EF < KH$, hai đường cao $EM, FN$.

a) Chứng minh $\Delta {EMH}=\Delta {FNK}$.

b) Chứng minh $EF=MN$.

c) Chỉ ra $NK = \dfrac { HK- EF}{2}$.
\end{minipage}\hfill
\begin{minipage}[t]{0.39\textwidth}
\vspace{0pt}
\centering

\resizebox{\linewidth}{!}{%
\begin{tikzpicture}
            % Define points
            \coordinate (D) at (-2, 0);
            \coordinate (C) at (6, 0);
            \coordinate (B) at (4, 3);
            \coordinate (A) at (0, 3);
            \coordinate (H) at (0, 0);
            \coordinate (K) at (4, 0);
            % Vẽ các cạnh của hình thang
            \draw[thick] (-2,0) -- (6,0) -- (4,3) -- (0,3) -- cycle;
                 \draw[thick] (A) -- (H);
                  \draw[thick] (B) -- (K);

    % Gán nhãn cho các điểm
    \node[above left] at (0,3) {$E$};
            \node[above right] at (4,3) {$F$};
            \node[below right] at (6,0) {$K$};
            \node[below left] at (-2,0) {$H$};
                \node[below] at (H) {$M$};
                  \node[below] at (K) {$N$};
    \def\khvuong[size=#1](#2,#3,#4){%
    \draw ($(#3)!#1!(#2)$) --
    ($($(#3)!#1!(#2)$)+($(#3)!#1!(#4)$)-(#3)$)--
    ($(#3)!#1!(#4)$);
    }
    \khvuong[size=5pt](A,H,D);
    \khvuong[size=5pt](B,K,C);
    \end{tikzpicture}
}
\end{minipage}

\vspace{3pt}\hrule\vspace{3pt}\footnotesize\textcolor{gray!80!black}{\textbf{Các mã dùng chung hình (HINH\_148):} \texttt{L08\_C3\_B2\_17}}
\end{tcolorbox}
\vspace{4pt}


\begin{tcolorbox}[
    breakable,
    colback=white,
    colframe=blue!60!black,
    boxrule=0.8pt,
    arc=3pt,
    title={\textbf{Câu \#190} \quad \textbf{MÃ CÂU: \texttt{[L08\_C3\_B2\_09]}} \quad | \quad \textbf{Nhóm: \texttt{HINH\_149}} \quad | \quad \textbf{MapID: \texttt{8H3V2}} \hfill \textcolor{yellow!90!white}{[Hình gốc .pyc]}}
]
\textbf{Dạng toán:} Bài tổng hợp 4 \hfill \textit{(Tự luận -- Vận dụng)}
\vspace{4pt}\hrule\vspace{6pt}
\begin{minipage}[t]{0.58\textwidth}
\vspace{0pt}
Cho hình thang cân ${IJQD}$ có ${IJ} // {QD}$ và ${IJ} < {QD}$. Gọi ${O}$ là giao điểm của ${ID}$ và ${JQ}$, ${E}$ là giao điểm của ${IQ}$ và ${JD}$.

a) Chứng minh rằng tam giác ${OIJ}$ là tam giác cân tại  ${O}$. 

b) Chứng minh  $\Delta{IJD}=\Delta{JIQ}$.

c) Chứng minh  ${EQ}={ED}$ 

d) Chứng minh ${O}$, ${E}$ và trung điểm của ${DQ}$ thẳng hàng.
\end{minipage}\hfill
\begin{minipage}[t]{0.39\textwidth}
\vspace{0pt}
\centering

\resizebox{\linewidth}{!}{%
\begin{tikzpicture}
    % Định nghĩa các điểm
    \coordinate (O) at (0, 4);
    \coordinate (D) at (-2, 0);
    \coordinate (C) at (2, 0);
    \coordinate (B) at (1, 2);
        \coordinate (A) at (-1, 2);
            \coordinate (E) at (0, 4/3);
    % Vẽ tứ giác
    \draw[thick] (O) -- (D) -- (C) -- cycle;
        \draw[thick] (D) -- (B);
        \draw[thick] (C) -- (A);
         \draw[thick] (B) -- (A);
    % Ghi nhãn các điểm
    \node[above] at (O) {$O$};
    \node[above right] at (B) {$J$};
    \node[below right] at (C) {$Q$};
    \node[below left] at (D) {$D$};
        \node[above left] at (A){$I$};
   \node[above] at (E) {$E$};
    \end{tikzpicture}
}
\end{minipage}

\vspace{3pt}\hrule\vspace{3pt}\footnotesize\textcolor{gray!80!black}{\textbf{Các mã dùng chung hình (HINH\_149):} \texttt{L08\_C3\_B2\_18}}
\end{tcolorbox}
\vspace{4pt}


\begin{tcolorbox}[
    breakable,
    colback=white,
    colframe=blue!60!black,
    boxrule=0.8pt,
    arc=3pt,
    title={\textbf{Câu \#191} \quad \textbf{MÃ CÂU: \texttt{[L08\_C3\_B3\_01]}} \quad | \quad \textbf{Nhóm: \texttt{HINH\_090}} \quad | \quad \textbf{MapID: \texttt{8H3N3}} \hfill \textcolor{yellow!90!white}{[Hình gốc .pyc]}}
]
\textbf{Dạng toán:} Nhận biết về hình bình hành \hfill \textit{(Trắc nghiệm -- Nhận biết)}
\vspace{4pt}\hrule\vspace{6pt}
\begin{minipage}[t]{0.58\textwidth}
\vspace{0pt}
Cho hình bình hành ${MNPQ}$ như hình vẽ. Biết $\widehat{ NMQ } =58^{\circ}$ và ${O}$ là trung điểm của ${NQ}$.

Chọn khẳng định sai trong các khẳng định sau?

\vspace{4pt}\noindent
\textbf{A.} $\widehat{MNP}=122^{ \circ}$ \quad \textbf{\textcolor{blue!75!black}{B. (Đ)}} $\Delta{PMN}= \Delta{NQP}$ \quad \textbf{C.} ${O}$, ${M}$, ${P}$ thẳng hàng \quad \textbf{D.} $\widehat{MQP}=122^{ \circ}$
\end{minipage}\hfill
\begin{minipage}[t]{0.39\textwidth}
\vspace{0pt}
\centering

\resizebox{\linewidth}{!}{%
\begin{tikzpicture}
            \coordinate (O) at (3/2, 3/2);
            \coordinate (D) at (2, 0);
            \coordinate (C) at (5, 3);
            \coordinate (B) at (1, 3);
            \coordinate (A) at (-2, 0);
            \
            % Vẽ tứ giác
            \draw[thick] (O) -- (A) -- (C) ;
            \draw[thick] (A) -- (B)--(C)--(D)--(A);
        \draw[thick] (O) -- (B) -- (D) ;
            % Ghi nhãn các điểm
            \node[above right] at (O) {$O$};
            \node[above left] at (B) {$N$};
            \node[above right] at (C) {$P$};
            \node[below right] at (D) {$Q$};
            \node[below left] at (A){$M$};
                \draw
            ($(B)!0.4!(O)$)node[rotate=80]{$/$} ($(O)!0.4!(D)$)node[rotate=80]{$/$};
    \end{tikzpicture}
}
\end{minipage}

\vspace{3pt}\hrule\vspace{3pt}\footnotesize\textcolor{gray!80!black}{\textbf{Các mã dùng chung hình (HINH\_090):} \texttt{L07\_C3\_B1\_07, L08\_C3\_B3\_09}}
\end{tcolorbox}
\vspace{4pt}


\begin{tcolorbox}[
    breakable,
    colback=white,
    colframe=blue!60!black,
    boxrule=0.8pt,
    arc=3pt,
    title={\textbf{Câu \#192} \quad \textbf{MÃ CÂU: \texttt{[L08\_C3\_B3\_02]}} \quad | \quad \textbf{Nhóm: \texttt{HINH\_153}} \quad | \quad \textbf{MapID: \texttt{8H3H3}} \hfill \textcolor{yellow!90!white}{[Hình gốc .pyc]}}
]
\textbf{Dạng toán:} Thông hiểu về hình bình hành 1 \hfill \textit{(Trắc nghiệm -- Thông hiểu)}
\vspace{4pt}\hrule\vspace{6pt}
\begin{minipage}[t]{0.58\textwidth}
\vspace{0pt}
Cho $\Delta{EFK}$ cân tại ${E}$ có điểm ${H}$ trên cạnh ${FK}$. Kẻ $HI//EK$, $HJ//EF$. Chọn khẳng định sai trong các khẳng định sau?

\vspace{4pt}\noindent
\textbf{A.} $\Delta{FHI}$ cân tại ${I}$ \quad \textbf{\textcolor{blue!75!black}{B. (Đ)}} $\Delta{JHK}$ là tam giác đều \quad \textbf{C.} $HI+HJ=FI+EI=EF$ \quad \textbf{D.} Tứ giác ${EIHJ}$ là hình bình hành.
\end{minipage}\hfill
\begin{minipage}[t]{0.39\textwidth}
\vspace{0pt}
\centering

\resizebox{\linewidth}{!}{%
\begin{tikzpicture}
            % Định nghĩa các điểm
            \coordinate (A) at (0, 4);
            \coordinate (B) at (-2, 0);
            \coordinate (C) at (2, 0);
            \coordinate (D) at (1, 0);
            \coordinate (M) at (-0.5, 3);
                \coordinate (N) at (1.5, 1);
            % Vẽ tứ giác
            \draw[thick] (A) -- (B) -- (C) -- cycle;
            \draw[thick] (D) -- (M);
            \draw[thick] (D) -- (N);
            % Ghi nhãn các điểm
            \node[above] at (A) {$E $};
            \node[above left] at (B) {$F $};
            \node[above right] at (C) {$K $};
            \node[below] at (D) {$H $};
            \node[above left] at (M){$I $};
            \node[above right] at (N){$J $};
        \end{tikzpicture}
}
\end{minipage}

\vspace{3pt}\hrule\vspace{3pt}\footnotesize\textcolor{gray!80!black}{\textbf{Các mã dùng chung hình (HINH\_153):} \texttt{L08\_C3\_B3\_10}}
\end{tcolorbox}
\vspace{4pt}


\begin{tcolorbox}[
    breakable,
    colback=white,
    colframe=blue!60!black,
    boxrule=0.8pt,
    arc=3pt,
    title={\textbf{Câu \#193} \quad \textbf{MÃ CÂU: \texttt{[L08\_C3\_B3\_03]}} \quad | \quad \textbf{Nhóm: \texttt{HINH\_154}} \quad | \quad \textbf{MapID: \texttt{8H3V3}} \hfill \textcolor{yellow!90!white}{[Hình gốc .pyc]}}
]
\textbf{Dạng toán:} Bài tổng hợp 1 \hfill \textit{(Trắc nghiệm -- Vận dụng)}
\vspace{4pt}\hrule\vspace{6pt}
\begin{minipage}[t]{0.58\textwidth}
\vspace{0pt}
Cho hình bình hành ${A B C D }$ có $A B  < C D $. Tia phân giác của góc $\widehat{B }$, $\widehat{D }$ lần lượt cắt ${A D , B C }$ tại ${M , N }$. Trong các khẳng định sau khẳng định nào sai?

\vspace{4pt}\noindent
\textbf{\textcolor{blue!75!black}{A. (Đ)}} Tứ giác ${B C D M }$ là hình thang cân \quad \textbf{B.} $\Delta{A B M }$ cân tại ${A }$ \quad \textbf{C.} $\Delta{C N D }$ cân tại ${C }$ \quad \textbf{D.} Tứ giác ${B M D N }$ là hình bình hành
\end{minipage}\hfill
\begin{minipage}[t]{0.39\textwidth}
\vspace{0pt}
\centering

\resizebox{\linewidth}{!}{%
\begin{tikzpicture}
            \coordinate (M) at (-0.76, 0);
            \coordinate (D) at (2, 0);
            \coordinate (C) at (3, 2);
            \coordinate (B) at (-2, 2);
            \coordinate (A) at (-3, 0);
            \coordinate (N) at (0.76, 2);
            % Vẽ tứ giác
            \draw[thick] (B) -- (M) ;
            \draw[thick] (A) -- (B)--(C)--(D)--(A);
        \draw[thick] (N) -- (D) ;
            % Ghi nhãn các điểm
            \node[below right] at (M) {$M $};
            \node[above left] at (B) {$B $};
            \node[above right] at (C) {$C $};
            \node[below right] at (D) {$D $};
            \node[below left] at (A){$A $};
            \node[above right] at (N) {$N $};
                \draw 
            pic[draw,blue,angle radius=4mm] {angle = M--B--C}; 
             \draw  pic[draw,blue,angle radius=4.7mm] {angle = A--B--M};
             \draw  pic[draw,blue,angle radius=3mm] {angle = C--D--N} ;
            \draw pic[draw,blue,angle radius=3.7mm] {angle = N--D--M};
    \end{tikzpicture}
}
\end{minipage}

\vspace{3pt}\hrule\vspace{3pt}\footnotesize\textcolor{gray!80!black}{\textbf{Các mã dùng chung hình (HINH\_154):} \texttt{L08\_C3\_B3\_11}}
\end{tcolorbox}
\vspace{4pt}


\begin{tcolorbox}[
    breakable,
    colback=white,
    colframe=blue!60!black,
    boxrule=0.8pt,
    arc=3pt,
    title={\textbf{Câu \#194} \quad \textbf{MÃ CÂU: \texttt{[L08\_C3\_B3\_04]}} \quad | \quad \textbf{Nhóm: \texttt{HINH\_155}} \quad | \quad \textbf{MapID: \texttt{8H3H3}} \hfill \textcolor{yellow!90!white}{[Hình gốc .pyc]}}
]
\textbf{Dạng toán:} Thông hiểu hình bình hành 2 \hfill \textit{(Trắc nghiệm -- Thông hiểu)}
\vspace{4pt}\hrule\vspace{6pt}
\begin{minipage}[t]{0.58\textwidth}
\vspace{0pt}
Cho hình bình hành ${QKHO}$. Trên cạnh ${QK}$ lấy điểm ${F}$ trên cạnh ${HO}$ lấy điểm ${E}$ sao cho $QF=HE$. Chọn khẳng sai trong các khẳng định sau?

\vspace{4pt}\noindent
\textbf{A.} ${KF=OE}$ \quad \textbf{B.} ${QFHE}$ là hình bình hành \quad \textbf{\textcolor{blue!75!black}{C. (Đ)}} $\Delta { KHE } =\Delta {FOQ}$ \quad \textbf{D.} ${OFKE}$ là hình bình hành
\end{minipage}\hfill
\begin{minipage}[t]{0.39\textwidth}
\vspace{0pt}
\centering

\resizebox{\linewidth}{!}{%
\begin{tikzpicture}
            \coordinate (M) at (-2.34, 1.31);
            \coordinate (D) at (2, 0);
            \coordinate (C) at (3, 2);
            \coordinate (B) at (-2, 2);
            \coordinate (A) at (-3, 0);
            \coordinate (N) at (2.34, 0.69);
            % Vẽ tứ giác
            \draw[thick] (D) -- (M) ;
            \draw[thick] (A) -- (B)--(C)--(D)--(A);
        \draw[thick] (N) -- (B) ;
            % Ghi nhãn các điểm
            \node[above left] at (M) {$F$};
            \node[above left] at (B) {$K$};
            \node[above right] at (C) {$H$};
            \node[below right] at (D) {$O$};
            \node[below left] at (A){$Q$};
            \node[below right] at (N) {$E$};
            \draw
        ($(A)!0.5!(M)$)node[rotate=20]{$/$} ($(N)!0.5!(C)$)node[rotate=20]{$/$};
    \end{tikzpicture}
}
\end{minipage}

\vspace{3pt}\hrule\vspace{3pt}\footnotesize\textcolor{gray!80!black}{\textbf{Các mã dùng chung hình (HINH\_155):} \texttt{L08\_C3\_B3\_12}}
\end{tcolorbox}
\vspace{4pt}


\begin{tcolorbox}[
    breakable,
    colback=white,
    colframe=blue!60!black,
    boxrule=0.8pt,
    arc=3pt,
    title={\textbf{Câu \#195} \quad \textbf{MÃ CÂU: \texttt{[L08\_C3\_B3\_05]}} \quad | \quad \textbf{Nhóm: \texttt{HINH\_156}} \quad | \quad \textbf{MapID: \texttt{8H3H3}} \hfill \textcolor{yellow!90!white}{[Hình gốc .pyc]}}
]
\textbf{Dạng toán:} Thông hiểu hình bình hành 3 \hfill \textit{(Trắc nghiệm -- Thông hiểu)}
\vspace{4pt}\hrule\vspace{6pt}
\begin{minipage}[t]{0.58\textwidth}
\vspace{0pt}
Cho hình bình hành ${MNPQ}$ có $MN>PQ$. Từ ${M}$ kẻ đường thẳng vuông góc với ${NQ}$ cắt ${QP}$ tại ${E}$, từ ${P}$ vẽ đường thẳng vuông góc với ${NQ}$ cắt ${MN}$ tại ${F}$. Gọi ${O}$ là trung điểm của ${NQ}$. Trong các khẳng định sau khẳng định nào sai?

\vspace{4pt}\noindent
\textbf{A.} ${O}$ là trung điểm của ${EF}$ \quad \textbf{B.} ${NFQE}$ là hình bình hành \quad \textbf{\textcolor{blue!75!black}{C. (Đ)}} ${MNPE}$ là hình thang cân \quad \textbf{D.} ${MEPF}$ là hình bình hành
\end{minipage}\hfill
\begin{minipage}[t]{0.39\textwidth}
\vspace{0pt}
\centering

\resizebox{\linewidth}{!}{%
\begin{tikzpicture}
            \coordinate (O) at (0, 1);
            \coordinate (D) at (-3, 0);
            \coordinate (C) at (2, 0);
            \coordinate (B) at (3, 2);
            \coordinate (A) at (-2, 2);
            \coordinate (H) at (-1.33, 0);
            \coordinate (K) at (1.33, 2);
            % Vẽ tứ giác
            \draw[thick] (D) -- (B);
            \draw[thick] (A) -- (B) -- (C) -- (D) -- (A);
            \draw[thick] (C) -- (K);
            \draw[thick] (A) -- (H);
            \draw[thick] (K) -- (H);
            % Ghi nhãn các điểm
            \node[below left] at (H) {$E$};
            \node[above right] at (B) {$N$};
            \node[below right] at (C) {$P$};
            \node[below left] at (D) {$Q$};
            \node[above left] at (A){$M$};
            \node[above right] at (K) {$F$};
            \node[below] at (O) {$O$};
    \end{tikzpicture}
}
\end{minipage}

\vspace{3pt}\hrule\vspace{3pt}\footnotesize\textcolor{gray!80!black}{\textbf{Các mã dùng chung hình (HINH\_156):} \texttt{L08\_C3\_B3\_13}}
\end{tcolorbox}
\vspace{4pt}


\begin{tcolorbox}[
    breakable,
    colback=white,
    colframe=blue!60!black,
    boxrule=0.8pt,
    arc=3pt,
    title={\textbf{Câu \#196} \quad \textbf{MÃ CÂU: \texttt{[L08\_C3\_B3\_06]}} \quad | \quad \textbf{Nhóm: \texttt{HINH\_157}} \quad | \quad \textbf{MapID: \texttt{8H3V3}} \hfill \textcolor{yellow!90!white}{[Hình gốc .pyc]}}
]
\textbf{Dạng toán:} Tổng hợp hình bình hành 1 \hfill \textit{(Trắc nghiệm -- Vận dụng)}
\vspace{4pt}\hrule\vspace{6pt}
\begin{minipage}[t]{0.58\textwidth}
\vspace{0pt}
Cho tam giác ${QKH}$ cân tại ${Q}$ lấy điểm ${G}$ bất kỳ trên ${QK}$, lấy điểm ${E}$ trên tia đối của tia ${HQ}$ sao cho $HE=KG$. Từ ${G}$ kẻ đường thẳng song song với ${QH}$ cắt ${KH}$ tại ${F}$. Chọn khẳng định sai trong các khẳng định sau?

\vspace{4pt}\noindent
\textbf{A.} Tam giác ${GKF}$ là tam giác cân \quad \textbf{B.} $GF=HE$ \quad \textbf{C.} Tứ giác ${GHEF}$ là hình bình hành \quad \textbf{\textcolor{blue!75!black}{D. (Đ)}} ${ QGFH }$ là hình thang cân
\end{minipage}\hfill
\begin{minipage}[t]{0.39\textwidth}
\vspace{0pt}
\centering

\resizebox{\linewidth}{!}{%
\begin{tikzpicture}
            \coordinate (O) at (0, 1);
            \coordinate (D) at (-1.2, 1.99);
            \coordinate (C) at (2, 0);
            \coordinate (B) at (-2, 0);
            \coordinate (A) at (0, 5);
            \coordinate (F) at (-0.41, 0);
            \coordinate (E) at (2.9, -1.99);
            % Vẽ tứ giác
            \draw[thick] (D) -- (C)--(E)--(F)--(D);
            \draw[thick] (A) -- (B) -- (C)-- (A);
         
            % Ghi nhãn các điểm
            \node[below right] at (E) {$E$};
            \node[below left] at (B) {$K$};
            \node[above right] at (C) {$H$};
            \node[above left] at (D) {$G$};
            \node[above] at (A){$Q$};
            \node[below left] at (F) {$F$};
           \draw
        ($(B)!0.5!(D)$)node[rotate=-40]{$/$};\draw ($(C)!0.5!(E)$)node[rotate=-40]{$/$};
    \end{tikzpicture}
}
\end{minipage}

\vspace{3pt}\hrule\vspace{3pt}\footnotesize\textcolor{gray!80!black}{\textbf{Các mã dùng chung hình (HINH\_157):} \texttt{L08\_C3\_B3\_14}}
\end{tcolorbox}
\vspace{4pt}


\begin{tcolorbox}[
    breakable,
    colback=white,
    colframe=blue!60!black,
    boxrule=0.8pt,
    arc=3pt,
    title={\textbf{Câu \#197} \quad \textbf{MÃ CÂU: \texttt{[L08\_C3\_B3\_07]}} \quad | \quad \textbf{Nhóm: \texttt{HINH\_158}} \quad | \quad \textbf{MapID: \texttt{8H3V3}} \hfill \textcolor{yellow!90!white}{[Hình gốc .pyc]}}
]
\textbf{Dạng toán:} Tổng hợp hình bình hành 2 \hfill \textit{(Trắc nghiệm -- Vận dụng)}
\vspace{4pt}\hrule\vspace{6pt}
\begin{minipage}[t]{0.58\textwidth}
\vspace{0pt}
Cho hình bình hành  ${GHFE}$. Có ${I, J, M, N}$ lần lượt là trung điểm của ${ GH, HF, FE, GE }$. Chọn khẳng định sai trong các khẳng định sau?

\vspace{4pt}\noindent
\textbf{A.} $IJ=MN$ \quad \textbf{\textcolor{blue!75!black}{B. (Đ)}} ${JMEH}$ là hình thang cân \quad \textbf{C.} ${IJMN }$ là hình bình hành \quad \textbf{D.} $IN=JM$
\end{minipage}\hfill
\begin{minipage}[t]{0.39\textwidth}
\vspace{0pt}
\centering

\resizebox{\linewidth}{!}{%
\begin{tikzpicture}
            \coordinate (P) at (2.5, 1);
            \coordinate (D) at (-3, 0);
            \coordinate (C) at (2, 0);
            \coordinate (B) at (3, 2);
            \coordinate (A) at (-2, 2);
            \coordinate (M) at (-5/2, 1);
            \coordinate (N) at (1/2, 2);
             \coordinate (Q) at (-1/2, 0);
            % Vẽ tứ giác
            \draw[thick] (M) -- (N)--(P) -- (Q)--(M);
            \draw[thick] (A) -- (B) -- (C) -- (D) -- (A);
        
            % Ghi nhãn các điểm
            \node[above left] at (M) {$I$};
            \node[above right] at (B) {$H$};
            \node[below right] at (C) {$F$};
            \node[below left] at (D) {$E$};
            \node[above left] at (A){$G$};
            \node[above] at (N) {$J$};
            \node[above right] at (P) {$M$};
             \node[below] at (Q) {$N$};
    \end{tikzpicture}
}
\end{minipage}

\vspace{3pt}\hrule\vspace{3pt}\footnotesize\textcolor{gray!80!black}{\textbf{Các mã dùng chung hình (HINH\_158):} \texttt{L08\_C3\_B3\_15}}
\end{tcolorbox}
\vspace{4pt}


\begin{tcolorbox}[
    breakable,
    colback=white,
    colframe=blue!60!black,
    boxrule=0.8pt,
    arc=3pt,
    title={\textbf{Câu \#198} \quad \textbf{MÃ CÂU: \texttt{[L08\_C3\_B3\_08]}} \quad | \quad \textbf{Nhóm: \texttt{HINH\_159}} \quad | \quad \textbf{MapID: \texttt{8H3C3}} \hfill \textcolor{yellow!90!white}{[Hình gốc .pyc]}}
]
\textbf{Dạng toán:} Vận dụng hình bình hành \hfill \textit{(Trắc nghiệm -- Vận dụng cao)}
\vspace{4pt}\hrule\vspace{6pt}
\begin{minipage}[t]{0.58\textwidth}
\vspace{0pt}
Cho hình bình hành  ${QKHG}$ có ${F, E}$ lần lượt là trung điểm của ${QK, HG }$. Đường ${QE}$ và ${HF}$ cắt ${KG}$ lần lượt tại ${I, J }$.  Từ ${J}$ kẻ đường thẳng song song với ${QK}$ cắt ${QE}$ tại ${O}$. Trong các khẳng định sau khẳng định nào sai?

\vspace{4pt}\noindent
\textbf{A.} $KJ=IJ=IG$ \quad \textbf{B.} ${ QFHE }$ là hình bình hành \quad \textbf{\textcolor{blue!75!black}{C. (Đ)}} ${ KJOQ }$ là hình thang cân \quad \textbf{D.} $\Delta {IJO}=\Delta {IGE }$
\end{minipage}\hfill
\begin{minipage}[t]{0.39\textwidth}
\vspace{0pt}
\centering

\resizebox{\linewidth}{!}{%
\begin{tikzpicture}
            \coordinate (A) at (-3, 0);
            \coordinate (D) at (2, 0);
            \coordinate (C) at (3, 2);
            \coordinate (B) at (-2, 2);
           \coordinate (E) at (0.67, 0.67);
            \coordinate (M) at (-2.5, 1);
             \coordinate (N) at (2.5, 1);
              \coordinate (F) at (-0.67, 1.33);
              \coordinate (G) at (-1.17, 0.33);
            % Vẽ tứ giác
            \draw[thick] (M) -- (C);
            \draw[thick] (A) -- (B) -- (C) -- (D) -- (A);
        \draw[thick] (A) -- (N);
        \draw[thick] (B) -- (D);
         \draw[thick] (F) -- (G);
            % Ghi nhãn các điểm
            \node[above left] at (M) {$F$};
            \node[above left] at (B) {$K$};
            \node[above right] at (C) {$H$};
            \node[below right] at (D) {$G$};
            \node[above left] at (A){$Q$};
            \node[above right] at (N) {$E$};
            \node[below] at (E) {$I$};
             \node[above] at (F) {$J$};
                \node[above right] at (G) {$O$};
    \end{tikzpicture}
}
\end{minipage}

\vspace{3pt}\hrule\vspace{3pt}\footnotesize\textcolor{gray!80!black}{\textbf{Các mã dùng chung hình (HINH\_159):} \texttt{L08\_C3\_B3\_16}}
\end{tcolorbox}
\vspace{4pt}


\begin{tcolorbox}[
    breakable,
    colback=white,
    colframe=blue!60!black,
    boxrule=0.8pt,
    arc=3pt,
    title={\textbf{Câu \#199} \quad \textbf{MÃ CÂU: \texttt{[L08\_C3\_B3\_09]}} \quad | \quad \textbf{Nhóm: \texttt{HINH\_090}} \quad | \quad \textbf{MapID: \texttt{8H3N3}} \hfill \textcolor{yellow!90!white}{[Hình gốc .pyc]}}
]
\textbf{Dạng toán:} Nhận biết về hình bình hành \hfill \textit{(Tự luận -- Nhận biết)}
\vspace{4pt}\hrule\vspace{6pt}
\begin{minipage}[t]{0.58\textwidth}
\vspace{0pt}
Cho hình bình hành ${PQRN}$ như hình vẽ. Biết $\widehat{ QPN } =74^{\circ}$ và ${O}$ là trung điểm của ${QN}$.

a) Tính số đo các góc còn lại của hình bình hành.

 b) Chứng minh ${O}$, ${P}$, ${R}$ thẳng hàng
\end{minipage}\hfill
\begin{minipage}[t]{0.39\textwidth}
\vspace{0pt}
\centering

\resizebox{\linewidth}{!}{%
\begin{tikzpicture}
            \coordinate (O) at (3/2, 3/2);
            \coordinate (D) at (2, 0);
            \coordinate (C) at (5, 3);
            \coordinate (B) at (1, 3);
            \coordinate (A) at (-2, 0);
            \
            % Vẽ tứ giác
            \draw[thick] (O) -- (A) -- (C) ;
            \draw[thick] (A) -- (B)--(C)--(D)--(A);
        \draw[thick] (O) -- (B) -- (D) ;
            % Ghi nhãn các điểm
            \node[above right] at (O) {$O$};
            \node[above left] at (B) {$Q$};
            \node[above right] at (C) {$R$};
            \node[below right] at (D) {$N$};
            \node[below left] at (A){$P$};
                \draw
            ($(B)!0.4!(O)$)node[rotate=80]{$/$} ($(O)!0.4!(D)$)node[rotate=80]{$/$};
    \end{tikzpicture}
}
\end{minipage}

\vspace{3pt}\hrule\vspace{3pt}\footnotesize\textcolor{gray!80!black}{\textbf{Các mã dùng chung hình (HINH\_090):} \texttt{L07\_C3\_B1\_07, L08\_C3\_B3\_01}}
\end{tcolorbox}
\vspace{4pt}


\begin{tcolorbox}[
    breakable,
    colback=white,
    colframe=blue!60!black,
    boxrule=0.8pt,
    arc=3pt,
    title={\textbf{Câu \#200} \quad \textbf{MÃ CÂU: \texttt{[L08\_C3\_B3\_10]}} \quad | \quad \textbf{Nhóm: \texttt{HINH\_153}} \quad | \quad \textbf{MapID: \texttt{8H3H3}} \hfill \textcolor{yellow!90!white}{[Hình gốc .pyc]}}
]
\textbf{Dạng toán:} Thông hiểu về hình bình hành \hfill \textit{(Tự luận -- Thông hiểu)}
\vspace{4pt}\hrule\vspace{6pt}
\begin{minipage}[t]{0.58\textwidth}
\vspace{0pt}
Cho $\Delta{EFK}$ cân tại ${E}$ có điểm ${H}$ trên cạnh ${FK}$. Kẻ $HI//EK$, $HJ//EF$. 

a) Chứng minh: Tứ giác ${EIHJ}$ là hình bình hành.

b) Chứng minh: $HI+HJ=FI+EI=EF$
\end{minipage}\hfill
\begin{minipage}[t]{0.39\textwidth}
\vspace{0pt}
\centering

\resizebox{\linewidth}{!}{%
\begin{tikzpicture}
            % Định nghĩa các điểm
            \coordinate (A) at (0, 4);
            \coordinate (B) at (-2, 0);
            \coordinate (C) at (2, 0);
            \coordinate (D) at (1, 0);
            \coordinate (M) at (-0.5, 3);
                \coordinate (N) at (1.5, 1);
            % Vẽ tứ giác
            \draw[thick] (A) -- (B) -- (C) -- cycle;
            \draw[thick] (D) -- (M);
            \draw[thick] (D) -- (N);
            % Ghi nhãn các điểm
            \node[above] at (A) {$E $};
            \node[above left] at (B) {$F $};
            \node[above right] at (C) {$K $};
            \node[below] at (D) {$H $};
            \node[above left] at (M){$I $};
            \node[above right] at (N){$J $};
        \end{tikzpicture}
}
\end{minipage}

\vspace{3pt}\hrule\vspace{3pt}\footnotesize\textcolor{gray!80!black}{\textbf{Các mã dùng chung hình (HINH\_153):} \texttt{L08\_C3\_B3\_02}}
\end{tcolorbox}
\vspace{4pt}


\begin{tcolorbox}[
    breakable,
    colback=white,
    colframe=blue!60!black,
    boxrule=0.8pt,
    arc=3pt,
    title={\textbf{Câu \#201} \quad \textbf{MÃ CÂU: \texttt{[L08\_C3\_B3\_11]}} \quad | \quad \textbf{Nhóm: \texttt{HINH\_154}} \quad | \quad \textbf{MapID: \texttt{8H3V3}} \hfill \textcolor{yellow!90!white}{[Hình gốc .pyc]}}
]
\textbf{Dạng toán:} Bài tổng hợp 1 \hfill \textit{(Tự luận -- Vận dụng)}
\vspace{4pt}\hrule\vspace{6pt}
\begin{minipage}[t]{0.58\textwidth}
\vspace{0pt}
Cho hình bình hành ${GHFE}$ có $GH < FE$. Tia phân giác của góc $\widehat{H}$, $\widehat{E}$ lần lượt cắt ${GE, HF}$ tại ${I, J}$. 

a) Chứng minh $\Delta{GHI}$ cân tại ${G}$.  

b) Chứng minh tứ giác ${HIEJ}$ là hình bình hành.
\end{minipage}\hfill
\begin{minipage}[t]{0.39\textwidth}
\vspace{0pt}
\centering

\resizebox{\linewidth}{!}{%
\begin{tikzpicture}
            \coordinate (M) at (-0.76, 0);
            \coordinate (D) at (2, 0);
            \coordinate (C) at (3, 2);
            \coordinate (B) at (-2, 2);
            \coordinate (A) at (-3, 0);
            \coordinate (N) at (0.76, 2);
            % Vẽ tứ giác
            \draw[thick] (B) -- (M) ;
            \draw[thick] (A) -- (B)--(C)--(D)--(A);
        \draw[thick] (N) -- (D) ;
            % Ghi nhãn các điểm
            \node[below right] at (M) {$I$};
            \node[above left] at (B) {$H$};
            \node[above right] at (C) {$F$};
            \node[below right] at (D) {$E$};
            \node[below left] at (A){$G$};
            \node[above right] at (N) {$J$};
                \draw 
            pic[draw,blue,angle radius=4mm] {angle = M--B--C}; 
             \draw  pic[draw,blue,angle radius=4.7mm] {angle = A--B--M};
             \draw  pic[draw,blue,angle radius=3mm] {angle = C--D--N} ;
            \draw pic[draw,blue,angle radius=3.7mm] {angle = N--D--M};
    \end{tikzpicture}
}
\end{minipage}

\vspace{3pt}\hrule\vspace{3pt}\footnotesize\textcolor{gray!80!black}{\textbf{Các mã dùng chung hình (HINH\_154):} \texttt{L08\_C3\_B3\_03}}
\end{tcolorbox}
\vspace{4pt}


\begin{tcolorbox}[
    breakable,
    colback=white,
    colframe=blue!60!black,
    boxrule=0.8pt,
    arc=3pt,
    title={\textbf{Câu \#202} \quad \textbf{MÃ CÂU: \texttt{[L08\_C3\_B3\_12]}} \quad | \quad \textbf{Nhóm: \texttt{HINH\_155}} \quad | \quad \textbf{MapID: \texttt{8H3H3}} \hfill \textcolor{yellow!90!white}{[Hình gốc .pyc]}}
]
\textbf{Dạng toán:} Thông hiểu hình bình hành \hfill \textit{(Tự luận -- Thông hiểu)}
\vspace{4pt}\hrule\vspace{6pt}
\begin{minipage}[t]{0.58\textwidth}
\vspace{0pt}
Cho hình bình hành ${MNPQ}$. Trên cạnh ${MN}$ lấy điểm ${E}$ trên cạnh ${PQ}$ lấy điểm ${F}$ sao cho $ME=PF$.

a) Chứng minh ${MEPF}$ là hình bình hành.

 b) Chứng minh ${QENF}$ là hình bình hành
\end{minipage}\hfill
\begin{minipage}[t]{0.39\textwidth}
\vspace{0pt}
\centering

\resizebox{\linewidth}{!}{%
\begin{tikzpicture}
            \coordinate (M) at (-2.34, 1.31);
            \coordinate (D) at (2, 0);
            \coordinate (C) at (3, 2);
            \coordinate (B) at (-2, 2);
            \coordinate (A) at (-3, 0);
            \coordinate (N) at (2.34, 0.69);
            % Vẽ tứ giác
            \draw[thick] (D) -- (M) ;
            \draw[thick] (A) -- (B)--(C)--(D)--(A);
        \draw[thick] (N) -- (B) ;
            % Ghi nhãn các điểm
            \node[above left] at (M) {$E$};
            \node[above left] at (B) {$N$};
            \node[above right] at (C) {$P$};
            \node[below right] at (D) {$Q$};
            \node[below left] at (A){$M$};
            \node[below right] at (N) {$F$};
            \draw
        ($(A)!0.5!(M)$)node[rotate=20]{$/$} ($(N)!0.5!(C)$)node[rotate=20]{$/$};
    \end{tikzpicture}
}
\end{minipage}

\vspace{3pt}\hrule\vspace{3pt}\footnotesize\textcolor{gray!80!black}{\textbf{Các mã dùng chung hình (HINH\_155):} \texttt{L08\_C3\_B3\_04}}
\end{tcolorbox}
\vspace{4pt}


\begin{tcolorbox}[
    breakable,
    colback=white,
    colframe=blue!60!black,
    boxrule=0.8pt,
    arc=3pt,
    title={\textbf{Câu \#203} \quad \textbf{MÃ CÂU: \texttt{[L08\_C3\_B3\_13]}} \quad | \quad \textbf{Nhóm: \texttt{HINH\_156}} \quad | \quad \textbf{MapID: \texttt{8H3H3}} \hfill \textcolor{yellow!90!white}{[Hình gốc .pyc]}}
]
\textbf{Dạng toán:} Thông hiểu hình bình hành 3 \hfill \textit{(Tự luận -- Thông hiểu)}
\vspace{4pt}\hrule\vspace{6pt}
\begin{minipage}[t]{0.58\textwidth}
\vspace{0pt}
Cho hình bình hành ${EFKH}$ có $EF>KH$. Từ ${E}$ kẻ đường thẳng vuông góc với ${FH}$ cắt ${HK}$ tại ${I}$, từ ${K}$ vẽ đường thẳng vuông góc với ${FH}$ cắt ${EF}$ tại ${J}$. Gọi ${O}$ là trung điểm của ${FH}$. 

a) Chứng minh ${EIKJ}$ là hình bình hành.

b)  Chứng minh ${O}$ là trung điểm của ${IJ}$
\end{minipage}\hfill
\begin{minipage}[t]{0.39\textwidth}
\vspace{0pt}
\centering

\resizebox{\linewidth}{!}{%
\begin{tikzpicture}
            \coordinate (O) at (0, 1);
            \coordinate (D) at (-3, 0);
            \coordinate (C) at (2, 0);
            \coordinate (B) at (3, 2);
            \coordinate (A) at (-2, 2);
            \coordinate (H) at (-1.33, 0);
            \coordinate (K) at (1.33, 2);
            % Vẽ tứ giác
            \draw[thick] (D) -- (B);
            \draw[thick] (A) -- (B) -- (C) -- (D) -- (A);
            \draw[thick] (C) -- (K);
            \draw[thick] (A) -- (H);
            \draw[thick] (K) -- (H);
            % Ghi nhãn các điểm
            \node[below left] at (H) {$I$};
            \node[above right] at (B) {$F$};
            \node[below right] at (C) {$K$};
            \node[below left] at (D) {$H$};
            \node[above left] at (A){$E$};
            \node[above right] at (K) {$J$};
            \node[below] at (O) {$O$};
    \end{tikzpicture}
}
\end{minipage}

\vspace{3pt}\hrule\vspace{3pt}\footnotesize\textcolor{gray!80!black}{\textbf{Các mã dùng chung hình (HINH\_156):} \texttt{L08\_C3\_B3\_05}}
\end{tcolorbox}
\vspace{4pt}


\begin{tcolorbox}[
    breakable,
    colback=white,
    colframe=blue!60!black,
    boxrule=0.8pt,
    arc=3pt,
    title={\textbf{Câu \#204} \quad \textbf{MÃ CÂU: \texttt{[L08\_C3\_B3\_14]}} \quad | \quad \textbf{Nhóm: \texttt{HINH\_157}} \quad | \quad \textbf{MapID: \texttt{8H3V3}} \hfill \textcolor{yellow!90!white}{[Hình gốc .pyc]}}
]
\textbf{Dạng toán:} Tổng hợp hình bình hành 1 \hfill \textit{(Tự luận -- Vận dụng)}
\vspace{4pt}\hrule\vspace{6pt}
\begin{minipage}[t]{0.58\textwidth}
\vspace{0pt}
Cho tam giác ${ABC}$ cân tại ${A}$ lấy điểm ${D}$ bất kỳ trên ${AB}$, lấy điểm ${K}$ trên tia đối của tia ${CA}$ sao cho $CK=BD$. Từ ${D}$ kẻ đường thẳng song song với ${AC}$ cắt ${BC}$ tại ${H}$. 

a) Chứng minh: $\Delta{DBH}$ là tam giác cân.

 b) Chứng minh: tứ giác ${DCKH}$ là hình bình hành.
\end{minipage}\hfill
\begin{minipage}[t]{0.39\textwidth}
\vspace{0pt}
\centering

\resizebox{\linewidth}{!}{%
\begin{tikzpicture}
            \coordinate (O) at (0, 1);
            \coordinate (D) at (-1.2, 1.99);
            \coordinate (C) at (2, 0);
            \coordinate (B) at (-2, 0);
            \coordinate (A) at (0, 5);
            \coordinate (F) at (-0.41, 0);
            \coordinate (E) at (2.9, -1.99);
            % Vẽ tứ giác
            \draw[thick] (D) -- (C)--(E)--(F)--(D);
            \draw[thick] (A) -- (B) -- (C)-- (A);
         
            % Ghi nhãn các điểm
            \node[below right] at (E) {$K$};
            \node[below left] at (B) {$B$};
            \node[above right] at (C) {$C$};
            \node[above left] at (D) {$D$};
            \node[above] at (A){$A$};
            \node[below left] at (F) {$H$};
           \draw
        ($(B)!0.5!(D)$)node[rotate=-40]{$/$};\draw ($(C)!0.5!(E)$)node[rotate=-40]{$/$};
    \end{tikzpicture}
}
\end{minipage}

\vspace{3pt}\hrule\vspace{3pt}\footnotesize\textcolor{gray!80!black}{\textbf{Các mã dùng chung hình (HINH\_157):} \texttt{L08\_C3\_B3\_06}}
\end{tcolorbox}
\vspace{4pt}


\begin{tcolorbox}[
    breakable,
    colback=white,
    colframe=blue!60!black,
    boxrule=0.8pt,
    arc=3pt,
    title={\textbf{Câu \#205} \quad \textbf{MÃ CÂU: \texttt{[L08\_C3\_B3\_15]}} \quad | \quad \textbf{Nhóm: \texttt{HINH\_158}} \quad | \quad \textbf{MapID: \texttt{8H3V3}} \hfill \textcolor{yellow!90!white}{[Hình gốc .pyc]}}
]
\textbf{Dạng toán:} Tổng hợp hình bình hành 2 \hfill \textit{(Tự luận -- Vận dụng)}
\vspace{4pt}\hrule\vspace{6pt}
\begin{minipage}[t]{0.58\textwidth}
\vspace{0pt}
Cho hình bình hành  ${PQRT}$. Có ${M, N, E, F}$ lần lượt là trung điểm của ${ PQ, QR, RT, PT }$.

a) Chứng minh $MN=EF$. 

b) Chứng minh ${MNEF }$ là hình bình hành.
\end{minipage}\hfill
\begin{minipage}[t]{0.39\textwidth}
\vspace{0pt}
\centering

\resizebox{\linewidth}{!}{%
\begin{tikzpicture}
            \coordinate (P) at (2.5, 1);
            \coordinate (D) at (-3, 0);
            \coordinate (C) at (2, 0);
            \coordinate (B) at (3, 2);
            \coordinate (A) at (-2, 2);
            \coordinate (M) at (-5/2, 1);
            \coordinate (N) at (1/2, 2);
             \coordinate (Q) at (-1/2, 0);
            % Vẽ tứ giác
            \draw[thick] (M) -- (N)--(P) -- (Q)--(M);
            \draw[thick] (A) -- (B) -- (C) -- (D) -- (A);
        
            % Ghi nhãn các điểm
            \node[above left] at (M) {$M$};
            \node[above right] at (B) {$Q$};
            \node[below right] at (C) {$R$};
            \node[below left] at (D) {$T$};
            \node[above left] at (A){$P$};
            \node[above] at (N) {$N$};
            \node[above right] at (P) {$E$};
             \node[below] at (Q) {$F$};
    \end{tikzpicture}
}
\end{minipage}

\vspace{3pt}\hrule\vspace{3pt}\footnotesize\textcolor{gray!80!black}{\textbf{Các mã dùng chung hình (HINH\_158):} \texttt{L08\_C3\_B3\_07}}
\end{tcolorbox}
\vspace{4pt}


\begin{tcolorbox}[
    breakable,
    colback=white,
    colframe=blue!60!black,
    boxrule=0.8pt,
    arc=3pt,
    title={\textbf{Câu \#206} \quad \textbf{MÃ CÂU: \texttt{[L08\_C3\_B3\_16]}} \quad | \quad \textbf{Nhóm: \texttt{HINH\_159}} \quad | \quad \textbf{MapID: \texttt{8H3C3}} \hfill \textcolor{yellow!90!white}{[Hình gốc .pyc]}}
]
\textbf{Dạng toán:} Bài toán về hình bình hành mức vận dụng \hfill \textit{(Tự luận -- Vận dụng cao)}
\vspace{4pt}\hrule\vspace{6pt}
\begin{minipage}[t]{0.58\textwidth}
\vspace{0pt}
Cho hình bình hành  ${PQRT}$ có ${M, N}$ lần lượt là trung điểm của ${PQ, RT }$. Đường ${PN}$ và ${RM}$ cắt ${QT}$ lần lượt tại ${E, F }$.

 a) Chứng minh rằng ${ PMRN }$ là hình bình hành,

 b) Từ ${F}$ kẻ đường thẳng song song với ${PQ}$ cắt ${PN}$ tại ${G}$. Chứng minh $QF=EF=ET$
\end{minipage}\hfill
\begin{minipage}[t]{0.39\textwidth}
\vspace{0pt}
\centering

\resizebox{\linewidth}{!}{%
\begin{tikzpicture}
            \coordinate (A) at (-3, 0);
            \coordinate (D) at (2, 0);
            \coordinate (C) at (3, 2);
            \coordinate (B) at (-2, 2);
           \coordinate (E) at (0.67, 0.67);
            \coordinate (M) at (-2.5, 1);
             \coordinate (N) at (2.5, 1);
              \coordinate (F) at (-0.67, 1.33);
              \coordinate (G) at (-1.17, 0.33);
            % Vẽ tứ giác
            \draw[thick] (M) -- (C);
            \draw[thick] (A) -- (B) -- (C) -- (D) -- (A);
        \draw[thick] (A) -- (N);
        \draw[thick] (B) -- (D);
         \draw[thick] (F) -- (G);
            % Ghi nhãn các điểm
            \node[above left] at (M) {$M$};
            \node[above left] at (B) {$Q$};
            \node[above right] at (C) {$R$};
            \node[below right] at (D) {$T$};
            \node[above left] at (A){$P$};
            \node[above right] at (N) {$N$};
            \node[below] at (E) {$E$};
             \node[above] at (F) {$F$};
                \node[above right] at (G) {$G$};
    \end{tikzpicture}
}
\end{minipage}

\vspace{3pt}\hrule\vspace{3pt}\footnotesize\textcolor{gray!80!black}{\textbf{Các mã dùng chung hình (HINH\_159):} \texttt{L08\_C3\_B3\_08}}
\end{tcolorbox}
\vspace{4pt}


\begin{tcolorbox}[
    breakable,
    colback=white,
    colframe=blue!60!black,
    boxrule=0.8pt,
    arc=3pt,
    title={\textbf{Câu \#207} \quad \textbf{MÃ CÂU: \texttt{[L08\_C3\_B4\_02]}} \quad | \quad \textbf{Nhóm: \texttt{HINH\_160}} \quad | \quad \textbf{MapID: \texttt{8H3N4}} \hfill \textcolor{yellow!90!white}{[Hình gốc .pyc]}}
]
\textbf{Dạng toán:} Nhận biết hình chữ nhật. \hfill \textit{(Tự luận -- Nhận biết)}
\vspace{4pt}\hrule\vspace{6pt}
\begin{minipage}[t]{0.58\textwidth}
\vspace{0pt}
Cho tam giác ${IJQ}$ vuông tại ${I}$ có ${M}$ là trung điểm ${JQ}$. Từ ${M}$ kẻ $ME//IQ \left (E \in IJ \right), MK//IJ \left(K \in IQ \right)$ 

a) Tứ giác ${JEKM, IEMK}$ là hình gì?

b) Gọi ${O}$ là trung điểm của ${IM}$, chứng minh $OE=OK$.
\end{minipage}\hfill
\begin{minipage}[t]{0.39\textwidth}
\vspace{0pt}
\centering

\resizebox{\linewidth}{!}{%
\begin{tikzpicture}
            \coordinate (A) at (-2, 0);
            \coordinate (C) at (-2, 4);
            \coordinate (B) at (1, 0);
           \coordinate (E) at (-0.5, 0);
            \coordinate (M) at (-0.5, 2);
              \coordinate (F) at (-2, 2);
             \coordinate (O) at (-1.25, 1);
            % Vẽ tứ giác
            \draw[thick] (M) -- (E);
             \draw[thick] (M) -- (F);
              \draw[thick] (M) -- (A);
               \draw[thick] (E) -- (F);
            \draw[thick] (A) -- (B) -- (C) -- (A);
        
            % Ghi nhãn các điểm
            \node[above right] at (M) {$M$};
            \node[below right] at (B) {$J$};
            \node[above] at (C) {$Q$};
            \node[above] at (O) {$O$};
            \node[below left] at (A){$I$};
            \node[below] at (E) {$E$};
             \node[above left ] at (F) {$K$};
             \def\khvuong[size=#1](#2,#3,#4){%
            \draw ($(#3)!#1!(#2)$) --
            ($($(#3)!#1!(#2)$)+($(#3)!#1!(#4)$)-(#3)$)--
            ($(#3)!#1!(#4)$);
            }
            \khvuong[size=6pt](B,A,C)
    \end{tikzpicture}
}
\end{minipage}

\end{tcolorbox}
\vspace{4pt}


\begin{tcolorbox}[
    breakable,
    colback=white,
    colframe=blue!60!black,
    boxrule=0.8pt,
    arc=3pt,
    title={\textbf{Câu \#208} \quad \textbf{MÃ CÂU: \texttt{[L08\_C3\_B4\_03]}} \quad | \quad \textbf{Nhóm: \texttt{HINH\_161}} \quad | \quad \textbf{MapID: \texttt{8H3N4}} \hfill \textcolor{yellow!90!white}{[Hình gốc .pyc]}}
]
\textbf{Dạng toán:} Nhận biết hình chữ nhật 2. \hfill \textit{(Tự luận -- Nhận biết)}
\vspace{4pt}\hrule\vspace{6pt}
\begin{minipage}[t]{0.58\textwidth}
\vspace{0pt}
Cho tam giác ${ABC}$vuông tại ${A}$, đường cao ${AQ}$, trung tuyến ${AE}$. Qua ${Q}$ kẻ $QH//AC \left(H \in AB \right)$ và $QK//AB \left( K \in AC \right )$. Đoạn ${HK}$ cắt ${AQ, AE}$ lần lượt tại ${I}$ và ${D}$.

 a) Chứng minh rằng $AQ=HK$ 

b) $\Delta {EAC}$ là tam giác gì? 

c) Chứng minh $\Delta {AKD}$ là tam giác vuông.
\end{minipage}\hfill
\begin{minipage}[t]{0.39\textwidth}
\vspace{0pt}
\centering

\resizebox{\linewidth}{!}{%
\begin{tikzpicture}
            \coordinate (A) at (0, 4);
            \coordinate (C) at (8, 0);
            \coordinate (B) at (-2, 0);
           \coordinate (D) at (-1.6, 0.8);
            \coordinate (M) at (3, 0);
              \coordinate (P) at (1.6, 3.2);
             \coordinate (H) at (0, 0);
               \coordinate (N) at (0.96, 2.72);
                 \coordinate (O) at (0, 2);
            % Vẽ tứ giác
             \draw[thick] (A) -- (H);
              \draw[thick] (M) -- (A);
               \draw[thick] (D) -- (P);
            \draw[thick] (A) -- (B) -- (C) -- (A);
         \draw[thick] (D) -- (H);
          \draw[thick] (P) -- (H);
            % Ghi nhãn các điểm
            \node[below] at (M) {$E$};
            \node[below left] at (B) {$B$};
            \node[below right] at (C) {$C$};
            \node[below right] at (O) {$I$};
            \node[above] at (A){$A$};
            \node[above] at (N) {$D$};
             \node[below ] at (H) {$Q$};
                 \node[above right] at (P) {$K$};
                  \node[above left] at (D) {$H$};
             \def\khvuong[size=#1](#2,#3,#4){%
            \draw ($(#3)!#1!(#2)$) --
            ($($(#3)!#1!(#2)$)+($(#3)!#1!(#4)$)-(#3)$)--
            ($(#3)!#1!(#4)$);
            }
            \khvuong[size=6pt](B,A,C);
             \khvuong[size=6pt](A,H,B);
    \end{tikzpicture}
}
\end{minipage}

\end{tcolorbox}
\vspace{4pt}


\begin{tcolorbox}[
    breakable,
    colback=white,
    colframe=blue!60!black,
    boxrule=0.8pt,
    arc=3pt,
    title={\textbf{Câu \#209} \quad \textbf{MÃ CÂU: \texttt{[L08\_C3\_B4\_04]}} \quad | \quad \textbf{Nhóm: \texttt{HINH\_162}} \quad | \quad \textbf{MapID: \texttt{8H3N4}} \hfill \textcolor{yellow!90!white}{[Hình gốc .pyc]}}
]
\textbf{Dạng toán:} Nhận biết hình chữ nhật 4. \hfill \textit{(Tự luận -- Nhận biết)}
\vspace{4pt}\hrule\vspace{6pt}
\begin{minipage}[t]{0.58\textwidth}
\vspace{0pt}
Cho tam giác ${GHF}$ vuông tại ${G}$ có $GH<GF$. Gọi ${I}$ là trung điểm của $ {HF}$. Trên tia đối của tia ${IG}$ lấy điểm ${M}$ sao cho $IM=IG$.

 a) Chứng minh rằng ${GHFM}$ là hình chữ nhật.

 b) Lấy điểm ${E}$ sao cho ${H}$ là trung điểm của ${GE}$. Chứng minh ${HEMF}$ là hình bình hành.

c) ${EI}$ cắt ${HM}$ tại ${P}$. Chứng minh ${EP=2PI}$.
\end{minipage}\hfill
\begin{minipage}[t]{0.39\textwidth}
\vspace{0pt}
\centering

\resizebox{\linewidth}{!}{%
\begin{tikzpicture}
            \coordinate (A) at (0, 4);
            \coordinate (C) at (8, 0);
            \coordinate (B) at (-2, 0);
           \coordinate (D) at (6, -4);
            \coordinate (M) at (3, 0);
              \coordinate (E) at (-4, -4);
             \coordinate (K) at (0.67, -1.33);
            % Vẽ tứ giác
             \draw[thick] (E) -- (M);
              \draw[thick] (D) -- (B);
               \draw[thick] (B) -- (E)--(D)--(M);
            \draw[thick] (A) -- (B) -- (C) -- (A);
         \draw[thick] (A) -- (D);
          \draw[thick] (D) -- (C);
            % Ghi nhãn các điểm
            \node[above right] at (M) {$I$};
            \node[below left] at (B) {$H$};
            \node[below right] at (C) {$F$};
            \node[below left] at (E) {$E$};
            \node[above] at (A) {$G$};
                 \node[below] at (K) {$P$};
                  \node[below right] at (D) {$M$};
             \def\khvuong[size=#1](#2,#3,#4){%
            \draw ($(#3)!#1!(#2)$) --
            ($($(#3)!#1!(#2)$)+($(#3)!#1!(#4)$)-(#3)$)--
            ($(#3)!#1!(#4)$);
            }
            \khvuong[size=6pt](B,A,C);
    \end{tikzpicture}
}
\end{minipage}

\end{tcolorbox}
\vspace{4pt}


\begin{tcolorbox}[
    breakable,
    colback=white,
    colframe=blue!60!black,
    boxrule=0.8pt,
    arc=3pt,
    title={\textbf{Câu \#210} \quad \textbf{MÃ CÂU: \texttt{[L08\_C3\_B4\_05]}} \quad | \quad \textbf{Nhóm: \texttt{HINH\_163}} \quad | \quad \textbf{MapID: \texttt{8H3H4}} \hfill \textcolor{yellow!90!white}{[Hình gốc .pyc]}}
]
\textbf{Dạng toán:} Thông hiểu hình chữ nhật 1. \hfill \textit{(Tự luận -- Thông hiểu)}
\vspace{4pt}\hrule\vspace{6pt}
\begin{minipage}[t]{0.58\textwidth}
\vspace{0pt}
Cho tam giác ${A B C }$ vuông tại ${A }$ có $A B <A C $. Cho ${N }$ là trung điểm của ${B C }$. Gọi ${M ,P }$ lần lượt là hình chiếu của ${N }$ lên ${A B , A C  }$. Lấy ${E }$ sao cho ${P }$ là trung điểm của ${N E }$.

a) Chứng minh $ {M , P }$ lần lượt là trung điểm của ${A B , A C  }$ 

b) Tứ giác ${A N C E }$ là hình gì?
\end{minipage}\hfill
\begin{minipage}[t]{0.39\textwidth}
\vspace{0pt}
\centering

\resizebox{\linewidth}{!}{%
\begin{tikzpicture}
            \coordinate (A) at (0, 4);
            \coordinate (C) at (8, 0);
            \coordinate (B) at (-2, 0);
         \coordinate (N) at (3, 0);
            \coordinate (M) at (-1, 2);
              \coordinate (E) at (5, 4);
             \coordinate (P) at (4, 2);
            % Vẽ tứ giác
             \draw[thick] (M) -- (N)--(E);
              \draw[thick] (A) -- (N);
            \draw[thick] (A) -- (B) -- (C) -- (A);
             \draw[thick] (C) -- (E);
             \draw[thick] (A) -- (E);
            % Ghi nhãn các điểm
            \node[above left] at (M) {$M $};
            \node[below left] at (B) {$B $};
            \node[below right] at (C) {$C $};
            \node[above right] at (E) {$E $};
            \node[above] at (A) {$A $};
                 \node[below] at (N) {$N $};
                  \node[below right] at (P) {$P $};
             \def\khvuong[size=#1](#2,#3,#4){%
            \draw ($(#3)!#1!(#2)$) --
            ($($(#3)!#1!(#2)$)+($(#3)!#1!(#4)$)-(#3)$)--
            ($(#3)!#1!(#4)$);
            }
            \khvuong[size=6pt](B,A,C);
             \khvuong[size=6pt](A,P,N);
              \khvuong[size=6pt](A,M,N);
                \draw
        ($(B)!0.5!(N)$)node[rotate=40]{$/$} ($(N)!0.5!(C)$)node[rotate=40]{$/$}
        ($(N)!0.5!(P)$)node[rotate=90]{$//$} ($(P)!0.5!(E)$)node[rotate=90]{$//$};
    \end{tikzpicture}
}
\end{minipage}

\end{tcolorbox}
\vspace{4pt}


\begin{tcolorbox}[
    breakable,
    colback=white,
    colframe=blue!60!black,
    boxrule=0.8pt,
    arc=3pt,
    title={\textbf{Câu \#211} \quad \textbf{MÃ CÂU: \texttt{[L08\_C3\_B4\_06]}} \quad | \quad \textbf{Nhóm: \texttt{HINH\_164}} \quad | \quad \textbf{MapID: \texttt{8H3H4}} \hfill \textcolor{yellow!90!white}{[Hình gốc .pyc]}}
]
\textbf{Dạng toán:} Thông hiểu hình chữ nhật 2. \hfill \textit{(Tự luận -- Thông hiểu)}
\vspace{4pt}\hrule\vspace{6pt}
\begin{minipage}[t]{0.58\textwidth}
\vspace{0pt}
Cho tam giác ${ABC}$ vuông tại ${A}$, có ${AD}$ là đường cao. Gọi ${F}$ và ${H}$ lần lượt là hình chiếu của ${D}$ xuống ${ AB, AC}$.  Gọi ${G}$ là trung điểm của ${DB}$, ${M}$ là trung điểm của ${DC}$, ${AD}$ cắt ${FH}$ tại ${E}$a) Tứ giác ${ AFDH}$ là hình gì? 

b) Chứng minh $\Delta{ MHD}$ là tam giác cân.

c) Chứng minh $\widehat { MHF}=90^{\circ}$ và $FG//HM$
\end{minipage}\hfill
\begin{minipage}[t]{0.39\textwidth}
\vspace{0pt}
\centering

\resizebox{\linewidth}{!}{%
\begin{tikzpicture}
            \coordinate (A) at (0, 4);
            \coordinate (C) at (8, 0);
            \coordinate (B) at (-2, 0);
         \coordinate (H) at (0, 0);
            \coordinate (P) at (-1.6, 0.8);
              \coordinate (Q) at (1.6, 3.2);
             \coordinate (I) at (-1, 0);
                \coordinate (K) at (4, 0);
                   \coordinate (O) at (0, 2);
            % Vẽ tứ giác
             \draw[thick] (A) -- (H);
              \draw[thick] (P) -- (Q);
            \draw[thick] (A) -- (B) -- (C) -- (A);
             \draw[thick] (P) -- (H);
             \draw[thick] (P) -- (I);
             \draw[thick] (H) -- (Q);
               \draw[thick] (K) -- (Q);
    % Ghi nhãn các điểm
            
            \node[below left] at (B) {$B$};
            \node[below right] at (C) {$C$};
            \node[above right] at (Q) {$H$};
            \node[above] at (A) {$A$};
                 \node[below] at (I) {$G$};
                  \node[above left] at (P) {$F$};
                  \node[below] at (K) {$M$};
                  \node[below] at (H) {$D$};
                  \node[above left] at (O) {$E$};
             \def\khvuong[size=#1](#2,#3,#4){%
            \draw ($(#3)!#1!(#2)$) --
            ($($(#3)!#1!(#2)$)+($(#3)!#1!(#4)$)-(#3)$)--
            ($(#3)!#1!(#4)$);
            }
            \khvuong[size=6pt](B,A,C);
             \khvuong[size=6pt](A,H,B);
              \khvuong[size=6pt](A,P,H);
                  \khvuong[size=6pt](A,Q,H);
                \draw
        ($(B)!0.5!(I)$)node[rotate=40]{$/$} ($(I)!0.5!(H)$)node[rotate=40]{$/$}
        ($(H)!0.5!(K)$)node[rotate=50]{$//$} ($(K)!0.5!(C)$)node[rotate=50]{$//$};
    \end{tikzpicture}
}
\end{minipage}

\end{tcolorbox}
\vspace{4pt}


\begin{tcolorbox}[
    breakable,
    colback=white,
    colframe=blue!60!black,
    boxrule=0.8pt,
    arc=3pt,
    title={\textbf{Câu \#212} \quad \textbf{MÃ CÂU: \texttt{[L08\_C3\_B4\_07]}} \quad | \quad \textbf{Nhóm: \texttt{HINH\_165}} \quad | \quad \textbf{MapID: \texttt{8H3H4}} \hfill \textcolor{yellow!90!white}{[Hình gốc .pyc]}}
]
\textbf{Dạng toán:} Thông hiểu hình chữ nhật 3. \hfill \textit{(Tự luận -- Thông hiểu)}
\vspace{4pt}\hrule\vspace{6pt}
\begin{minipage}[t]{0.58\textwidth}
\vspace{0pt}
Cho tam giác $\Delta{PQR}$ vuông tại ${P}$, gọi ${M}$ là trung điểm của ${QR}$. Gọi ${E, T}$ lần lượt là chân đường vuông góc kẻ từ ${M}$ đến ${PQ, PR}$. Gọi ${J}, {O}$ lần lượt là trung điểm của ${MQ, MR}$.

Tứ giác ${EJOT}$ là hình gì?

$\Delta {PQR}$ cần thêm điều kiện gì để ${EJOT}$ là hình chữ nhật.
\end{minipage}\hfill
\begin{minipage}[t]{0.39\textwidth}
\vspace{0pt}
\centering

\resizebox{\linewidth}{!}{%
\begin{tikzpicture}
            \coordinate (A) at (0, 4);
            \coordinate (C) at (8, 0);
            \coordinate (B) at (-2, 0);
         \coordinate (D) at (-1, 2);
            \coordinate (I) at (0.5, 0);
              \coordinate (M) at (3, 0);
             \coordinate (E) at (4, 2);
                \coordinate (K) at (5.5, 0);
                 
            % Vẽ tứ giác
             \draw[thick] (A) -- (M);
              \draw[thick] (D) -- (I);
            \draw[thick] (A) -- (B) -- (C) -- (A);
             \draw[thick] (E) -- (K);
             \draw[thick] (D) -- (M)--(E);
    % Ghi nhãn các điểm
            
            \node[below left] at (B) {$Q$};
            \node[below right] at (C) {$R$};
            \node[above left] at (D) {$E$};
            \node[above] at (A) {$P$};
                 \node[below] at (I) {$J$};
                  \node[below] at (M) {$M$};
                  \node[below] at (K) {$O$};
                  \node[above right] at (E) {$T$};
              
             \def\khvuong[size=#1](#2,#3,#4){%
            \draw ($(#3)!#1!(#2)$) --
            ($($(#3)!#1!(#2)$)+($(#3)!#1!(#4)$)-(#3)$)--
            ($(#3)!#1!(#4)$);
            }
            \khvuong[size=6pt](B,A,C);
             \khvuong[size=6pt](A,D,M);
              \khvuong[size=6pt](A,E,M);
                
       \draw ($(B)!0.5!(I)$)node[rotate=40]{$/$};
        \draw ($(I)!0.5!(M)$)node[rotate=40]{$/$};
       \draw ($(M)!0.5!(K)$)node[rotate=40]{$//$} ;
       \draw ($(K)!0.5!(C)$)node[rotate=40]{$//$};
    \end{tikzpicture}
}
\end{minipage}

\end{tcolorbox}
\vspace{4pt}


\begin{tcolorbox}[
    breakable,
    colback=white,
    colframe=blue!60!black,
    boxrule=0.8pt,
    arc=3pt,
    title={\textbf{Câu \#213} \quad \textbf{MÃ CÂU: \texttt{[L08\_C3\_B4\_08]}} \quad | \quad \textbf{Nhóm: \texttt{HINH\_166}} \quad | \quad \textbf{MapID: \texttt{8H3V4}} \hfill \textcolor{yellow!90!white}{[Hình gốc .pyc]}}
]
\textbf{Dạng toán:} Bài tổng hợp 1 \hfill \textit{(Tự luận -- Vận dụng)}
\vspace{4pt}\hrule\vspace{6pt}
\begin{minipage}[t]{0.58\textwidth}
\vspace{0pt}
Cho tam giác $\Delta{QKH}$ vuông tại ${Q}$, có ${F}$ là trung điểm của ${KH}$. Gọi ${I, G }$ lần lượt là hình chiếu của ${F}$ lên ${QK, QH }$ Gọi ${I, G }$ lần lượt là hình chiếu của ${F}$ trên ${ QK, QH}$.

 a) Chứng minh ${I; G}$ lần lượt là trung điểm của ${QK, QH}$.

b) Chứng minh ${KIGF }$ là hình bình hành.

c) Lấy ${N}$ sao cho ${F}$ là trung điểm của $ {NG }$. Hạ $GO \bot KH$. Chứng minh $QO \bot ON$.
\end{minipage}\hfill
\begin{minipage}[t]{0.39\textwidth}
\vspace{0pt}
\centering

\resizebox{\linewidth}{!}{%
\begin{tikzpicture}
            \coordinate (A) at (0, 4);
            \coordinate (C) at (8, 0);
            \coordinate (B) at (-2, 0);
         \coordinate (D) at (-1, 2);
            \coordinate (K) at (4, 0);
              \coordinate (M) at (3, 0);
             \coordinate (E) at (4, 2);
                \coordinate (N) at (2, -2);
                 
            % Vẽ tứ giác
             \draw[thick] (A) -- (K);
              \draw[thick] (D) -- (M)--(E);
            \draw[thick] (A) -- (B) -- (C) -- (A);
             \draw[thick] (E) -- (D);
             \draw[thick] (B) -- (N)--(K)--(E);
               \draw[thick] (N) -- (M);
    % Ghi nhãn các điểm
            
            \node[below left] at (B) {$K$};
            \node[below right] at (C) {$H$};
            \node[above left] at (D) {$I$};
            \node[above] at (A) {$Q$};
                 \node[below] at (N) {$N$};
                  \node[below] at (M) {$F$};
                  \node[below] at (K) {$O$};
                  \node[above right] at (E) {$G$};
              
             \def\khvuong[size=#1](#2,#3,#4){%
            \draw ($(#3)!#1!(#2)$) --
            ($($(#3)!#1!(#2)$)+($(#3)!#1!(#4)$)-(#3)$)--
            ($(#3)!#1!(#4)$);
            }
            \khvuong[size=6pt](B,A,C);
             \khvuong[size=6pt](A,D,M);
              \khvuong[size=6pt](A,E,M);
                 \khvuong[size=6pt](E,K,C);
       \draw ($(M)!0.5!(N)$)node[rotate=40]{$/$};
        \draw ($(E)!0.5!(M)$)node[rotate=40]{$/$};
       \draw ($(M)!0.5!(B)$)node[rotate=40]{$//$} ;
       \draw ($(M)!0.5!(C)$)node[rotate=40]{$//$};
    \end{tikzpicture}
}
\end{minipage}

\end{tcolorbox}
\vspace{4pt}


\begin{tcolorbox}[
    breakable,
    colback=white,
    colframe=blue!60!black,
    boxrule=0.8pt,
    arc=3pt,
    title={\textbf{Câu \#214} \quad \textbf{MÃ CÂU: \texttt{[L08\_C3\_B4\_09]}} \quad | \quad \textbf{Nhóm: \texttt{HINH\_167}} \quad | \quad \textbf{MapID: \texttt{8H3V4}} \hfill \textcolor{yellow!90!white}{[Hình gốc .pyc]}}
]
\textbf{Dạng toán:} Bài tổng hợp 2 \hfill \textit{(Tự luận -- Vận dụng)}
\vspace{4pt}\hrule\vspace{6pt}
\begin{minipage}[t]{0.58\textwidth}
\vspace{0pt}
Cho $\Delta{ABC}$ vuông tại ${A}$. Điểm ${M}$ trên cạnh ${BC}$. Hạ $ ME \bot AB, MD \bot AC$.

 Tứ giác ${AEMD}$ là hình gì?

 Gọi ${AH }$ là đường cao $\Delta{ABC}$. Tính $\widehat{EHD} $
\end{minipage}\hfill
\begin{minipage}[t]{0.39\textwidth}
\vspace{0pt}
\centering

\resizebox{\linewidth}{!}{%
\begin{tikzpicture}
            \coordinate (A) at (0, 4);
            \coordinate (C) at (8, 0);
            \coordinate (B) at (-2, 0);
         \coordinate (D) at (4, 0);
            \coordinate (H) at (0, 0);
              \coordinate (M) at (-0.8, 2.4);
             \coordinate (I) at (2, 2);
                \coordinate (N) at (4.8, 1.6);
                 
            % Vẽ tứ giác
             \draw[thick] (A) -- (H);
              \draw[thick] (D) -- (M)--(N)--(D);
            \draw[thick] (A) -- (B) -- (C) -- (A);
    % Ghi nhãn các điểm
            
            \node[below left] at (B) {$B$};
            \node[below right] at (C) {$C$};
            \node[below] at (D) {$M$};
            \node[above] at (A) {$A$};
                 \node[above right] at (N) {$D$};
                  \node[above left] at (M) {$E$};
                  \node[below] at (H) {$H$};
                  \node[above right] at (I) {$I$};
              
             \def\khvuong[size=#1](#2,#3,#4){%
            \draw ($(#3)!#1!(#2)$) --
            ($($(#3)!#1!(#2)$)+($(#3)!#1!(#4)$)-(#3)$)--
            ($(#3)!#1!(#4)$);
            }
            \khvuong[size=6pt](B,A,C);
             \khvuong[size=6pt](A,M,D);
              \khvuong[size=6pt](A,N,D);
                 \khvuong[size=6pt](A,H,C);
       \draw ($(M)!0.5!(I)$)node[rotate=40]{$/$};
        \draw ($(I)!0.5!(N)$)node[rotate=40]{$/$};
         \fill [black] (I) circle (1pt);
    \end{tikzpicture}
}
\end{minipage}

\end{tcolorbox}
\vspace{4pt}


\begin{tcolorbox}[
    breakable,
    colback=white,
    colframe=blue!60!black,
    boxrule=0.8pt,
    arc=3pt,
    title={\textbf{Câu \#215} \quad \textbf{MÃ CÂU: \texttt{[L08\_C3\_B5\_01]}} \quad | \quad \textbf{Nhóm: \texttt{HINH\_168}} \quad | \quad \textbf{MapID: \texttt{8H3N5}} \hfill \textcolor{yellow!90!white}{[Hình gốc .pyc]}}
]
\textbf{Dạng toán:} Nhận biết 1 \hfill \textit{(Tự luận -- Nhận biết)}
\vspace{4pt}\hrule\vspace{6pt}
\begin{minipage}[t]{0.58\textwidth}
\vspace{0pt}
Cho hình bình hành ${QKHF}$ có $QF\bot QH$. Gọi ${M, N}$ lần lượt là trung điểm của ${ QK, HF }$ 

a) Chứng minh $MN \bot QH $.

b) Tứ giác ${QMHN}$ là hình gì?
\end{minipage}\hfill
\begin{minipage}[t]{0.39\textwidth}
\vspace{0pt}
\centering

\resizebox{\linewidth}{!}{%
\begin{tikzpicture}
            \coordinate (A) at (-1, 3);
            \coordinate (C) at (8, 0);
            \coordinate (B) at (9, 3);
         \coordinate (D) at (-2, 0);
              \coordinate (M) at (4, 3);
                \coordinate (N) at (3, 0);
                 
            % Vẽ tứ giác
             \draw[thick] (A) -- (N);
              \draw[thick] (M) -- (C);
               \draw[thick] (M) -- (N);
              \draw[thick] (A) -- (C);
            \draw[thick] (A) -- (B) -- (C) -- (D)--(A);
    % Ghi nhãn các điểm
            
            \node[above] at (B) {$K$};
            \node[below] at (C) {$H$};
            \node[below] at (D) {$F$};
            \node[above] at (A) {$Q$};
                 \node[below] at (N) {$N$};
                  \node[above ] at (M) {$M$};
             \def\khvuong[size=#1](#2,#3,#4){%
            \draw ($(#3)!#1!(#2)$) --
            ($($(#3)!#1!(#2)$)+($(#3)!#1!(#4)$)-(#3)$)--
            ($(#3)!#1!(#4)$);
            }
            \khvuong[size=6pt](D,A,C);
       \draw ($(M)!0.5!(A)$)node[rotate=40]{$/$};
        \draw ($(M)!0.5!(B)$)node[rotate=40]{$/$};
        \draw ($(N)!0.5!(D)$)node[rotate=40]{$/$};
        \draw ($(N)!0.5!(C)$)node[rotate=40]{$/$};
    \end{tikzpicture}
}
\end{minipage}

\end{tcolorbox}
\vspace{4pt}


\begin{tcolorbox}[
    breakable,
    colback=white,
    colframe=blue!60!black,
    boxrule=0.8pt,
    arc=3pt,
    title={\textbf{Câu \#216} \quad \textbf{MÃ CÂU: \texttt{[L08\_C3\_B5\_02]}} \quad | \quad \textbf{Nhóm: \texttt{HINH\_169}} \quad | \quad \textbf{MapID: \texttt{8H3N5}} \hfill \textcolor{yellow!90!white}{[Hình gốc .pyc]}}
]
\textbf{Dạng toán:} Nhận biết 2 \hfill \textit{(Tự luận -- Nhận biết)}
\vspace{4pt}\hrule\vspace{6pt}
\begin{minipage}[t]{0.58\textwidth}
\vspace{0pt}
Cho hình vuông ${PQRT}$. Trên các cạnh $AB, BC, CD, DA$ lấy lần lượt các điểm $M,N,P,Q$ sao cho $PM=QN=RP=TQ $ 

a) Chứng minh rằng: $MQ=NR=PT=QP $.

 b) Chứng minh $\Delta{QPM}= \Delta{NRP}$ 

c) Chứng minh ${MNPQ}$ là hình vuông.
\end{minipage}\hfill
\begin{minipage}[t]{0.39\textwidth}
\vspace{0pt}
\centering

\resizebox{\linewidth}{!}{%
\begin{tikzpicture}
            \coordinate (A) at (0, 0);
            \coordinate (C) at (3, 3);
            \coordinate (B) at (0, 3);
         \coordinate (D) at (3, 0);
              \coordinate (M) at (0, 1);
                \coordinate (N) at (1, 3);
                   \coordinate (P) at (3, 2);
                \coordinate (Q) at (2, 0);
            % Vẽ tứ giác
               \draw[thick] (M) -- (N)--(P)--(Q)--(M);
            \draw[thick] (A) -- (B) -- (C) -- (D)--(A);
    % Ghi nhãn các điểm
            
            \node[above left] at (B) {$Q$};
            \node[above right] at (C) {$R$};
            \node[below right] at (D) {$T$};
            \node[below left] at (A) {$P$};
                 \node[above] at (N) {$N$};
                  \node[above left ] at (M) {$M$};
                      \node[below right] at (P) {$P$};
                  \node[below] at (Q) {$Q$};
             \def\khvuong[size=#1](#2,#3,#4){%
            \draw ($(#3)!#1!(#2)$) --
            ($($(#3)!#1!(#2)$)+($(#3)!#1!(#4)$)-(#3)$)--
            ($(#3)!#1!(#4)$);
            }
            \khvuong[size=6pt](A,B,C);
             \khvuong[size=6pt](B,C,D);
              \khvuong[size=6pt](C,D,A);
               \khvuong[size=6pt](D,A,B);
       \draw ($(M)!0.5!(A)$)node[rotate=40]{$/$};
        \draw ($(N)!0.5!(B)$)node[rotate=40]{$/$};
        \draw ($(Q)!0.5!(D)$)node[rotate=40]{$/$};
        \draw ($(P)!0.5!(C)$)node[rotate=40]{$/$};
    \end{tikzpicture}
}
\end{minipage}

\end{tcolorbox}
\vspace{4pt}


\begin{tcolorbox}[
    breakable,
    colback=white,
    colframe=blue!60!black,
    boxrule=0.8pt,
    arc=3pt,
    title={\textbf{Câu \#217} \quad \textbf{MÃ CÂU: \texttt{[L08\_C3\_B5\_03]}} \quad | \quad \textbf{Nhóm: \texttt{HINH\_170}} \quad | \quad \textbf{MapID: \texttt{8H3N5}} \hfill \textcolor{yellow!90!white}{[Hình gốc .pyc]}}
]
\textbf{Dạng toán:} Nhận biết 3 \hfill \textit{(Tự luận -- Nhận biết)}
\vspace{4pt}\hrule\vspace{6pt}
\begin{minipage}[t]{0.58\textwidth}
\vspace{0pt}
Cho hình vuông ${IJQK }$. Trên cạnh ${JQ}$ lấy điểm ${E}$, trên cạnh ${QK}$ lấy điểm ${F}$ sao cho ${JE=QF }$

 a) Chứng minh $IE=JF$ 

b) Chứng minh $IE\bot JF$
\end{minipage}\hfill
\begin{minipage}[t]{0.39\textwidth}
\vspace{0pt}
\centering

\resizebox{\linewidth}{!}{%
\begin{tikzpicture}
            \coordinate (A) at (0, 0);
            \coordinate (C) at (3, 3);
            \coordinate (B) at (0, 3);
         \coordinate (D) at (3, 0);
              \coordinate (N) at (3, 2);
                \coordinate (M) at (1, 3);
               
            % Vẽ tứ giác
            \draw[thick] (A) -- (B) -- (C) -- (D)--(A);
             \draw[thick] (A) -- (M);
              \draw[thick] (B) -- (N);
    % Ghi nhãn các điểm
            
            \node[above left] at (B) {$J$};
            \node[above right] at (C) {$Q$};
            \node[below right] at (D) {$K$};
            \node[below left] at (A) {$I$};
                 \node[above right] at (N) {$F$};
                  \node[above] at (M) {$E$};
                  
             \def\khvuong[size=#1](#2,#3,#4){%
            \draw ($(#3)!#1!(#2)$) --
            ($($(#3)!#1!(#2)$)+($(#3)!#1!(#4)$)-(#3)$)--
            ($(#3)!#1!(#4)$);
            }
            \khvuong[size=6pt](A,B,C);
             \khvuong[size=6pt](B,C,D);
              \khvuong[size=6pt](C,D,A);
               \khvuong[size=6pt](D,A,B);
       \draw ($(M)!0.5!(B)$)node[rotate=40]{$/$};
        \draw ($(N)!0.5!(C)$)node[rotate=40]{$/$};
      
    \end{tikzpicture}
}
\end{minipage}

\end{tcolorbox}
\vspace{4pt}


\begin{tcolorbox}[
    breakable,
    colback=white,
    colframe=blue!60!black,
    boxrule=0.8pt,
    arc=3pt,
    title={\textbf{Câu \#218} \quad \textbf{MÃ CÂU: \texttt{[L08\_C3\_B5\_04]}} \quad | \quad \textbf{Nhóm: \texttt{HINH\_171}} \quad | \quad \textbf{MapID: \texttt{8H3H5}} \hfill \textcolor{yellow!90!white}{[Hình gốc .pyc]}}
]
\textbf{Dạng toán:} Thông hiểu 1 \hfill \textit{(Tự luận -- Thông hiểu)}
\vspace{4pt}\hrule\vspace{6pt}
\begin{minipage}[t]{0.58\textwidth}
\vspace{0pt}
Cho hình vuông ${QKHF }$. Trên cạnh ${QF}$ lấy điểm ${N}$, trên cạnh ${HF}$ lấy điểm ${M}$ sao cho ${QN=FM }$

 a) Chứng minh $\Delta{ QKN }=\Delta{ QFM }$ 

b) Chứng minh $\widehat{N Q M}+\widehat{Q N K} =90^{\circ} $ 

c) Chứng minh $QM \bot KN$.
\end{minipage}\hfill
\begin{minipage}[t]{0.39\textwidth}
\vspace{0pt}
\centering

\resizebox{\linewidth}{!}{%
\begin{tikzpicture}
            \coordinate (A) at (0, 0);
            \coordinate (C) at (3, 3);
            \coordinate (B) at (0, 3);
         \coordinate (D) at (3, 0);
              \coordinate (E) at (3, 2);
                \coordinate (F) at (2, 0);
    
            % Vẽ tứ giác
           \draw[thick] (A) -- (E);
            \draw[thick] (B) -- (F);
            \draw[thick] (A) -- (B) -- (C) -- (D)--(A);
    % Ghi nhãn các điểm
            
            \node[above left] at (B) {$K$};
            \node[above right] at (C) {$H$};
            \node[below right] at (D) {$F$};
            \node[below left] at (A) {$Q$};
                 \node[below] at (F) {$N$};
                      \node[below right] at (E) {$M$};
                 
             \def\khvuong[size=#1](#2,#3,#4){%
            \draw ($(#3)!#1!(#2)$) --
            ($($(#3)!#1!(#2)$)+($(#3)!#1!(#4)$)-(#3)$)--
            ($(#3)!#1!(#4)$);
            }
            \khvuong[size=6pt](A,B,C);
             \khvuong[size=6pt](B,C,D);
              \khvuong[size=6pt](C,D,A);
               \khvuong[size=6pt](D,A,B);
       \draw ($(F)!0.5!(A)$)node[rotate=40]{$/$};
        \draw ($(E)!0.5!(D)$)node[rotate=40]{$/$};
     
    \end{tikzpicture}
}
\end{minipage}

\end{tcolorbox}
\vspace{4pt}


\begin{tcolorbox}[
    breakable,
    colback=white,
    colframe=blue!60!black,
    boxrule=0.8pt,
    arc=3pt,
    title={\textbf{Câu \#219} \quad \textbf{MÃ CÂU: \texttt{[L08\_C3\_B5\_05]}} \quad | \quad \textbf{Nhóm: \texttt{HINH\_172}} \quad | \quad \textbf{MapID: \texttt{8H3H5}} \hfill \textcolor{yellow!90!white}{[Hình gốc .pyc]}}
]
\textbf{Dạng toán:} Thông hiểu 2 \hfill \textit{(Tự luận -- Thông hiểu)}
\vspace{4pt}\hrule\vspace{6pt}
\begin{minipage}[t]{0.58\textwidth}
\vspace{0pt}
Cho hình vuông ${QKHF }$, gọi ${M, N, O}$ lần lượt là trung điểm của $ {QK, KH, HF }$

 a) ${QMHO }$ là hình bình hành.

 b) Chứng minh $FN \bot HM$ tại ${I}$. c) ${QO}$ cắt ${FN}$ tại ${J}$. Chứng minh $JF=JI$
\end{minipage}\hfill
\begin{minipage}[t]{0.39\textwidth}
\vspace{0pt}
\centering

\resizebox{\linewidth}{!}{%
\begin{tikzpicture}
            \coordinate (A) at (0, 0);
            \coordinate (C) at (3, 3);
            \coordinate (B) at (0, 3);
         \coordinate (D) at (3, 0);
              \coordinate (E) at (0, 1.5);
                \coordinate (F) at (1.5, 3);
     \coordinate (K) at (3, 1.5);
              \coordinate (M) at (1.8, 2.4);
                \coordinate (N) at (2.4, 1.2);
            % Vẽ tứ giác
           \draw[thick] (C) -- (E);
            \draw[thick] (A) -- (K);
             \draw[thick] (D) -- (F);
            \draw[thick] (A) -- (B) -- (C) -- (D)--(A);
    % Ghi nhãn các điểm
            
            \node[above left] at (B) {$K$};
            \node[above right] at (C) {$H$};
            \node[below right] at (D) {$F$};
            \node[below left] at (A) {$Q$};
                 \node[above] at (F) {$N$};
                      \node[below left] at (E) {$M$};
                   \node[above left] at (M) {$I$};
                      \node[above right] at (N) {$J$};
                       \node[below right] at (K) {$O$};
             \def\khvuong[size=#1](#2,#3,#4){%
            \draw ($(#3)!#1!(#2)$) --
            ($($(#3)!#1!(#2)$)+($(#3)!#1!(#4)$)-(#3)$)--
            ($(#3)!#1!(#4)$);
            }
            \khvuong[size=6pt](A,B,C);
             \khvuong[size=6pt](B,C,D);
              \khvuong[size=6pt](C,D,A);
               \khvuong[size=6pt](D,A,B);
       \draw ($(E)!0.5!(A)$)node[rotate=50]{$/$};
        \draw ($(E)!0.5!(B)$)node[rotate=50]{$/$};
     \draw ($(F)!0.5!(B)$)node[rotate=50]{$/$};
        \draw ($(K)!0.5!(C)$)node[rotate=50]{$/$};
        \draw ($(F)!0.5!(C)$)node[rotate=50]{$/$};
        \draw ($(K)!0.5!(D)$)node[rotate=50]{$/$};
    \end{tikzpicture}
}
\end{minipage}

\end{tcolorbox}
\vspace{4pt}


\begin{tcolorbox}[
    breakable,
    colback=white,
    colframe=blue!60!black,
    boxrule=0.8pt,
    arc=3pt,
    title={\textbf{Câu \#220} \quad \textbf{MÃ CÂU: \texttt{[L08\_C3\_B5\_06]}} \quad | \quad \textbf{Nhóm: \texttt{HINH\_173}} \quad | \quad \textbf{MapID: \texttt{8H3H5}} \hfill \textcolor{yellow!90!white}{[Hình gốc .pyc]}}
]
\textbf{Dạng toán:} Thông hiểu 3 \hfill \textit{(Tự luận -- Thông hiểu)}
\vspace{4pt}\hrule\vspace{6pt}
\begin{minipage}[t]{0.58\textwidth}
\vspace{0pt}
Cho hình chữ nhật ${ABCD }$, tia phân giác $\widehat{C}$ và $\widehat{D}$ cắt nhau tại ${N}$, tia phân giác $\widehat{A}$ và $\widehat{B}$ cắt nhau tại ${M}$

 a) Tính các góc trong $\Delta {DNC} $.

 b) Chứng minh $\Delta {AMB}= \Delta {CND}$.

 c) ${BM}$ cắt ${CN}$ tại ${F}$.${AM}$ cắt ${DN}$ tại ${E}$. Chứng minh $ {EMFN}$ là hình vuông.
\end{minipage}\hfill
\begin{minipage}[t]{0.39\textwidth}
\vspace{0pt}
\centering

\resizebox{\linewidth}{!}{%
\begin{tikzpicture}
            \coordinate (A) at (0, 3);
            \coordinate (C) at (5, 0);
            \coordinate (B) at (5, 3);
         \coordinate (D) at (0, 0);
              \coordinate (E) at (2.5, 0.5);
                \coordinate (F) at (2.5, 2.5);
     \coordinate (H) at (3.5, 1.5);
              \coordinate (G) at (1.5, 1.5);
              
            % Vẽ tứ giác
           \draw[thick] (B) -- (E);
            \draw[thick] (A) -- (E);
             \draw[thick] (D) -- (F);
             \draw[thick] (C) -- (F);
            \draw[thick] (A) -- (B) -- (C) -- (D)--(A);
    % Ghi nhãn các điểm
            
            \node[above right] at (B) {$B$};
            \node[below right] at (C) {$C$};
            \node[below left] at (D) {$D$};
            \node[above left] at (A) {$A$};
                 \node[above] at (F) {$N$};
                      \node[below ] at (E) {$M$};
                   \node[left] at (G) {$E$};
                      \node[right] at (H) {$F$};
                    
             \def\khvuong[size=#1](#2,#3,#4){%
            \draw ($(#3)!#1!(#2)$) --
            ($($(#3)!#1!(#2)$)+($(#3)!#1!(#4)$)-(#3)$)--
            ($(#3)!#1!(#4)$);
            }
            \khvuong[size=6pt](A,B,C);
             \khvuong[size=6pt](B,C,D);
              \khvuong[size=6pt](C,D,A);
               \khvuong[size=6pt](D,A,B);
      
    \end{tikzpicture}
}
\end{minipage}

\end{tcolorbox}
\vspace{4pt}


\begin{tcolorbox}[
    breakable,
    colback=white,
    colframe=blue!60!black,
    boxrule=0.8pt,
    arc=3pt,
    title={\textbf{Câu \#221} \quad \textbf{MÃ CÂU: \texttt{[L08\_C3\_B5\_07]}} \quad | \quad \textbf{Nhóm: \texttt{HINH\_174}} \quad | \quad \textbf{MapID: \texttt{8H3V5}} \hfill \textcolor{yellow!90!white}{[Hình gốc .pyc]}}
]
\textbf{Dạng toán:} Bài tổng hợp 1 \hfill \textit{(Tự luận -- Vận dụng)}
\vspace{4pt}\hrule\vspace{6pt}
\begin{minipage}[t]{0.58\textwidth}
\vspace{0pt}
Cho hình chữ nhật ${EFGH}$ có $EF=2FG$. Gọi ${M}$ là trung điểm của ${EF}$ và ${N}$ là trung điểm của ${HG}$. 

 a) Chứng minh ${EMNH}$ và ${FMNG}$ là hình vuông.

 b) Chứng minh $\Delta{HMG}$ vuông cân. 

c) Gọi ${Q}$ và ${P}$ lần lượt là tâm các hình vuông ${EMNH, FMNG}$. Chứng minh ${MQNP}$ là hình vuông.
\end{minipage}\hfill
\begin{minipage}[t]{0.39\textwidth}
\vspace{0pt}
\centering

\resizebox{\linewidth}{!}{%
\begin{tikzpicture}
            \coordinate (A) at (0, 3);
            \coordinate (C) at (6, 0);
            \coordinate (B) at (6, 3);
         \coordinate (D) at (0, 0);
              \coordinate (K) at (3, 0);
                \coordinate (I) at (3, 3);
     \coordinate (R) at (4.5, 1.5);
              \coordinate (S) at (1.5, 1.5);
              
            % Vẽ tứ giác
           \draw[thick] (B) -- (K);
           \draw[thick] (I) -- (K);
            \draw[thick] (A) -- (K);
             \draw[thick] (D) -- (I);
             \draw[thick] (C) -- (I);
            \draw[thick] (A) -- (B) -- (C) -- (D)--(A);
    % Ghi nhãn các điểm
            
            \node[above right] at (B) {$F$};
            \node[below right] at (C) {$G$};
            \node[below left] at (D) {$H$};
            \node[above left] at (A) {$E$};
                 \node[above] at (S) {$Q$};
                      \node[above ] at (I) {$M$};
                   \node[below] at (K) {$N$};
                      \node[right] at (R) {$P$};
                    
             \def\khvuong[size=#1](#2,#3,#4){%
            \draw ($(#3)!#1!(#2)$) --
            ($($(#3)!#1!(#2)$)+($(#3)!#1!(#4)$)-(#3)$)--
            ($(#3)!#1!(#4)$);
            }
            \khvuong[size=6pt](A,B,C);
             \khvuong[size=6pt](B,C,D);
              \khvuong[size=6pt](C,D,A);
               \khvuong[size=6pt](D,A,B);
      
    \end{tikzpicture}
}
\end{minipage}

\end{tcolorbox}
\vspace{4pt}


\begin{tcolorbox}[
    breakable,
    colback=white,
    colframe=blue!60!black,
    boxrule=0.8pt,
    arc=3pt,
    title={\textbf{Câu \#222} \quad \textbf{MÃ CÂU: \texttt{[L08\_C3\_B5\_12]}} \quad | \quad \textbf{Nhóm: \texttt{HINH\_175}} \quad | \quad \textbf{MapID: \texttt{8H3H5}} \hfill \textcolor{yellow!90!white}{[Hình gốc .pyc]}}
]
\textbf{Dạng toán:} Chứng minh tứ giác là HBH, hình thoi \hfill \textit{(Tự luận -- Thông hiểu)}
\vspace{4pt}\hrule\vspace{6pt}
\begin{minipage}[t]{0.58\textwidth}
\vspace{0pt}
Cho tam giác ${CDF}$ vuông tại ${C}$ có ${CL}$ là đường trung tuyến. Kẻ ${LH\perp CD}$ tại ${H}$ và ${LJ\perp CF}$ tại ${J}$.

a) Chứng minh rằng tứ giác ${CHLJ}$ là hình chữ nhật.

b) Gọi ${N}$ là điểm đối xứng với ${L}$ qua ${J}$. Chứng minh rằng tứ giác ${CLFN}$ là hình thoi.\\
\end{minipage}\hfill
\begin{minipage}[t]{0.39\textwidth}
\vspace{0pt}
\centering

\resizebox{\linewidth}{!}{%
\begin{tikzpicture}[scale=0.6,line cap=round,line join=round,font=\footnotesize]
            \path
            (0,0) coordinate (A)
            (0,5) coordinate (B)
            (7,0) coordinate (C)
            ($(B)!.5!(C)$) coordinate (M)
            ($(A)!(M)!(B)$) coordinate (D)
            ($(A)!(M)!(C)$) coordinate (E)
            ($(M)!2!(E)$) coordinate (K)% M la tâm đối xứng
            ;
            \draw (A)--(B)--(C)--cycle
            (D)--(M) (M)--(K) (A)--(M) (A)--(K)
            ;
            \path (B)--(M) node[midway, sloped]{$/$} (M)--(C)node[midway, sloped]{$/$} ;
            \foreach \x/\y/\z in {B/A/C,A/E/M,A/D/M}\draw pic[draw,angle radius =2mm] {right angle = \x--\y--\z};
           \filldraw[black] (A) circle (2pt) node[above left ] {$C$};
    \filldraw[black] (B) circle (2pt) node[above right] {$D$};
    \filldraw[black] (C) circle (2pt) node[above ] {$F$};
   \filldraw[black] (D) circle (2pt) node[below left]{$H$};
 \filldraw[black] (E) circle (2pt) node[above left ] {$J$};
    \filldraw[black] (K) circle (2pt) node[above right] {$N$};
     \filldraw[black] (M) circle (2pt) node[above right] {$L$};
            \end{tikzpicture}
}
\end{minipage}

\end{tcolorbox}
\vspace{4pt}


\subsection{Chương 4 - Định lí Talet}

\begin{tcolorbox}[
    breakable,
    colback=white,
    colframe=blue!60!black,
    boxrule=0.8pt,
    arc=3pt,
    title={\textbf{Câu \#223} \quad \textbf{MÃ CÂU: \texttt{[L08\_C4\_B1\_01]}} \quad | \quad \textbf{Nhóm: \texttt{HINH\_176}} \quad | \quad \textbf{MapID: \texttt{8H4N1}} \hfill \textcolor{yellow!90!white}{[Hình gốc .pyc]}}
]
\textbf{Dạng toán:} Áp dụng dl THALES tìm x \hfill \textit{(Trắc nghiệm -- Nhận biết)}
\vspace{4pt}\hrule\vspace{6pt}
\begin{minipage}[t]{0.58\textwidth}
\vspace{0pt}
Tính độ dài ${x}$ trong hình sau biết ${EF//BC}$ (làm tròn kết quả đến chữ số thập phân thứ nhất)

\vspace{4pt}\noindent
\textbf{A.} 20,2 \quad \textbf{B.} 25,2 \quad \textbf{\textcolor{blue!75!black}{C. (Đ)}} $22,2$ \quad \textbf{D.} 23,2
\end{minipage}\hfill
\begin{minipage}[t]{0.39\textwidth}
\vspace{0pt}
\centering

\resizebox{\linewidth}{!}{%
\begin{tikzpicture}
    % Coordinates of the triangle
    \coordinate [label=above:$A$] (A) at (0,3);
    \coordinate [label=left:$B$] (B) at (-2,0);
    \coordinate [label=right:$C$] (C) at (4,0);
    
    % Coordinates for points E and F
    \coordinate [label=left:$E$] (E) at (-4/3, 1);
    \coordinate [label=right:$F$] (F) at (8/3, 1);
    
    % Draw the triangle ABC
    \draw [thick] (A) -- (B) -- (C) -- cycle;
    
    % Draw line EF
    \draw [thick] (E) -- (F);
    
    % Label side lengths
    \node at (-0.9, 2.3) {$20$}; % AB
    \node at (1.2, 2.3) {$x$}; % AC
    \node at (-2, 0.5) {$9$}; % BE
    \node at (3.7, 0.5) {$10$}; % CF
     \fill (A) circle (1.2pt);
    \fill (B) circle (1.2pt);
    \fill (C) circle (1.2pt);
    \fill (E) circle (1.2pt);
    \fill (F) circle (1.2pt);
\end{tikzpicture}
}
\end{minipage}

\end{tcolorbox}
\vspace{4pt}


\begin{tcolorbox}[
    breakable,
    colback=white,
    colframe=blue!60!black,
    boxrule=0.8pt,
    arc=3pt,
    title={\textbf{Câu \#224} \quad \textbf{MÃ CÂU: \texttt{[L08\_C4\_B1\_02]}} \quad | \quad \textbf{Nhóm: \texttt{HINH\_177}} \quad | \quad \textbf{MapID: \texttt{8H4H1}} \hfill \textcolor{yellow!90!white}{[Hình gốc .pyc]}}
]
\textbf{Dạng toán:} Chỉ ra cặp dt song song. Áp dụng dl THALES tìm x \hfill \textit{(Trắc nghiệm -- Thông hiểu)}
\vspace{4pt}\hrule\vspace{6pt}
\begin{minipage}[t]{0.58\textwidth}
\vspace{0pt}
Tính độ dài ${x}$ trong hình (làm tròn kết quả đến chữ số thập phân thứ nhất)

\vspace{4pt}\noindent
\textbf{A.} 22,1 \quad \textbf{\textcolor{blue!75!black}{B. (Đ)}} $17,1$ \quad \textbf{C.} 19,1 \quad \textbf{D.} 16,1
\end{minipage}\hfill
\begin{minipage}[t]{0.39\textwidth}
\vspace{0pt}
\centering

\resizebox{\linewidth}{!}{%
\begin{tikzpicture}
        % Coordinates
        \coordinate (A) at (0, 0);
        \coordinate (H) at (0, 2);
        \coordinate (B) at (0, 4);
        \coordinate (C) at (4, 0);
        \coordinate (M) at (2, 2);
        
        % Draw the triangle
        \draw [thick] (A) -- (B) -- (C) -- cycle;
        \draw [thick] (H) -- (M);
        
        % Right angle marks
        \draw (A) rectangle +(0.3,0.3);
        \draw (H) rectangle +(0.3,-0.3);
        
        % Labels
        \node [left] at (A) {$A$};
        \node [left] at (H) {$H$};
        \node [left] at (B) {$B$};
        \node [right] at (C) {$C$};
        \node [above right] at (M) {$M$};
        
        % Side lengths
        \node [left] at (0, 3) {$9$};
        \node [left] at (0, 1) {$10$};
        \node [above] at (1, 3) {$x$};
        \node [right] at (3, 1) {$19$};
    \end{tikzpicture}
}
\end{minipage}

\end{tcolorbox}
\vspace{4pt}


\begin{tcolorbox}[
    breakable,
    colback=white,
    colframe=blue!60!black,
    boxrule=0.8pt,
    arc=3pt,
    title={\textbf{Câu \#225} \quad \textbf{MÃ CÂU: \texttt{[L08\_C4\_B1\_03]}} \quad | \quad \textbf{Nhóm: \texttt{HINH\_178}} \quad | \quad \textbf{MapID: \texttt{8H4H1}} \hfill \textcolor{yellow!90!white}{[Hình gốc .pyc]}}
]
\textbf{Dạng toán:} Toán thực tế áp dụng thalet \hfill \textit{(Trắc nghiệm -- Thông hiểu)}
\vspace{4pt}\hrule\vspace{6pt}
\begin{minipage}[t]{0.58\textwidth}
\vspace{0pt}
Để đo khoảng cách giữa hai vị trí ${B}$ và ${E}$ ở hai bên bờ sông, bác An chọn ba vị trí ${A}$, ${F}$, ${C}$ cùng nằm ở một bên bờ sông sao cho ba điểm ${C}$, ${E}$, ${B}$ thẳng hàng, ba điểm ${C}$, ${F}$, ${A}$ thẳng hàng và ${AB\parallel FE}$ (như hình vẽ). Sau đó bác An đo được ${AF=14}$ m, ${FC=7}$ m, ${EC=17}$ m. Hỏi khoảng cách giữa hai vị trí ${B}$ và ${E}$ bằng bao nhiêu ?

\vspace{4pt}\noindent
\textbf{\textcolor{blue!75!black}{A. (Đ)}} $34$ \quad \textbf{B.} $33$ \quad \textbf{C.} $38$ \quad \textbf{D.} $36$
\end{minipage}\hfill
\begin{minipage}[t]{0.39\textwidth}
\vspace{0pt}
\centering

\resizebox{\linewidth}{!}{%
\begin{tikzpicture}
        \def\cao{4}
        \path
        (1,3.5)coordinate (B)
        (1,1.4)coordinate (A)
        (1,.5)coordinate (H)
        +(-8:3)coordinate (C)
        ($(B)!0.65!(C)$)coordinate (E)
        ($(A)!0.65!(C)$)coordinate (F)
        ;
        %           \pgfresetboundingbox
        \clip (0,0)rectangle(5,\cao);
        \fill[lime!70!olive](0,0)rectangle(5,\cao);
        \pgfmathsetseed{14159}
        \draw[decorate, decoration={snake,post length=0.5 mm,random steps, segment length=3mm, amplitude=3pt}, lime!70!black, line width=.5mm, double=cyan!60!white, double distance=2cm](0,2.5)--(5,2.5);
        \draw (B)--(C)--(A);
        \draw[dashed] (A)--(B) (E)--(F);
        \foreach \p/\g/\pp in{
            A/-120/A,
            B/150/B,
            C/20/C,
            E/20/E,
            F/-130/F
        } \shade[shading=ball,fill=cyan](\p)circle(1.5pt)node[shift={(\g:.35)}]{$\pp$};
        \path (H)--(C);
        \end{tikzpicture}
}
\end{minipage}

\vspace{3pt}\hrule\vspace{3pt}\footnotesize\textcolor{gray!80!black}{\textbf{Các mã dùng chung hình (HINH\_178):} \texttt{L08\_C4\_B1\_05, L08\_C4\_B1\_04}}
\end{tcolorbox}
\vspace{4pt}


\begin{tcolorbox}[
    breakable,
    colback=white,
    colframe=blue!60!black,
    boxrule=0.8pt,
    arc=3pt,
    title={\textbf{Câu \#226} \quad \textbf{MÃ CÂU: \texttt{[L08\_C4\_B1\_05]}} \quad | \quad \textbf{Nhóm: \texttt{HINH\_178}} \quad | \quad \textbf{MapID: \texttt{8H4H1}} \hfill \textcolor{yellow!90!white}{[Hình gốc .pyc]}}
]
\textbf{Dạng toán:} SA-M2. Toán thực tế áp dụng thalet \hfill \textit{(Trả lời ngắn -- Thông hiểu)}
\vspace{4pt}\hrule\vspace{6pt}
\begin{minipage}[t]{0.58\textwidth}
\vspace{0pt}
Để đo khoảng cách giữa hai vị trí ${B}$ và ${E}$ ở hai bên bờ sông, bác An chọn ba vị trí ${A}$, ${F}$, ${C}$ cùng nằm ở một bên bờ sông sao cho ba điểm ${C}$, ${E}$, ${B}$ thẳng hàng, ba điểm ${C}$, ${F}$, ${A}$ thẳng hàng và ${AB\parallel FE}$ (như hình vẽ). Sau đó bác An đo được ${AF=30}$ m, ${FC=15}$ m, ${EC=12}$ m. Hỏi khoảng cách giữa hai vị trí ${B}$ và ${E}$ bằng bao nhiêu ?


\textbf{Đáp án điền số:} 	exttt{24}
\end{minipage}\hfill
\begin{minipage}[t]{0.39\textwidth}
\vspace{0pt}
\centering

\resizebox{\linewidth}{!}{%
\begin{tikzpicture}
        \def\cao{4}
        \path
        (1,3.5)coordinate (B)
        (1,1.4)coordinate (A)
        (1,.5)coordinate (H)
        +(-8:3)coordinate (C)
        ($(B)!0.65!(C)$)coordinate (E)
        ($(A)!0.65!(C)$)coordinate (F)
        ;
        %           \pgfresetboundingbox
        \clip (0,0)rectangle(5,\cao);
        \fill[lime!70!olive](0,0)rectangle(5,\cao);
        \pgfmathsetseed{14159}
        \draw[decorate, decoration={snake,post length=0.5 mm,random steps, segment length=3mm, amplitude=3pt}, lime!70!black, line width=.5mm, double=cyan!60!white, double distance=2cm](0,2.5)--(5,2.5);
        \draw (B)--(C)--(A);
        \draw[dashed] (A)--(B) (E)--(F);
        \foreach \p/\g/\pp in{
            A/-120/A,
            B/150/B,
            C/20/C,
            E/20/E,
            F/-130/F
        } \shade[shading=ball,fill=cyan](\p)circle(1.5pt)node[shift={(\g:.35)}]{$\pp$};
        \path (H)--(C);
        \end{tikzpicture}
}
\end{minipage}

\vspace{3pt}\hrule\vspace{3pt}\footnotesize\textcolor{gray!80!black}{\textbf{Các mã dùng chung hình (HINH\_178):} \texttt{L08\_C4\_B1\_03, L08\_C4\_B1\_04}}
\end{tcolorbox}
\vspace{4pt}


\begin{tcolorbox}[
    breakable,
    colback=white,
    colframe=blue!60!black,
    boxrule=0.8pt,
    arc=3pt,
    title={\textbf{Câu \#227} \quad \textbf{MÃ CÂU: \texttt{[L08\_C4\_B1\_04]}} \quad | \quad \textbf{Nhóm: \texttt{HINH\_178}} \quad | \quad \textbf{MapID: \texttt{8H4H1}} \hfill \textcolor{yellow!90!white}{[Hình gốc .pyc]}}
]
\textbf{Dạng toán:} Toán thực tế áp dụng thalet \hfill \textit{(Tự luận -- Thông hiểu)}
\vspace{4pt}\hrule\vspace{6pt}
\begin{minipage}[t]{0.58\textwidth}
\vspace{0pt}
Để đo khoảng cách giữa hai vị trí ${B}$ và ${E}$ ở hai bên bờ sông, bác An chọn ba vị trí ${A}$, ${F}$, ${C}$ cùng nằm ở một bên bờ sông sao cho ba điểm ${C}$, ${E}$, ${B}$ thẳng hàng, ba điểm ${C}$, ${F}$, ${A}$ thẳng hàng và ${AB\parallel FE}$ (như hình vẽ). Sau đó bác An đo được ${AF=28}$ m, ${FC=14}$ m, ${EC=14}$ m. Hỏi khoảng cách giữa hai vị trí ${B}$ và ${E}$ bằng bao nhiêu ?\\
\end{minipage}\hfill
\begin{minipage}[t]{0.39\textwidth}
\vspace{0pt}
\centering

\resizebox{\linewidth}{!}{%
\begin{tikzpicture}
        \def\cao{4}
        \path
        (1,3.5)coordinate (B)
        (1,1.4)coordinate (A)
        (1,.5)coordinate (H)
        +(-8:3)coordinate (C)
        ($(B)!0.65!(C)$)coordinate (E)
        ($(A)!0.65!(C)$)coordinate (F)
        ;
        %           \pgfresetboundingbox
        \clip (0,0)rectangle(5,\cao);
        \fill[lime!70!olive](0,0)rectangle(5,\cao);
        \pgfmathsetseed{14159}
        \draw[decorate, decoration={snake,post length=0.5 mm,random steps, segment length=3mm, amplitude=3pt}, lime!70!black, line width=.5mm, double=cyan!60!white, double distance=2cm](0,2.5)--(5,2.5);
        \draw (B)--(C)--(A);
        \draw[dashed] (A)--(B) (E)--(F);
        \foreach \p/\g/\pp in{
            A/-120/A,
            B/150/B,
            C/20/C,
            E/20/E,
            F/-130/F
        } \shade[shading=ball,fill=cyan](\p)circle(1.5pt)node[shift={(\g:.35)}]{$\pp$};
        \path (H)--(C);
        \end{tikzpicture}
}
\end{minipage}

\vspace{3pt}\hrule\vspace{3pt}\footnotesize\textcolor{gray!80!black}{\textbf{Các mã dùng chung hình (HINH\_178):} \texttt{L08\_C4\_B1\_03, L08\_C4\_B1\_05}}
\end{tcolorbox}
\vspace{4pt}


\begin{tcolorbox}[
    breakable,
    colback=white,
    colframe=blue!60!black,
    boxrule=0.8pt,
    arc=3pt,
    title={\textbf{Câu \#228} \quad \textbf{MÃ CÂU: \texttt{[L08\_C4\_B2\_01]}} \quad | \quad \textbf{Nhóm: \texttt{HINH\_179}} \quad | \quad \textbf{MapID: \texttt{8H4H2}} \hfill \textcolor{yellow!90!white}{[Hình gốc .pyc]}}
]
\textbf{Dạng toán:} Tìm số đo x \hfill \textit{(Trắc nghiệm -- Thông hiểu)}
\vspace{4pt}\hrule\vspace{6pt}
\begin{minipage}[t]{0.58\textwidth}
\vspace{0pt}
Tìm số đo ${x}$ trong hình

\vspace{4pt}\noindent
\textbf{\textcolor{blue!75!black}{A. (Đ)}} $14$ \quad \textbf{B.} $15$ \quad \textbf{C.} $13$ \quad \textbf{D.} $12$
\end{minipage}\hfill
\begin{minipage}[t]{0.39\textwidth}
\vspace{0pt}
\centering

\resizebox{\linewidth}{!}{%
\begin{tikzpicture}[scale=1.5]
    % Vẽ tam giác ABC
    \coordinate [label=below left: $B$] (B) at (0,0);
    \coordinate [label=below right: $C$] (C) at (4,0);
    \coordinate [label=above: $A$] (A) at (4,3);
    
    \draw [thick] (B) -- (C) -- (A) -- cycle;
    
    % Điểm M và N
    \coordinate [label=above left: $M$] (M) at (1.8,1.35);
    \coordinate [label=right: $N$] (N) at (4,1.35);
    
    % Vẽ đoạn thẳng MN
    \draw [thick] (M) -- (N);
    
    % Ký hiệu độ dài
    \node at (3,1.5) {$7$};
    \node at (3,0.1) {$x$};
    
    % Ký hiệu tam giác đồng dạng và độ dài các đoạn thẳng
    \draw [thick] (B) -- (M);
    \draw [thick] (A) -- (N);
        \draw
    ($(B)!0.5!(M)$)node[rotate=20]{$/$} ($(M)!0.5!(A)$)node[rotate=20]{$/$}
    ($(A)!0.5!(N)$)node[rotate=60]{$//$} ($(N)!0.5!(C)$)node[rotate=60]{$//$};
    % Ký hiệu các đoạn bằng nhau

\end{tikzpicture}
}
\end{minipage}

\vspace{3pt}\hrule\vspace{3pt}\footnotesize\textcolor{gray!80!black}{\textbf{Các mã dùng chung hình (HINH\_179):} \texttt{L08\_C4\_B2\_08, L08\_C4\_B2\_07}}
\end{tcolorbox}
\vspace{4pt}


\begin{tcolorbox}[
    breakable,
    colback=white,
    colframe=blue!60!black,
    boxrule=0.8pt,
    arc=3pt,
    title={\textbf{Câu \#229} \quad \textbf{MÃ CÂU: \texttt{[L08\_C4\_B2\_02]}} \quad | \quad \textbf{Nhóm: \texttt{HINH\_180}} \quad | \quad \textbf{MapID: \texttt{8H4H2}} \hfill \textcolor{yellow!90!white}{[Hình gốc .pyc]}}
]
\textbf{Dạng toán:} Chứng mình đường trung bình từ đó tìm số đo x \hfill \textit{(Trắc nghiệm -- Thông hiểu)}
\vspace{4pt}\hrule\vspace{6pt}
\begin{minipage}[t]{0.58\textwidth}
\vspace{0pt}
Tìm số đo ${x}$ trong hình

\vspace{4pt}\noindent
\textbf{A.} $14$ \quad \textbf{\textcolor{blue!75!black}{B. (Đ)}} $16$ \quad \textbf{C.} $17$ \quad \textbf{D.} $15$
\end{minipage}\hfill
\begin{minipage}[t]{0.39\textwidth}
\vspace{0pt}
\centering

\resizebox{\linewidth}{!}{%
\begin{tikzpicture}
    
    % Coordinates for the vertices
    \coordinate [label=above:$A$] (A) at (0,4);
    \coordinate [label=below left:$B$] (B) at (-2,0);
    \coordinate [label=below right:$C$] (C) at (4,0);
    \coordinate [label=above left:$K$] (K) at (-1,2);
        \coordinate [label=above right:$I$] (I) at (2,2);
    % Draw the triangle
    \draw (A) -- (B) -- (C) -- cycle;
    
    % Draw the segment AK
    \draw (I) -- (K);
    
 % Ký hiệu độ dài
\node at (1,1.6) {$8$};
\node at (1,-0.4) {$x$};
 \draw 
(C) pic[draw,angle radius = 15] {angle = A--C--B}  ;
 \draw 
(I) pic[draw,angle radius = 15] {angle = A--I--K}  ;
\draw
($(A)!0.5!(K)$)node[rotate=40]{$/$} ($(B)!0.5!(K)$)node[rotate=40]{$/$};

\end{tikzpicture}
}
\end{minipage}

\vspace{3pt}\hrule\vspace{3pt}\footnotesize\textcolor{gray!80!black}{\textbf{Các mã dùng chung hình (HINH\_180):} \texttt{L08\_C4\_B2\_09, L08\_C4\_B2\_10}}
\end{tcolorbox}
\vspace{4pt}


\begin{tcolorbox}[
    breakable,
    colback=white,
    colframe=blue!60!black,
    boxrule=0.8pt,
    arc=3pt,
    title={\textbf{Câu \#230} \quad \textbf{MÃ CÂU: \texttt{[L08\_C4\_B2\_06]}} \quad | \quad \textbf{Nhóm: \texttt{HINH\_181}} \quad | \quad \textbf{MapID: \texttt{8H4H2}} \hfill \textcolor{yellow!90!white}{[Hình gốc .pyc]}}
]
\textbf{Dạng toán:} SA-M2. Toán thực tế về đường TB \hfill \textit{(Trả lời ngắn -- Thông hiểu)}
\vspace{4pt}\hrule\vspace{6pt}
\begin{minipage}[t]{0.58\textwidth}
\vspace{0pt}
Giữa hai điểm B và C bị ngăn cách bởi hồ nước (như hình vẽ). Xác định độ dài BC mà không cần phải di chuyển qua hồ nước. Biết rằng đoạn thẳng KI dài $18$ m và K là trung điểm của AB, I là trung điểm của AC.


\textbf{Đáp án điền số:} 	exttt{36}
\end{minipage}\hfill
\begin{minipage}[t]{0.39\textwidth}
\vspace{0pt}
\centering

\resizebox{\linewidth}{!}{%
\begin{tikzpicture}[scale=1, font=\footnotesize, line join=round, line cap=round, >=stealth]
            \tikzset{co/.pic={
                    \fill[green!80!teal,out=110,in=-20] (0,0)to(-1,1)[out=-10,in=120]to(-0.2,0.7)[out=110,in=-20]to(0,1.5)[out=-40,in=100]to(0.5,0.5)[out=40,in=-120]to(1,1.3)--(0.75,0)--cycle;
            }}
            \tikzset{Ho/.pic={
                    \fill[green!30](-2.2,1.5)--(5,1.8)..controls++(0:1) and++(-10:1)..(3,1)..controls++(-30:3) and++(-5:1)..(0.3,0.5)..controls++(175:0.5) and++(0:2)..(0,0)..controls++(180:1) and++(-120:1)..(-1.3,0.8)..controls++(60:0.5) and++(-90:1)..(-2.2,1.5);
                    \fill[cyan](-1,0.8)..controls++(20:2) and++(185:1)..(2,1.5)..controls++(5:1) and++(50:0.5)..(3,1.5)..controls++(-130:0.25)and++(15:2.5)..(2.5,1)..controls++(-165:1) and++(-60:0.5)..(-1,0.8);
                    \draw[red](-1,0.8)coordinate(A)--(3.7,1.5)coordinate(B)node[above]{$C$}--(1.5,-1.2)coordinate(C)node[below]{$A$}--(A)node[above]{$B$};
            }}
            \begin{scope}[scale=.75]
                \path(0,0)pic{Ho};
                \foreach \x/\y in{-1.5/1,-0.5/0.1,3/0.5,4/1.2,-0.5/1.2}
                \path(\x,\y)pic[scale=0.25]{co};
                (A)--(B)--(C)--cycle                
                \path ($(A)!.5!(C)$)coordinate(K) node[below]{$K$}
                ($(B)!.5!(C)$)coordinate(I)node[below]{$I$};
                \draw   (C)++(90:0.5) 
                (I)--(K);
                \foreach \y/\g in {A/90,B/40,C/-90,I/-90,K/90}       
                \fill (\y) circle (1pt)+(\g:.5);
            \end{scope}     
    \end{tikzpicture}
}
\end{minipage}

\vspace{3pt}\hrule\vspace{3pt}\footnotesize\textcolor{gray!80!black}{\textbf{Các mã dùng chung hình (HINH\_181):} \texttt{L08\_C4\_B2\_05}}
\end{tcolorbox}
\vspace{4pt}


\begin{tcolorbox}[
    breakable,
    colback=white,
    colframe=blue!60!black,
    boxrule=0.8pt,
    arc=3pt,
    title={\textbf{Câu \#231} \quad \textbf{MÃ CÂU: \texttt{[L08\_C4\_B2\_08]}} \quad | \quad \textbf{Nhóm: \texttt{HINH\_179}} \quad | \quad \textbf{MapID: \texttt{8H4N2}} \hfill \textcolor{yellow!90!white}{[Hình gốc .pyc]}}
]
\textbf{Dạng toán:} SA-M1. Tìm số đo x \hfill \textit{(Trả lời ngắn -- Nhận biết)}
\vspace{4pt}\hrule\vspace{6pt}
\begin{minipage}[t]{0.58\textwidth}
\vspace{0pt}
Tìm số đo ${x}$ trong hình


\textbf{Đáp án điền số:} 	exttt{4}
\end{minipage}\hfill
\begin{minipage}[t]{0.39\textwidth}
\vspace{0pt}
\centering

\resizebox{\linewidth}{!}{%
\begin{tikzpicture}[scale=1.5]
    % Vẽ tam giác ABC
    \coordinate [label=below left: $B$] (B) at (0,0);
    \coordinate [label=below right: $C$] (C) at (4,0);
    \coordinate [label=above: $A$] (A) at (4,3);
    
    \draw [thick] (B) -- (C) -- (A) -- cycle;
    
    % Điểm M và N
    \coordinate [label=above left: $M$] (M) at (1.8,1.35);
    \coordinate [label=right: $N$] (N) at (4,1.35);
    
    % Vẽ đoạn thẳng MN
    \draw [thick] (M) -- (N);
    
    % Ký hiệu độ dài
    \node at (3,1.5) {$2$};
    \node at (3,0.1) {$x$};
    
    % Ký hiệu tam giác đồng dạng và độ dài các đoạn thẳng
    \draw [thick] (B) -- (M);
    \draw [thick] (A) -- (N);
        \draw
    ($(B)!0.5!(M)$)node[rotate=20]{$/$} ($(M)!0.5!(A)$)node[rotate=20]{$/$}
    ($(A)!0.5!(N)$)node[rotate=60]{$//$} ($(N)!0.5!(C)$)node[rotate=60]{$//$};
    % Ký hiệu các đoạn bằng nhau

\end{tikzpicture}
}
\end{minipage}

\vspace{3pt}\hrule\vspace{3pt}\footnotesize\textcolor{gray!80!black}{\textbf{Các mã dùng chung hình (HINH\_179):} \texttt{L08\_C4\_B2\_01, L08\_C4\_B2\_07}}
\end{tcolorbox}
\vspace{4pt}


\begin{tcolorbox}[
    breakable,
    colback=white,
    colframe=blue!60!black,
    boxrule=0.8pt,
    arc=3pt,
    title={\textbf{Câu \#232} \quad \textbf{MÃ CÂU: \texttt{[L08\_C4\_B2\_09]}} \quad | \quad \textbf{Nhóm: \texttt{HINH\_180}} \quad | \quad \textbf{MapID: \texttt{8H4H2}} \hfill \textcolor{yellow!90!white}{[Hình gốc .pyc]}}
]
\textbf{Dạng toán:} SA-M2. Chứng mình đường trung bình từ đó tìm số đo x \hfill \textit{(Trả lời ngắn -- Thông hiểu)}
\vspace{4pt}\hrule\vspace{6pt}
\begin{minipage}[t]{0.58\textwidth}
\vspace{0pt}
Tìm số đo ${x}$ trong hình


\textbf{Đáp án điền số:} 	exttt{6}
\end{minipage}\hfill
\begin{minipage}[t]{0.39\textwidth}
\vspace{0pt}
\centering

\resizebox{\linewidth}{!}{%
\begin{tikzpicture}
    
    % Coordinates for the vertices
    \coordinate [label=above:$A$] (A) at (0,4);
    \coordinate [label=below left:$B$] (B) at (-2,0);
    \coordinate [label=below right:$C$] (C) at (4,0);
    \coordinate [label=above left:$K$] (K) at (-1,2);
        \coordinate [label=above right:$I$] (I) at (2,2);
    % Draw the triangle
    \draw (A) -- (B) -- (C) -- cycle;
    
    % Draw the segment AK
    \draw (I) -- (K);
    
 % Ký hiệu độ dài
\node at (1,1.6) {$3$};
\node at (1,-0.4) {$x$};
 \draw 
(C) pic[draw,angle radius = 15] {angle = A--C--B}  ;
 \draw 
(I) pic[draw,angle radius = 15] {angle = A--I--K}  ;
\draw
($(A)!0.5!(K)$)node[rotate=40]{$/$} ($(B)!0.5!(K)$)node[rotate=40]{$/$};

\end{tikzpicture}
}
\end{minipage}

\vspace{3pt}\hrule\vspace{3pt}\footnotesize\textcolor{gray!80!black}{\textbf{Các mã dùng chung hình (HINH\_180):} \texttt{L08\_C4\_B2\_02, L08\_C4\_B2\_10}}
\end{tcolorbox}
\vspace{4pt}


\begin{tcolorbox}[
    breakable,
    colback=white,
    colframe=blue!60!black,
    boxrule=0.8pt,
    arc=3pt,
    title={\textbf{Câu \#233} \quad \textbf{MÃ CÂU: \texttt{[L08\_C4\_B2\_04]}} \quad | \quad \textbf{Nhóm: \texttt{HINH\_182}} \quad | \quad \textbf{MapID: \texttt{8H4N2}} \hfill \textcolor{yellow!90!white}{[Hình gốc .pyc]}}
]
\textbf{Dạng toán:} Bài tổng hợp 1 \hfill \textit{(Tự luận -- Nhận biết)}
\vspace{4pt}\hrule\vspace{6pt}
\begin{minipage}[t]{0.58\textwidth}
\vspace{0pt}
Cho tam giác ${GHF}$ cân tại ${G}$ đường cao ${GJ}$ , ${I}$ là trung điểm của ${GF}$. Từ ${G}$ kẻ tia ${Gy}$ song song với ${HF}$ cắt ${JI}$ tại ${E}$ 

a) Chứng minh ${JH=JF}$.

b) Chứng minh $JE //HF$.

c) Chứng minh $GE=JF$
\end{minipage}\hfill
\begin{minipage}[t]{0.39\textwidth}
\vspace{0pt}
\centering

\resizebox{\linewidth}{!}{%
\begin{tikzpicture}
    % Tọa độ của các điểm
    \coordinate [label=above:{$G$}] (A) at (0,3);
    \coordinate [label=below left:{$H$}] (B) at (-2,0);
    \coordinate [label=below right:{$F$}] (C) at (2,0);
    \coordinate [label=below:{$J$}] (M) at (0,0); % Trung điểm của BC
    \coordinate [label=right:{$I$}] (N) at (1,1.5); % Trung điểm của AC
    \coordinate [label=above:{$E$}] (E) at (2,3);
    % Vẽ tam giác ABC cân tại A
    \draw [thick] (B) -- (A) -- (C) -- cycle;
    
    % Vẽ trung điểm M của BC và N của AC
    \draw [thick] (A) -- (M);
    % Vẽ tia Ax song song với BC
    \draw [thick, ->] (A) -- ++(3,0) coordinate (X) node[above] {$y$};
    
    % Vẽ điểm E trên MN
    \draw [thick] (M) -- (E);
    \fill (E) circle (2pt); % Ký hiệu điểm E
    
    % Ký hiệu các điểm khác
    \fill (A) circle (2pt);
    \fill (B) circle (2pt);
    \fill (C) circle (2pt);
    \fill (M) circle (2pt);
    \fill (N) circle (2pt);
\end{tikzpicture}
}
\end{minipage}

\end{tcolorbox}
\vspace{4pt}


\begin{tcolorbox}[
    breakable,
    colback=white,
    colframe=blue!60!black,
    boxrule=0.8pt,
    arc=3pt,
    title={\textbf{Câu \#234} \quad \textbf{MÃ CÂU: \texttt{[L08\_C4\_B2\_05]}} \quad | \quad \textbf{Nhóm: \texttt{HINH\_181}} \quad | \quad \textbf{MapID: \texttt{8H4H2}} \hfill \textcolor{yellow!90!white}{[Hình gốc .pyc]}}
]
\textbf{Dạng toán:} Toán thực tế về đường TB \hfill \textit{(Tự luận -- Thông hiểu)}
\vspace{4pt}\hrule\vspace{6pt}
\begin{minipage}[t]{0.58\textwidth}
\vspace{0pt}
Giữa hai điểm B và C bị ngăn cách bởi hồ nước (như hình vẽ). Xác định độ dài BC mà không cần phải di chuyển qua hồ nước. Biết rằng đoạn thẳng KI dài $19$ m và K là trung điểm của AB, I là trung điểm của AC.
\end{minipage}\hfill
\begin{minipage}[t]{0.39\textwidth}
\vspace{0pt}
\centering

\resizebox{\linewidth}{!}{%
\begin{tikzpicture}[scale=1, font=\footnotesize, line join=round, line cap=round, >=stealth]
            \tikzset{co/.pic={
                    \fill[green!80!teal,out=110,in=-20] (0,0)to(-1,1)[out=-10,in=120]to(-0.2,0.7)[out=110,in=-20]to(0,1.5)[out=-40,in=100]to(0.5,0.5)[out=40,in=-120]to(1,1.3)--(0.75,0)--cycle;
            }}
            \tikzset{Ho/.pic={
                    \fill[green!30](-2.2,1.5)--(5,1.8)..controls++(0:1) and++(-10:1)..(3,1)..controls++(-30:3) and++(-5:1)..(0.3,0.5)..controls++(175:0.5) and++(0:2)..(0,0)..controls++(180:1) and++(-120:1)..(-1.3,0.8)..controls++(60:0.5) and++(-90:1)..(-2.2,1.5);
                    \fill[cyan](-1,0.8)..controls++(20:2) and++(185:1)..(2,1.5)..controls++(5:1) and++(50:0.5)..(3,1.5)..controls++(-130:0.25)and++(15:2.5)..(2.5,1)..controls++(-165:1) and++(-60:0.5)..(-1,0.8);
                    \draw[red](-1,0.8)coordinate(A)--(3.7,1.5)coordinate(B)node[above]{$C$}--(1.5,-1.2)coordinate(C)node[below]{$A$}--(A)node[above]{$B$};
            }}
            \begin{scope}[scale=.75]
                \path(0,0)pic{Ho};
                \foreach \x/\y in{-1.5/1,-0.5/0.1,3/0.5,4/1.2,-0.5/1.2}
                \path(\x,\y)pic[scale=0.25]{co};
                (A)--(B)--(C)--cycle                
                \path ($(A)!.5!(C)$)coordinate(K) node[below]{$K$}
                ($(B)!.5!(C)$)coordinate(I)node[below]{$I$};
                \draw   (C)++(90:0.5) 
                (I)--(K);
                \foreach \y/\g in {A/90,B/40,C/-90,I/-90,K/90}       
                \fill (\y) circle (1pt)+(\g:.5);
            \end{scope}     
    \end{tikzpicture}
}
\end{minipage}

\vspace{3pt}\hrule\vspace{3pt}\footnotesize\textcolor{gray!80!black}{\textbf{Các mã dùng chung hình (HINH\_181):} \texttt{L08\_C4\_B2\_06}}
\end{tcolorbox}
\vspace{4pt}


\begin{tcolorbox}[
    breakable,
    colback=white,
    colframe=blue!60!black,
    boxrule=0.8pt,
    arc=3pt,
    title={\textbf{Câu \#235} \quad \textbf{MÃ CÂU: \texttt{[L08\_C4\_B2\_07]}} \quad | \quad \textbf{Nhóm: \texttt{HINH\_179}} \quad | \quad \textbf{MapID: \texttt{8H4N2}} \hfill \textcolor{yellow!90!white}{[Hình gốc .pyc]}}
]
\textbf{Dạng toán:} Tìm số đo x \hfill \textit{(Tự luận -- Nhận biết)}
\vspace{4pt}\hrule\vspace{6pt}
\begin{minipage}[t]{0.58\textwidth}
\vspace{0pt}
Tìm số đo ${x}$ trong hình
\end{minipage}\hfill
\begin{minipage}[t]{0.39\textwidth}
\vspace{0pt}
\centering

\resizebox{\linewidth}{!}{%
\begin{tikzpicture}[scale=1.5]
    % Vẽ tam giác ABC
    \coordinate [label=below left: $B$] (B) at (0,0);
    \coordinate [label=below right: $C$] (C) at (4,0);
    \coordinate [label=above: $A$] (A) at (4,3);
    
    \draw [thick] (B) -- (C) -- (A) -- cycle;
    
    % Điểm M và N
    \coordinate [label=above left: $M$] (M) at (1.8,1.35);
    \coordinate [label=right: $N$] (N) at (4,1.35);
    
    % Vẽ đoạn thẳng MN
    \draw [thick] (M) -- (N);
    
    % Ký hiệu độ dài
    \node at (3,1.5) {$20$};
    \node at (3,0.1) {$x$};
    
    % Ký hiệu tam giác đồng dạng và độ dài các đoạn thẳng
    \draw [thick] (B) -- (M);
    \draw [thick] (A) -- (N);
        \draw
    ($(B)!0.5!(M)$)node[rotate=20]{$/$} ($(M)!0.5!(A)$)node[rotate=20]{$/$}
    ($(A)!0.5!(N)$)node[rotate=60]{$//$} ($(N)!0.5!(C)$)node[rotate=60]{$//$};
    % Ký hiệu các đoạn bằng nhau

\end{tikzpicture}
}
\end{minipage}

\vspace{3pt}\hrule\vspace{3pt}\footnotesize\textcolor{gray!80!black}{\textbf{Các mã dùng chung hình (HINH\_179):} \texttt{L08\_C4\_B2\_01, L08\_C4\_B2\_08}}
\end{tcolorbox}
\vspace{4pt}


\begin{tcolorbox}[
    breakable,
    colback=white,
    colframe=blue!60!black,
    boxrule=0.8pt,
    arc=3pt,
    title={\textbf{Câu \#236} \quad \textbf{MÃ CÂU: \texttt{[L08\_C4\_B2\_10]}} \quad | \quad \textbf{Nhóm: \texttt{HINH\_180}} \quad | \quad \textbf{MapID: \texttt{8H4H2}} \hfill \textcolor{yellow!90!white}{[Hình gốc .pyc]}}
]
\textbf{Dạng toán:} Chứng mình đường trung bình từ đó tìm số đo x \hfill \textit{(Tự luận -- Thông hiểu)}
\vspace{4pt}\hrule\vspace{6pt}
\begin{minipage}[t]{0.58\textwidth}
\vspace{0pt}
Tìm số đo ${x}$ trong hình
\end{minipage}\hfill
\begin{minipage}[t]{0.39\textwidth}
\vspace{0pt}
\centering

\resizebox{\linewidth}{!}{%
\begin{tikzpicture}
    
    % Coordinates for the vertices
    \coordinate [label=above:$A$] (A) at (0,4);
    \coordinate [label=below left:$B$] (B) at (-2,0);
    \coordinate [label=below right:$C$] (C) at (4,0);
    \coordinate [label=above left:$K$] (K) at (-1,2);
        \coordinate [label=above right:$I$] (I) at (2,2);
    % Draw the triangle
    \draw (A) -- (B) -- (C) -- cycle;
    
    % Draw the segment AK
    \draw (I) -- (K);
    
 % Ký hiệu độ dài
\node at (1,1.6) {$2$};
\node at (1,-0.4) {$x$};
 \draw 
(C) pic[draw,angle radius = 15] {angle = A--C--B}  ;
 \draw 
(I) pic[draw,angle radius = 15] {angle = A--I--K}  ;
\draw
($(A)!0.5!(K)$)node[rotate=40]{$/$} ($(B)!0.5!(K)$)node[rotate=40]{$/$};

\end{tikzpicture}
}
\end{minipage}

\vspace{3pt}\hrule\vspace{3pt}\footnotesize\textcolor{gray!80!black}{\textbf{Các mã dùng chung hình (HINH\_180):} \texttt{L08\_C4\_B2\_02, L08\_C4\_B2\_09}}
\end{tcolorbox}
\vspace{4pt}


\begin{tcolorbox}[
    breakable,
    colback=white,
    colframe=blue!60!black,
    boxrule=0.8pt,
    arc=3pt,
    title={\textbf{Câu \#237} \quad \textbf{MÃ CÂU: \texttt{[L08\_C4\_B2\_03]}} \quad | \quad \textbf{Nhóm: \texttt{HINH\_183}} \quad | \quad \textbf{MapID: \texttt{8H4N2}} \hfill \textcolor{yellow!90!white}{[Hình gốc .pyc]}}
]
\textbf{Dạng toán:} Chỉ ra các cặp đường thẳng song song \hfill \textit{(Tự luận -- Nhận biết)}
\vspace{4pt}\hrule\vspace{6pt}
\begin{minipage}[t]{0.58\textwidth}
\vspace{0pt}
Chỉ ra các cặp đường thẳng song song trong hình.
\end{minipage}\hfill
\begin{minipage}[t]{0.39\textwidth}
\vspace{0pt}
\centering

\resizebox{\linewidth}{!}{%
\begin{tikzpicture}
    
    % Tọa độ các điểm của tam giác
    \coordinate [label=left: {$I$}] (A) at (0, 0);
    \coordinate [label=right:{$J$}] (B) at (4, 0);
    \coordinate [label=above:{$Q$}] (C) at (1, 5);
    % Tạo điểm M thuộc AB (chia đoạn AB theo tỉ lệ 1:2)
    \coordinate [label=below:M] (M) at ($(A)!0.33!(B)$);
    % Tạo điểm N thuộc AC sao cho MN song song với BC
    \coordinate [label=above left:N] (N) at ($(A)!0.33!(C)$);
    % Vẽ các cạnh của tam giác ABC
    \draw (A) -- (B) -- (C) -- cycle;
    % Vẽ đoạn thẳng MN song song với BC
    \draw (M) -- (N);
     \fill (A) circle (1.2pt);
    \fill (B) circle (1.2pt);
    \fill (C) circle (1.2pt);
    \fill (M) circle (1.2pt);
    \fill (N) circle (1.2pt);
      \draw[thick] (A) -- (N) node[midway, left] {$a$};
        \draw[thick] (C) -- (N) node[midway, left] {$2a$};
          \draw[thick] (A) -- (M) node[midway, below] {$b$};
            \draw[thick] (M) -- (B) node[midway, below] {$2b$};
            % Tọa độ các điểm của tam giác thứ hai A'B'C' (bên phải tam giác ABC)
            \coordinate [label=left:{$I'$}] (A') at (6, 0);
            \coordinate [label=right:{$J'$}] (B') at (10, 0);
            \coordinate [label=above:{$Q'$}] (C') at (12, 3);
            
            % Vẽ tam giác thứ hai A'B'C'
            \draw (A') -- (B') -- (C') -- cycle;
            
            % Tạo điểm M' thuộc A'B' và N' thuộc A'C' sao cho M'N' song song với B'C'
            \coordinate [label=below right:M'] (M') at ($(A')!0.33!(B')$);
            \coordinate [label=above:N'] (N') at ($(A')!0.33!(C')$);
            \draw (M') -- (N'); % Vẽ đoạn thẳng M'N'
            
            % Vẽ đoạn thẳng MN song song với BC
            \draw (M) -- (N);
            \fill (A') circle (1.2pt);
            \fill (B') circle (1.2pt);
            \fill (C') circle (1.2pt);
            \fill (M') circle (1.2pt);
            \fill (N') circle (1.2pt);
            \draw[thick] (A') -- (N') node[midway, above] {$6$};
            \draw[thick] (A') -- (M') node[midway, below] {$5$};
            \draw[dashed] (A') to[bend left] node[midway, above, sloped] {$18$} (C');
              \draw[dashed] (A') to[bend right] node[midway, below] {$15$} (B');
              
              % Tọa độ các điểm của tam giác thứ ba A''B''C'' (bên phải tam giác thứ hai)
              \coordinate [label=left:A''] (A'') at (12, 0);
              \coordinate [label=right:B''] (B'') at (16, 0);
              \coordinate [label=above:C''] (C'') at (14, 3);
              
              % Vẽ tam giác thứ ba A''B''C''
              \draw (A'') -- (B'') -- (C'') -- cycle;
              
              % Tạo điểm M'' thuộc A''B'' và N'' thuộc A''C'' sao cho M''N'' song song với B''C''
              \coordinate [label=below:M''] (M'') at ($(A'')!0.33!(B'')$);
              \coordinate [label=above left:N''] (N'') at ($(A'')!0.33!(C'')$);
              \draw (M'') -- (N''); % Vẽ đoạn thẳng M''N''
              \draw[thick] (A'') -- (N'') node[midway, above] {$a$};
              \draw[thick] (N'') -- (C'') node[midway, above left] {$ 2a$};
               \draw[thick] (A'') -- (M'') node[midway, below] {$c$};
              % Vẽ cung nét đứt ký hiệu độ dài từ A'' đến B'' (nằm dưới tam giác A''B''C'')
              \draw[dashed] (A'') to[bend right] node[midway, below] {$3c$} (B'');
                \fill (A'') circle (1.2pt);
              \fill (B'') circle (1.2pt);
              \fill (C'') circle (1.2pt);
              \fill (M'') circle (1.2pt);
              \fill (N'') circle (1.2pt);
\end{tikzpicture}
}
\end{minipage}

\end{tcolorbox}
\vspace{4pt}


\begin{tcolorbox}[
    breakable,
    colback=white,
    colframe=blue!60!black,
    boxrule=0.8pt,
    arc=3pt,
    title={\textbf{Câu \#238} \quad \textbf{MÃ CÂU: \texttt{[L08\_C4\_B2\_11]}} \quad | \quad \textbf{Nhóm: \texttt{HINH\_184}} \quad | \quad \textbf{MapID: \texttt{8H4H2}} \hfill \textcolor{yellow!90!white}{[Hình gốc .pyc]}}
]
\textbf{Dạng toán:} Toán thực tế đường TB 2 \hfill \textit{(Tự luận -- Thông hiểu)}
\vspace{4pt}\hrule\vspace{6pt}
\begin{minipage}[t]{0.58\textwidth}
\vspace{0pt}
Ông An muốn làm một tấm phủ che hồ bơi. Ông cần tính độ rộng của hồ bơi (đoạn ${AB}$) nhưng ông chỉ đo được độ dài đoạn $EF$ là $5,5$ m. Biết ${E, F}$ lần lượt là trung điểm của ${AC}$ và ${BC}$. Em hãy tính độ rộng của hồ bơi giúp ông An.
\end{minipage}\hfill
\begin{minipage}[t]{0.39\textwidth}
\vspace{0pt}
\centering

\resizebox{\linewidth}{!}{%
\begin{tikzpicture}[scale=0.6, font=\footnotesize, line join=round, line cap=round, >=stealth]
            \def \bk{2}
            \path 
                (-3,-3) coordinate (I)
                (3,3) coordinate (J)
                ($(I)+(-135:\bk)$) coordinate (B)
                ($(J)+(45:\bk)$) coordinate (A)
                (3,-3) coordinate (C)
                ($(A)!0.5!(C)$) coordinate (E)
                ($(B)!0.5!(C)$) coordinate (F)
            ;
            \fill[cyan!50,rotate=45] (0,0) ellipse (4 and 2);
            \fill[cyan!50] (J) circle (\bk);
            \fill[red] (2.7,4.3) rectangle (3.3,4.35); 
            \fill[red] (2.7,4.5) rectangle (3.3,4.55); 
            \fill[red] (2.7,4.7) rectangle (3.3,4.75);
            \fill[red] (2.7,4.9) rectangle (3.3,4.95);  
            \fill[red] (2.65,4.3) rectangle (2.7,5.3);
            \fill[red] (3.3,4.3) rectangle (3.35,5.3);
            \fill[cyan!50] (-3,-3) circle (\bk);
            \fill[white] (I) circle (0.5);
            \fill[white] (J) circle (0.5);
            \draw (A) -- (B) -- (C) -- (A) (E) -- (F);
            \foreach \x/\g in {A/00,B/180,C/0,E/0,F/-90} \fill[black](\x) circle (1pt)+(\g:3mm) node{$\x$};
%           \draw[white,thick] (0,0) to[out=-70,in=130] (0.5,0);
        \end{tikzpicture}
}
\end{minipage}

\end{tcolorbox}
\vspace{4pt}


\begin{tcolorbox}[
    breakable,
    colback=white,
    colframe=blue!60!black,
    boxrule=0.8pt,
    arc=3pt,
    title={\textbf{Câu \#239} \quad \textbf{MÃ CÂU: \texttt{[L08\_C4\_B3\_01]}} \quad | \quad \textbf{Nhóm: \texttt{HINH\_185}} \quad | \quad \textbf{MapID: \texttt{8H4H3}} \hfill \textcolor{yellow!90!white}{[Hình gốc .pyc]}}
]
\textbf{Dạng toán:} Sử dụng tính chất đường phân giác tìm x \hfill \textit{(Trắc nghiệm -- Thông hiểu)}
\vspace{4pt}\hrule\vspace{6pt}
\begin{minipage}[t]{0.58\textwidth}
\vspace{0pt}
Tìm số đo ${x}$ trong hình vẽ sau:

\vspace{4pt}\noindent
\textbf{A.} ${\dfrac{16}{3}}$ \quad \textbf{\textcolor{blue!75!black}{B. (Đ)}} $12$ \quad \textbf{C.} $3$ \quad \textbf{D.} $13$
\end{minipage}\hfill
\begin{minipage}[t]{0.39\textwidth}
\vspace{0pt}
\centering

\resizebox{\linewidth}{!}{%
\begin{tikzpicture}
        % Coordinates
        \coordinate [label=below left:$C$] (C) at (0,0);
        \coordinate [label=below right:$A$] (A) at (8,0);
        \coordinate [label=above:$B$] (B) at (7,5);
        \coordinate [label=below:$D$] (D) at (5,0);
        % Draw the triangle
        \draw[thick] (D) -- (A) -- (B) -- cycle;
        
        % Draw the altitude and mark lengths
        \draw[thick] (B) -- (D);
        \draw[thick] (D) -- (A) node[midway, below] {$4$};
        \draw[thick] (B) -- (A) node[midway, right] {$6$};
            \draw[thick] (B) -- (C);
                \draw[thick] (A) -- (C);
        % Add additional labels and lengths
        \coordinate [label=below:{$8$}] (E) at (3,0);
        \draw[thick] (C) -- (B) node[midway, left] {$x$};
         \draw 
        (B) pic[draw,angle radius = 15] {angle = D--B--A} 
        ;
         \draw 
        (B) pic[draw,angle radius = 16] {angle = C--B--D} 
        ;
        
    \end{tikzpicture}
}
\end{minipage}

\end{tcolorbox}
\vspace{4pt}


\subsection{Chương 7 - Phương trình bậc nhất và Hàm số bậc nhất}

\begin{tcolorbox}[
    breakable,
    colback=white,
    colframe=blue!60!black,
    boxrule=0.8pt,
    arc=3pt,
    title={\textbf{Câu \#240} \quad \textbf{MÃ CÂU: \texttt{[L08\_C7\_B1\_05]}} \quad | \quad \textbf{Nhóm: \texttt{HINH\_186}} \quad | \quad \textbf{MapID: \texttt{8D7H1}} \hfill \textcolor{yellow!90!white}{[Hình gốc .pyc]}}
]
\textbf{Dạng toán:} Tìm x để tứ giác là hình vuông \hfill \textit{(Trắc nghiệm -- Thông hiểu)}
\vspace{4pt}\hrule\vspace{6pt}
\begin{minipage}[t]{0.58\textwidth}
\vspace{0pt}
Biết tứ giác ${ABCD}$ là hình vuông. Tìm x?

\vspace{4pt}\noindent
\textbf{\textcolor{blue!75!black}{A. (Đ)}} $3$ \quad \textbf{B.} $5$ \quad \textbf{C.} $7$ \quad \textbf{D.} $2$
\end{minipage}\hfill
\begin{minipage}[t]{0.39\textwidth}
\vspace{0pt}
\centering

\resizebox{\linewidth}{!}{%
\begin{tikzpicture}
        % Các đỉnh của hình vuông
        \coordinate (A) at (0, 4);
        \coordinate (B) at (4, 4);
        \coordinate (C) at (4, 0);
        \coordinate (D) at (0, 0);
        
        % Vẽ hình vuông
        \draw [thick] (A) -- (B) -- (C) -- (D) -- cycle;
        
        % Ghi nhãn các đỉnh
        \node at (-0.3, 4) {$A$};
        \node at (4.3, 4) {$B$};
        \node at (4.3, 0) {$C$};
        \node at (-0.3, 0) {$D$};
        
        % Ghi nhãn các cạnh
        \node at (2, 4.5) {$1x+9$};
        \node at (4.7, 2) {$14x+-30$};
    
    \end{tikzpicture}
}
\end{minipage}

\end{tcolorbox}
\vspace{4pt}


\begin{tcolorbox}[
    breakable,
    colback=white,
    colframe=blue!60!black,
    boxrule=0.8pt,
    arc=3pt,
    title={\textbf{Câu \#241} \quad \textbf{MÃ CÂU: \texttt{[L08\_C7\_B1\_06]}} \quad | \quad \textbf{Nhóm: \texttt{HINH\_187}} \quad | \quad \textbf{MapID: \texttt{8D7V1}} \hfill \textcolor{yellow!90!white}{[Hình gốc .pyc]}}
]
\textbf{Dạng toán:} Tìm x để chu vi hình chữ nhật và tam giác bằng nhau \hfill \textit{(Trắc nghiệm -- Vận dụng)}
\vspace{4pt}\hrule\vspace{6pt}
\begin{minipage}[t]{0.58\textwidth}
\vspace{0pt}
Biết chu vi tam giác và hình chữ nhật bằng nhau. Tìm x?

\vspace{4pt}\noindent
\textbf{A.} $5$ \quad \textbf{B.} $6$ \quad \textbf{\textcolor{blue!75!black}{C. (Đ)}} $3$ \quad \textbf{D.} $2$
\end{minipage}\hfill
\begin{minipage}[t]{0.39\textwidth}
\vspace{0pt}
\centering

\resizebox{\linewidth}{!}{%
\begin{tikzpicture}
    
    % Vẽ tam giác
    \coordinate (A) at (-1, 0);
    \coordinate (B) at (3.5, 0);
    \coordinate (C) at (0.5, 2);
    
    % Vẽ tam giác
    \draw [thick] (A) -- (B) -- (C) -- cycle;
    
    % Ghi nhãn các cạnh của tam giác
    \node at (1.3, -0.3) {$2 x + 5$};
    \node at (-0.7, 1.4) {$x + 1$};
    \node at (2.4, 1.4) {$4 x + 63$};
    
    % Vẽ hình chữ nhật
    \coordinate (D) at (5, 0);
    \coordinate (E) at (9, 0);
    \coordinate (F) at (9, 2);
    \coordinate (G) at (5, 2);
    
    % Vẽ hình chữ nhật
    \draw [thick] (D) -- (E) -- (F) -- (G) -- cycle;
    
    % Ghi nhãn các cạnh của hình chữ nhật
    \node at (7, -0.3) {$7 x + 5$};
    \node at (10, 1) {$4 x + 7$};
\end{tikzpicture}
}
\end{minipage}

\end{tcolorbox}
\vspace{4pt}


\begin{tcolorbox}[
    breakable,
    colback=white,
    colframe=blue!60!black,
    boxrule=0.8pt,
    arc=3pt,
    title={\textbf{Câu \#242} \quad \textbf{MÃ CÂU: \texttt{[L08\_C7\_B3\_01]}} \quad | \quad \textbf{Nhóm: \texttt{HINH\_188}} \quad | \quad \textbf{MapID: \texttt{8D7N3}} \hfill \textcolor{yellow!90!white}{[Hình gốc .pyc]}}
]
\textbf{Dạng toán:} Đọc toạ độ điểm \hfill \textit{(Trắc nghiệm -- Nhận biết)}
\vspace{4pt}\hrule\vspace{6pt}
\begin{minipage}[t]{0.58\textwidth}
\vspace{0pt}
Toạ độ điểm N là

\vspace{4pt}\noindent
\textbf{A.} $(-4;5)$ \quad \textbf{B.} $(-1;3)$ \quad \textbf{C.} $(-3;3)$ \quad \textbf{\textcolor{blue!75!black}{D. (Đ)}} $(-2;4)$
\end{minipage}\hfill
\begin{minipage}[t]{0.39\textwidth}
\vspace{0pt}
\centering

\resizebox{\linewidth}{!}{%
\begin{tikzpicture}[line join=round, line cap=round,>=stealth,thick,scale=0.6]
    \tikzset{every node/.style={scale=0.9}}
    \draw[gray!20](-6,-6.5)grid(6,6.5);
    \draw[->] (-6,0)--(6,0) node[below left] {$x$};
    \draw[->] (0,-6.5)--(0,6.5) node[below left] {$y$};
    \draw (0,0) node [below left] {\footnotesize $O$};
     \foreach \x in {-5,-4,-3,-2,-1,1,2,3,4,5}
    \draw[thin] (\x,1pt)--(\x,-1pt) node [below] {\footnotesize$\x$};
    \foreach \y in {-4,-3,-2,-1,1,2,3,4}
    \draw[thin] (1pt,\y)--(-1pt,\y) node [left] {\footnotesize$\y$};
    \coordinate (N) at (-2,4);
    \foreach \i/\g in {N/90}{\draw[fill=white](\i) circle (1.5pt) ($(\i)+(\g:3mm)$) node[scale=1]{$\i$};}
    \clip (-6.5,-6.5) rectangle (6.5,6.5);
    \end{tikzpicture}
}
\end{minipage}

\end{tcolorbox}
\vspace{4pt}


\subsection{Chương 9 -  Tam giác đồng dạng}

\begin{tcolorbox}[
    breakable,
    colback=white,
    colframe=blue!60!black,
    boxrule=0.8pt,
    arc=3pt,
    title={\textbf{Câu \#243} \quad \textbf{MÃ CÂU: \texttt{[L08\_C9\_B1\_08]}} \quad | \quad \textbf{Nhóm: \texttt{HINH\_189}} \quad | \quad \textbf{MapID: \texttt{8H9N1}} \hfill \textcolor{yellow!90!white}{[Hình gốc .pyc]}}
]
\textbf{Dạng toán:} Cho 2 tam giác đồng dạng, viết các cặp góc, tỉ só các cạnh tương ứng \hfill \textit{(Trắc nghiệm -- Nhận biết)}
\vspace{4pt}\hrule\vspace{6pt}
\begin{minipage}[t]{0.58\textwidth}
\vspace{0pt}
Trong hình bên, cho biết: $\Delta {BHI} \backsim \Delta {MNO}$. Chọn khẳng định đúng trong các khẳng định sau.

\vspace{4pt}\noindent
\textbf{\textcolor{blue!75!black}{A. (Đ)}} $\dfrac{BH}{MN}= \dfrac{HI}{NO}=\dfrac{BI}{MO}$ \quad \textbf{B.} $\dfrac{BH}{ON}= \dfrac{HI}{NM}=\dfrac{BI}{MO}$ \quad \textbf{C.} $\dfrac{BH}{MN}= \dfrac{HI}{MO}=\dfrac{BI}{NO}$ \quad \textbf{D.} $\dfrac{BI}{MN}= \dfrac{HI}{NO}=\dfrac{BH}{MO}$
\end{minipage}\hfill
\begin{minipage}[t]{0.39\textwidth}
\vspace{0pt}
\centering

\resizebox{\linewidth}{!}{%
\begin{tikzpicture}[scale=1.5, font=\footnotesize, line join=round, line cap=round, >=stealth]
                    \path
                    (1.4,1.25) coordinate (A)
                    (0,0) coordinate (B)
                    (2,0) coordinate (C)
                    ;
                    \draw [line width = .3pt,blue] (A)--(B)--(C)--cycle;
                     \filldraw[black] (A) circle (1pt) node[above right] {$B$};
    \filldraw[black] (B) circle (1pt) node[below left] {$H$};
    \filldraw[black] (C) circle (1pt) node[below right] {$I$};
  
                \end{tikzpicture}
                \begin{tikzpicture}[scale=1, font=\footnotesize, line join=round, line cap=round, >=stealth]
                    \path
                    (2.8,2.5) coordinate (D)
                    (0,0) coordinate (E)
                    (4,0) coordinate (F)
                    ;
                    \draw [line width = .3pt,blue] (D)--(E)--(F)--cycle;
                     \filldraw[black] (D) circle (1pt) node[above right]{$M$};
    \filldraw[black] (E) circle (1pt) node[below left] {$N$};
    \filldraw[black] (F) circle (1pt) node[below right] {$O$};
                \end{tikzpicture}
}
\end{minipage}

\vspace{3pt}\hrule\vspace{3pt}\footnotesize\textcolor{gray!80!black}{\textbf{Các mã dùng chung hình (HINH\_189):} \texttt{L08\_C9\_B1\_02}}
\end{tcolorbox}
\vspace{4pt}


\begin{tcolorbox}[
    breakable,
    colback=white,
    colframe=blue!60!black,
    boxrule=0.8pt,
    arc=3pt,
    title={\textbf{Câu \#244} \quad \textbf{MÃ CÂU: \texttt{[L08\_C9\_B1\_09]}} \quad | \quad \textbf{Nhóm: \texttt{HINH\_190}} \quad | \quad \textbf{MapID: \texttt{8H9H1}} \hfill \textcolor{yellow!90!white}{[Hình gốc .pyc]}}
]
\textbf{Dạng toán:} Cho 2 tam giác đồng dạng, Tính góc \hfill \textit{(Trắc nghiệm -- Thông hiểu)}
\vspace{4pt}\hrule\vspace{6pt}
\begin{minipage}[t]{0.58\textwidth}
\vspace{0pt}
Trong hình bên, cho biết $\Delta {BHI} \backsim \Delta KMO$. Tính số đo góc $\widehat{K}$, góc $\widehat{O}$.

\vspace{4pt}\noindent
\textbf{A.} $\widehat{K}=42^{\circ} ; \widehat{O}=84^{\circ }$ \quad \textbf{B.} $\widehat{K}=38^{\circ} ; \widehat{O}=80^{\circ }$ \quad \textbf{C.} $\widehat{K}=42^{\circ} ; \widehat{O}=82^{\circ }$ \quad \textbf{\textcolor{blue!75!black}{D. (Đ)}} $\widehat{K}=40^{\circ} ; \widehat{O}=82^{\circ }$
\end{minipage}\hfill
\begin{minipage}[t]{0.39\textwidth}
\vspace{0pt}
\centering

\resizebox{\linewidth}{!}{%
\begin{tikzpicture}[scale=1.2, font=\footnotesize, line join=round, line cap=round, >=stealth]
                    \path
                    (70:2.5) coordinate (D)
                    (0,0) coordinate (M)
                    (1,0) coordinate (E)
                    ;
                    \node (B) at ($(D)!1!57:(M)$){};
                    \path
                    (intersection of D--B and M--E) coordinate (N)
                    ;
                    \draw [line width = .3pt,blue] (D)--(M)--(N)--cycle;
                   \filldraw[black] (D) circle (1pt) node[above right]{$K$};
    \filldraw[black] (M) circle (1pt) node[below left] {$M$};
    \filldraw[black] (N) circle (1pt) node[below right] {$O$};
                \end{tikzpicture}
                \begin{tikzpicture}[scale=1, font=\footnotesize, line join=round, line cap=round, >=stealth]
                    \path
                    (70:4) coordinate (A) node[xshift=1.4mm,yshift=-6mm]{$40^\circ$}
                    (0,0) coordinate (B) node[xshift=5mm,yshift=2.25mm]{$58^\circ$}
                    (1,0) coordinate (E)
                    ;
                    \node (D) at ($(A)!1!57:(B)$){};
                    \path
                    (intersection of A--D and B--E) coordinate (C)
                    ;
                    \draw [line width = .3pt,blue] (A)--(B)--(C)--cycle;
                      
            \filldraw[black] (A) circle (1pt) node[above right] {$B$};
    \filldraw[black] (B) circle (1pt) node[below left] {$H$};
    \filldraw[black] (C) circle (1pt) node[below right] {$I$};
                \end{tikzpicture}
}
\end{minipage}

\vspace{3pt}\hrule\vspace{3pt}\footnotesize\textcolor{gray!80!black}{\textbf{Các mã dùng chung hình (HINH\_190):} \texttt{L08\_C9\_B1\_03}}
\end{tcolorbox}
\vspace{4pt}


\begin{tcolorbox}[
    breakable,
    colback=white,
    colframe=blue!60!black,
    boxrule=0.8pt,
    arc=3pt,
    title={\textbf{Câu \#245} \quad \textbf{MÃ CÂU: \texttt{[L08\_C9\_B1\_10]}} \quad | \quad \textbf{Nhóm: \texttt{HINH\_191}} \quad | \quad \textbf{MapID: \texttt{8H9H1}} \hfill \textcolor{yellow!90!white}{[Hình gốc .pyc]}}
]
\textbf{Dạng toán:} Cho 2 tam giác đồng dạng, Tính cạnh \hfill \textit{(Trắc nghiệm -- Thông hiểu)}
\vspace{4pt}\hrule\vspace{6pt}
\begin{minipage}[t]{0.58\textwidth}
\vspace{0pt}
Trong hình bên, cho biết $\Delta {EFH} \backsim \Delta {KMQ}$. Tính độ dài các đoạn thẳng ${EF}$, ${KQ}$.

\vspace{4pt}\noindent
\textbf{\textcolor{blue!75!black}{A. (Đ)}} ${EF=18; KQ=64}$ \quad \textbf{B.} ${EF=22; KQ=63}$ \quad \textbf{C.} ${EF=20; KQ=62}$ \quad \textbf{D.} ${EF=16; KQ=62}$
\end{minipage}\hfill
\begin{minipage}[t]{0.39\textwidth}
\vspace{0pt}
\centering

\resizebox{\linewidth}{!}{%
\begin{tikzpicture}[scale=1.2, font=\footnotesize, line join=round, line cap=round, >=stealth]
   
                    \path
                    (-2.4,1.2) coordinate (A)
                    (0,0) coordinate (B1)
                    (0.5,1.5) coordinate (C)
                    ($(B1)!(A)!(C)$) coordinate (B)
                    ($(A)!0.5!(C)$)node[above]{$16$}
                    ($(C)!0.5!(B)$)node[xshift=3mm, yshift=-1mm]{$13$}
                    ;

                    \draw [line width = .3pt,blue] (A)--(B)--(C)--cycle;
                    \filldraw[black] (A) circle (1pt) node[above right] {$E$};
    \filldraw[black] (B) circle (1pt) node[below left] {$F$};
    \filldraw[black] (C) circle (1pt) node[below right] {$H$};
                \end{tikzpicture}
                \begin{tikzpicture}[scale=1, font=\footnotesize, line join=round, line cap=round, >=stealth]
                    \path
                    (0,4) coordinate (D)
                    (0,0) coordinate (E)
                    (3.2,0) coordinate (F)
                    ($(D)!0.5!(E)$)node[xshift=-3mm]{$72$}
                    ($(F)!0.5!(E)$)node[below]{$52$}
                    ;
                    \draw [line width = .3pt,blue] (D)--(E)--(F)--cycle;
                   \filldraw[black] (D) circle (1pt) node[above right]{$K$};
    \filldraw[black] (E) circle (1pt) node[below left] {$M$};
    \filldraw[black] (F) circle (1pt) node[below right] {$Q$};
                \end{tikzpicture}
}
\end{minipage}

\vspace{3pt}\hrule\vspace{3pt}\footnotesize\textcolor{gray!80!black}{\textbf{Các mã dùng chung hình (HINH\_191):} \texttt{L08\_C9\_B1\_04}}
\end{tcolorbox}
\vspace{4pt}


\begin{tcolorbox}[
    breakable,
    colback=white,
    colframe=blue!60!black,
    boxrule=0.8pt,
    arc=3pt,
    title={\textbf{Câu \#246} \quad \textbf{MÃ CÂU: \texttt{[L08\_C9\_B1\_02]}} \quad | \quad \textbf{Nhóm: \texttt{HINH\_189}} \quad | \quad \textbf{MapID: \texttt{8H9N1}} \hfill \textcolor{yellow!90!white}{[Hình gốc .pyc]}}
]
\textbf{Dạng toán:} Cho 2 tam giác đồng dạng, viết các cặp góc, tỉ só các cạnh tương ứng \hfill \textit{(Tự luận -- Nhận biết)}
\vspace{4pt}\hrule\vspace{6pt}
\begin{minipage}[t]{0.58\textwidth}
\vspace{0pt}
Trong hình bên, cho biết: $\Delta {CEF} \backsim \Delta {GJM}$. Viết các tỉ số của các cạnh tương ứng và chỉ ra các cặp góc tương ứng.
\end{minipage}\hfill
\begin{minipage}[t]{0.39\textwidth}
\vspace{0pt}
\centering

\resizebox{\linewidth}{!}{%
\begin{tikzpicture}[scale=1.5, font=\footnotesize, line join=round, line cap=round, >=stealth]
                    \path
                    (1.4,1.25) coordinate (A)
                    (0,0) coordinate (B)
                    (2,0) coordinate (C)
                    ;
                    \draw [line width = .3pt,blue] (A)--(B)--(C)--cycle;
                     \filldraw[black] (A) circle (1pt) node[above right] {$C$};
    \filldraw[black] (B) circle (1pt) node[below left] {$E$};
    \filldraw[black] (C) circle (1pt) node[below right] {$F$};
  
                \end{tikzpicture}
                \begin{tikzpicture}[scale=1, font=\footnotesize, line join=round, line cap=round, >=stealth]
                    \path
                    (2.8,2.5) coordinate (D)
                    (0,0) coordinate (E)
                    (4,0) coordinate (F)
                    ;
                    \draw [line width = .3pt,blue] (D)--(E)--(F)--cycle;
                     \filldraw[black] (D) circle (1pt) node[above right]{$G$};
    \filldraw[black] (E) circle (1pt) node[below left] {$J$};
    \filldraw[black] (F) circle (1pt) node[below right] {$M$};
                \end{tikzpicture}
}
\end{minipage}

\vspace{3pt}\hrule\vspace{3pt}\footnotesize\textcolor{gray!80!black}{\textbf{Các mã dùng chung hình (HINH\_189):} \texttt{L08\_C9\_B1\_08}}
\end{tcolorbox}
\vspace{4pt}


\begin{tcolorbox}[
    breakable,
    colback=white,
    colframe=blue!60!black,
    boxrule=0.8pt,
    arc=3pt,
    title={\textbf{Câu \#247} \quad \textbf{MÃ CÂU: \texttt{[L08\_C9\_B1\_03]}} \quad | \quad \textbf{Nhóm: \texttt{HINH\_190}} \quad | \quad \textbf{MapID: \texttt{8H9H1}} \hfill \textcolor{yellow!90!white}{[Hình gốc .pyc]}}
]
\textbf{Dạng toán:} Cho 2 tam giác đồng dạng, Tính góc \hfill \textit{(Tự luận -- Thông hiểu)}
\vspace{4pt}\hrule\vspace{6pt}
\begin{minipage}[t]{0.58\textwidth}
\vspace{0pt}
Trong hình bên, cho biết $\Delta {BDK} \backsim \Delta LMN$. Tính số đo góc $\widehat{L}$, góc $\widehat{N}$.
\end{minipage}\hfill
\begin{minipage}[t]{0.39\textwidth}
\vspace{0pt}
\centering

\resizebox{\linewidth}{!}{%
\begin{tikzpicture}[scale=1.2, font=\footnotesize, line join=round, line cap=round, >=stealth]
                    \path
                    (70:2.5) coordinate (D)
                    (0,0) coordinate (M)
                    (1,0) coordinate (E)
                    ;
                    \node (B) at ($(D)!1!57:(M)$){};
                    \path
                    (intersection of D--B and M--E) coordinate (N)
                    ;
                    \draw [line width = .3pt,blue] (D)--(M)--(N)--cycle;
                   \filldraw[black] (D) circle (1pt) node[above right]{$L$};
    \filldraw[black] (M) circle (1pt) node[below left] {$M$};
    \filldraw[black] (N) circle (1pt) node[below right] {$N$};
                \end{tikzpicture}
                \begin{tikzpicture}[scale=1, font=\footnotesize, line join=round, line cap=round, >=stealth]
                    \path
                    (70:4) coordinate (A) node[xshift=1.4mm,yshift=-6mm]{$41^\circ$}
                    (0,0) coordinate (B) node[xshift=5mm,yshift=2.25mm]{$60^\circ$}
                    (1,0) coordinate (E)
                    ;
                    \node (D) at ($(A)!1!57:(B)$){};
                    \path
                    (intersection of A--D and B--E) coordinate (C)
                    ;
                    \draw [line width = .3pt,blue] (A)--(B)--(C)--cycle;
                      
            \filldraw[black] (A) circle (1pt) node[above right] {$B$};
    \filldraw[black] (B) circle (1pt) node[below left] {$D$};
    \filldraw[black] (C) circle (1pt) node[below right] {$K$};
                \end{tikzpicture}
}
\end{minipage}

\vspace{3pt}\hrule\vspace{3pt}\footnotesize\textcolor{gray!80!black}{\textbf{Các mã dùng chung hình (HINH\_190):} \texttt{L08\_C9\_B1\_09}}
\end{tcolorbox}
\vspace{4pt}


\begin{tcolorbox}[
    breakable,
    colback=white,
    colframe=blue!60!black,
    boxrule=0.8pt,
    arc=3pt,
    title={\textbf{Câu \#248} \quad \textbf{MÃ CÂU: \texttt{[L08\_C9\_B1\_04]}} \quad | \quad \textbf{Nhóm: \texttt{HINH\_191}} \quad | \quad \textbf{MapID: \texttt{8H9H1}} \hfill \textcolor{yellow!90!white}{[Hình gốc .pyc]}}
]
\textbf{Dạng toán:} Cho 2 tam giác đồng dạng, Tính cạnh \hfill \textit{(Tự luận -- Thông hiểu)}
\vspace{4pt}\hrule\vspace{6pt}
\begin{minipage}[t]{0.58\textwidth}
\vspace{0pt}
Trong hình bên, cho biết $\Delta {CGH} \backsim \Delta {JLN}$. Tính độ dài các đoạn thẳng ${CG}$, ${JN}$.
\end{minipage}\hfill
\begin{minipage}[t]{0.39\textwidth}
\vspace{0pt}
\centering

\resizebox{\linewidth}{!}{%
\begin{tikzpicture}[scale=1.2, font=\footnotesize, line join=round, line cap=round, >=stealth]
   
                    \path
                    (-2.4,1.2) coordinate (A)
                    (0,0) coordinate (B1)
                    (0.5,1.5) coordinate (C)
                    ($(B1)!(A)!(C)$) coordinate (B)
                    ($(A)!0.5!(C)$)node[above]{$17$}
                    ($(C)!0.5!(B)$)node[xshift=3mm, yshift=-1mm]{$12$}
                    ;

                    \draw [line width = .3pt,blue] (A)--(B)--(C)--cycle;
                    \filldraw[black] (A) circle (1pt) node[above right] {$C$};
    \filldraw[black] (B) circle (1pt) node[below left] {$G$};
    \filldraw[black] (C) circle (1pt) node[below right] {$H$};
                \end{tikzpicture}
                \begin{tikzpicture}[scale=1, font=\footnotesize, line join=round, line cap=round, >=stealth]
                    \path
                    (0,4) coordinate (D)
                    (0,0) coordinate (E)
                    (3.2,0) coordinate (F)
                    ($(D)!0.5!(E)$)node[xshift=-3mm]{$88$}
                    ($(F)!0.5!(E)$)node[below]{$48$}
                    ;
                    \draw [line width = .3pt,blue] (D)--(E)--(F)--cycle;
                   \filldraw[black] (D) circle (1pt) node[above right]{$J$};
    \filldraw[black] (E) circle (1pt) node[below left] {$L$};
    \filldraw[black] (F) circle (1pt) node[below right] {$N$};
                \end{tikzpicture}
}
\end{minipage}

\vspace{3pt}\hrule\vspace{3pt}\footnotesize\textcolor{gray!80!black}{\textbf{Các mã dùng chung hình (HINH\_191):} \texttt{L08\_C9\_B1\_10}}
\end{tcolorbox}
\vspace{4pt}


\begin{tcolorbox}[
    breakable,
    colback=white,
    colframe=blue!60!black,
    boxrule=0.8pt,
    arc=3pt,
    title={\textbf{Câu \#249} \quad \textbf{MÃ CÂU: \texttt{[L08\_C9\_B1\_05]}} \quad | \quad \textbf{Nhóm: \texttt{HINH\_192}} \quad | \quad \textbf{MapID: \texttt{8H9H1}} \hfill \textcolor{yellow!90!white}{[Hình gốc .pyc]}}
]
\textbf{Dạng toán:} Cho biết AB // DE, Cm tam giác đồng dạng, tìm x \hfill \textit{(Tự luận -- Thông hiểu)}
\vspace{4pt}\hrule\vspace{6pt}
\begin{minipage}[t]{0.58\textwidth}
\vspace{0pt}
Trong hình bên, cho biết $AE \parallel KM$.

 a) Chứng minh: $\Delta {AEF} \backsim \Delta KMF$.

b) Tìm ${x}$.
\end{minipage}\hfill
\begin{minipage}[t]{0.39\textwidth}
\vspace{0pt}
\centering

\resizebox{\linewidth}{!}{%
\begin{tikzpicture}[scale=1.2, font=\footnotesize, line join=round, line cap=round, >=stealth]
            \path
            (-1,2) coordinate (A)
            (0,0) coordinate (B)
            (4,2.3) coordinate (E)
            (6,-1.7) coordinate (D)
            (intersection of A--D and B--E) coordinate (C)
            ($(A)!0.5!(B)$)node[xshift=-2mm, yshift=-1mm]{$12$}
            ($(A)!0.5!(C)$)node[above,rotate={atan(-0.5)}]{$x-15$}
            ($(C)!0.5!(D)$)node[xshift=-3mm, yshift=-1mm]{$39$}
            ($(E)!0.5!(D)$)node[xshift=3mm, yshift=1mm]{$36$}
            ;
            \draw [line width = .3pt,blue] (A)--(B)--(E)--(D)--cycle;
               \filldraw[black] (A) circle (1pt) node[above right] {$A$};
    \filldraw[black] (B) circle (1pt) node[below left] {$E$};
    \filldraw[black] (C) circle (1pt) node[above] {$F$};
   \filldraw[black] (D) circle (1pt) node[below right]{$K$};
    \filldraw[black] (E) circle (1pt) node[above left] {$M$};
                \end{tikzpicture}
}
\end{minipage}

\end{tcolorbox}
\vspace{4pt}


\begin{tcolorbox}[
    breakable,
    colback=white,
    colframe=blue!60!black,
    boxrule=0.8pt,
    arc=3pt,
    title={\textbf{Câu \#250} \quad \textbf{MÃ CÂU: \texttt{[L08\_C9\_B1\_07]}} \quad | \quad \textbf{Nhóm: \texttt{HINH\_193}} \quad | \quad \textbf{MapID: \texttt{8H9V1}} \hfill \textcolor{yellow!90!white}{[Hình gốc .pyc]}}
]
\textbf{Dạng toán:} Toán thực tế ứng dụng tam giác đồng dạng \hfill \textit{(Tự luận -- Vận dụng)}
\vspace{4pt}\hrule\vspace{6pt}
\begin{minipage}[t]{0.58\textwidth}
\vspace{0pt}
Người ta ứng dụng hai tam giác đồng dạng để đo khoảng cách ${DE}$ ở hai địa điểm không thể đến được (Hình bên). Biết $FG \parallel DE$.

a) Chứng minh: $\Delta AFG \backsim \Delta {ADE}$.b) Tính khoảng cách ${DE}$. (Làm tròn kết quả đến hàng phần mười)
\end{minipage}\hfill
\begin{minipage}[t]{0.39\textwidth}
\vspace{0pt}
\centering

\resizebox{\linewidth}{!}{%
\begin{tikzpicture}[scale=1, font=\footnotesize, line join=round, line cap=round, >=stealth]
  
    \def\r{3}
    \def\gocC{160}
    \path (0,0) coordinate (O)
    (\gocC:6) coordinate (A)
    (210:\r) coordinate (B)
    (-30:\r) coordinate (C)
    ($(A)!0.5!(B)$) coordinate (D)
    ($(A)!0.5!(C)$) coordinate (E)
    ;
  
    \draw (B)--(C)--(A)--cycle (D)--(E);
    \draw[<->] ($(A)+(-1,0)$)--($(B)+(-1,0)$);
    \draw ($(D)+(-1.2,0)$)node[left,scale=0.7]{$30$ m};
    \draw ($(D)!0.5!(E)$)node[below,scale=0.7]{$16$ m};
    \draw ($(D)!0.5!(A)$)node[right,scale=0.7]{$15$ m};

               \filldraw[black] (A) circle (1pt) node[above right] {$A$};
    \filldraw[black] (B) circle (1pt) node[below left] {$D$};
    \filldraw[black] (C) circle (1pt) node[below right] {$E$};
   \filldraw[black] (D) circle (1pt) node[below left]{$F$};
    \filldraw[black] (E) circle (1pt) node[above right] {$G$};
                \end{tikzpicture}
}
\end{minipage}

\end{tcolorbox}
\vspace{4pt}


\begin{tcolorbox}[
    breakable,
    colback=white,
    colframe=blue!60!black,
    boxrule=0.8pt,
    arc=3pt,
    title={\textbf{Câu \#251} \quad \textbf{MÃ CÂU: \texttt{[L08\_C9\_B2\_01]}} \quad | \quad \textbf{Nhóm: \texttt{HINH\_194}} \quad | \quad \textbf{MapID: \texttt{8H9N2}} \hfill \textcolor{yellow!90!white}{[Hình gốc .pyc]}}
]
\textbf{Dạng toán:} Cho hình vẽ, chỉ ra cặp tam giác đồng dạng theo c.c.c \hfill \textit{(Trắc nghiệm -- Nhận biết)}
\vspace{4pt}\hrule\vspace{6pt}
\begin{minipage}[t]{0.58\textwidth}
\vspace{0pt}
Cho $\Delta {ADG}$ và $\Delta {HKQ}$ có kích thước các cạnh như hình bên. Trong các khẳng định sau khẳng định nào đúng.

\vspace{4pt}\noindent
\textbf{A.} $ \Delta AGD \backsim \Delta HKQ$ \quad \textbf{B.} $ \Delta ADG \backsim \Delta QKH$ \quad \textbf{\textcolor{blue!75!black}{C. (Đ)}} $ \Delta ADG \backsim \Delta HKQ$ \quad \textbf{D.} $ \Delta ADG \backsim \Delta HQK$
\end{minipage}\hfill
\begin{minipage}[t]{0.39\textwidth}
\vspace{0pt}
\centering

\resizebox{\linewidth}{!}{%
\begin{tikzpicture}[scale=1, font=\footnotesize, line join=round, line cap=round, >=stealth]
                \path
                (1,3) coordinate (D)
                (0,0) coordinate (E)
                (4,0) coordinate (F)
                ($(D)!0.5!(E)$)node[xshift=-3mm]{$18$}
                ($(D)!0.5!(F)$)node[xshift=4mm]{$21$}
                ($(F)!0.5!(E)$)node[below]{$27$}
                ;
                \draw [line width = .3pt,blue] (D)--(E)--(F)--cycle;
                \filldraw[black] (D) circle (1pt) node[above right]{$H$};
    \filldraw[black] (E) circle (1pt) node[below left] {$K$};
    \filldraw[black] (F) circle (1pt) node[below right] {$Q$};
            \end{tikzpicture}
            \begin{tikzpicture}[scale=1, font=\footnotesize, line join=round, line cap=round, >=stealth]
                \path
                (0.6,1.8) coordinate (A)
                (0,0) coordinate (B)
                (2.4,0) coordinate (C)
                ($(A)!0.5!(B)$)node[xshift=-3mm]{$6$}
                ($(A)!0.5!(C)$)node[xshift=4mm]{$7$}
                ($(B)!0.5!(C)$)node[below]{$9$}
                ;

                \draw [line width = .3pt,blue] (A)--(B)--(C)--cycle;
                   \filldraw[black] (A) circle (1pt) node[above left] {$A$};
    \filldraw[black] (B) circle (1pt) node[below left] {$D$};
    \filldraw[black] (C) circle (1pt) node[below right] {$G$};
            \end{tikzpicture}
}
\end{minipage}

\vspace{3pt}\hrule\vspace{3pt}\footnotesize\textcolor{gray!80!black}{\textbf{Các mã dùng chung hình (HINH\_194):} \texttt{L08\_C9\_B2\_05}}
\end{tcolorbox}
\vspace{4pt}


\begin{tcolorbox}[
    breakable,
    colback=white,
    colframe=blue!60!black,
    boxrule=0.8pt,
    arc=3pt,
    title={\textbf{Câu \#252} \quad \textbf{MÃ CÂU: \texttt{[L08\_C9\_B2\_02]}} \quad | \quad \textbf{Nhóm: \texttt{HINH\_195}} \quad | \quad \textbf{MapID: \texttt{8H9N2}} \hfill \textcolor{yellow!90!white}{[Hình gốc .pyc]}}
]
\textbf{Dạng toán:} Cho hình vẽ, chỉ ra cặp tam giác đồng dạng theo c.g.c \hfill \textit{(Trắc nghiệm -- Nhận biết)}
\vspace{4pt}\hrule\vspace{6pt}
\begin{minipage}[t]{0.58\textwidth}
\vspace{0pt}
Cho hai tam giác $\Delta ACD$ và $\Delta GHK$ có ${ GH=4}$ cm, ${ GK=5 }$ cm, ${ AC=12}$ cm, ${AD=15}$ cm, $\widehat{A}=\widehat{G}$ (Hình bên). Trong các khẳng định sau khẳng định nào đúng.

\vspace{4pt}\noindent
\textbf{A.} $ \Delta ACD \backsim \Delta GKH$ \quad \textbf{B.} $ \Delta DCA \backsim \Delta GHK$ \quad \textbf{C.} $ \Delta ACD \backsim \Delta KHG$ \quad \textbf{\textcolor{blue!75!black}{D. (Đ)}} $ \Delta ACD \backsim \Delta GHK$
\end{minipage}\hfill
\begin{minipage}[t]{0.39\textwidth}
\vspace{0pt}
\centering

\resizebox{\linewidth}{!}{%
\begin{tikzpicture}[scale=1.2, font=\footnotesize, line join=round, line cap=round, >=stealth]
                \path
                (0.6,1) coordinate (D)
                (0,0) coordinate (E)
                (1.9,0) coordinate (F)
                ($(D)!0.5!(E)$)node[xshift=-2mm, yshift=1mm]{$4$}
                ($(D)!0.5!(F)$)node[xshift=2mm, yshift=1mm]{$5$}
                ;
                \draw [line width = .3pt] (D)--(E)--(F)--cycle;
                \draw pic[draw, angle radius=3mm] {angle=E--D--F};
                \filldraw[black] (D) circle (1pt) node[above right]{$G$};
    \filldraw[black] (E) circle (1pt) node[below left] {$H$};
    \filldraw[black] (F) circle (1pt) node[below] {$K$};
            \end{tikzpicture}
            \begin{tikzpicture}[scale=1.2, font=\footnotesize, line join=round, line cap=round, >=stealth]
                \path
                (4,0) coordinate (A)
                (2.3,2) coordinate (B)
                (0,-1) coordinate (C)
                ($(A)!0.5!(C)$)node[xshift=1mm,yshift=-3mm]{$15$}
                ($(A)!0.5!(B)$)node[xshift=1mm,yshift=2mm]{$12$}
                ;
                \draw [line width = .3pt] (A)--(B)--(C)--cycle;
                \draw pic[draw, angle radius=3mm] {angle=B--A--C};
               
                   \filldraw[black] (A) circle (1pt) node[below right] {$A$};
    \filldraw[black] (B) circle (1pt) node[above] {$C$};
    \filldraw[black] (C) circle (1pt) node[below right] {$D$};
            \end{tikzpicture}
}
\end{minipage}

\end{tcolorbox}
\vspace{4pt}


\begin{tcolorbox}[
    breakable,
    colback=white,
    colframe=blue!60!black,
    boxrule=0.8pt,
    arc=3pt,
    title={\textbf{Câu \#253} \quad \textbf{MÃ CÂU: \texttt{[L08\_C9\_B2\_03]}} \quad | \quad \textbf{Nhóm: \texttt{HINH\_196}} \quad | \quad \textbf{MapID: \texttt{8H9N2}} \hfill \textcolor{yellow!90!white}{[Hình gốc .pyc]}}
]
\textbf{Dạng toán:} Cho hình vẽ, chỉ ra cặp tam giác đồng dạng theo g.g \hfill \textit{(Trắc nghiệm -- Nhận biết)}
\vspace{4pt}\hrule\vspace{6pt}
\begin{minipage}[t]{0.58\textwidth}
\vspace{0pt}
Trong hình bên, cho biết $BG \parallel HL$; $LM \parallel GH$. Trong các khẳng định sau khẳng định nào đúng.

\vspace{4pt}\noindent
\textbf{A.} $ \Delta HGB \backsim \Delta HLM$ \quad \textbf{B.} $ \Delta GBH \backsim \Delta HLM$ \quad \textbf{C.} $ \Delta BHG \backsim \Delta HLM$ \quad \textbf{\textcolor{blue!75!black}{D. (Đ)}} $ \Delta BGH \backsim \Delta HLM$
\end{minipage}\hfill
\begin{minipage}[t]{0.39\textwidth}
\vspace{0pt}
\centering

\resizebox{\linewidth}{!}{%
\begin{tikzpicture}[scale=1.2, font=\footnotesize, line join=round, line cap=round, >=stealth]
            \path
            (1,3) coordinate (A)
            (0,0) coordinate (B)
            (4,0) coordinate (C)
            ($(A)!0.65!(C)$) coordinate (E)
            ($2/5*(A)+(C)$) coordinate (D2)
            ($(E)+(2,0)$) coordinate (D1)
            (intersection of C--D2 and E--D1) coordinate (D)
            ;
            \draw [line width =.5pt,blue] (A)--(B)--(C)--cycle (D)--(C) (D)--(E);
                   \filldraw[black] (A) circle (1pt) node[above] {$B$};
    \filldraw[black] (B) circle (1pt) node[below left] {$G$};
    \filldraw[black] (C) circle (1pt) node[below right] {$H$};
     \filldraw[black] (D) circle (1pt) node[above right]{$L$};
    \filldraw[black] (E) circle (1pt) node[below left] {$M$};
            \end{tikzpicture}
}
\end{minipage}

\vspace{3pt}\hrule\vspace{3pt}\footnotesize\textcolor{gray!80!black}{\textbf{Các mã dùng chung hình (HINH\_196):} \texttt{L08\_C9\_B2\_04}}
\end{tcolorbox}
\vspace{4pt}


\begin{tcolorbox}[
    breakable,
    colback=white,
    colframe=blue!60!black,
    boxrule=0.8pt,
    arc=3pt,
    title={\textbf{Câu \#254} \quad \textbf{MÃ CÂU: \texttt{[L08\_C9\_B2\_04]}} \quad | \quad \textbf{Nhóm: \texttt{HINH\_196}} \quad | \quad \textbf{MapID: \texttt{8H9H2}} \hfill \textcolor{yellow!90!white}{[Hình gốc .pyc]}}
]
\textbf{Dạng toán:} Phát hiện 2 tam giác đồng dạng, tính chu vi của 1 tam giác khi biết chu vi tam giác kia \hfill \textit{(Trắc nghiệm -- Thông hiểu)}
\vspace{4pt}\hrule\vspace{6pt}
\begin{minipage}[t]{0.58\textwidth}
\vspace{0pt}
Trong hình bên, cho biết $AG \parallel HJ$; $JL \parallel GH$. Trong các khẳng định sau khẳng định nào đúng.

\vspace{4pt}\noindent
\textbf{A.} $ \Delta AGH \backsim \Delta JHL$ \quad \textbf{\textcolor{blue!75!black}{B. (Đ)}} $ \Delta AGH \backsim \Delta HJL$ \quad \textbf{C.} $ \Delta GAH \backsim \Delta HJL$ \quad \textbf{D.} $ \Delta AHG \backsim \Delta HJL$
\end{minipage}\hfill
\begin{minipage}[t]{0.39\textwidth}
\vspace{0pt}
\centering

\resizebox{\linewidth}{!}{%
\begin{tikzpicture}[scale=1.2, font=\footnotesize, line join=round, line cap=round, >=stealth]
            \path
            (1,3) coordinate (A)
            (0,0) coordinate (B)
            (4,0) coordinate (C)
            ($(A)!0.65!(C)$) coordinate (E)
            ($2/5*(A)+(C)$) coordinate (D2)
            ($(E)+(2,0)$) coordinate (D1)
            (intersection of C--D2 and E--D1) coordinate (D)
            ;
            \draw [line width =.5pt,blue] (A)--(B)--(C)--cycle (D)--(C) (D)--(E);
                   \filldraw[black] (A) circle (1pt) node[above] {$A$};
    \filldraw[black] (B) circle (1pt) node[below left] {$G$};
    \filldraw[black] (C) circle (1pt) node[below right] {$H$};
     \filldraw[black] (D) circle (1pt) node[above right]{$J$};
    \filldraw[black] (E) circle (1pt) node[below left] {$L$};
            \end{tikzpicture}
}
\end{minipage}

\vspace{3pt}\hrule\vspace{3pt}\footnotesize\textcolor{gray!80!black}{\textbf{Các mã dùng chung hình (HINH\_196):} \texttt{L08\_C9\_B2\_03}}
\end{tcolorbox}
\vspace{4pt}


\begin{tcolorbox}[
    breakable,
    colback=white,
    colframe=blue!60!black,
    boxrule=0.8pt,
    arc=3pt,
    title={\textbf{Câu \#255} \quad \textbf{MÃ CÂU: \texttt{[L08\_C9\_B2\_05]}} \quad | \quad \textbf{Nhóm: \texttt{HINH\_194}} \quad | \quad \textbf{MapID: \texttt{8H9N2}} \hfill \textcolor{yellow!90!white}{[Hình gốc .pyc]}}
]
\textbf{Dạng toán:} Cho hình vẽ, chỉ ra cặp tam giác đồng dạng theo c.c.c \hfill \textit{(Tự luận -- Nhận biết)}
\vspace{4pt}\hrule\vspace{6pt}
\begin{minipage}[t]{0.58\textwidth}
\vspace{0pt}
Cho $\Delta {AGH}$ và $\Delta {KOQ}$ có kích thước các cạnh như hình bên. Chứng minh $ \Delta AGH \backsim \Delta KOQ$
\end{minipage}\hfill
\begin{minipage}[t]{0.39\textwidth}
\vspace{0pt}
\centering

\resizebox{\linewidth}{!}{%
\begin{tikzpicture}[scale=1, font=\footnotesize, line join=round, line cap=round, >=stealth]
                \path
                (1,3) coordinate (D)
                (0,0) coordinate (E)
                (4,0) coordinate (F)
                ($(D)!0.5!(E)$)node[xshift=-3mm]{$24$}
                ($(D)!0.5!(F)$)node[xshift=4mm]{$26$}
                ($(F)!0.5!(E)$)node[below]{$30$}
                ;
                \draw [line width = .3pt,blue] (D)--(E)--(F)--cycle;
                \filldraw[black] (D) circle (1pt) node[above right]{$K$};
    \filldraw[black] (E) circle (1pt) node[below left] {$O$};
    \filldraw[black] (F) circle (1pt) node[below right] {$Q$};
            \end{tikzpicture}
            \begin{tikzpicture}[scale=1, font=\footnotesize, line join=round, line cap=round, >=stealth]
                \path
                (0.6,1.8) coordinate (A)
                (0,0) coordinate (B)
                (2.4,0) coordinate (C)
                ($(A)!0.5!(B)$)node[xshift=-3mm]{$12$}
                ($(A)!0.5!(C)$)node[xshift=4mm]{$13$}
                ($(B)!0.5!(C)$)node[below]{$15$}
                ;

                \draw [line width = .3pt,blue] (A)--(B)--(C)--cycle;
                   \filldraw[black] (A) circle (1pt) node[above left] {$A$};
    \filldraw[black] (B) circle (1pt) node[below left] {$G$};
    \filldraw[black] (C) circle (1pt) node[below right] {$H$};
            \end{tikzpicture}
}
\end{minipage}

\vspace{3pt}\hrule\vspace{3pt}\footnotesize\textcolor{gray!80!black}{\textbf{Các mã dùng chung hình (HINH\_194):} \texttt{L08\_C9\_B2\_01}}
\end{tcolorbox}
\vspace{4pt}


\begin{tcolorbox}[
    breakable,
    colback=white,
    colframe=blue!60!black,
    boxrule=0.8pt,
    arc=3pt,
    title={\textbf{Câu \#256} \quad \textbf{MÃ CÂU: \texttt{[L08\_C9\_B2\_08]}} \quad | \quad \textbf{Nhóm: \texttt{HINH\_197}} \quad | \quad \textbf{MapID: \texttt{8H9V2}} \hfill \textcolor{yellow!90!white}{[Hình gốc .pyc]}}
]
\textbf{Dạng toán:} Toán thực tế ứng dụng tam giác đồng dạng \hfill \textit{(Tự luận -- Vận dụng)}
\vspace{4pt}\hrule\vspace{6pt}
\begin{minipage}[t]{0.58\textwidth}
\vspace{0pt}
Một công viên có hai đường chạy bộ hình tam giác đồng dạng như hình bên. Kích thước của con đường bên trong lần lượt là $331$ m, $376$ m, $575$ m. Cạnh ngắn nhất của con đường bên ngoài là $1324$ m. Huệ chạy $5$ vòng trên con đường bên trong, Công chạy $4$ vòng trên con đường bên ngoài. So sánh quãng đường chạy được của hai bạn.
\end{minipage}\hfill
\begin{minipage}[t]{0.39\textwidth}
\vspace{0pt}
\centering

\resizebox{\linewidth}{!}{%
\begin{tikzpicture}[scale=0.6,line join=round, line cap=round, >=stealth]
    \path (0,0) coordinate (A) (-5,-4) coordinate (B)  (5,-3) coordinate (C);
    \draw (A)--(B)--(C)--cycle;
    \begin{scope}[scale=.5,shift={(0,-2)}]
        \path (0,0) coordinate (D) (-5,-4) coordinate (E)  (5,-3) coordinate (F);
        \draw (D)--(E)--(F)--cycle;
    \end{scope}
    \draw ($(A)!.3!(B)$)--(-5,1)--(-3.7,2);
    \foreach \x/\g in {A/90,B/180,C/0,D/90,E/180,F/0} \fill[black] (\x)+(\g:.3) node {$\x$};
    
    \foreach \i/\j in {-5/.6,7/.5,8/-10,0/-11,-8/-7}{
        \begin{scope}[scale=.5,shift={(\i,\j)}]%cây 2
            \draw[fill=black!70] plot[smooth cycle] coordinates{(-2.14,0.58)(-2.09,0.24)(-2.33,-0.13) (-1.19,-0.13) (-1.38,-0.13) (-1.69,0.21)(-1.75,0.53)};
            \draw[fill=blue!70] plot[smooth cycle] coordinates{(-1.75,0.53)(-1.46,0.79) (-1.01,0.58) (-1.08,1.06) (-0.64,1.08)(-0.93,1.59)(-0.53,1.85)(-0.93,2.2)(-1.38,2.7)(-2.04,3.18)(-2.49,2.91)(-3.07,2.91)(-3.2,2.22)(-3.73,1.96)(-3.1,1.3)(-3.31,1.01)(-2.83,1.01)(-2.99,0.69)(-2.14,0.58)};
    \end{scope}}
    \foreach \i/\j in {-8/-2,-12/-5,3/-4}{
        \begin{scope}[scale=.5,shift={(\i,\j)}]%cây 1
            \draw[fill=black!70] plot[smooth ] coordinates{(5.3,0)(5.25,1.14)(4.95,1.19) (5,0) };
            \draw[fill=blue!30] plot[smooth cycle] coordinates{(5.11,1.14)(5.87,1.01) (6.09,1.77) (5.93,2.62) (6.03,3.31)(5.5,3.65)(5.37,4.45)(5,5)(4.66,4.29)(4.29,3.55)(4.5,2.88)(3.94,2.62)(3.81,1.8)(4.23,1.24)(4.82,1.19)};
    \end{scope}}
    \foreach \i/\j in {3/1,17/-7,-20/-16}{
        \begin{scope}[scale=.3,shift={(\i,\j)}]%đèn
            \fill[blue] (0,0)--(.4,.1)--(.5,.4)--(.35,.55)--(.35,1.5)--(.2,1.6)--(.1,4.7)--(.3,4.8)--(.6,6.1)--(-.6,6.1)--(-.3,4.8)--(-.1,4.7)--(-.2,1.6)--(-.35,1.5)--(-.35,.55)--(-.5,.4)--(-.4,.1)--(0,0)
            (0,7.5).. controls +(30:0) and +(160:.9) ..(.8,6.1).. controls +(-90:.1) and +(-90:.1) ..(-.8,6.1).. controls +(20:.9) and +(20:0) ..(0,7.5);
            \fill[white] (.54,5.9)--(.35,5.87)--(.2,4.85)--(.3,4.85)--cycle (.2,5.85)--(-.2,5.86)--(-.1,4.83)--(.15,4.83)--cycle
            (-.54,5.9)--(-.35,5.87)--(-.2,4.85)--(-.3,4.85)--cycle;
    \end{scope}}

\end{tikzpicture}
}
\end{minipage}

\end{tcolorbox}
\vspace{4pt}


\begin{tcolorbox}[
    breakable,
    colback=white,
    colframe=blue!60!black,
    boxrule=0.8pt,
    arc=3pt,
    title={\textbf{Câu \#257} \quad \textbf{MÃ CÂU: \texttt{[L08\_C9\_B2\_10]}} \quad | \quad \textbf{Nhóm: \texttt{HINH\_198}} \quad | \quad \textbf{MapID: \texttt{8H9N2}} \hfill \textcolor{yellow!90!white}{[Hình gốc .pyc]}}
]
\textbf{Dạng toán:} Cho hvẽ chỉ ra cặp tam giác đồng dạng theo c.g.c \hfill \textit{(Tự luận -- Nhận biết)}
\vspace{4pt}\hrule\vspace{6pt}
\begin{minipage}[t]{0.58\textwidth}
\vspace{0pt}
Xét xem các cặp tam giác nào trong Hình a, b có đồng dạng hay không?
\end{minipage}\hfill
\begin{minipage}[t]{0.39\textwidth}
\vspace{0pt}
\centering

\resizebox{\linewidth}{!}{%
\begin{tikzpicture}[scale=1,line join=round,font=\footnotesize ,line cap=round,thick]
\coordinate (A) at (0,0);
\coordinate (B) at (2,-2);
\coordinate (C) at (6,-0.5);
\coordinate (D) at (4,1.5);
\coordinate (E) at (6,0.5);
\coordinate (F) at (2,0.5);
\draw[thick](A)--(B)--(C)--cycle;
\draw[thick](D)--(E)--(F)--cycle;

\draw(3.9,1.2)node[]{$134^\circ$};
\draw(2.1,-1.7)node[]{$134^\circ$};

   \filldraw[black] (A) circle (1pt) node[above left] {$D$};
    \filldraw[black] (B) circle (1pt) node[below] {$F$};
    \filldraw[black] (C) circle (1pt) node[above right] {$G$};
   \filldraw[black] (D) circle (1pt) node[above right]{$I$};
    \filldraw[black] (E) circle (1pt) node[below right] {$J$};
    \filldraw[black] (F) circle (1pt) node[below left] {$K$};


\path ($(A)!0.5!(B)$) node [below]{$12$};
\path ($(B)!0.5!(C)$) node [below]{$16$};
\path ($(D)!0.5!(F)$) node [above]{$4$};
\path ($(D)!0.5!(E)$) node [above]{$3$};
\draw(3,-2.5)node[]{$a)$};
\end{tikzpicture}\hspace{0.5cm}
\begin{tikzpicture}[scale=1,line join=round,font=\footnotesize ,line cap=round,thick]
\coordinate (M) at (0,4);
\coordinate (N) at (-1,0);
\coordinate (P) at (2,0);
\draw(M)--(N)--(P)--cycle;
 \filldraw[black] (M) circle (1pt) node[above] {$N$};
    \filldraw[black] (N) circle (1pt) node[below] {$O$};
     \filldraw[black] (P) circle (1pt) node[below] {$Q$};



    \draw(0.1,3.3)node[]{$48^\circ$};

\path ($(M)!0.5!(P)$) node [above right]{$5$};
\coordinate (T) at (4,2.5);
\coordinate (Q) at (3,0);
\coordinate (K) at (4.5,0);
\draw(T)--(Q)--(K)--cycle;

 \filldraw[black] (T) circle (1pt) node[above] {$T$};
    \filldraw[black] (Q) circle (1pt) node[below] {$V$};
     \filldraw[black] (K) circle (1pt) node[below] {$W$};


\path ($(T)!0.5!(K)$) node [right]{$10$};
\draw(3.9,1.6)node[]{$48^\circ$};



\draw(2.5,-1)node[]{$b)$};
\end{tikzpicture}
}
\end{minipage}

\end{tcolorbox}
\vspace{4pt}


\begin{tcolorbox}[
    breakable,
    colback=white,
    colframe=blue!60!black,
    boxrule=0.8pt,
    arc=3pt,
    title={\textbf{Câu \#258} \quad \textbf{MÃ CÂU: \texttt{[L08\_C9\_B2\_11]}} \quad | \quad \textbf{Nhóm: \texttt{HINH\_199}} \quad | \quad \textbf{MapID: \texttt{8H9N2}} \hfill \textcolor{yellow!90!white}{[Hình gốc .pyc]}}
]
\textbf{Dạng toán:} Cho h.vẽ, cm tam giác đồng dạng theo c.c.c \hfill \textit{(Tự luận -- Nhận biết)}
\vspace{4pt}\hrule\vspace{6pt}
\begin{minipage}[t]{0.58\textwidth}
\vspace{0pt}
a) $\Delta {DEF}$ và $\Delta {HIM}$ ở hình bên có đồng dạng không? Vì sao?

b) Biết $\Delta {DEF}$ có chu vi bằng ${18}cm$. Tính chu vi $\Delta {HIM}$.
\end{minipage}\hfill
\begin{minipage}[t]{0.39\textwidth}
\vspace{0pt}
\centering

\resizebox{\linewidth}{!}{%
\begin{tikzpicture}[scale=1, font=\footnotesize, line join=round, line cap=round, >=stealth]
                \path
                (2,3) coordinate (A)
                (0,0) coordinate (E)
                (4,0) coordinate (F)
                ($(A)!0.5!(E)$)node[xshift=-3mm,yshift=1mm]{$3a$}
                ($(F)!0.5!(E)$)node[below]{$3b$}
                ($(A)!0.5!(F)$)node[xshift=3mm,yshift=1mm]{$3c$}
                ;
                \draw [line width = .3pt,blue] (A)--(E)--(F)--cycle;
                  \filldraw[black] (A) circle (1pt) node[above] {$D$};
    \filldraw[black] (E) circle (1pt) node[left] {$E$};
    \filldraw[black] (F) circle (1pt) node[right] {$F$};v
            \end{tikzpicture}
            \begin{tikzpicture}[scale=1, font=\footnotesize, line join=round, line cap=round, >=stealth]
                \path
                (1,1.5) coordinate (M)
                (0,0) coordinate (N)
                (2,0) coordinate (G)
                ($(M)!0.5!(N)$)node[xshift=-3mm,yshift=1mm]{$a$}
                ($(N)!0.5!(G)$)node[below]{$b$}
                ($(M)!0.5!(G)$)node[xshift=3mm,yshift=1mm]{$c$}
                ;
                \draw [line width = .3pt,blue] (M)--(N)--(G)--cycle;
                    \filldraw[black] (M) circle (1pt) node[above] {$H$};
            \filldraw[black] (N) circle (1pt) node[left] {$I$};
            \filldraw[black] (G) circle (1pt) node[right] {$M$};
            \end{tikzpicture}
}
\end{minipage}

\end{tcolorbox}
\vspace{4pt}


\begin{tcolorbox}[
    breakable,
    colback=white,
    colframe=blue!60!black,
    boxrule=0.8pt,
    arc=3pt,
    title={\textbf{Câu \#259} \quad \textbf{MÃ CÂU: \texttt{[L08\_C9\_B3\_08]}} \quad | \quad \textbf{Nhóm: \texttt{HINH\_200}} \quad | \quad \textbf{MapID: \texttt{8H9N3}} \hfill \textcolor{yellow!90!white}{[Hình gốc .pyc]}}
]
\textbf{Dạng toán:} Cho h.vẽ,tìm cạnh huyền của tam giác vuông \hfill \textit{(Trắc nghiệm -- Nhận biết)}
\vspace{4pt}\hrule\vspace{6pt}
\begin{minipage}[t]{0.58\textwidth}
\vspace{0pt}
Tìm ${x}$.

\vspace{4pt}\noindent
\textbf{A.} $65$ cm \quad \textbf{B.} $59$ cm \quad \textbf{\textcolor{blue!75!black}{C. (Đ)}} $61$ cm \quad \textbf{D.} $64$ cm
\end{minipage}\hfill
\begin{minipage}[t]{0.39\textwidth}
\vspace{0pt}
\centering

\resizebox{\linewidth}{!}{%
\begin{tikzpicture}[scale=0.8, font=\footnotesize,>=stealth]
        % Tọa độ cho tam giác vuông
        \coordinate (A) at (0, 0);
        \coordinate (B) at (4, 0);
        \coordinate (C) at (0, 3);

        % Vẽ các cạnh của tam giác
        \draw[thick] (A) -- (B) -- (C) -- cycle;

        % Ghi nhãn các điểm
        \node[below left] at (A) {$A$};
        \node[below right] at (B) {$C$};
        \node[above left] at (C) {$D$};
\node[left] at (-0.1,1.3) {$60$cm};
    \node[left] at (2.5,-0.3) {$11$cm};
    \node[left] at (2.3,2) {$x$};
        % Vẽ dấu vuông tại điểm A
        \draw (0,0) rectangle +(0.3,0.3);

    \end{tikzpicture}
}
\end{minipage}

\end{tcolorbox}
\vspace{4pt}


\begin{tcolorbox}[
    breakable,
    colback=white,
    colframe=blue!60!black,
    boxrule=0.8pt,
    arc=3pt,
    title={\textbf{Câu \#260} \quad \textbf{MÃ CÂU: \texttt{[L08\_C9\_B3\_09]}} \quad | \quad \textbf{Nhóm: \texttt{HINH\_201}} \quad | \quad \textbf{MapID: \texttt{8H9N3}} \hfill \textcolor{yellow!90!white}{[Hình gốc .pyc]}}
]
\textbf{Dạng toán:} Cho h.vẽ,tìm cạnh góc vuông của tam giác vuông \hfill \textit{(Trắc nghiệm -- Nhận biết)}
\vspace{4pt}\hrule\vspace{6pt}
\begin{minipage}[t]{0.58\textwidth}
\vspace{0pt}
Tìm ${x}$.

\vspace{4pt}\noindent
\textbf{A.} $26$ cm \quad \textbf{B.} $25$ cm \quad \textbf{\textcolor{blue!75!black}{C. (Đ)}} $24$ cm \quad \textbf{D.} $23$ cm
\end{minipage}\hfill
\begin{minipage}[t]{0.39\textwidth}
\vspace{0pt}
\centering

\resizebox{\linewidth}{!}{%
\begin{tikzpicture}[scale=0.8, font=\footnotesize,>=stealth]
        % Tọa độ cho tam giác vuông
        \coordinate (A) at (0, 0);
        \coordinate (B) at (4, 0);
        \coordinate (C) at (0, 3);

        % Vẽ các cạnh của tam giác
        \draw[thick] (A) -- (B) -- (C) -- cycle;

        % Ghi nhãn các điểm
        \node[below left] at (A) {$A$};
        \node[below right] at (B) {$B$};
        \node[above left] at (C) {$D$};
\node[left] at (-0.1,1.3) {$x$};
    \node[left] at (2.5,-0.3) {$7$cm};
    \node[left] at (2.9,1.8) {$25$cm};
        % Vẽ dấu vuông tại điểm A
        \draw (0,0) rectangle +(0.3,0.3);

    \end{tikzpicture}
}
\end{minipage}

\end{tcolorbox}
\vspace{4pt}


\begin{tcolorbox}[
    breakable,
    colback=white,
    colframe=blue!60!black,
    boxrule=0.8pt,
    arc=3pt,
    title={\textbf{Câu \#261} \quad \textbf{MÃ CÂU: \texttt{[L08\_C9\_B3\_11]}} \quad | \quad \textbf{Nhóm: \texttt{HINH\_202}} \quad | \quad \textbf{MapID: \texttt{8H9H3}} \hfill \textcolor{yellow!90!white}{[Hình gốc .pyc]}}
]
\textbf{Dạng toán:} Toán thực tế tính chiều cao của bức tường \hfill \textit{(Trắc nghiệm -- Thông hiểu)}
\vspace{4pt}\hrule\vspace{6pt}
\begin{minipage}[t]{0.58\textwidth}
\vspace{0pt}
Một cái thang dài 4,8 m tựa vào mét tường, khoảng cách từ chân thang đến chân tường là $1,3$ m. Tính chiều cao của tường (làm tròn đến hàng phần mười).

\vspace{4pt}\noindent
\textbf{A.} ${5,1 \left( m \right) } $ \quad \textbf{B.} ${3,9 \left( m \right) }$ \quad \textbf{C.} ${4,1 \left(m \right) }$ \quad \textbf{\textcolor{blue!75!black}{D. (Đ)}} ${4,6  \left( m \right) } $
\end{minipage}\hfill
\begin{minipage}[t]{0.39\textwidth}
\vspace{0pt}
\centering

\resizebox{\linewidth}{!}{%
\begin{tikzpicture}[scale=0.8, line join=round, line cap=round,>=stealth]   
    \path
    (-1,0) coordinate (P)+(-80:1) 
    (P)+(60:5) coordinate (Q)
    (P)+(-30:0.2) coordinate (P1)+(-90:.3) node{$B$}
    (Q)+(-30:0.2) coordinate (Q1)+(0:.3) node{$A$}
    (P)+(150:0.9) coordinate (PP)
    (Q)+(150:0.9) coordinate (QQ)
    (P1)+(150:0.9) coordinate (PP1)
    (Q1)+(150:0.9) coordinate (QQ1)
    (4,0) coordinate (SD)
    (SD)+(150:10) coordinate (ST)
    (Q1)+(-90:1) coordinate (QT)
    (Q)+(0:1) coordinate (Q2)
    (intersection of SD--ST and Q1--QT) coordinate (C)+(-90:.3) node{$C$}
    ;
    
    \foreach \b in {0.1,0.2,...,0.9} {
        \path 
        ($(P)!\b!(Q)$) coordinate (m)
        (m)+(50:0.2)coordinate (m1)
        (m)+(150:0.7)coordinate (m2)
        (m1)+(150:0.7)coordinate (m3);
        \draw[fill=cyan!50](m)--(m1)--(m3)--(m2)--(m);}
    \draw[fill=cyan!50] (P)--(Q)--(Q1)--(P1)--(P);
    \draw[fill=cyan!50] (PP)--(QQ)--(QQ1)--(PP1)--(PP);
    \draw  (P1)--(C)--(Q1);
    \draw (Q1)--++(-30:2) (Q1)--++(150:3);
    \draw (Q2)--++(-30:2) (Q2)--++(150:3);
   
\end{tikzpicture}
}
\end{minipage}

\vspace{3pt}\hrule\vspace{3pt}\footnotesize\textcolor{gray!80!black}{\textbf{Các mã dùng chung hình (HINH\_202):} \texttt{L08\_C9\_B3\_10}}
\end{tcolorbox}
\vspace{4pt}


\begin{tcolorbox}[
    breakable,
    colback=white,
    colframe=blue!60!black,
    boxrule=0.8pt,
    arc=3pt,
    title={\textbf{Câu \#262} \quad \textbf{MÃ CÂU: \texttt{[L08\_C9\_B3\_12]}} \quad | \quad \textbf{Nhóm: \texttt{HINH\_203}} \quad | \quad \textbf{MapID: \texttt{8H9V3}} \hfill \textcolor{yellow!90!white}{[Hình gốc .pyc]}}
]
\textbf{Dạng toán:} Toán thực tế tính chiều cao của cây \hfill \textit{(Trắc nghiệm -- Vận dụng)}
\vspace{4pt}\hrule\vspace{6pt}
\begin{minipage}[t]{0.58\textwidth}
\vspace{0pt}
Sau một trận bão lớn, một cái cây bị gãy ngang (như hình vẽ). Ngọn cây chạm mặt đất cách gốc $8$ m. Đoạn thân cây còn lại người ta đo được làm $15$ m. Hỏi lúc đầu cây cao bao nhiêu

\vspace{4pt}\noindent
\textbf{A.} ${30 \left( m \right) }$ \quad \textbf{B.} ${33 \left( m \right) } $ \quad \textbf{\textcolor{blue!75!black}{C. (Đ)}} ${32  \left( m \right) } $ \quad \textbf{D.} ${31 \left( m \right) }$
\end{minipage}\hfill
\begin{minipage}[t]{0.39\textwidth}
\vspace{0pt}
\centering

\resizebox{\linewidth}{!}{%
\begin{tikzpicture}[font=\footnotesize, line join=round, line cap=round, >=stealth]
    \tikzset{Cay-tre/.pic={
            \tikzset{la/.pic={
                    \fill[green!80!white](0,0)..controls++(-20:0.5) and ++(-175:1)..(2,0)..controls++(170:1) and ++(20:0.5)..(0,0);
            } }
            \def\cao{4}
            \def\a{0.1*\cao}
            \def\k{0.5}% Hệ số góc
            \pgfmathsetmacro{\dai}{\cao*\k}
            \draw[fill=green!75!blue,draw=teal](0,0)rectangle(\a,\cao)pic[rotate=20]{la};
            \draw[fill=green!70!blue,draw=teal](-\dai,0)--(0,\cao)coordinate(cao)--([turn]90:\a)--([turn]90:{2/.9*\dai})--cycle;
            \draw[teal](-\dai,0)coordinate(dai)--(0,0)coordinate(oo);
            \begin{scope}
                \clip(oo)--(cao)--(dai)--cycle;
            \end{scope}
            \draw[line width=3pt,green!65!blue](0,\cao)pic[rotate=-130]{la}--(\a,\cao)(cao)--++(120:\a)pic[rotate=150]{la}pic[rotate=90]{la}(dai)--++(150:\a)(-0.5*\dai,0.5*\cao)pic[rotate=-40]{la}--++(150:\a)pic[rotate=90]{la}pic[rotate=20]{la}pic[rotate=100]{la};
            \path(dai)pic[rotate=-30]{la}(dai)pic[rotate=-60]{la}(dai)++(120:\a)pic[rotate=60]{la}pic[rotate=90]{la};
            \path(dai)--(oo)node[pos=0.5,below, teal]{$8$ m};
            \path(cao)--($(oo)+(.8,0)$)node[pos=0.5,above right, teal]{$15$ m};
    } }
    \path(0,0)pic[scale=1]{Cay-tre};
\end{tikzpicture}
}
\end{minipage}

\end{tcolorbox}
\vspace{4pt}


\begin{tcolorbox}[
    breakable,
    colback=white,
    colframe=blue!60!black,
    boxrule=0.8pt,
    arc=3pt,
    title={\textbf{Câu \#263} \quad \textbf{MÃ CÂU: \texttt{[L08\_C9\_B3\_15]}} \quad | \quad \textbf{Nhóm: \texttt{HINH\_204}} \quad | \quad \textbf{MapID: \texttt{8H9H3}} \hfill \textcolor{yellow!90!white}{[Hình gốc .pyc]}}
]
\textbf{Dạng toán:} Toán thực tế tính k.cách giữa ngọn hải đăng và con tàu \hfill \textit{(Trắc nghiệm -- Thông hiểu)}
\vspace{4pt}\hrule\vspace{6pt}
\begin{minipage}[t]{0.58\textwidth}
\vspace{0pt}
Một con tàu cách chân ngọn hải đăng $265$ m, biết hải đăng cao $47$ m. Tính khoảng cách từ con tàu đến đỉnh ngọn hải đăng. (Làm tròn kết quả đến hàng đơn vị)

\vspace{4pt}\noindent
\textbf{A.} ${259 \left( m \right) }$ \quad \textbf{B.} ${249 \left( m \right) }$ \quad \textbf{C.} ${279 \left( m \right) } $ \quad \textbf{\textcolor{blue!75!black}{D. (Đ)}} ${269  \left( m \right) } $
\end{minipage}\hfill
\begin{minipage}[t]{0.39\textwidth}
\vspace{0pt}
\centering

\resizebox{\linewidth}{!}{%
\begin{tikzpicture}
    % Sea (slightly wavy line)
    \fill[blue!30] (0,0) rectangle (10,-2);
    \draw[blue, thick] (0,-0.1) .. controls (1,-0.3) and (2,0.1) .. (3,-0.2)
    .. controls (4,-0.5) and (5,-0.1) .. (6,-0.3)
    .. controls (7,-0.1) and (8,-0.3) .. (9,-0.2)
    .. controls (10,-0.5) and (11,0.1) .. (12,-0.2);
    
    % Enhanced lighthouse (cartoonish style with more details)
    
    % Lighthouse base (stone)
    \fill[gray!50] (0.7,-0.5) rectangle (1.3,0.3);
    \draw[thick] (0.7,-0.5) -- (0.7,0.3) -- (1.3,0.3) -- (1.3,-0.5) -- cycle;
    
    % Lighthouse body (stripes for a more lively look)
    \fill[white] (0.7,0.3) rectangle (1.3,1);
    \fill[red!80] (0.7,1) rectangle (1.3,1.5);
    \fill[white] (0.7,1.5) rectangle (1.3,2);
    \fill[red!80] (0.7,2) rectangle (1.3,2.5);
    \fill[white] (0.7,2.5) rectangle (1.3,3);
    
    \draw[thick] (0.7,0.3) -- (0.7,3) -- (1.3,3) -- (1.3,0.3) -- cycle;
    
    % Windows on the lighthouse
    \fill[black!70] (0.9,1.2) rectangle (1.1,1.4); % first window
    \fill[black!70] (0.9,2.2) rectangle (1.1,2.4); % second window
    
    % Lighthouse top (lantern room)
    \fill[gray!70] (0.6,3) rectangle (1.4,3.3);
    \draw[thick] (0.6,3) -- (1.4,3);
    \draw[thick] (0.6,3.3) -- (1.4,3.3);
    
    % Lighthouse roof (dome-shaped top)
    \fill[red!80] (0.65,3.3) -- (1.35,3.3) -- (1,3.8) -- cycle;
    
    % Lighthouse light
    \fill[yellow] (1,3.15) circle (0.1);
    
    % Lighthouse light beams (enhanced beams with glow effect)
    \draw[yellow!80, thick] (1,3.15) -- (0,4);
    \draw[yellow!80, thick] (1,3.15) -- (2,4);
    
    % Lighthouse door
    \fill[brown!80] (0.9,-0.5) rectangle (1.1,-0.2);
    \draw[thick] (0.9,-0.5) -- (0.9,-0.2) -- (1.1,-0.2) -- (1.1,-0.5) -- cycle;
    
    \foreach \x in {3.5, 5.5, 7.5} {
        \draw[blue!60,thick] (\x,-0.5) .. controls (\x+0.3,-0.6) and (\x+0.7,-0.4) .. (\x+1,-0.5);
    }
    
    % Second boat (smaller and farther away)
    \fill[brown!70] (7,-0.5) .. controls (7.1,-0.9) and (7.9,-0.9) .. (8,-0.5) -- cycle; % hull
    \draw[thick] (7,-0.5) .. controls (7.1,-0.9) and (7.9,-0.9) .. (8,-0.5) -- cycle;
    
    \draw[thick] (7.5,-0.5) -- (7.5,0.2); % mast
    \fill[white] (7.5,-0.1) .. controls (8,0) and (8,0.1) .. (7.5,0.2) -- cycle; % sail
    \draw[thick] (7.5,-0.1) .. controls (8,0) and (8,0.1) .. (7.5,0.2) -- cycle;
    
    % Flag for the second boat
    \fill[red] (7.5,0.2) -- (7.7,0.3) -- (7.5,0.4) -- cycle;
    
    % Marking points A (top of lighthouse), B (bottom of lighthouse), C (first boat), D (second boat)
    \fill[black] (1,3.5) circle (0.05) node[above] {A};  % Point A (top of lighthouse)
    \fill[black] (1,-0.5) circle (0.05) node[below] {B}; % Point B (bottom of lighthouse)
 % Point C (position of boat 1)
    \fill[black] (7.5,-0.5) circle (0.05) node[below] {D}; % Point D (position of boat 2)
    
    % Draw lines connecting A-B, A-C, A-D
    \draw[thick] (1,3.5) -- (1,-0.5);  % Line AB% Line AC
    \draw[thick] (1,3.5) -- (7.5,-0.5); % Line AD
    
    % Draw line B-C-D
    \draw[thick] (1,-0.5) -- (5,-0.5) -- (7.5,-0.5);  % Line BCD
    
    % Add a line x parallel to BC through A
% Dashed line x, parallel to BC
    
    % Add a smiling sun in the sky
    \fill[yellow] (9,3.5) circle (0.5); % sun body
    \draw (9,3.7) arc[start angle=45, end angle=135, radius=0.2]; % smiling face
    \fill[black] (8.8,3.6) circle (0.05); % left eye
    \fill[black] (9.2,3.6) circle (0.05); % right eye
    \foreach \angle in {0,45,...,315} {
        \draw[thick, yellow] (9,3.5) -- ++(\angle:0.8);
    }
    
\end{tikzpicture}
}
\end{minipage}

\vspace{3pt}\hrule\vspace{3pt}\footnotesize\textcolor{gray!80!black}{\textbf{Các mã dùng chung hình (HINH\_204):} \texttt{L08\_C9\_B3\_14}}
\end{tcolorbox}
\vspace{4pt}


\begin{tcolorbox}[
    breakable,
    colback=white,
    colframe=blue!60!black,
    boxrule=0.8pt,
    arc=3pt,
    title={\textbf{Câu \#264} \quad \textbf{MÃ CÂU: \texttt{[L08\_C9\_B3\_01]}} \quad | \quad \textbf{Nhóm: \texttt{HINH\_205}} \quad | \quad \textbf{MapID: \texttt{8H9V3}} \hfill \textcolor{yellow!90!white}{[Hình gốc .pyc]}}
]
\textbf{Dạng toán:} Toán thực tế tính chiều cao của cây \hfill \textit{(Tự luận -- Vận dụng)}
\vspace{4pt}\hrule\vspace{6pt}
\begin{minipage}[t]{0.58\textwidth}
\vspace{0pt}
Sau một trận bão lớn, một cái cây bị gãy ngang (như hình vẽ). Ngọn cây chạm mặt đất cách gốc $9$ m. Đoạn thân cây còn lại người ta đo được làm $12$ m. Hỏi lúc đầu cây cao bao nhiêu
\end{minipage}\hfill
\begin{minipage}[t]{0.39\textwidth}
\vspace{0pt}
\centering

\resizebox{\linewidth}{!}{%
\begin{tikzpicture}[font=\footnotesize, line join=round, line cap=round, >=stealth]
    \tikzset{Cay-tre/.pic={
            \tikzset{la/.pic={
                    \fill[green!80!white](0,0)..controls++(-20:0.5) and ++(-175:1)..(2,0)..controls++(170:1) and ++(20:0.5)..(0,0);
            } }
            \def\cao{4}
            \def\a{0.1*\cao}
            \def\k{0.5}% Hệ số góc
            \pgfmathsetmacro{\dai}{\cao*\k}
            \draw[fill=green!75!blue,draw=teal](0,0)rectangle(\a,\cao)pic[rotate=20]{la};
            \draw[fill=green!70!blue,draw=teal](-\dai,0)--(0,\cao)coordinate(cao)--([turn]90:\a)--([turn]90:{2/.9*\dai})--cycle;
            \draw[teal](-\dai,0)coordinate(dai)--(0,0)coordinate(oo);
            \begin{scope}
                \clip(oo)--(cao)--(dai)--cycle;
            \end{scope}
            \draw[line width=3pt,green!65!blue](0,\cao)pic[rotate=-130]{la}--(\a,\cao)(cao)--++(120:\a)pic[rotate=150]{la}pic[rotate=90]{la}(dai)--++(150:\a)(-0.5*\dai,0.5*\cao)pic[rotate=-40]{la}--++(150:\a)pic[rotate=90]{la}pic[rotate=20]{la}pic[rotate=100]{la};
            \path(dai)pic[rotate=-30]{la}(dai)pic[rotate=-60]{la}(dai)++(120:\a)pic[rotate=60]{la}pic[rotate=90]{la};
            \path(dai)--(oo)node[pos=0.5,below, teal]{$9$ m};
            \path(cao)--($(oo)+(.8,0)$)node[pos=0.5,above right, teal]{$12$ m};
    } }
    \path(0,0)pic[scale=1]{Cay-tre};
\end{tikzpicture}
}
\end{minipage}

\end{tcolorbox}
\vspace{4pt}


\begin{tcolorbox}[
    breakable,
    colback=white,
    colframe=blue!60!black,
    boxrule=0.8pt,
    arc=3pt,
    title={\textbf{Câu \#265} \quad \textbf{MÃ CÂU: \texttt{[L08\_C9\_B3\_02]}} \quad | \quad \textbf{Nhóm: \texttt{HINH\_206}} \quad | \quad \textbf{MapID: \texttt{8H9V3}} \hfill \textcolor{yellow!90!white}{[Hình gốc .pyc]}}
]
\textbf{Dạng toán:} Toán thực tế xử lý tình huống. \hfill \textit{(Tự luận -- Vận dụng)}
\vspace{4pt}\hrule\vspace{6pt}
\begin{minipage}[t]{0.58\textwidth}
\vspace{0pt}
Trong lúc anh Thịnh dựng tủ cho đứng thẳng, tủ có bị vướng vào trần nhà không? Tính toán và giải thích lý do?
\end{minipage}\hfill
\begin{minipage}[t]{0.39\textwidth}
\vspace{0pt}
\centering
\definecolor{bananamania}{rgb}{0.98, 0.91, 0.71}
\definecolor{amber}{rgb}{1.0, 0.49, 0.0}
\definecolor{anti-flashwhite}{rgb}{0.95, 0.95, 0.96}
\definecolor{brown}{rgb}{0.59, 0.29, 0.0}
\definecolor{azure}{rgb}{0.0, 0.5, 1.0}
\resizebox{\linewidth}{!}{%
\begin{tikzpicture}[line join=round, line cap=round,scale=1,transform shape,scale=0.8]
    \clip (-5,-2) rectangle (15,3);
    \tikzset{hoi/.pic={
            \def\H{ 
                (3.4,1.8)
                ..controls +(50:.2) and + (0:.2) ..(3.4,2.15)
                ..controls +(-180:.2) and + (-150:.2) ..(3.4,2)
                ;
            }
            \draw[red]\H;
            \fill[red] (3.4,1.7) circle (.5pt);
    }}
    
    \tikzset{h1/.pic={
            \def\tay1{ %tay trái
                (3.45,.55)
                ..controls +(160:.05) and +(60:.05) ..(3.2,.3)
                ..controls +(160:.05) and +(60:.05) ..(3.25,.7)
                ..controls +(-120:.05) and +(20:.05) ..(3.4,.95)
                ..controls +(40:.3) and +(120:.2) ..(3.15,.7)
                ..controls +(-150:.15) and +(120:.15) ..(3,.05)
                ..controls +(-60:.1) and +(-140:.25) ..(3.4,.15)
                ;
            }
            \draw[black]\tay1;
            \fill[bananamania] \tay1;
            
            \def\A{ %áo
                (3.8,.7)--(3.87,.67)
                ..controls +(-100:.3) and +(160:.3) ..(4,.05)--(3.98,-.2)
                ..controls +(-160:.2) and +(-20:.2) ..(3.5,-.25)
                ..controls +(140:.2) and +(-120:.2) ..(3.47,.57)--cycle
                
                ;
            }
            \draw[black]\A;
            \fill[anti-flashwhite] \A;
            
            
            \def\tay{ %tay
                (3.87,.67)
                ..controls +(-100:.3) and +(160:.3) ..(4,.05)--(4.03,.2)
                ..controls +(-30:.1) and +(160:.1) ..(4.2,.05)%khuỷu
                ..controls +(-150:.1) and +(60:.1) ..(4.1,-.15)
                ..controls +(-150:.1) and +(60:.1) ..(4,-.2)%tay túi
                ..controls +(-150:.1) and +(160:.1) ..(4.1,-.4)
                ..controls +(30:.05) and +(-80:.05) ..(4.2,-.3)
                ..controls +(30:.05) and +(-80:.05) ..(4.4,.07)
                ..controls +(60:.05) and +(0:.2) ..cycle
                
                ;
            }
            \draw[black]\tay;
            \fill[bananamania] \tay;
            \draw[black] (4.03,.2)
            ..controls +(150:.1) and + (60:.1) ..(3.85,.3);
            
            
            \def\G{%giày
                (4.15,-1.75)--(4.17,-1.9)
                ..controls +(-100:.2) and +(-150:.3) ..(3.9,-1.85)
                (3.7,-1.75)--(3.7,-1.9)
                ..controls +(-160:.35) and +(-150:.45) ..(3.4,-1.8)
                ;
            }
            \draw[black]\G;
            \fill[brown] \G;
            
            \def\Q{ %quần
                (4,-.2)%tay túi
                ..controls +(-150:.1) and +(160:.1) ..(4.1,-.4)
                ..controls +(30:.05) and +(-80:.05) ..(4.15,-1.2)%gối
                ..controls +(30:.05) and +(-80:.05) ..(4.2,-1.8)
                ..controls +(-160:.05) and +(-10:.05) ..(3.9,-1.85)
                ..controls +(-80:.05) and +(130:.05) ..(3.8,-1.3)
                ..controls +(80:.05) and +(-90:.05) ..(3.7,-.7)
                ..controls +(80:.05) and +(-90:.05) ..(3.75,-1.8)
                ..controls +(-160:.05) and +(-10:.05) ..(3.4,-1.8)
                ..controls +(130:.05) and +(-70:.05) ..(3.35,-1.2)
                ..controls +(130:.05) and +(-70:.05) ..(3.5,-.25)
                ..controls +(-20:.1) and +(-175:.05) ..cycle
                ;
            }
            \draw[black]\Q;
            \fill[azure] \Q;
            
            \def\Th{ %thân
                (3.7,.9)%tai dưới
                ..controls +(-100:.2) and +(160:.1) ..(3.8,.7)
                ..controls +(-100:.3) and +(-10:.4) ..(3.45,.3)
                ..controls +(90:.2) and +(-90:.2) ..(3.5,.6)
                ..controls +(90:.2) and +(-90:.2) ..(3.45,.9)
                ;
            }
            \draw[black]\Th;
            \fill[bananamania] \Th;
            
            
            \def\N{ %mặt
                (3.3,1.1)
                ..controls +(-90:.4) and +(-125:.32) ..(3.7,.9)%tai dưới
                ..controls +(-20:.1) and +(65:.1) ..(3.7,1)--(3.65,1.1)--cycle
                ;
            }
            \draw[black]\N;
            \fill[bananamania] \N;
            
            \def\M{ %nón
                (3.3,1.3)
                ..controls +(-160:.4) and +(-165:.3) ..(3.65,1.1)
                ..controls +(-20:.2) and +(-140:.2) ..(3.85,1.1)
                ..controls +(80:.45) and +(70:.4) ..cycle
                ;
            }
            \draw[black]\M;
            \fill[amber] \M;
            
            \def\T{ %tóc
                (3.7,.9)%tai dưới
                ..controls +(-20:.1) and +(65:.1) ..(3.7,1)--(3.65,1.1)--(3.85,1.1)
                ..controls +(-50:.1) and +(-45:.1) ..cycle
                ;
            }
            \draw[black]\T;
            \fill[black] \T;
            
    }}
    
    \path
    (0,0)pic[scale=1]{h1};
    \path
    (0,0)pic[scale=1]{hoi}
    (-1,-1)pic[rotate=4,scale=1.2]{hoi}
    (0,-.1)pic[rotate=-4,scale=1.1]{hoi};
    
    %mắt
    \draw[black] (3.5,1)
    ..controls +(-30:.05) and + (-150:.05) ..(3.6,1);
    \draw[black] (3.3,1.02)
    ..controls +(-30:.02) and + (-150:.02) ..(3.37,1.02);
    %miệng
    \draw[black] (3.4,.85)
    ..controls +(-30:.05) and + (-150:.05) ..(3.5,.85);
    
    \draw (-2.5,-2) rectangle (5,2.5);
    \fill[brown!40] (-1.5,-2) rectangle (2.8,-1.1);
    \draw[pattern=north east lines,pattern color=brown!90]  (-1.5,-2) rectangle (2.8,-1.1);
    
    \draw[<->] (-2.7,-2)--(-2.7,2.5);
    \draw[<->] (-1.7,-2)--(-1.7,-1);
    \draw[<->] (-1.5,-.9)--(2.8,-.9);
    \node at (-2,-0.5) [left,rotate=90]{$4,4$ dm};
    \node at (-3,1) [left,rotate=90]{$16$ dm};
    \node at (1,-.5) [left]{$30$ dm};
    
    %%%%%%%%%%%%%%%%%%%%%%%%%%%%%%
    \tikzset{h2/.pic={
            \def\m{ %mặt
                (11.1,1.23)
                %..controls +(-150:.4) and +(-125:.32) ..(10.7,1)
                ..controls +(-130:.35) and +(-170:.3) ..(11.2,.8)
                --(11.15,.7)
                ..controls +(-170:.1) and +(-30:.1) ..(10.9,.65)
                ..controls +(150:.1) and +(-60:.4) ..(10.4,1)
                ..controls +(70:.1) and +(-60:.2) ..(10.4,1.25)--(10.25,1)
                ..controls +(-140:.15) and +(140:.1) ..(10.6,.56)
                ..controls +(160:.15) and +(-40:.1) ..(10.15,.9)--(10.05,.75)
                ..controls +(-60:.15) and +(150:.2) ..(10.45,.35)
                ..controls +(-20:.15) and +(-170:.15) ..(11,.35)
                ..controls +(-60:.45) and +(40:.45) ..(11,.53)
                ..controls +(40:.2) and +(-140:.1) ..(11.16,.7)%cổ áo 
                ..controls +(-60:.5) and +(-40:.5) ..(11.35,.72)--(11.4,.8)--(11.34,.87)
                ..controls +(20:.15) and +(40:.15) ..(11.34,.95)
                ..controls +(60:.15) and +(40:.15) ..cycle
                ;
            }
            \draw[black]\m;
            \fill[bananamania] \m;
            \draw (10.6,.56)
            ..controls +(20:.15) and +(-160:.1) ..(11,.55);
            
            \def\toc{ %toc
                (11.4,.8)--(11.34,.87)
                ..controls +(20:.15) and +(40:.15) ..(11.34,.95)
                ..controls +(60:.15) and +(40:.15) ..(11.1,1.23)
                ..controls +(40:.55) and +(20:.3) ..cycle
                ;
            }
            \draw[black]\toc;
            \fill[black] \toc;
            
            \def\N2{ %nón
                (11.3,1.4)
                ..controls +(170:.6) and +(160:.3) ..(11.35,1.1)
                ..controls +(-40:.2) and +(-120:.3) ..(11.7,1.1)
                ..controls +(60:.1) and +(30:.2) ..cycle
                ;
            }
            \draw[black]\N2;
            \fill[amber] \N2;
            
            \def\ao{ %áo
                (11,.35)
                ..controls +(-60:.45) and +(40:.45) ..(11,.53)
                ..controls +(40:.2) and +(-140:.1) ..(11.16,.7)%cổ áo 
                ..controls +(-60:.5) and +(-40:.5) ..(11.35,.72)
                ..controls +(-30:.4) and +(100:.25) ..(11.6,0)
                ..controls +(-40:.15) and +(120:.15) ..(11.8,-.3)
                ..controls +(-160:.15) and +(20:.2) ..(11.15,-.35)
                ..controls +(80:.2) and +(-100:.2) ..cycle
                ;
            }
            \draw[black]\ao;
            \fill[anti-flashwhite] \ao;
            
            \def\qu{ %quần
                (11.8,-.3)
                ..controls +(-160:.15) and +(20:.2) ..(11.17,-.35)
                ..controls +(-70:.1) and +(100:.1) ..(11.18,-.6)
                ..controls +(-140:.1) and +(90:.1) ..(10.9,-1)
                ..controls +(-80:.2) and +(50:.2) ..(10.55,-1.75)%giày
                ..controls +(-30:.1) and +(-170:.1) ..(11,-1.83)
                ..controls +(40:.1) and +(-120:.1) ..(11.3,-1)
                ..controls +(10:.1) and +(-150:.1) ..(11.5,-.9)
                ..controls +(-100:.15) and +(110:.1) ..(11.5,-1.3)
                ..controls +(-40:.15) and +(130:.2) ..(12,-1.75)
                ..controls +(10:.15) and +(-130:.1) ..(12.4,-1.7)
                ..controls +(120:.2) and +(-40:.2) ..(11.85,-1.15)
                ..controls +(100:.2) and +(-60:.2) ..cycle
                ;
            }
            \draw[black]\qu;
            \fill[azure] \qu;
            \draw[black] (11.5,-.9)..controls +(10:.1) and +(-80:.1) ..(11.7,-.55)
            ;
    }}
    
    \tikzset{gi1/.pic={
            \def\gi{ %giày
                (10.55,-1.75)
                ..controls +(-30:.1) and +(-30:.1) ..(10.94,-1.82)
                ..controls +(-100:.05) and +(30:.1) ..(10.92,-1.95)
                ..controls +(-150:.1) and +(-60:.1) ..(10.3,-1.85)
                ..controls +(100:.1) and +(-160:.1) ..cycle
                ;
            }
            \draw[black]\gi;
            \fill[brown] \gi;
    }}
    
    \path
    (0,0)pic[scale=1]{gi1}
    (2.3,-4.83)pic[scale=.95,rotate=27]{gi1}
    (0,0)pic[scale=1]{h2}
    ;
    %mắt
    \draw[black] (11.15,1.07)
    ..controls +(-30:.05) and + (-150:.05) ..(11.25,1.03);
    %miệng
    \draw[black] (11,.95)
    ..controls +(-30:.02) and + (-150:.02) ..(11.06,.92);
    \fill[brown!40,rotate=-38] (6.6,3.4) rectangle (7.45,7.6);
    \draw[rotate=-38,pattern=north east lines,pattern color=brown!90]  (6.6,3.4) rectangle (7.45,7.6);
    \draw (7,-2) rectangle (14.5,2.5);
\end{tikzpicture}
}
\end{minipage}

\end{tcolorbox}
\vspace{4pt}


\begin{tcolorbox}[
    breakable,
    colback=white,
    colframe=blue!60!black,
    boxrule=0.8pt,
    arc=3pt,
    title={\textbf{Câu \#266} \quad \textbf{MÃ CÂU: \texttt{[L08\_C9\_B3\_10]}} \quad | \quad \textbf{Nhóm: \texttt{HINH\_202}} \quad | \quad \textbf{MapID: \texttt{8H9H3}} \hfill \textcolor{yellow!90!white}{[Hình gốc .pyc]}}
]
\textbf{Dạng toán:} Toán thực tế tính chiều cao của bức tường \hfill \textit{(Tự luận -- Thông hiểu)}
\vspace{4pt}\hrule\vspace{6pt}
\begin{minipage}[t]{0.58\textwidth}
\vspace{0pt}
Một cái thang dài 4 m tựa vào mét tường, khoảng cách từ chân thang đến chân tường là $1,4$ m. Tính chiều cao của tường (làm tròn đến hàng phần mười).
\end{minipage}\hfill
\begin{minipage}[t]{0.39\textwidth}
\vspace{0pt}
\centering

\resizebox{\linewidth}{!}{%
\begin{tikzpicture}[scale=0.8, line join=round, line cap=round,>=stealth]   
    \path
    (-1,0) coordinate (P)+(-80:1) 
    (P)+(60:5) coordinate (Q)
    (P)+(-30:0.2) coordinate (P1)+(-90:.3) node{$B$}
    (Q)+(-30:0.2) coordinate (Q1)+(0:.3) node{$A$}
    (P)+(150:0.9) coordinate (PP)
    (Q)+(150:0.9) coordinate (QQ)
    (P1)+(150:0.9) coordinate (PP1)
    (Q1)+(150:0.9) coordinate (QQ1)
    (4,0) coordinate (SD)
    (SD)+(150:10) coordinate (ST)
    (Q1)+(-90:1) coordinate (QT)
    (Q)+(0:1) coordinate (Q2)
    (intersection of SD--ST and Q1--QT) coordinate (C)+(-90:.3) node{$C$}
    ;
    
    \foreach \b in {0.1,0.2,...,0.9} {
        \path 
        ($(P)!\b!(Q)$) coordinate (m)
        (m)+(50:0.2)coordinate (m1)
        (m)+(150:0.7)coordinate (m2)
        (m1)+(150:0.7)coordinate (m3);
        \draw[fill=cyan!50](m)--(m1)--(m3)--(m2)--(m);}
    \draw[fill=cyan!50] (P)--(Q)--(Q1)--(P1)--(P);
    \draw[fill=cyan!50] (PP)--(QQ)--(QQ1)--(PP1)--(PP);
    \draw  (P1)--(C)--(Q1);
    \draw (Q1)--++(-30:2) (Q1)--++(150:3);
    \draw (Q2)--++(-30:2) (Q2)--++(150:3);
   
\end{tikzpicture}
}
\end{minipage}

\vspace{3pt}\hrule\vspace{3pt}\footnotesize\textcolor{gray!80!black}{\textbf{Các mã dùng chung hình (HINH\_202):} \texttt{L08\_C9\_B3\_11}}
\end{tcolorbox}
\vspace{4pt}


\begin{tcolorbox}[
    breakable,
    colback=white,
    colframe=blue!60!black,
    boxrule=0.8pt,
    arc=3pt,
    title={\textbf{Câu \#267} \quad \textbf{MÃ CÂU: \texttt{[L08\_C9\_B3\_14]}} \quad | \quad \textbf{Nhóm: \texttt{HINH\_204}} \quad | \quad \textbf{MapID: \texttt{8H9H3}} \hfill \textcolor{yellow!90!white}{[Hình gốc .pyc]}}
]
\textbf{Dạng toán:} Toán thực tế tính k.cách giữa ngọn hải đăng và con tàu \hfill \textit{(Tự luận -- Thông hiểu)}
\vspace{4pt}\hrule\vspace{6pt}
\begin{minipage}[t]{0.58\textwidth}
\vspace{0pt}
Một con tàu cách chân ngọn hải đăng $279$ m, biết hải đăng cao $72$ m. Tính khoảng cách từ con tàu đến đỉnh ngọn hải đăng. (Làm tròn kết quả đến hàng đơn vị)
\end{minipage}\hfill
\begin{minipage}[t]{0.39\textwidth}
\vspace{0pt}
\centering

\resizebox{\linewidth}{!}{%
\begin{tikzpicture}
    % Sea (slightly wavy line)
    \fill[blue!30] (0,0) rectangle (10,-2);
    \draw[blue, thick] (0,-0.1) .. controls (1,-0.3) and (2,0.1) .. (3,-0.2)
    .. controls (4,-0.5) and (5,-0.1) .. (6,-0.3)
    .. controls (7,-0.1) and (8,-0.3) .. (9,-0.2)
    .. controls (10,-0.5) and (11,0.1) .. (12,-0.2);
    
    % Enhanced lighthouse (cartoonish style with more details)
    
    % Lighthouse base (stone)
    \fill[gray!50] (0.7,-0.5) rectangle (1.3,0.3);
    \draw[thick] (0.7,-0.5) -- (0.7,0.3) -- (1.3,0.3) -- (1.3,-0.5) -- cycle;
    
    % Lighthouse body (stripes for a more lively look)
    \fill[white] (0.7,0.3) rectangle (1.3,1);
    \fill[red!80] (0.7,1) rectangle (1.3,1.5);
    \fill[white] (0.7,1.5) rectangle (1.3,2);
    \fill[red!80] (0.7,2) rectangle (1.3,2.5);
    \fill[white] (0.7,2.5) rectangle (1.3,3);
    
    \draw[thick] (0.7,0.3) -- (0.7,3) -- (1.3,3) -- (1.3,0.3) -- cycle;
    
    % Windows on the lighthouse
    \fill[black!70] (0.9,1.2) rectangle (1.1,1.4); % first window
    \fill[black!70] (0.9,2.2) rectangle (1.1,2.4); % second window
    
    % Lighthouse top (lantern room)
    \fill[gray!70] (0.6,3) rectangle (1.4,3.3);
    \draw[thick] (0.6,3) -- (1.4,3);
    \draw[thick] (0.6,3.3) -- (1.4,3.3);
    
    % Lighthouse roof (dome-shaped top)
    \fill[red!80] (0.65,3.3) -- (1.35,3.3) -- (1,3.8) -- cycle;
    
    % Lighthouse light
    \fill[yellow] (1,3.15) circle (0.1);
    
    % Lighthouse light beams (enhanced beams with glow effect)
    \draw[yellow!80, thick] (1,3.15) -- (0,4);
    \draw[yellow!80, thick] (1,3.15) -- (2,4);
    
    % Lighthouse door
    \fill[brown!80] (0.9,-0.5) rectangle (1.1,-0.2);
    \draw[thick] (0.9,-0.5) -- (0.9,-0.2) -- (1.1,-0.2) -- (1.1,-0.5) -- cycle;
    
    \foreach \x in {3.5, 5.5, 7.5} {
        \draw[blue!60,thick] (\x,-0.5) .. controls (\x+0.3,-0.6) and (\x+0.7,-0.4) .. (\x+1,-0.5);
    }
    
    % Second boat (smaller and farther away)
    \fill[brown!70] (7,-0.5) .. controls (7.1,-0.9) and (7.9,-0.9) .. (8,-0.5) -- cycle; % hull
    \draw[thick] (7,-0.5) .. controls (7.1,-0.9) and (7.9,-0.9) .. (8,-0.5) -- cycle;
    
    \draw[thick] (7.5,-0.5) -- (7.5,0.2); % mast
    \fill[white] (7.5,-0.1) .. controls (8,0) and (8,0.1) .. (7.5,0.2) -- cycle; % sail
    \draw[thick] (7.5,-0.1) .. controls (8,0) and (8,0.1) .. (7.5,0.2) -- cycle;
    
    % Flag for the second boat
    \fill[red] (7.5,0.2) -- (7.7,0.3) -- (7.5,0.4) -- cycle;
    
    % Marking points A (top of lighthouse), B (bottom of lighthouse), C (first boat), D (second boat)
    \fill[black] (1,3.5) circle (0.05) node[above] {A};  % Point A (top of lighthouse)
    \fill[black] (1,-0.5) circle (0.05) node[below] {B}; % Point B (bottom of lighthouse)
 % Point C (position of boat 1)
    \fill[black] (7.5,-0.5) circle (0.05) node[below] {D}; % Point D (position of boat 2)
    
    % Draw lines connecting A-B, A-C, A-D
    \draw[thick] (1,3.5) -- (1,-0.5);  % Line AB% Line AC
    \draw[thick] (1,3.5) -- (7.5,-0.5); % Line AD
    
    % Draw line B-C-D
    \draw[thick] (1,-0.5) -- (5,-0.5) -- (7.5,-0.5);  % Line BCD
    
    % Add a line x parallel to BC through A
% Dashed line x, parallel to BC
    
    % Add a smiling sun in the sky
    \fill[yellow] (9,3.5) circle (0.5); % sun body
    \draw (9,3.7) arc[start angle=45, end angle=135, radius=0.2]; % smiling face
    \fill[black] (8.8,3.6) circle (0.05); % left eye
    \fill[black] (9.2,3.6) circle (0.05); % right eye
    \foreach \angle in {0,45,...,315} {
        \draw[thick, yellow] (9,3.5) -- ++(\angle:0.8);
    }
    
\end{tikzpicture}
}
\end{minipage}

\vspace{3pt}\hrule\vspace{3pt}\footnotesize\textcolor{gray!80!black}{\textbf{Các mã dùng chung hình (HINH\_204):} \texttt{L08\_C9\_B3\_15}}
\end{tcolorbox}
\vspace{4pt}


\begin{tcolorbox}[
    breakable,
    colback=white,
    colframe=blue!60!black,
    boxrule=0.8pt,
    arc=3pt,
    title={\textbf{Câu \#268} \quad \textbf{MÃ CÂU: \texttt{[L08\_C9\_B3\_16]}} \quad | \quad \textbf{Nhóm: \texttt{HINH\_207}} \quad | \quad \textbf{MapID: \texttt{8H9V3}} \hfill \textcolor{yellow!90!white}{[Hình gốc .pyc]}}
]
\textbf{Dạng toán:} Toán thực tế tính chiều cao của thang cứu hoả \hfill \textit{(Tự luận -- Vận dụng)}
\vspace{4pt}\hrule\vspace{6pt}
\begin{minipage}[t]{0.58\textwidth}
\vspace{0pt}
Cho biết chiếc thang của một xe cứu hỏa có chiều dài $AB = 17$ m, chân thang cách mặt đất một khoảng $BE = 2,2$ m và cách tường tòa nhà một khoảng $BC = 5$ m. Tính chiều cao $AD$ mà thang có thể vươn tới (kết quả làm tròn đến hàng phần mười).
\end{minipage}\hfill
\begin{minipage}[t]{0.39\textwidth}
\vspace{0pt}
\centering

\resizebox{\linewidth}{!}{%
\begin{tikzpicture}[line cap=round, line join=round, scale=0.3]
    % Định nghĩa các điểm (Dùng b_val để tính toán tọa độ)
    \coordinate (E) at (0,0);
    \coordinate (D) at (5,0);
    \coordinate (B) at (0,2.2);
    \coordinate (C) at (5,2.2);
    \coordinate (A) at (5, 18.44807680927192);

    % Vẽ các đoạn thẳng
    \draw[thick] (A) -- (B) node[midway, left, xshift=-2pt] {17m};
    \draw[thick] (B) -- (C) node[midway, above] {5m};
    \draw[thick] (A) -- (D);
    \draw[thick] (B) -- (E) node[midway, left] {2,2m}; % Hiển thị dấu phẩy ở đây
    \draw[thick] (E) -- (D);

    % Ký hiệu góc vuông tại C
    \draw (4.2, 2.2) -- (4.2, 3.0) -- (5, 3.0);

    % Gán nhãn các đỉnh
    \node[above] at (A) {$A$};
    \node[left] at (B) {$B$};
    \node[right] at (C) {$C$};
    \node[below] at (D) {$D$};
    \node[below] at (E) {$E$};
\end{tikzpicture}
}
\end{minipage}

\end{tcolorbox}
\vspace{4pt}


\begin{tcolorbox}[
    breakable,
    colback=white,
    colframe=blue!60!black,
    boxrule=0.8pt,
    arc=3pt,
    title={\textbf{Câu \#269} \quad \textbf{MÃ CÂU: \texttt{[L08\_C9\_B4\_01]}} \quad | \quad \textbf{Nhóm: \texttt{HINH\_208}} \quad | \quad \textbf{MapID: \texttt{8H9H4}} \hfill \textcolor{yellow!90!white}{[Hình gốc .pyc]}}
]
\textbf{Dạng toán:} Chỉ ra các tam giác đồng dạng 1 \hfill \textit{(Tự luận -- Thông hiểu)}
\vspace{4pt}\hrule\vspace{6pt}
\begin{minipage}[t]{0.58\textwidth}
\vspace{0pt}
a)  Cho $\triangle {HLM}$ vuông tại ${H}$, đường cao ${HO}$ (Hình a). Chứng minh rằng $\triangle {HLM} \backsim \triangle {OLH}$.

b) Tam giác vuông ${QTU}$ và tam giác vuông ${VYX}$ có các kích thước như Hình b có đồng dạng không?
\end{minipage}\hfill
\begin{minipage}[t]{0.39\textwidth}
\vspace{0pt}
\centering

\resizebox{\linewidth}{!}{%
\begin{tikzpicture}[scale=1, font=\footnotesize,>=stealth]
            %Gán số liệu.
            \def\canhAB{3};\def\canhCA{4};
            %Gán tọa độ.
            \coordinate (A) at (0,0);
            \coordinate (B) at (90:\canhAB);
            \coordinate (C) at (0:\canhCA);
            \coordinate (H) at ($(B)!(A)!(C)$);
            %Vẽ tam giác vuông BAC.
            \draw (A)--(B)--(C)--cycle (A)--(H);
            \draw pic[draw, angle radius=3mm, angle eccentricity=1.5]{right angle = B--A--C}
            pic[draw, angle radius=3mm, angle eccentricity=1.5]{right angle = B--H--A}
            ;
            %Gán nhãn.
                   \filldraw[black] (A) circle (1pt) node[below] {$H$};
    \filldraw[black] (B) circle (1pt) node[above] {$L$};
    \filldraw[black] (C) circle (1pt) node[below right] {$M$};
   \filldraw[black] (H) circle (1pt) node[right]{$O$};
            \draw (2,-0.5) node {Hình a};
        \end{tikzpicture}
        \begin{tikzpicture}[scale=1, font=\footnotesize,>=stealth]
            %Gán số liệu.
            \def\canhDE{2};\def\canhFD{4};
            %Gán tọa độ.
            \coordinate (D) at (0,0);
            \coordinate (E) at (90:\canhDE);
            \coordinate (F) at (0:\canhFD);
            %Gán số liệu.
            \def\canhMP{4};\def\canhNM{2};
            %Gán tọa độ.
            \coordinate (M) at (6,0);
            \coordinate (P) at ($(M)+(0,5)$);
            \coordinate (N) at ($(M)+(3,0)$);
            %Vẽ tam giác vuông EDF.
            \draw (D)--node[left]{$8$}(E)--(F)--node[below]{$5$}cycle;
            \draw pic[draw, angle radius=3mm, angle eccentricity=1.5]{right angle = E--D--F};
            %Vẽ tam giác vuông PMN.
            \draw (M)--node[left]{$10$}(P)--(N)--node[below]{$16$}cycle;
            \draw pic[draw, angle radius=3mm, angle eccentricity=1.5]{right angle = P--M--N};
            %Gán nhãn.
               \filldraw[black] (D) circle (1pt) node[below]{$Q$};
    \filldraw[black] (E) circle (1pt) node[above] {$T$};
    \filldraw[black] (F) circle (1pt) node[below] {$U$};
     \filldraw[black] (M) circle (1pt) node[below] {$V$};
        \filldraw[black] (P) circle (1pt) node[above right]{$X$};
    \filldraw[black] (N) circle (1pt) node[below left] {$Y$};
            \draw (5,-0.5) node {Hình b};
        \end{tikzpicture}
}
\end{minipage}

\end{tcolorbox}
\vspace{4pt}


\begin{tcolorbox}[
    breakable,
    colback=white,
    colframe=blue!60!black,
    boxrule=0.8pt,
    arc=3pt,
    title={\textbf{Câu \#270} \quad \textbf{MÃ CÂU: \texttt{[L08\_C9\_B4\_02]}} \quad | \quad \textbf{Nhóm: \texttt{HINH\_209}} \quad | \quad \textbf{MapID: \texttt{8H9H4}} \hfill \textcolor{yellow!90!white}{[Hình gốc .pyc]}}
]
\textbf{Dạng toán:} Chỉ ra các tam giác đồng dạng 2 \hfill \textit{(Tự luận -- Thông hiểu)}
\vspace{4pt}\hrule\vspace{6pt}
\begin{minipage}[t]{0.58\textwidth}
\vspace{0pt}
Trong hình bên dưới, tam giác nào đồng dạng với tam giác ${BEF}$?
\end{minipage}\hfill
\begin{minipage}[t]{0.39\textwidth}
\vspace{0pt}
\centering

\resizebox{\linewidth}{!}{%
\begin{tikzpicture}[scale=1, font=\footnotesize,>=stealth]%<DTools>
            %Gán số liệu.
            \def\canhAB{1};\def\canhCA{2};
            %Gán tọa độ.
            \coordinate (A) at (0,0);
            \coordinate (B) at (90:\canhAB);
            \coordinate (C) at (0:\canhCA);
            %Vẽ tam giác vuông BAC.
            \draw (A)--node[left]{$5$}(B)--(C)--node[below]{$15$}cycle;
            \draw pic[draw, angle radius=3mm, angle eccentricity=1.5]{right angle = B--A--C};
            %Gán nhãn.
                            \filldraw[black] (A) circle (1pt) node[below] {$B$};
    \filldraw[black] (B) circle (1pt) node[above] {$E$};
    \filldraw[black] (C) circle (1pt) node[below right] {$F$};
   
        \end{tikzpicture}
        \begin{tikzpicture}[scale=1, font=\footnotesize,>=stealth]
            %Gán số liệu.
            \def\canhDE{2};\def\canhFD{3};
            %Gán tọa độ.
            \coordinate (D) at (0,0);
            \coordinate (E) at (90:\canhDE);
            \coordinate (F) at (0:\canhFD);
            %Vẽ tam giác vuông EDF.
            \draw (D)--node[left]{$12$}(E)--(F)--node[below]{$36$}cycle;
            \draw pic[draw, angle radius=3mm, angle eccentricity=1.5]{right angle = E--D--F};
            %Gán nhãn.
                                   \filldraw[black] (D) circle (1pt) node[below] {$M$};
    \filldraw[black] (E) circle (1pt) node[above] {$O$};
    \filldraw[black] (F) circle (1pt) node[below right] {$R$};
        \end{tikzpicture}
        \begin{tikzpicture}[scale=1, font=\footnotesize,>=stealth]%<DTools>
            %Gán số liệu.
            \def\canhPQ{1};\def\canhKP{2.5};
            %Gán tọa độ.
            \coordinate (P) at (0,0);
            \coordinate (Q) at (90:\canhPQ);
            \coordinate (K) at (0:\canhKP);
            %Vẽ tam giác vuông QPK.
            \draw (P)--node[left]{$13$}(Q)--(K)--node[below]{$182$}cycle;
            \draw pic[draw, angle radius=3mm, angle eccentricity=1.5]{right angle = Q--P--K};
            %Gán nhãn.
                                            \filldraw[black] (P) circle (1pt) node[below] {$T$};
    \filldraw[black] (Q) circle (1pt) node[above] {$U$};
    \filldraw[black] (K) circle (1pt) node[below right] {$V$};
        \end{tikzpicture}
        \begin{tikzpicture}[scale=1, font=\footnotesize,>=stealth]
            %Gán số liệu.
            \def\canhHG{1};\def\canhIH{2};
            %Gán tọa độ.
            \coordinate (H) at (0,0);
            \coordinate (G) at (90:\canhHG);
            \coordinate (I) at (0:\canhIH);
            %Vẽ tam giác vuông GHI.
            \draw (H)--node[left]{$14$}(G)--(I)--node[below]{$42$}cycle;
            \draw pic[draw, angle radius=3mm, angle eccentricity=1.5]{right angle = G--H--I};
            %Gán nhãn.
                                      \filldraw[black] (H) circle (1pt) node[below] {$I$};
    \filldraw[black] (G) circle (1pt) node[above] {$W$};
    \filldraw[black] (I) circle (1pt) node[below right] {$Y$};
        \end{tikzpicture}
}
\end{minipage}

\end{tcolorbox}
\vspace{4pt}


\begin{tcolorbox}[
    breakable,
    colback=white,
    colframe=blue!60!black,
    boxrule=0.8pt,
    arc=3pt,
    title={\textbf{Câu \#271} \quad \textbf{MÃ CÂU: \texttt{[L08\_C9\_B4\_04]}} \quad | \quad \textbf{Nhóm: \texttt{HINH\_210}} \quad | \quad \textbf{MapID: \texttt{8H9V4}} \hfill \textcolor{yellow!90!white}{[Hình gốc .pyc]}}
]
\textbf{Dạng toán:} Chứng minh t.giác đ.dạng hệ thức lượng trong t.giác vuông tính cạnh. \hfill \textit{(Tự luận -- Vận dụng)}
\vspace{4pt}\hrule\vspace{6pt}
\begin{minipage}[t]{0.58\textwidth}
\vspace{0pt}
Quan sát hình bên.

 Chứng minh ${ \triangle BEL\backsim \triangle OEB}$;

  Chứng minh $BE^{2}=EO\cdot EL$;

 Biết ${EL=65}$ cm, ${EO=\dfrac{125}{13}}$ cm. Tính ${BL}$.
\end{minipage}\hfill
\begin{minipage}[t]{0.39\textwidth}
\vspace{0pt}
\centering

\resizebox{\linewidth}{!}{%
\begin{tikzpicture}[scale=1, font=\footnotesize,>=stealth]
            %Gán số liệu.
            \def\canhAB{3};\def\canhCA{4};
            %Gán tọa độ.
            \coordinate (A) at (0,0);
            \coordinate (B) at (90:\canhAB);
            \coordinate (C) at (0:\canhCA);
            \coordinate (H) at ($(B)!(A)!(C)$);
            %Vẽ tam giác vuông BAC.
            \draw (A)--(B)--(C)--cycle (A)--(H);
            \draw pic[draw, angle radius=3mm, angle eccentricity=1.5]{right angle = B--A--C}
            pic[draw, angle radius=3mm, angle eccentricity=1.5]{right angle = B--H--A}
            ;
           
    \filldraw[black] (A) circle (1pt) node[below left] {$B$};
    \filldraw[black] (B) circle (1pt) node[above left] {$E$};
    \filldraw[black] (C) circle (1pt) node[below right] {$L$};
   \filldraw[black] (H) circle (1pt) node[above right]{$O$};
        \end{tikzpicture}
}
\end{minipage}

\end{tcolorbox}
\vspace{4pt}


\begin{tcolorbox}[
    breakable,
    colback=white,
    colframe=blue!60!black,
    boxrule=0.8pt,
    arc=3pt,
    title={\textbf{Câu \#272} \quad \textbf{MÃ CÂU: \texttt{[L08\_C9\_B6\_01]}} \quad | \quad \textbf{Nhóm: \texttt{HINH\_211}} \quad | \quad \textbf{MapID: \texttt{8H9V6}} \hfill \textcolor{yellow!90!white}{[Hình gốc .pyc]}}
]
\textbf{Dạng toán:} Bài tổng hợp 1 \hfill \textit{(Tự luận -- Vận dụng)}
\vspace{4pt}\hrule\vspace{6pt}
\begin{minipage}[t]{0.58\textwidth}
\vspace{0pt}
Cho tam giác ${QKH}$ vuông tại ${Q}$ $(QK < QH)$, có đường cao $QE$ $(E \in KH)$.

 a) Chứng minh $\triangle {EKQ} \backsim \triangle {QKH}$.

b) Kẻ ${KF}$ là đường phân giác của tam giác ${QKH}$. Giả sử ${QK}={54}$ cm, $QH=72$ cm. Tính độ dài các đoạn thẳng ${KH},{FQ},{FH}$.

c) Gọi ${I}$ là giao điểm của ${QE}$ và ${KF}$. Chứng minh: $QI=QF$.
\end{minipage}\hfill
\begin{minipage}[t]{0.39\textwidth}
\vspace{0pt}
\centering

\resizebox{\linewidth}{!}{%
\begin{tikzpicture}[line join = round, line cap=round,>=stealth,font=\footnotesize,scale=1]
                \path 
                (0,0) coordinate (A)
                (4,0) coordinate (C)
                (0,3) coordinate (B)    
                ($(B)!(A)!(C)$) coordinate (H)           
                ($(B)!1cm!(A)$) coordinate (A1)
                ($(B)!1cm!(C)$) coordinate (C1)
                ($(A1)!0.5!(C1)$) coordinate (Bt)
                (intersection of B--Bt and A--C) coordinate (D) 
                (intersection of B--D and A--H) coordinate (I)  ;   
                \draw (B)--(A)--(C)--cycle (B)--(D) (A)--(H)        
                ;
                   
                \draw pic[draw, angle radius=3mm, angle eccentricity=1.2]{right angle = B--A--C};
                \draw pic[draw, angle radius=3mm, angle eccentricity=1.2]{right angle = A--H--C};
                \draw pic[draw, angle radius=5mm, angle eccentricity=1.2]{angle = A--B--C};
                \path 
                (A)--(B)node[left,pos=0.5,scale=1.2]{$54$}
                (C)--(A)node[below,pos=0.5,scale=1.2]{$72$}              
                ;       
                   \filldraw[black] (A) circle (2pt) node[below left] {$Q$};
    \filldraw[black] (B) circle (2pt) node[above left] {$K$};
    \filldraw[black] (C) circle (2pt) node[below right] {$H$};
   \filldraw[black] (D) circle (2pt) node[below left]{$F$};
    \filldraw[black] (H) circle (2pt) node[above right] {$E$};
    \filldraw[black] (I) circle (2pt) node[above left] {$I$};
            \end{tikzpicture}
}
\end{minipage}

\end{tcolorbox}
\vspace{4pt}


\begin{tcolorbox}[
    breakable,
    colback=white,
    colframe=blue!60!black,
    boxrule=0.8pt,
    arc=3pt,
    title={\textbf{Câu \#273} \quad \textbf{MÃ CÂU: \texttt{[L08\_C9\_B6\_02]}} \quad | \quad \textbf{Nhóm: \texttt{HINH\_212}} \quad | \quad \textbf{MapID: \texttt{8H9V6}} \hfill \textcolor{yellow!90!white}{[Hình gốc .pyc]}}
]
\textbf{Dạng toán:} Bài tổng hợp 2 \hfill \textit{(Tự luận -- Vận dụng)}
\vspace{4pt}\hrule\vspace{6pt}
\begin{minipage}[t]{0.58\textwidth}
\vspace{0pt}
Cho $\triangle {IJQ}$ vuông cân tại ${I}$ có đường cao ${IE}$.

 a) Chứng minh $\triangle {EJI}\backsim\triangle {IJQ}$ Từ đó suy ra $IJ^{2}=JE\cdot JQ$.

 b)  Gọi ${M}$, ${N}$ lần lượt là trung điểm của ${IJ}$ và ${IQ}$. Chứng minh tứ giác ${JMNQ}$ là hình thang cân và $MN=IE$.

 c) Gọi ${F}$ là giao điểm của ${IE}$ và ${QM}$. Chứng minh $JQ=3IF$.
\end{minipage}\hfill
\begin{minipage}[t]{0.39\textwidth}
\vspace{0pt}
\centering

\resizebox{\linewidth}{!}{%
\begin{tikzpicture}[line join = round, line cap=round,>=stealth,font=\footnotesize,scale=1]
    \path 
    (0,0) coordinate (A)
    (6,0) coordinate (C)
    (0,6) coordinate (B)
    ($(C)!(A)!(B)$) coordinate (H)
    ($(A)!.5!(B)$) coordinate (M)
    ($(A)!.5!(C)$) coordinate (N)
    (intersection of A--H and C--M) coordinate (K)
    ;
    
    \draw 
    (A)--(B)--(C)--(A)--(H) (M)--(N) (C)--(M) (B)--(M)node[midway]{$-$}--(A)node[midway]{$-$} (A)--(N)node[midway,scale=0.7]{$|$}--(C)node[midway,scale=0.7]{$|$}
    ;
    \draw   pic[draw, angle radius=0.5cm] {right angle=C--A--B};
    \draw   pic[draw, angle radius=.5cm] {right angle=A--H--C} ;
     \filldraw[black] (A) circle (2pt) node[above left] {$I$};
    \filldraw[black] (B) circle (2pt) node[above left] {$J$};
    \filldraw[black] (C) circle (2pt) node[above] {$Q$};
   \filldraw[black] (M) circle (2pt) node[above left]{$M$};
    \filldraw[black] (N) circle (2pt) node[below left] {$N$};
    \filldraw[black] (H) circle (2pt) node[above] {$E$};
     \filldraw[black] (K) circle (2pt) node[above] {$F$};
\end{tikzpicture}
}
\end{minipage}

\end{tcolorbox}
\vspace{4pt}


\section{Toán 9}

\subsection{Chương 1 -  Phương trình và HPT Bậc Nhất 2 ẩn}

\begin{tcolorbox}[
    breakable,
    colback=white,
    colframe=blue!60!black,
    boxrule=0.8pt,
    arc=3pt,
    title={\textbf{Câu \#274} \quad \textbf{MÃ CÂU: \texttt{[L09\_C1\_B1\_18]}} \quad | \quad \textbf{Nhóm: \texttt{HINH\_213}} \quad | \quad \textbf{MapID: \texttt{9D1V1-5}} \hfill \textcolor{yellow!90!white}{[Hình gốc .pyc]}}
]
\textbf{Dạng toán:} Đường thằng biểu diễn miền nghiệm của phương trình ax+0y=c. \hfill \textit{(Trắc nghiệm -- Vận dụng)}
\vspace{4pt}\hrule\vspace{6pt}
\begin{minipage}[t]{0.58\textwidth}
\vspace{0pt}
Cho đồ thị như hình dưới, hỏi đường thẳng trên đồ thị biểu diễn miền nghiệm của phương trình nào?

\vspace{4pt}\noindent
\textbf{A.} $0x+2 - y=0$ \quad \textbf{\textcolor{blue!75!black}{B. (Đ)}} $ -3x+0y=6$ \quad \textbf{C.} $ - 5 x + y=-1$ \quad \textbf{D.} $ 2x+0y=-2$
\end{minipage}\hfill
\begin{minipage}[t]{0.39\textwidth}
\vspace{0pt}
\centering

\resizebox{\linewidth}{!}{%
\begin{tikzpicture}[>=stealth]
    \draw[->](-6,0)--(6,0);% Vẽ trục hoành
    \draw[->](0,-3)--(0,3);%Vẽ trục tung

    \draw[thick](-2,-3)--(-2,3);
        \draw (6,0) node[above]{$x$} (0,3) node[left]{$y$}
        (0,0) node[below left]{$O$}
            (-2,0) node[below left]{$-2$};
    \end{tikzpicture}
}
\end{minipage}

\end{tcolorbox}
\vspace{4pt}


\subsection{Chương 2 -  Phương trình và bất phương trình bậc nhất 1 ẩn}

\begin{tcolorbox}[
    breakable,
    colback=white,
    colframe=blue!60!black,
    boxrule=0.8pt,
    arc=3pt,
    title={\textbf{Câu \#275} \quad \textbf{MÃ CÂU: \texttt{[L09\_C2\_B1\_29]}} \quad | \quad \textbf{Nhóm: \texttt{HINH\_214}} \quad | \quad \textbf{MapID: \texttt{9D2V1-3}} \hfill \textcolor{yellow!90!white}{[Hình gốc .pyc]}}
]
\textbf{Dạng toán:} SA-M3. bài toán thực tế về diện tích đất \hfill \textit{(Trả lời ngắn -- Vận dụng)}
\vspace{4pt}\hrule\vspace{6pt}
\begin{minipage}[t]{0.58\textwidth}
\vspace{0pt}
Cô Dung có một mảnh đất hình chữ nhật với chiều dài $33$ m và chiều rộng $25$ m. Cô dự định xây nhà trên mảnh đất đó và dành một phần diện tích đất để làm sân vườn như Hình dưới đây. Biết diện tích đất làm nhà là ${304}m^{2}$. Hỏi ${x}$ bằng bao nhiêu mét?


\textbf{Đáp án điền số:} 	exttt{9}
\end{minipage}\hfill
\begin{minipage}[t]{0.39\textwidth}
\vspace{0pt}
\centering

\resizebox{\linewidth}{!}{%
\begin{tikzpicture}
    % Vẽ hình sân vườn ngoài cùng
    \draw[thick] (0, 0) rectangle (10, 7.5);
    \node at (5, 6.5) {Sân vườn};
    
    % Vẽ hình nhà bên trong, góc dưới bên trái của nhà trùng với góc dưới bên trái của sân vườn
    \draw[thick] (0, 0) rectangle (7, 5.5);
    \node at (3.5, 2.25) {Nhà};
    
    % Ghi chú các chiều
    \draw[<->] (0, 8) -- (10, 8);
    \node at (5, 8.3) {$33$m};
    
    \draw[<->] (10.5, 0) -- (10.5, 7.5);
    \node at (11, 3.75) {$25$ m};
    
    \draw[<->] (0, 5.5) -- (0, 7.5);
    \node at (-0.4, 6.25) {x};
    
    \draw[<->] (7, -0.5) -- (10, -0.5);
    \node at (8.5, -0.8) {x+5};
    
\end{tikzpicture}
}
\end{minipage}

\end{tcolorbox}
\vspace{4pt}


\begin{tcolorbox}[
    breakable,
    colback=white,
    colframe=blue!60!black,
    boxrule=0.8pt,
    arc=3pt,
    title={\textbf{Câu \#276} \quad \textbf{MÃ CÂU: \texttt{[L09\_C2\_B1\_19]}} \quad | \quad \textbf{Nhóm: \texttt{HINH\_215}} \quad | \quad \textbf{MapID: \texttt{9D2V1-3}} \hfill \textcolor{yellow!90!white}{[Hình gốc .pyc]}}
]
\textbf{Dạng toán:} bài toán thực tế về diện tích đẩt \hfill \textit{(Tự luận -- Vận dụng)}
\vspace{4pt}\hrule\vspace{6pt}
\begin{minipage}[t]{0.58\textwidth}
\vspace{0pt}
Cô Dung có một mảnh đất hình chữ nhật với chiều dài $29$ m và chiều rộng $24$ m. Cô dự định xây nhà trên mảnh đất đó và dành một phần diện tích đất để làm sân vườn như Hình dưới đây. Biết diện tích đất làm nhà là ${340}m^{2}$. Hỏi ${x}$ bằng bao nhiêu mét?
\end{minipage}\hfill
\begin{minipage}[t]{0.39\textwidth}
\vspace{0pt}
\centering

\resizebox{\linewidth}{!}{%
\begin{tikzpicture}
    % Vẽ hình sân vườn ngoài cùng
    \draw[thick] (0, 0) rectangle (10, 7.5);
    \node at (5, 6.5) {Sân vườn};
    
    % Vẽ hình nhà bên trong, góc dưới bên trái của nhà trùng với góc dưới bên trái của sân vườn
    \draw[thick] (0, 0) rectangle (7, 5.5);
    \node at (3.5, 2.25) {Nhà};
    
    % Ghi chú các chiều
    \draw[<->] (0, 8) -- (10, 8);
    \node at (5, 8.3) {$29$m};
    
    \draw[<->] (10.5, 0) -- (10.5, 7.5);
    \node at (11, 3.75) {$24$ m};
    
    \draw[<->] (0, 5.5) -- (0, 7.5);
    \node at (-0.4, 6.25) {x};
    
    \draw[<->] (7, -0.5) -- (10, -0.5);
    \node at (8.5, -0.8) {x+2};
    
\end{tikzpicture}
}
\end{minipage}

\end{tcolorbox}
\vspace{4pt}


\subsection{Chương 4 -  Hệ thức lượng trong tam giác vuông}

\begin{tcolorbox}[
    breakable,
    colback=white,
    colframe=blue!60!black,
    boxrule=0.8pt,
    arc=3pt,
    title={\textbf{Câu \#277} \quad \textbf{MÃ CÂU: \texttt{[L09\_C4\_B1\_01]}} \quad | \quad \textbf{Nhóm: \texttt{HINH\_216}} \quad | \quad \textbf{MapID: \texttt{9H4N1-1}} \hfill \textcolor{yellow!90!white}{[Hình gốc .pyc]}}
]
\textbf{Dạng toán:} Nhận biết tỉ số lượng giác trong tam giác vuông ABC \hfill \textit{(Trắc nghiệm -- Nhận biết)}
\vspace{4pt}\hrule\vspace{6pt}
\begin{minipage}[t]{0.58\textwidth}
\vspace{0pt}
Cho hình vẽ sau: $\tan C$ bằng:

\vspace{4pt}\noindent
\textbf{A.} $\dfrac{AC}{BC}$ \quad \textbf{B.} $\dfrac{AC}{AB}$ \quad \textbf{\textcolor{blue!75!black}{C. (Đ)}} $\dfrac{AB}{AC}$ \quad \textbf{D.} $\dfrac{AB}{BC}$
\end{minipage}\hfill
\begin{minipage}[t]{0.39\textwidth}
\vspace{0pt}
\centering

\resizebox{\linewidth}{!}{%
\begin{tikzpicture}
        % Tọa độ cho tam giác vuông
        \coordinate (A) at (0, 0);
        \coordinate (B) at (4, 0);
        \coordinate (C) at (0, 3);

        % Vẽ các cạnh của tam giác
        \draw[thick] (A) -- (B) -- (C) -- cycle;

        % Ghi nhãn các điểm
        \node[below left] at (A) {A};
        \node[below right] at (B) {B};
        \node[above left] at (C) {C};

        % Vẽ dấu vuông tại điểm A
        \draw (0,0) rectangle +(0.3,0.3);
    \end{tikzpicture}
}
\end{minipage}

\vspace{3pt}\hrule\vspace{3pt}\footnotesize\textcolor{gray!80!black}{\textbf{Các mã dùng chung hình (HINH\_216):} \texttt{L09\_C4\_B2\_01}}
\end{tcolorbox}
\vspace{4pt}


\begin{tcolorbox}[
    breakable,
    colback=white,
    colframe=blue!60!black,
    boxrule=0.8pt,
    arc=3pt,
    title={\textbf{Câu \#278} \quad \textbf{MÃ CÂU: \texttt{[L09\_C4\_B1\_02]}} \quad | \quad \textbf{Nhóm: \texttt{HINH\_217}} \quad | \quad \textbf{MapID: \texttt{9H4N1-1}} \hfill \textcolor{yellow!90!white}{[Hình gốc .pyc]}}
]
\textbf{Dạng toán:} Nhận biết tỉ số lượng giác trong tam giác vuông ABC có đường cao AH \hfill \textit{(Trắc nghiệm -- Nhận biết)}
\vspace{4pt}\hrule\vspace{6pt}
\begin{minipage}[t]{0.58\textwidth}
\vspace{0pt}
Cho hình vẽ $\sin B$ bằng bao nhiêu?

\vspace{4pt}\noindent
\textbf{A.} $\dfrac{BH}{AB}$ \quad \textbf{\textcolor{blue!75!black}{B. (Đ)}} $\dfrac{AH}{AB}$ \quad \textbf{C.} $\dfrac{AB}{BC}$ \quad \textbf{D.} $\dfrac{AC}{BH}$
\end{minipage}\hfill
\begin{minipage}[t]{0.39\textwidth}
\vspace{0pt}
\centering

\resizebox{\linewidth}{!}{%
\begin{tikzpicture}[scale=0.6]  % Thu nhỏ toàn bộ hình vẽ
            % Tạo các điểm A, B, C, H
            \coordinate [label=above:$A$] (A) at (0, 4);
            \coordinate [label=below left:$B$] (B) at (-3, 0);
            \coordinate [label=below right:$C$] (C) at (4, 0);
            \coordinate [label=below:$H$] (H) at (0, 0);
            
            % Vẽ các cạnh tam giác
            \draw (B) -- (A) -- (C) -- cycle;
            
            % Vẽ đường cao từ A xuống H
            \draw  (A) -- (H);
            
            % Vẽ góc vuông tại H
            \draw (H) rectangle ++(0.3, 0.3);   

            %Vẽ góc vuông tại A
            \pic[draw,angle radius=10]{right angle=B--A--C};
        \end{tikzpicture}
}
\end{minipage}

\end{tcolorbox}
\vspace{4pt}


\begin{tcolorbox}[
    breakable,
    colback=white,
    colframe=blue!60!black,
    boxrule=0.8pt,
    arc=3pt,
    title={\textbf{Câu \#279} \quad \textbf{MÃ CÂU: \texttt{[L09\_C4\_B1\_16]}} \quad | \quad \textbf{Nhóm: \texttt{HINH\_218}} \quad | \quad \textbf{MapID: \texttt{9H4V1-5}} \hfill \textcolor{yellow!90!white}{[Hình gốc .pyc]}}
]
\textbf{Dạng toán:} Bài toán thực tế tính chiều cao của cái thang. \hfill \textit{(Trắc nghiệm -- Vận dụng)}
\vspace{4pt}\hrule\vspace{6pt}
\begin{minipage}[t]{0.58\textwidth}
\vspace{0pt}
Một cái thang dài 4,1 m tựa vào tường làm thành góc $\alpha = 62^{\circ}$ với mặt đất. Khoảng cách từ chân thang đến tường (làm tròn đến 2 chữ số thập phân) là:

\vspace{4pt}\noindent
\textbf{A.} 1,93m \quad \textbf{B.} 1,94m \quad \textbf{C.} 1,91m \quad \textbf{\textcolor{blue!75!black}{D. (Đ)}} $1,92$m
\end{minipage}\hfill
\begin{minipage}[t]{0.39\textwidth}
\vspace{0pt}
\centering

\resizebox{\linewidth}{!}{%
\begin{tikzpicture}[scale=0.8, line join=round, line cap=round,>=stealth]   
    \path
    (-1,0) coordinate (P)+(-80:1) 
    (P)+(60:5) coordinate (Q)
    (P)+(-30:0.2) coordinate (P1)+(-90:.3) node{$B$}
    (Q)+(-30:0.2) coordinate (Q1)+(0:.3) node{$A$}
    (P)+(150:0.9) coordinate (PP)
    (Q)+(150:0.9) coordinate (QQ)
    (P1)+(150:0.9) coordinate (PP1)
    (Q1)+(150:0.9) coordinate (QQ1)
    (4,0) coordinate (SD)
    (SD)+(150:10) coordinate (ST)
    (Q1)+(-90:1) coordinate (QT)
    (Q)+(0:1) coordinate (Q2)
    (intersection of SD--ST and Q1--QT) coordinate (C)+(-90:.3) node{$C$}
    ;
    
    \foreach \b in {0.1,0.2,...,0.9} {
        \path 
        ($(P)!\b!(Q)$) coordinate (m)
        (m)+(50:0.2)coordinate (m1)
        (m)+(150:0.7)coordinate (m2)
        (m1)+(150:0.7)coordinate (m3);
        \draw[fill=cyan!50](m)--(m1)--(m3)--(m2)--(m);}
    \draw[fill=cyan!50] (P)--(Q)--(Q1)--(P1)--(P);
    \draw[fill=cyan!50] (PP)--(QQ)--(QQ1)--(PP1)--(PP);
    \draw  (P1)--(C)--(Q1);
    \draw (Q1)--++(-30:2) (Q1)--++(150:3);
    \draw (Q2)--++(-30:2) (Q2)--++(150:3);
    \pic[draw,angle radius=10]{right angle=P1--C--Q1};
    \pic [draw,angle eccentricity=1.5, angle radius=1.2cm] {angle=C--P1--Q1} ;
    \draw(0,0.25)node[above]{$62^{\circ}$};
\end{tikzpicture}
}
\end{minipage}

\vspace{3pt}\hrule\vspace{3pt}\footnotesize\textcolor{gray!80!black}{\textbf{Các mã dùng chung hình (HINH\_218):} \texttt{L09\_C4\_B1\_29}}
\end{tcolorbox}
\vspace{4pt}


\begin{tcolorbox}[
    breakable,
    colback=white,
    colframe=blue!60!black,
    boxrule=0.8pt,
    arc=3pt,
    title={\textbf{Câu \#280} \quad \textbf{MÃ CÂU: \texttt{[L09\_C4\_B1\_17]}} \quad | \quad \textbf{Nhóm: \texttt{HINH\_219}} \quad | \quad \textbf{MapID: \texttt{9H4V1-5}} \hfill \textcolor{yellow!90!white}{[Hình gốc .pyc]}}
]
\textbf{Dạng toán:} Bài toán thực tế tính chiều cao của toà nhà \hfill \textit{(Trắc nghiệm -- Vận dụng)}
\vspace{4pt}\hrule\vspace{6pt}
\begin{minipage}[t]{0.58\textwidth}
\vspace{0pt}
Khi tia nắng mặt trời tạo với mặt đất một góc $84^{\circ}$ thì bóng của một tòa nhà cao tầng trên mặt đất dài $41$ m (hình vẽ). Chiều cao của tòa nhà (làm tròn đến hàng đơn vị) là bao nhiêu?

\vspace{4pt}\noindent
\textbf{A.} $385$ m \quad \textbf{\textcolor{blue!75!black}{B. (Đ)}} $390$ m \quad \textbf{C.} $410$ m \quad \textbf{D.} $400$ m
\end{minipage}\hfill
\begin{minipage}[t]{0.39\textwidth}
\vspace{0pt}
\centering

\resizebox{\linewidth}{!}{%
\begin{tikzpicture}[scale=0.8]  % Thu nhỏ toàn bộ hình vẽ

        % Tạo các điểm A, B, C
        \coordinate [label=below left:$A$] (A) at (0, 0);
        \coordinate [label=above:$B$] (B) at (0, 7);  % Đặt điểm B ở cao hơn A một khoảng cố định
        \coordinate [label=below right:$C$] (C) at (4, 0);  % Đặt điểm C ở bên phải A một khoảng cố định

        % Vẽ tam giác ABC
        \draw [thick] (A) -- (B) -- (C) -- cycle;

        % Vẽ góc vuông tại A
        \draw (A) rectangle +(0.5, 0.5);

        % Thêm nhãn độ dài cạnh
        \node at (2, -0.3) {41 m};  % Nhãn độ dài cạnh AC với vị trí cố định

        \draw 

    (C) pic[draw,angle radius = 15] {angle = B--C--A} node[shift={(-25:-25pt)}]{$ 84^\circ $}   
    ;

        % Vẽ tòa nhà bám sát cạnh AB
        \fill[pattern=bricks] (A) rectangle ++(-1.3, 7); % Tòa nhà với chiều cao cố định

        \end{tikzpicture}
}
\end{minipage}

\end{tcolorbox}
\vspace{4pt}


\begin{tcolorbox}[
    breakable,
    colback=white,
    colframe=blue!60!black,
    boxrule=0.8pt,
    arc=3pt,
    title={\textbf{Câu \#281} \quad \textbf{MÃ CÂU: \texttt{[L09\_C4\_B1\_18]}} \quad | \quad \textbf{Nhóm: \texttt{HINH\_220}} \quad | \quad \textbf{MapID: \texttt{9H4V1-5}} \hfill \textcolor{yellow!90!white}{[Hình gốc .pyc]}}
]
\textbf{Dạng toán:} Bài toán thực tế tính góc \hfill \textit{(Trắc nghiệm -- Vận dụng)}
\vspace{4pt}\hrule\vspace{6pt}
\begin{minipage}[t]{0.58\textwidth}
\vspace{0pt}
Tại một thời điểm trong ngày, một cột bê tông AB bên đường cao $4,3$m có bóng in trên mặt đất là AC có chiều dài là $4$m (như hình vẽ). Số đo góc tạo bởi tia nắng mặt trời tạo với mặt đất thời điểm đó là:

\vspace{4pt}\noindent
\textbf{A.} $44,06^{\circ}$ \quad \textbf{B.} $53,08^{\circ}$ \quad \textbf{\textcolor{blue!75!black}{C. (Đ)}} $47,07^{\circ}$ \quad \textbf{D.} $50,09^{\circ}$
\end{minipage}\hfill
\begin{minipage}[t]{0.39\textwidth}
\vspace{0pt}
\centering

\resizebox{\linewidth}{!}{%
\begin{tikzpicture}[scale=0.8]
            % Tạo các điểm A, B, C
            \coordinate [label=below left:$A$] (A) at (0, 0);
            \coordinate [label=above:$B$] (B) at (0, 7);
            \coordinate [label=below right:$C$] (C) at (4, 0);

            % Vẽ tam giác ABC
            \draw [thick] (A) -- (B) -- (C) -- cycle;

            % Vẽ góc vuông tại A
            \draw (A) rectangle +(0.5, 0.5);

            % Vẽ cột bê tông bám sát cạnh AB và màu nâu
            \fill[brown!70] (-0.5, 0) rectangle +(0.5, 7); % Cột bê tông với chiều cao cố định và màu nâu
            \draw[thick] (-0.5, 0) rectangle +(0.5, 7); % Vẽ đường viền cho cột bê tông, dày hơn

            % Draw details on the concrete column
        \draw[gray, thick] (-0.5, 1) -- (0, 1); % Horizontal line at y = 1
        \draw[gray, thick] (-0.5, 2) -- (0, 2); % Horizontal line at y = 2
        \draw[gray, thick] (-0.5, 3) -- (0, 3); % Horizontal line at y = 3
        \draw[gray, thick] (-0.5, 4) -- (0, 4); % Horizontal line at y = 4
        \draw[gray, thick] (-0.5, 5) -- (0, 5); % Horizontal line at y = 5
        \draw[gray, thick] (-0.5, 6) -- (0, 6); % Horizontal line at y = 6

            % Tạo hiệu ứng bóng cho cột bê tông
            \fill[black, opacity=0.3] (-0.45, 0) rectangle +(0.5, 7); % Bóng mờ bên trái cột bê tông

            % Thêm nhãn độ dài cạnh
            \node at (2, -0.3) {4 m};  % Nhãn độ dài cạnh AC với vị trí cố định

            % Thêm nhãn độ dài cạnh AB (di chuyển ra xa cột bê tông)
            \node at (0.75, 3.5) {4,3 m};  
        \end{tikzpicture}
}
\end{minipage}

\end{tcolorbox}
\vspace{4pt}


\begin{tcolorbox}[
    breakable,
    colback=white,
    colframe=blue!60!black,
    boxrule=0.8pt,
    arc=3pt,
    title={\textbf{Câu \#282} \quad \textbf{MÃ CÂU: \texttt{[L09\_C4\_B1\_19]}} \quad | \quad \textbf{Nhóm: \texttt{HINH\_221}} \quad | \quad \textbf{MapID: \texttt{9H4V1-5}} \hfill \textcolor{yellow!90!white}{[Hình gốc .pyc]}}
]
\textbf{Dạng toán:} Bài toán thực tế tính chiều cao của cây \hfill \textit{(Trắc nghiệm -- Vận dụng)}
\vspace{4pt}\hrule\vspace{6pt}
\begin{minipage}[t]{0.58\textwidth}
\vspace{0pt}
Khi tia sáng mặt trời tạo với mặt đất một góc $\alpha={51}^{\circ}$ thì bóng của một cái cây trên mặt đất có độ dài $x={6}$m. Chiều cao của cái cây (kết quả làm tròn đến hàng đơn vị) là

\vspace{4pt}\noindent
\textbf{A.} $6$m \quad \textbf{B.} $9$m \quad \textbf{C.} $8$m \quad \textbf{\textcolor{blue!75!black}{D. (Đ)}} ${7}m$
\end{minipage}\hfill
\begin{minipage}[t]{0.39\textwidth}
\vspace{0pt}
\centering

\resizebox{\linewidth}{!}{%
\begin{tikzpicture}[scale=1, font=\footnotesize, line cap=round, line join=round]

    % Vẽ ông mặt trời
    \fill[yellow] (-3,7) circle (0.5cm); % Di chuyển mặt trời ra xa một chút

    % Vẽ các tia nắng sinh động bao quanh ông mặt trời (từng tia riêng lẻ)
    \pgfmathsetmacro{\radius}{0.8} % Khoảng cách từ tâm mặt trời đến điểm bắt đầu của tia nắng
    
    % Vẽ từng tia nắng riêng lẻ
    \pgfmathsetmacro{\lengthA}{0.5 + 0.4 * sin(0 * 8.8)}
    \draw[orange, thick] (-3,7) ++(0:\radius) -- ++(0:\lengthA);

    \pgfmathsetmacro{\lengthB}{0.5 + 0.4 * sin(15 * 8.8)}
    \draw[orange, thick] (-3,7) ++(15:\radius) -- ++(15:\lengthB);

    \pgfmathsetmacro{\lengthC}{0.5 + 0.4 * sin(30 * 8.8)}
    \draw[orange, thick] (-3,7) ++(30:\radius) -- ++(30:\lengthC);

    \pgfmathsetmacro{\lengthD}{0.5 + 0.4 * sin(45 * 8.8)}
    \draw[orange, thick] (-3,7) ++(45:\radius) -- ++(45:\lengthD);

    \pgfmathsetmacro{\lengthE}{0.5 + 0.4 * sin(60 * 8.8)}
    \draw[orange, thick] (-3,7) ++(60:\radius) -- ++(60:\lengthE);

    \pgfmathsetmacro{\lengthF}{0.5 + 0.4 * sin(75 * 8.8)}
    \draw[orange, thick] (-3,7) ++(75:\radius) -- ++(75:\lengthF);

    \pgfmathsetmacro{\lengthG}{0.5 + 0.4 * sin(90 * 8.8)}
    \draw[orange, thick] (-3,7) ++(90:\radius) -- ++(90:\lengthG);

    \pgfmathsetmacro{\lengthH}{0.5 + 0.4 * sin(105 * 8.8)}
    \draw[orange, thick] (-3,7) ++(105:\radius) -- ++(105:\lengthH);

    \pgfmathsetmacro{\lengthI}{0.5 + 0.4 * sin(120 * 8.8)}
    \draw[orange, thick] (-3,7) ++(120:\radius) -- ++(120:\lengthI);

    \pgfmathsetmacro{\lengthJ}{0.5 + 0.4 * sin(135 * 8.8)}
    \draw[orange, thick] (-3,7) ++(135:\radius) -- ++(135:\lengthJ);

    \pgfmathsetmacro{\lengthK}{0.5 + 0.4 * sin(150 * 8.8)}
    \draw[orange, thick] (-3,7) ++(150:\radius) -- ++(150:\lengthK);

    \pgfmathsetmacro{\lengthL}{0.5 + 0.4 * sin(165 * 8.8)}
    \draw[orange, thick] (-3,7) ++(165:\radius) -- ++(165:\lengthL);

    \pgfmathsetmacro{\lengthM}{0.5 + 0.4 * sin(180 * 8.8)}
    \draw[orange, thick] (-3,7) ++(180:\radius) -- ++(180:\lengthM);

    \pgfmathsetmacro{\lengthN}{0.5 + 0.4 * sin(195 * 8.8)}
    \draw[orange, thick] (-3,7) ++(195:\radius) -- ++(195:\lengthN);

    \pgfmathsetmacro{\lengthO}{0.5 + 0.4 * sin(210 * 8.8)}
    \draw[orange, thick] (-3,7) ++(210:\radius) -- ++(210:\lengthO);

    \pgfmathsetmacro{\lengthP}{0.5 + 0.4 * sin(225 * 8.8)}
    \draw[orange, thick] (-3,7) ++(225:\radius) -- ++(225:\lengthP);

    \pgfmathsetmacro{\lengthQ}{0.5 + 0.4 * sin(240 * 8.8)}
    \draw[orange, thick] (-3,7) ++(240:\radius) -- ++(240:\lengthQ);

    \pgfmathsetmacro{\lengthR}{0.5 + 0.4 * sin(255 * 8.8)}
    \draw[orange, thick] (-3,7) ++(255:\radius) -- ++(255:\lengthR);

    \pgfmathsetmacro{\lengthS}{0.5 + 0.4 * sin(270 * 8.8)}
    \draw[orange, thick] (-3,7) ++(270:\radius) -- ++(270:\lengthS);

    \pgfmathsetmacro{\lengthT}{0.5 + 0.4 * sin(285 * 8.8)}
    \draw[orange, thick] (-3,7) ++(285:\radius) -- ++(285:\lengthT);

    \pgfmathsetmacro{\lengthU}{0.5 + 0.4 * sin(300 * 8.8)}
    \draw[orange, thick] (-3,7) ++(300:\radius) -- ++(300:\lengthU);

    \pgfmathsetmacro{\lengthV}{0.5 + 0.4 * sin(315 * 8.8)}
    \draw[orange, thick] (-3,7) ++(315:\radius) -- ++(315:\lengthV);

    \pgfmathsetmacro{\lengthW}{0.5 + 0.4 * sin(330 * 8.8)}
    \draw[orange, thick] (-3,7) ++(330:\radius) -- ++(330:\lengthW);

    \pgfmathsetmacro{\lengthX}{0.5 + 0.4 * sin(345 * 8.8)}
    \draw[orange, thick] (-3,7) ++(345:\radius) -- ++(345:\lengthX);

    % Vẽ cái cây cao hơn
    \fill[brown] (0,0) rectangle (0.5,4); % Tăng chiều cao của cái cây
    
    % Vẽ tán cây bằng các hình tròn chồng lên nhau, thẳng hàng với thân cây và xếp cao dần
    \fill[green] (0.25,3.5) circle (0.8cm); % Vòng tròn lớn dưới cùng
    \fill[green] (0.25,4.5) circle (1.2cm); % Vòng tròn trung bình ở giữa
    \fill[green] (0.25,4) circle (1cm); % Vòng tròn nhỏ ở trên cùng
    \fill[brown] (0.2,0) ellipse (0.8 and 0.24); % Elip với kích thước tuỳ chỉnh

    % Định nghĩa các điểm A, B, C
    \path
    (0.25,0) coordinate (A) % Gốc cây
    (0.25,5.6) coordinate (B) % Ngọn cây
    (4,0) coordinate (C); % Điểm trên mặt đất, đã di chuyển ra xa
    
    % Vẽ tam giác ABC
    \draw (A) -- (B) -- (C) -- cycle; % Nối các điểm A, B, C

    % Gắn nhãn các đoạn thẳng và góc
    \draw ($(A)!.5!(C)$) node[below] {$x$}; % Nhãn đoạn thẳng AC
    \node at ($(C)+(155:1)$) {$\alpha$}; % Ký hiệu góc BCA

    % Đánh dấu các điểm A, B, C
    \foreach \i/\p in {A/-90,B/90,C/0}  
    \fill[black] (\i) circle (1.5pt)  
    node[shift={(\p:3mm)},black] {$\i$}; % Chữ cái màu đen

\end{tikzpicture}
}
\end{minipage}

\vspace{3pt}\hrule\vspace{3pt}\footnotesize\textcolor{gray!80!black}{\textbf{Các mã dùng chung hình (HINH\_221):} \texttt{L09\_C4\_B2\_11, L09\_C4\_B2\_48, L09\_C4\_B2\_33}}
\end{tcolorbox}
\vspace{4pt}


\begin{tcolorbox}[
    breakable,
    colback=white,
    colframe=blue!60!black,
    boxrule=0.8pt,
    arc=3pt,
    title={\textbf{Câu \#283} \quad \textbf{MÃ CÂU: \texttt{[L09\_C4\_B1\_20]}} \quad | \quad \textbf{Nhóm: \texttt{HINH\_222}} \quad | \quad \textbf{MapID: \texttt{9H4V1-5}} \hfill \textcolor{yellow!90!white}{[Hình gốc .pyc]}}
]
\textbf{Dạng toán:} Bài toán thực tế tính vận tốc của xe khi nhìn từ một toà nhà \hfill \textit{(Trắc nghiệm -- Vận dụng)}
\vspace{4pt}\hrule\vspace{6pt}
\begin{minipage}[t]{0.58\textwidth}
\vspace{0pt}
Một người cao ${1,8}m$ đứng trên sân thượng của một toà nhà cao ${79}m$ đặt ống nhòm hạ xuống dưới hợp với phương ngang một góc $4^{\circ}$ thì nhìn thấy một chiếc xe đang tiến lại gần toà nhà. Sau $2$ phút thì xe đó đến cổng toà nhà, hỏi xe đã đi với vận tốc bao nhiêu? (làm tròn một chữ số thập phân sau dấu phẩy).

\vspace{4pt}\noindent
\textbf{\textcolor{blue!75!black}{A. (Đ)}} ${34,7}km/h$ \quad \textbf{B.} $32,7km/h$ \quad \textbf{C.} $37,7km/h$ \quad \textbf{D.} $36,7km/h$
\end{minipage}\hfill
\begin{minipage}[t]{0.39\textwidth}
\vspace{0pt}
\centering

\resizebox{\linewidth}{!}{%
\begin{tikzpicture}[scale=1,font=\footnotesize,line cap=round,line join=round]
    %Tòa nhà
    \path
    (3,0) coordinate (B)
    (-1,0) coordinate (B')
    (3,4) coordinate (A')
    (3,4.4) coordinate (A)
    (15,0) coordinate (C)
    (-1,4.4) coordinate (O)
    (7,4.4) coordinate (D)
    (-0.1,4) coordinate (D')
    (-1,4) coordinate (O')
    (-0.5,4) coordinate (C')
    (-0.5,0) coordinate (E)
    (2.7,4) coordinate (E')
    (2.7,4.4) coordinate (F)
    ;
    
    \foreach \i in {0,0.5,1,1.5,2,2.5} {
        \foreach \j in {0,0.8,1.6,2.4,3.2} {
            \draw[fill=magenta!10!white] (\i,\j) rectangle +(0.5,0.8);
        }
    }
    
    %ban công 
    \foreach \i in {0.1,0.6,1.1,1.6,2.1,2.6} {
        \foreach \j in {0,0.8,1.6,2.4,3.2} {
            \draw[fill=gray] (\i,\j) rectangle +(0.3,0.6);
        }
    }
    
    \foreach \j in {0,0.8,1.6,2.4,3.2} 
        \draw[fill=cyan!10!white] (0,\j) rectangle +(3,0.3);
    
    \draw[ultra thick,cyan](-0.1,4)--(3.1,4);
    \draw (A)--(B)--(C)--cycle (O)--(D) node[above]{$x$}
          (B)--(B') (D')--(O') (C')--(E) node[pos=.5,left]{$79m$ } 
          (E')--(F) node[pos=.5,left]{$1,8m$ } 
          (15.1,.4) node{Xe};
    
    \foreach \i/\p in {A/90,B/-90,C/0}  
        \fill[black] (\i) circle (1.5pt) node[shift={(\p:3mm)}]{$\i$};
        \draw 

    (A) pic[draw,angle radius = 15] {angle = C--A--D} node[shift={(-25:-25pt)}]{$ 4^\circ $}   
    ;
    %\path pic[draw,angle radius=4mm,opacity=.6,"$4^{\circ}$",angle eccentricity=1.7] {angle = C--A--D};
    \newcommand{\gocv}[4][black]{\draw[#1] ($(#3)!5pt!(#2)$)--($(#3)!2!($($(#3)!5pt!(#2)$)!.5!($(#3)!5pt!(#4)$)$)$)--($(#3)!5pt!(#4)$);}
    \gocv{C}{B}{A}  
    \end{tikzpicture}
}
\end{minipage}

\end{tcolorbox}
\vspace{4pt}


\begin{tcolorbox}[
    breakable,
    colback=white,
    colframe=blue!60!black,
    boxrule=0.8pt,
    arc=3pt,
    title={\textbf{Câu \#284} \quad \textbf{MÃ CÂU: \texttt{[L09\_C4\_B1\_29]}} \quad | \quad \textbf{Nhóm: \texttt{HINH\_218}} \quad | \quad \textbf{MapID: \texttt{9H4V1-5}} \hfill \textcolor{yellow!90!white}{[Hình gốc .pyc]}}
]
\textbf{Dạng toán:} Bài toán thực tế tính chiều cao của cái thang. \hfill \textit{(Tự luận -- Vận dụng)}
\vspace{4pt}\hrule\vspace{6pt}
\begin{minipage}[t]{0.58\textwidth}
\vspace{0pt}
Một cái thang dài 2,9 m tựa vào tường làm thành góc $\alpha = 60^{\circ}$ với mặt đất. Tính khoảng cách từ chân thang đến tường (làm tròn đến 2 chữ số thập phân).
\end{minipage}\hfill
\begin{minipage}[t]{0.39\textwidth}
\vspace{0pt}
\centering

\resizebox{\linewidth}{!}{%
\begin{tikzpicture}[scale=0.8, line join=round, line cap=round,>=stealth]   
    \path
    (-1,0) coordinate (P)+(-80:1) 
    (P)+(60:5) coordinate (Q)
    (P)+(-30:0.2) coordinate (P1)+(-90:.3) node{$B$}
    (Q)+(-30:0.2) coordinate (Q1)+(0:.3) node{$A$}
    (P)+(150:0.9) coordinate (PP)
    (Q)+(150:0.9) coordinate (QQ)
    (P1)+(150:0.9) coordinate (PP1)
    (Q1)+(150:0.9) coordinate (QQ1)
    (4,0) coordinate (SD)
    (SD)+(150:10) coordinate (ST)
    (Q1)+(-90:1) coordinate (QT)
    (Q)+(0:1) coordinate (Q2)
    (intersection of SD--ST and Q1--QT) coordinate (C)+(-90:.3) node{$C$}
    ;
    
    \foreach \b in {0.1,0.2,...,0.9} {
        \path 
        ($(P)!\b!(Q)$) coordinate (m)
        (m)+(50:0.2)coordinate (m1)
        (m)+(150:0.7)coordinate (m2)
        (m1)+(150:0.7)coordinate (m3);
        \draw[fill=cyan!50](m)--(m1)--(m3)--(m2)--(m);}
    \draw[fill=cyan!50] (P)--(Q)--(Q1)--(P1)--(P);
    \draw[fill=cyan!50] (PP)--(QQ)--(QQ1)--(PP1)--(PP);
    \draw  (P1)--(C)--(Q1);
    \draw (Q1)--++(-30:2) (Q1)--++(150:3);
    \draw (Q2)--++(-30:2) (Q2)--++(150:3);
    \pic[draw,angle radius=10]{right angle=P1--C--Q1};
    \pic [draw,angle eccentricity=1.5, angle radius=1.2cm] {angle=C--P1--Q1} ;
    \draw(0,0.25)node[above]{$\alpha$};
\end{tikzpicture}
}
\end{minipage}

\vspace{3pt}\hrule\vspace{3pt}\footnotesize\textcolor{gray!80!black}{\textbf{Các mã dùng chung hình (HINH\_218):} \texttt{L09\_C4\_B1\_16}}
\end{tcolorbox}
\vspace{4pt}


\begin{tcolorbox}[
    breakable,
    colback=white,
    colframe=blue!60!black,
    boxrule=0.8pt,
    arc=3pt,
    title={\textbf{Câu \#285} \quad \textbf{MÃ CÂU: \texttt{[L09\_C4\_B2\_01]}} \quad | \quad \textbf{Nhóm: \texttt{HINH\_216}} \quad | \quad \textbf{MapID: \texttt{9H4N2-4}} \hfill \textcolor{yellow!90!white}{[Hình gốc .pyc]}}
]
\textbf{Dạng toán:} Nhận biết hệ thức giữa cạnh và góc trong tam giác vuông. \hfill \textit{(Trắc nghiệm -- Nhận biết)}
\vspace{4pt}\hrule\vspace{6pt}
\begin{minipage}[t]{0.58\textwidth}
\vspace{0pt}
Cho hình vẽ sau. Khẳng định nào sau đây là đúng

\vspace{4pt}\noindent
\textbf{\textcolor{blue!75!black}{A. (Đ)}} $AB=BC.\cos B$ \quad \textbf{B.} $AC=BC.\sin C$ \quad \textbf{C.} $AB=AC.\cot C$ \quad \textbf{D.} $AC=AB.\sin B$
\end{minipage}\hfill
\begin{minipage}[t]{0.39\textwidth}
\vspace{0pt}
\centering

\resizebox{\linewidth}{!}{%
\begin{tikzpicture}
        % Tọa độ cho tam giác vuông
        \coordinate (A) at (0, 0);
        \coordinate (B) at (4, 0);
        \coordinate (C) at (0, 3);

        % Vẽ các cạnh của tam giác
        \draw[thick] (A) -- (B) -- (C) -- cycle;

        % Ghi nhãn các điểm
        \node[below left] at (A) {A};
        \node[below right] at (B) {B};
        \node[above left] at (C) {C};

        % Vẽ dấu vuông tại điểm A
        \draw (0,0) rectangle +(0.3,0.3);
    \end{tikzpicture}
}
\end{minipage}

\vspace{3pt}\hrule\vspace{3pt}\footnotesize\textcolor{gray!80!black}{\textbf{Các mã dùng chung hình (HINH\_216):} \texttt{L09\_C4\_B1\_01}}
\end{tcolorbox}
\vspace{4pt}


\begin{tcolorbox}[
    breakable,
    colback=white,
    colframe=blue!60!black,
    boxrule=0.8pt,
    arc=3pt,
    title={\textbf{Câu \#286} \quad \textbf{MÃ CÂU: \texttt{[L09\_C4\_B2\_08]}} \quad | \quad \textbf{Nhóm: \texttt{HINH\_223}} \quad | \quad \textbf{MapID: \texttt{9H4V2-5}} \hfill \textcolor{yellow!90!white}{[Hình gốc .pyc]}}
]
\textbf{Dạng toán:} Bài toán thực tế tính vận tốc của máy bay khi bay lên \hfill \textit{(Trắc nghiệm -- Vận dụng)}
\vspace{4pt}\hrule\vspace{6pt}
\begin{minipage}[t]{0.58\textwidth}
\vspace{0pt}
Một máy bay vừa phi lên theo đường bay hợp với mặt đất một góc $\widehat{ATH} = {30}^{\circ}$. Sau $1,3$ phút máy bay ở độ cao $a={6}$km so với mặt đất. Hỏi vận tốc của máy bay khi bay lên?

\vspace{4pt}\noindent
\textbf{A.} $564$km/h \quad \textbf{B.} $559$km/h \quad \textbf{C.} $544$km/h \quad \textbf{\textcolor{blue!75!black}{D. (Đ)}} $554$km/h
\end{minipage}\hfill
\begin{minipage}[t]{0.39\textwidth}
\vspace{0pt}
\centering

\resizebox{\linewidth}{!}{%
\begin{tikzpicture}[declare function={r=4;rh=r*sin(30);}]
        \path (0,0) coordinate (T)
        --(30:r) coordinate (A)
        node[above,shift={(0.2,-0.05)}]{
            \resizebox{1cm}{!}{\begin{tikzpicture}[line join=round]
                    \draw[line width=2pt,line cap=butt,rounded corners=2pt] 
                    (0.3,0) arc(270:85:0.35 and 0.4)--(0.8,1.5)..controls +(50:0.2) and +(180:0.2)..(1.4,1.8)--(2.5,1.6)..controls +(0:0.35) and +(140:0.15)..(3,1.45)--(3.8,1.45)--(4.3,2.15)--(4.85,2.15)
                    --(4.3,0.95)--(3.1,0.7)..controls +(-120:0.65) and +(0:0.1)..(2.3,0)--cycle;
                    \draw[line width=2pt,line cap=butt] 
                    (4.7,1.25) arc (57:-150:0.45)
                    (0.1,-0.4) ..controls +(-30:0.15) and +(180:0.1)..
                    (0.5,-0.475)--(2.1,-0.475)..controls +(0:0.2) and +(-100:0.1)..(2.7,-0.375)
                    (-0.05,2.45)--(3.4,2.45);
                    \path[fill] (1.6,1.9)--++(0.18,0)--+(0,0.35) arc (0:180:0.09 and 0.13)--cycle;
                    \path[fill] (1.8,0.8){[rounded corners=3pt]--++(0.5,0)--++(0,0.7)--+(-0.5,0)--cycle}
                    {[rounded corners=2pt](0.8,0.8)--(1.55,0.8)--++(0,0.7)--+(-0.2,0)}
                    ..controls +(200:0.325) and +(85:0.15)..(0.8,0.8)--cycle;
                \end{tikzpicture}
            }
        }--([turn]-120:rh) coordinate (H)--cycle;
        \draw 
            (A)--(H) node[pos=0.5,right]{$a$ km}--(T)--(A)
          
            %(O) pic[draw,angle radius = 6] {right angle = B--O--C}
        ;
        \foreach \t/\g in {T/-90,H/-90,A/-30}{
            \draw[fill=white] (\t) circle (1pt) node[shift={(\g:7pt)},font=\scriptsize]{$ \t $};
        }
    \end{tikzpicture}
}
\end{minipage}

\vspace{3pt}\hrule\vspace{3pt}\footnotesize\textcolor{gray!80!black}{\textbf{Các mã dùng chung hình (HINH\_223):} \texttt{L09\_C4\_B2\_46, L09\_C4\_B2\_31}}
\end{tcolorbox}
\vspace{4pt}


\begin{tcolorbox}[
    breakable,
    colback=white,
    colframe=blue!60!black,
    boxrule=0.8pt,
    arc=3pt,
    title={\textbf{Câu \#287} \quad \textbf{MÃ CÂU: \texttt{[L09\_C4\_B2\_14]}} \quad | \quad \textbf{Nhóm: \texttt{HINH\_224}} \quad | \quad \textbf{MapID: \texttt{9H4V2-5}} \hfill \textcolor{yellow!90!white}{[Hình gốc .pyc]}}
]
\textbf{Dạng toán:} Bài toán thực tế tính khoảng cách cù lao đến bờ sông. \hfill \textit{(Trắc nghiệm -- Vận dụng)}
\vspace{4pt}\hrule\vspace{6pt}
\begin{minipage}[t]{0.58\textwidth}
\vspace{0pt}
Hai ngư dân đứng ở một bên bờ sông cách nhau $AB={200}m$ cùng nhìn thấy một cù lao trên sông với các góc nhìn lần lượt là  $\beta=49^{\circ}$ và $\alpha=28^{\circ}$. Khoảng cách ${h}$ từ cù lao ${C}$ đến bờ sông ${AB}$ (làm tròn đến hàng đơn vị) là:

\vspace{4pt}\noindent
\textbf{\textcolor{blue!75!black}{A. (Đ)}} ${73}m$ \quad \textbf{B.} ${75}m$ \quad \textbf{C.} ${71}m$ \quad \textbf{D.} ${77}m$
\end{minipage}\hfill
\begin{minipage}[t]{0.39\textwidth}
\vspace{0pt}
\centering

\resizebox{\linewidth}{!}{%
\begin{tikzpicture}[scale=0.7, font=\footnotesize, line join = round, line cap = round,>=stealth]
    \clip (-4.39,3.2) rectangle (4,-4);
    \draw[pattern=north east lines,opacity=0.3] plot[smooth] coordinates{(-4.39,1.43)(-2.43,1.46) (-1.48,1.67)(1.43,2.12)(2.01,2.09)(4,1.64) (4.1,-1.06)(3.02,-0.77)(1.85,-1.38)(0.95,-1.48)(0.13,-1.72)(-1.48,-1.67)(-2.46,-1.51)(-3.57,-0.98)(-4.37,-1.08)};
    \draw[opacity=3] plot[smooth] coordinates{(-4.39,1.43)(-2.43,1.46) (-1.48,1.67)(1.43,2.12)(2.01,2.09)(4,1.64) (4.1,-1.06)(3.02,-0.77)(1.85,-1.38)(0.95,-1.48)(0.13,-1.72)(-1.48,-1.67)(-2.46,-1.51)(-3.57,-0.98)(-4.37,-1.08)};
    \draw[fill=white] plot[smooth  cycle] coordinates{(-2.75,-0.08)(-1.83,0.32) (-0.03,0.42) (1,0) (0.24,-0.48)(-0.58,-0.79)(-1.38,-0.77)(-1.91,-0.71)};
    \draw[fill=black!70]  plot[smooth  cycle] coordinates{(-2.14,0.58)(-2.09,0.24)(-2.33,-0.13) (-1.19,-0.13) (-1.38,-0.13) (-1.69,0.21)(-1.75,0.53)};
    \draw[fill=blue!30]  plot[smooth  cycle] coordinates{(-1.75,0.53)(-1.46,0.79) (-1.01,0.58) (-1.08,1.06) (-0.64,1.08)(-0.93,1.59)(-0.53,1.85)(-0.93,2.2)(-1.38,2.7)(-2.04,3.18)(-2.49,2.91)(-3.07,2.91)(-3.2,2.22)(-3.73,1.96)(-3.1,1.3)(-3.31,1.01)(-2.83,1.01)(-2.99,0.69)(-2.14,0.58)};
    \tkzDefPoints{-1.38/-0.13/C,-2/-3/B,3/-2/A}
    \tkzDrawPoints[fill=black](A,B,C)
    \tkzDrawPolygon(A,B,C)
    \tkzDefMidPoint(A,B) 
    \tkzLabelPoints[below](A,B)
    \tkzLabelPoints[above right](C)
    \tkzMarkAngles[size=1cm,arc=l,mark=|](C,A,B)
    \tkzMarkAngles[size=1cm,arc=l,mark=||](A,B,C)
    \tkzLabelAngles[left=1cm,pos=0.5,rotate=-20](B,A,C){\footnotesize$\alpha$}
    \tkzLabelAngles[right,pos=1](A,B,C){\footnotesize$\beta$}
 \end{tikzpicture}
}
\end{minipage}

\vspace{3pt}\hrule\vspace{3pt}\footnotesize\textcolor{gray!80!black}{\textbf{Các mã dùng chung hình (HINH\_224):} \texttt{L09\_C4\_B2\_39, L09\_C4\_B2\_45, L09\_C4\_B2\_30}}
\end{tcolorbox}
\vspace{4pt}


\begin{tcolorbox}[
    breakable,
    colback=white,
    colframe=blue!60!black,
    boxrule=0.8pt,
    arc=3pt,
    title={\textbf{Câu \#288} \quad \textbf{MÃ CÂU: \texttt{[L09\_C4\_B2\_25]}} \quad | \quad \textbf{Nhóm: \texttt{HINH\_225}} \quad | \quad \textbf{MapID: \texttt{9H4V2-5}} \hfill \textcolor{yellow!90!white}{[Hình gốc .pyc]}}
]
\textbf{Dạng toán:} Bài toán thực tế hải đăng và con tàu \hfill \textit{(Trắc nghiệm -- Vận dụng)}
\vspace{4pt}\hrule\vspace{6pt}
\begin{minipage}[t]{0.58\textwidth}
\vspace{0pt}
Một người đứng trên tháp tại ${A}$ của ngọn hải đăng cao $67$ m quan sát hai lần một con tàu đang hướng về ngọn hải đăng. Lần thứ nhất người đó nhìn thấy  tàu tại ${D}$ với góc hạ $\alpha= 21^{\circ}$, lần thứ hai người đó nhìn thấy tàu tại ${C}$ với góc hạ $\beta=30^{\circ}$.  Hỏi con tàu đã đi được bao nhiêu mét giữa hai lần quan sát.( Kết quả làm tròn đến hàng đơn vị).

\vspace{4pt}\noindent
\textbf{\textcolor{blue!75!black}{A. (Đ)}} $58$ m \quad \textbf{B.} $56$ m \quad \textbf{C.} $68$ m \quad \textbf{D.} $63$ m
\end{minipage}\hfill
\begin{minipage}[t]{0.39\textwidth}
\vspace{0pt}
\centering

\resizebox{\linewidth}{!}{%
\begin{tikzpicture}
    % Sea (slightly wavy line)
    \fill[blue!30] (0,0) rectangle (10,-2);
    \draw[blue, thick] (0,-0.1) .. controls (1,-0.3) and (2,0.1) .. (3,-0.2)
    .. controls (4,-0.5) and (5,-0.1) .. (6,-0.3)
    .. controls (7,-0.1) and (8,-0.3) .. (9,-0.2)
    .. controls (10,-0.5) and (11,0.1) .. (12,-0.2);
    
    % Enhanced lighthouse (cartoonish style with more details)
    
    % Lighthouse base (stone)
    \fill[gray!50] (0.7,-0.5) rectangle (1.3,0.3);
    \draw[thick] (0.7,-0.5) -- (0.7,0.3) -- (1.3,0.3) -- (1.3,-0.5) -- cycle;
    
    % Lighthouse body (stripes for a more lively look)
    \fill[white] (0.7,0.3) rectangle (1.3,1);
    \fill[red!80] (0.7,1) rectangle (1.3,1.5);
    \fill[white] (0.7,1.5) rectangle (1.3,2);
    \fill[red!80] (0.7,2) rectangle (1.3,2.5);
    \fill[white] (0.7,2.5) rectangle (1.3,3);
    
    \draw[thick] (0.7,0.3) -- (0.7,3) -- (1.3,3) -- (1.3,0.3) -- cycle;
    
    % Windows on the lighthouse
    \fill[black!70] (0.9,1.2) rectangle (1.1,1.4); % first window
    \fill[black!70] (0.9,2.2) rectangle (1.1,2.4); % second window
    
    % Lighthouse top (lantern room)
    \fill[gray!70] (0.6,3) rectangle (1.4,3.3);
    \draw[thick] (0.6,3) -- (1.4,3);
    \draw[thick] (0.6,3.3) -- (1.4,3.3);
    
    % Lighthouse roof (dome-shaped top)
    \fill[red!80] (0.65,3.3) -- (1.35,3.3) -- (1,3.8) -- cycle;
    
    % Lighthouse light
    \fill[yellow] (1,3.15) circle (0.1);
    
    % Lighthouse light beams (enhanced beams with glow effect)
    \draw[yellow!80, thick] (1,3.15) -- (0,4);
    \draw[yellow!80, thick] (1,3.15) -- (2,4);
    
    % Lighthouse door
    \fill[brown!80] (0.9,-0.5) rectangle (1.1,-0.2);
    \draw[thick] (0.9,-0.5) -- (0.9,-0.2) -- (1.1,-0.2) -- (1.1,-0.5) -- cycle;
    
    % Boat's hull (with rounded edges for cartoon style)
    \fill[brown!70] (4,-0.5) .. controls (4.2,-1.5) and (5.8,-1.5) .. (6,-0.5) -- cycle;
    \draw[thick] (4,-0.5) .. controls (4.2,-1.5) and (5.8,-1.5) .. (6,-0.5) -- cycle;
    
    % Boat's deck
    \fill[brown!50] (4.2,-0.5) rectangle (5.8,0);
    
    % Mast of the boat
    \draw[thick] (5,0) -- (5,1.5);
    
    % Sail (cartoonish curved sail)
    \fill[white] (5,0.2) .. controls (5.8,0.7) and (5.8,1.3) .. (5,1.4) -- cycle;
    \draw[thick] (5,0.2) .. controls (5.8,0.7) and (5.8,1.3) .. (5,1.4) -- cycle;
    
    % Little flag on the top of the mast
    \fill[red] (5,1.5) -- (5.4,1.6) -- (5,1.7) -- cycle;
    
    % Add some cartoonish waves
    \foreach \x in {3.5, 5.5, 7.5} {
        \draw[blue!60,thick] (\x,-0.5) .. controls (\x+0.3,-0.6) and (\x+0.7,-0.4) .. (\x+1,-0.5);
    }
    
    % Second boat (smaller and farther away)
    \fill[brown!70] (7,-0.5) .. controls (7.1,-0.9) and (7.9,-0.9) .. (8,-0.5) -- cycle; % hull
    \draw[thick] (7,-0.5) .. controls (7.1,-0.9) and (7.9,-0.9) .. (8,-0.5) -- cycle;
    
    \draw[thick] (7.5,-0.5) -- (7.5,0.2); % mast
    \fill[white] (7.5,-0.1) .. controls (8,0) and (8,0.1) .. (7.5,0.2) -- cycle; % sail
    \draw[thick] (7.5,-0.1) .. controls (8,0) and (8,0.1) .. (7.5,0.2) -- cycle;
    
    % Flag for the second boat
    \fill[red] (7.5,0.2) -- (7.7,0.3) -- (7.5,0.4) -- cycle;
    
    % Marking points A (top of lighthouse), B (bottom of lighthouse), C (first boat), D (second boat)
    \fill[black] (1,3.5) circle (0.05) node[above] {A};  % Point A (top of lighthouse)
    \fill[black] (1,-0.5) circle (0.05) node[below] {B}; % Point B (bottom of lighthouse)
    \fill[black] (5,-0.5) circle (0.05) node[below] {C}; % Point C (position of boat 1)
    \fill[black] (7.5,-0.5) circle (0.05) node[below] {D}; % Point D (position of boat 2)
    
    % Draw lines connecting A-B, A-C, A-D
    \draw[thick] (1,3.5) -- (1,-0.5);  % Line AB
    \draw[thick] (1,3.5) -- (5,-0.5);  % Line AC
    \draw[thick] (1,3.5) -- (7.5,-0.5); % Line AD
    
    % Draw line B-C-D
    \draw[thick] (1,-0.5) -- (5,-0.5) -- (7.5,-0.5);  % Line BCD
    
    % Add a line x parallel to BC through A
    \draw[dashed] (1,3.5) -- (6,3.5) node[right] {$x$}; % Dashed line x, parallel to BC
    
    % Add angle alpha for xAD
    \draw[thick] (1.6,2.9) arc[start angle=-60, end angle= -10, radius=0.8];
    \node at (1.6,3.2) {$\beta$};  % Label for angle alpha
    
    % Add angle beta for xAC
    \draw[thick] (1.9,2.9) arc[start angle=-25, end angle= -13, radius=3.1];
    \node at (2.3,3.3) {$\alpha$};  % Label for angle beta
    
    % Add a smiling sun in the sky
    \fill[yellow] (9,3.5) circle (0.5); % sun body
    \draw (9,3.7) arc[start angle=45, end angle=135, radius=0.2]; % smiling face
    \fill[black] (8.8,3.6) circle (0.05); % left eye
    \fill[black] (9.2,3.6) circle (0.05); % right eye
    \foreach \angle in {0,45,...,315} {
        \draw[thick, yellow] (9,3.5) -- ++(\angle:0.8);
    }
    
\end{tikzpicture}
}
\end{minipage}

\vspace{3pt}\hrule\vspace{3pt}\footnotesize\textcolor{gray!80!black}{\textbf{Các mã dùng chung hình (HINH\_225):} \texttt{L09\_C4\_B2\_42, L09\_C4\_B2\_24}}
\end{tcolorbox}
\vspace{4pt}


\begin{tcolorbox}[
    breakable,
    colback=white,
    colframe=blue!60!black,
    boxrule=0.8pt,
    arc=3pt,
    title={\textbf{Câu \#289} \quad \textbf{MÃ CÂU: \texttt{[L09\_C4\_B2\_10]}} \quad | \quad \textbf{Nhóm: \texttt{HINH\_226}} \quad | \quad \textbf{MapID: \texttt{9H4V2-5}} \hfill \textcolor{yellow!90!white}{[Hình gốc .pyc]}}
]
\textbf{Dạng toán:} Bài toán thực tế tính chiều cao toà nhà \hfill \textit{(Trắc nghiệm -- Vận dụng)}
\vspace{4pt}\hrule\vspace{6pt}
\begin{minipage}[t]{0.58\textwidth}
\vspace{0pt}
Khi tia sáng mặt trời tạo với mặt đất một góc ${67}^{\circ}$ thì bóng của một tòa nhà trên mặt đất có độ dài $b={34}$m. Chiều cao của tòa nhà (kết quả làm tròn đến mét) là

\vspace{4pt}\noindent
\textbf{A.} $84$m \quad \textbf{B.} $90$m \quad \textbf{C.} $76$m \quad \textbf{\textcolor{blue!75!black}{D. (Đ)}} ${80}m$
\end{minipage}\hfill
\begin{minipage}[t]{0.39\textwidth}
\vspace{0pt}
\centering

\resizebox{\linewidth}{!}{%
\begin{tikzpicture}[scale=1,font=\footnotesize,line cap=round,line join=round]
    % Tòa nhà ----------------------------------------------
    \path
    (3,0) coordinate (A)
    (3,5) coordinate (B)
    (8,0) coordinate (C)
    ($(A)!.5!(C)$) node[below] {$34$ m}
    ($(C)+(162:1)$) node {$67^{\circ}$}
    ;
    \foreach \i in {0,0.5,1,1.5,2,2.5}{\foreach \j in {0,0.8,1.6,2.4,3.2,4.0}
        \draw[fill=yellow!50!white] (\i,\j) rectangle +(0.5,0.8);
    } 
    % Ban công 
    \foreach \i in {0.1,0.6,1.1,1.6,2.1,2.6}{\foreach \j in {0,0.8,1.6,2.4,3.2,4.0}
        \draw[fill=orange!60!white] (\i,\j) rectangle +(0.3,0.6);
    }
    \foreach \j in {0,0.8,1.6,2.4,3.2,4.0,4.8}
        \draw[fill=cyan!30!white] (0,\j) rectangle +(3,0.3);
    
    \draw (A)--(B)--(C)--cycle;    %       \tkzMarkRightAngles(B,A,C)
    %       \tkzMarkAngles[size=.7](B,C,A)
    
    \foreach \i/\p in {A/-90,B/90,C/0}  
        \fill[black] (\i) circle (1.5pt)  % Điểm đánh dấu màu đen
        node[shift={(\p:3mm)},black] {$\i$}; % Chữ cái màu đen
\end{tikzpicture}
}
\end{minipage}

\end{tcolorbox}
\vspace{4pt}


\begin{tcolorbox}[
    breakable,
    colback=white,
    colframe=blue!60!black,
    boxrule=0.8pt,
    arc=3pt,
    title={\textbf{Câu \#290} \quad \textbf{MÃ CÂU: \texttt{[L09\_C4\_B2\_11]}} \quad | \quad \textbf{Nhóm: \texttt{HINH\_221}} \quad | \quad \textbf{MapID: \texttt{9H4V2-5}} \hfill \textcolor{yellow!90!white}{[Hình gốc .pyc]}}
]
\textbf{Dạng toán:} Bài toán thực tế tính chiều cao của cái cây \hfill \textit{(Trắc nghiệm -- Vận dụng)}
\vspace{4pt}\hrule\vspace{6pt}
\begin{minipage}[t]{0.58\textwidth}
\vspace{0pt}
Khi tia sáng mặt trời tạo với mặt đất một góc $\alpha={58}^{\circ}$ thì bóng của một cái cây trên mặt đất có độ dài $x={6}$m. Chiều cao của cái cây (kết quả làm tròn đến hàng đơn vị) là

\vspace{4pt}\noindent
\textbf{A.} $11$m \quad \textbf{\textcolor{blue!75!black}{B. (Đ)}} ${10}m$ \quad \textbf{C.} $9$m \quad \textbf{D.} $12$m
\end{minipage}\hfill
\begin{minipage}[t]{0.39\textwidth}
\vspace{0pt}
\centering

\resizebox{\linewidth}{!}{%
\begin{tikzpicture}[scale=1, font=\footnotesize, line cap=round, line join=round]

    % Vẽ ông mặt trời
    \fill[yellow] (-3,7) circle (0.5cm); % Di chuyển mặt trời ra xa một chút

    % Vẽ các tia nắng sinh động bao quanh ông mặt trời (từng tia riêng lẻ)
    \pgfmathsetmacro{\radius}{0.8} % Khoảng cách từ tâm mặt trời đến điểm bắt đầu của tia nắng
    
    % Vẽ từng tia nắng riêng lẻ
    \pgfmathsetmacro{\lengthA}{0.5 + 0.4 * sin(0 * 8.8)}
    \draw[orange, thick] (-3,7) ++(0:\radius) -- ++(0:\lengthA);

    \pgfmathsetmacro{\lengthB}{0.5 + 0.4 * sin(15 * 8.8)}
    \draw[orange, thick] (-3,7) ++(15:\radius) -- ++(15:\lengthB);

    \pgfmathsetmacro{\lengthC}{0.5 + 0.4 * sin(30 * 8.8)}
    \draw[orange, thick] (-3,7) ++(30:\radius) -- ++(30:\lengthC);

    \pgfmathsetmacro{\lengthD}{0.5 + 0.4 * sin(45 * 8.8)}
    \draw[orange, thick] (-3,7) ++(45:\radius) -- ++(45:\lengthD);

    \pgfmathsetmacro{\lengthE}{0.5 + 0.4 * sin(60 * 8.8)}
    \draw[orange, thick] (-3,7) ++(60:\radius) -- ++(60:\lengthE);

    \pgfmathsetmacro{\lengthF}{0.5 + 0.4 * sin(75 * 8.8)}
    \draw[orange, thick] (-3,7) ++(75:\radius) -- ++(75:\lengthF);

    \pgfmathsetmacro{\lengthG}{0.5 + 0.4 * sin(90 * 8.8)}
    \draw[orange, thick] (-3,7) ++(90:\radius) -- ++(90:\lengthG);

    \pgfmathsetmacro{\lengthH}{0.5 + 0.4 * sin(105 * 8.8)}
    \draw[orange, thick] (-3,7) ++(105:\radius) -- ++(105:\lengthH);

    \pgfmathsetmacro{\lengthI}{0.5 + 0.4 * sin(120 * 8.8)}
    \draw[orange, thick] (-3,7) ++(120:\radius) -- ++(120:\lengthI);

    \pgfmathsetmacro{\lengthJ}{0.5 + 0.4 * sin(135 * 8.8)}
    \draw[orange, thick] (-3,7) ++(135:\radius) -- ++(135:\lengthJ);

    \pgfmathsetmacro{\lengthK}{0.5 + 0.4 * sin(150 * 8.8)}
    \draw[orange, thick] (-3,7) ++(150:\radius) -- ++(150:\lengthK);

    \pgfmathsetmacro{\lengthL}{0.5 + 0.4 * sin(165 * 8.8)}
    \draw[orange, thick] (-3,7) ++(165:\radius) -- ++(165:\lengthL);

    \pgfmathsetmacro{\lengthM}{0.5 + 0.4 * sin(180 * 8.8)}
    \draw[orange, thick] (-3,7) ++(180:\radius) -- ++(180:\lengthM);

    \pgfmathsetmacro{\lengthN}{0.5 + 0.4 * sin(195 * 8.8)}
    \draw[orange, thick] (-3,7) ++(195:\radius) -- ++(195:\lengthN);

    \pgfmathsetmacro{\lengthO}{0.5 + 0.4 * sin(210 * 8.8)}
    \draw[orange, thick] (-3,7) ++(210:\radius) -- ++(210:\lengthO);

    \pgfmathsetmacro{\lengthP}{0.5 + 0.4 * sin(225 * 8.8)}
    \draw[orange, thick] (-3,7) ++(225:\radius) -- ++(225:\lengthP);

    \pgfmathsetmacro{\lengthQ}{0.5 + 0.4 * sin(240 * 8.8)}
    \draw[orange, thick] (-3,7) ++(240:\radius) -- ++(240:\lengthQ);

    \pgfmathsetmacro{\lengthR}{0.5 + 0.4 * sin(255 * 8.8)}
    \draw[orange, thick] (-3,7) ++(255:\radius) -- ++(255:\lengthR);

    \pgfmathsetmacro{\lengthS}{0.5 + 0.4 * sin(270 * 8.8)}
    \draw[orange, thick] (-3,7) ++(270:\radius) -- ++(270:\lengthS);

    \pgfmathsetmacro{\lengthT}{0.5 + 0.4 * sin(285 * 8.8)}
    \draw[orange, thick] (-3,7) ++(285:\radius) -- ++(285:\lengthT);

    \pgfmathsetmacro{\lengthU}{0.5 + 0.4 * sin(300 * 8.8)}
    \draw[orange, thick] (-3,7) ++(300:\radius) -- ++(300:\lengthU);

    \pgfmathsetmacro{\lengthV}{0.5 + 0.4 * sin(315 * 8.8)}
    \draw[orange, thick] (-3,7) ++(315:\radius) -- ++(315:\lengthV);

    \pgfmathsetmacro{\lengthW}{0.5 + 0.4 * sin(330 * 8.8)}
    \draw[orange, thick] (-3,7) ++(330:\radius) -- ++(330:\lengthW);

    \pgfmathsetmacro{\lengthX}{0.5 + 0.4 * sin(345 * 8.8)}
    \draw[orange, thick] (-3,7) ++(345:\radius) -- ++(345:\lengthX);

    % Vẽ cái cây cao hơn
    \fill[brown] (0,0) rectangle (0.5,4); % Tăng chiều cao của cái cây
    
    % Vẽ tán cây bằng các hình tròn chồng lên nhau, thẳng hàng với thân cây và xếp cao dần
    \fill[green] (0.25,3.5) circle (0.8cm); % Vòng tròn lớn dưới cùng
    \fill[green] (0.25,4.5) circle (1.2cm); % Vòng tròn trung bình ở giữa
    \fill[green] (0.25,4) circle (1cm); % Vòng tròn nhỏ ở trên cùng
    \fill[brown] (0.2,0) ellipse (0.8 and 0.24); % Elip với kích thước tuỳ chỉnh

    % Định nghĩa các điểm A, B, C
    \path
    (0.25,0) coordinate (A) % Gốc cây
    (0.25,5.6) coordinate (B) % Ngọn cây
    (4,0) coordinate (C); % Điểm trên mặt đất, đã di chuyển ra xa
    
    % Vẽ tam giác ABC
    \draw (A) -- (B) -- (C) -- cycle; % Nối các điểm A, B, C

    % Gắn nhãn các đoạn thẳng và góc
    \draw ($(A)!.5!(C)$) node[below] {$x$}; % Nhãn đoạn thẳng AC
    \node at ($(C)+(155:1)$) {$\alpha$}; % Ký hiệu góc BCA

    % Đánh dấu các điểm A, B, C
    \foreach \i/\p in {A/-90,B/90,C/0}  
    \fill[black] (\i) circle (1.5pt)  
    node[shift={(\p:3mm)},black] {$\i$}; % Chữ cái màu đen

\end{tikzpicture}
}
\end{minipage}

\vspace{3pt}\hrule\vspace{3pt}\footnotesize\textcolor{gray!80!black}{\textbf{Các mã dùng chung hình (HINH\_221):} \texttt{L09\_C4\_B1\_19, L09\_C4\_B2\_48, L09\_C4\_B2\_33}}
\end{tcolorbox}
\vspace{4pt}


\begin{tcolorbox}[
    breakable,
    colback=white,
    colframe=blue!60!black,
    boxrule=0.8pt,
    arc=3pt,
    title={\textbf{Câu \#291} \quad \textbf{MÃ CÂU: \texttt{[L09\_C4\_B2\_09]}} \quad | \quad \textbf{Nhóm: \texttt{HINH\_227}} \quad | \quad \textbf{MapID: \texttt{9H4C2-5}} \hfill \textcolor{yellow!90!white}{[Hình gốc .pyc]}}
]
\textbf{Dạng toán:} Bài toán thực tế tính thời gian của máy bay khi hạ cánh \hfill \textit{(Trắc nghiệm -- Vận dụng cao)}
\vspace{4pt}\hrule\vspace{6pt}
\begin{minipage}[t]{0.58\textwidth}
\vspace{0pt}
Quãng đường bay khi hạ cánh của máy bay, hay còn gọi là quãng đường tiếp cận cuối cùng (landing distance), là khoảng cách mà máy bay bay từ lúc bắt đầu tiếp cận cuối cùng cho đến khi chạm đất. Đây là giai đoạn máy bay đã hạ thấp gần mặt đất và đang hướng thẳng tới đường băng với góc hạ cánh cố định (thường là $\widehat{ATH}={3}^{\circ}$), chuẩn bị tiếp xúc với đường băng để hoàn tất quá trình hạ cánh. Một máy bay ở giai đoạn tiếp cận cuối cùng cách mặt đất $a={450}$m, tính thời gian hạ cánh ở giai đoạn này biết vận tốc hạ cánh là $250$km/h.

\vspace{4pt}\noindent
\textbf{A.} $2,01$ phút \quad \textbf{\textcolor{blue!75!black}{B. (Đ)}} $2,06$ phút \quad \textbf{C.} $2,26$ phút \quad \textbf{D.} $2,16$ phút
\end{minipage}\hfill
\begin{minipage}[t]{0.39\textwidth}
\vspace{0pt}
\centering

\resizebox{\linewidth}{!}{%
\begin{tikzpicture}[declare function={r=4;rh=r*sin(30);}]
        \path (0,0) coordinate (T)
        --(30:r) coordinate (A)
        node[above,shift={(0.2,-0.05)}]{
            \resizebox{1cm}{!}{\begin{tikzpicture}[line join=round]
                    \draw[line width=2pt,line cap=butt,rounded corners=2pt] 
                    (0.3,0) arc(270:85:0.35 and 0.4)--(0.8,1.5)..controls +(50:0.2) and +(180:0.2)..(1.4,1.8)--(2.5,1.6)..controls +(0:0.35) and +(140:0.15)..(3,1.45)--(3.8,1.45)--(4.3,2.15)--(4.85,2.15)
                    --(4.3,0.95)--(3.1,0.7)..controls +(-120:0.65) and +(0:0.1)..(2.3,0)--cycle;
                    \draw[line width=2pt,line cap=butt] 
                    (4.7,1.25) arc (57:-150:0.45)
                    (0.1,-0.4) ..controls +(-30:0.15) and +(180:0.1)..
                    (0.5,-0.475)--(2.1,-0.475)..controls +(0:0.2) and +(-100:0.1)..(2.7,-0.375)
                    (-0.05,2.45)--(3.4,2.45);
                    \path[fill] (1.6,1.9)--++(0.18,0)--+(0,0.35) arc (0:180:0.09 and 0.13)--cycle;
                    \path[fill] (1.8,0.8){[rounded corners=3pt]--++(0.5,0)--++(0,0.7)--+(-0.5,0)--cycle}
                    {[rounded corners=2pt](0.8,0.8)--(1.55,0.8)--++(0,0.7)--+(-0.2,0)}
                    ..controls +(200:0.325) and +(85:0.15)..(0.8,0.8)--cycle;
                \end{tikzpicture}
            }
        }--([turn]-120:rh) coordinate (H)--cycle;
        %\draw[-]  (-r/2,0) --(3/2*r,0)node[below];
        \draw 
            (A)--(H) node[pos=0.5,right]{$450$ m}--(T)--(A)
        ;
        \foreach \t/\g in {T/-90,H/-90,A/-30}{
            \draw[fill=white] (\t) circle (1pt) node[shift={(\g:7pt)},font=\scriptsize]{$ \t $};
        }
    \end{tikzpicture}
}
\end{minipage}

\end{tcolorbox}
\vspace{4pt}


\begin{tcolorbox}[
    breakable,
    colback=white,
    colframe=blue!60!black,
    boxrule=0.8pt,
    arc=3pt,
    title={\textbf{Câu \#292} \quad \textbf{MÃ CÂU: \texttt{[L09\_C4\_B2\_12]}} \quad | \quad \textbf{Nhóm: \texttt{HINH\_228}} \quad | \quad \textbf{MapID: \texttt{9H4C2-5}} \hfill \textcolor{yellow!90!white}{[Hình gốc .pyc]}}
]
\textbf{Dạng toán:} Bài toán thực tế tính chiều cao của hải đăng \hfill \textit{(Trắc nghiệm -- Vận dụng cao)}
\vspace{4pt}\hrule\vspace{6pt}
\begin{minipage}[t]{0.58\textwidth}
\vspace{0pt}
Ở hai địa điểm ${A}$ và ${B}$ trên mặt đất cách nhau $22$m, người ta đặt ống nhòm nhìn thấy vị trí ${D}$ cao nhất của ngọn hải đăng, góc đặt ống nhóm tại ${B}$ hợp với mặt đất là $40^{\circ}$, góc đặt ống nhòm tại vị trí ${A}$ hợp với mặt đất là $60^{\circ}$. Chiều cao của hải đăng (làm tròn đến hàng đơn vị) là:

\vspace{4pt}\noindent
\textbf{A.} $44$m \quad \textbf{B.} $40$m \quad \textbf{C.} $46$m \quad \textbf{\textcolor{blue!75!black}{D. (Đ)}} $36$m
\end{minipage}\hfill
\begin{minipage}[t]{0.39\textwidth}
\vspace{0pt}
\centering

\resizebox{\linewidth}{!}{%
\begin{tikzpicture}[scale=.5,declare function={h=13.4;a=h/tan(63);b=h/tan(48);}]
        \fill[brown] (0,0) -- (-5,0) 
        to[out=20,in=180] (-4.3,.15)
        to[out=40,in=180] (-3.05,.65)
        to[out=80,in=200] (-2.8,.85)
        to[out=70,in=160] (-2.3,1.25)
        to[out=50,in=190] (-2,1.45)
        --(-1.6,1.45)--(-1.35,3.2)--(-1.3,2.15)
        to[out=50,in=210] (0,3.2)
        to[out=28,in=222] (1.25,4.05)
        --(1.5,1.5)--(1.7,1.5)
        to[out=-40,in=180] (2,1.33)
        to[out=-10,in=110] (2.45,1)
        to[out=-40,in=160] (3,.8)
        to[out=-40,in=150] (3.2,.6)
        to[out=10,in=130] (4.4,.1)
        to[out=0,in=160] (5,0)
        --cycle;
        \fill[brown] (-.85,3.6) rectangle (-.53,4.1);
        \fill[brown] (-1.3,4.2)
        to[out=50,in=208] (0,5.15)
        to[out=25,in=225] (1,5.75)
        to[out=50,in=275] (1.15,6.3)
        --(1.08,8)
        to[out=-130,in=37] (-1,6.6)
        to[out=-135,in=88] (-1.25,5.5)
        --cycle;
        \fill[brown] (.25,8.45) rectangle (.55,9);
        \fill[brown] (-1.1,8.3)
        to[out=55,in=245] (1.05,9.85)--(1.2,9.85)
        to[out=10,in=263] (1.45,10.1)
        to[out=100,in=-10] (.96,10.35)--(.93,11.1)--(1.05,11.15)
        to[out=155,in=20] (-1.07,11.2)--(-1.07,10.3)
        to[out=-175,in=135] (-1.4,10)
        to[out=-55,in=190] (-1.08,9.85)--cycle;
        \draw[brown,line width= 2pt] (.9,11.2)--(.9,11.7)
        (.25,11.2)--(.25,11.9)
        (-.35,11.2)--(-.35,11.9)
        (-1.03,11.2)--(-1.03,11.7)
        to[out=20,in=160] (0.94,11.7)
        (.55,11.8)--(.54,12.7)
        (-.6,11.8)--(-.59,12.7)
        (-.04,11.9)--(-.04,12.7);
        \draw[brown,line width=4pt] (.6,11.2)--(.6,11.8)
        (-.75,11.2)--(-.75,11.8);
        \fill[brown] (-.04,12.15) circle (7pt)
        (.6,12.7) arc (0:180:.65 and .7) -- cycle
        (-.04,13.4) circle (5pt);
        \path 
        (0,0) coordinate (C)
        (a,0) coordinate (A)
        (b,0) coordinate (B)
        (0,h) coordinate (D);
        \foreach \i/\g in {C/-90, A/-90, B/-90, D/90} \fill (\i)circle(.1) node[shift={(\g:.25)}] {$\i$};
        \draw (C)--(B)node[pos=.8, below]{$22m$}--(D)--cycle (A)--(D)
        (A)+(-1,0) arc (180:117:1)node[pos=.7, left]{$\alpha=60^\circ$};
        \draw[double] (B)+(-1,0) arc (180:132:1)node[pos=.7, left]{$\beta=40^\circ$};
        \draw[dashed] (-5,h)--(D);
        \draw[<->,-stealth] (-4,0) -- (-4,h)node[pos=.5, left]{$h$};
    \end{tikzpicture}
}
\end{minipage}

\vspace{3pt}\hrule\vspace{3pt}\footnotesize\textcolor{gray!80!black}{\textbf{Các mã dùng chung hình (HINH\_228):} \texttt{L09\_C4\_B2\_38, L09\_C4\_B2\_52, L09\_C4\_B2\_35}}
\end{tcolorbox}
\vspace{4pt}


\begin{tcolorbox}[
    breakable,
    colback=white,
    colframe=blue!60!black,
    boxrule=0.8pt,
    arc=3pt,
    title={\textbf{Câu \#293} \quad \textbf{MÃ CÂU: \texttt{[L09\_C4\_B2\_13]}} \quad | \quad \textbf{Nhóm: \texttt{HINH\_229}} \quad | \quad \textbf{MapID: \texttt{9H4C2-5}} \hfill \textcolor{yellow!90!white}{[Hình gốc .pyc]}}
]
\textbf{Dạng toán:} Bài toán thực tế tính chiều cao của con dốc \hfill \textit{(Trắc nghiệm -- Vận dụng cao)}
\vspace{4pt}\hrule\vspace{6pt}
\begin{minipage}[t]{0.58\textwidth}
\vspace{0pt}
Lúc 6h sáng, bạn Anh đi xe đạp từ nhà (điểm ${A}$) đến trường (điểm ${B}$) phải leo và xuống một con dốc (như hình vẽ sau) biết ${AB}=1,0km$, $\widehat{A}=11^{\circ}$, $\widehat{B}=6^{\circ}$. Tính chiều cao ${h}$ của con dốc (theo mét làm tròn đến chữ số hàng đơn vị)

\vspace{4pt}\noindent
\textbf{\textcolor{blue!75!black}{A. (Đ)}} ${68}m$ \quad \textbf{B.} ${72}m$ \quad \textbf{C.} ${66}m$ \quad \textbf{D.} ${70}m$
\end{minipage}\hfill
\begin{minipage}[t]{0.39\textwidth}
\vspace{0pt}
\centering

\resizebox{\linewidth}{!}{%
\begin{tikzpicture}
    
    % Vẽ tam giác ABC
    \coordinate [label=left:A] (A) at (0,0);
    \coordinate [label=right:B] (B) at (8,0);
    \coordinate [label=above:C] (C) at (2,1);
    
    % Đặt các góc và chiều cao
    \draw (A) -- (B) -- (C) -- cycle;
    
    % Vẽ chiều cao từ C tới AB
\draw (C) -- (2,0) coordinate (H);
\node[below] at (H) {H};
      % Vẽ dấu vuông tại điểm H
        \draw (2,0) rectangle +(0.1,0.1);
    % Nhãn chiều cao
    \draw[<-] (C) -- ($(C)!0.5!(H)$) node[midway, right] {$h$};
    
    % Vẽ các cung góc
    
    \node at ($(A)+(0.6,0.15)$) {$11^{\circ}$};
    

    \node at ($(B)+(-1.9,0.15)$) {$6^{\circ}$};
    
\end{tikzpicture}
}
\end{minipage}

\vspace{3pt}\hrule\vspace{3pt}\footnotesize\textcolor{gray!80!black}{\textbf{Các mã dùng chung hình (HINH\_229):} \texttt{L09\_C4\_B2\_37, L09\_C4\_B2\_50, L09\_C4\_B2\_36}}
\end{tcolorbox}
\vspace{4pt}


\begin{tcolorbox}[
    breakable,
    colback=white,
    colframe=blue!60!black,
    boxrule=0.8pt,
    arc=3pt,
    title={\textbf{Câu \#294} \quad \textbf{MÃ CÂU: \texttt{[L09\_C4\_B2\_15]}} \quad | \quad \textbf{Nhóm: \texttt{HINH\_230}} \quad | \quad \textbf{MapID: \texttt{9H4C2-5}} \hfill \textcolor{yellow!90!white}{[Hình gốc .pyc]}}
]
\textbf{Dạng toán:} Bài toán thực tế khó \hfill \textit{(Trắc nghiệm -- Vận dụng cao)}
\vspace{4pt}\hrule\vspace{6pt}
\begin{minipage}[t]{0.58\textwidth}
\vspace{0pt}
Một người đang ở trên sân thượng ${A}$ của một cái toà nhà nhìn xuống thấy một con đường chạy thẳng đến chân tường. Anh ta nhìn thấy một chiếc xe máy với góc hạ $\widehat{xAC}=30^{\circ}$. Sau $8$ phút lại nhìn thấy xe đó với góc hạ $\widehat{xAD}=65^{\circ}$. Hỏi sau bao nhiêu phút thì xe máy chạy đến chân tường biết vận tốc của xe máy không đổi (làm tròn kết quả đến hàng phần mười).

\vspace{4pt}\noindent
\textbf{A.} $3,1$ phút \quad \textbf{\textcolor{blue!75!black}{B. (Đ)}} $2,9$ phút \quad \textbf{C.} $2,7$ phút \quad \textbf{D.} $3,0$ phút
\end{minipage}\hfill
\begin{minipage}[t]{0.39\textwidth}
\vspace{0pt}
\centering

\resizebox{\linewidth}{!}{%
\begin{tikzpicture}
    \def\hh{6cm}
    \def\aa{5cm}
    \path (0,0) coordinate (B) (0,\hh) coordinate (A) (\aa,0) coordinate (C)
    (3*\aa/5,0) coordinate (D)      (3*\aa/5,\hh) coordinate (x);
        \draw[->] (0,\hh) -- (3*\aa/5,\hh) node[above] {$x$};
    \draw 
(A) pic[draw,angle radius = 6mm] {angle =D--A--x} 
(A) pic[draw,angle radius = 5mm] {angle =C--A--x}   ;
        % Hiển thị nhãn cho các điểm
    \draw (A) node [above] {$A$};
    \draw (B) node [below] {$B$};
    \draw (C) node [below] {$C$};
    \draw (D) node [below right] {$D$};
    % Vẽ tia Ax song song với BC

    
    
    % Vẽ hình tứ giác với các điểm A, B, C, D
    \draw (A) -- (B) -- (C) -- (D) -- cycle;
    
    % Vẽ đường nối A và C bằng nét liền
    \draw (A) -- (C);
    
    % Vẽ tường mô phòng
    \path[fill=red!30] (B) -- ++(180:5) 
    to[out=100,in=-90] (-4.2,1.5) 
    to[out=100,in=-90] (-3.7,3.5) 
    to[out=100,in=-90] (-2.8,4.8) 
    to[out=120,in=-110] (-2,6) -- (A) -- cycle;
    
    % Cắt vùng mô phỏng bức tường
    \clip (B) -- ++(180:5) 
    to[out=100,in=-90] (-4.2,1.5) 
    to[out=100,in=-90] (-3.7,3.5) 
    to[out=100,in=-90] (-2.8,4.8) 
    to[out=120,in=-110] (-2,6) -- (A) -- cycle;
    
    % Tô màu các viên gạch hình chữ nhật
    \fill[pattern=bricks, pattern color=red] (B) -- ++(180:5) 
    to[out=100,in=-90] (-4.2,1.5) 
    to[out=100,in=-90] (-3.7,3.5) 
    to[out=100,in=-90] (-2.8,4.8) 
    to[out=120,in=-110] (-2,6) -- (A) -- cycle;
    
    % Vẽ lại hình tứ giác và đường nối
    \draw (A)--(B)--(C)--(D)--cycle;
    \draw[dashed] (B)--(A);
    \draw (B)--++(180:5) to[out=100,in=-90](-4.2,1.5) to[out=100,in=-90](-3.7,3.5) to[out=100,in=-90](-2.8,4.8) to[out=120,in=-110](-2,6)--(A)--cycle;
    

\end{tikzpicture}
}
\end{minipage}

\vspace{3pt}\hrule\vspace{3pt}\footnotesize\textcolor{gray!80!black}{\textbf{Các mã dùng chung hình (HINH\_230):} \texttt{L09\_C4\_B2\_51}}
\end{tcolorbox}
\vspace{4pt}


\begin{tcolorbox}[
    breakable,
    colback=white,
    colframe=blue!60!black,
    boxrule=0.8pt,
    arc=3pt,
    title={\textbf{Câu \#295} \quad \textbf{MÃ CÂU: \texttt{[L09\_C4\_B2\_16]}} \quad | \quad \textbf{Nhóm: \texttt{HINH\_231}} \quad | \quad \textbf{MapID: \texttt{9H4C2-5}} \hfill \textcolor{yellow!90!white}{[Hình gốc .pyc]}}
]
\textbf{Dạng toán:} Bài toán thực tế tính chiều cao của lâu đài. \hfill \textit{(Trắc nghiệm -- Vận dụng cao)}
\vspace{4pt}\hrule\vspace{6pt}
\begin{minipage}[t]{0.58\textwidth}
\vspace{0pt}
Để đo chiều cao của một toà lâu đài người ta đặt giác kế thẳng đứng tại vị trí ${D}$. Quay ống ngắm của giác kế sao cho nhìn thấy đỉnh ${A}$ của toà lâu đài dưới góc nhọn $\alpha=28^{\circ}$.  Sau đó đặt giác kế thẳng đứng tại điểm ${C}$, biết $CD=20m$ thì nhìn thấy đỉnh ${A}$ dưới góc $\beta={21^{\circ}}$. Biết chiều cao của giác kế là ${1,4}m$. Hãy tính chiều cao của lâu đài.

\vspace{4pt}\noindent
\textbf{A.} $30,21m$ \quad \textbf{B.} $31,01m$ \quad \textbf{C.} $27,81m$ \quad \textbf{\textcolor{blue!75!black}{D. (Đ)}} ${29,01}m$
\end{minipage}\hfill
\begin{minipage}[t]{0.39\textwidth}
\vspace{0pt}
\centering

\resizebox{\linewidth}{!}{%
\begin{tikzpicture}[scale=1]
    
    % Định nghĩa các điểm
    \path
    (1,5) coordinate (A)
    (1,0) coordinate (B)
    (8,0) coordinate (C)
    (6,1) coordinate (E)
    (8,1) coordinate (F)
    (6,0) coordinate (D)
    (1,1) coordinate (H);
    
    % Vẽ tháp
    \fill[pink!50] (0,0) rectangle (2,4);
    \fill[red!30] (0,4) -- (1,5) -- (2,4) -- cycle;
    
    % Cửa sổ
    \fill[blue!30] (0.5,1) rectangle (1.5,2.5);
    
    % Điểm A trên đỉnh tháp
    \node[circle, fill=red, inner sep=2pt] at (1,5) {};
    \node[above] at (1,5) {A};
    
    % Điểm B dưới chân tháp
    \node[circle, fill=blue, inner sep=2pt] at (1,0) {};
    \node[below] at (1,0) {B};
    
    % Điểm C dưới chân tháp
    \node[circle, fill=blue, inner sep=2pt] at (8,0) {};
    \node[below] at (8,0) {C};
    
    % Điểm D dưới chân tháp
    \node[circle, fill=blue, inner sep=2pt] at (6,0) {};
    \node[below] at (6,0) {D};
    
    % Điểm F dưới chân tháp
    \node[circle, fill=blue, inner sep=2pt] at (8,1) {};
    \node[below left] at (8,1) {F};
    
    % Điểm E dưới chân tháp
    \node[circle, fill=blue, inner sep=2pt] at (6,1) {};
    \node[below left] at (6,1) {E};
    
    % Điểm H dưới chân tháp
    \node[circle, fill=blue, inner sep=2pt] at (1,1) {};
    \node[below left] at (1,1) {H};
    
    % Vẽ đường viền của tháp
    \draw[black, thick] (0,0) rectangle (2,4);
    \draw[black, thick] (0,4) -- (1,5) -- (2,4) -- cycle;
    \draw[black, thick] (0.5,1) rectangle (1.5,2.5);
    
    % Nối các điểm bằng nét liền
        \draw[black, thick] (A) -- (F) -- (C) --(B)-- cycle;
    \draw[black, thick] (A) -- (E) -- (D) -- (C) -- (B) -- cycle;
    \draw[black, thick] (A) -- (E) -- (F) -- (H) -- cycle;
        \path (E) pic[draw=black,fill=gray!25, angle radius = 7] {angle =A--E--H}
    node[shift={(-0.5,2.9pt)},scale=0.8](GA){$\alpha $};
        \path (F) pic[draw=black,fill=gray!25, angle radius = 7] {angle =A--F--H}
    node[shift={(-0.6,2.9pt)},scale=0.8](GA){$\beta $};
\end{tikzpicture}
}
\end{minipage}

\end{tcolorbox}
\vspace{4pt}


\begin{tcolorbox}[
    breakable,
    colback=white,
    colframe=blue!60!black,
    boxrule=0.8pt,
    arc=3pt,
    title={\textbf{Câu \#296} \quad \textbf{MÃ CÂU: \texttt{[L09\_C4\_B2\_27]}} \quad | \quad \textbf{Nhóm: \texttt{HINH\_232}} \quad | \quad \textbf{MapID: \texttt{9H4V2-5}} \hfill \textcolor{yellow!90!white}{[Hình gốc .pyc]}}
]
\textbf{Dạng toán:} Bài toán đo khoảng cách khinh khí cầu đến mặt đất \hfill \textit{(Trắc nghiệm -- Vận dụng)}
\vspace{4pt}\hrule\vspace{6pt}
\begin{minipage}[t]{0.58\textwidth}
\vspace{0pt}
Một cậu bé cao ${1,4}m$ đứng ngắm khinh khí cầu ở một vị trí cách hình chiếu của nó trên một đất tức điểm ${E}$ một khoảng là $330$m. Cậu bé đặt ống nhòm hợp với phương ngang một góc $57^{\circ}$ thì nhìn thấy khinh khí cầu. Hỏi khi đó khinh khí cầu cách mặt đất bao nhiêu mét?(Kết quả làm tròn đến chữ số thập phân thứ nhất)

\vspace{4pt}\noindent
\textbf{A.} $519,6$ m \quad \textbf{B.} $499,6$ m \quad \textbf{\textcolor{blue!75!black}{C. (Đ)}} $509,6$ m \quad \textbf{D.} $529,6$ m
\end{minipage}\hfill
\begin{minipage}[t]{0.39\textwidth}
\vspace{0pt}
\centering

\resizebox{\linewidth}{!}{%
\begin{tikzpicture}[line join = round, line cap=round,>=stealth,font=\footnotesize,scale=1]
\path
(0,0) coordinate (D)
($(D)+(0,1)$) coordinate (B)
($(D)+(4,0)$) coordinate (E)
($(E)+(0,1)$) coordinate (A)
($(E)+(0,4)$) coordinate (C)
($(C)+(-2,0)$) coordinate (x)
;
\draw (D)--(E)--(C)--(B)--cycle (B)--(A);
\foreach \p/\r in {A/0,B/90,C/0,D/-90,E/-90}
\fill (\p) circle (1.2pt) node[shift={(\r:3mm)}]{$\p$};
\draw[pattern=north west lines] ($(D)+(-1,0)$) rectangle ($(E)+(1,-0.5)$);
\path
(D)--(E) node[above,pos=0.5]{$330$ m} ;
\draw[<->] ($(D)+(-0.75,0)$)--($(B)+(-0.75,0)$) node[left,pos=0.5]{$1,4$ m};
\draw pic[draw,angle radius=3mm] {right angle = B--A--C};
\draw pic[draw,angle radius=8mm] {angle = A--B--C};
\draw ($(B)+(1.25,0)$) node[above]{$57^\circ$};
\draw[dashed] (C)--($(x)+(-2,0)$)node[above]{$x$};
\begin{scope}[line width=3pt]
\foreach \yA/\xB/\yB in {0.6/-0.3/0.05, 0.6/-0.15/0.05} {
\draw (-0.25, 0.45) -- (-0.25, \yA) -- (\xB, \yB);
}
\foreach \yA/\xB/\yB/\xC/\yC in {
0.9/-0.1/0.8/0.0/0.9,
0.9/-0.4/0.8/-0.3/0.7
} {
\draw (-0.25, 0.45) -- (-0.25, \yA) -- (\xB, \yB) -- (\xC, \yC);
}
\draw (-0.25, 0.45) -- (-0.25, 1.0)circle(2pt);
\end{scope}
\fill [scale=.1,black,yshift=40cm,xshift=40cm]
(-1,0) rectangle (1,1)
(-1,2).. controls +(135:1) and +(180:4) .. (0,7)
.. controls +(180:2) and +(135:2).. (-.5,2)
.. controls +(120:2) and +(180:1) .. (0,6.99)
.. controls +(0:1) and +(60:2) .. (.5,2)
.. controls +(45:2) and +(0:2) .. (0,7)
.. controls +(0:4) and +(45:1).. (1,2) ;
\end{tikzpicture}
}
\end{minipage}

\end{tcolorbox}
\vspace{4pt}


\begin{tcolorbox}[
    breakable,
    colback=white,
    colframe=blue!60!black,
    boxrule=0.8pt,
    arc=3pt,
    title={\textbf{Câu \#297} \quad \textbf{MÃ CÂU: \texttt{[L09\_C4\_B2\_37]}} \quad | \quad \textbf{Nhóm: \texttt{HINH\_229}} \quad | \quad \textbf{MapID: \texttt{9H4C2-5}} \hfill \textcolor{yellow!90!white}{[Hình gốc .pyc]}}
]
\textbf{Dạng toán:} TF-M4. Bài toán thực tế tính chiều cao của con dốc \hfill \textit{(Đúng-Sai -- Vận dụng cao)}
\vspace{4pt}\hrule\vspace{6pt}
\begin{minipage}[t]{0.58\textwidth}
\vspace{0pt}
Lúc 6h sáng, bạn Anh đi xe đạp từ nhà (điểm ${A}$) đến trường (điểm ${B}$) phải leo và xuống một con dốc (như hình vẽ sau) biết ${AB}=1,6km$, $\widehat{A}=10^{\circ}$, $\widehat{B}=5^{\circ}$. Gọi chiều cao của con dốc là ${h}$. Xét tính đúng-sai của các khẳng định sau.

\vspace{4pt}\noindent
\textbf{a)} $AH=h \cdot \tan{5^{\circ}}$\\[2pt]
\textbf{\textcolor{blue!75!black}{b) (Đ)}} $BH=h \cdot \cot{10^{\circ}}$\\[2pt]
\textbf{\textcolor{blue!75!black}{c) (Đ)}} $AB = h \left(\cot{5^{\circ}} +\cot{10^{\circ}} \right)$\\[2pt]
\textbf{d)} $h\approx 94m$
\end{minipage}\hfill
\begin{minipage}[t]{0.39\textwidth}
\vspace{0pt}
\centering

\resizebox{\linewidth}{!}{%
\begin{tikzpicture}
    
    % Vẽ tam giác ABC
    \coordinate [label=left:A] (A) at (0,0);
    \coordinate [label=right:B] (B) at (8,0);
    \coordinate [label=above:C] (C) at (2,1);
    
    % Đặt các góc và chiều cao
    \draw (A) -- (B) -- (C) -- cycle;
    
    % Vẽ chiều cao từ C tới AB
\draw (C) -- (2,0) coordinate (H);
\node[below] at (H) {H};
      % Vẽ dấu vuông tại điểm H
        \draw (2,0) rectangle +(0.1,0.1);
    % Nhãn chiều cao
    \draw[<-] (C) -- ($(C)!0.5!(H)$) node[midway, right] {$h$};
    
    % Vẽ các cung góc
    
    \node at ($(A)+(0.6,0.15)$) {$10^{\circ}$};
    

    \node at ($(B)+(-1.9,0.15)$) {$5^{\circ}$};
    
\end{tikzpicture}
}
\end{minipage}

\vspace{3pt}\hrule\vspace{3pt}\footnotesize\textcolor{gray!80!black}{\textbf{Các mã dùng chung hình (HINH\_229):} \texttt{L09\_C4\_B2\_13, L09\_C4\_B2\_50, L09\_C4\_B2\_36}}
\end{tcolorbox}
\vspace{4pt}


\begin{tcolorbox}[
    breakable,
    colback=white,
    colframe=blue!60!black,
    boxrule=0.8pt,
    arc=3pt,
    title={\textbf{Câu \#298} \quad \textbf{MÃ CÂU: \texttt{[L09\_C4\_B2\_38]}} \quad | \quad \textbf{Nhóm: \texttt{HINH\_228}} \quad | \quad \textbf{MapID: \texttt{9H4C2-5}} \hfill \textcolor{yellow!90!white}{[Hình gốc .pyc]}}
]
\textbf{Dạng toán:} TF-M4. Bài toán thực tế tính chiều cao của hải đăng \hfill \textit{(Đúng-Sai -- Vận dụng cao)}
\vspace{4pt}\hrule\vspace{6pt}
\begin{minipage}[t]{0.58\textwidth}
\vspace{0pt}
Ở hai địa điểm ${A}$ và ${B}$ trên mặt đất cách nhau $22$m, người ta đặt ống nhòm nhìn thấy vị trí ${D}$ cao nhất của ngọn hải đăng, góc đặt ống nhóm tại ${B}$ hợp với mặt đất là $49^{\circ}$, góc đặt ống nhòm tại vị trí ${A}$ hợp với mặt đất là $67^{\circ}$. Gọi chiều cao của hải đăng là ${h}$. Xét tính đúng-sai của các khẳng định sau.

\vspace{4pt}\noindent
\textbf{a)} $AC=h\cdot \tan{67^{\circ}}$\\[2pt]
\textbf{\textcolor{blue!75!black}{b) (Đ)}} $BC=h\cdot \cot{49^{\circ}}$\\[2pt]
\textbf{\textcolor{blue!75!black}{c) (Đ)}} $h= \dfrac{22}{\cot{49^{\circ}}-\cot{67^{\circ}}}$\\[2pt]
\textbf{\textcolor{blue!75!black}{d) (Đ)}} $h\approx 49$m
\end{minipage}\hfill
\begin{minipage}[t]{0.39\textwidth}
\vspace{0pt}
\centering

\resizebox{\linewidth}{!}{%
\begin{tikzpicture}[scale=.5,declare function={h=13.4;a=h/tan(63);b=h/tan(48);}]
        \fill[brown] (0,0) -- (-5,0) 
        to[out=20,in=180] (-4.3,.15)
        to[out=40,in=180] (-3.05,.65)
        to[out=80,in=200] (-2.8,.85)
        to[out=70,in=160] (-2.3,1.25)
        to[out=50,in=190] (-2,1.45)
        --(-1.6,1.45)--(-1.35,3.2)--(-1.3,2.15)
        to[out=50,in=210] (0,3.2)
        to[out=28,in=222] (1.25,4.05)
        --(1.5,1.5)--(1.7,1.5)
        to[out=-40,in=180] (2,1.33)
        to[out=-10,in=110] (2.45,1)
        to[out=-40,in=160] (3,.8)
        to[out=-40,in=150] (3.2,.6)
        to[out=10,in=130] (4.4,.1)
        to[out=0,in=160] (5,0)
        --cycle;
        \fill[brown] (-.85,3.6) rectangle (-.53,4.1);
        \fill[brown] (-1.3,4.2)
        to[out=50,in=208] (0,5.15)
        to[out=25,in=225] (1,5.75)
        to[out=50,in=275] (1.15,6.3)
        --(1.08,8)
        to[out=-130,in=37] (-1,6.6)
        to[out=-135,in=88] (-1.25,5.5)
        --cycle;
        \fill[brown] (.25,8.45) rectangle (.55,9);
        \fill[brown] (-1.1,8.3)
        to[out=55,in=245] (1.05,9.85)--(1.2,9.85)
        to[out=10,in=263] (1.45,10.1)
        to[out=100,in=-10] (.96,10.35)--(.93,11.1)--(1.05,11.15)
        to[out=155,in=20] (-1.07,11.2)--(-1.07,10.3)
        to[out=-175,in=135] (-1.4,10)
        to[out=-55,in=190] (-1.08,9.85)--cycle;
        \draw[brown,line width= 2pt] (.9,11.2)--(.9,11.7)
        (.25,11.2)--(.25,11.9)
        (-.35,11.2)--(-.35,11.9)
        (-1.03,11.2)--(-1.03,11.7)
        to[out=20,in=160] (0.94,11.7)
        (.55,11.8)--(.54,12.7)
        (-.6,11.8)--(-.59,12.7)
        (-.04,11.9)--(-.04,12.7);
        \draw[brown,line width=4pt] (.6,11.2)--(.6,11.8)
        (-.75,11.2)--(-.75,11.8);
        \fill[brown] (-.04,12.15) circle (7pt)
        (.6,12.7) arc (0:180:.65 and .7) -- cycle
        (-.04,13.4) circle (5pt);
        \path 
        (0,0) coordinate (C)
        (a,0) coordinate (A)
        (b,0) coordinate (B)
        (0,h) coordinate (D);
        \foreach \i/\g in {C/-90, A/-90, B/-90, D/90} \fill (\i)circle(.1) node[shift={(\g:.25)}] {$\i$};
        \draw (C)--(B)node[pos=.8, below]{$22m$}--(D)--cycle (A)--(D)
        (A)+(-1,0) arc (180:117:1)node[pos=.7, left]{$\alpha=67^\circ$};
        \draw[double] (B)+(-1,0) arc (180:132:1)node[pos=.7, left]{$\beta=49^\circ$};
        \draw[dashed] (-5,h)--(D);
        \draw[<->,-stealth] (-4,0) -- (-4,h)node[pos=.5, left]{$h$};
    \end{tikzpicture}
}
\end{minipage}

\vspace{3pt}\hrule\vspace{3pt}\footnotesize\textcolor{gray!80!black}{\textbf{Các mã dùng chung hình (HINH\_228):} \texttt{L09\_C4\_B2\_12, L09\_C4\_B2\_52, L09\_C4\_B2\_35}}
\end{tcolorbox}
\vspace{4pt}


\begin{tcolorbox}[
    breakable,
    colback=white,
    colframe=blue!60!black,
    boxrule=0.8pt,
    arc=3pt,
    title={\textbf{Câu \#299} \quad \textbf{MÃ CÂU: \texttt{[L09\_C4\_B2\_39]}} \quad | \quad \textbf{Nhóm: \texttt{HINH\_224}} \quad | \quad \textbf{MapID: \texttt{9H4C2-5}} \hfill \textcolor{yellow!90!white}{[Hình gốc .pyc]}}
]
\textbf{Dạng toán:} TF-M4. Bài toán thực tế tính khoảng cách cù lao đến bờ sông. \hfill \textit{(Đúng-Sai -- Vận dụng cao)}
\vspace{4pt}\hrule\vspace{6pt}
\begin{minipage}[t]{0.58\textwidth}
\vspace{0pt}
Hai ngư dân đứng ở một bên bờ sông cách nhau $AB={200}m$ cùng nhìn thấy một cù lao trên sông với các góc nhìn lần lượt là $\beta=51^{\circ}$ và $\alpha=35^{\circ}$. Gọi khoảng cách từ cù lao ${C}$ đến bờ sông ${AB}$ là ${h}$.

\vspace{4pt}\noindent
\textbf{a)} $BH=h \cdot \tan{51^{\circ}}$\\[2pt]
\textbf{\textcolor{blue!75!black}{b) (Đ)}} $AH=h \cdot \cot{35^{\circ}}$\\[2pt]
\textbf{\textcolor{blue!75!black}{c) (Đ)}} $AB= h \left( \cot{51^{\circ}} +  \cot{35^{\circ}} \right)$\\[2pt]
\textbf{d)} $h\approx 99$m
\end{minipage}\hfill
\begin{minipage}[t]{0.39\textwidth}
\vspace{0pt}
\centering

\resizebox{\linewidth}{!}{%
\begin{tikzpicture}[scale=0.7, font=\footnotesize, line join = round, line cap = round,>=stealth]
    \clip (-4.39,3.2) rectangle (4,-4);
    \draw[pattern=north east lines,opacity=0.3] plot[smooth] coordinates{(-4.39,1.43)(-2.43,1.46) (-1.48,1.67)(1.43,2.12)(2.01,2.09)(4,1.64) (4.1,-1.06)(3.02,-0.77)(1.85,-1.38)(0.95,-1.48)(0.13,-1.72)(-1.48,-1.67)(-2.46,-1.51)(-3.57,-0.98)(-4.37,-1.08)};
    \draw[opacity=3] plot[smooth] coordinates{(-4.39,1.43)(-2.43,1.46) (-1.48,1.67)(1.43,2.12)(2.01,2.09)(4,1.64) (4.1,-1.06)(3.02,-0.77)(1.85,-1.38)(0.95,-1.48)(0.13,-1.72)(-1.48,-1.67)(-2.46,-1.51)(-3.57,-0.98)(-4.37,-1.08)};
    \draw[fill=white] plot[smooth  cycle] coordinates{(-2.75,-0.08)(-1.83,0.32) (-0.03,0.42) (1,0) (0.24,-0.48)(-0.58,-0.79)(-1.38,-0.77)(-1.91,-0.71)};
    \draw[fill=black!70]  plot[smooth  cycle] coordinates{(-2.14,0.58)(-2.09,0.24)(-2.33,-0.13) (-1.19,-0.13) (-1.38,-0.13) (-1.69,0.21)(-1.75,0.53)};
    \draw[fill=blue!30]  plot[smooth  cycle] coordinates{(-1.75,0.53)(-1.46,0.79) (-1.01,0.58) (-1.08,1.06) (-0.64,1.08)(-0.93,1.59)(-0.53,1.85)(-0.93,2.2)(-1.38,2.7)(-2.04,3.18)(-2.49,2.91)(-3.07,2.91)(-3.2,2.22)(-3.73,1.96)(-3.1,1.3)(-3.31,1.01)(-2.83,1.01)(-2.99,0.69)(-2.14,0.58)};
    \tkzDefPoints{-1.38/-0.13/C,-2/-3/B,3/-2/A}
    \tkzDrawPoints[fill=black](A,B,C)
    \tkzDrawPolygon(A,B,C)
    \tkzDefMidPoint(A,B) 
    \tkzLabelPoints[below](A,B)
    \tkzLabelPoints[above right](C)
    \tkzMarkAngles[size=1cm,arc=l,mark=|](C,A,B)
    \tkzMarkAngles[size=1cm,arc=l,mark=||](A,B,C)
    \tkzLabelAngles[left=1cm,pos=0.5,rotate=-20](B,A,C){\footnotesize$\alpha$}
    \tkzLabelAngles[right,pos=1](A,B,C){\footnotesize$\beta$}
 \end{tikzpicture}
}
\end{minipage}

\vspace{3pt}\hrule\vspace{3pt}\footnotesize\textcolor{gray!80!black}{\textbf{Các mã dùng chung hình (HINH\_224):} \texttt{L09\_C4\_B2\_14, L09\_C4\_B2\_45, L09\_C4\_B2\_30}}
\end{tcolorbox}
\vspace{4pt}


\begin{tcolorbox}[
    breakable,
    colback=teal!3!white,
    colframe=teal!75!black,
    boxrule=0.8pt,
    arc=3pt,
    title={\textbf{Câu \#300} \quad \textbf{MÃ CÂU: \texttt{[L09\_C4\_B2\_57]}} \quad | \quad \textbf{Nhóm: \texttt{HINH\_233}} \quad | \quad \textbf{MapID: \texttt{9H4V2-5}} \hfill \textbf{\textcolor{green!40!black}{[ĐÃ THAY MẪU TIKZ MỚI]}}}
]
\textbf{Dạng toán:} TF-M3. Bài toán đo chiều cao của nhà 2 khi đứng ở nhà 1 \hfill \textit{(Đúng-Sai -- Vận dụng)}
\vspace{4pt}\hrule\vspace{6pt}
\begin{minipage}[t]{0.58\textwidth}
\vspace{0pt}
Cho hai tòa nhà 1 và tòa nhà 2 như hình vẽ bên. Trên nóc tòa nhà 2 có một cột ăng-ten thẳng cao $1$ m. Từ vị trí quan sát ${A}$ (trên nóc tòa nhà 1) cao $8$ m so với mặt đất có thể nhìn thấy đỉnh ${B}$ và chân ${C}$ của cột ăng-ten lần lượt dưới góc $ \widehat{BAD}= 54^{\circ}$ và $ \widehat{CAD}= 48^{\circ}$ so với phương nằm ngang.

\vspace{4pt}\noindent
\textbf{\textcolor{blue!75!black}{a) (Đ)}} $DC=AD\cdot \tan 48^{\circ}$\\[2pt]
\textbf{b)} $AD=\dfrac{1}{\sin 54^{\circ}-\sin 48^{\circ}} $\\[2pt]
\textbf{c)} $DC=\dfrac{1\tan 48^{\circ}}{\sin 54^{\circ}-\sin 48^{\circ} } $\\[2pt]
\textbf{d)} Chiều cao của tòa nhà 2 khoảng $14,2$ m. (Làm tròn kết quả đến hàng phần mười)
\end{minipage}\hfill
\begin{minipage}[t]{0.39\textwidth}
\vspace{0pt}
\centering
\definecolor{skyb}{RGB}{102, 178, 240}
\definecolor{buildblue}{RGB}{30, 115, 230}
\definecolor{gridblue}{RGB}{185, 225, 255}
\definecolor{roofblue}{RGB}{15, 75, 170}
\definecolor{grassgreen}{RGB}{125, 185, 45}
\resizebox{\linewidth}{!}{%
\begin{tikzpicture}[
    >=stealth, 
    line join=round, 
    line cap=round, 
    every node/.style={font=\normalsize},
    declare function={
        w1 = 2.4;
        h1 = 3.6;
        d_space = 4.2;
        w2 = 2.6;
        ang_C = 48;
        ang_B = 54;
        dist_AD = d_space + 0.5*w2; 
        yC = h1 + dist_AD * tan(ang_C);
        yB = h1 + dist_AD * tan(ang_B);
        h_ant = yB - yC;
        r_pt = 0.4pt;
    },scale=0.55
]
    \draw[line width=0.8pt] (-0.8,0) -- (w1 + d_space + w2 + 0.8, 0);
    \draw[fill=grassgreen!85, draw=grassgreen!60!black, rounded corners=5pt, line width=0.6pt] 
        (-0.3,-0.05) rectangle (w1 + 0.3, 0.25);
    \draw[fill=grassgreen!85, draw=grassgreen!60!black, rounded corners=5pt, line width=0.6pt] 
        (w1 + d_space - 0.3, -0.05) rectangle (w1 + d_space + w2 + 0.3, 0.25);
    \begin{scope}
        \draw[fill=buildblue!70, draw=white, line width=0.8pt] 
            (0.1, 0.2) rectangle ($(w1-0.1, 0.9)$);
        \draw[fill=buildblue, draw=white, line width=1pt] 
            (0.2, 0.9) rectangle ($(w1-0.2, {h1-0.4})$);
        \foreach \i in {1,...,10} {
            \draw[gridblue!60, line width=0.3pt] 
                (0.2, {0.9 + \i*(h1-1.3)/11}) -- ($(w1-0.2, {0.9 + \i*(h1-1.3)/11})$);
        }
        \foreach \j in {1,...,4} {
            \draw[gridblue!60, line width=0.3pt] 
                ({0.2 + \j*(w1-0.4)/5}, 0.9) -- ({0.2 + \j*(w1-0.4)/5}, {h1-0.4});
        }
        \draw[fill=roofblue, draw=white, line width=0.8pt] 
            (0.4, {h1-0.4}) rectangle ($(w1-0.4, {h1-0.15})$);
        \draw[fill=roofblue!90, draw=white, line width=0.8pt] 
            (0.6, {h1-0.15}) rectangle ($(w1-0.6, h1)$);
        \draw[line width=1pt, roofblue!50!black] (0.3, h1) -- ($(w1-0.3, h1)$);
    \end{scope}
    \begin{scope}[shift={(w1 + d_space, 0)}]
        \draw[fill=white, draw=gray!40, line width=0.6pt] 
            (0.1, 0.2) rectangle ($(w2-0.1, 1.3)$);
        \foreach \k in {0,1,2} {
            \draw[fill=buildblue!90, draw=white, line width=0.6pt] 
                ({0.3 + \k*0.75}, 0.3) rectangle ({0.85 + \k*0.75}, 1.15);
        }
        \draw[fill=buildblue, draw=white, line width=1.2pt] 
            (0.2, 1.3) rectangle ($(w2-0.2, {yC-2.2})$);
        \foreach \i in {1,...,6} {
            \draw[gridblue!60, line width=0.3pt] 
                (0.2, {1.3 + \i*(yC-3.5)/7}) -- ($(w2-0.2, {1.3 + \i*(yC-3.5)/7})$);
        }
        \foreach \j in {1,...,5} {
            \draw[gridblue!60, line width=0.3pt] 
                ({0.2 + \j*(w2-0.4)/6}, 1.3) -- ({0.2 + \j*(w2-0.4)/6}, {yC-2.2});
        }
        \draw[fill=buildblue, draw=white, line width=1.2pt] 
            (0.45, {yC-2.2}) rectangle ($(w2-0.45, {yC-0.5})$);
        \foreach \i in {1,...,4} {
            \draw[gridblue!60, line width=0.3pt] 
                (0.45, {yC-2.2 + \i*1.7/5}) -- ($(w2-0.45, {yC-2.2 + \i*1.7/5})$);
        }
        \foreach \j in {1,...,3} {
            \draw[gridblue!60, line width=0.3pt] 
                ({0.45 + \j*(w2-0.9)/4}, {yC-2.2}) -- ({0.45 + \j*(w2-0.9)/4}, {yC-0.5});
        }
        \draw[fill=buildblue!90, draw=white, line width=1pt] 
            (0.75, {yC-0.5}) rectangle ($(w2-0.75, yC)$);
        \draw[line width=1.2pt, roofblue] (0.15, {yC-2.2}) -- ($(w2-0.15, {yC-2.2})$);
        \draw[line width=1.2pt, roofblue] (0.4, {yC-0.5}) -- ($(w2-0.4, {yC-0.5})$);
        \draw[line width=1.2pt, roofblue] (0.7, yC) -- ($(w2-0.7, yC)$);
        \draw[line width=1.3pt, roofblue!90!black] (0.5*w2, yC) -- (0.5*w2, yB);
        \foreach \h in {0.25, 0.5, 0.75} {
            \draw[line width=0.7pt, roofblue!90!black] 
                ({0.5*w2 - 0.35*(1-\h*0.5)}, {yC + \h*h_ant}) -- ({0.5*w2 + 0.35*(1-\h*0.5)}, {yC + \h*h_ant});
        }
        \fill[roofblue!90!black] (0.5*w2, yB) circle (1.5pt);
    \end{scope}
    \coordinate (A) at (w1, h1);
    \coordinate (D) at (w1 + d_space + 0.5*w2, h1);
    \coordinate (C) at (w1 + d_space + 0.5*w2, yC);
    \coordinate (B) at (w1 + d_space + 0.5*w2, yB);
    \coordinate (H1) at (w1, 0);
    \coordinate (H2) at (w1 + d_space + 0.5*w2, 0);
    \draw[dashed, line width=0.6pt, black!70] (A) -- (D);
    \draw[dashed, line width=0.6pt, black!70] (B) -- (H2);
    \draw[line width=0.7pt, black!85] (A) -- (C);
    \draw[line width=0.7pt, black!85] (A) -- (B);
    \draw (D) ++(-0.25, 0) -- ++(0, 0.25) -- ++(0.25, 0);
    \draw[line width=0.5pt] (A) ++(0:1.25) arc (0:ang_C:1.25);
    \node at ($(A) + ({0.5*ang_C}:1.75)$) {$ 48^\circ $};
    \draw[line width=0.5pt] (A) ++(0:2.35) arc (0:ang_B:2.35);
    \draw[->, >=stealth, line width=0.4pt, black!80] 
        ($(A) + ({ang_B + 9}:2.85)$) to[bend left=20] ($(A) + ({0.5*(ang_C + ang_B)}:2.25)$);
    \node[above left, inner sep=1pt] at ($(A) + ({ang_B + 9}:2.85)$) {$ 54^\circ $};
    \draw[<->, >=stealth, line width=0.6pt, black!80] (0.65, 0) -- (0.65, h1) 
        node[midway, fill=buildblue!90, text=white, inner sep=1.5pt, rounded corners=1pt] {$ 8\,\text{m} $};
    \draw[line width=0.4pt, gray!80] (C) -- ++(0.65, 0);
    \draw[line width=0.4pt, gray!80] (B) -- ++(0.65, 0);
    \draw[<->, >=stealth, line width=0.6pt, black!80] ($(C)+(0.5,0)$) -- ($(B)+(0.5,0)$)
        node[midway, right, inner sep=2pt] {$ 1\,\text{m} $};
    \foreach \pt in {A, B, C, D} {
        \fill[black] (\pt) circle (1.2pt);
    }
    \node[above left, inner sep=2pt] at (A) {$ A $};
    \node[above left, inner sep=2pt] at (B) {$ B $};
    \node[above left, inner sep=2pt] at (C) {$ C $};
    \node[below right, inner sep=2pt, fill=buildblue!85, text=white, rounded corners=1pt] at ($(D)+(0.05,-0.05)$) {$ D $};
    \node[below, font=\bfseries\small] at (0.5*w1, -0.1) {Tòa nhà 1};
    \node[below, font=\bfseries\small] at (w1 + d_space + 0.5*w2, -0.1) {Tòa nhà 2};
\end{tikzpicture}
}
\end{minipage}

\end{tcolorbox}
\vspace{4pt}


\begin{tcolorbox}[
    breakable,
    colback=white,
    colframe=blue!60!black,
    boxrule=0.8pt,
    arc=3pt,
    title={\textbf{Câu \#301} \quad \textbf{MÃ CÂU: \texttt{[L09\_C4\_B2\_40]}} \quad | \quad \textbf{Nhóm: \texttt{HINH\_234}} \quad | \quad \textbf{MapID: \texttt{9H4V2-5}} \hfill \textcolor{yellow!90!white}{[Hình gốc .pyc]}}
]
\textbf{Dạng toán:} SA-M3. Toán thực tế tính chiều cao của bức tường. \hfill \textit{(Trả lời ngắn -- Vận dụng)}
\vspace{4pt}\hrule\vspace{6pt}
\begin{minipage}[t]{0.58\textwidth}
\vspace{0pt}
Một cậu bé cao $a={1,0}m$ đứng cách một bức tường cao $b={1,5}m$ sao cho $\widehat {BAC}$ vuông tại ${A}$. Hỏi toà nhà cao bao nhiêu mét? (Làm tròn đến chữ số thập phân thứ nhất)


\textbf{Đáp án điền số:} 	exttt{3,2}
\end{minipage}\hfill
\begin{minipage}[t]{0.39\textwidth}
\vspace{0pt}
\centering

\resizebox{\linewidth}{!}{%
\begin{tikzpicture}[line join=round, line cap=round, >=stealth, font=\footnotesize, scale=2]
    \tikzset{
        man/.pic={
            \begin{scope}[shift={(-0.3,0)}]
                \fill [rounded corners=1.5] (0,0.45) -- (0,0.9) -- (0.4,0.9) -- (0.4,0.5) --
                (0.325,0.5) -- (0.325,0.8) -- (0.3,0.8) -- (0.3,0) -- (0.225,0) --
                (0.225,0.5) -- (0.175,0.5) -- (0.175,0) -- (0.1,0) -- (0.1,0.8) --
                (0.075,0.8) -- (0.075,0.5) -- cycle;
                \fill (0.2,1.05) circle (0.12);
                \coordinate (-head) at (0.2,1);
                \coordinate (-foot) at (0.2,0);
            \end{scope}
        }
    }
    % \clip (1.7,-.4) rectangle (6.1,5.45);
    \def\c{1.1} \def\h{1.6}
    \pgfmathsetmacro{\b}{((\h)^2)/(\c)+\c}
    \path
    (0,0) coordinate (O)
    (\h,0) coordinate (C)
    (0,\c) coordinate (A)
    (\h,\b)coordinate (B)
    (\h,\c)coordinate (H);
    \pic[scale=2,pink] at (O) {man};
    \draw[pattern=bricks,pattern color=brown] (C) rectangle ($(B)+(.5,0)$);
    \draw[thick] 
    (A)--(C)--(B)--cycle
    (A)--(H) node[midway,above]{$b\ \mathrm{m}$}
    (O)--(A) node[midway,above=.5,sloped]{$a\ \mathrm{m}$};
    \draw
    pic[draw,angle radius=3mm]{right angle=C--A--B}
    pic[draw,angle radius=3mm]{right angle=A--H--C};
    \foreach \d/\g in {A/100,B/75,C/-75,H/135,O/-75} {\fill (\d) circle(1pt) node[shift=(\g:10pt)]{$\d$};}
    \end{tikzpicture}
}
\end{minipage}

\vspace{3pt}\hrule\vspace{3pt}\footnotesize\textcolor{gray!80!black}{\textbf{Các mã dùng chung hình (HINH\_234):} \texttt{L09\_C4\_B2\_22}}
\end{tcolorbox}
\vspace{4pt}


\begin{tcolorbox}[
    breakable,
    colback=white,
    colframe=blue!60!black,
    boxrule=0.8pt,
    arc=3pt,
    title={\textbf{Câu \#302} \quad \textbf{MÃ CÂU: \texttt{[L09\_C4\_B2\_41]}} \quad | \quad \textbf{Nhóm: \texttt{HINH\_235}} \quad | \quad \textbf{MapID: \texttt{9H4C2-5}} \hfill \textcolor{yellow!90!white}{[Hình gốc .pyc]}}
]
\textbf{Dạng toán:} SA-M4. Toán thực tế tính chiều cao của lâu đài. \hfill \textit{(Trả lời ngắn -- Vận dụng cao)}
\vspace{4pt}\hrule\vspace{6pt}
\begin{minipage}[t]{0.58\textwidth}
\vspace{0pt}
Để đo chiều cao của một toà lâu đài người ta đặt giác kế thẳng đứng tại vị trí ${D}$. Quay ống ngắm của giác kế sao cho nhìn thấy đỉnh ${A}$ của toà lâu đài dưới góc nhọn $\alpha=28^{\circ}$.  Sau đó đặt giác kế thẳng đứng tại điểm ${C}$, biết $CD=20m$ thì nhìn thấy đỉnh ${A}$ dưới góc $\beta={17^{\circ}}$. Biết chiều cao của giác kế là ${1,5}m$. Hãy tính chiều cao của lâu đài. (Chỉ làm tròn kết quả cuối cùng đến hàng phần mười)


\textbf{Đáp án điền số:} 	exttt{15,9}
\end{minipage}\hfill
\begin{minipage}[t]{0.39\textwidth}
\vspace{0pt}
\centering

\resizebox{\linewidth}{!}{%
\begin{tikzpicture}[scale=1]
    
    % Định nghĩa các điểm
    \path
    (1,5) coordinate (A)
    (1,0) coordinate (B)
    (8,0) coordinate (C)
    (6,1) coordinate (E)
    (8,1) coordinate (F)
    (6,0) coordinate (D)
    (1,1) coordinate (H);
    
    % Vẽ tháp
    \fill[pink!50] (0,0) rectangle (2,4);
    \fill[red!30] (0,4) -- (1,5) -- (2,4) -- cycle;
    
    % Cửa sổ
    \fill[blue!30] (0.5,1) rectangle (1.5,2.5);
    
    % Điểm A trên đỉnh tháp
    \node[circle, fill=red, inner sep=2pt] at (1,5) {};
    \node[above] at (1,5) {A};
    
    % Điểm B dưới chân tháp
    \node[circle, fill=blue, inner sep=2pt] at (1,0) {};
    \node[below] at (1,0) {B};
    
    % Điểm C dưới chân tháp
    \node[circle, fill=blue, inner sep=2pt] at (8,0) {};
    \node[below] at (8,0) {C};
    
    % Điểm D dưới chân tháp
    \node[circle, fill=blue, inner sep=2pt] at (6,0) {};
    \node[below] at (6,0) {D};
    
    % Điểm F dưới chân tháp
    \node[circle, fill=blue, inner sep=2pt] at (8,1) {};
    \node[below left] at (8,1) {F};
    
    % Điểm E dưới chân tháp
    \node[circle, fill=blue, inner sep=2pt] at (6,1) {};
    \node[below left] at (6,1) {E};
    
    % Điểm H dưới chân tháp
    \node[circle, fill=blue, inner sep=2pt] at (1,1) {};
    \node[below left] at (1,1) {H};
    
    % Vẽ đường viền của tháp
    \draw[black, thick] (0,0) rectangle (2,4);
    \draw[black, thick] (0,4) -- (1,5) -- (2,4) -- cycle;
    \draw[black, thick] (0.5,1) rectangle (1.5,2.5);
    
    % Nối các điểm bằng nét liền
    \draw[black, thick] (A) -- (F) -- (C) -- (B) -- cycle;
    \draw[black, thick] (A) -- (E) -- (D) -- (C) -- (B) -- cycle;
    \draw[black, thick] (A) -- (E) -- (F) -- (H) -- cycle;
    \path (E) pic[draw=black,fill=gray!25, angle radius = 7] {angle =A--E--H}
    node[shift={(-0.5,2.9pt)},scale=0.8](GA){$\alpha$};
    \path (F) pic[draw=black,fill=gray!25, angle radius = 7] {angle =A--F--H}
    node[shift={(-0.6,2.9pt)},scale=0.8](GA){$\beta$};
    \end{tikzpicture}
}
\end{minipage}

\vspace{3pt}\hrule\vspace{3pt}\footnotesize\textcolor{gray!80!black}{\textbf{Các mã dùng chung hình (HINH\_235):} \texttt{L09\_C4\_B2\_23}}
\end{tcolorbox}
\vspace{4pt}


\begin{tcolorbox}[
    breakable,
    colback=white,
    colframe=blue!60!black,
    boxrule=0.8pt,
    arc=3pt,
    title={\textbf{Câu \#303} \quad \textbf{MÃ CÂU: \texttt{[L09\_C4\_B2\_42]}} \quad | \quad \textbf{Nhóm: \texttt{HINH\_225}} \quad | \quad \textbf{MapID: \texttt{9H4V2-5}} \hfill \textcolor{yellow!90!white}{[Hình gốc .pyc]}}
]
\textbf{Dạng toán:} SA-M3. Toán thực tế hải đăng và con tàu \hfill \textit{(Trả lời ngắn -- Vận dụng)}
\vspace{4pt}\hrule\vspace{6pt}
\begin{minipage}[t]{0.58\textwidth}
\vspace{0pt}
Một người đứng trên tháp tại ${A}$ của ngọn hải đăng cao $85$ m quan sát hai lần một con tàu đang hướng về ngọn hải đăng. Lần thứ nhất người đó nhìn thấy  tàu tại ${D}$ với góc hạ $\alpha= 24^{\circ}$, lần thứ hai người đó nhìn thấy tàu tại ${C}$ với góc hạ $\beta=35^{\circ}$.  Hỏi con tàu đã đi được bao nhiêu mét giữa hai lần quan sát.( Kết quả làm tròn đến hàng đơn vị).


\textbf{Đáp án điền số:} 	exttt{70}
\end{minipage}\hfill
\begin{minipage}[t]{0.39\textwidth}
\vspace{0pt}
\centering

\resizebox{\linewidth}{!}{%
\begin{tikzpicture}
    % Sea (slightly wavy line)
    \fill[blue!30] (0,0) rectangle (10,-2);
    \draw[blue, thick] (0,-0.1) .. controls (1,-0.3) and (2,0.1) .. (3,-0.2)
    .. controls (4,-0.5) and (5,-0.1) .. (6,-0.3)
    .. controls (7,-0.1) and (8,-0.3) .. (9,-0.2)
    .. controls (10,-0.5) and (11,0.1) .. (12,-0.2);
    
    % Enhanced lighthouse (cartoonish style with more details)
    
    % Lighthouse base (stone)
    \fill[gray!50] (0.7,-0.5) rectangle (1.3,0.3);
    \draw[thick] (0.7,-0.5) -- (0.7,0.3) -- (1.3,0.3) -- (1.3,-0.5) -- cycle;
    
    % Lighthouse body (stripes for a more lively look)
    \fill[white] (0.7,0.3) rectangle (1.3,1);
    \fill[red!80] (0.7,1) rectangle (1.3,1.5);
    \fill[white] (0.7,1.5) rectangle (1.3,2);
    \fill[red!80] (0.7,2) rectangle (1.3,2.5);
    \fill[white] (0.7,2.5) rectangle (1.3,3);
    
    \draw[thick] (0.7,0.3) -- (0.7,3) -- (1.3,3) -- (1.3,0.3) -- cycle;
    
    % Windows on the lighthouse
    \fill[black!70] (0.9,1.2) rectangle (1.1,1.4); % first window
    \fill[black!70] (0.9,2.2) rectangle (1.1,2.4); % second window
    
    % Lighthouse top (lantern room)
    \fill[gray!70] (0.6,3) rectangle (1.4,3.3);
    \draw[thick] (0.6,3) -- (1.4,3);
    \draw[thick] (0.6,3.3) -- (1.4,3.3);
    
    % Lighthouse roof (dome-shaped top)
    \fill[red!80] (0.65,3.3) -- (1.35,3.3) -- (1,3.8) -- cycle;
    
    % Lighthouse light
    \fill[yellow] (1,3.15) circle (0.1);
    
    % Lighthouse light beams (enhanced beams with glow effect)
    \draw[yellow!80, thick] (1,3.15) -- (0,4);
    \draw[yellow!80, thick] (1,3.15) -- (2,4);
    
    % Lighthouse door
    \fill[brown!80] (0.9,-0.5) rectangle (1.1,-0.2);
    \draw[thick] (0.9,-0.5) -- (0.9,-0.2) -- (1.1,-0.2) -- (1.1,-0.5) -- cycle;
    
    % Boat's hull (with rounded edges for cartoon style)
    \fill[brown!70] (4,-0.5) .. controls (4.2,-1.5) and (5.8,-1.5) .. (6,-0.5) -- cycle;
    \draw[thick] (4,-0.5) .. controls (4.2,-1.5) and (5.8,-1.5) .. (6,-0.5) -- cycle;
    
    % Boat's deck
    \fill[brown!50] (4.2,-0.5) rectangle (5.8,0);
    
    % Mast of the boat
    \draw[thick] (5,0) -- (5,1.5);
    
    % Sail (cartoonish curved sail)
    \fill[white] (5,0.2) .. controls (5.8,0.7) and (5.8,1.3) .. (5,1.4) -- cycle;
    \draw[thick] (5,0.2) .. controls (5.8,0.7) and (5.8,1.3) .. (5,1.4) -- cycle;
    
    % Little flag on the top of the mast
    \fill[red] (5,1.5) -- (5.4,1.6) -- (5,1.7) -- cycle;
    
    % Add some cartoonish waves
    \foreach \x in {3.5, 5.5, 7.5} {
        \draw[blue!60,thick] (\x,-0.5) .. controls (\x+0.3,-0.6) and (\x+0.7,-0.4) .. (\x+1,-0.5);
    }
    
    % Second boat (smaller and farther away)
    \fill[brown!70] (7,-0.5) .. controls (7.1,-0.9) and (7.9,-0.9) .. (8,-0.5) -- cycle; % hull
    \draw[thick] (7,-0.5) .. controls (7.1,-0.9) and (7.9,-0.9) .. (8,-0.5) -- cycle;
    
    \draw[thick] (7.5,-0.5) -- (7.5,0.2); % mast
    \fill[white] (7.5,-0.1) .. controls (8,0) and (8,0.1) .. (7.5,0.2) -- cycle; % sail
    \draw[thick] (7.5,-0.1) .. controls (8,0) and (8,0.1) .. (7.5,0.2) -- cycle;
    
    % Flag for the second boat
    \fill[red] (7.5,0.2) -- (7.7,0.3) -- (7.5,0.4) -- cycle;
    
    % Marking points A (top of lighthouse), B (bottom of lighthouse), C (first boat), D (second boat)
    \fill[black] (1,3.5) circle (0.05) node[above] {A};  % Point A (top of lighthouse)
    \fill[black] (1,-0.5) circle (0.05) node[below] {B}; % Point B (bottom of lighthouse)
    \fill[black] (5,-0.5) circle (0.05) node[below] {C}; % Point C (position of boat 1)
    \fill[black] (7.5,-0.5) circle (0.05) node[below] {D}; % Point D (position of boat 2)
    
    % Draw lines connecting A-B, A-C, A-D
    \draw[thick] (1,3.5) -- (1,-0.5);  % Line AB
    \draw[thick] (1,3.5) -- (5,-0.5);  % Line AC
    \draw[thick] (1,3.5) -- (7.5,-0.5); % Line AD
    
    % Draw line B-C-D
    \draw[thick] (1,-0.5) -- (5,-0.5) -- (7.5,-0.5);  % Line BCD
    
    % Add a line x parallel to BC through A
    \draw[dashed] (1,3.5) -- (6,3.5) node[right] {$x$}; % Dashed line x, parallel to BC
    
    % Add angle alpha for xAD
    \draw[thick] (1.6,2.9) arc[start angle=-60, end angle= -10, radius=0.8];
    \node at (1.6,3.2) {$\beta$};  % Label for angle alpha
    
    % Add angle beta for xAC
    \draw[thick] (1.9,2.9) arc[start angle=-25, end angle= -13, radius=3.1];
    \node at (2.3,3.3) {$\alpha$};  % Label for angle beta
    
    % Add a smiling sun in the sky
    \fill[yellow] (9,3.5) circle (0.5); % sun body
    \draw (9,3.7) arc[start angle=45, end angle=135, radius=0.2]; % smiling face
    \fill[black] (8.8,3.6) circle (0.05); % left eye
    \fill[black] (9.2,3.6) circle (0.05); % right eye
    \foreach \angle in {0,45,...,315} {
        \draw[thick, yellow] (9,3.5) -- ++(\angle:0.8);
    }
    
\end{tikzpicture}
}
\end{minipage}

\vspace{3pt}\hrule\vspace{3pt}\footnotesize\textcolor{gray!80!black}{\textbf{Các mã dùng chung hình (HINH\_225):} \texttt{L09\_C4\_B2\_25, L09\_C4\_B2\_24}}
\end{tcolorbox}
\vspace{4pt}


\begin{tcolorbox}[
    breakable,
    colback=white,
    colframe=blue!60!black,
    boxrule=0.8pt,
    arc=3pt,
    title={\textbf{Câu \#304} \quad \textbf{MÃ CÂU: \texttt{[L09\_C4\_B2\_43]}} \quad | \quad \textbf{Nhóm: \texttt{HINH\_236}} \quad | \quad \textbf{MapID: \texttt{9H4V2-5}} \hfill \textcolor{yellow!90!white}{[Hình gốc .pyc]}}
]
\textbf{Dạng toán:} SA-M3. Toán đo khoảng cách kinh khí cầu đến mặt đất \hfill \textit{(Trả lời ngắn -- Vận dụng)}
\vspace{4pt}\hrule\vspace{6pt}
\begin{minipage}[t]{0.58\textwidth}
\vspace{0pt}
Một cậu bé cao ${1,5}m$ đứng ngắm khinh khí cầu ở một vị trí cách hình chiếu của nó trên một đất tức điểm ${E}$ một khoảng là $110$m. Cậu bé đặt ống nhòm hợp với phương ngang một góc $40^{\circ}$ thì nhìn thấy khinh khí cầu. Hỏi khi đó khinh khí cầu cách mặt đất bao nhiêu mét?(Kết quả làm tròn đến hàng đơn vị)


\textbf{Đáp án điền số:} 	exttt{94}
\end{minipage}\hfill
\begin{minipage}[t]{0.39\textwidth}
\vspace{0pt}
\centering

\resizebox{\linewidth}{!}{%
\begin{tikzpicture}[line join = round, line cap=round,>=stealth,font=\footnotesize,scale=1]
\path
(0,0) coordinate (D)
($(D)+(0,1)$) coordinate (B)
($(D)+(4,0)$) coordinate (E)
($(E)+(0,1)$) coordinate (A)
($(E)+(0,4)$) coordinate (C)
($(C)+(-2,0)$) coordinate (x)
;
\draw (D)--(E)--(C)--(B)--cycle (B)--(A);
\foreach \p/\r in {A/0,B/90,C/0,D/-90,E/-90}
\fill (\p) circle (1.2pt) node[shift={(\r:3mm)}]{$\p$};
\draw[pattern=north west lines] ($(D)+(-1,0)$) rectangle ($(E)+(1,-0.5)$);
\path
(D)--(E) node[above,pos=0.5]{$110$ m} ;
\draw[<->] ($(D)+(-0.75,0)$)--($(B)+(-0.75,0)$) node[left,pos=0.5]{$1,5$ m};
\draw pic[draw,angle radius=3mm] {right angle = B--A--C};
\draw pic[draw,angle radius=8mm] {angle = A--B--C};
\draw ($(B)+(1.25,0)$) node[above]{$40^\circ$};
\draw[dashed] (C)--($(x)+(-2,0)$)node[above]{$x$};
\begin{scope}[line width=3pt]
\foreach \yA/\xB/\yB in {0.6/-0.3/0.05, 0.6/-0.15/0.05} {
\draw (-0.25, 0.45) -- (-0.25, \yA) -- (\xB, \yB);
}
\foreach \yA/\xB/\yB/\xC/\yC in {
0.9/-0.1/0.8/0.0/0.9,
0.9/-0.4/0.8/-0.3/0.7
} {
\draw (-0.25, 0.45) -- (-0.25, \yA) -- (\xB, \yB) -- (\xC, \yC);
}
\draw (-0.25, 0.45) -- (-0.25, 1.0)circle(2pt);
\end{scope}
\fill [scale=.1,black,yshift=40cm,xshift=40cm]
(-1,0) rectangle (1,1)
(-1,2).. controls +(135:1) and +(180:4) .. (0,7)
.. controls +(180:2) and +(135:2).. (-.5,2)
.. controls +(120:2) and +(180:1) .. (0,6.99)
.. controls +(0:1) and +(60:2) .. (.5,2)
.. controls +(45:2) and +(0:2) .. (0,7)
.. controls +(0:4) and +(45:1).. (1,2) ;
\end{tikzpicture}
}
\end{minipage}

\end{tcolorbox}
\vspace{4pt}


\begin{tcolorbox}[
    breakable,
    colback=white,
    colframe=blue!60!black,
    boxrule=0.8pt,
    arc=3pt,
    title={\textbf{Câu \#305} \quad \textbf{MÃ CÂU: \texttt{[L09\_C4\_B2\_44]}} \quad | \quad \textbf{Nhóm: \texttt{HINH\_237}} \quad | \quad \textbf{MapID: \texttt{9H4V2-5}} \hfill \textcolor{yellow!90!white}{[Hình gốc .pyc]}}
]
\textbf{Dạng toán:} SA-M3. Toán tính chiều cao của quạt gió \hfill \textit{(Trả lời ngắn -- Vận dụng)}
\vspace{4pt}\hrule\vspace{6pt}
\begin{minipage}[t]{0.58\textwidth}
\vspace{0pt}
Một cậu bé cao ${1,4}m$ đứng cách thân một cái quạt gió $17$ m và nhìn thấy tâm của cánh quạt với góc nâng $47^{\circ}$ (Hình vẽ bên). Tính khoảng cách từ tâm của cánh quạt đến mặt đất.(Kết quả làm tròn đến chữ số thập phân thứ nhất)


\textbf{Đáp án điền số:} 	exttt{19,6}
\end{minipage}\hfill
\begin{minipage}[t]{0.39\textwidth}
\vspace{0pt}
\centering

\resizebox{\linewidth}{!}{%
\begin{tikzpicture}[smooth,font=\footnotesize,scale=1]
    \coordinate (M) at ($(4,4.16)+(-60:2cm)$);
    \coordinate (N) at ($(4,4.16)+(60:2cm)$);
    \coordinate (K) at ($(4,4.16)+(180:2cm)$);
    \fill[color=red](-0.05,1)--(0.05,1)--(0.05,0)--(-0.05,0)--cycle;
    \fill[color=blue](3.95,0)--(4.05,0)--(4.05,4)--(3.95,4)--cycle;
    \fill[ball color=red](0,1.16) circle (0.16);
    \fill[ball color=blue](4,4.16) circle (0.16);
    \draw[dashed,blue] (0.16,1.16)--(3.95,1.16) (0.16,1.16)--(3.84,4.16);
    \draw[very thick,blue] (4,4.16)--(M) (4,4.16)--(N) (4,4.16)--(K);
    \path
    (3.95,1.16)--(3.95,0) node[below,midway,sloped,scale=1]{$1,4m$ }
    (0.16,1.16)--(3.95,1.16) node[below,midway,sloped,scale=1]{$17m$ };
    \node[right] at (0.9,1.45) {$47^\circ$};
    \begin{scope}
        \clip (3.95,1.16)--(0.16,1.16)--(3.84,4.16);
        \draw[double] (0.16,1.16) circle(0.7cm);
    \end{scope}
    
    \end{tikzpicture}
}
\end{minipage}

\end{tcolorbox}
\vspace{4pt}


\begin{tcolorbox}[
    breakable,
    colback=white,
    colframe=blue!60!black,
    boxrule=0.8pt,
    arc=3pt,
    title={\textbf{Câu \#306} \quad \textbf{MÃ CÂU: \texttt{[L09\_C4\_B2\_45]}} \quad | \quad \textbf{Nhóm: \texttt{HINH\_224}} \quad | \quad \textbf{MapID: \texttt{9H4V2-5}} \hfill \textcolor{yellow!90!white}{[Hình gốc .pyc]}}
]
\textbf{Dạng toán:} SA-M3. Toán thực tế tính khoảng cách cù lao đến bờ sông. \hfill \textit{(Trả lời ngắn -- Vận dụng)}
\vspace{4pt}\hrule\vspace{6pt}
\begin{minipage}[t]{0.58\textwidth}
\vspace{0pt}
Hai ngư dân đứng ở một bên bờ sông cách nhau $AB={290}m$ cùng nhìn thấy một cù lao trên sông với các góc nhìn lần lượt là  $\beta=54^{\circ}$ và $\alpha=27^{\circ}$. Tính khoảng cách từ cù lao ${C}$ đến bờ sông ${AB}$ (làm tròn đến hàng đơn vị).


\textbf{Đáp án điền số:} 	exttt{108}
\end{minipage}\hfill
\begin{minipage}[t]{0.39\textwidth}
\vspace{0pt}
\centering

\resizebox{\linewidth}{!}{%
\begin{tikzpicture}[scale=0.7, font=\footnotesize, line join = round, line cap = round,>=stealth]
    \clip (-4.39,3.2) rectangle (4,-4);
    \draw[pattern=north east lines,opacity=0.3] plot[smooth] coordinates{(-4.39,1.43)(-2.43,1.46) (-1.48,1.67)(1.43,2.12)(2.01,2.09)(4,1.64) (4.1,-1.06)(3.02,-0.77)(1.85,-1.38)(0.95,-1.48)(0.13,-1.72)(-1.48,-1.67)(-2.46,-1.51)(-3.57,-0.98)(-4.37,-1.08)};
    \draw[opacity=3] plot[smooth] coordinates{(-4.39,1.43)(-2.43,1.46) (-1.48,1.67)(1.43,2.12)(2.01,2.09)(4,1.64) (4.1,-1.06)(3.02,-0.77)(1.85,-1.38)(0.95,-1.48)(0.13,-1.72)(-1.48,-1.67)(-2.46,-1.51)(-3.57,-0.98)(-4.37,-1.08)};
    \draw[fill=white] plot[smooth  cycle] coordinates{(-2.75,-0.08)(-1.83,0.32) (-0.03,0.42) (1,0) (0.24,-0.48)(-0.58,-0.79)(-1.38,-0.77)(-1.91,-0.71)};
    \draw[fill=black!70]  plot[smooth  cycle] coordinates{(-2.14,0.58)(-2.09,0.24)(-2.33,-0.13) (-1.19,-0.13) (-1.38,-0.13) (-1.69,0.21)(-1.75,0.53)};
    \draw[fill=blue!30]  plot[smooth  cycle] coordinates{(-1.75,0.53)(-1.46,0.79) (-1.01,0.58) (-1.08,1.06) (-0.64,1.08)(-0.93,1.59)(-0.53,1.85)(-0.93,2.2)(-1.38,2.7)(-2.04,3.18)(-2.49,2.91)(-3.07,2.91)(-3.2,2.22)(-3.73,1.96)(-3.1,1.3)(-3.31,1.01)(-2.83,1.01)(-2.99,0.69)(-2.14,0.58)};
    \tkzDefPoints{-1.38/-0.13/C,-2/-3/B,3/-2/A}
    \tkzDrawPoints[fill=black](A,B,C)
    \tkzDrawPolygon(A,B,C)
    \tkzDefMidPoint(A,B) 
    \tkzLabelPoints[below](A,B)
    \tkzLabelPoints[above right](C)
    \tkzMarkAngles[size=1cm,arc=l,mark=|](C,A,B)
    \tkzMarkAngles[size=1cm,arc=l,mark=||](A,B,C)
    \tkzLabelAngles[left=1cm,pos=0.5,rotate=-20](B,A,C){\footnotesize$\alpha$}
    \tkzLabelAngles[right,pos=1](A,B,C){\footnotesize$\beta$}
 \end{tikzpicture}
}
\end{minipage}

\vspace{3pt}\hrule\vspace{3pt}\footnotesize\textcolor{gray!80!black}{\textbf{Các mã dùng chung hình (HINH\_224):} \texttt{L09\_C4\_B2\_14, L09\_C4\_B2\_39, L09\_C4\_B2\_30}}
\end{tcolorbox}
\vspace{4pt}


\begin{tcolorbox}[
    breakable,
    colback=white,
    colframe=blue!60!black,
    boxrule=0.8pt,
    arc=3pt,
    title={\textbf{Câu \#307} \quad \textbf{MÃ CÂU: \texttt{[L09\_C4\_B2\_46]}} \quad | \quad \textbf{Nhóm: \texttt{HINH\_223}} \quad | \quad \textbf{MapID: \texttt{9H4V2-5}} \hfill \textcolor{yellow!90!white}{[Hình gốc .pyc]}}
]
\textbf{Dạng toán:} SA-M3. Toán thực tế tính vận tốc của máy bay khi bay lên \hfill \textit{(Trả lời ngắn -- Vận dụng)}
\vspace{4pt}\hrule\vspace{6pt}
\begin{minipage}[t]{0.58\textwidth}
\vspace{0pt}
Một máy bay vừa phi lên theo đường bay hợp với mặt đất một góc $\widehat{ATH} = {30}^{\circ}$. Sau $1,2$ phút máy bay ở độ cao $a={5}$km so với mặt đất. Hỏi vận tốc của máy bay khi bay lên? (Kết quả làm tròn đến hàng đơn vị).


\textbf{Đáp án điền số:} 	exttt{500}
\end{minipage}\hfill
\begin{minipage}[t]{0.39\textwidth}
\vspace{0pt}
\centering

\resizebox{\linewidth}{!}{%
\begin{tikzpicture}[declare function={r=4;rh=r*sin(30);}]
        \path (0,0) coordinate (T)
        --(30:r) coordinate (A)
        node[above,shift={(0.2,-0.05)}]{
            \resizebox{1cm}{!}{\begin{tikzpicture}[line join=round]
                    \draw[line width=2pt,line cap=butt,rounded corners=2pt] 
                    (0.3,0) arc(270:85:0.35 and 0.4)--(0.8,1.5)..controls +(50:0.2) and +(180:0.2)..(1.4,1.8)--(2.5,1.6)..controls +(0:0.35) and +(140:0.15)..(3,1.45)--(3.8,1.45)--(4.3,2.15)--(4.85,2.15)
                    --(4.3,0.95)--(3.1,0.7)..controls +(-120:0.65) and +(0:0.1)..(2.3,0)--cycle;
                    \draw[line width=2pt,line cap=butt] 
                    (4.7,1.25) arc (57:-150:0.45)
                    (0.1,-0.4) ..controls +(-30:0.15) and +(180:0.1)..
                    (0.5,-0.475)--(2.1,-0.475)..controls +(0:0.2) and +(-100:0.1)..(2.7,-0.375)
                    (-0.05,2.45)--(3.4,2.45);
                    \path[fill] (1.6,1.9)--++(0.18,0)--+(0,0.35) arc (0:180:0.09 and 0.13)--cycle;
                    \path[fill] (1.8,0.8){[rounded corners=3pt]--++(0.5,0)--++(0,0.7)--+(-0.5,0)--cycle}
                    {[rounded corners=2pt](0.8,0.8)--(1.55,0.8)--++(0,0.7)--+(-0.2,0)}
                    ..controls +(200:0.325) and +(85:0.15)..(0.8,0.8)--cycle;
                \end{tikzpicture}
            }
        }--([turn]-120:rh) coordinate (H)--cycle;
        \draw 
            (A)--(H) node[pos=0.5,right]{$a$ km}--(T)--(A)
          
            %(O) pic[draw,angle radius = 6] {right angle = B--O--C}
        ;
        \foreach \t/\g in {T/-90,H/-90,A/-30}{
            \draw[fill=white] (\t) circle (1pt) node[shift={(\g:7pt)},font=\scriptsize]{$ \t $};
        }
    \end{tikzpicture}
}
\end{minipage}

\vspace{3pt}\hrule\vspace{3pt}\footnotesize\textcolor{gray!80!black}{\textbf{Các mã dùng chung hình (HINH\_223):} \texttt{L09\_C4\_B2\_08, L09\_C4\_B2\_31}}
\end{tcolorbox}
\vspace{4pt}


\begin{tcolorbox}[
    breakable,
    colback=white,
    colframe=blue!60!black,
    boxrule=0.8pt,
    arc=3pt,
    title={\textbf{Câu \#308} \quad \textbf{MÃ CÂU: \texttt{[L09\_C4\_B2\_47]}} \quad | \quad \textbf{Nhóm: \texttt{HINH\_238}} \quad | \quad \textbf{MapID: \texttt{9H4V2-5}} \hfill \textcolor{yellow!90!white}{[Hình gốc .pyc]}}
]
\textbf{Dạng toán:} SA-M3. Toán thực tế tính chiều cao toà nhà \hfill \textit{(Trả lời ngắn -- Vận dụng)}
\vspace{4pt}\hrule\vspace{6pt}
\begin{minipage}[t]{0.58\textwidth}
\vspace{0pt}
Khi tia sáng mặt trời tạo với mặt đất một góc ${63}^{\circ}$ thì bóng của một tòa nhà trên mặt đất có độ dài $b={30}$m. Tính chiều cao của tòa nhà (kết quả làm tròn đến hàng đơn vị).


\textbf{Đáp án điền số:} 	exttt{59}
\end{minipage}\hfill
\begin{minipage}[t]{0.39\textwidth}
\vspace{0pt}
\centering

\resizebox{\linewidth}{!}{%
\begin{tikzpicture}[scale=1,font=\footnotesize,line cap=round,line join=round]
    % Tòa nhà ----------------------------------------------
    \path
    (3,0) coordinate (A)
    (3,5) coordinate (B)
    (8,0) coordinate (C)
    ($(A)!.5!(C)$) node[below] {$30$ m}
    ($(C)+(162:1)$) node {$63^{\circ}$}
    ;
    \foreach \i in {0,0.5,1,1.5,2,2.5}{\foreach \j in {0,0.8,1.6,2.4,3.2,4.0}
        \draw[fill=yellow!50!white] (\i,\j) rectangle +(0.5,0.8);
    } 
    % Ban công 
    \foreach \i in {0.1,0.6,1.1,1.6,2.1,2.6}{\foreach \j in {0,0.8,1.6,2.4,3.2,4.0}
        \draw[fill=orange!60!white] (\i,\j) rectangle +(0.3,0.6);
    }
    \foreach \j in {0,0.8,1.6,2.4,3.2,4.0,4.8}
        \draw[fill=cyan!30!white] (0,\j) rectangle +(3,0.3);
    
    \draw (A)--(B)--(C)--cycle;    %       \tkzMarkRightAngles(B,A,C)
    %       \tkzMarkAngles[size=.7](B,C,A)
    
    \foreach \i/\p in {A/-90,B/90,C/0}  
        \fill[black] (\i) circle (1.5pt)  % Điểm đánh dấu màu đen
        node[shift={(\p:3mm)},black] {$\i$}; % Chữ cái màu đen
\end{tikzpicture}
}
\end{minipage}

\end{tcolorbox}
\vspace{4pt}


\begin{tcolorbox}[
    breakable,
    colback=white,
    colframe=blue!60!black,
    boxrule=0.8pt,
    arc=3pt,
    title={\textbf{Câu \#309} \quad \textbf{MÃ CÂU: \texttt{[L09\_C4\_B2\_48]}} \quad | \quad \textbf{Nhóm: \texttt{HINH\_221}} \quad | \quad \textbf{MapID: \texttt{9H4V2-5}} \hfill \textcolor{yellow!90!white}{[Hình gốc .pyc]}}
]
\textbf{Dạng toán:} SA-M3. Toán thực tế tính chiều cao của cái cây \hfill \textit{(Trả lời ngắn -- Vận dụng)}
\vspace{4pt}\hrule\vspace{6pt}
\begin{minipage}[t]{0.58\textwidth}
\vspace{0pt}
Khi tia sáng mặt trời tạo với mặt đất một góc $\alpha={53}^{\circ}$ thì bóng của một cái cây trên mặt đất có độ dài $x={5}$m. Tính chiều cao của cái cây (kết quả làm tròn đến hàng đơn vị).


\textbf{Đáp án điền số:} 	exttt{7}
\end{minipage}\hfill
\begin{minipage}[t]{0.39\textwidth}
\vspace{0pt}
\centering

\resizebox{\linewidth}{!}{%
\begin{tikzpicture}[scale=1, font=\footnotesize, line cap=round, line join=round]

    % Vẽ ông mặt trời
    \fill[yellow] (-3,7) circle (0.5cm); % Di chuyển mặt trời ra xa một chút

    % Vẽ các tia nắng sinh động bao quanh ông mặt trời (từng tia riêng lẻ)
    \pgfmathsetmacro{\radius}{0.8} % Khoảng cách từ tâm mặt trời đến điểm bắt đầu của tia nắng
    
    % Vẽ từng tia nắng riêng lẻ
    \pgfmathsetmacro{\lengthA}{0.5 + 0.4 * sin(0 * 8.8)}
    \draw[orange, thick] (-3,7) ++(0:\radius) -- ++(0:\lengthA);

    \pgfmathsetmacro{\lengthB}{0.5 + 0.4 * sin(15 * 8.8)}
    \draw[orange, thick] (-3,7) ++(15:\radius) -- ++(15:\lengthB);

    \pgfmathsetmacro{\lengthC}{0.5 + 0.4 * sin(30 * 8.8)}
    \draw[orange, thick] (-3,7) ++(30:\radius) -- ++(30:\lengthC);

    \pgfmathsetmacro{\lengthD}{0.5 + 0.4 * sin(45 * 8.8)}
    \draw[orange, thick] (-3,7) ++(45:\radius) -- ++(45:\lengthD);

    \pgfmathsetmacro{\lengthE}{0.5 + 0.4 * sin(60 * 8.8)}
    \draw[orange, thick] (-3,7) ++(60:\radius) -- ++(60:\lengthE);

    \pgfmathsetmacro{\lengthF}{0.5 + 0.4 * sin(75 * 8.8)}
    \draw[orange, thick] (-3,7) ++(75:\radius) -- ++(75:\lengthF);

    \pgfmathsetmacro{\lengthG}{0.5 + 0.4 * sin(90 * 8.8)}
    \draw[orange, thick] (-3,7) ++(90:\radius) -- ++(90:\lengthG);

    \pgfmathsetmacro{\lengthH}{0.5 + 0.4 * sin(105 * 8.8)}
    \draw[orange, thick] (-3,7) ++(105:\radius) -- ++(105:\lengthH);

    \pgfmathsetmacro{\lengthI}{0.5 + 0.4 * sin(120 * 8.8)}
    \draw[orange, thick] (-3,7) ++(120:\radius) -- ++(120:\lengthI);

    \pgfmathsetmacro{\lengthJ}{0.5 + 0.4 * sin(135 * 8.8)}
    \draw[orange, thick] (-3,7) ++(135:\radius) -- ++(135:\lengthJ);

    \pgfmathsetmacro{\lengthK}{0.5 + 0.4 * sin(150 * 8.8)}
    \draw[orange, thick] (-3,7) ++(150:\radius) -- ++(150:\lengthK);

    \pgfmathsetmacro{\lengthL}{0.5 + 0.4 * sin(165 * 8.8)}
    \draw[orange, thick] (-3,7) ++(165:\radius) -- ++(165:\lengthL);

    \pgfmathsetmacro{\lengthM}{0.5 + 0.4 * sin(180 * 8.8)}
    \draw[orange, thick] (-3,7) ++(180:\radius) -- ++(180:\lengthM);

    \pgfmathsetmacro{\lengthN}{0.5 + 0.4 * sin(195 * 8.8)}
    \draw[orange, thick] (-3,7) ++(195:\radius) -- ++(195:\lengthN);

    \pgfmathsetmacro{\lengthO}{0.5 + 0.4 * sin(210 * 8.8)}
    \draw[orange, thick] (-3,7) ++(210:\radius) -- ++(210:\lengthO);

    \pgfmathsetmacro{\lengthP}{0.5 + 0.4 * sin(225 * 8.8)}
    \draw[orange, thick] (-3,7) ++(225:\radius) -- ++(225:\lengthP);

    \pgfmathsetmacro{\lengthQ}{0.5 + 0.4 * sin(240 * 8.8)}
    \draw[orange, thick] (-3,7) ++(240:\radius) -- ++(240:\lengthQ);

    \pgfmathsetmacro{\lengthR}{0.5 + 0.4 * sin(255 * 8.8)}
    \draw[orange, thick] (-3,7) ++(255:\radius) -- ++(255:\lengthR);

    \pgfmathsetmacro{\lengthS}{0.5 + 0.4 * sin(270 * 8.8)}
    \draw[orange, thick] (-3,7) ++(270:\radius) -- ++(270:\lengthS);

    \pgfmathsetmacro{\lengthT}{0.5 + 0.4 * sin(285 * 8.8)}
    \draw[orange, thick] (-3,7) ++(285:\radius) -- ++(285:\lengthT);

    \pgfmathsetmacro{\lengthU}{0.5 + 0.4 * sin(300 * 8.8)}
    \draw[orange, thick] (-3,7) ++(300:\radius) -- ++(300:\lengthU);

    \pgfmathsetmacro{\lengthV}{0.5 + 0.4 * sin(315 * 8.8)}
    \draw[orange, thick] (-3,7) ++(315:\radius) -- ++(315:\lengthV);

    \pgfmathsetmacro{\lengthW}{0.5 + 0.4 * sin(330 * 8.8)}
    \draw[orange, thick] (-3,7) ++(330:\radius) -- ++(330:\lengthW);

    \pgfmathsetmacro{\lengthX}{0.5 + 0.4 * sin(345 * 8.8)}
    \draw[orange, thick] (-3,7) ++(345:\radius) -- ++(345:\lengthX);

    % Vẽ cái cây cao hơn
    \fill[brown] (0,0) rectangle (0.5,4); % Tăng chiều cao của cái cây
    
    % Vẽ tán cây bằng các hình tròn chồng lên nhau, thẳng hàng với thân cây và xếp cao dần
    \fill[green] (0.25,3.5) circle (0.8cm); % Vòng tròn lớn dưới cùng
    \fill[green] (0.25,4.5) circle (1.2cm); % Vòng tròn trung bình ở giữa
    \fill[green] (0.25,4) circle (1cm); % Vòng tròn nhỏ ở trên cùng
    \fill[brown] (0.2,0) ellipse (0.8 and 0.24); % Elip với kích thước tuỳ chỉnh

    % Định nghĩa các điểm A, B, C
    \path
    (0.25,0) coordinate (A) % Gốc cây
    (0.25,5.6) coordinate (B) % Ngọn cây
    (4,0) coordinate (C); % Điểm trên mặt đất, đã di chuyển ra xa
    
    % Vẽ tam giác ABC
    \draw (A) -- (B) -- (C) -- cycle; % Nối các điểm A, B, C

    % Gắn nhãn các đoạn thẳng và góc
    \draw ($(A)!.5!(C)$) node[below] {$x$}; % Nhãn đoạn thẳng AC
    \node at ($(C)+(155:1)$) {$\alpha$}; % Ký hiệu góc BCA

    % Đánh dấu các điểm A, B, C
    \foreach \i/\p in {A/-90,B/90,C/0}  
    \fill[black] (\i) circle (1.5pt)  
    node[shift={(\p:3mm)},black] {$\i$}; % Chữ cái màu đen

\end{tikzpicture}
}
\end{minipage}

\vspace{3pt}\hrule\vspace{3pt}\footnotesize\textcolor{gray!80!black}{\textbf{Các mã dùng chung hình (HINH\_221):} \texttt{L09\_C4\_B1\_19, L09\_C4\_B2\_11, L09\_C4\_B2\_33}}
\end{tcolorbox}
\vspace{4pt}


\begin{tcolorbox}[
    breakable,
    colback=white,
    colframe=blue!60!black,
    boxrule=0.8pt,
    arc=3pt,
    title={\textbf{Câu \#310} \quad \textbf{MÃ CÂU: \texttt{[L09\_C4\_B2\_49]}} \quad | \quad \textbf{Nhóm: \texttt{HINH\_239}} \quad | \quad \textbf{MapID: \texttt{9H4C2-5}} \hfill \textcolor{yellow!90!white}{[Hình gốc .pyc]}}
]
\textbf{Dạng toán:} SA-M4. Toán thực tế tính thời gian của máy bay khi hạ cánh \hfill \textit{(Trả lời ngắn -- Vận dụng cao)}
\vspace{4pt}\hrule\vspace{6pt}
\begin{minipage}[t]{0.58\textwidth}
\vspace{0pt}
Quãng đường bay khi hạ cánh của máy bay, hay còn gọi là quãng đường tiếp cận cuối cùng (landing distance), là khoảng cách mà máy bay bay từ lúc bắt đầu tiếp cận cuối cùng cho đến khi chạm đất. Đây là giai đoạn máy bay đã hạ thấp gần mặt đất và đang hướng thẳng tới đường băng với góc hạ cánh cố định (thường là $\widehat{ATH}={3}^{\circ}$), chuẩn bị tiếp xúc với đường băng để hoàn tất quá trình hạ cánh. Một máy bay ở giai đoạn tiếp cận cuối cùng cách mặt đất $a={400}$m, tính thời gian (đơn vị phút) hạ cánh ở giai đoạn này biết vận tốc hạ cánh là $260$km/h.


\textbf{Đáp án điền số:} 	exttt{1,76}
\end{minipage}\hfill
\begin{minipage}[t]{0.39\textwidth}
\vspace{0pt}
\centering

\resizebox{\linewidth}{!}{%
\begin{tikzpicture}[declare function={r=4;rh=r*sin(30);}]
        \path (0,0) coordinate (T)
        --(30:r) coordinate (A)
        node[above,shift={(0.2,-0.05)}]{
            \resizebox{1cm}{!}{\begin{tikzpicture}[line join=round]
                    \draw[line width=2pt,line cap=butt,rounded corners=2pt] 
                    (0.3,0) arc(270:85:0.35 and 0.4)--(0.8,1.5)..controls +(50:0.2) and +(180:0.2)..(1.4,1.8)--(2.5,1.6)..controls +(0:0.35) and +(140:0.15)..(3,1.45)--(3.8,1.45)--(4.3,2.15)--(4.85,2.15)
                    --(4.3,0.95)--(3.1,0.7)..controls +(-120:0.65) and +(0:0.1)..(2.3,0)--cycle;
                    \draw[line width=2pt,line cap=butt] 
                    (4.7,1.25) arc (57:-150:0.45)
                    (0.1,-0.4) ..controls +(-30:0.15) and +(180:0.1)..
                    (0.5,-0.475)--(2.1,-0.475)..controls +(0:0.2) and +(-100:0.1)..(2.7,-0.375)
                    (-0.05,2.45)--(3.4,2.45);
                    \path[fill] (1.6,1.9)--++(0.18,0)--+(0,0.35) arc (0:180:0.09 and 0.13)--cycle;
                    \path[fill] (1.8,0.8){[rounded corners=3pt]--++(0.5,0)--++(0,0.7)--+(-0.5,0)--cycle}
                    {[rounded corners=2pt](0.8,0.8)--(1.55,0.8)--++(0,0.7)--+(-0.2,0)}
                    ..controls +(200:0.325) and +(85:0.15)..(0.8,0.8)--cycle;
                \end{tikzpicture}
            }
        }--([turn]-120:rh) coordinate (H)--cycle;
        %\draw[-]  (-r/2,0) --(3/2*r,0)node[below];
        \draw 
            (A)--(H) node[pos=0.5,right]{$400$ m}--(T)--(A)
        ;
        \foreach \t/\g in {T/-90,H/-90,A/-30}{
            \draw[fill=white] (\t) circle (1pt) node[shift={(\g:7pt)},font=\scriptsize]{$ \t $};
        }
    \end{tikzpicture}
}
\end{minipage}

\end{tcolorbox}
\vspace{4pt}


\begin{tcolorbox}[
    breakable,
    colback=white,
    colframe=blue!60!black,
    boxrule=0.8pt,
    arc=3pt,
    title={\textbf{Câu \#311} \quad \textbf{MÃ CÂU: \texttt{[L09\_C4\_B2\_50]}} \quad | \quad \textbf{Nhóm: \texttt{HINH\_229}} \quad | \quad \textbf{MapID: \texttt{9H4C2-5}} \hfill \textcolor{yellow!90!white}{[Hình gốc .pyc]}}
]
\textbf{Dạng toán:} SA-M4. Bài toán thực tế tính chiều cao của con dốc \hfill \textit{(Trả lời ngắn -- Vận dụng cao)}
\vspace{4pt}\hrule\vspace{6pt}
\begin{minipage}[t]{0.58\textwidth}
\vspace{0pt}
Lúc 6h sáng, bạn Anh đi xe đạp từ nhà (điểm ${A}$) đến trường (điểm ${B}$) phải leo và xuống một con dốc (như hình vẽ sau) biết ${AB}=1,6km$, $\widehat{A}=6^{\circ}$, $\widehat{B}=4^{\circ}$. Tính chiều cao ${h}$ của con dốc (theo mét làm tròn đến chữ số hàng đơn vị)


\textbf{Đáp án điền số:} 	exttt{67}
\end{minipage}\hfill
\begin{minipage}[t]{0.39\textwidth}
\vspace{0pt}
\centering

\resizebox{\linewidth}{!}{%
\begin{tikzpicture}
    
    % Vẽ tam giác ABC
    \coordinate [label=left:A] (A) at (0,0);
    \coordinate [label=right:B] (B) at (8,0);
    \coordinate [label=above:C] (C) at (2,1);
    
    % Đặt các góc và chiều cao
    \draw (A) -- (B) -- (C) -- cycle;
    
    % Vẽ chiều cao từ C tới AB
\draw (C) -- (2,0) coordinate (H);
\node[below] at (H) {H};
      % Vẽ dấu vuông tại điểm H
        \draw (2,0) rectangle +(0.1,0.1);
    % Nhãn chiều cao
    \draw[<-] (C) -- ($(C)!0.5!(H)$) node[midway, right] {$h$};
    
    % Vẽ các cung góc
    
    \node at ($(A)+(0.6,0.15)$) {$6^{\circ}$};
    

    \node at ($(B)+(-1.9,0.15)$) {$4^{\circ}$};
    
\end{tikzpicture}
}
\end{minipage}

\vspace{3pt}\hrule\vspace{3pt}\footnotesize\textcolor{gray!80!black}{\textbf{Các mã dùng chung hình (HINH\_229):} \texttt{L09\_C4\_B2\_13, L09\_C4\_B2\_37, L09\_C4\_B2\_36}}
\end{tcolorbox}
\vspace{4pt}


\begin{tcolorbox}[
    breakable,
    colback=white,
    colframe=blue!60!black,
    boxrule=0.8pt,
    arc=3pt,
    title={\textbf{Câu \#312} \quad \textbf{MÃ CÂU: \texttt{[L09\_C4\_B2\_51]}} \quad | \quad \textbf{Nhóm: \texttt{HINH\_230}} \quad | \quad \textbf{MapID: \texttt{9H4C2-5}} \hfill \textcolor{yellow!90!white}{[Hình gốc .pyc]}}
]
\textbf{Dạng toán:} SA-M4. Bài toán thực tế khó \hfill \textit{(Trả lời ngắn -- Vận dụng cao)}
\vspace{4pt}\hrule\vspace{6pt}
\begin{minipage}[t]{0.58\textwidth}
\vspace{0pt}
Một người đang ở trên sân thượng ${A}$ của một cái toà nhà nhìn xuống thấy một con đường chạy thẳng đến chân tường. Anh ta nhìn thấy một chiếc xe máy với góc hạ $\widehat{xAC}=34^{\circ}$. Sau $8$ phút lại nhìn thấy xe đó với góc hạ $\widehat{xAD}=65^{\circ}$. Hỏi sau bao nhiêu phút thì xe máy chạy đến chân tường biết vận tốc của xe máy không đổi (làm tròn kết quả đến hàng phần mười).


\textbf{Đáp án điền số:} 	exttt{3,7}
\end{minipage}\hfill
\begin{minipage}[t]{0.39\textwidth}
\vspace{0pt}
\centering

\resizebox{\linewidth}{!}{%
\begin{tikzpicture}
    \def\hh{6cm}
    \def\aa{5cm}
    \path (0,0) coordinate (B) (0,\hh) coordinate (A) (\aa,0) coordinate (C)
    (3*\aa/5,0) coordinate (D)      (3*\aa/5,\hh) coordinate (x);
        \draw[->] (0,\hh) -- (3*\aa/5,\hh) node[above] {$x$};
    \draw 
(A) pic[draw,angle radius = 6mm] {angle =D--A--x} 
(A) pic[draw,angle radius = 5mm] {angle =C--A--x}   ;
        % Hiển thị nhãn cho các điểm
    \draw (A) node [above] {$A$};
    \draw (B) node [below] {$B$};
    \draw (C) node [below] {$C$};
    \draw (D) node [below right] {$D$};
    % Vẽ tia Ax song song với BC

    
    
    % Vẽ hình tứ giác với các điểm A, B, C, D
    \draw (A) -- (B) -- (C) -- (D) -- cycle;
    
    % Vẽ đường nối A và C bằng nét liền
    \draw (A) -- (C);
    
    % Vẽ tường mô phòng
    \path[fill=red!30] (B) -- ++(180:5) 
    to[out=100,in=-90] (-4.2,1.5) 
    to[out=100,in=-90] (-3.7,3.5) 
    to[out=100,in=-90] (-2.8,4.8) 
    to[out=120,in=-110] (-2,6) -- (A) -- cycle;
    
    % Cắt vùng mô phỏng bức tường
    \clip (B) -- ++(180:5) 
    to[out=100,in=-90] (-4.2,1.5) 
    to[out=100,in=-90] (-3.7,3.5) 
    to[out=100,in=-90] (-2.8,4.8) 
    to[out=120,in=-110] (-2,6) -- (A) -- cycle;
    
    % Tô màu các viên gạch hình chữ nhật
    \fill[pattern=bricks, pattern color=red] (B) -- ++(180:5) 
    to[out=100,in=-90] (-4.2,1.5) 
    to[out=100,in=-90] (-3.7,3.5) 
    to[out=100,in=-90] (-2.8,4.8) 
    to[out=120,in=-110] (-2,6) -- (A) -- cycle;
    
    % Vẽ lại hình tứ giác và đường nối
    \draw (A)--(B)--(C)--(D)--cycle;
    \draw[dashed] (B)--(A);
    \draw (B)--++(180:5) to[out=100,in=-90](-4.2,1.5) to[out=100,in=-90](-3.7,3.5) to[out=100,in=-90](-2.8,4.8) to[out=120,in=-110](-2,6)--(A)--cycle;
    

\end{tikzpicture}
}
\end{minipage}

\vspace{3pt}\hrule\vspace{3pt}\footnotesize\textcolor{gray!80!black}{\textbf{Các mã dùng chung hình (HINH\_230):} \texttt{L09\_C4\_B2\_15}}
\end{tcolorbox}
\vspace{4pt}


\begin{tcolorbox}[
    breakable,
    colback=white,
    colframe=blue!60!black,
    boxrule=0.8pt,
    arc=3pt,
    title={\textbf{Câu \#313} \quad \textbf{MÃ CÂU: \texttt{[L09\_C4\_B2\_52]}} \quad | \quad \textbf{Nhóm: \texttt{HINH\_228}} \quad | \quad \textbf{MapID: \texttt{9H4C2-5}} \hfill \textcolor{yellow!90!white}{[Hình gốc .pyc]}}
]
\textbf{Dạng toán:} SA-M4. Bài toán thực tế tính chiều cao của hải đăng \hfill \textit{(Trả lời ngắn -- Vận dụng cao)}
\vspace{4pt}\hrule\vspace{6pt}
\begin{minipage}[t]{0.58\textwidth}
\vspace{0pt}
Ở hai địa điểm ${A}$ và ${B}$ trên mặt đất cách nhau $20$m, người ta đặt ống nhòm nhìn thấy vị trí ${D}$ cao nhất của ngọn hải đăng, góc đặt ống nhóm tại ${B}$ hợp với mặt đất là $45^{\circ}$, góc đặt ống nhòm tại vị trí ${A}$ hợp với mặt đất là $55^{\circ}$. Tính chiều cao của hải đăng theo đơn vị mét (làm tròn đến hàng đơn vị).


\textbf{Đáp án điền số:} 	exttt{67}
\end{minipage}\hfill
\begin{minipage}[t]{0.39\textwidth}
\vspace{0pt}
\centering

\resizebox{\linewidth}{!}{%
\begin{tikzpicture}[scale=.5,declare function={h=13.4;a=h/tan(63);b=h/tan(48);}]
        \fill[brown] (0,0) -- (-5,0) 
        to[out=20,in=180] (-4.3,.15)
        to[out=40,in=180] (-3.05,.65)
        to[out=80,in=200] (-2.8,.85)
        to[out=70,in=160] (-2.3,1.25)
        to[out=50,in=190] (-2,1.45)
        --(-1.6,1.45)--(-1.35,3.2)--(-1.3,2.15)
        to[out=50,in=210] (0,3.2)
        to[out=28,in=222] (1.25,4.05)
        --(1.5,1.5)--(1.7,1.5)
        to[out=-40,in=180] (2,1.33)
        to[out=-10,in=110] (2.45,1)
        to[out=-40,in=160] (3,.8)
        to[out=-40,in=150] (3.2,.6)
        to[out=10,in=130] (4.4,.1)
        to[out=0,in=160] (5,0)
        --cycle;
        \fill[brown] (-.85,3.6) rectangle (-.53,4.1);
        \fill[brown] (-1.3,4.2)
        to[out=50,in=208] (0,5.15)
        to[out=25,in=225] (1,5.75)
        to[out=50,in=275] (1.15,6.3)
        --(1.08,8)
        to[out=-130,in=37] (-1,6.6)
        to[out=-135,in=88] (-1.25,5.5)
        --cycle;
        \fill[brown] (.25,8.45) rectangle (.55,9);
        \fill[brown] (-1.1,8.3)
        to[out=55,in=245] (1.05,9.85)--(1.2,9.85)
        to[out=10,in=263] (1.45,10.1)
        to[out=100,in=-10] (.96,10.35)--(.93,11.1)--(1.05,11.15)
        to[out=155,in=20] (-1.07,11.2)--(-1.07,10.3)
        to[out=-175,in=135] (-1.4,10)
        to[out=-55,in=190] (-1.08,9.85)--cycle;
        \draw[brown,line width= 2pt] (.9,11.2)--(.9,11.7)
        (.25,11.2)--(.25,11.9)
        (-.35,11.2)--(-.35,11.9)
        (-1.03,11.2)--(-1.03,11.7)
        to[out=20,in=160] (0.94,11.7)
        (.55,11.8)--(.54,12.7)
        (-.6,11.8)--(-.59,12.7)
        (-.04,11.9)--(-.04,12.7);
        \draw[brown,line width=4pt] (.6,11.2)--(.6,11.8)
        (-.75,11.2)--(-.75,11.8);
        \fill[brown] (-.04,12.15) circle (7pt)
        (.6,12.7) arc (0:180:.65 and .7) -- cycle
        (-.04,13.4) circle (5pt);
        \path 
        (0,0) coordinate (C)
        (a,0) coordinate (A)
        (b,0) coordinate (B)
        (0,h) coordinate (D);
        \foreach \i/\g in {C/-90, A/-90, B/-90, D/90} \fill (\i)circle(.1) node[shift={(\g:.25)}] {$\i$};
        \draw (C)--(B)node[pos=.8, below]{$20m$}--(D)--cycle (A)--(D)
        (A)+(-1,0) arc (180:117:1)node[pos=.7, left]{$\alpha=55^\circ$};
        \draw[double] (B)+(-1,0) arc (180:132:1)node[pos=.7, left]{$\beta=45^\circ$};
        \draw[dashed] (-5,h)--(D);
        \draw[<->,-stealth] (-4,0) -- (-4,h)node[pos=.5, left]{$h$};
    \end{tikzpicture}
}
\end{minipage}

\vspace{3pt}\hrule\vspace{3pt}\footnotesize\textcolor{gray!80!black}{\textbf{Các mã dùng chung hình (HINH\_228):} \texttt{L09\_C4\_B2\_12, L09\_C4\_B2\_38, L09\_C4\_B2\_35}}
\end{tcolorbox}
\vspace{4pt}


\begin{tcolorbox}[
    breakable,
    colback=white,
    colframe=blue!60!black,
    boxrule=0.8pt,
    arc=3pt,
    title={\textbf{Câu \#314} \quad \textbf{MÃ CÂU: \texttt{[L09\_C4\_B2\_54]}} \quad | \quad \textbf{Nhóm: \texttt{HINH\_240}} \quad | \quad \textbf{MapID: \texttt{9H4V2-5}} \hfill \textcolor{yellow!90!white}{[Hình gốc .pyc]}}
]
\textbf{Dạng toán:} SA-M3. Tính chiều dài cây bị đổ gãy \hfill \textit{(Trả lời ngắn -- Vận dụng)}
\vspace{4pt}\hrule\vspace{6pt}
\begin{minipage}[t]{0.58\textwidth}
\vspace{0pt}
Một cây cao bị bão quật gãy, ngọn cây đổ xuống mặt đất. Gốc cây, điểm gãy và ngọn cây tạo thành một tam giác vuông. Đoạn cây gãy tạo với mặt đất một góc $\alpha = 15^{\circ} $ và chắn ngang lối đi một đoạn ${a=3,4} $ m (hình vẽ). Hãy tính chiều cao của cây trước khi bị gãy (với độ chính xác $0,05$).\\


\textbf{Đáp án điền số:} 	exttt{4,4}
\end{minipage}\hfill
\begin{minipage}[t]{0.39\textwidth}
\vspace{0pt}
\centering

\resizebox{\linewidth}{!}{%
\begin{tikzpicture}[line join=round, line cap=round, >=stealth, font=\footnotesize, scale=1]
    \path
    (5.65,-1.35)coordinate(A)
    (5.65,2)coordinate(B)
    (-4.4,-1.35)coordinate(C)
    (5.65,-1.6)coordinate(D)
    (-4.4,-1.6)coordinate(E);
    \draw[very thick](B)--(C);
    \draw pic[draw,angle radius=0.6cm,thick]{angle=A--C--B};
    \path(C)--++(60:0.5)node{$\alpha^\circ$};
    \draw[<->](D)--node[midway,below]{$a$ m}(E);
    \begin{scope}[shift={(0,0)},rotate=109]
        \foreach \w/\f in {0.3/30,0.2/50,0.1/70} {
            \fill [brown!\f!black] (0,-6) --(\w/2,4)--(\w/2,-6);}
        %Cành to
        \draw [line width=1.5pt,brown!50!black] plot[smooth, tension=.7] coordinates {(0.1,-1.77) (0.3,-1.65) (0.41,-1.55) (0.6,-1.48) (0.72,-1.45) (0.84,-1.34)};
        \draw  [line width=1.5pt,brown!50!black] plot[smooth, tension=.7] coordinates {(0.1,-1.8) (-0.13,-1.7) (-0.26,-1.62) (-0.4,-1.51) (-0.54,-1.46) (-0.67,-1.45) (-0.88,-1.52) (-1.05,-1.48)};
        \draw  [line width=1.5pt,brown!50!black] plot[smooth, tension=.7] coordinates {(0.18,-0.84) (0.26,-0.76) (0.4,-0.68) (0.55,-0.59) (0.75,-0.5)};
        \draw  [line width=1.5pt,brown!50!black] plot[smooth, tension=.7] coordinates {(0.02,-0.83) (-0.11,-0.71) (-0.23,-0.65) (-0.47,-0.56)};
        \draw  [line width=1.5pt,brown!50!black] plot[smooth, tension=.7] coordinates { (0.27,-0.11) (0.42,-0.01) (0.7,0.08)};
        \draw  [line width=1.5pt,brown!50!black] plot[smooth, tension=.7] coordinates {(0.05,-0.33)  (-0.37,0.02) (-0.58,0.04) (-0.69,0.15)};
        \draw  [line width=1.5pt,brown!50!black] plot[smooth, tension=.7] coordinates {(0.16,0.48) (0.3,0.56) (0.44,0.68) (0.59,0.73)};
        \draw  [line width=1.5pt,brown!50!black] plot[smooth, tension=.7] coordinates {(0.03,0.51) (-0.07,0.58) (-0.2,0.65) (-0.43,0.75)};
        \draw  [line width=1.5pt,brown!50!black] plot[smooth, tension=.7] coordinates {(0.19,1.18) (0.29,1.25) (0.39,1.3) (0.51,1.44)};
        \draw  [line width=1.5pt,brown!50!black] plot[smooth, tension=.7] coordinates {(0.05,1.21) (-0.02,1.31) (-0.15,1.44)};
        \draw  [line width=1.5pt,brown!50!black] plot[smooth, tension=.7] coordinates {(0.21,1.76) (0.31,1.86) (0.41,1.92) (0.55,2.01) (0.69,2.05)};
        \draw  [line width=1.5pt,brown!50!black] plot[smooth, tension=.7] coordinates {(0.08,1.8) (-0.03,1.89) (-0.14,1.96) (-0.34,2.06)};
        \draw  [line width=1.5pt,brown!50!black] plot[smooth, tension=.7] coordinates {(0.2,2.35) (0.3,2.4) (0.4,2.49) (0.49,2.53) (0.56,2.62)};
        \draw  [line width=1.5pt,brown!50!black] plot[smooth, tension=.7] coordinates {(0.1,2.41) (0.02,2.47) (-0.07,2.52) (-0.2,2.61)};
        \draw  [line width=1.5pt,brown!50!black] plot[smooth, tension=.7] coordinates {(0.2,3.06) (0.28,3.21) (0.39,3.25) (0.51,3.29)};
        %Tán lá
        \fill [green!40!black] (0.05,-1.41) -- (0.22,-1.33) -- (0.32,-1.33) -- (0.41,-1.38) -- (0.55,-1.35) -- (0.7,-1.34) -- (0.84,-1.39)  -- (0.87,-1.52) -- (1.02,-1.56) -- (1.17,-1.58) -- (1.34,-1.5) -- (1.4,-1.44) -- (1.54,-1.39) -- (1.62,-1.26) -- (1.69,-1.16) -- (1.68,-1.02) -- (1.54,-0.96) -- (1.39,-0.91) -- (1.2,-0.87) -- (0.97,-0.87) -- (0.81,-0.93) -- (0.67,-0.93) -- (0.54,-1.02) -- (0.42,-1.07) -- (0.24,-1.08) -- (0.1,-1.02) -- (0.01,-1.05) -- (-0.04,-0.99) -- (-0.15,-0.92) -- (-0.24,-0.94) -- (-0.34,-1) -- (-0.4,-1.01) -- (-0.46,-0.93) -- (-0.61,-0.91) -- (-0.78,-0.83) -- (-0.99,-0.83) -- (-1.2,-0.83) -- (-1.33,-0.89) -- (-1.51,-0.92) -- (-1.62,-0.92) -- (-1.73,-1.01) -- (-1.77,-1.17) -- (-1.66,-1.26) -- (-1.62,-1.34) -- (-1.63,-1.47) -- (-1.48,-1.51) -- (-1.23,-1.51) -- (-0.98,-1.5) -- (-0.93,-1.37) -- (-0.82,-1.26) -- (-0.62,-1.31) -- (-0.44,-1.35) -- (-0.24,-1.4) -- (-0.11,-1.46) -- (-0.01,-1.43);
        \fill [green!40!black] (0.12,3.52) -- (0.2,3.58) -- (0.29,3.59) -- (0.39,3.59) -- (0.46,3.63) -- (0.48,3.73) -- (0.48,3.83) -- (0.4,3.85) -- (0.36,3.9) -- (0.3,3.89) -- (0.22,3.9) -- (0.1,3.9) -- (-0.01,3.9) -- (-0.1,3.83) -- (-0.1,3.7) -- (-0.02,3.58) -- (0.06,3.55);
        \fill [green!40!black] (0.08,3.23) -- (0.19,3.23) -- (0.26,3.32) -- (0.36,3.38) -- (0.45,3.33) -- (0.5,3.22) -- (0.61,3.21) -- (0.67,3.27) -- (0.72,3.33) -- (0.67,3.42) -- (0.59,3.44) -- (0.52,3.48) -- (0.41,3.5) -- (0.3,3.49) -- (0.17,3.41) -- (0.05,3.41) -- (-0.03,3.43) -- (-0.19,3.46) -- (-0.25,3.47) -- (-0.26,3.35) -- (-0.2,3.28) -- (-0.06,3.26) -- (0.03,3.25);
        \fill [green!40!black]  (0.1,2.84) -- (0.26,2.84) -- (0.31,2.91) -- (0.42,2.82) -- (0.6,2.78) -- (0.66,2.85) -- (0.67,2.98) -- (0.59,3.07) -- (0.42,3.11) -- (0.3,3.13) -- (0.14,3.13) -- (-0.01,3.13) -- (-0.08,3.1) -- (-0.12,3.11) -- (-0.23,3.13) -- (-0.27,3.02) -- (-0.29,2.9) -- (-0.13,2.84) -- (-0.01,2.84);
        \fill [green!40!black] (0.13,2.5) -- (0.24,2.5) -- (0.33,2.54) -- (0.39,2.59) -- (0.51,2.6) -- (0.6,2.53) -- (0.71,2.49) -- (0.79,2.52) -- (0.86,2.6) -- (0.81,2.69) -- (0.67,2.71) -- (0.53,2.79) -- (0.37,2.75) -- (0.28,2.71) -- (0.19,2.69) -- (0.04,2.69) -- (-0.06,2.71) -- (-0.18,2.71) -- (-0.26,2.74) -- (-0.36,2.73) -- (-0.4,2.64) -- (-0.34,2.53) -- (-0.2,2.54) -- (-0.12,2.57) -- (-0.02,2.56) -- (0.07,2.51);
        \fill [green!40!black] (0.15,1.9) -- (0.27,1.94) -- (0.35,2) -- (0.47,2.06) -- (0.61,2.06) -- (0.62,1.96) -- (0.68,1.88) -- (0.8,1.87) -- (0.93,1.89) -- (1.03,1.93) -- (1.12,2) -- (1.14,2.11) -- (1.07,2.17) -- (0.93,2.22) -- (0.8,2.21) -- (0.73,2.24) -- (0.52,2.25) -- (0.38,2.26) -- (0.27,2.25) -- (0.2,2.18) -- (0.06,2.18) -- (-0.05,2.21) -- (-0.16,2.26) -- (-0.28,2.29) -- (-0.41,2.29) -- (-0.49,2.26) -- (-0.54,2.21) -- (-0.56,2.08) -- (-0.53,1.97) -- (-0.46,1.93) -- (-0.33,1.99) -- (-0.21,2.02) -- (-0.08,2.03) -- (0.03,2) -- (0.1,1.94);
        \fill [green!40!black] (0.1,1.32) -- (0.23,1.34) -- (0.32,1.39) -- (0.43,1.39) -- (0.49,1.29) -- (0.61,1.25) -- (0.77,1.23) -- (0.88,1.23) -- (1.05,1.29) -- (1.11,1.35) -- (1.18,1.4) -- (1.22,1.47) -- (1.19,1.6) -- (1.1,1.66) -- (0.91,1.71) -- (0.77,1.71) -- (0.66,1.66) -- (0.55,1.68) -- (0.44,1.64) -- (0.31,1.63) -- (0.15,1.64) -- (0.02,1.63) -- (-0.1,1.61) -- (-0.25,1.68) -- (-0.37,1.75) -- (-0.49,1.78) -- (-0.61,1.79) -- (-0.71,1.78) -- (-0.8,1.7) -- (-0.84,1.58) -- (-0.83,1.46) -- (-0.8,1.34) -- (-0.68,1.3) -- (-0.62,1.25) -- (-0.54,1.28) -- (-0.47,1.3) -- (-0.35,1.3) -- (-0.25,1.3) -- (-0.16,1.31) -- (-0.09,1.4) -- (0.03,1.35);
        \fill [green!40!black] (0.09,0.65) -- (0.28,0.66) -- (0.38,0.74) -- (0.48,0.77) -- (0.5,0.64) -- (0.53,0.55) -- (0.66,0.52) -- (0.83,0.54) -- (0.95,0.57) -- (1.05,0.59) -- (1.12,0.66) -- (1.2,0.67) -- (1.26,0.77) -- (1.33,0.91) -- (1.22,0.99) -- (1.09,0.98) -- (0.96,0.97) -- (0.89,1.01) -- (0.79,0.98) -- (0.65,1.01) -- (0.52,1.02) -- (0.36,1.02) -- (0.25,1.02) -- (0.13,0.98) -- (0.11,0.89) -- (0.05,0.83) -- (-0.03,0.86) -- (-0.01,0.97) -- (-0.03,1.05) -- (-0.15,1.05) -- (-0.24,1.02) -- (-0.33,1.04) -- (-0.5,1.03) -- (-0.65,1.05) -- (-0.79,1.05) -- (-0.9,1.04) -- (-0.97,0.91) -- (-1.05,0.75) -- (-1.02,0.64) -- (-0.86,0.62) -- (-0.72,0.65) -- (-0.61,0.63) -- (-0.46,0.63) -- (-0.35,0.69) -- (-0.21,0.76) -- (-0.1,0.72) -- (0.03,0.67);
        \fill [green!40!black] (0.1,0.02) -- (0.25,0.06) -- (0.34,0.07) -- (0.45,0.1) -- (0.55,0.08) -- (0.66,0.01) -- (0.75,-0.05) -- (0.86,-0.08) -- (0.98,-0.06) -- (1.07,0.01) -- (1.16,0.05) -- (1.24,0.11) -- (1.29,0.19) -- (1.33,0.24) -- (1.3,0.3) -- (1.31,0.36) -- (1.19,0.43) -- (0.97,0.48) -- (0.75,0.45) -- (0.58,0.42) -- (0.44,0.4) -- (0.29,0.33) -- (0.11,0.35) -- (0.03,0.32) -- (-0.07,0.41) -- (-0.15,0.45) -- (-0.27,0.43) -- (-0.43,0.44) -- (-0.65,0.5) -- (-0.86,0.5) -- (-1.05,0.45) -- (-1.18,0.37) -- (-1.26,0.24) -- (-1.25,0.11) -- (-1.17,-0.03) -- (-1.01,-0.06) -- (-0.81,-0.02) -- (-0.73,0.1) -- (-0.63,0.19) -- (-0.51,0.15) -- (-0.37,0.13) -- (-0.26,0.07) -- (-0.13,0.07);
        \fill [green!40!black] (0.06,-0.67)  -- (0.23,-0.64) -- (0.38,-0.6) -- (0.48,-0.66) -- (0.57,-0.75) -- (0.8,-0.8) -- (1.02,-0.8) -- (1.18,-0.76) -- (1.31,-0.71) -- (1.4,-0.6) -- (1.53,-0.6) -- (1.57,-0.5) -- (1.62,-0.38) -- (1.63,-0.25) -- (1.45,-0.23) -- (1.32,-0.23) -- (1.07,-0.17) -- (0.96,-0.19) -- (0.77,-0.19) -- (0.62,-0.15) -- (0.43,-0.15) -- (0.31,-0.16) -- (0.15,-0.23) -- (0.05,-0.27) -- (-0.06,-0.17)  -- (-0.18,-0.11)  -- (-0.3,-0.14) -- (-0.35,-0.19) -- (-0.52,-0.15) -- (-0.61,-0.17) -- (-0.75,-0.17) -- (-0.88,-0.15) -- (-1.05,-0.12) -- (-1.19,-0.14) -- (-1.3,-0.15) -- (-1.43,-0.2) -- (-1.48,-0.31) -- (-1.53,-0.46) -- (-1.46,-0.59) -- (-1.3,-0.68) -- (-1.1,-0.68) -- (-0.94,-0.59) -- (-0.85,-0.66) -- (-0.64,-0.66) -- (-0.52,-0.63) -- (-0.32,-0.56) -- (-0.2,-0.56);
    \end{scope}
    \foreach \w/\f in {0.3/30,0.2/50,0.1/70} {
        \fill [brown!\f!black] (5.65,-1.35) rectangle(5.65+\w/2,2);}
    \draw[very thick](6.5,-1.35)--(-4.5,-1.35);
    \fill[pattern=north east lines] (6.5,-1.35)rectangle(-4.5,-1.5);
\end{tikzpicture}
}
\end{minipage}

\vspace{3pt}\hrule\vspace{3pt}\footnotesize\textcolor{gray!80!black}{\textbf{Các mã dùng chung hình (HINH\_240):} \texttt{L09\_C4\_B2\_53}}
\end{tcolorbox}
\vspace{4pt}


\begin{tcolorbox}[
    breakable,
    colback=white,
    colframe=blue!60!black,
    boxrule=0.8pt,
    arc=3pt,
    title={\textbf{Câu \#315} \quad \textbf{MÃ CÂU: \texttt{[L09\_C4\_B2\_53]}} \quad | \quad \textbf{Nhóm: \texttt{HINH\_240}} \quad | \quad \textbf{MapID: \texttt{9H4V2-5}} \hfill \textcolor{yellow!90!white}{[Hình gốc .pyc]}}
]
\textbf{Dạng toán:} Tính chiều dài cây bị đổ gãy \hfill \textit{(Trả lời ngắn -- Vận dụng)}
\vspace{4pt}\hrule\vspace{6pt}
\begin{minipage}[t]{0.58\textwidth}
\vspace{0pt}
Một cây cao bị bão quật gãy, ngọn cây đổ xuống mặt đất. Gốc cây, điểm gãy và ngọn cây tạo thành một tam giác vuông. Đoạn cây gãy tạo với mặt đất một góc $\alpha = 20^{\circ} $ và chắn ngang lối đi một đoạn ${a=5,4} $ m (hình vẽ). Hãy tính chiều cao của cây trước khi bị gãy (với độ chính xác $0,05$).\\


\textbf{Đáp án điền số:} 	exttt{}
\end{minipage}\hfill
\begin{minipage}[t]{0.39\textwidth}
\vspace{0pt}
\centering

\resizebox{\linewidth}{!}{%
\begin{tikzpicture}[line join=round, line cap=round, >=stealth, font=\footnotesize, scale=1]
    \path
    (5.65,-1.35)coordinate(A)
    (5.65,2)coordinate(B)
    (-4.4,-1.35)coordinate(C)
    (5.65,-1.6)coordinate(D)
    (-4.4,-1.6)coordinate(E);
    \draw[very thick](B)--(C);
    \draw pic[draw,angle radius=0.6cm,thick]{angle=A--C--B};
    \path(C)--++(60:0.5)node{$\alpha^\circ$};
    \draw[<->](D)--node[midway,below]{$a$ m}(E);
    \begin{scope}[shift={(0,0)},rotate=109]
        \foreach \w/\f in {0.3/30,0.2/50,0.1/70} {
            \fill [brown!\f!black] (0,-6) --(\w/2,4)--(\w/2,-6);}
        %Cành to
        \draw [line width=1.5pt,brown!50!black] plot[smooth, tension=.7] coordinates {(0.1,-1.77) (0.3,-1.65) (0.41,-1.55) (0.6,-1.48) (0.72,-1.45) (0.84,-1.34)};
        \draw  [line width=1.5pt,brown!50!black] plot[smooth, tension=.7] coordinates {(0.1,-1.8) (-0.13,-1.7) (-0.26,-1.62) (-0.4,-1.51) (-0.54,-1.46) (-0.67,-1.45) (-0.88,-1.52) (-1.05,-1.48)};
        \draw  [line width=1.5pt,brown!50!black] plot[smooth, tension=.7] coordinates {(0.18,-0.84) (0.26,-0.76) (0.4,-0.68) (0.55,-0.59) (0.75,-0.5)};
        \draw  [line width=1.5pt,brown!50!black] plot[smooth, tension=.7] coordinates {(0.02,-0.83) (-0.11,-0.71) (-0.23,-0.65) (-0.47,-0.56)};
        \draw  [line width=1.5pt,brown!50!black] plot[smooth, tension=.7] coordinates { (0.27,-0.11) (0.42,-0.01) (0.7,0.08)};
        \draw  [line width=1.5pt,brown!50!black] plot[smooth, tension=.7] coordinates {(0.05,-0.33)  (-0.37,0.02) (-0.58,0.04) (-0.69,0.15)};
        \draw  [line width=1.5pt,brown!50!black] plot[smooth, tension=.7] coordinates {(0.16,0.48) (0.3,0.56) (0.44,0.68) (0.59,0.73)};
        \draw  [line width=1.5pt,brown!50!black] plot[smooth, tension=.7] coordinates {(0.03,0.51) (-0.07,0.58) (-0.2,0.65) (-0.43,0.75)};
        \draw  [line width=1.5pt,brown!50!black] plot[smooth, tension=.7] coordinates {(0.19,1.18) (0.29,1.25) (0.39,1.3) (0.51,1.44)};
        \draw  [line width=1.5pt,brown!50!black] plot[smooth, tension=.7] coordinates {(0.05,1.21) (-0.02,1.31) (-0.15,1.44)};
        \draw  [line width=1.5pt,brown!50!black] plot[smooth, tension=.7] coordinates {(0.21,1.76) (0.31,1.86) (0.41,1.92) (0.55,2.01) (0.69,2.05)};
        \draw  [line width=1.5pt,brown!50!black] plot[smooth, tension=.7] coordinates {(0.08,1.8) (-0.03,1.89) (-0.14,1.96) (-0.34,2.06)};
        \draw  [line width=1.5pt,brown!50!black] plot[smooth, tension=.7] coordinates {(0.2,2.35) (0.3,2.4) (0.4,2.49) (0.49,2.53) (0.56,2.62)};
        \draw  [line width=1.5pt,brown!50!black] plot[smooth, tension=.7] coordinates {(0.1,2.41) (0.02,2.47) (-0.07,2.52) (-0.2,2.61)};
        \draw  [line width=1.5pt,brown!50!black] plot[smooth, tension=.7] coordinates {(0.2,3.06) (0.28,3.21) (0.39,3.25) (0.51,3.29)};
        %Tán lá
        \fill [green!40!black] (0.05,-1.41) -- (0.22,-1.33) -- (0.32,-1.33) -- (0.41,-1.38) -- (0.55,-1.35) -- (0.7,-1.34) -- (0.84,-1.39)  -- (0.87,-1.52) -- (1.02,-1.56) -- (1.17,-1.58) -- (1.34,-1.5) -- (1.4,-1.44) -- (1.54,-1.39) -- (1.62,-1.26) -- (1.69,-1.16) -- (1.68,-1.02) -- (1.54,-0.96) -- (1.39,-0.91) -- (1.2,-0.87) -- (0.97,-0.87) -- (0.81,-0.93) -- (0.67,-0.93) -- (0.54,-1.02) -- (0.42,-1.07) -- (0.24,-1.08) -- (0.1,-1.02) -- (0.01,-1.05) -- (-0.04,-0.99) -- (-0.15,-0.92) -- (-0.24,-0.94) -- (-0.34,-1) -- (-0.4,-1.01) -- (-0.46,-0.93) -- (-0.61,-0.91) -- (-0.78,-0.83) -- (-0.99,-0.83) -- (-1.2,-0.83) -- (-1.33,-0.89) -- (-1.51,-0.92) -- (-1.62,-0.92) -- (-1.73,-1.01) -- (-1.77,-1.17) -- (-1.66,-1.26) -- (-1.62,-1.34) -- (-1.63,-1.47) -- (-1.48,-1.51) -- (-1.23,-1.51) -- (-0.98,-1.5) -- (-0.93,-1.37) -- (-0.82,-1.26) -- (-0.62,-1.31) -- (-0.44,-1.35) -- (-0.24,-1.4) -- (-0.11,-1.46) -- (-0.01,-1.43);
        \fill [green!40!black] (0.12,3.52) -- (0.2,3.58) -- (0.29,3.59) -- (0.39,3.59) -- (0.46,3.63) -- (0.48,3.73) -- (0.48,3.83) -- (0.4,3.85) -- (0.36,3.9) -- (0.3,3.89) -- (0.22,3.9) -- (0.1,3.9) -- (-0.01,3.9) -- (-0.1,3.83) -- (-0.1,3.7) -- (-0.02,3.58) -- (0.06,3.55);
        \fill [green!40!black] (0.08,3.23) -- (0.19,3.23) -- (0.26,3.32) -- (0.36,3.38) -- (0.45,3.33) -- (0.5,3.22) -- (0.61,3.21) -- (0.67,3.27) -- (0.72,3.33) -- (0.67,3.42) -- (0.59,3.44) -- (0.52,3.48) -- (0.41,3.5) -- (0.3,3.49) -- (0.17,3.41) -- (0.05,3.41) -- (-0.03,3.43) -- (-0.19,3.46) -- (-0.25,3.47) -- (-0.26,3.35) -- (-0.2,3.28) -- (-0.06,3.26) -- (0.03,3.25);
        \fill [green!40!black]  (0.1,2.84) -- (0.26,2.84) -- (0.31,2.91) -- (0.42,2.82) -- (0.6,2.78) -- (0.66,2.85) -- (0.67,2.98) -- (0.59,3.07) -- (0.42,3.11) -- (0.3,3.13) -- (0.14,3.13) -- (-0.01,3.13) -- (-0.08,3.1) -- (-0.12,3.11) -- (-0.23,3.13) -- (-0.27,3.02) -- (-0.29,2.9) -- (-0.13,2.84) -- (-0.01,2.84);
        \fill [green!40!black] (0.13,2.5) -- (0.24,2.5) -- (0.33,2.54) -- (0.39,2.59) -- (0.51,2.6) -- (0.6,2.53) -- (0.71,2.49) -- (0.79,2.52) -- (0.86,2.6) -- (0.81,2.69) -- (0.67,2.71) -- (0.53,2.79) -- (0.37,2.75) -- (0.28,2.71) -- (0.19,2.69) -- (0.04,2.69) -- (-0.06,2.71) -- (-0.18,2.71) -- (-0.26,2.74) -- (-0.36,2.73) -- (-0.4,2.64) -- (-0.34,2.53) -- (-0.2,2.54) -- (-0.12,2.57) -- (-0.02,2.56) -- (0.07,2.51);
        \fill [green!40!black] (0.15,1.9) -- (0.27,1.94) -- (0.35,2) -- (0.47,2.06) -- (0.61,2.06) -- (0.62,1.96) -- (0.68,1.88) -- (0.8,1.87) -- (0.93,1.89) -- (1.03,1.93) -- (1.12,2) -- (1.14,2.11) -- (1.07,2.17) -- (0.93,2.22) -- (0.8,2.21) -- (0.73,2.24) -- (0.52,2.25) -- (0.38,2.26) -- (0.27,2.25) -- (0.2,2.18) -- (0.06,2.18) -- (-0.05,2.21) -- (-0.16,2.26) -- (-0.28,2.29) -- (-0.41,2.29) -- (-0.49,2.26) -- (-0.54,2.21) -- (-0.56,2.08) -- (-0.53,1.97) -- (-0.46,1.93) -- (-0.33,1.99) -- (-0.21,2.02) -- (-0.08,2.03) -- (0.03,2) -- (0.1,1.94);
        \fill [green!40!black] (0.1,1.32) -- (0.23,1.34) -- (0.32,1.39) -- (0.43,1.39) -- (0.49,1.29) -- (0.61,1.25) -- (0.77,1.23) -- (0.88,1.23) -- (1.05,1.29) -- (1.11,1.35) -- (1.18,1.4) -- (1.22,1.47) -- (1.19,1.6) -- (1.1,1.66) -- (0.91,1.71) -- (0.77,1.71) -- (0.66,1.66) -- (0.55,1.68) -- (0.44,1.64) -- (0.31,1.63) -- (0.15,1.64) -- (0.02,1.63) -- (-0.1,1.61) -- (-0.25,1.68) -- (-0.37,1.75) -- (-0.49,1.78) -- (-0.61,1.79) -- (-0.71,1.78) -- (-0.8,1.7) -- (-0.84,1.58) -- (-0.83,1.46) -- (-0.8,1.34) -- (-0.68,1.3) -- (-0.62,1.25) -- (-0.54,1.28) -- (-0.47,1.3) -- (-0.35,1.3) -- (-0.25,1.3) -- (-0.16,1.31) -- (-0.09,1.4) -- (0.03,1.35);
        \fill [green!40!black] (0.09,0.65) -- (0.28,0.66) -- (0.38,0.74) -- (0.48,0.77) -- (0.5,0.64) -- (0.53,0.55) -- (0.66,0.52) -- (0.83,0.54) -- (0.95,0.57) -- (1.05,0.59) -- (1.12,0.66) -- (1.2,0.67) -- (1.26,0.77) -- (1.33,0.91) -- (1.22,0.99) -- (1.09,0.98) -- (0.96,0.97) -- (0.89,1.01) -- (0.79,0.98) -- (0.65,1.01) -- (0.52,1.02) -- (0.36,1.02) -- (0.25,1.02) -- (0.13,0.98) -- (0.11,0.89) -- (0.05,0.83) -- (-0.03,0.86) -- (-0.01,0.97) -- (-0.03,1.05) -- (-0.15,1.05) -- (-0.24,1.02) -- (-0.33,1.04) -- (-0.5,1.03) -- (-0.65,1.05) -- (-0.79,1.05) -- (-0.9,1.04) -- (-0.97,0.91) -- (-1.05,0.75) -- (-1.02,0.64) -- (-0.86,0.62) -- (-0.72,0.65) -- (-0.61,0.63) -- (-0.46,0.63) -- (-0.35,0.69) -- (-0.21,0.76) -- (-0.1,0.72) -- (0.03,0.67);
        \fill [green!40!black] (0.1,0.02) -- (0.25,0.06) -- (0.34,0.07) -- (0.45,0.1) -- (0.55,0.08) -- (0.66,0.01) -- (0.75,-0.05) -- (0.86,-0.08) -- (0.98,-0.06) -- (1.07,0.01) -- (1.16,0.05) -- (1.24,0.11) -- (1.29,0.19) -- (1.33,0.24) -- (1.3,0.3) -- (1.31,0.36) -- (1.19,0.43) -- (0.97,0.48) -- (0.75,0.45) -- (0.58,0.42) -- (0.44,0.4) -- (0.29,0.33) -- (0.11,0.35) -- (0.03,0.32) -- (-0.07,0.41) -- (-0.15,0.45) -- (-0.27,0.43) -- (-0.43,0.44) -- (-0.65,0.5) -- (-0.86,0.5) -- (-1.05,0.45) -- (-1.18,0.37) -- (-1.26,0.24) -- (-1.25,0.11) -- (-1.17,-0.03) -- (-1.01,-0.06) -- (-0.81,-0.02) -- (-0.73,0.1) -- (-0.63,0.19) -- (-0.51,0.15) -- (-0.37,0.13) -- (-0.26,0.07) -- (-0.13,0.07);
        \fill [green!40!black] (0.06,-0.67)  -- (0.23,-0.64) -- (0.38,-0.6) -- (0.48,-0.66) -- (0.57,-0.75) -- (0.8,-0.8) -- (1.02,-0.8) -- (1.18,-0.76) -- (1.31,-0.71) -- (1.4,-0.6) -- (1.53,-0.6) -- (1.57,-0.5) -- (1.62,-0.38) -- (1.63,-0.25) -- (1.45,-0.23) -- (1.32,-0.23) -- (1.07,-0.17) -- (0.96,-0.19) -- (0.77,-0.19) -- (0.62,-0.15) -- (0.43,-0.15) -- (0.31,-0.16) -- (0.15,-0.23) -- (0.05,-0.27) -- (-0.06,-0.17)  -- (-0.18,-0.11)  -- (-0.3,-0.14) -- (-0.35,-0.19) -- (-0.52,-0.15) -- (-0.61,-0.17) -- (-0.75,-0.17) -- (-0.88,-0.15) -- (-1.05,-0.12) -- (-1.19,-0.14) -- (-1.3,-0.15) -- (-1.43,-0.2) -- (-1.48,-0.31) -- (-1.53,-0.46) -- (-1.46,-0.59) -- (-1.3,-0.68) -- (-1.1,-0.68) -- (-0.94,-0.59) -- (-0.85,-0.66) -- (-0.64,-0.66) -- (-0.52,-0.63) -- (-0.32,-0.56) -- (-0.2,-0.56);
    \end{scope}
    \foreach \w/\f in {0.3/30,0.2/50,0.1/70} {
        \fill [brown!\f!black] (5.65,-1.35) rectangle(5.65+\w/2,2);}
    \draw[very thick](6.5,-1.35)--(-4.5,-1.35);
    \fill[pattern=north east lines] (6.5,-1.35)rectangle(-4.5,-1.5);
\end{tikzpicture}
}
\end{minipage}

\vspace{3pt}\hrule\vspace{3pt}\footnotesize\textcolor{gray!80!black}{\textbf{Các mã dùng chung hình (HINH\_240):} \texttt{L09\_C4\_B2\_54}}
\end{tcolorbox}
\vspace{4pt}


\begin{tcolorbox}[
    breakable,
    colback=white,
    colframe=blue!60!black,
    boxrule=0.8pt,
    arc=3pt,
    title={\textbf{Câu \#316} \quad \textbf{MÃ CÂU: \texttt{[L09\_C4\_B2\_22]}} \quad | \quad \textbf{Nhóm: \texttt{HINH\_234}} \quad | \quad \textbf{MapID: \texttt{9H4V2-5}} \hfill \textcolor{yellow!90!white}{[Hình gốc .pyc]}}
]
\textbf{Dạng toán:} Bài toán thực tế tính chiều cao của bức tường. \hfill \textit{(Tự luận -- Vận dụng)}
\vspace{4pt}\hrule\vspace{6pt}
\begin{minipage}[t]{0.58\textwidth}
\vspace{0pt}
Một cậu bé cao $a={1,2}m$ đứng cách một bức tường cao $b={1,6}m$ sao cho $\widehat {BAC}$ vuông tại ${A}$. Hỏi toà nhà cao bao nhiêu mét? (Làm tròn đến chữ số thập phân thứ nhất)
\end{minipage}\hfill
\begin{minipage}[t]{0.39\textwidth}
\vspace{0pt}
\centering

\resizebox{\linewidth}{!}{%
\begin{tikzpicture}[line join=round, line cap=round, >=stealth, font=\footnotesize, scale=2]
    \tikzset{
        man/.pic={
            \begin{scope}[shift={(-0.3,0)}]
                \fill [rounded corners=1.5] (0,0.45) -- (0,0.9) -- (0.4,0.9) -- (0.4,0.5) --
                (0.325,0.5) -- (0.325,0.8) -- (0.3,0.8) -- (0.3,0) -- (0.225,0) --
                (0.225,0.5) -- (0.175,0.5) -- (0.175,0) -- (0.1,0) -- (0.1,0.8) --
                (0.075,0.8) -- (0.075,0.5) -- cycle;
                \fill (0.2,1.05) circle (0.12);
                \coordinate (-head) at (0.2,1);
                \coordinate (-foot) at (0.2,0);
            \end{scope}
        }
    }
    % \clip (1.7,-.4) rectangle (6.1,5.45);
    \def\c{1.1} \def\h{1.6}
    \pgfmathsetmacro{\b}{((\h)^2)/(\c)+\c}
    \path
    (0,0) coordinate (O)
    (\h,0) coordinate (C)
    (0,\c) coordinate (A)
    (\h,\b)coordinate (B)
    (\h,\c)coordinate (H);
    \pic[scale=2,pink] at (O) {man};
    \draw[pattern=bricks,pattern color=brown] (C) rectangle ($(B)+(.5,0)$);
    \draw[thick] 
    (A)--(C)--(B)--cycle
    (A)--(H) node[midway,above]{$b\ \mathrm{m}$}
    (O)--(A) node[midway,above=.5,sloped]{$a\ \mathrm{m}$};
    \draw
    pic[draw,angle radius=3mm]{right angle=C--A--B}
    pic[draw,angle radius=3mm]{right angle=A--H--C};
    \foreach \d/\g in {A/100,B/75,C/-75,H/135,O/-75} {\fill (\d) circle(1pt) node[shift=(\g:10pt)]{$\d$};}
    \end{tikzpicture}
}
\end{minipage}

\vspace{3pt}\hrule\vspace{3pt}\footnotesize\textcolor{gray!80!black}{\textbf{Các mã dùng chung hình (HINH\_234):} \texttt{L09\_C4\_B2\_40}}
\end{tcolorbox}
\vspace{4pt}


\begin{tcolorbox}[
    breakable,
    colback=white,
    colframe=blue!60!black,
    boxrule=0.8pt,
    arc=3pt,
    title={\textbf{Câu \#317} \quad \textbf{MÃ CÂU: \texttt{[L09\_C4\_B2\_23]}} \quad | \quad \textbf{Nhóm: \texttt{HINH\_235}} \quad | \quad \textbf{MapID: \texttt{9H4C2-5}} \hfill \textcolor{yellow!90!white}{[Hình gốc .pyc]}}
]
\textbf{Dạng toán:} Bài toán thực tế tính chiều cao của lâu đài. \hfill \textit{(Tự luận -- Vận dụng cao)}
\vspace{4pt}\hrule\vspace{6pt}
\begin{minipage}[t]{0.58\textwidth}
\vspace{0pt}
Để đo chiều cao của một toà lâu đài người ta đặt giác kế thẳng đứng tại vị trí ${D}$. Quay ống ngắm của giác kế sao cho nhìn thấy đỉnh ${A}$ của toà lâu đài dưới góc nhọn $\alpha=28^{\circ}$.  Sau đó đặt giác kế thẳng đứng tại điểm ${C}$, biết $CD=21m$ thì nhìn thấy đỉnh ${A}$ dưới góc $\beta={17^{\circ}}$. Biết chiều cao của giác kế là ${1,2}m$. Hãy tính chiều cao của lâu đài.
\end{minipage}\hfill
\begin{minipage}[t]{0.39\textwidth}
\vspace{0pt}
\centering

\resizebox{\linewidth}{!}{%
\begin{tikzpicture}[scale=1]
    
    % Định nghĩa các điểm
    \path
    (1,5) coordinate (A)
    (1,0) coordinate (B)
    (8,0) coordinate (C)
    (6,1) coordinate (E)
    (8,1) coordinate (F)
    (6,0) coordinate (D)
    (1,1) coordinate (H);
    
    % Vẽ tháp
    \fill[pink!50] (0,0) rectangle (2,4);
    \fill[red!30] (0,4) -- (1,5) -- (2,4) -- cycle;
    
    % Cửa sổ
    \fill[blue!30] (0.5,1) rectangle (1.5,2.5);
    
    % Điểm A trên đỉnh tháp
    \node[circle, fill=red, inner sep=2pt] at (1,5) {};
    \node[above] at (1,5) {A};
    
    % Điểm B dưới chân tháp
    \node[circle, fill=blue, inner sep=2pt] at (1,0) {};
    \node[below] at (1,0) {B};
    
    % Điểm C dưới chân tháp
    \node[circle, fill=blue, inner sep=2pt] at (8,0) {};
    \node[below] at (8,0) {C};
    
    % Điểm D dưới chân tháp
    \node[circle, fill=blue, inner sep=2pt] at (6,0) {};
    \node[below] at (6,0) {D};
    
    % Điểm F dưới chân tháp
    \node[circle, fill=blue, inner sep=2pt] at (8,1) {};
    \node[below left] at (8,1) {F};
    
    % Điểm E dưới chân tháp
    \node[circle, fill=blue, inner sep=2pt] at (6,1) {};
    \node[below left] at (6,1) {E};
    
    % Điểm H dưới chân tháp
    \node[circle, fill=blue, inner sep=2pt] at (1,1) {};
    \node[below left] at (1,1) {H};
    
    % Vẽ đường viền của tháp
    \draw[black, thick] (0,0) rectangle (2,4);
    \draw[black, thick] (0,4) -- (1,5) -- (2,4) -- cycle;
    \draw[black, thick] (0.5,1) rectangle (1.5,2.5);
    
    % Nối các điểm bằng nét liền
    \draw[black, thick] (A) -- (F) -- (C) -- (B) -- cycle;
    \draw[black, thick] (A) -- (E) -- (D) -- (C) -- (B) -- cycle;
    \draw[black, thick] (A) -- (E) -- (F) -- (H) -- cycle;
    \path (E) pic[draw=black,fill=gray!25, angle radius = 7] {angle =A--E--H}
    node[shift={(-0.5,2.9pt)},scale=0.8](GA){$\alpha$};
    \path (F) pic[draw=black,fill=gray!25, angle radius = 7] {angle =A--F--H}
    node[shift={(-0.6,2.9pt)},scale=0.8](GA){$\beta$};
    \end{tikzpicture}
}
\end{minipage}

\vspace{3pt}\hrule\vspace{3pt}\footnotesize\textcolor{gray!80!black}{\textbf{Các mã dùng chung hình (HINH\_235):} \texttt{L09\_C4\_B2\_41}}
\end{tcolorbox}
\vspace{4pt}


\begin{tcolorbox}[
    breakable,
    colback=white,
    colframe=blue!60!black,
    boxrule=0.8pt,
    arc=3pt,
    title={\textbf{Câu \#318} \quad \textbf{MÃ CÂU: \texttt{[L09\_C4\_B2\_24]}} \quad | \quad \textbf{Nhóm: \texttt{HINH\_225}} \quad | \quad \textbf{MapID: \texttt{9H4V2-5}} \hfill \textcolor{yellow!90!white}{[Hình gốc .pyc]}}
]
\textbf{Dạng toán:} Bài toán thực tế hải đăng và con tàu \hfill \textit{(Tự luận -- Vận dụng)}
\vspace{4pt}\hrule\vspace{6pt}
\begin{minipage}[t]{0.58\textwidth}
\vspace{0pt}
Một người đứng trên tháp tại ${A}$ của ngọn hải đăng cao $70$ m quan sát hai lần một con tàu đang hướng về ngọn hải đăng. Lần thứ nhất người đó nhìn thấy  tàu tại ${D}$ với góc hạ $\alpha= 14^{\circ}$, lần thứ hai người đó nhìn thấy tàu tại ${C}$ với góc hạ $\beta=29^{\circ}$.  Hỏi con tàu đã đi được bao nhiêu mét giữa hai lần quan sát.( Kết quả làm tròn đến hàng đơn vị).
\end{minipage}\hfill
\begin{minipage}[t]{0.39\textwidth}
\vspace{0pt}
\centering

\resizebox{\linewidth}{!}{%
\begin{tikzpicture}
    % Sea (slightly wavy line)
    \fill[blue!30] (0,0) rectangle (10,-2);
    \draw[blue, thick] (0,-0.1) .. controls (1,-0.3) and (2,0.1) .. (3,-0.2)
    .. controls (4,-0.5) and (5,-0.1) .. (6,-0.3)
    .. controls (7,-0.1) and (8,-0.3) .. (9,-0.2)
    .. controls (10,-0.5) and (11,0.1) .. (12,-0.2);
    
    % Enhanced lighthouse (cartoonish style with more details)
    
    % Lighthouse base (stone)
    \fill[gray!50] (0.7,-0.5) rectangle (1.3,0.3);
    \draw[thick] (0.7,-0.5) -- (0.7,0.3) -- (1.3,0.3) -- (1.3,-0.5) -- cycle;
    
    % Lighthouse body (stripes for a more lively look)
    \fill[white] (0.7,0.3) rectangle (1.3,1);
    \fill[red!80] (0.7,1) rectangle (1.3,1.5);
    \fill[white] (0.7,1.5) rectangle (1.3,2);
    \fill[red!80] (0.7,2) rectangle (1.3,2.5);
    \fill[white] (0.7,2.5) rectangle (1.3,3);
    
    \draw[thick] (0.7,0.3) -- (0.7,3) -- (1.3,3) -- (1.3,0.3) -- cycle;
    
    % Windows on the lighthouse
    \fill[black!70] (0.9,1.2) rectangle (1.1,1.4); % first window
    \fill[black!70] (0.9,2.2) rectangle (1.1,2.4); % second window
    
    % Lighthouse top (lantern room)
    \fill[gray!70] (0.6,3) rectangle (1.4,3.3);
    \draw[thick] (0.6,3) -- (1.4,3);
    \draw[thick] (0.6,3.3) -- (1.4,3.3);
    
    % Lighthouse roof (dome-shaped top)
    \fill[red!80] (0.65,3.3) -- (1.35,3.3) -- (1,3.8) -- cycle;
    
    % Lighthouse light
    \fill[yellow] (1,3.15) circle (0.1);
    
    % Lighthouse light beams (enhanced beams with glow effect)
    \draw[yellow!80, thick] (1,3.15) -- (0,4);
    \draw[yellow!80, thick] (1,3.15) -- (2,4);
    
    % Lighthouse door
    \fill[brown!80] (0.9,-0.5) rectangle (1.1,-0.2);
    \draw[thick] (0.9,-0.5) -- (0.9,-0.2) -- (1.1,-0.2) -- (1.1,-0.5) -- cycle;
    
    % Boat's hull (with rounded edges for cartoon style)
    \fill[brown!70] (4,-0.5) .. controls (4.2,-1.5) and (5.8,-1.5) .. (6,-0.5) -- cycle;
    \draw[thick] (4,-0.5) .. controls (4.2,-1.5) and (5.8,-1.5) .. (6,-0.5) -- cycle;
    
    % Boat's deck
    \fill[brown!50] (4.2,-0.5) rectangle (5.8,0);
    
    % Mast of the boat
    \draw[thick] (5,0) -- (5,1.5);
    
    % Sail (cartoonish curved sail)
    \fill[white] (5,0.2) .. controls (5.8,0.7) and (5.8,1.3) .. (5,1.4) -- cycle;
    \draw[thick] (5,0.2) .. controls (5.8,0.7) and (5.8,1.3) .. (5,1.4) -- cycle;
    
    % Little flag on the top of the mast
    \fill[red] (5,1.5) -- (5.4,1.6) -- (5,1.7) -- cycle;
    
    % Add some cartoonish waves
    \foreach \x in {3.5, 5.5, 7.5} {
        \draw[blue!60,thick] (\x,-0.5) .. controls (\x+0.3,-0.6) and (\x+0.7,-0.4) .. (\x+1,-0.5);
    }
    
    % Second boat (smaller and farther away)
    \fill[brown!70] (7,-0.5) .. controls (7.1,-0.9) and (7.9,-0.9) .. (8,-0.5) -- cycle; % hull
    \draw[thick] (7,-0.5) .. controls (7.1,-0.9) and (7.9,-0.9) .. (8,-0.5) -- cycle;
    
    \draw[thick] (7.5,-0.5) -- (7.5,0.2); % mast
    \fill[white] (7.5,-0.1) .. controls (8,0) and (8,0.1) .. (7.5,0.2) -- cycle; % sail
    \draw[thick] (7.5,-0.1) .. controls (8,0) and (8,0.1) .. (7.5,0.2) -- cycle;
    
    % Flag for the second boat
    \fill[red] (7.5,0.2) -- (7.7,0.3) -- (7.5,0.4) -- cycle;
    
    % Marking points A (top of lighthouse), B (bottom of lighthouse), C (first boat), D (second boat)
    \fill[black] (1,3.5) circle (0.05) node[above] {A};  % Point A (top of lighthouse)
    \fill[black] (1,-0.5) circle (0.05) node[below] {B}; % Point B (bottom of lighthouse)
    \fill[black] (5,-0.5) circle (0.05) node[below] {C}; % Point C (position of boat 1)
    \fill[black] (7.5,-0.5) circle (0.05) node[below] {D}; % Point D (position of boat 2)
    
    % Draw lines connecting A-B, A-C, A-D
    \draw[thick] (1,3.5) -- (1,-0.5);  % Line AB
    \draw[thick] (1,3.5) -- (5,-0.5);  % Line AC
    \draw[thick] (1,3.5) -- (7.5,-0.5); % Line AD
    
    % Draw line B-C-D
    \draw[thick] (1,-0.5) -- (5,-0.5) -- (7.5,-0.5);  % Line BCD
    
    % Add a line x parallel to BC through A
    \draw[dashed] (1,3.5) -- (6,3.5) node[right] {$x$}; % Dashed line x, parallel to BC
    
    % Add angle alpha for xAD
    \draw[thick] (1.6,2.9) arc[start angle=-60, end angle= -10, radius=0.8];
    \node at (1.6,3.2) {$\beta$};  % Label for angle alpha
    
    % Add angle beta for xAC
    \draw[thick] (1.9,2.9) arc[start angle=-25, end angle= -13, radius=3.1];
    \node at (2.3,3.3) {$\alpha$};  % Label for angle beta
    
    % Add a smiling sun in the sky
    \fill[yellow] (9,3.5) circle (0.5); % sun body
    \draw (9,3.7) arc[start angle=45, end angle=135, radius=0.2]; % smiling face
    \fill[black] (8.8,3.6) circle (0.05); % left eye
    \fill[black] (9.2,3.6) circle (0.05); % right eye
    \foreach \angle in {0,45,...,315} {
        \draw[thick, yellow] (9,3.5) -- ++(\angle:0.8);
    }
    
\end{tikzpicture}
}
\end{minipage}

\vspace{3pt}\hrule\vspace{3pt}\footnotesize\textcolor{gray!80!black}{\textbf{Các mã dùng chung hình (HINH\_225):} \texttt{L09\_C4\_B2\_25, L09\_C4\_B2\_42}}
\end{tcolorbox}
\vspace{4pt}


\begin{tcolorbox}[
    breakable,
    colback=white,
    colframe=blue!60!black,
    boxrule=0.8pt,
    arc=3pt,
    title={\textbf{Câu \#319} \quad \textbf{MÃ CÂU: \texttt{[L09\_C4\_B2\_26]}} \quad | \quad \textbf{Nhóm: \texttt{HINH\_241}} \quad | \quad \textbf{MapID: \texttt{9H4V2-5}} \hfill \textcolor{yellow!90!white}{[Hình gốc .pyc]}}
]
\textbf{Dạng toán:} Bài toán đo khoảng cách khinh khí cầu đến mặt đất \hfill \textit{(Tự luận -- Vận dụng)}
\vspace{4pt}\hrule\vspace{6pt}
\begin{minipage}[t]{0.58\textwidth}
\vspace{0pt}
Một cậu bé cao ${1,4}m$ đứng ngắm khinh khí cầu ở một vị trí cách hình chiếu của nó trên một đất tức điểm ${E}$ một khoảng là $220$m. Cậu bé đặt ống nhòm hợp với phương ngang một góc $52^{\circ}$ thì nhìn thấy khinh khí cầu. Hỏi khi đó khinh khí cầu cách mặt đất bao nhiêu mét?(Kết quả làm tròn đến chữ số thập phân thứ nhất)
\end{minipage}\hfill
\begin{minipage}[t]{0.39\textwidth}
\vspace{0pt}
\centering

\resizebox{\linewidth}{!}{%
\begin{tikzpicture}[line join = round, line cap=round,>=stealth,font=\footnotesize,scale=1]
\path
(0,0) coordinate (D)
($(D)+(0,1)$) coordinate (B)
($(D)+(4,0)$) coordinate (E)
($(E)+(0,1)$) coordinate (A)
($(E)+(0,4)$) coordinate (C)
($(C)+(-2,0)$) coordinate (x)
;
\draw (D)--(E)--(C)--(B)--cycle (B)--(A);
\foreach \p/\r in {A/0,B/90,C/0,D/-90,E/-90}
\fill (\p) circle (1.2pt) node[shift={(\r:3mm)}]{$\p$};
\draw[pattern=north west lines] ($(D)+(-1,0)$) rectangle ($(E)+(1,-0.5)$);
\path
(D)--(E) node[above,pos=0.5]{$220$ m} ;
\draw[<->] ($(D)+(-0.75,0)$)--($(B)+(-0.75,0)$) node[left,pos=0.5]{$1,4$ m};
\draw pic[draw,angle radius=3mm] {right angle = B--A--C};
\draw pic[draw,angle radius=8mm] {angle = A--B--C};
\draw ($(B)+(1.25,0)$) node[above]{$52^\circ$};
\draw[dashed] (C)--($(x)+(-2,0)$)node[above]{$x$};
\begin{scope}[line width=3pt]
\foreach \yA/\xB/\yB in {0.6/-0.3/0.05, 0.6/-0.15/0.05} {
\draw (-0.25, 0.45) -- (-0.25, \yA) -- (\xB, \yB);
}
\foreach \yA/\xB/\yB/\xC/\yC in {
0.9/-0.1/0.8/0.0/0.9,
0.9/-0.4/0.8/-0.3/0.7
} {
\draw (-0.25, 0.45) -- (-0.25, \yA) -- (\xB, \yB) -- (\xC, \yC);
}
\draw (-0.25, 0.45) -- (-0.25, 1.0)circle(2pt);
\end{scope}
\fill [scale=.1,black,yshift=40cm,xshift=40cm]
(-1,0) rectangle (1,1)
(-1,2).. controls +(135:1) and +(180:4) .. (0,7)
.. controls +(180:2) and +(135:2).. (-.5,2)
.. controls +(120:2) and +(180:1) .. (0,6.99)
.. controls +(0:1) and +(60:2) .. (.5,2)
.. controls +(45:2) and +(0:2) .. (0,7)
.. controls +(0:4) and +(45:1).. (1,2) ;
\end{tikzpicture}
}
\end{minipage}

\end{tcolorbox}
\vspace{4pt}


\begin{tcolorbox}[
    breakable,
    colback=white,
    colframe=blue!60!black,
    boxrule=0.8pt,
    arc=3pt,
    title={\textbf{Câu \#320} \quad \textbf{MÃ CÂU: \texttt{[L09\_C4\_B2\_28]}} \quad | \quad \textbf{Nhóm: \texttt{HINH\_242}} \quad | \quad \textbf{MapID: \texttt{9H4V2-5}} \hfill \textcolor{yellow!90!white}{[Hình gốc .pyc]}}
]
\textbf{Dạng toán:} Bài toán tính chiều cao của quạt gió \hfill \textit{(Tự luận -- Vận dụng)}
\vspace{4pt}\hrule\vspace{6pt}
\begin{minipage}[t]{0.58\textwidth}
\vspace{0pt}
Một cậu bé cao ${1,5}m$ đứng cách thân một cái quạt gió $31$ m và nhìn thấy tâm của cánh quạt với góc nâng $56^{\circ}$ (Hình vẽ bên). Tính khoảng cách từ tâm của cánh quạt đến mặt đất.(Kết quả làm tròn đến chữ số thập phân thứ nhất)
\end{minipage}\hfill
\begin{minipage}[t]{0.39\textwidth}
\vspace{0pt}
\centering

\resizebox{\linewidth}{!}{%
\begin{tikzpicture}[smooth,font=\footnotesize,scale=1]
    \coordinate (M) at ($(4,4.16)+(-60:2cm)$);
    \coordinate (N) at ($(4,4.16)+(60:2cm)$);
    \coordinate (K) at ($(4,4.16)+(180:2cm)$);
    \fill[color=red](-0.05,1)--(0.05,1)--(0.05,0)--(-0.05,0)--cycle;
    \fill[color=blue](3.95,0)--(4.05,0)--(4.05,4)--(3.95,4)--cycle;
    \fill[ball color=red](0,1.16) circle (0.16);
    \fill[ball color=blue](4,4.16) circle (0.16);
    \draw[dashed,blue] (0.16,1.16)--(3.95,1.16) (0.16,1.16)--(3.84,4.16);
    \draw[very thick,blue] (4,4.16)--(M) (4,4.16)--(N) (4,4.16)--(K);
    \path
    (3.95,1.16)--(3.95,0) node[below,midway,sloped,scale=1]{$1,5m$ }
    (0.16,1.16)--(3.95,1.16) node[below,midway,sloped,scale=1]{$31m$ };
    \node[right] at (0.9,1.45) {$56^\circ$};
    \begin{scope}
        \clip (3.95,1.16)--(0.16,1.16)--(3.84,4.16);
        \draw[double] (0.16,1.16) circle(0.7cm);
    \end{scope}
    
    \end{tikzpicture}
}
\end{minipage}

\end{tcolorbox}
\vspace{4pt}


\begin{tcolorbox}[
    breakable,
    colback=white,
    colframe=blue!60!black,
    boxrule=0.8pt,
    arc=3pt,
    title={\textbf{Câu \#321} \quad \textbf{MÃ CÂU: \texttt{[L09\_C4\_B2\_30]}} \quad | \quad \textbf{Nhóm: \texttt{HINH\_224}} \quad | \quad \textbf{MapID: \texttt{9H4V2-5}} \hfill \textcolor{yellow!90!white}{[Hình gốc .pyc]}}
]
\textbf{Dạng toán:} Bài toán thực tế tính khoảng cách cù lao đến bờ sông. \hfill \textit{(Tự luận -- Vận dụng)}
\vspace{4pt}\hrule\vspace{6pt}
\begin{minipage}[t]{0.58\textwidth}
\vspace{0pt}
Hai ngư dân đứng ở một bên bờ sông cách nhau $AB={240}m$ cùng nhìn thấy một cù lao trên sông với các góc nhìn lần lượt là  $\beta=58^{\circ}$ và $\alpha=25^{\circ}$. Tính khoảng cách từ cù lao ${C}$ đến bờ sông ${AB}$ (làm tròn đến hàng đơn vị).
\end{minipage}\hfill
\begin{minipage}[t]{0.39\textwidth}
\vspace{0pt}
\centering

\resizebox{\linewidth}{!}{%
\begin{tikzpicture}[scale=0.7, font=\footnotesize, line join = round, line cap = round,>=stealth]
    \clip (-4.39,3.2) rectangle (4,-4);
    \draw[pattern=north east lines,opacity=0.3] plot[smooth] coordinates{(-4.39,1.43)(-2.43,1.46) (-1.48,1.67)(1.43,2.12)(2.01,2.09)(4,1.64) (4.1,-1.06)(3.02,-0.77)(1.85,-1.38)(0.95,-1.48)(0.13,-1.72)(-1.48,-1.67)(-2.46,-1.51)(-3.57,-0.98)(-4.37,-1.08)};
    \draw[opacity=3] plot[smooth] coordinates{(-4.39,1.43)(-2.43,1.46) (-1.48,1.67)(1.43,2.12)(2.01,2.09)(4,1.64) (4.1,-1.06)(3.02,-0.77)(1.85,-1.38)(0.95,-1.48)(0.13,-1.72)(-1.48,-1.67)(-2.46,-1.51)(-3.57,-0.98)(-4.37,-1.08)};
    \draw[fill=white] plot[smooth  cycle] coordinates{(-2.75,-0.08)(-1.83,0.32) (-0.03,0.42) (1,0) (0.24,-0.48)(-0.58,-0.79)(-1.38,-0.77)(-1.91,-0.71)};
    \draw[fill=black!70]  plot[smooth  cycle] coordinates{(-2.14,0.58)(-2.09,0.24)(-2.33,-0.13) (-1.19,-0.13) (-1.38,-0.13) (-1.69,0.21)(-1.75,0.53)};
    \draw[fill=blue!30]  plot[smooth  cycle] coordinates{(-1.75,0.53)(-1.46,0.79) (-1.01,0.58) (-1.08,1.06) (-0.64,1.08)(-0.93,1.59)(-0.53,1.85)(-0.93,2.2)(-1.38,2.7)(-2.04,3.18)(-2.49,2.91)(-3.07,2.91)(-3.2,2.22)(-3.73,1.96)(-3.1,1.3)(-3.31,1.01)(-2.83,1.01)(-2.99,0.69)(-2.14,0.58)};
    \tkzDefPoints{-1.38/-0.13/C,-2/-3/B,3/-2/A}
    \tkzDrawPoints[fill=black](A,B,C)
    \tkzDrawPolygon(A,B,C)
    \tkzDefMidPoint(A,B) 
    \tkzLabelPoints[below](A,B)
    \tkzLabelPoints[above right](C)
    \tkzMarkAngles[size=1cm,arc=l,mark=|](C,A,B)
    \tkzMarkAngles[size=1cm,arc=l,mark=||](A,B,C)
    \tkzLabelAngles[left=1cm,pos=0.5,rotate=-20](B,A,C){\footnotesize$\alpha$}
    \tkzLabelAngles[right,pos=1](A,B,C){\footnotesize$\beta$}
 \end{tikzpicture}
}
\end{minipage}

\vspace{3pt}\hrule\vspace{3pt}\footnotesize\textcolor{gray!80!black}{\textbf{Các mã dùng chung hình (HINH\_224):} \texttt{L09\_C4\_B2\_14, L09\_C4\_B2\_39, L09\_C4\_B2\_45}}
\end{tcolorbox}
\vspace{4pt}


\begin{tcolorbox}[
    breakable,
    colback=white,
    colframe=blue!60!black,
    boxrule=0.8pt,
    arc=3pt,
    title={\textbf{Câu \#322} \quad \textbf{MÃ CÂU: \texttt{[L09\_C4\_B2\_31]}} \quad | \quad \textbf{Nhóm: \texttt{HINH\_223}} \quad | \quad \textbf{MapID: \texttt{9H4V2-5}} \hfill \textcolor{yellow!90!white}{[Hình gốc .pyc]}}
]
\textbf{Dạng toán:} Bài toán thực tế tính vận tốc của máy bay khi bay lên \hfill \textit{(Tự luận -- Vận dụng)}
\vspace{4pt}\hrule\vspace{6pt}
\begin{minipage}[t]{0.58\textwidth}
\vspace{0pt}
Một máy bay vừa phi lên theo đường bay hợp với mặt đất một góc $\widehat{ATH} = {30}^{\circ}$. Sau $1,2$ phút máy bay ở độ cao $a={5}$km so với mặt đất. Hỏi vận tốc của máy bay khi bay lên?
\end{minipage}\hfill
\begin{minipage}[t]{0.39\textwidth}
\vspace{0pt}
\centering

\resizebox{\linewidth}{!}{%
\begin{tikzpicture}[declare function={r=4;rh=r*sin(30);}]
        \path (0,0) coordinate (T)
        --(30:r) coordinate (A)
        node[above,shift={(0.2,-0.05)}]{
            \resizebox{1cm}{!}{\begin{tikzpicture}[line join=round]
                    \draw[line width=2pt,line cap=butt,rounded corners=2pt] 
                    (0.3,0) arc(270:85:0.35 and 0.4)--(0.8,1.5)..controls +(50:0.2) and +(180:0.2)..(1.4,1.8)--(2.5,1.6)..controls +(0:0.35) and +(140:0.15)..(3,1.45)--(3.8,1.45)--(4.3,2.15)--(4.85,2.15)
                    --(4.3,0.95)--(3.1,0.7)..controls +(-120:0.65) and +(0:0.1)..(2.3,0)--cycle;
                    \draw[line width=2pt,line cap=butt] 
                    (4.7,1.25) arc (57:-150:0.45)
                    (0.1,-0.4) ..controls +(-30:0.15) and +(180:0.1)..
                    (0.5,-0.475)--(2.1,-0.475)..controls +(0:0.2) and +(-100:0.1)..(2.7,-0.375)
                    (-0.05,2.45)--(3.4,2.45);
                    \path[fill] (1.6,1.9)--++(0.18,0)--+(0,0.35) arc (0:180:0.09 and 0.13)--cycle;
                    \path[fill] (1.8,0.8){[rounded corners=3pt]--++(0.5,0)--++(0,0.7)--+(-0.5,0)--cycle}
                    {[rounded corners=2pt](0.8,0.8)--(1.55,0.8)--++(0,0.7)--+(-0.2,0)}
                    ..controls +(200:0.325) and +(85:0.15)..(0.8,0.8)--cycle;
                \end{tikzpicture}
            }
        }--([turn]-120:rh) coordinate (H)--cycle;
        \draw 
            (A)--(H) node[pos=0.5,right]{$a$ km}--(T)--(A)
          
            %(O) pic[draw,angle radius = 6] {right angle = B--O--C}
        ;
        \foreach \t/\g in {T/-90,H/-90,A/-30}{
            \draw[fill=white] (\t) circle (1pt) node[shift={(\g:7pt)},font=\scriptsize]{$ \t $};
        }
    \end{tikzpicture}
}
\end{minipage}

\vspace{3pt}\hrule\vspace{3pt}\footnotesize\textcolor{gray!80!black}{\textbf{Các mã dùng chung hình (HINH\_223):} \texttt{L09\_C4\_B2\_08, L09\_C4\_B2\_46}}
\end{tcolorbox}
\vspace{4pt}


\begin{tcolorbox}[
    breakable,
    colback=white,
    colframe=blue!60!black,
    boxrule=0.8pt,
    arc=3pt,
    title={\textbf{Câu \#323} \quad \textbf{MÃ CÂU: \texttt{[L09\_C4\_B2\_32]}} \quad | \quad \textbf{Nhóm: \texttt{HINH\_243}} \quad | \quad \textbf{MapID: \texttt{9H4V2-5}} \hfill \textcolor{yellow!90!white}{[Hình gốc .pyc]}}
]
\textbf{Dạng toán:} Bài toán thực tế tính chiều cao toà nhà \hfill \textit{(Tự luận -- Vận dụng)}
\vspace{4pt}\hrule\vspace{6pt}
\begin{minipage}[t]{0.58\textwidth}
\vspace{0pt}
Khi tia sáng mặt trời tạo với mặt đất một góc ${64}^{\circ}$ thì bóng của một tòa nhà trên mặt đất có độ dài $b={35}$m. Tính chiều cao của tòa nhà (kết quả làm tròn đến hàng đơn vị).
\end{minipage}\hfill
\begin{minipage}[t]{0.39\textwidth}
\vspace{0pt}
\centering

\resizebox{\linewidth}{!}{%
\begin{tikzpicture}[scale=1,font=\footnotesize,line cap=round,line join=round]
    % Tòa nhà ----------------------------------------------
    \path
    (3,0) coordinate (A)
    (3,5) coordinate (B)
    (8,0) coordinate (C)
    ($(A)!.5!(C)$) node[below] {$35$ m}
    ($(C)+(162:1)$) node {$64^{\circ}$}
    ;
    \foreach \i in {0,0.5,1,1.5,2,2.5}{\foreach \j in {0,0.8,1.6,2.4,3.2,4.0}
        \draw[fill=yellow!50!white] (\i,\j) rectangle +(0.5,0.8);
    } 
    % Ban công 
    \foreach \i in {0.1,0.6,1.1,1.6,2.1,2.6}{\foreach \j in {0,0.8,1.6,2.4,3.2,4.0}
        \draw[fill=orange!60!white] (\i,\j) rectangle +(0.3,0.6);
    }
    \foreach \j in {0,0.8,1.6,2.4,3.2,4.0,4.8}
        \draw[fill=cyan!30!white] (0,\j) rectangle +(3,0.3);
    
    \draw (A)--(B)--(C)--cycle;    %       \tkzMarkRightAngles(B,A,C)
    %       \tkzMarkAngles[size=.7](B,C,A)
    
    \foreach \i/\p in {A/-90,B/90,C/0}  
        \fill[black] (\i) circle (1.5pt)  % Điểm đánh dấu màu đen
        node[shift={(\p:3mm)},black] {$\i$}; % Chữ cái màu đen
\end{tikzpicture}
}
\end{minipage}

\end{tcolorbox}
\vspace{4pt}


\begin{tcolorbox}[
    breakable,
    colback=white,
    colframe=blue!60!black,
    boxrule=0.8pt,
    arc=3pt,
    title={\textbf{Câu \#324} \quad \textbf{MÃ CÂU: \texttt{[L09\_C4\_B2\_33]}} \quad | \quad \textbf{Nhóm: \texttt{HINH\_221}} \quad | \quad \textbf{MapID: \texttt{9H4V2-5}} \hfill \textcolor{yellow!90!white}{[Hình gốc .pyc]}}
]
\textbf{Dạng toán:} Bài toán thực tế tính chiều cao của cái cây \hfill \textit{(Tự luận -- Vận dụng)}
\vspace{4pt}\hrule\vspace{6pt}
\begin{minipage}[t]{0.58\textwidth}
\vspace{0pt}
Khi tia sáng mặt trời tạo với mặt đất một góc $\alpha={49}^{\circ}$ thì bóng của một cái cây trên mặt đất có độ dài $x={6}$m. Tính chiều cao của cái cây (kết quả làm tròn đến hàng đơn vị).
\end{minipage}\hfill
\begin{minipage}[t]{0.39\textwidth}
\vspace{0pt}
\centering

\resizebox{\linewidth}{!}{%
\begin{tikzpicture}[scale=1, font=\footnotesize, line cap=round, line join=round]

    % Vẽ ông mặt trời
    \fill[yellow] (-3,7) circle (0.5cm); % Di chuyển mặt trời ra xa một chút

    % Vẽ các tia nắng sinh động bao quanh ông mặt trời (từng tia riêng lẻ)
    \pgfmathsetmacro{\radius}{0.8} % Khoảng cách từ tâm mặt trời đến điểm bắt đầu của tia nắng
    
    % Vẽ từng tia nắng riêng lẻ
    \pgfmathsetmacro{\lengthA}{0.5 + 0.4 * sin(0 * 8.8)}
    \draw[orange, thick] (-3,7) ++(0:\radius) -- ++(0:\lengthA);

    \pgfmathsetmacro{\lengthB}{0.5 + 0.4 * sin(15 * 8.8)}
    \draw[orange, thick] (-3,7) ++(15:\radius) -- ++(15:\lengthB);

    \pgfmathsetmacro{\lengthC}{0.5 + 0.4 * sin(30 * 8.8)}
    \draw[orange, thick] (-3,7) ++(30:\radius) -- ++(30:\lengthC);

    \pgfmathsetmacro{\lengthD}{0.5 + 0.4 * sin(45 * 8.8)}
    \draw[orange, thick] (-3,7) ++(45:\radius) -- ++(45:\lengthD);

    \pgfmathsetmacro{\lengthE}{0.5 + 0.4 * sin(60 * 8.8)}
    \draw[orange, thick] (-3,7) ++(60:\radius) -- ++(60:\lengthE);

    \pgfmathsetmacro{\lengthF}{0.5 + 0.4 * sin(75 * 8.8)}
    \draw[orange, thick] (-3,7) ++(75:\radius) -- ++(75:\lengthF);

    \pgfmathsetmacro{\lengthG}{0.5 + 0.4 * sin(90 * 8.8)}
    \draw[orange, thick] (-3,7) ++(90:\radius) -- ++(90:\lengthG);

    \pgfmathsetmacro{\lengthH}{0.5 + 0.4 * sin(105 * 8.8)}
    \draw[orange, thick] (-3,7) ++(105:\radius) -- ++(105:\lengthH);

    \pgfmathsetmacro{\lengthI}{0.5 + 0.4 * sin(120 * 8.8)}
    \draw[orange, thick] (-3,7) ++(120:\radius) -- ++(120:\lengthI);

    \pgfmathsetmacro{\lengthJ}{0.5 + 0.4 * sin(135 * 8.8)}
    \draw[orange, thick] (-3,7) ++(135:\radius) -- ++(135:\lengthJ);

    \pgfmathsetmacro{\lengthK}{0.5 + 0.4 * sin(150 * 8.8)}
    \draw[orange, thick] (-3,7) ++(150:\radius) -- ++(150:\lengthK);

    \pgfmathsetmacro{\lengthL}{0.5 + 0.4 * sin(165 * 8.8)}
    \draw[orange, thick] (-3,7) ++(165:\radius) -- ++(165:\lengthL);

    \pgfmathsetmacro{\lengthM}{0.5 + 0.4 * sin(180 * 8.8)}
    \draw[orange, thick] (-3,7) ++(180:\radius) -- ++(180:\lengthM);

    \pgfmathsetmacro{\lengthN}{0.5 + 0.4 * sin(195 * 8.8)}
    \draw[orange, thick] (-3,7) ++(195:\radius) -- ++(195:\lengthN);

    \pgfmathsetmacro{\lengthO}{0.5 + 0.4 * sin(210 * 8.8)}
    \draw[orange, thick] (-3,7) ++(210:\radius) -- ++(210:\lengthO);

    \pgfmathsetmacro{\lengthP}{0.5 + 0.4 * sin(225 * 8.8)}
    \draw[orange, thick] (-3,7) ++(225:\radius) -- ++(225:\lengthP);

    \pgfmathsetmacro{\lengthQ}{0.5 + 0.4 * sin(240 * 8.8)}
    \draw[orange, thick] (-3,7) ++(240:\radius) -- ++(240:\lengthQ);

    \pgfmathsetmacro{\lengthR}{0.5 + 0.4 * sin(255 * 8.8)}
    \draw[orange, thick] (-3,7) ++(255:\radius) -- ++(255:\lengthR);

    \pgfmathsetmacro{\lengthS}{0.5 + 0.4 * sin(270 * 8.8)}
    \draw[orange, thick] (-3,7) ++(270:\radius) -- ++(270:\lengthS);

    \pgfmathsetmacro{\lengthT}{0.5 + 0.4 * sin(285 * 8.8)}
    \draw[orange, thick] (-3,7) ++(285:\radius) -- ++(285:\lengthT);

    \pgfmathsetmacro{\lengthU}{0.5 + 0.4 * sin(300 * 8.8)}
    \draw[orange, thick] (-3,7) ++(300:\radius) -- ++(300:\lengthU);

    \pgfmathsetmacro{\lengthV}{0.5 + 0.4 * sin(315 * 8.8)}
    \draw[orange, thick] (-3,7) ++(315:\radius) -- ++(315:\lengthV);

    \pgfmathsetmacro{\lengthW}{0.5 + 0.4 * sin(330 * 8.8)}
    \draw[orange, thick] (-3,7) ++(330:\radius) -- ++(330:\lengthW);

    \pgfmathsetmacro{\lengthX}{0.5 + 0.4 * sin(345 * 8.8)}
    \draw[orange, thick] (-3,7) ++(345:\radius) -- ++(345:\lengthX);

    % Vẽ cái cây cao hơn
    \fill[brown] (0,0) rectangle (0.5,4); % Tăng chiều cao của cái cây
    
    % Vẽ tán cây bằng các hình tròn chồng lên nhau, thẳng hàng với thân cây và xếp cao dần
    \fill[green] (0.25,3.5) circle (0.8cm); % Vòng tròn lớn dưới cùng
    \fill[green] (0.25,4.5) circle (1.2cm); % Vòng tròn trung bình ở giữa
    \fill[green] (0.25,4) circle (1cm); % Vòng tròn nhỏ ở trên cùng
    \fill[brown] (0.2,0) ellipse (0.8 and 0.24); % Elip với kích thước tuỳ chỉnh

    % Định nghĩa các điểm A, B, C
    \path
    (0.25,0) coordinate (A) % Gốc cây
    (0.25,5.6) coordinate (B) % Ngọn cây
    (4,0) coordinate (C); % Điểm trên mặt đất, đã di chuyển ra xa
    
    % Vẽ tam giác ABC
    \draw (A) -- (B) -- (C) -- cycle; % Nối các điểm A, B, C

    % Gắn nhãn các đoạn thẳng và góc
    \draw ($(A)!.5!(C)$) node[below] {$x$}; % Nhãn đoạn thẳng AC
    \node at ($(C)+(155:1)$) {$\alpha$}; % Ký hiệu góc BCA

    % Đánh dấu các điểm A, B, C
    \foreach \i/\p in {A/-90,B/90,C/0}  
    \fill[black] (\i) circle (1.5pt)  
    node[shift={(\p:3mm)},black] {$\i$}; % Chữ cái màu đen

\end{tikzpicture}
}
\end{minipage}

\vspace{3pt}\hrule\vspace{3pt}\footnotesize\textcolor{gray!80!black}{\textbf{Các mã dùng chung hình (HINH\_221):} \texttt{L09\_C4\_B1\_19, L09\_C4\_B2\_11, L09\_C4\_B2\_48}}
\end{tcolorbox}
\vspace{4pt}


\begin{tcolorbox}[
    breakable,
    colback=white,
    colframe=blue!60!black,
    boxrule=0.8pt,
    arc=3pt,
    title={\textbf{Câu \#325} \quad \textbf{MÃ CÂU: \texttt{[L09\_C4\_B2\_55]}} \quad | \quad \textbf{Nhóm: \texttt{HINH\_244}} \quad | \quad \textbf{MapID: \texttt{9H4V2-5}} \hfill \textcolor{yellow!90!white}{[Hình gốc .pyc]}}
]
\textbf{Dạng toán:} Tính góc trượt của thang \hfill \textit{(Tự luận -- Vận dụng)}
\vspace{4pt}\hrule\vspace{6pt}
\begin{minipage}[t]{0.58\textwidth}
\vspace{0pt}
Một chiếc thang ${AC}$ được dựng vào một bức tường thẳng đứng (hình vẽ). 

 – Ban đầu khoảng cách từ chân thang đến tường là $BC=5,4 m$ và góc tạo bởi thang và phương nằm ngang là $\widehat{ACB}=32^{\circ}$.

 – Sau đó, đầu ${A}$ của thang bị trượt xuống $ 40cm$ đến vị trí ${D}$. Khi đó, góc $\widehat{DEB}$ tạo bởi thang và phương nằm ngang bằng bao nhiêu (Kết quả số đo góc làm tròn đến phút)?
\end{minipage}\hfill
\begin{minipage}[t]{0.39\textwidth}
\vspace{0pt}
\centering

\resizebox{\linewidth}{!}{%
\begin{tikzpicture}[line cap=round,line join=round]
                \path 
                (0,0) coordinate (B)
                ($(B)+(90:5)$) coordinate (A)
                ($(B)+(0:3)$) coordinate (C)
                ($(C)+(0:1)$) coordinate (E)
                ($(A)+(-90:.5)$) coordinate (D)
                ($(E)+(0:1.4)$) coordinate (N);
                \draw[dashed] (E)--(D);
                \draw[line width=3pt,brown] (A)--(C);
                \draw ($(A)+(90:.5)$) rectangle ($(B)+(180:1.2)$);
                \fill[pattern=bricks] ($(A)+(90:.5)$) rectangle ($(B)+(180:1.2)$);
                \draw ($(B)+(180:2.2)$) rectangle ($(N)+(-90:.8)$);
                \fill[yellow!30] ($(B)+(180:2.2)$) rectangle ($(N)+(-90:.8)$);
                \foreach \t/\g in {A/60,B/-100,C/-70,E/30}{\fill[black] (\t)circle(.5pt) node[shift={(\g:3mm)}] {$\t$};}
                \node[shift={(-50:6mm)}] at (D) {$D$};
            \end{tikzpicture}
}
\end{minipage}

\end{tcolorbox}
\vspace{4pt}


\begin{tcolorbox}[
    breakable,
    colback=teal!3!white,
    colframe=teal!75!black,
    boxrule=0.8pt,
    arc=3pt,
    title={\textbf{Câu \#326} \quad \textbf{MÃ CÂU: \texttt{[L09\_C4\_B2\_56]}} \quad | \quad \textbf{Nhóm: \texttt{HINH\_245}} \quad | \quad \textbf{MapID: \texttt{9H4V2-5}} \hfill \textbf{\textcolor{green!40!black}{[ĐÃ THAY MẪU TIKZ MỚI]}}}
]
\textbf{Dạng toán:} Bài toán đo chiều cao của nhà 2 khi đứng ở nhà 1 \hfill \textit{(Tự luận -- Vận dụng)}
\vspace{4pt}\hrule\vspace{6pt}
\begin{minipage}[t]{0.58\textwidth}
\vspace{0pt}
Cho hai tòa nhà 1 và tòa nhà 2 như hình vẽ bên. Trên nóc tòa nhà 2 có một cột ăng-ten thẳng cao $1$ m. Từ vị trí quan sát ${A}$ (trên nóc tòa nhà 1) cao $9$ m so với mặt đất có thể nhìn thấy đỉnh ${B}$ và chân ${C}$ của cột ăng-ten lần lượt dưới góc $ \widehat{BAD}= 47^{\circ}$ và $ \widehat{CAD}= 39^{\circ}$ so với phương nằm ngang. Tính chiều cao ${CH}$ của tòa nhà 2 (làm tròn kết quả đến hàng phần mười).
\end{minipage}\hfill
\begin{minipage}[t]{0.39\textwidth}
\vspace{0pt}
\centering
\definecolor{skyb}{RGB}{102, 178, 240}
\definecolor{buildblue}{RGB}{30, 115, 230}
\definecolor{gridblue}{RGB}{185, 225, 255}
\definecolor{roofblue}{RGB}{15, 75, 170}
\definecolor{grassgreen}{RGB}{125, 185, 45}
\resizebox{\linewidth}{!}{%
\begin{tikzpicture}[
    >=stealth, 
    line join=round, 
    line cap=round, 
    every node/.style={font=\normalsize},
    declare function={
        w1 = 2.4;
        h1 = 3.6;
        d_space = 4.2;
        w2 = 2.6;
        ang_C = 39;
        ang_B = 47;
        dist_AD = d_space + 0.5*w2; 
        yC = h1 + dist_AD * tan(ang_C);
        yB = h1 + dist_AD * tan(ang_B);
        h_ant = yB - yC;
        r_pt = 0.4pt;
    },scale=0.55
]
    \draw[line width=0.8pt] (-0.8,0) -- (w1 + d_space + w2 + 0.8, 0);
    \draw[fill=grassgreen!85, draw=grassgreen!60!black, rounded corners=5pt, line width=0.6pt] 
        (-0.3,-0.05) rectangle (w1 + 0.3, 0.25);
    \draw[fill=grassgreen!85, draw=grassgreen!60!black, rounded corners=5pt, line width=0.6pt] 
        (w1 + d_space - 0.3, -0.05) rectangle (w1 + d_space + w2 + 0.3, 0.25);
    \begin{scope}
        \draw[fill=buildblue!70, draw=white, line width=0.8pt] 
            (0.1, 0.2) rectangle ($(w1-0.1, 0.9)$);
        \draw[fill=buildblue, draw=white, line width=1pt] 
            (0.2, 0.9) rectangle ($(w1-0.2, {h1-0.4})$);
        \foreach \i in {1,...,10} {
            \draw[gridblue!60, line width=0.3pt] 
                (0.2, {0.9 + \i*(h1-1.3)/11}) -- ($(w1-0.2, {0.9 + \i*(h1-1.3)/11})$);
        }
        \foreach \j in {1,...,4} {
            \draw[gridblue!60, line width=0.3pt] 
                ({0.2 + \j*(w1-0.4)/5}, 0.9) -- ({0.2 + \j*(w1-0.4)/5}, {h1-0.4});
        }
        \draw[fill=roofblue, draw=white, line width=0.8pt] 
            (0.4, {h1-0.4}) rectangle ($(w1-0.4, {h1-0.15})$);
        \draw[fill=roofblue!90, draw=white, line width=0.8pt] 
            (0.6, {h1-0.15}) rectangle ($(w1-0.6, h1)$);
        \draw[line width=1pt, roofblue!50!black] (0.3, h1) -- ($(w1-0.3, h1)$);
    \end{scope}
    \begin{scope}[shift={(w1 + d_space, 0)}]
        \draw[fill=white, draw=gray!40, line width=0.6pt] 
            (0.1, 0.2) rectangle ($(w2-0.1, 1.3)$);
        \foreach \k in {0,1,2} {
            \draw[fill=buildblue!90, draw=white, line width=0.6pt] 
                ({0.3 + \k*0.75}, 0.3) rectangle ({0.85 + \k*0.75}, 1.15);
        }
        \draw[fill=buildblue, draw=white, line width=1.2pt] 
            (0.2, 1.3) rectangle ($(w2-0.2, {yC-2.2})$);
        \foreach \i in {1,...,6} {
            \draw[gridblue!60, line width=0.3pt] 
                (0.2, {1.3 + \i*(yC-3.5)/7}) -- ($(w2-0.2, {1.3 + \i*(yC-3.5)/7})$);
        }
        \foreach \j in {1,...,5} {
            \draw[gridblue!60, line width=0.3pt] 
                ({0.2 + \j*(w2-0.4)/6}, 1.3) -- ({0.2 + \j*(w2-0.4)/6}, {yC-2.2});
        }
        \draw[fill=buildblue, draw=white, line width=1.2pt] 
            (0.45, {yC-2.2}) rectangle ($(w2-0.45, {yC-0.5})$);
        \foreach \i in {1,...,4} {
            \draw[gridblue!60, line width=0.3pt] 
                (0.45, {yC-2.2 + \i*1.7/5}) -- ($(w2-0.45, {yC-2.2 + \i*1.7/5})$);
        }
        \foreach \j in {1,...,3} {
            \draw[gridblue!60, line width=0.3pt] 
                ({0.45 + \j*(w2-0.9)/4}, {yC-2.2}) -- ({0.45 + \j*(w2-0.9)/4}, {yC-0.5});
        }
        \draw[fill=buildblue!90, draw=white, line width=1pt] 
            (0.75, {yC-0.5}) rectangle ($(w2-0.75, yC)$);
        \draw[line width=1.2pt, roofblue] (0.15, {yC-2.2}) -- ($(w2-0.15, {yC-2.2})$);
        \draw[line width=1.2pt, roofblue] (0.4, {yC-0.5}) -- ($(w2-0.4, {yC-0.5})$);
        \draw[line width=1.2pt, roofblue] (0.7, yC) -- ($(w2-0.7, yC)$);
        \draw[line width=1.3pt, roofblue!90!black] (0.5*w2, yC) -- (0.5*w2, yB);
        \foreach \h in {0.25, 0.5, 0.75} {
            \draw[line width=0.7pt, roofblue!90!black] 
                ({0.5*w2 - 0.35*(1-\h*0.5)}, {yC + \h*h_ant}) -- ({0.5*w2 + 0.35*(1-\h*0.5)}, {yC + \h*h_ant});
        }
        \fill[roofblue!90!black] (0.5*w2, yB) circle (1.5pt);
    \end{scope}
    \coordinate (A) at (w1, h1);
    \coordinate (D) at (w1 + d_space + 0.5*w2, h1);
    \coordinate (C) at (w1 + d_space + 0.5*w2, yC);
    \coordinate (B) at (w1 + d_space + 0.5*w2, yB);
    \coordinate (H1) at (w1, 0);
    \coordinate (H2) at (w1 + d_space + 0.5*w2, 0);
    \draw[dashed, line width=0.6pt, black!70] (A) -- (D);
    \draw[dashed, line width=0.6pt, black!70] (B) -- (H2);
    \draw[line width=0.7pt, black!85] (A) -- (C);
    \draw[line width=0.7pt, black!85] (A) -- (B);
    \draw (D) ++(-0.25, 0) -- ++(0, 0.25) -- ++(0.25, 0);
    \draw[line width=0.5pt] (A) ++(0:1.25) arc (0:ang_C:1.25);
    \node at ($(A) + ({0.5*ang_C}:1.75)$) {$ 39^\circ $};
    \draw[line width=0.5pt] (A) ++(0:2.35) arc (0:ang_B:2.35);
    \draw[->, >=stealth, line width=0.4pt, black!80] 
        ($(A) + ({ang_B + 9}:2.85)$) to[bend left=20] ($(A) + ({0.5*(ang_C + ang_B)}:2.25)$);
    \node[above left, inner sep=1pt] at ($(A) + ({ang_B + 9}:2.85)$) {$ 47^\circ $};
    \draw[<->, >=stealth, line width=0.6pt, black!80] (0.65, 0) -- (0.65, h1) 
        node[midway, fill=buildblue!90, text=white, inner sep=1.5pt, rounded corners=1pt] {$ 9\,\text{m} $};
    \draw[line width=0.4pt, gray!80] (C) -- ++(0.65, 0);
    \draw[line width=0.4pt, gray!80] (B) -- ++(0.65, 0);
    \draw[<->, >=stealth, line width=0.6pt, black!80] ($(C)+(0.5,0)$) -- ($(B)+(0.5,0)$)
        node[midway, right, inner sep=2pt] {$ 1\,\text{m} $};
    \foreach \pt in {A, B, C, D} {
        \fill[black] (\pt) circle (1.2pt);
    }
    \node[above left, inner sep=2pt] at (A) {$ A $};
    \node[above left, inner sep=2pt] at (B) {$ B $};
    \node[above left, inner sep=2pt] at (C) {$ C $};
    \node[below right, inner sep=2pt, fill=buildblue!85, text=white, rounded corners=1pt] at ($(D)+(0.05,-0.05)$) {$ D $};
    \node[below, font=\bfseries\small] at (0.5*w1, -0.1) {Tòa nhà 1};
    \node[below, font=\bfseries\small] at (w1 + d_space + 0.5*w2, -0.1) {Tòa nhà 2};
\end{tikzpicture}
}
\end{minipage}

\end{tcolorbox}
\vspace{4pt}


\begin{tcolorbox}[
    breakable,
    colback=white,
    colframe=blue!60!black,
    boxrule=0.8pt,
    arc=3pt,
    title={\textbf{Câu \#327} \quad \textbf{MÃ CÂU: \texttt{[L09\_C4\_B2\_34]}} \quad | \quad \textbf{Nhóm: \texttt{HINH\_246}} \quad | \quad \textbf{MapID: \texttt{9H4C2-5}} \hfill \textcolor{yellow!90!white}{[Hình gốc .pyc]}}
]
\textbf{Dạng toán:} Bài toán thực tế tính thời gian của máy bay khi hạ cánh \hfill \textit{(Tự luận -- Vận dụng cao)}
\vspace{4pt}\hrule\vspace{6pt}
\begin{minipage}[t]{0.58\textwidth}
\vspace{0pt}
Quãng đường bay khi hạ cánh của máy bay, hay còn gọi là quãng đường tiếp cận cuối cùng (landing distance), là khoảng cách mà máy bay bay từ lúc bắt đầu tiếp cận cuối cùng cho đến khi chạm đất. Đây là giai đoạn máy bay đã hạ thấp gần mặt đất và đang hướng thẳng tới đường băng với góc hạ cánh cố định (thường là $\widehat{ATH}={3}^{\circ}$), chuẩn bị tiếp xúc với đường băng để hoàn tất quá trình hạ cánh. Một máy bay ở giai đoạn tiếp cận cuối cùng cách mặt đất $a={500}$m, tính thời gian hạ cánh ở giai đoạn này biết vận tốc hạ cánh là $280$km/h.
\end{minipage}\hfill
\begin{minipage}[t]{0.39\textwidth}
\vspace{0pt}
\centering

\resizebox{\linewidth}{!}{%
\begin{tikzpicture}[declare function={r=4;rh=r*sin(30);}]
        \path (0,0) coordinate (T)
        --(30:r) coordinate (A)
        node[above,shift={(0.2,-0.05)}]{
            \resizebox{1cm}{!}{\begin{tikzpicture}[line join=round]
                    \draw[line width=2pt,line cap=butt,rounded corners=2pt] 
                    (0.3,0) arc(270:85:0.35 and 0.4)--(0.8,1.5)..controls +(50:0.2) and +(180:0.2)..(1.4,1.8)--(2.5,1.6)..controls +(0:0.35) and +(140:0.15)..(3,1.45)--(3.8,1.45)--(4.3,2.15)--(4.85,2.15)
                    --(4.3,0.95)--(3.1,0.7)..controls +(-120:0.65) and +(0:0.1)..(2.3,0)--cycle;
                    \draw[line width=2pt,line cap=butt] 
                    (4.7,1.25) arc (57:-150:0.45)
                    (0.1,-0.4) ..controls +(-30:0.15) and +(180:0.1)..
                    (0.5,-0.475)--(2.1,-0.475)..controls +(0:0.2) and +(-100:0.1)..(2.7,-0.375)
                    (-0.05,2.45)--(3.4,2.45);
                    \path[fill] (1.6,1.9)--++(0.18,0)--+(0,0.35) arc (0:180:0.09 and 0.13)--cycle;
                    \path[fill] (1.8,0.8){[rounded corners=3pt]--++(0.5,0)--++(0,0.7)--+(-0.5,0)--cycle}
                    {[rounded corners=2pt](0.8,0.8)--(1.55,0.8)--++(0,0.7)--+(-0.2,0)}
                    ..controls +(200:0.325) and +(85:0.15)..(0.8,0.8)--cycle;
                \end{tikzpicture}
            }
        }--([turn]-120:rh) coordinate (H)--cycle;
        %\draw[-]  (-r/2,0) --(3/2*r,0)node[below];
        \draw 
            (A)--(H) node[pos=0.5,right]{$500$ m}--(T)--(A)
        ;
        \foreach \t/\g in {T/-90,H/-90,A/-30}{
            \draw[fill=white] (\t) circle (1pt) node[shift={(\g:7pt)},font=\scriptsize]{$ \t $};
        }
    \end{tikzpicture}
}
\end{minipage}

\end{tcolorbox}
\vspace{4pt}


\begin{tcolorbox}[
    breakable,
    colback=white,
    colframe=blue!60!black,
    boxrule=0.8pt,
    arc=3pt,
    title={\textbf{Câu \#328} \quad \textbf{MÃ CÂU: \texttt{[L09\_C4\_B2\_35]}} \quad | \quad \textbf{Nhóm: \texttt{HINH\_228}} \quad | \quad \textbf{MapID: \texttt{9H4C2-5}} \hfill \textcolor{yellow!90!white}{[Hình gốc .pyc]}}
]
\textbf{Dạng toán:} Bài toán thực tế tính chiều cao của hải đăng \hfill \textit{(Tự luận -- Vận dụng cao)}
\vspace{4pt}\hrule\vspace{6pt}
\begin{minipage}[t]{0.58\textwidth}
\vspace{0pt}
Ở hai địa điểm ${A}$ và ${B}$ trên mặt đất cách nhau $25$m, người ta đặt ống nhòm nhìn thấy vị trí ${D}$ cao nhất của ngọn hải đăng, góc đặt ống nhóm tại ${B}$ hợp với mặt đất là $50^{\circ}$, góc đặt ống nhòm tại vị trí ${A}$ hợp với mặt đất là $61^{\circ}$. Tính chiều cao của hải đăng (làm tròn đến hàng đơn vị).
\end{minipage}\hfill
\begin{minipage}[t]{0.39\textwidth}
\vspace{0pt}
\centering

\resizebox{\linewidth}{!}{%
\begin{tikzpicture}[scale=.5,declare function={h=13.4;a=h/tan(63);b=h/tan(48);}]
        \fill[brown] (0,0) -- (-5,0) 
        to[out=20,in=180] (-4.3,.15)
        to[out=40,in=180] (-3.05,.65)
        to[out=80,in=200] (-2.8,.85)
        to[out=70,in=160] (-2.3,1.25)
        to[out=50,in=190] (-2,1.45)
        --(-1.6,1.45)--(-1.35,3.2)--(-1.3,2.15)
        to[out=50,in=210] (0,3.2)
        to[out=28,in=222] (1.25,4.05)
        --(1.5,1.5)--(1.7,1.5)
        to[out=-40,in=180] (2,1.33)
        to[out=-10,in=110] (2.45,1)
        to[out=-40,in=160] (3,.8)
        to[out=-40,in=150] (3.2,.6)
        to[out=10,in=130] (4.4,.1)
        to[out=0,in=160] (5,0)
        --cycle;
        \fill[brown] (-.85,3.6) rectangle (-.53,4.1);
        \fill[brown] (-1.3,4.2)
        to[out=50,in=208] (0,5.15)
        to[out=25,in=225] (1,5.75)
        to[out=50,in=275] (1.15,6.3)
        --(1.08,8)
        to[out=-130,in=37] (-1,6.6)
        to[out=-135,in=88] (-1.25,5.5)
        --cycle;
        \fill[brown] (.25,8.45) rectangle (.55,9);
        \fill[brown] (-1.1,8.3)
        to[out=55,in=245] (1.05,9.85)--(1.2,9.85)
        to[out=10,in=263] (1.45,10.1)
        to[out=100,in=-10] (.96,10.35)--(.93,11.1)--(1.05,11.15)
        to[out=155,in=20] (-1.07,11.2)--(-1.07,10.3)
        to[out=-175,in=135] (-1.4,10)
        to[out=-55,in=190] (-1.08,9.85)--cycle;
        \draw[brown,line width= 2pt] (.9,11.2)--(.9,11.7)
        (.25,11.2)--(.25,11.9)
        (-.35,11.2)--(-.35,11.9)
        (-1.03,11.2)--(-1.03,11.7)
        to[out=20,in=160] (0.94,11.7)
        (.55,11.8)--(.54,12.7)
        (-.6,11.8)--(-.59,12.7)
        (-.04,11.9)--(-.04,12.7);
        \draw[brown,line width=4pt] (.6,11.2)--(.6,11.8)
        (-.75,11.2)--(-.75,11.8);
        \fill[brown] (-.04,12.15) circle (7pt)
        (.6,12.7) arc (0:180:.65 and .7) -- cycle
        (-.04,13.4) circle (5pt);
        \path 
        (0,0) coordinate (C)
        (a,0) coordinate (A)
        (b,0) coordinate (B)
        (0,h) coordinate (D);
        \foreach \i/\g in {C/-90, A/-90, B/-90, D/90} \fill (\i)circle(.1) node[shift={(\g:.25)}] {$\i$};
        \draw (C)--(B)node[pos=.8, below]{$25m$}--(D)--cycle (A)--(D)
        (A)+(-1,0) arc (180:117:1)node[pos=.7, left]{$\alpha=61^\circ$};
        \draw[double] (B)+(-1,0) arc (180:132:1)node[pos=.7, left]{$\beta=50^\circ$};
        \draw[dashed] (-5,h)--(D);
        \draw[<->,-stealth] (-4,0) -- (-4,h)node[pos=.5, left]{$h$};
    \end{tikzpicture}
}
\end{minipage}

\vspace{3pt}\hrule\vspace{3pt}\footnotesize\textcolor{gray!80!black}{\textbf{Các mã dùng chung hình (HINH\_228):} \texttt{L09\_C4\_B2\_12, L09\_C4\_B2\_38, L09\_C4\_B2\_52}}
\end{tcolorbox}
\vspace{4pt}


\begin{tcolorbox}[
    breakable,
    colback=white,
    colframe=blue!60!black,
    boxrule=0.8pt,
    arc=3pt,
    title={\textbf{Câu \#329} \quad \textbf{MÃ CÂU: \texttt{[L09\_C4\_B2\_36]}} \quad | \quad \textbf{Nhóm: \texttt{HINH\_229}} \quad | \quad \textbf{MapID: \texttt{9H4C2-5}} \hfill \textcolor{yellow!90!white}{[Hình gốc .pyc]}}
]
\textbf{Dạng toán:} Bài toán thực tế tính chiều cao của con dốc \hfill \textit{(Tự luận -- Vận dụng cao)}
\vspace{4pt}\hrule\vspace{6pt}
\begin{minipage}[t]{0.58\textwidth}
\vspace{0pt}
Lúc 6h sáng, bạn Anh đi xe đạp từ nhà (điểm ${A}$) đến trường (điểm ${B}$) phải leo và xuống một con dốc (như hình vẽ sau) biết ${AB}=1,6km$, $\widehat{A}=6^{\circ}$, $\widehat{B}=4^{\circ}$. Tính chiều cao ${h}$ của con dốc (theo mét làm tròn đến chữ số hàng đơn vị)
\end{minipage}\hfill
\begin{minipage}[t]{0.39\textwidth}
\vspace{0pt}
\centering

\resizebox{\linewidth}{!}{%
\begin{tikzpicture}
    
    % Vẽ tam giác ABC
    \coordinate [label=left:A] (A) at (0,0);
    \coordinate [label=right:B] (B) at (8,0);
    \coordinate [label=above:C] (C) at (2,1);
    
    % Đặt các góc và chiều cao
    \draw (A) -- (B) -- (C) -- cycle;
    
    % Vẽ chiều cao từ C tới AB
\draw (C) -- (2,0) coordinate (H);
\node[below] at (H) {H};
      % Vẽ dấu vuông tại điểm H
        \draw (2,0) rectangle +(0.1,0.1);
    % Nhãn chiều cao
    \draw[<-] (C) -- ($(C)!0.5!(H)$) node[midway, right] {$h$};
    
    % Vẽ các cung góc
    
    \node at ($(A)+(0.6,0.15)$) {$6^{\circ}$};
    

    \node at ($(B)+(-1.9,0.15)$) {$4^{\circ}$};
    
\end{tikzpicture}
}
\end{minipage}

\vspace{3pt}\hrule\vspace{3pt}\footnotesize\textcolor{gray!80!black}{\textbf{Các mã dùng chung hình (HINH\_229):} \texttt{L09\_C4\_B2\_13, L09\_C4\_B2\_37, L09\_C4\_B2\_50}}
\end{tcolorbox}
\vspace{4pt}


\subsection{Chương 5 -  Đường Tròn}

\begin{tcolorbox}[
    breakable,
    colback=white,
    colframe=blue!60!black,
    boxrule=0.8pt,
    arc=3pt,
    title={\textbf{Câu \#330} \quad \textbf{MÃ CÂU: \texttt{[L09\_C5\_B1\_12]}} \quad | \quad \textbf{Nhóm: \texttt{HINH\_247}} \quad | \quad \textbf{MapID: \texttt{9H5V1-1}} \hfill \textcolor{yellow!90!white}{[Hình gốc .pyc]}}
]
\textbf{Dạng toán:} Tìm bán kính đường tròn ngoại tiếp tam giác đều \hfill \textit{(Trắc nghiệm -- Vận dụng)}
\vspace{4pt}\hrule\vspace{6pt}
\begin{minipage}[t]{0.58\textwidth}
\vspace{0pt}
Cho tam giác đều cạnh bằng ${10}cm$. Bán kính đường tròn ngoại tiếp tam giác đó là bao nhiêu? \\

\vspace{4pt}\noindent
\textbf{\textcolor{blue!75!black}{A. (Đ)}} ${\dfrac{10 \sqrt{3}}{3}}$ \quad \textbf{B.} ${\dfrac{5 \sqrt{3}}{3}}$ \quad \textbf{C.} ${\dfrac{10 \sqrt{6}}{3}}$ \quad \textbf{D.} ${5 \sqrt{3}}$
\end{minipage}\hfill
\begin{minipage}[t]{0.39\textwidth}
\vspace{0pt}
\centering

\resizebox{\linewidth}{!}{%
\begin{tikzpicture}
    % Vẽ đường tròn tâm O bán kính 3
    \draw[thick] (0,0) circle (3);
    
    % Định nghĩa điểm O
    \coordinate [label=below right:$O$] (O) at (0,0);
    
    % Định nghĩa các đỉnh của tam giác đều ABC
    \coordinate [label=above:$A$] (A) at (90:3);
    \coordinate [label=below left:$B$] (B) at (210:3);
    \coordinate [label=below right:$C$] (C) at (330:3);

    % Vẽ tam giác ABC
    \draw[thick] (A) -- (B) -- (C) -- cycle;
    \coordinate [label=below:$H$] (H) at ($(B)!0.5!(C)$);
        
        % Vẽ đường cao AH
        \draw[ thick] (A) -- (H);
            \draw[ thick] (B) -- (O);
                \draw[ thick] (C) -- (O);
\end{tikzpicture}
}
\end{minipage}

\vspace{3pt}\hrule\vspace{3pt}\footnotesize\textcolor{gray!80!black}{\textbf{Các mã dùng chung hình (HINH\_247):} \texttt{L09\_C5\_B1\_13, L09\_C5\_B1\_16, L09\_C5\_B1\_25, L09\_C5\_B1\_26, L09\_C5\_B1\_29}}
\end{tcolorbox}
\vspace{4pt}


\begin{tcolorbox}[
    breakable,
    colback=white,
    colframe=blue!60!black,
    boxrule=0.8pt,
    arc=3pt,
    title={\textbf{Câu \#331} \quad \textbf{MÃ CÂU: \texttt{[L09\_C5\_B1\_13]}} \quad | \quad \textbf{Nhóm: \texttt{HINH\_247}} \quad | \quad \textbf{MapID: \texttt{9H5C1-1}} \hfill \textcolor{yellow!90!white}{[Hình gốc .pyc]}}
]
\textbf{Dạng toán:} Tính độ dài cạnh của tam giác đều khi biết bán kính đường tròn ngoại tiếp \hfill \textit{(Trắc nghiệm -- Vận dụng cao)}
\vspace{4pt}\hrule\vspace{6pt}
\begin{minipage}[t]{0.58\textwidth}
\vspace{0pt}
Tìm số đo cạnh của một tam giác đều, biết bán kính đường tròn ngoại tiếp tam giác đó là ${\sqrt{3}}cm$. \\

\vspace{4pt}\noindent
\textbf{A.} ${2}cm$ \quad \textbf{\textcolor{blue!75!black}{B. (Đ)}} ${3}cm$ \quad \textbf{C.} ${4}cm$ \quad \textbf{D.} ${5}cm$
\end{minipage}\hfill
\begin{minipage}[t]{0.39\textwidth}
\vspace{0pt}
\centering

\resizebox{\linewidth}{!}{%
\begin{tikzpicture}
    % Vẽ đường tròn tâm O bán kính 3
    \draw[thick] (0,0) circle (3);
    
    % Định nghĩa điểm O
    \coordinate [label=below right:$O$] (O) at (0,0);
    
    % Định nghĩa các đỉnh của tam giác đều ABC
    \coordinate [label=above:$A$] (A) at (90:3);
    \coordinate [label=below left:$B$] (B) at (210:3);
    \coordinate [label=below right:$C$] (C) at (330:3);

    % Vẽ tam giác ABC
    \draw[thick] (A) -- (B) -- (C) -- cycle;
    \coordinate [label=below:$H$] (H) at ($(B)!0.5!(C)$);
        
        % Vẽ đường cao AH
        \draw[ thick] (A) -- (H);
            \draw[ thick] (B) -- (O);
                \draw[ thick] (C) -- (O);
\end{tikzpicture}
}
\end{minipage}

\vspace{3pt}\hrule\vspace{3pt}\footnotesize\textcolor{gray!80!black}{\textbf{Các mã dùng chung hình (HINH\_247):} \texttt{L09\_C5\_B1\_12, L09\_C5\_B1\_16, L09\_C5\_B1\_25, L09\_C5\_B1\_26, L09\_C5\_B1\_29}}
\end{tcolorbox}
\vspace{4pt}


\begin{tcolorbox}[
    breakable,
    colback=white,
    colframe=blue!60!black,
    boxrule=0.8pt,
    arc=3pt,
    title={\textbf{Câu \#332} \quad \textbf{MÃ CÂU: \texttt{[L09\_C5\_B1\_14]}} \quad | \quad \textbf{Nhóm: \texttt{HINH\_248}} \quad | \quad \textbf{MapID: \texttt{9H5V1-1}} \hfill \textcolor{yellow!90!white}{[Hình gốc .pyc]}}
]
\textbf{Dạng toán:} Bài toán tổng hợp 1 \hfill \textit{(Trắc nghiệm -- Vận dụng)}
\vspace{4pt}\hrule\vspace{6pt}
\begin{minipage}[t]{0.58\textwidth}
\vspace{0pt}
Cho đường tròn ${(O;7cm)}$. Vẽ dây cung $AB=5cm$, ${C}$ là điểm trên dây cung ${AB}$ sao cho ${AC=4cm}$. Vẽ ${CD}$ vuông góc với ${OA}$ tại ${D}$. Độ dài đoạn thẳng ${AD}$ là:

\vspace{4pt}\noindent
\textbf{A.} ${\dfrac{20}{21}}cm$ \quad \textbf{\textcolor{blue!75!black}{B. (Đ)}} ${\dfrac{10}{7}}cm$ \quad \textbf{C.} ${\dfrac{2}{7}}cm$ \quad \textbf{D.} ${\dfrac{20}{7}}cm$
\end{minipage}\hfill
\begin{minipage}[t]{0.39\textwidth}
\vspace{0pt}
\centering

\resizebox{\linewidth}{!}{%
\begin{tikzpicture}
    % Vẽ đường tròn tâm O bán kính 3
    \coordinate (O) at (0, 0);
    \draw[thick] (O) circle(3);
    
    % Đặt các điểm A, E sao cho AE là đường kính
    \coordinate (A) at (-3, 0);
    \coordinate (E) at (3, 0);
    
    % Đặt điểm B sao cho tam giác ABE vuông tại B
    \coordinate (B) at (0, 3);
    
    % Tính tọa độ điểm C trên đoạn AB
    \coordinate (C) at ($(A)!0.6!(B)$); % Điểm C nằm trên đoạn AB, cách A 60% đoạn AB
    
    % Tính tọa độ điểm D, với CD vuông góc với OA tại D
    \coordinate (D) at ($(O)!(C)!(A)$); % D là hình chiếu vuông góc của C lên OA
    
    % Vẽ tam giác ABE
    \draw[thick] (A) -- (B) -- (E) -- cycle;
    
    % Vẽ đoạn CD
    \draw[thick] (C) -- (D);
        \path
    (B) pic[draw=black,fill=white, angle radius = 5] {right angle = A--B--E};
    
    % Vẽ dấu vuông tại D để chỉ rõ CD vuông góc với OA
    \draw (D) rectangle +(0.3,0.3);
    
    % Vẽ đường kính AE
    \draw[dashed] (A) -- (E);
    
    % Vẽ các điểm A, B, C, D, E, O
    \fill (O) circle(1.5pt) node[above left] {$O$};
    \fill (A) circle(1.5pt) node[left] {$A$};
    \fill (B) circle(1.5pt) node[above] {$B$};
    \fill (E) circle(1.5pt) node[right] {$E$};
    \fill (C) circle(1.5pt) node[above left] {$C$};
    \fill (D) circle(1.5pt) node[below right] {$D$};
    
    % Vẽ góc vuông tại B

    \end{tikzpicture}
}
\end{minipage}

\vspace{3pt}\hrule\vspace{3pt}\footnotesize\textcolor{gray!80!black}{\textbf{Các mã dùng chung hình (HINH\_248):} \texttt{L09\_C5\_B1\_27}}
\end{tcolorbox}
\vspace{4pt}


\begin{tcolorbox}[
    breakable,
    colback=white,
    colframe=blue!60!black,
    boxrule=0.8pt,
    arc=3pt,
    title={\textbf{Câu \#333} \quad \textbf{MÃ CÂU: \texttt{[L09\_C5\_B1\_15]}} \quad | \quad \textbf{Nhóm: \texttt{HINH\_249}} \quad | \quad \textbf{MapID: \texttt{9H5V1-1}} \hfill \textcolor{yellow!90!white}{[Hình gốc .pyc]}}
]
\textbf{Dạng toán:} Bài toán tổng hợp 2 \hfill \textit{(Trắc nghiệm -- Vận dụng)}
\vspace{4pt}\hrule\vspace{6pt}
\begin{minipage}[t]{0.58\textwidth}
\vspace{0pt}
Cho tam giác ${ABC}$ cân tại ${A}$ có $\widehat{A}=120^{\circ}$ nội tiếp đường tròn ${(O; 5cm)}$. Khi đó diện tích tam giác ${ABC}$ là:

\vspace{4pt}\noindent
\textbf{A.} $\dfrac{25 \sqrt{6}}{2}cm^{2}$ \quad \textbf{B.} $\dfrac{5 \sqrt{3}}{4}cm^{2}$ \quad \textbf{C.} $\dfrac{25 \sqrt{3}}{2}cm^{2}$ \quad \textbf{\textcolor{blue!75!black}{D. (Đ)}} $\dfrac{25 \sqrt{3}}{4} cm^{2}$
\end{minipage}\hfill
\begin{minipage}[t]{0.39\textwidth}
\vspace{0pt}
\centering

\resizebox{\linewidth}{!}{%
\begin{tikzpicture}
    
    % Vẽ đường tròn tâm O bán kính 3cm
    \draw[thick] (0,0) circle (3cm);
    
    % Định nghĩa các điểm A, O
    \coordinate [label=above:$A$] (A) at (0,3);
    \coordinate [label=below:$O$] (O) at (0,0);
    
    % Xác định tọa độ điểm B và C trên đường tròn sao cho tam giác ABC cân tại A với góc BAC = 120 độ
    \coordinate [label=below left:$B$] (B) at ({3*cos(150)},{3*sin(150)});
    \coordinate [label=below right:$C$] (C) at ({3*cos(30)},{3*sin(30)});
    
    % Xác định điểm H là trung điểm của BC
    \coordinate [label=below:$H$] (H) at ($(B)!0.5!(C)$);
    
    % Vẽ tam giác ABC
    \draw[thick] (B) -- (A) -- (C) -- cycle;
    
    % Nối A, H, O
    \draw[thick] (A) -- (H);
    \draw[thick] (O) -- (H);
    
    % Vẽ các đoạn thẳng nối các điểm với tâm đường tròn O
    \draw[dashed] (O) -- (A);
    \draw[dashed] (O) -- (B);
    \draw[dashed] (O) -- (C);
    
    % Vẽ các điểm A, B, C, O và H
    \fill (A) circle (2pt);
    \fill (B) circle (2pt);
    \fill (C) circle (2pt);
    \fill (O) circle (2pt);
    \fill (H) circle (2pt);
    
\end{tikzpicture}
}
\end{minipage}

\vspace{3pt}\hrule\vspace{3pt}\footnotesize\textcolor{gray!80!black}{\textbf{Các mã dùng chung hình (HINH\_249):} \texttt{L09\_C5\_B1\_28}}
\end{tcolorbox}
\vspace{4pt}


\begin{tcolorbox}[
    breakable,
    colback=white,
    colframe=blue!60!black,
    boxrule=0.8pt,
    arc=3pt,
    title={\textbf{Câu \#334} \quad \textbf{MÃ CÂU: \texttt{[L09\_C5\_B1\_16]}} \quad | \quad \textbf{Nhóm: \texttt{HINH\_247}} \quad | \quad \textbf{MapID: \texttt{9H5V1-1}} \hfill \textcolor{yellow!90!white}{[Hình gốc .pyc]}}
]
\textbf{Dạng toán:} Tính diện tích của tam giác đều khi biết bán kính đường tròn ngoại tiếp \hfill \textit{(Trắc nghiệm -- Vận dụng)}
\vspace{4pt}\hrule\vspace{6pt}
\begin{minipage}[t]{0.58\textwidth}
\vspace{0pt}
Tính diện tích của một tam giác đều, biết bán kính đường tròn ngoại tiếp tam giác đó là ${5}cm$.

\vspace{4pt}\noindent
\textbf{A.} ${\dfrac{15 \sqrt{3}}{4}}cm$ \quad \textbf{B.} ${\dfrac{25 \sqrt{3}}{2}}cm$ \quad \textbf{\textcolor{blue!75!black}{C. (Đ)}} ${\dfrac{75 \sqrt{3}}{4}}cm$ \quad \textbf{D.} ${\dfrac{75 \sqrt{3}}{2}}cm$
\end{minipage}\hfill
\begin{minipage}[t]{0.39\textwidth}
\vspace{0pt}
\centering

\resizebox{\linewidth}{!}{%
\begin{tikzpicture}
    % Vẽ đường tròn tâm O bán kính 3
    \draw[thick] (0,0) circle (3);
    
    % Định nghĩa điểm O
    \coordinate [label=below right:$O$] (O) at (0,0);
    
    % Định nghĩa các đỉnh của tam giác đều ABC
    \coordinate [label=above:$A$] (A) at (90:3);
    \coordinate [label=below left:$B$] (B) at (210:3);
    \coordinate [label=below right:$C$] (C) at (330:3);

    % Vẽ tam giác ABC
    \draw[thick] (A) -- (B) -- (C) -- cycle;
    \coordinate [label=below:$H$] (H) at ($(B)!0.5!(C)$);
        
        % Vẽ đường cao AH
        \draw[ thick] (A) -- (H);
            \draw[ thick] (B) -- (O);
                \draw[ thick] (C) -- (O);
\end{tikzpicture}
}
\end{minipage}

\vspace{3pt}\hrule\vspace{3pt}\footnotesize\textcolor{gray!80!black}{\textbf{Các mã dùng chung hình (HINH\_247):} \texttt{L09\_C5\_B1\_12, L09\_C5\_B1\_13, L09\_C5\_B1\_25, L09\_C5\_B1\_26, L09\_C5\_B1\_29}}
\end{tcolorbox}
\vspace{4pt}


\begin{tcolorbox}[
    breakable,
    colback=white,
    colframe=blue!60!black,
    boxrule=0.8pt,
    arc=3pt,
    title={\textbf{Câu \#335} \quad \textbf{MÃ CÂU: \texttt{[L09\_C5\_B1\_17]}} \quad | \quad \textbf{Nhóm: \texttt{HINH\_250}} \quad | \quad \textbf{MapID: \texttt{9H5V1-1}} \hfill \textcolor{yellow!90!white}{[Hình gốc .pyc]}}
]
\textbf{Dạng toán:} Bài toán tổng hợp 3 \hfill \textit{(Trắc nghiệm -- Vận dụng)}
\vspace{4pt}\hrule\vspace{6pt}
\begin{minipage}[t]{0.58\textwidth}
\vspace{0pt}
Cho tam giác ${ABC}$ có $AB=5cm$, $AC=4cm$, $BC=3cm$. Đường tròn ${(O)}$ là đường tròn có bán kính nhỏ nhất đi qua ${C}$ và tiếp xúc với ${AB}$. Gọi ${Q}$, ${P}$ là giao điểm khác ${C}$ của đường tròn ${(O)}$ và cạnh ${CA}$, ${CB}$. Độ dài đoạn ${PQ}$ là:

\vspace{4pt}\noindent
\textbf{A.} ${\dfrac{20}{3}}cm$ \quad \textbf{\textcolor{blue!75!black}{B. (Đ)}} ${\dfrac{12}{5}}cm$ \quad \textbf{C.} ${\dfrac{4}{5}}cm$ \quad \textbf{D.} ${\dfrac{5}{3}}cm$
\end{minipage}\hfill
\begin{minipage}[t]{0.39\textwidth}
\vspace{0pt}
\centering

\resizebox{\linewidth}{!}{%
\begin{tikzpicture}[scale=2]
    % Định nghĩa các điểm A, B, C
    \coordinate [label=left:A] (A) at (0, 0);
    \coordinate [label=below right:C] (C) at (2, 0);
    \coordinate [label=above:B] (B) at (2, 3);
    % Tính tọa độ điểm H là hình chiếu vuông góc từ C lên AB
    \coordinate [label=left:H] (H) at (0.62, 0.92);
    \coordinate [label=left:O] (O) at (1.31, 0.46);
    
        % Vẽ đường tròn đường kính CH
    \draw [blue] (O) circle [radius=0.83];
    % Vẽ tam giác ABC
    \draw [thick] (A) -- (B) -- (C) -- cycle;
    
    % Vẽ đường cao CH
    \draw [thick] (C) --  (H);
    \coordinate (O) at ($(C)!0.5!(H)$);
    

    
    % Vẽ góc vuông tại C
    \draw (C) rectangle +(-0.2, 0.2);

    % Tính trung điểm của CH để xác định tâm đường tròn

    % Tìm và vẽ giao điểm P và Q
    \path [name path=CA] (C) -- (A);
    \path [name path=CB] (C) -- (B);
    \path [name path=circle] (O) circle [radius=0.83];
    
    % Tính toán giao điểm P và Q
    \path [name intersections={of=circle and CA, by=P}];
    \path [name intersections={of=circle and CB, by=Q}];
        \draw [thick] (Q) --  (H);
            \draw [thick] (P) --  (H);
                \draw [thick] (P) --  (Q);
    % Vẽ các giao điểm P và Q
    \fill [black] (P) circle (1pt) node[below left] {P};
    \fill [black] (Q) circle (1pt) node[right] {Q};

\end{tikzpicture}
}
\end{minipage}

\vspace{3pt}\hrule\vspace{3pt}\footnotesize\textcolor{gray!80!black}{\textbf{Các mã dùng chung hình (HINH\_250):} \texttt{L09\_C5\_B1\_30}}
\end{tcolorbox}
\vspace{4pt}


\begin{tcolorbox}[
    breakable,
    colback=white,
    colframe=blue!60!black,
    boxrule=0.8pt,
    arc=3pt,
    title={\textbf{Câu \#336} \quad \textbf{MÃ CÂU: \texttt{[L09\_C5\_B1\_18]}} \quad | \quad \textbf{Nhóm: \texttt{HINH\_251}} \quad | \quad \textbf{MapID: \texttt{9H5C1-1}} \hfill \textcolor{yellow!90!white}{[Hình gốc .pyc]}}
]
\textbf{Dạng toán:} Bài toán min, max \hfill \textit{(Trắc nghiệm -- Vận dụng cao)}
\vspace{4pt}\hrule\vspace{6pt}
\begin{minipage}[t]{0.58\textwidth}
\vspace{0pt}
Cho đường tròn ${(O; 9cm)}$. Gọi ${AB}$ và ${CD}$ là hai đường kính có vị trí thay đổi. Diện tích lớn nhất của tứ giác ${ACBD}$ là

\vspace{4pt}\noindent
\textbf{A.} ${161}cm^{2}$ \quad \textbf{\textcolor{blue!75!black}{B. (Đ)}} ${162}cm^{2}$ \quad \textbf{C.} ${82}cm^{2}$ \quad \textbf{D.} ${81}cm^{2}$
\end{minipage}\hfill
\begin{minipage}[t]{0.39\textwidth}
\vspace{0pt}
\centering

\resizebox{\linewidth}{!}{%
\begin{tikzpicture}
    % Vẽ lưới để hỗ trợ vẽ

    
    % Vẽ đường tròn tâm O
    \draw[thick] (0,0) circle(3);
    
    % Tâm O
    \node at (0,0) [circle,fill,inner sep=1.5pt, label=above right:{$O$}] {};
    
    % Tọa độ các điểm của hình chữ nhật nội tiếp
    \coordinate (A) at (2.54,1.59);
    \coordinate (D) at (-2.54,1.59);
    \coordinate (C) at (2.54,-1.59);
    \coordinate (B) at (-2.54,-1.59);
    
    % Vẽ các điểm A, B, C, D
    \foreach \p/\pos in {A/right, B/left, C/right, D/left}
    \node at (\p) [circle,fill,inner sep=1.5pt, label=\pos:{$\p$}] {};
    
    % Vẽ hình chữ nhật ACBD
    \draw[thick] (A) -- (D) -- (B) -- (C) -- (A)--cycle;
    \draw[thick] (A) -- (B);
        \draw[thick] (D) -- (C);
    % Xác định và vẽ điểm H trên AB sao cho CH vuông góc với AB
    \coordinate (H) at (1.11,0.69); % Điểm H trên đoạn AB
    \draw[thick] (C) -- (H);
            \path
    (H) pic[draw=black,fill=white!25, angle radius = 5] {right angle = C--H--A};
    % Vẽ điểm H
    \node at (H) [circle,fill,inner sep=1.5pt, label=above:{$H$}] {};

\end{tikzpicture}
}
\end{minipage}

\vspace{3pt}\hrule\vspace{3pt}\footnotesize\textcolor{gray!80!black}{\textbf{Các mã dùng chung hình (HINH\_251):} \texttt{L09\_C5\_B1\_31}}
\end{tcolorbox}
\vspace{4pt}


\begin{tcolorbox}[
    breakable,
    colback=white,
    colframe=blue!60!black,
    boxrule=0.8pt,
    arc=3pt,
    title={\textbf{Câu \#337} \quad \textbf{MÃ CÂU: \texttt{[L09\_C5\_B1\_19]}} \quad | \quad \textbf{Nhóm: \texttt{HINH\_252}} \quad | \quad \textbf{MapID: \texttt{9H5C1-1}} \hfill \textcolor{yellow!90!white}{[Hình gốc .pyc]}}
]
\textbf{Dạng toán:} Bài toán tổng hợp 4 \hfill \textit{(Trắc nghiệm -- Vận dụng cao)}
\vspace{4pt}\hrule\vspace{6pt}
\begin{minipage}[t]{0.58\textwidth}
\vspace{0pt}
Cho hình thoi ${ABCD}$ có cạnh ${AB}=9cm$, hai đường chéo cắt nhau tại ${(O)}$. Gọi $R_{1}, R_{2}$ lần lượt là các bán kính các đường tròn ngoại tiếp các tam giác ${ABC}$ và tam giác ${ABD}$. Giá trị của biểu thức $\dfrac{1}{R_{1}^{2}}+ \dfrac{1}{R_{2}^{2}}$ là

\vspace{4pt}\noindent
\textbf{\textcolor{blue!75!black}{A. (Đ)}} ${\dfrac{4}{81}}$ \quad \textbf{B.} ${\dfrac{1}{81}}$ \quad \textbf{C.} ${\dfrac{2}{81}}$ \quad \textbf{D.} ${\dfrac{2}{81}}$
\end{minipage}\hfill
\begin{minipage}[t]{0.39\textwidth}
\vspace{0pt}
\centering

\resizebox{\linewidth}{!}{%
\begin{tikzpicture}
    \def\khvuong[size=#1](#2,#3,#4){%
        \draw ($(#3)!#1!(#2)$) --
        ($($(#3)!#1!(#2)$)+($(#3)!#1!(#4)$)-(#3)$)--
        ($(#3)!#1!(#4)$);
    }
    % Định nghĩa các điểm với nhãn
    \coordinate[label=above left:A] (a) at (-4,0);
    \coordinate[label=above:B] (b) at (0,2);
    \coordinate[label=above right:C] (c) at (4,0);
    \coordinate[label=below right:D] (d) at (0,-2);
    \coordinate[label=above right:O] (o) at (0,0);
    \coordinate[label=above left:M] (m) at (-2,1);
    \coordinate[label=below left:K] (k) at (-1.5,0);
    \coordinate[label=below:I] (i) at (0,-3);
    
    % Vẽ các đoạn thẳng giữa các điểm
    \draw (m) -- (i);
    \draw (a) -- (c);
    \draw (b) -- (i);
    \draw (a) -- (b) -- (c) -- (d) -- cycle;
    
\khvuong[size=6pt](b,m,k);
\khvuong[size=6pt](b,o,a);
\end{tikzpicture}
}
\end{minipage}

\vspace{3pt}\hrule\vspace{3pt}\footnotesize\textcolor{gray!80!black}{\textbf{Các mã dùng chung hình (HINH\_252):} \texttt{L09\_C5\_B1\_32}}
\end{tcolorbox}
\vspace{4pt}


\begin{tcolorbox}[
    breakable,
    colback=white,
    colframe=blue!60!black,
    boxrule=0.8pt,
    arc=3pt,
    title={\textbf{Câu \#338} \quad \textbf{MÃ CÂU: \texttt{[L09\_C5\_B1\_25]}} \quad | \quad \textbf{Nhóm: \texttt{HINH\_247}} \quad | \quad \textbf{MapID: \texttt{9H5V1-1}} \hfill \textcolor{yellow!90!white}{[Hình gốc .pyc]}}
]
\textbf{Dạng toán:} Tìm bán kính đường tròn ngoại tiếp tam giác đều \hfill \textit{(Tự luận -- Vận dụng)}
\vspace{4pt}\hrule\vspace{6pt}
\begin{minipage}[t]{0.58\textwidth}
\vspace{0pt}
Cho tam giác đều cạnh bằng ${10}cm$. Bán kính đường tròn ngoại tiếp tam giác đó là: \\
\end{minipage}\hfill
\begin{minipage}[t]{0.39\textwidth}
\vspace{0pt}
\centering

\resizebox{\linewidth}{!}{%
\begin{tikzpicture}
    % Vẽ đường tròn tâm O bán kính 3
    \draw[thick] (0,0) circle (3);
    
    % Định nghĩa điểm O
    \coordinate [label=below right:$O$] (O) at (0,0);
    
    % Định nghĩa các đỉnh của tam giác đều ABC
    \coordinate [label=above:$A$] (A) at (90:3);
    \coordinate [label=below left:$B$] (B) at (210:3);
    \coordinate [label=below right:$C$] (C) at (330:3);

    % Vẽ tam giác ABC
    \draw[thick] (A) -- (B) -- (C) -- cycle;
    \coordinate [label=below:$H$] (H) at ($(B)!0.5!(C)$);
        
        % Vẽ đường cao AH
        \draw[ thick] (A) -- (H);
            \draw[ thick] (B) -- (O);
                \draw[ thick] (C) -- (O);
\end{tikzpicture}
}
\end{minipage}

\vspace{3pt}\hrule\vspace{3pt}\footnotesize\textcolor{gray!80!black}{\textbf{Các mã dùng chung hình (HINH\_247):} \texttt{L09\_C5\_B1\_12, L09\_C5\_B1\_13, L09\_C5\_B1\_16, L09\_C5\_B1\_26, L09\_C5\_B1\_29}}
\end{tcolorbox}
\vspace{4pt}


\begin{tcolorbox}[
    breakable,
    colback=white,
    colframe=blue!60!black,
    boxrule=0.8pt,
    arc=3pt,
    title={\textbf{Câu \#339} \quad \textbf{MÃ CÂU: \texttt{[L09\_C5\_B1\_26]}} \quad | \quad \textbf{Nhóm: \texttt{HINH\_247}} \quad | \quad \textbf{MapID: \texttt{9H5C1-1}} \hfill \textcolor{yellow!90!white}{[Hình gốc .pyc]}}
]
\textbf{Dạng toán:} Tính độ dài cạnh của tam giác đều khi biết bán kính đường tròn ngoại tiếp \hfill \textit{(Tự luận -- Vận dụng cao)}
\vspace{4pt}\hrule\vspace{6pt}
\begin{minipage}[t]{0.58\textwidth}
\vspace{0pt}
Tìm số đo cạnh của một tam giác đều, biết bán kính đường tròn ngoại tiếp tam giác đó là ${\dfrac{5 \sqrt{3}}{3}}cm$. \\
\end{minipage}\hfill
\begin{minipage}[t]{0.39\textwidth}
\vspace{0pt}
\centering

\resizebox{\linewidth}{!}{%
\begin{tikzpicture}
    % Vẽ đường tròn tâm O bán kính 3
    \draw[thick] (0,0) circle (3);
    
    % Định nghĩa điểm O
    \coordinate [label=below right:$O$] (O) at (0,0);
    
    % Định nghĩa các đỉnh của tam giác đều ABC
    \coordinate [label=above:$A$] (A) at (90:3);
    \coordinate [label=below left:$B$] (B) at (210:3);
    \coordinate [label=below right:$C$] (C) at (330:3);

    % Vẽ tam giác ABC
    \draw[thick] (A) -- (B) -- (C) -- cycle;
    \coordinate [label=below:$H$] (H) at ($(B)!0.5!(C)$);
        
        % Vẽ đường cao AH
        \draw[ thick] (A) -- (H);
            \draw[ thick] (B) -- (O);
                \draw[ thick] (C) -- (O);
\end{tikzpicture}
}
\end{minipage}

\vspace{3pt}\hrule\vspace{3pt}\footnotesize\textcolor{gray!80!black}{\textbf{Các mã dùng chung hình (HINH\_247):} \texttt{L09\_C5\_B1\_12, L09\_C5\_B1\_13, L09\_C5\_B1\_16, L09\_C5\_B1\_25, L09\_C5\_B1\_29}}
\end{tcolorbox}
\vspace{4pt}


\begin{tcolorbox}[
    breakable,
    colback=white,
    colframe=blue!60!black,
    boxrule=0.8pt,
    arc=3pt,
    title={\textbf{Câu \#340} \quad \textbf{MÃ CÂU: \texttt{[L09\_C5\_B1\_27]}} \quad | \quad \textbf{Nhóm: \texttt{HINH\_248}} \quad | \quad \textbf{MapID: \texttt{9H5V1-1}} \hfill \textcolor{yellow!90!white}{[Hình gốc .pyc]}}
]
\textbf{Dạng toán:} Bài toán tổng hợp 1 \hfill \textit{(Tự luận -- Vận dụng)}
\vspace{4pt}\hrule\vspace{6pt}
\begin{minipage}[t]{0.58\textwidth}
\vspace{0pt}
Cho đường tròn ${(O; 10cm)}$. Vẽ dây cung $AB=10cm$, ${C}$ là điểm trên dây cung ${AB}$ sao cho $AC=7cm$. Vẽ ${CD}$ vuông góc với ${OA}$ tại ${D}$. Độ dài đoạn thẳng ${AD}$ là:\\
\end{minipage}\hfill
\begin{minipage}[t]{0.39\textwidth}
\vspace{0pt}
\centering

\resizebox{\linewidth}{!}{%
\begin{tikzpicture}
    % Vẽ đường tròn tâm O bán kính 3
    \coordinate (O) at (0, 0);
    \draw[thick] (O) circle(3);
    
    % Đặt các điểm A, E sao cho AE là đường kính
    \coordinate (A) at (-3, 0);
    \coordinate (E) at (3, 0);
    
    % Đặt điểm B sao cho tam giác ABE vuông tại B
    \coordinate (B) at (0, 3);
    
    % Tính tọa độ điểm C trên đoạn AB
    \coordinate (C) at ($(A)!0.6!(B)$); % Điểm C nằm trên đoạn AB, cách A 60% đoạn AB
    
    % Tính tọa độ điểm D, với CD vuông góc với OA tại D
    \coordinate (D) at ($(O)!(C)!(A)$); % D là hình chiếu vuông góc của C lên OA
    
    % Vẽ tam giác ABE
    \draw[thick] (A) -- (B) -- (E) -- cycle;
    
    % Vẽ đoạn CD
    \draw[thick] (C) -- (D);
        \path
    (B) pic[draw=black,fill=white, angle radius = 5] {right angle = A--B--E};
    
    % Vẽ dấu vuông tại D để chỉ rõ CD vuông góc với OA
    \draw (D) rectangle +(0.3,0.3);
    
    % Vẽ đường kính AE
    \draw[dashed] (A) -- (E);
    
    % Vẽ các điểm A, B, C, D, E, O
    \fill (O) circle(1.5pt) node[above left] {$O$};
    \fill (A) circle(1.5pt) node[left] {$A$};
    \fill (B) circle(1.5pt) node[above] {$B$};
    \fill (E) circle(1.5pt) node[right] {$E$};
    \fill (C) circle(1.5pt) node[above left] {$C$};
    \fill (D) circle(1.5pt) node[below right] {$D$};
    
    % Vẽ góc vuông tại B

    \end{tikzpicture}
}
\end{minipage}

\vspace{3pt}\hrule\vspace{3pt}\footnotesize\textcolor{gray!80!black}{\textbf{Các mã dùng chung hình (HINH\_248):} \texttt{L09\_C5\_B1\_14}}
\end{tcolorbox}
\vspace{4pt}


\begin{tcolorbox}[
    breakable,
    colback=white,
    colframe=blue!60!black,
    boxrule=0.8pt,
    arc=3pt,
    title={\textbf{Câu \#341} \quad \textbf{MÃ CÂU: \texttt{[L09\_C5\_B1\_28]}} \quad | \quad \textbf{Nhóm: \texttt{HINH\_249}} \quad | \quad \textbf{MapID: \texttt{9H5V1-1}} \hfill \textcolor{yellow!90!white}{[Hình gốc .pyc]}}
]
\textbf{Dạng toán:} Bài toán tổng hợp 2 \hfill \textit{(Tự luận -- Vận dụng)}
\vspace{4pt}\hrule\vspace{6pt}
\begin{minipage}[t]{0.58\textwidth}
\vspace{0pt}
Cho tam giác ${ABC}$ cân tại ${A}$ có $\widehat{A}=120^{\circ}$ nội tiếp đường tròn ${(O; 8cm)}$. Khi đó diện tích tam giác ${ABC}$ là: \\
\end{minipage}\hfill
\begin{minipage}[t]{0.39\textwidth}
\vspace{0pt}
\centering

\resizebox{\linewidth}{!}{%
\begin{tikzpicture}
    
    % Vẽ đường tròn tâm O bán kính 3cm
    \draw[thick] (0,0) circle (3cm);
    
    % Định nghĩa các điểm A, O
    \coordinate [label=above:$A$] (A) at (0,3);
    \coordinate [label=below:$O$] (O) at (0,0);
    
    % Xác định tọa độ điểm B và C trên đường tròn sao cho tam giác ABC cân tại A với góc BAC = 120 độ
    \coordinate [label=below left:$B$] (B) at ({3*cos(150)},{3*sin(150)});
    \coordinate [label=below right:$C$] (C) at ({3*cos(30)},{3*sin(30)});
    
    % Xác định điểm H là trung điểm của BC
    \coordinate [label=below:$H$] (H) at ($(B)!0.5!(C)$);
    
    % Vẽ tam giác ABC
    \draw[thick] (B) -- (A) -- (C) -- cycle;
    
    % Nối A, H, O
    \draw[thick] (A) -- (H);
    \draw[thick] (O) -- (H);
    
    % Vẽ các đoạn thẳng nối các điểm với tâm đường tròn O
    \draw[dashed] (O) -- (A);
    \draw[dashed] (O) -- (B);
    \draw[dashed] (O) -- (C);
    
    % Vẽ các điểm A, B, C, O và H
    \fill (A) circle (2pt);
    \fill (B) circle (2pt);
    \fill (C) circle (2pt);
    \fill (O) circle (2pt);
    \fill (H) circle (2pt);
    
\end{tikzpicture}
}
\end{minipage}

\vspace{3pt}\hrule\vspace{3pt}\footnotesize\textcolor{gray!80!black}{\textbf{Các mã dùng chung hình (HINH\_249):} \texttt{L09\_C5\_B1\_15}}
\end{tcolorbox}
\vspace{4pt}


\begin{tcolorbox}[
    breakable,
    colback=white,
    colframe=blue!60!black,
    boxrule=0.8pt,
    arc=3pt,
    title={\textbf{Câu \#342} \quad \textbf{MÃ CÂU: \texttt{[L09\_C5\_B1\_29]}} \quad | \quad \textbf{Nhóm: \texttt{HINH\_247}} \quad | \quad \textbf{MapID: \texttt{9H5V1-1}} \hfill \textcolor{yellow!90!white}{[Hình gốc .pyc]}}
]
\textbf{Dạng toán:} Tính diện tích của tam giác đều khi biết bán kính đường tròn ngoại tiếp \hfill \textit{(Tự luận -- Vận dụng)}
\vspace{4pt}\hrule\vspace{6pt}
\begin{minipage}[t]{0.58\textwidth}
\vspace{0pt}
Tính diện tích của một tam giác đều, biết bán kính đường tròn ngoại tiếp tam giác đó là ${9}cm$. \\
\end{minipage}\hfill
\begin{minipage}[t]{0.39\textwidth}
\vspace{0pt}
\centering

\resizebox{\linewidth}{!}{%
\begin{tikzpicture}
    % Vẽ đường tròn tâm O bán kính 3
    \draw[thick] (0,0) circle (3);
    
    % Định nghĩa điểm O
    \coordinate [label=below right:$O$] (O) at (0,0);
    
    % Định nghĩa các đỉnh của tam giác đều ABC
    \coordinate [label=above:$A$] (A) at (90:3);
    \coordinate [label=below left:$B$] (B) at (210:3);
    \coordinate [label=below right:$C$] (C) at (330:3);

    % Vẽ tam giác ABC
    \draw[thick] (A) -- (B) -- (C) -- cycle;
    \coordinate [label=below:$H$] (H) at ($(B)!0.5!(C)$);
        
        % Vẽ đường cao AH
        \draw[ thick] (A) -- (H);
            \draw[ thick] (B) -- (O);
                \draw[ thick] (C) -- (O);
\end{tikzpicture}
}
\end{minipage}

\vspace{3pt}\hrule\vspace{3pt}\footnotesize\textcolor{gray!80!black}{\textbf{Các mã dùng chung hình (HINH\_247):} \texttt{L09\_C5\_B1\_12, L09\_C5\_B1\_13, L09\_C5\_B1\_16, L09\_C5\_B1\_25, L09\_C5\_B1\_26}}
\end{tcolorbox}
\vspace{4pt}


\begin{tcolorbox}[
    breakable,
    colback=white,
    colframe=blue!60!black,
    boxrule=0.8pt,
    arc=3pt,
    title={\textbf{Câu \#343} \quad \textbf{MÃ CÂU: \texttt{[L09\_C5\_B1\_30]}} \quad | \quad \textbf{Nhóm: \texttt{HINH\_250}} \quad | \quad \textbf{MapID: \texttt{9H5V1-1}} \hfill \textcolor{yellow!90!white}{[Hình gốc .pyc]}}
]
\textbf{Dạng toán:} Bài toán tổng hợp 3 \hfill \textit{(Tự luận -- Vận dụng)}
\vspace{4pt}\hrule\vspace{6pt}
\begin{minipage}[t]{0.58\textwidth}
\vspace{0pt}
Cho tam giác ${ABC}$ có $AB=40cm$, $AC=32cm$, $BC=24cm$. Đường tròn ${(O)}$ là đường tròn có bán kính nhỏ nhất đi qua ${C}$ và tiếp xúc với ${BC}$. Gọi ${Q}$, ${R}$ là giao điểm khác ${C}$ của đường tròn ${(O)}$ và cạnh ${CA}$, ${CB}$. Độ dài đoạn ${PQ}$ là: \\
\end{minipage}\hfill
\begin{minipage}[t]{0.39\textwidth}
\vspace{0pt}
\centering

\resizebox{\linewidth}{!}{%
\begin{tikzpicture}[scale=2]
    % Định nghĩa các điểm A, B, C
    \coordinate [label=left:A] (A) at (0, 0);
    \coordinate [label=below right:C] (C) at (2, 0);
    \coordinate [label=above:B] (B) at (2, 3);
    % Tính tọa độ điểm H là hình chiếu vuông góc từ C lên AB
    \coordinate [label=left:H] (H) at (0.62, 0.92);
    \coordinate [label=left:O] (O) at (1.31, 0.46);
    
        % Vẽ đường tròn đường kính CH
    \draw [blue] (O) circle [radius=0.83];
    % Vẽ tam giác ABC
    \draw [thick] (A) -- (B) -- (C) -- cycle;
    
    % Vẽ đường cao CH
    \draw [thick] (C) --  (H);
    \coordinate (O) at ($(C)!0.5!(H)$);
    

    
    % Vẽ góc vuông tại C
    \draw (C) rectangle +(-0.2, 0.2);

    % Tính trung điểm của CH để xác định tâm đường tròn

    % Tìm và vẽ giao điểm P và Q
    \path [name path=CA] (C) -- (A);
    \path [name path=CB] (C) -- (B);
    \path [name path=circle] (O) circle [radius=0.83];
    
    % Tính toán giao điểm P và Q
    \path [name intersections={of=circle and CA, by=P}];
    \path [name intersections={of=circle and CB, by=Q}];
        \draw [thick] (Q) --  (H);
            \draw [thick] (P) --  (H);
                \draw [thick] (P) --  (Q);
    % Vẽ các giao điểm P và Q
    \fill [black] (P) circle (1pt) node[below left] {P};
    \fill [black] (Q) circle (1pt) node[right] {Q};

\end{tikzpicture}
}
\end{minipage}

\vspace{3pt}\hrule\vspace{3pt}\footnotesize\textcolor{gray!80!black}{\textbf{Các mã dùng chung hình (HINH\_250):} \texttt{L09\_C5\_B1\_17}}
\end{tcolorbox}
\vspace{4pt}


\begin{tcolorbox}[
    breakable,
    colback=white,
    colframe=blue!60!black,
    boxrule=0.8pt,
    arc=3pt,
    title={\textbf{Câu \#344} \quad \textbf{MÃ CÂU: \texttt{[L09\_C5\_B1\_31]}} \quad | \quad \textbf{Nhóm: \texttt{HINH\_251}} \quad | \quad \textbf{MapID: \texttt{9H5C1-1}} \hfill \textcolor{yellow!90!white}{[Hình gốc .pyc]}}
]
\textbf{Dạng toán:} Bài toán min, max \hfill \textit{(Tự luận -- Vận dụng cao)}
\vspace{4pt}\hrule\vspace{6pt}
\begin{minipage}[t]{0.58\textwidth}
\vspace{0pt}
Cho đường tròn ${(O; 9cm)}$. Gọi ${AB}$ và ${CD}$ là hai đường kính có vị trí thay đổi. Diện tích lớn nhất của tứ giác ${ACBD}$ là \\
\end{minipage}\hfill
\begin{minipage}[t]{0.39\textwidth}
\vspace{0pt}
\centering

\resizebox{\linewidth}{!}{%
\begin{tikzpicture}
    % Vẽ lưới để hỗ trợ vẽ

    
    % Vẽ đường tròn tâm O
    \draw[thick] (0,0) circle(3);
    
    % Tâm O
    \node at (0,0) [circle,fill,inner sep=1.5pt, label=above right:{$O$}] {};
    
    % Tọa độ các điểm của hình chữ nhật nội tiếp
    \coordinate (A) at (2.54,1.59);
    \coordinate (D) at (-2.54,1.59);
    \coordinate (C) at (2.54,-1.59);
    \coordinate (B) at (-2.54,-1.59);
    
    % Vẽ các điểm A, B, C, D
    \foreach \p/\pos in {A/right, B/left, C/right, D/left}
    \node at (\p) [circle,fill,inner sep=1.5pt, label=\pos:{$\p$}] {};
    
    % Vẽ hình chữ nhật ACBD
    \draw[thick] (A) -- (D) -- (B) -- (C) -- (A)--cycle;
    \draw[thick] (A) -- (B);
        \draw[thick] (D) -- (C);
    % Xác định và vẽ điểm H trên AB sao cho CH vuông góc với AB
    \coordinate (H) at (1.11,0.69); % Điểm H trên đoạn AB
    \draw[thick] (C) -- (H);
            \path
    (H) pic[draw=black,fill=white!25, angle radius = 5] {right angle = C--H--A};
    % Vẽ điểm H
    \node at (H) [circle,fill,inner sep=1.5pt, label=above:{$H$}] {};

\end{tikzpicture}
}
\end{minipage}

\vspace{3pt}\hrule\vspace{3pt}\footnotesize\textcolor{gray!80!black}{\textbf{Các mã dùng chung hình (HINH\_251):} \texttt{L09\_C5\_B1\_18}}
\end{tcolorbox}
\vspace{4pt}


\begin{tcolorbox}[
    breakable,
    colback=white,
    colframe=blue!60!black,
    boxrule=0.8pt,
    arc=3pt,
    title={\textbf{Câu \#345} \quad \textbf{MÃ CÂU: \texttt{[L09\_C5\_B1\_32]}} \quad | \quad \textbf{Nhóm: \texttt{HINH\_252}} \quad | \quad \textbf{MapID: \texttt{9H5C1-1}} \hfill \textcolor{yellow!90!white}{[Hình gốc .pyc]}}
]
\textbf{Dạng toán:} Bài toán tổng hợp 4 \hfill \textit{(Tự luận -- Vận dụng cao)}
\vspace{4pt}\hrule\vspace{6pt}
\begin{minipage}[t]{0.58\textwidth}
\vspace{0pt}
Cho hình thoi ${ABCD}$ có cạnh ${AB}=2cm$, hai đường chéo cắt nhau tại ${(O)}$. Gọi $R_{1}, R_{2}$ lần lượt là các bán kính các đường tròn ngoại tiếp các tam giác ${ABC}$ và tam giác ${ABD}$. Giá trị của biểu thức $\dfrac{1}{R_{1}^{2}}+ \dfrac{1}{R_{2}^{2}}$ là \\
\end{minipage}\hfill
\begin{minipage}[t]{0.39\textwidth}
\vspace{0pt}
\centering

\resizebox{\linewidth}{!}{%
\begin{tikzpicture}
    \def\khvuong[size=#1](#2,#3,#4){%
        \draw ($(#3)!#1!(#2)$) --
        ($($(#3)!#1!(#2)$)+($(#3)!#1!(#4)$)-(#3)$)--
        ($(#3)!#1!(#4)$);
    }
    % Định nghĩa các điểm với nhãn
    \coordinate[label=above left:A] (a) at (-4,0);
    \coordinate[label=above:B] (b) at (0,2);
    \coordinate[label=above right:C] (c) at (4,0);
    \coordinate[label=below right:D] (d) at (0,-2);
    \coordinate[label=above right:O] (o) at (0,0);
    \coordinate[label=above left:M] (m) at (-2,1);
    \coordinate[label=below left:K] (k) at (-1.5,0);
    \coordinate[label=below:I] (i) at (0,-3);
    
    % Vẽ các đoạn thẳng giữa các điểm
    \draw (m) -- (i);
    \draw (a) -- (c);
    \draw (b) -- (i);
    \draw (a) -- (b) -- (c) -- (d) -- cycle;
    
\khvuong[size=6pt](b,m,k);
\khvuong[size=6pt](b,o,a);
\end{tikzpicture}
}
\end{minipage}

\vspace{3pt}\hrule\vspace{3pt}\footnotesize\textcolor{gray!80!black}{\textbf{Các mã dùng chung hình (HINH\_252):} \texttt{L09\_C5\_B1\_19}}
\end{tcolorbox}
\vspace{4pt}


\begin{tcolorbox}[
    breakable,
    colback=white,
    colframe=blue!60!black,
    boxrule=0.8pt,
    arc=3pt,
    title={\textbf{Câu \#346} \quad \textbf{MÃ CÂU: \texttt{[L09\_C5\_B2\_02]}} \quad | \quad \textbf{Nhóm: \texttt{HINH\_253}} \quad | \quad \textbf{MapID: \texttt{9H5H2-1}} \hfill \textcolor{yellow!90!white}{[Hình gốc .pyc]}}
]
\textbf{Dạng toán:} Thông hiểu cung, dây cung, đường kính trong tam giác vuông \hfill \textit{(Trắc nghiệm -- Thông hiểu)}
\vspace{4pt}\hrule\vspace{6pt}
\begin{minipage}[t]{0.58\textwidth}
\vspace{0pt}
Cho tam giác ${ABC}$ vuông tại ${A}$ có $AB=3cm$; $AC=4cm$. Khẳng định sai là \\

\vspace{4pt}\noindent
\textbf{\textcolor{blue!75!black}{A. (Đ)}} Bán kính của đường ngoại tiếp tam giác ${ABC}$ là ${5}cm$ \quad \textbf{B.} Khoảng cách từ ${O}$ đến dây ${AC}$ bằng ${\dfrac{3}{2}}cm$ \quad \textbf{C.} Khoảng cách từ ${O}$ đến dây ${AB}$ bằng ${2}cm$ \quad \textbf{D.} Ba điểm ${A, B, C}$ cùng thuộc đường tròn đường kính ${BC}$
\end{minipage}\hfill
\begin{minipage}[t]{0.39\textwidth}
\vspace{0pt}
\centering

\resizebox{\linewidth}{!}{%
\begin{tikzpicture}
    % Định nghĩa các điểm A, B, C
    \coordinate [label=left:A] (A) at (0, 0);
    \coordinate [label=right:B] (B) at (4, 0);
    \coordinate [label=above:C] (C) at (0, 3);
    
    % Vẽ tam giác ABC
    \draw[thick] (A) -- (B) -- (C) -- cycle;
    
    % Đánh dấu các điểm A, B, C
    \fill (A) circle (2pt);
    \fill (B) circle (2pt);
    \fill (C) circle (2pt);
    
    % Vẽ góc vuông tại A
    \draw (0, 0) rectangle (0.4, 0.4);
    
    % Tính trung điểm M của AB và N của AC
    \coordinate [label=below:M] (M) at ($(A)!0.5!(B)$);
    \coordinate [label=left:N] (N) at ($(A)!0.5!(C)$);
    
    % Đánh dấu các điểm M, N
    \fill (M) circle (2pt);
    \fill (N) circle (2pt);
    
    % Vẽ đường tròn đường kính BC
    \coordinate [label=above:O] (O) at (2,3/2);
    \fill (O) circle (2pt);
    \draw[thick] (O) circle ({5/2});
    \draw[thick] (O) -- (M);
    \draw[thick] (O) -- (N);
\end{tikzpicture}
}
\end{minipage}

\end{tcolorbox}
\vspace{4pt}


\begin{tcolorbox}[
    breakable,
    colback=white,
    colframe=blue!60!black,
    boxrule=0.8pt,
    arc=3pt,
    title={\textbf{Câu \#347} \quad \textbf{MÃ CÂU: \texttt{[L09\_C5\_B2\_14]}} \quad | \quad \textbf{Nhóm: \texttt{HINH\_254}} \quad | \quad \textbf{MapID: \texttt{9H5V2-1}} \hfill \textcolor{yellow!90!white}{[Hình gốc .pyc]}}
]
\textbf{Dạng toán:} Bài tập tổng hợp 1 \hfill \textit{(Trắc nghiệm -- Vận dụng)}
\vspace{4pt}\hrule\vspace{6pt}
\begin{minipage}[t]{0.58\textwidth}
\vspace{0pt}
Cho đường tròn ${(O)}$ và hai dây cung ${CD}$, ${EF}$ bằng nhau và vuông góc với nhau tại ${K}$. Giả sử $KE=7cm$, $KF=9cm$. Tính tổng khoảng cách từ tâm ${O}$ đến mỗi dây. \\

\vspace{4pt}\noindent
\textbf{A.} $3$ \quad \textbf{B.} $0$ \quad \textbf{\textcolor{blue!75!black}{C. (Đ)}} $2$ \quad \textbf{D.} $1$
\end{minipage}\hfill
\begin{minipage}[t]{0.39\textwidth}
\vspace{0pt}
\centering

\resizebox{\linewidth}{!}{%
\begin{tikzpicture}
\def\khvuong[size=#1](#2,#3,#4){%
    \draw ($(#3)!#1!(#2)$) --
    ($($(#3)!#1!(#2)$)+($(#3)!#1!(#4)$)-(#3)$)--
    ($(#3)!#1!(#4)$);
}
    % Vẽ đường tròn tâm O
    \coordinate [label=left:$O$] (O) at (0,0);
    \draw [thick] (O) circle(3);
    
    % Đặt điểm K trên đường tròn
    \coordinate [label=above right:$K$] (K) at (1,1);
    
    % Tính toán vị trí các điểm E, F, C, D nằm trên đường tròn
    \coordinate [label=above right:$E$] (E) at (1,2.82);
    \coordinate [label=below right:$F$] (F) at (1,-2.82);
    \coordinate [label=above left:$C$] (C) at (-2.82,1);
    \coordinate [label=above right:$D$] (D) at (2.82,1);
        \coordinate [label=above right:$H$] (H) at (1,0);
            \coordinate [label=above right:$M$] (M) at (0,1);
    % Vẽ các dây cung EF và CD
    \draw [thick] (E) -- (F);
    \draw [thick] (C) -- (D);
        \draw [thick] (O) -- (M);
            \draw [thick] (O) -- (H);
    % Vẽ các điểm và chú thích
    \fill (O) circle(2pt);
    \fill (K) circle(2pt);
    \fill (E) circle(2pt);
    \fill (F) circle(2pt);
    \fill (C) circle(2pt);
    \fill (D) circle(2pt);
        \fill (H) circle(2pt);
            \fill (M) circle(2pt);
            \khvuong[size=6pt](O,M,K);
                \khvuong[size=6pt](O,H,K);
\end{tikzpicture}
}
\end{minipage}

\end{tcolorbox}
\vspace{4pt}


\begin{tcolorbox}[
    breakable,
    colback=white,
    colframe=blue!60!black,
    boxrule=0.8pt,
    arc=3pt,
    title={\textbf{Câu \#348} \quad \textbf{MÃ CÂU: \texttt{[L09\_C5\_B2\_16]}} \quad | \quad \textbf{Nhóm: \texttt{HINH\_255}} \quad | \quad \textbf{MapID: \texttt{9H5C2-1}} \hfill \textcolor{yellow!90!white}{[Hình gốc .pyc]}}
]
\textbf{Dạng toán:} Bài tập tổng hợp 3 \hfill \textit{(Trắc nghiệm -- Vận dụng cao)}
\vspace{4pt}\hrule\vspace{6pt}
\begin{minipage}[t]{0.58\textwidth}
\vspace{0pt}
Cho nửa đường tròn ${(O)}$. Hai dây cung ${AB}$ và ${CD}$ song song với nhau có độ dài lần lượt là ${60}cm$ và ${144}cm$ và khoảng cách giữa hai dây là ${42}cm$. Tính bán kính đường tròn.\\

\vspace{4pt}\noindent
\textbf{A.} ${80} cm$ \quad \textbf{B.} ${77} cm$ \quad \textbf{C.} ${79} cm$ \quad \textbf{\textcolor{blue!75!black}{D. (Đ)}} ${78} cm$
\end{minipage}\hfill
\begin{minipage}[t]{0.39\textwidth}
\vspace{0pt}
\centering

\resizebox{\linewidth}{!}{%
\begin{tikzpicture}
    \def\khvuong[size=#1](#2,#3,#4){%
        \draw ($(#3)!#1!(#2)$) --
        ($($(#3)!#1!(#2)$)+($(#3)!#1!(#4)$)-(#3)$)--
        ($(#3)!#1!(#4)$);
    }
    
    \coordinate (o) at (0,2);
    % Vẽ nửa đường tròn tâm O trong lưới
    \draw (3,2) arc (0:180:3cm); % Vẽ nửa đường tròn với tâm tại (0,4) và bán kính 2cm
    \filldraw (0,2) circle (1pt) node[below] {$O$}; % Vẽ điểm O tại tâm và ghi nhãn
    \coordinate (c) at (-2.85,2.99);
    \coordinate (d) at (2.85,2.99);
    \coordinate (a) at (-1.65,4.53);
    \coordinate (b) at (1.65,4.53);
    \coordinate (h) at (0,4.53);
    \coordinate (k) at (0,2.99);
    \draw (o)--(d)  (a)--(b);
    \draw (o)--(b)  (c)--(d);
    \draw (o)--(h)  (o)--(k);
    \fill (a) circle (1pt) node[above left] {$A$}; 
    \fill (b) circle (1pt) node[above right] {$B$}; 
    \fill (c) circle (1pt) node[above left] {$C$}; 
    \fill (d) circle (1pt) node[above right] {$D$};
    \fill (k) circle (1pt) node[above left] {$K$};
    \fill (h) circle (1pt) node[above left] {$H$};
    \khvuong[size=6pt](d,k,h);
    \khvuong[size=6pt](b,h,o);
    \end{tikzpicture}
}
\end{minipage}

\end{tcolorbox}
\vspace{4pt}


\begin{tcolorbox}[
    breakable,
    colback=white,
    colframe=blue!60!black,
    boxrule=0.8pt,
    arc=3pt,
    title={\textbf{Câu \#349} \quad \textbf{MÃ CÂU: \texttt{[L09\_C5\_B2\_17]}} \quad | \quad \textbf{Nhóm: \texttt{HINH\_256}} \quad | \quad \textbf{MapID: \texttt{9H5N2-1}} \hfill \textcolor{yellow!90!white}{[Hình gốc .pyc]}}
]
\textbf{Dạng toán:} Cho số đo góc ở tâm tính số đo cung bị chắn. \hfill \textit{(Trắc nghiệm -- Nhận biết)}
\vspace{4pt}\hrule\vspace{6pt}
\begin{minipage}[t]{0.58\textwidth}
\vspace{0pt}
Trong hình vẽ sau biết $\widehat{COD}=138^{\circ}$. Số đo cung nhỏ $\stackrel\frown{CD}$ bằng \\

\vspace{4pt}\noindent
\textbf{A.} ${148}^{\circ}$ \quad \textbf{B.} ${143}^{\circ}$ \quad \textbf{C.} ${133}^{\circ}$ \quad \textbf{\textcolor{blue!75!black}{D. (Đ)}} ${138}^{\circ}$
\end{minipage}\hfill
\begin{minipage}[t]{0.39\textwidth}
\vspace{0pt}
\centering

\resizebox{\linewidth}{!}{%
\begin{tikzpicture}
    % Định nghĩa tâm của đường tròn
    \coordinate (o) at (0,2);
    
    % Vẽ đường tròn đầy đủ với tâm O và bán kính 3cm, nét vẽ đậm
    \draw[line width=1.2pt] (o) circle (3cm);
    
    % Vẽ các đoạn thẳng từ tâm O tới các điểm trên đường tròn
    \coordinate (c) at (-2.85,2.99);  % Điểm C trên đường tròn
    \coordinate (d) at (2.85,2.99);   % Điểm D trên đường tròn
    
    \draw[line width=1.3pt] (o) -- (d);
    \draw[line width=1.3pt] (o) -- (c);
    \draw[line width=1.3pt] (c) -- (d);
    
    % Vẽ và ghi nhãn các điểm O, C, D
    \filldraw (o) circle (1pt) node[below] {$O$}; 
    \fill (c) circle (1pt) node[above left] {$C$}; 
    \fill (d) circle (1pt) node[above right] {$D$};
    \end{tikzpicture}
}
\end{minipage}

\end{tcolorbox}
\vspace{4pt}


\begin{tcolorbox}[
    breakable,
    colback=white,
    colframe=blue!60!black,
    boxrule=0.8pt,
    arc=3pt,
    title={\textbf{Câu \#350} \quad \textbf{MÃ CÂU: \texttt{[L09\_C5\_B2\_18]}} \quad | \quad \textbf{Nhóm: \texttt{HINH\_257}} \quad | \quad \textbf{MapID: \texttt{9H5N2-1}} \hfill \textcolor{yellow!90!white}{[Hình gốc .pyc]}}
]
\textbf{Dạng toán:} Cho số đo của góc nội tiếp tính số đo của cung bị chắn. \hfill \textit{(Trắc nghiệm -- Nhận biết)}
\vspace{4pt}\hrule\vspace{6pt}
\begin{minipage}[t]{0.58\textwidth}
\vspace{0pt}
Trong hình vẽ sau, biết $\widehat{ABC}=33^{\circ}$. Số đo cung nhỏ $\stackrel\frown{AC}$ bằng \\

\vspace{4pt}\noindent
\textbf{A.} ${33}^{\circ}$ \quad \textbf{B.} ${76}^{\circ}$ \quad \textbf{\textcolor{blue!75!black}{C. (Đ)}} ${66}^{\circ}$ \quad \textbf{D.} ${71}^{\circ}$
\end{minipage}\hfill
\begin{minipage}[t]{0.39\textwidth}
\vspace{0pt}
\centering

\resizebox{\linewidth}{!}{%
\begin{tikzpicture}
            % Định nghĩa tâm của đường tròn
            \coordinate (o) at (0,2);
            
            % Vẽ đường tròn đầy đủ với tâm O và bán kính 3cm, nét vẽ đậm
            \draw[line width=1.2pt] (o) circle (3cm);
            
            % Vẽ các đoạn thẳng từ tâm O tới các điểm trên đường tròn
            \coordinate (c) at (-2.85,2.99);  % Điểm C trên đường tròn
            \coordinate (b) at (2.85,2.99);   % Điểm D trên đường tròn
            \coordinate (a) at (-1.65,4.53);
            \draw[line width=1.3pt] (a) -- (b);
            \draw[line width=1.3pt] (a) -- (c);
            \draw[line width=1.3pt] (c) -- (b);
            
            % Vẽ và ghi nhãn các điểm O, C, D
            \filldraw (o) circle (1pt) node[below] {$O$}; 
            \fill (c) circle (1pt) node[above left] {$C$}; 
            \fill (b) circle (1pt) node[above right] {$B$};
            \fill (a) circle (1pt) node[above left] {$A$}; 
        \end{tikzpicture}
}
\end{minipage}

\end{tcolorbox}
\vspace{4pt}


\begin{tcolorbox}[
    breakable,
    colback=white,
    colframe=blue!60!black,
    boxrule=0.8pt,
    arc=3pt,
    title={\textbf{Câu \#351} \quad \textbf{MÃ CÂU: \texttt{[L09\_C5\_B2\_24]}} \quad | \quad \textbf{Nhóm: \texttt{HINH\_258}} \quad | \quad \textbf{MapID: \texttt{9H5C2-1}} \hfill \textcolor{yellow!90!white}{[Hình gốc .pyc]}}
]
\textbf{Dạng toán:} Tính số đo của góc nội tiếp \hfill \textit{(Trắc nghiệm -- Vận dụng cao)}
\vspace{4pt}\hrule\vspace{6pt}
\begin{minipage}[t]{0.58\textwidth}
\vspace{0pt}
Cho tam giác ${ABC}$ vuông tại ${A}$ biết $\widehat{B}=14\widehat{C}$. Đường tròn ${(O)}$ tiếp xúc với cả ba cạnh ${BC,CA, AB}$ của tam giác ${ABC}$ lần lượt tại ${D, E, F}$. Tính $sd\stackrel\frown{EF}:sd\stackrel\frown{DF}:sd\stackrel\frown{DE}$\\

\vspace{4pt}\noindent
\textbf{A.} $17:16:29$ \quad \textbf{\textcolor{blue!75!black}{B. (Đ)}} $15:16:29$ \quad \textbf{C.} $15:17:29$ \quad \textbf{D.} $16:16:29$
\end{minipage}\hfill
\begin{minipage}[t]{0.39\textwidth}
\vspace{0pt}
\centering

\resizebox{\linewidth}{!}{%
\begin{tikzpicture}
    % Vẽ lưới
    
    % Định nghĩa các điểm A, B, C, O, D, E, F
    \coordinate [label=below left:{$A$}] (a) at (-2,-3); 
    \coordinate [label=below right:{$B$}] (b) at (2,-3); 
    \coordinate [label=above left:{$C$}] (c) at (-2,3); 
    \coordinate [label=above left:{$O$}] (o) at (-0.61,-1.61); 
    \coordinate [label=above left:{$F$}] (f) at (-2,-1.61); 
    \coordinate [label=above right:{$D$}] (d) at (0.55,-0.83); 
    \coordinate [label=below:{$E$}] (e) at (-0.61,-3); 
    
    % Vẽ tam giác ABC
    \draw[thick] (a)--(b)--(c)--cycle; 
    
    % Vẽ các điểm
    \fill (a) circle (1.2pt); % Đánh dấu điểm A
    \fill (b) circle (1.2pt); % Đánh dấu điểm B
    \fill (c) circle (1.2pt); % Đánh dấu điểm C
    \fill (o) circle (1.2pt); % Đánh dấu điểm O
    \fill (f) circle (1.2pt); % Đánh dấu điểm F
    \fill (d) circle (1.2pt); % Đánh dấu điểm D
    \fill (e) circle (1.2pt); % Đánh dấu điểm E
    
    % Vẽ đường tròn nội tiếp tam giác ABC
    \draw [black] (o) circle [radius=1.39cm]; 
    \def\khvuong[size=#1](#2,#3,#4){%
        \draw ($(#3)!#1!(#2)$) --
        ($($(#3)!#1!(#2)$)+($(#3)!#1!(#4)$)-(#3)$)--
        ($(#3)!#1!(#4)$);
    }
    \khvuong[size=5pt](c,a,b);%
    \end{tikzpicture}
}
\end{minipage}

\end{tcolorbox}
\vspace{4pt}


\begin{tcolorbox}[
    breakable,
    colback=white,
    colframe=blue!60!black,
    boxrule=0.8pt,
    arc=3pt,
    title={\textbf{Câu \#352} \quad \textbf{MÃ CÂU: \texttt{[L09\_C5\_B2\_28]}} \quad | \quad \textbf{Nhóm: \texttt{HINH\_259}} \quad | \quad \textbf{MapID: \texttt{9H5N2-1}} \hfill \textcolor{yellow!90!white}{[Hình gốc .pyc]}}
]
\textbf{Dạng toán:} Tổng hợp 1 \hfill \textit{(Tự luận -- Nhận biết)}
\vspace{4pt}\hrule\vspace{6pt}
\begin{minipage}[t]{0.58\textwidth}
\vspace{0pt}
Hai góc đối diện ${G}$ và ${F}$ của tứ giác ${GHFE}$ là những góc vuông.

a) Chứng minh rằng $GF\leq HE$.

b) Nếu $GF=HE$. Chứng minh rằng tứ giác ${GHFE}$ là hình chữ nhật.
\end{minipage}\hfill
\begin{minipage}[t]{0.39\textwidth}
\vspace{0pt}
\centering

\resizebox{\linewidth}{!}{%
\begin{tikzpicture}[scale=1, font=\footnotesize, line join=round, line cap=round, >=stealth] 
    \coordinate[label=above:{$I$}] (O) at (0,0);
    \coordinate[label=below:{$F$}] (C) at ($(0,0)+(-65:2cm)$);
    \coordinate[label=left:{$H$}] (B) at ($(0,0)+(180:2cm)$);
    \coordinate[label=above:{$G$}] (A) at ($(0,0)+(100:2cm)$);
    \coordinate[label=right:{$E$}] (D) at ($(0,0)+(0:2cm)$);
    \draw (0,0) circle[radius=2cm] (A)--(B)--(C)--(D)--(A) (B)--(D);
    \foreach \a/\b/\c in {B/A/D,B/C/D} \draw pic[angle radius=2mm,draw=black] {right angle =\a--\b--\c};
    \foreach \diem in {A,B,C,D,O}\fill (\diem)circle(1pt); 
\end{tikzpicture}
}
\end{minipage}

\end{tcolorbox}
\vspace{4pt}


\begin{tcolorbox}[
    breakable,
    colback=white,
    colframe=blue!60!black,
    boxrule=0.8pt,
    arc=3pt,
    title={\textbf{Câu \#353} \quad \textbf{MÃ CÂU: \texttt{[L09\_C5\_B2\_29]}} \quad | \quad \textbf{Nhóm: \texttt{HINH\_260}} \quad | \quad \textbf{MapID: \texttt{9H5V2-1}} \hfill \textcolor{yellow!90!white}{[Hình gốc .pyc]}}
]
\textbf{Dạng toán:} Tổng hợp 2 \hfill \textit{(Tự luận -- Vận dụng)}
\vspace{4pt}\hrule\vspace{6pt}
\begin{minipage}[t]{0.58\textwidth}
\vspace{0pt}
Cho tam giác cân ${MNE}$ có $\widehat{M}=49^{\circ}$. Đường tròn đường kính ${MN}$ cắt ${NE}$ tại ${F}$ (khác ${N}$).

 a) Chứng minh rằng ${F}$ là trung điểm của đoạn ${NE}$.

 b) Tính số đo của cung nhỏ ${MF}$.
\end{minipage}\hfill
\begin{minipage}[t]{0.39\textwidth}
\vspace{0pt}
\centering

\resizebox{\linewidth}{!}{%
\begin{tikzpicture}[scale=1, font=\footnotesize, line join=round, line cap=round, >=stealth] 
    \coordinate[label=below:{$E$}] (C) at ($(0,0)+(-70:3cm)$);
    \coordinate[label=below:{$N$}] (B) at ($(0,0)+(-110:3cm)$);
    \coordinate[label=above:{$M$}] (A) at (0,0);
    \coordinate[label=left:{$I$}] (O) at ($(A)!0.5!(B)$);
    \coordinate[label=below:{$F$}] (D) at ($(C)!0.5!(B)$);
    \draw let \p1=($(O)-(A)$) in (O) circle ({veclen(\x1,\y1)});
    \draw  (A)--(B)--(C)--(A)--(D)--(O);
    \foreach \diem in {A,B,C,D,O}\fill (\diem)circle(1pt);  
\end{tikzpicture}
}
\end{minipage}

\end{tcolorbox}
\vspace{4pt}


\begin{tcolorbox}[
    breakable,
    colback=white,
    colframe=blue!60!black,
    boxrule=0.8pt,
    arc=3pt,
    title={\textbf{Câu \#354} \quad \textbf{MÃ CÂU: \texttt{[L09\_C5\_B2\_30]}} \quad | \quad \textbf{Nhóm: \texttt{HINH\_261}} \quad | \quad \textbf{MapID: \texttt{9H5V2-1}} \hfill \textcolor{yellow!90!white}{[Hình gốc .pyc]}}
]
\textbf{Dạng toán:} Tổng hợp 3 \hfill \textit{(Tự luận -- Vận dụng)}
\vspace{4pt}\hrule\vspace{6pt}
\begin{minipage}[t]{0.58\textwidth}
\vspace{0pt}
Cho tam giác nhọn ${EFK}$, hai đường cao ${FH}$ và ${KM}$ cắt nhau tại ${I}$. Chứng minh rằng:

a) Bốn điểm ${E,M,I,H}$ cùng thuộc một đường tròn.

b) $EI> HM$.
\end{minipage}\hfill
\begin{minipage}[t]{0.39\textwidth}
\vspace{0pt}
\centering

\resizebox{\linewidth}{!}{%
\begin{tikzpicture}[scale=1, font=\footnotesize, line join=round, line cap=round, >=stealth] 
    \coordinate[label=above:{$O$}] (O) at (0,0);
    \coordinate[label=right:{$K$}] (C) at ($(0,0)+(0:2cm)$);
    \coordinate[label=left:{$F$}] (B) at ($(0,0)+(180:2cm)$);
    \coordinate[label=above left:{$M$}] (E) at ($(0,0)+(130:2cm)$);
    \coordinate[label=above right:{$H$}] (D) at ($(0,0)+(70:2cm)$);
    \coordinate[label=above:{$E$}] (A) at (intersection of B--E and D--C); 
    \coordinate[label=below:{$I$}] (H) at (intersection of B--D and E--C); 
    \coordinate[label=above right:{$N$}] (I) at ($(A)!0.5!(H)$);
    \draw (0,0) circle[radius=2cm] (A)--(B)--(C)--(D)--(A)--(H) (B)--(D)--(E) (C)--(E)--(I)--(D);
    \foreach \diem in {A,B,C,D,O,E,H,I}\fill (\diem)circle(1pt); 
    \foreach \a/\b/\c in {B/E/C,B/D/C} \draw pic[angle radius=2mm,draw=black] {right angle =\a--\b--\c};
\end{tikzpicture}
}
\end{minipage}

\end{tcolorbox}
\vspace{4pt}


\begin{tcolorbox}[
    breakable,
    colback=white,
    colframe=blue!60!black,
    boxrule=0.8pt,
    arc=3pt,
    title={\textbf{Câu \#355} \quad \textbf{MÃ CÂU: \texttt{[L09\_C5\_B2\_31]}} \quad | \quad \textbf{Nhóm: \texttt{HINH\_262}} \quad | \quad \textbf{MapID: \texttt{9H5V2-1}} \hfill \textcolor{yellow!90!white}{[Hình gốc .pyc]}}
]
\textbf{Dạng toán:} Tổng hợp 4 \hfill \textit{(Tự luận -- Vận dụng)}
\vspace{4pt}\hrule\vspace{6pt}
\begin{minipage}[t]{0.58\textwidth}
\vspace{0pt}
Cho tam giác cân ${EFK}$ $(EF=EK)$. Gọi $({O})$ là đường tròn đi qua ba điểm ${E}$, ${F, K}$ và ${M}$ là điểm trên cung nhỏ ${FK}$ sao cho $\stackrel\frown{FM}=\stackrel\frown{MK}$.

a) Chứng minh rằng ba điểm ${E, O, M}$ thẳng hàng.

b) Gọi ${I}$ là chân đường cao hạ từ ${E}$ xuống ${FK}$. Chứng minh rằng $EI< EF< EM$.
\end{minipage}\hfill
\begin{minipage}[t]{0.39\textwidth}
\vspace{0pt}
\centering

\resizebox{\linewidth}{!}{%
\begin{tikzpicture}[scale=1, font=\footnotesize, line join=round, line cap=round, >=stealth] 
    \coordinate[label=above right:{$O$}] (O) at (0,0);
    \coordinate[label=right:{$K$}] (C) at ($(0,0)+(-30:2cm)$);
    \coordinate[label=below:{$M$}] (E) at ($(0,0)+(-90:2cm)$);
    \coordinate[label=left:{$F$}] (B) at ($(0,0)+(-150:2cm)$);
    \coordinate[label=above:{$E$}] (A) at ($(0,0)+(90:2cm)$);
    \coordinate[label=below right:{$I$}] (H) at ($(B)!(O)!(C)$);
    \draw (0,0) circle[radius=2cm] (A)--(B)--(O)--(C)--(A)--(E) (O)--(H) (B)--(C);
    \foreach \diem in {A,B,C,O,H,E} \fill (\diem) circle(1pt); 
    \foreach \a/\b/\c in {O/H/C} \draw pic[angle radius=2mm, draw=black] {right angle=\a--\b--\c};
\end{tikzpicture}
}
\end{minipage}

\end{tcolorbox}
\vspace{4pt}


\begin{tcolorbox}[
    breakable,
    colback=white,
    colframe=blue!60!black,
    boxrule=0.8pt,
    arc=3pt,
    title={\textbf{Câu \#356} \quad \textbf{MÃ CÂU: \texttt{[L09\_C5\_B3\_03]}} \quad | \quad \textbf{Nhóm: \texttt{HINH\_263}} \quad | \quad \textbf{MapID: \texttt{9H5N3-1}} \hfill \textcolor{yellow!90!white}{[Hình gốc .pyc]}}
]
\textbf{Dạng toán:} Tính độ dài cung tròn của miếng Pizza \hfill \textit{(Trắc nghiệm -- Nhận biết)}
\vspace{4pt}\hrule\vspace{6pt}
\begin{minipage}[t]{0.58\textwidth}
\vspace{0pt}
Một miếng bánh pizza có bán kính $10$ cm và số đo góc như hình vẽ. Độ dài cung tròn của miếng Pizza là

\vspace{4pt}\noindent
\textbf{A.} $2,8\pi$ cm \quad \textbf{B.} $1,3\pi$ cm \quad \textbf{\textcolor{blue!75!black}{C. (Đ)}} ${0,8}\pi$ cm \quad \textbf{D.} $2,0\pi$ cm
\end{minipage}\hfill
\begin{minipage}[t]{0.39\textwidth}
\vspace{0pt}
\centering

\resizebox{\linewidth}{!}{%
\begin{tikzpicture}
       % Đặt các điểm A, B, C trên đường tròn
            \coordinate [label=above left:$A$] (A) at (120:3);
            \coordinate [label=below left:$B$] (B) at (170:3);
            \coordinate [label=right:$O$] (O) at (0,0);
            % Vẽ cung AB
    \draw [thick] (A) arc[start angle=120, end angle=170, radius=3];
     % Tô màu vàng cho phần cung AB
    \fill[yellow, opacity=0.3] (O) -- (A) arc[start angle=120, end angle=170, radius=3] -- cycle;
            % Vẽ tam giác ABC
            \draw [thick] (A) -- (O) -- (B);
            \draw 
        (O) pic[draw,angle radius = 13] {angle = A--O--B} node[shift={(-30:-20pt)}]{$ 15^\circ $}   
        ;
            % Vẽ các điểm
            \fill (A) circle (2pt);
            \fill (B) circle (2pt);
            \fill (O) circle (2pt);
        \end{tikzpicture}
}
\end{minipage}

\end{tcolorbox}
\vspace{4pt}


\begin{tcolorbox}[
    breakable,
    colback=white,
    colframe=blue!60!black,
    boxrule=0.8pt,
    arc=3pt,
    title={\textbf{Câu \#357} \quad \textbf{MÃ CÂU: \texttt{[L09\_C5\_B3\_14]}} \quad | \quad \textbf{Nhóm: \texttt{HINH\_264}} \quad | \quad \textbf{MapID: \texttt{9H5H3-7}} \hfill \textcolor{yellow!90!white}{[Hình gốc .pyc]}}
]
\textbf{Dạng toán:} Bài toán thực tế. Tính diện tích cái quạt \hfill \textit{(Trắc nghiệm -- Thông hiểu)}
\vspace{4pt}\hrule\vspace{6pt}
\begin{minipage}[t]{0.58\textwidth}
\vspace{0pt}
Hình quạt tô màu đỏ ở hình vẽ bên dưới có bán kính bằng $7$ dm và góc ở tâm bằng $120^{\circ}$. Tính diện tích của hình quạt đó.

\vspace{4pt}\noindent
\textbf{A.} $ \dfrac{191}{12} \pi dm^{2}$ \quad \textbf{B.} $ \dfrac{5897}{360} \pi dm^{2}$ \quad \textbf{\textcolor{blue!75!black}{C. (Đ)}} $ \dfrac{49}{3} \pi cm^{2}$ \quad \textbf{D.} $ \dfrac{1909}{120} \pi dm^{2}$
\end{minipage}\hfill
\begin{minipage}[t]{0.39\textwidth}
\vspace{0pt}
\centering

\resizebox{\linewidth}{!}{%
\begin{tikzpicture}
    
    % Vẽ hình quạt đỏ với góc ở tâm là 150 độ và trục tung làm trục đối xứng
    \fill[red] (0,0) -- (4,0) arc[start angle=0, end angle=150, radius=4cm] -- cycle;
    
    % Vẽ tay cầm ở dưới và đối xứng qua trục tung
    \fill[orange] (0,-0.2) -- (0.5,-0.8) arc[start angle=-60, end angle=-120, radius=0.5cm] -- cycle;
    
\end{tikzpicture}
}
\end{minipage}

\vspace{3pt}\hrule\vspace{3pt}\footnotesize\textcolor{gray!80!black}{\textbf{Các mã dùng chung hình (HINH\_264):} \texttt{L09\_C5\_B3\_44}}
\end{tcolorbox}
\vspace{4pt}


\begin{tcolorbox}[
    breakable,
    colback=white,
    colframe=blue!60!black,
    boxrule=0.8pt,
    arc=3pt,
    title={\textbf{Câu \#358} \quad \textbf{MÃ CÂU: \texttt{[L09\_C5\_B3\_15]}} \quad | \quad \textbf{Nhóm: \texttt{HINH\_265}} \quad | \quad \textbf{MapID: \texttt{9H5V3-7}} \hfill \textcolor{yellow!90!white}{[Hình gốc .pyc]}}
]
\textbf{Dạng toán:} Bài toán thực tế liên quan đến quãng đường đi của xe \hfill \textit{(Trắc nghiệm -- Vận dụng)}
\vspace{4pt}\hrule\vspace{6pt}
\begin{minipage}[t]{0.58\textwidth}
\vspace{0pt}
Bánh xe (khi bơm căng) của một chiếc xe đạp có đường kính $650$ mm. Biết rằng khi giò đĩa quay một vòng thì bánh xe đạp quay được khoảng $4$ vòng (hình vẽ). Hỏi chiếc xe đạp di chuyển được quãng đường dài bao nhiêu mét sau khi người đi xe đạp $27$ vòng liên tục?

\vspace{4pt}\noindent
\textbf{A.} $ 223,54 m$ \quad \textbf{\textcolor{blue!75!black}{B. (Đ)}} $ 220,54  m$ \quad \textbf{C.} $ 225,54 m$ \quad \textbf{D.} $ 216,54 m$
\end{minipage}\hfill
\begin{minipage}[t]{0.39\textwidth}
\vspace{0pt}
\centering

\resizebox{\linewidth}{!}{%
\begin{tikzpicture}[scale=0.8, line cap=round, line join=round, thick]
    
    % Bánh xe sau
    \draw[thick] (0,0) circle(1.5cm);             % Bánh xe lớn
    \draw[fill=black] (0,0) circle(0.1cm);         % Tâm bánh xe
    % Vẽ nan hoa cho bánh xe sau
    \draw (0,0) -- (0,1.5);
    \draw (0,0) -- (0,-1.5);
    \draw (0,0) -- (1.3,0.75);
    \draw (0,0) -- (-1.3,-0.75);
    \draw (0,0) -- (1.3,-0.75);
    \draw (0,0) -- (-1.3,0.75);
    \draw (0,0) -- (0.75,1.3);
    \draw (0,0) -- (-0.75,-1.3);
    \draw (0,0) -- (0.75,-1.3);
    \draw (0,0) -- (-0.75,1.3);
    
    % Bánh xe trước
    \draw[thick] (8,0) circle(1.5cm);             % Bánh xe lớn
    \draw[fill=black] (8,0) circle(0.1cm);         % Tâm bánh xe
    % Vẽ nan hoa cho bánh xe trước
    \draw (8,0) -- (8,1.5);
    \draw (8,0) -- (8,-1.5);
    \draw (8,0) -- (9.3,0.75);
    \draw (8,0) -- (6.7,-0.75);
    \draw (8,0) -- (9.3,-0.75);
    \draw (8,0) -- (6.7,0.75);
    \draw (8,0) -- (8.75,1.3);
    \draw (8,0) -- (7.25,-1.3);
    \draw (8,0) -- (8.75,-1.3);
    \draw (8,0) -- (7.25,1.3);
    
    % Khung xe
    \draw[thick,blue] (0,0) -- (3,2.5);            % Khung từ bánh sau đến khung trên
    \draw[thick,blue] (3,2.5) -- (7,2.5);          % Khung trên từ yên đến bánh trước
    \draw[thick,blue] (7,2.5) -- (8,0);            % Khung từ khung trên xuống bánh trước
    \draw[thick,blue] (0,0) -- (8,0);              % Thanh dưới khung xe
    \draw[thick,blue] (3,2.5) -- (5.5,0);          % Thanh nối từ yên xuống giò đĩa
    
    % Ghi đông và tay lái
    \draw[thick,blue] (7,2.5) -- (7,3.5);          % Thanh đứng nối ghi đông
    \draw[thick,blue] (6.5,3.5) -- (7.5,3.5);      % Tay lái ngang
    \draw[thick,blue] (6.5,3.5) -- (6,3);          % Tay cầm trái
    \draw[thick,blue] (7.5,3.5) -- (8,3);          % Tay cầm phải
    
    % Yên xe
    \draw[thick,blue] (3,2.5) -- (3,2.8);          % Thanh nối yên xe
    \draw[thick,red] (2.7,2.8) -- (3.3,2.8);       % Yên xe
    
    % Giò đĩa và bàn đạp
    \draw[thick,blue] (5.5,0) circle(0.3cm);       % Đĩa xích
    \draw (5.5,0) -- ++(0:0.3cm);                  % Nối giò đĩa ngang
    \draw (5.5,0) -- ++(90:0.3cm);                 % Nối giò đĩa dọc trên
    \draw (5.5,0) -- ++(180:0.3cm);                % Nối giò đĩa ngang trái
    \draw (5.5,0) -- ++(270:0.3cm);                % Nối giò đĩa dọc dưới
    \draw[thick,blue] (5.5,0) -- ++(45:1cm);       % Bàn đạp trái
    \draw[thick,blue] (5.5,0) -- ++(225:1cm);      % Bàn đạp phải
    
\end{tikzpicture}
}
\end{minipage}

\vspace{3pt}\hrule\vspace{3pt}\footnotesize\textcolor{gray!80!black}{\textbf{Các mã dùng chung hình (HINH\_265):} \texttt{L09\_C5\_B3\_45}}
\end{tcolorbox}
\vspace{4pt}


\begin{tcolorbox}[
    breakable,
    colback=white,
    colframe=blue!60!black,
    boxrule=0.8pt,
    arc=3pt,
    title={\textbf{Câu \#359} \quad \textbf{MÃ CÂU: \texttt{[L09\_C5\_B3\_18]}} \quad | \quad \textbf{Nhóm: \texttt{HINH\_266}} \quad | \quad \textbf{MapID: \texttt{9H5C3-1}} \hfill \textcolor{yellow!90!white}{[Hình gốc .pyc]}}
]
\textbf{Dạng toán:} Tính bán kính đường tròn ngoại tiếp tam giác khi biết 2 cạnh và góc xen giữa của tam giác. \hfill \textit{(Trắc nghiệm -- Vận dụng cao)}
\vspace{4pt}\hrule\vspace{6pt}
\begin{minipage}[t]{0.58\textwidth}
\vspace{0pt}
Cho tam giác ${ABC}$ có $AB=AC=100cm$, $\widehat{A} =154^{\circ}$. Khi đó bán kính đường tròn đi qua ba đỉnh của tam giác (làm tròn kết quả đến hàng phần trăm) bằng

\vspace{4pt}\noindent
\textbf{A.} ${221,7} cm$ \quad \textbf{B.} ${230,3} cm$ \quad \textbf{C.} ${224,3} cm$ \quad \textbf{\textcolor{blue!75!black}{D. (Đ)}} ${222,3} cm$
\end{minipage}\hfill
\begin{minipage}[t]{0.39\textwidth}
\vspace{0pt}
\centering

\resizebox{\linewidth}{!}{%
\begin{tikzpicture}
    
    % Tọa độ của các điểm
    \def\r{3}  % Bán kính của đường tròn
    \coordinate (O) at (0,0);         % Tâm O
    \coordinate (A) at (0,\r);        % Điểm A trên đường tròn
    \coordinate (B) at ({\r*sqrt(2)/2},{\r*sqrt(2)/2}); % Điểm B trên đường tròn
    \coordinate (C) at ({-\r*sqrt(2)/2},{\r*sqrt(2)/2}); % Điểm C trên đường tròn
    \coordinate (I) at (0,\r/2);      % Trung điểm của OA là điểm I
    
    % Vẽ đường tròn
    \draw (O) circle (\r);
    
    % Vẽ các đoạn thẳng
    \draw (O) -- (A);  % OA
    \draw (C) -- (A);  % CA
    \draw (B) -- (A);  % BA
    \draw (C) -- (B);  % CB
    \draw (C) -- (O);  % CO
    \draw (B) -- (O);  % BO
    
    % Đánh dấu các điểm
    \node[above] at (A) {A};
    \node[below] at (O) {O};
    \node[left] at (C) {C};
    \node[right] at (B) {B};
    \node[right] at (I) {I};
    
    % Vẽ các dấu hiệu góc bằng nhau
     \draw 
    (A) pic[draw,angle radius = 10] {{angle = O--A--B}} ;
 \draw 
(A) pic[draw,angle radius = 10] {{angle = C--A--O}} ;
    
\end{tikzpicture}
}
\end{minipage}

\vspace{3pt}\hrule\vspace{3pt}\footnotesize\textcolor{gray!80!black}{\textbf{Các mã dùng chung hình (HINH\_266):} \texttt{L09\_C5\_B3\_48}}
\end{tcolorbox}
\vspace{4pt}


\begin{tcolorbox}[
    breakable,
    colback=white,
    colframe=blue!60!black,
    boxrule=0.8pt,
    arc=3pt,
    title={\textbf{Câu \#360} \quad \textbf{MÃ CÂU: \texttt{[L09\_C5\_B3\_23]}} \quad | \quad \textbf{Nhóm: \texttt{HINH\_267}} \quad | \quad \textbf{MapID: \texttt{9H5V3-7}} \hfill \textcolor{yellow!90!white}{[Hình gốc .pyc]}}
]
\textbf{Dạng toán:} Bài toán thực tế tính diện tích sân vận động \hfill \textit{(Trắc nghiệm -- Vận dụng)}
\vspace{4pt}\hrule\vspace{6pt}
\begin{minipage}[t]{0.58\textwidth}
\vspace{0pt}
Một sân vận động có đường chạy đua như hình vẽ. Đường chạy gồm các đoạn ${AB}$, ${CD}$ là chiều dài của hình chữ nhật và hai nửa đường tròn. Biết ${AB}=130 (m)$, $CB= 120 (m)$. Tính diện tích sân vận động? (Làm tròn kết quả đến chữ số thập phân thứ nhất)

\vspace{4pt}\noindent
\textbf{\textcolor{blue!75!black}{A. (Đ)}} $26909,7 (m^{2})$ \quad \textbf{B.} $26919,7 (m^{2})$ \quad \textbf{C.} $26900,7 (m^{2})$ \quad \textbf{D.} $26911,7 (m^{2})$
\end{minipage}\hfill
\begin{minipage}[t]{0.39\textwidth}
\vspace{0pt}
\centering

\resizebox{\linewidth}{!}{%
\begin{tikzpicture}
    % Các tọa độ
    \coordinate (A) at (0,2);  % Điểm A
    \coordinate (B) at (6,2);  % Điểm B
    \coordinate (D) at (0,0);  % Điểm D
    \coordinate (C) at (6,0);  % Điểm C
    \coordinate (center1) at (0,1); % Tâm đường tròn bên trái
    \coordinate (center2) at (6,1); % Tâm đường tròn bên phải
    
    % Vẽ hình chữ nhật
    \draw (A) -- (B) -- (C) -- (D) -- cycle;
    
    % Vẽ nửa hình tròn bên trái
    \draw (A) arc[start angle=90, end angle=270, radius=1cm];
    
    % Vẽ nửa hình tròn bên phải
    \draw (B) arc[start angle=90, end angle=-90, radius=1cm];
    
    % Đánh dấu các điểm
    \filldraw (A) circle (1pt) node[above] {$A$};
    \filldraw (B) circle (1pt) node[above] {$B$};
    \filldraw (C) circle (1pt) node[below] {$C$};
    \filldraw (D) circle (1pt) node[below] {$D$};
\end{tikzpicture}
}
\end{minipage}

\vspace{3pt}\hrule\vspace{3pt}\footnotesize\textcolor{gray!80!black}{\textbf{Các mã dùng chung hình (HINH\_267):} \texttt{L09\_C5\_B3\_53}}
\end{tcolorbox}
\vspace{4pt}


\begin{tcolorbox}[
    breakable,
    colback=white,
    colframe=blue!60!black,
    boxrule=0.8pt,
    arc=3pt,
    title={\textbf{Câu \#361} \quad \textbf{MÃ CÂU: \texttt{[L09\_C5\_B3\_24]}} \quad | \quad \textbf{Nhóm: \texttt{HINH\_268}} \quad | \quad \textbf{MapID: \texttt{9H5V3-1}} \hfill \textcolor{yellow!90!white}{[Hình gốc .pyc]}}
]
\textbf{Dạng toán:} Tính diện tích các hình phức tạp \hfill \textit{(Trắc nghiệm -- Vận dụng)}
\vspace{4pt}\hrule\vspace{6pt}
\begin{minipage}[t]{0.58\textwidth}
\vspace{0pt}
Tính diện tích phần kẻ sọc ở hình sau, giới hạn bởi nửa đường tròn đường kính ${LM}$ và hai nửa đường tròn có đường kính tương ứng là $NM={36}$ cm và $LN= {18}$ cm.

\vspace{4pt}\noindent
\textbf{A.} $241 \pi (cm^{2})$ \quad \textbf{B.} $244 \pi (cm^{2})$ \quad \textbf{\textcolor{blue!75!black}{C. (Đ)}} $243 \pi (cm^{2})$ \quad \textbf{D.} $254 \pi (cm^{2})$
\end{minipage}\hfill
\begin{minipage}[t]{0.39\textwidth}
\vspace{0pt}
\centering

\resizebox{\linewidth}{!}{%
\begin{tikzpicture}
    
    % Khởi tạo các điểm
    \coordinate (L) at (0, 0);           % Điểm L
    \coordinate (N) at (2, 0);           % Điểm N
    \coordinate (M) at (6, 0);          % Điểm M
    \coordinate (centerLN) at (1, 0);    % Trung điểm đoạn LN (tâm của nửa đường tròn mới phía dưới)
    
    % Vẽ nửa đường tròn lớn với đường kính 12 cm (phía trên)
    \draw (L) arc[start angle=180, end angle=0, radius=3cm];
    
    % Vẽ nửa đường tròn nhỏ với đường kính 8 cm (phía trên)
    \draw (N) arc[start angle=180, end angle=0, radius=2cm];
    
    % Vẽ đoạn thẳng đáy LM
    \draw (L) -- (M);
        \draw (2,0) arc[start angle=0, end angle=-180, radius=1cm];
    % Tô phần gạch sọc chỉ thuộc nửa đường tròn lớn mà không thuộc nửa đường tròn nhỏ
    \begin{scope}
        % Giới hạn vùng tô chỉ thuộc nửa đường tròn lớn mà không thuộc nửa đường tròn nhỏ
        \clip (L) arc[start angle=180, end angle=0, radius=3cm] -- (M) -- (L) -- cycle;
        \fill[pattern=north east lines] (L) arc[start angle=180, end angle=0, radius=3cm];
        
        % Loại bỏ vùng thuộc nửa đường tròn nhỏ ra khỏi vùng tô bằng một vùng trắng
        \fill[white] (N) arc[start angle=180, end angle=0, radius=2cm] -- (N) -- cycle;
    \end{scope}
    
    % Vẽ nửa đường tròn đối diện phía dưới đoạn LN
    % Tâm của nửa đường tròn là trung điểm của đoạn LN, với bán kính 2 cm

    
    % Đánh dấu các điểm
    \filldraw (L) circle (1pt) node[below left] {$L$};        % Điểm L
    \filldraw (N) circle (1pt) node[below] {$N$};             % Điểm N
    \filldraw (M) circle (1pt) node[below right] {$M$};       % Điểm M
    
    % Đánh dấu độ dài
    \node[below] at (1, 0) {$18$ cm};   % Độ dài LN
    \node[below] at (4, 0) {$36$ cm} ;% Độ dài NM
        \begin{scope}
        % Giới hạn vùng tô chỉ thuộc nửa đường tròn lớn mà không thuộc nửa đường tròn nhỏ
        \clip (2,0) arc[start angle=0, end angle=-180, radius=1cm] -- (N) -- (L) -- cycle;
        \fill[pattern=north east lines] (2,0) arc[start angle=0, end angle=-180, radius=1cm];
        
    \end{scope}
    \draw[fill=black] (4,0) circle (1pt);
\end{tikzpicture}
}
\end{minipage}

\end{tcolorbox}
\vspace{4pt}


\begin{tcolorbox}[
    breakable,
    colback=white,
    colframe=blue!60!black,
    boxrule=0.8pt,
    arc=3pt,
    title={\textbf{Câu \#362} \quad \textbf{MÃ CÂU: \texttt{[L09\_C5\_B3\_28]}} \quad | \quad \textbf{Nhóm: \texttt{HINH\_269}} \quad | \quad \textbf{MapID: \texttt{9H5V3-1}} \hfill \textcolor{yellow!90!white}{[Hình gốc .pyc]}}
]
\textbf{Dạng toán:} bài tập tổng hợp \hfill \textit{(Tự luận -- Vận dụng)}
\vspace{4pt}\hrule\vspace{6pt}
\begin{minipage}[t]{0.58\textwidth}
\vspace{0pt}
Cho đường tròn ${(O)}$ đường kính ${AB}$. Lấy ${M}$ thuộc ${AB}$. Vẽ dây ${CD}$ vuông góc với ${AB}$ tại ${M}$. Giả sử ${AM =3cm}$, ${CD=6}cm$. Tính 

a) Độ dài đường tròn ${(O)}$. 

b) Độ dài cung ${CAD}$ 

(Các kết quả đều làm tròn đến hàng phần mười, lấy $\pi \approx 3,14$)
\end{minipage}\hfill
\begin{minipage}[t]{0.39\textwidth}
\vspace{0pt}
\centering

\resizebox{\linewidth}{!}{%
\begin{tikzpicture}
    % Vẽ đường tròn và các điểm
    \draw[thick] (0,0) circle(3cm); % Đường tròn tâm O bán kính 3cm
    \coordinate [label=below left: $A$] (A) at (-3,0); % Điểm A
    \coordinate [label=below right: $B$] (B) at (3,0);  % Điểm B
    \coordinate [label=below: $D$] (D) at (-1.5,-2.6);  % Trung điểm M của AB
    \coordinate [label=below right: $M$] (M) at (-1.5,0);  % Điểm D trên AB
    \coordinate [label=above: $C$] (C) at (-1.5,2.6);  % Điểm C
    
    % Tâm O của đường tròn
    \coordinate [label=below right: $O$] (O) at (0,0);
    
    % Vẽ các đoạn thẳng
    \draw[thick] (A) -- (B);  % Đoạn AB
    \draw[thick] (A) -- (C);  % Đoạn AC
    \draw[thick] (B) -- (C);  % Đoạn BC

    \draw[thick] (M) -- (C);  % Đoạn MC
    \draw[thick] (M) -- (D);  % Đoạn MD
    
    % Đánh dấu các điểm
    \filldraw[black] (A) circle(1.5pt);
    \filldraw[black] (B) circle(1.5pt);
    \filldraw[black] (M) circle(1.5pt);
    \filldraw[black] (O) circle(1.5pt);
    \filldraw[black] (C) circle(1.5pt);
    \filldraw[black] (D) circle(1.5pt);
\end{tikzpicture}
}
\end{minipage}

\end{tcolorbox}
\vspace{4pt}


\begin{tcolorbox}[
    breakable,
    colback=white,
    colframe=blue!60!black,
    boxrule=0.8pt,
    arc=3pt,
    title={\textbf{Câu \#363} \quad \textbf{MÃ CÂU: \texttt{[L09\_C5\_B3\_44]}} \quad | \quad \textbf{Nhóm: \texttt{HINH\_264}} \quad | \quad \textbf{MapID: \texttt{9H5H3-7}} \hfill \textcolor{yellow!90!white}{[Hình gốc .pyc]}}
]
\textbf{Dạng toán:} SA-M2. Bài toán thực tế. Tính diện tích cái quạt \hfill \textit{(Trả lời ngắn -- Thông hiểu)}
\vspace{4pt}\hrule\vspace{6pt}
\begin{minipage}[t]{0.58\textwidth}
\vspace{0pt}
Hình quạt tô màu đỏ ở hình vẽ bên dưới có bán kính bằng $15$ dm và góc ở tâm bằng $150^{\circ}$. Tính diện tích của hình quạt đó. Lấy $\pi=3,14$ làm tròn kết quả đến hàng đơn vị.


\textbf{Đáp án điền số:} 	exttt{14130}
\end{minipage}\hfill
\begin{minipage}[t]{0.39\textwidth}
\vspace{0pt}
\centering

\resizebox{\linewidth}{!}{%
\begin{tikzpicture}
    
    % Vẽ hình quạt đỏ với góc ở tâm là 150 độ và trục tung làm trục đối xứng
    \fill[red] (0,0) -- (4,0) arc[start angle=0, end angle=150, radius=4cm] -- cycle;
    
    % Vẽ tay cầm ở dưới và đối xứng qua trục tung
    \fill[orange] (0,-0.2) -- (0.5,-0.8) arc[start angle=-60, end angle=-120, radius=0.5cm] -- cycle;
    
\end{tikzpicture}
}
\end{minipage}

\vspace{3pt}\hrule\vspace{3pt}\footnotesize\textcolor{gray!80!black}{\textbf{Các mã dùng chung hình (HINH\_264):} \texttt{L09\_C5\_B3\_14}}
\end{tcolorbox}
\vspace{4pt}


\begin{tcolorbox}[
    breakable,
    colback=white,
    colframe=blue!60!black,
    boxrule=0.8pt,
    arc=3pt,
    title={\textbf{Câu \#364} \quad \textbf{MÃ CÂU: \texttt{[L09\_C5\_B3\_45]}} \quad | \quad \textbf{Nhóm: \texttt{HINH\_265}} \quad | \quad \textbf{MapID: \texttt{9H5V3-7}} \hfill \textcolor{yellow!90!white}{[Hình gốc .pyc]}}
]
\textbf{Dạng toán:} SA-M3. Bài toán thực tế liên quan đến quãng đường đi của xe \hfill \textit{(Trả lời ngắn -- Vận dụng)}
\vspace{4pt}\hrule\vspace{6pt}
\begin{minipage}[t]{0.58\textwidth}
\vspace{0pt}
Bánh xe (khi bơm căng) của một chiếc xe đạp có đường kính $660$ mm. Biết rằng khi giò đĩa quay một vòng thì bánh xe đạp quay được khoảng $4$ vòng (hình vẽ). Hỏi chiếc xe đạp di chuyển được quãng đường dài bao nhiêu mét sau khi người đi xe đạp $14$ vòng liên tục? Làm tròn kết quả đến hàng đơn vị.


\textbf{Đáp án điền số:} 	exttt{116}
\end{minipage}\hfill
\begin{minipage}[t]{0.39\textwidth}
\vspace{0pt}
\centering

\resizebox{\linewidth}{!}{%
\begin{tikzpicture}[scale=0.8, line cap=round, line join=round, thick]
    
    % Bánh xe sau
    \draw[thick] (0,0) circle(1.5cm);             % Bánh xe lớn
    \draw[fill=black] (0,0) circle(0.1cm);         % Tâm bánh xe
    % Vẽ nan hoa cho bánh xe sau
    \draw (0,0) -- (0,1.5);
    \draw (0,0) -- (0,-1.5);
    \draw (0,0) -- (1.3,0.75);
    \draw (0,0) -- (-1.3,-0.75);
    \draw (0,0) -- (1.3,-0.75);
    \draw (0,0) -- (-1.3,0.75);
    \draw (0,0) -- (0.75,1.3);
    \draw (0,0) -- (-0.75,-1.3);
    \draw (0,0) -- (0.75,-1.3);
    \draw (0,0) -- (-0.75,1.3);
    
    % Bánh xe trước
    \draw[thick] (8,0) circle(1.5cm);             % Bánh xe lớn
    \draw[fill=black] (8,0) circle(0.1cm);         % Tâm bánh xe
    % Vẽ nan hoa cho bánh xe trước
    \draw (8,0) -- (8,1.5);
    \draw (8,0) -- (8,-1.5);
    \draw (8,0) -- (9.3,0.75);
    \draw (8,0) -- (6.7,-0.75);
    \draw (8,0) -- (9.3,-0.75);
    \draw (8,0) -- (6.7,0.75);
    \draw (8,0) -- (8.75,1.3);
    \draw (8,0) -- (7.25,-1.3);
    \draw (8,0) -- (8.75,-1.3);
    \draw (8,0) -- (7.25,1.3);
    
    % Khung xe
    \draw[thick,blue] (0,0) -- (3,2.5);            % Khung từ bánh sau đến khung trên
    \draw[thick,blue] (3,2.5) -- (7,2.5);          % Khung trên từ yên đến bánh trước
    \draw[thick,blue] (7,2.5) -- (8,0);            % Khung từ khung trên xuống bánh trước
    \draw[thick,blue] (0,0) -- (8,0);              % Thanh dưới khung xe
    \draw[thick,blue] (3,2.5) -- (5.5,0);          % Thanh nối từ yên xuống giò đĩa
    
    % Ghi đông và tay lái
    \draw[thick,blue] (7,2.5) -- (7,3.5);          % Thanh đứng nối ghi đông
    \draw[thick,blue] (6.5,3.5) -- (7.5,3.5);      % Tay lái ngang
    \draw[thick,blue] (6.5,3.5) -- (6,3);          % Tay cầm trái
    \draw[thick,blue] (7.5,3.5) -- (8,3);          % Tay cầm phải
    
    % Yên xe
    \draw[thick,blue] (3,2.5) -- (3,2.8);          % Thanh nối yên xe
    \draw[thick,red] (2.7,2.8) -- (3.3,2.8);       % Yên xe
    
    % Giò đĩa và bàn đạp
    \draw[thick,blue] (5.5,0) circle(0.3cm);       % Đĩa xích
    \draw (5.5,0) -- ++(0:0.3cm);                  % Nối giò đĩa ngang
    \draw (5.5,0) -- ++(90:0.3cm);                 % Nối giò đĩa dọc trên
    \draw (5.5,0) -- ++(180:0.3cm);                % Nối giò đĩa ngang trái
    \draw (5.5,0) -- ++(270:0.3cm);                % Nối giò đĩa dọc dưới
    \draw[thick,blue] (5.5,0) -- ++(45:1cm);       % Bàn đạp trái
    \draw[thick,blue] (5.5,0) -- ++(225:1cm);      % Bàn đạp phải
    
\end{tikzpicture}
}
\end{minipage}

\vspace{3pt}\hrule\vspace{3pt}\footnotesize\textcolor{gray!80!black}{\textbf{Các mã dùng chung hình (HINH\_265):} \texttt{L09\_C5\_B3\_15}}
\end{tcolorbox}
\vspace{4pt}


\begin{tcolorbox}[
    breakable,
    colback=white,
    colframe=blue!60!black,
    boxrule=0.8pt,
    arc=3pt,
    title={\textbf{Câu \#365} \quad \textbf{MÃ CÂU: \texttt{[L09\_C5\_B3\_48]}} \quad | \quad \textbf{Nhóm: \texttt{HINH\_266}} \quad | \quad \textbf{MapID: \texttt{9H5C3-1}} \hfill \textcolor{yellow!90!white}{[Hình gốc .pyc]}}
]
\textbf{Dạng toán:} SA-M4. Tính bkính đtròn ngtiếp tam giác khi biết 2 cạnh và góc xen giữa của tam giác. \hfill \textit{(Trả lời ngắn -- Vận dụng cao)}
\vspace{4pt}\hrule\vspace{6pt}
\begin{minipage}[t]{0.58\textwidth}
\vspace{0pt}
Cho tam giác ${ABC}$ có $AB=AC=95cm$, $\widehat{A} =142^{\circ}$. Khi đó bán kính đường tròn đi qua ba đỉnh của tam giác bằng bao nhiêu? (làm tròn kết quả đến hàng đơn vị).


\textbf{Đáp án điền số:} 	exttt{146}
\end{minipage}\hfill
\begin{minipage}[t]{0.39\textwidth}
\vspace{0pt}
\centering

\resizebox{\linewidth}{!}{%
\begin{tikzpicture}
    
    % Tọa độ của các điểm
    \def\r{3}  % Bán kính của đường tròn
    \coordinate (O) at (0,0);         % Tâm O
    \coordinate (A) at (0,\r);        % Điểm A trên đường tròn
    \coordinate (B) at ({\r*sqrt(2)/2},{\r*sqrt(2)/2}); % Điểm B trên đường tròn
    \coordinate (C) at ({-\r*sqrt(2)/2},{\r*sqrt(2)/2}); % Điểm C trên đường tròn
    \coordinate (I) at (0,\r/2);      % Trung điểm của OA là điểm I
    
    % Vẽ đường tròn
    \draw (O) circle (\r);
    
    % Vẽ các đoạn thẳng
    \draw (O) -- (A);  % OA
    \draw (C) -- (A);  % CA
    \draw (B) -- (A);  % BA
    \draw (C) -- (B);  % CB
    \draw (C) -- (O);  % CO
    \draw (B) -- (O);  % BO
    
    % Đánh dấu các điểm
    \node[above] at (A) {A};
    \node[below] at (O) {O};
    \node[left] at (C) {C};
    \node[right] at (B) {B};
    \node[right] at (I) {I};
    
    % Vẽ các dấu hiệu góc bằng nhau
     \draw 
    (A) pic[draw,angle radius = 10] {{angle = O--A--B}} ;
 \draw 
(A) pic[draw,angle radius = 10] {{angle = C--A--O}} ;
    
\end{tikzpicture}
}
\end{minipage}

\vspace{3pt}\hrule\vspace{3pt}\footnotesize\textcolor{gray!80!black}{\textbf{Các mã dùng chung hình (HINH\_266):} \texttt{L09\_C5\_B3\_18}}
\end{tcolorbox}
\vspace{4pt}


\begin{tcolorbox}[
    breakable,
    colback=white,
    colframe=blue!60!black,
    boxrule=0.8pt,
    arc=3pt,
    title={\textbf{Câu \#366} \quad \textbf{MÃ CÂU: \texttt{[L09\_C5\_B3\_53]}} \quad | \quad \textbf{Nhóm: \texttt{HINH\_267}} \quad | \quad \textbf{MapID: \texttt{9H5V3-7}} \hfill \textcolor{yellow!90!white}{[Hình gốc .pyc]}}
]
\textbf{Dạng toán:} SA-M3. Bài toán thực tế tính diện tích sân vận động \hfill \textit{(Trả lời ngắn -- Vận dụng)}
\vspace{4pt}\hrule\vspace{6pt}
\begin{minipage}[t]{0.58\textwidth}
\vspace{0pt}
Một sân vận động có đường chạy đua như hình vẽ. Đường chạy gồm các đoạn ${AB}$, ${CD}$ là chiều dài của hình chữ nhật và hai nửa đường tròn. Biết ${AB}=55 (m)$, $CB= 44 (m)$. Diện tích sân vận động bằng a $(m^{2})$. Tìm a/100 (Làm tròn kết quả đến hàng đươn vị).


\textbf{Đáp án điền số:} 	exttt{39}
\end{minipage}\hfill
\begin{minipage}[t]{0.39\textwidth}
\vspace{0pt}
\centering

\resizebox{\linewidth}{!}{%
\begin{tikzpicture}
    % Các tọa độ
    \coordinate (A) at (0,2);  % Điểm A
    \coordinate (B) at (6,2);  % Điểm B
    \coordinate (D) at (0,0);  % Điểm D
    \coordinate (C) at (6,0);  % Điểm C
    \coordinate (center1) at (0,1); % Tâm đường tròn bên trái
    \coordinate (center2) at (6,1); % Tâm đường tròn bên phải
    
    % Vẽ hình chữ nhật
    \draw (A) -- (B) -- (C) -- (D) -- cycle;
    
    % Vẽ nửa hình tròn bên trái
    \draw (A) arc[start angle=90, end angle=270, radius=1cm];
    
    % Vẽ nửa hình tròn bên phải
    \draw (B) arc[start angle=90, end angle=-90, radius=1cm];
    
    % Đánh dấu các điểm
    \filldraw (A) circle (1pt) node[above] {$A$};
    \filldraw (B) circle (1pt) node[above] {$B$};
    \filldraw (C) circle (1pt) node[below] {$C$};
    \filldraw (D) circle (1pt) node[below] {$D$};
\end{tikzpicture}
}
\end{minipage}

\vspace{3pt}\hrule\vspace{3pt}\footnotesize\textcolor{gray!80!black}{\textbf{Các mã dùng chung hình (HINH\_267):} \texttt{L09\_C5\_B3\_23}}
\end{tcolorbox}
\vspace{4pt}


\begin{tcolorbox}[
    breakable,
    colback=white,
    colframe=blue!60!black,
    boxrule=0.8pt,
    arc=3pt,
    title={\textbf{Câu \#367} \quad \textbf{MÃ CÂU: \texttt{[L09\_C5\_B3\_54]}} \quad | \quad \textbf{Nhóm: \texttt{HINH\_270}} \quad | \quad \textbf{MapID: \texttt{9H5V3-1}} \hfill \textcolor{yellow!90!white}{[Hình gốc .pyc]}}
]
\textbf{Dạng toán:} SA-M3. Tính diện tích các hình phức tạp \hfill \textit{(Trả lời ngắn -- Vận dụng)}
\vspace{4pt}\hrule\vspace{6pt}
\begin{minipage}[t]{0.58\textwidth}
\vspace{0pt}
Diện tích phần kẻ sọc ở hình sau là $a \pi$ $(cm^{2})$, giới hạn bởi nửa đường tròn đường kính ${LM}$ và hai nửa đường tròn có đường kính tương ứng là $NM={20}$ cm và $LN= {10}$ cm. Tìm a.


\textbf{Đáp án điền số:} 	exttt{75}
\end{minipage}\hfill
\begin{minipage}[t]{0.39\textwidth}
\vspace{0pt}
\centering

\resizebox{\linewidth}{!}{%
\begin{tikzpicture}
    
    % Khởi tạo các điểm
    \coordinate (L) at (0, 0);           % Điểm L
    \coordinate (N) at (2, 0);           % Điểm N
    \coordinate (M) at (6, 0);          % Điểm M
    \coordinate (centerLN) at (1, 0);    % Trung điểm đoạn LN (tâm của nửa đường tròn mới phía dưới)
    
    % Vẽ nửa đường tròn lớn với đường kính 12 cm (phía trên)
    \draw (L) arc[start angle=180, end angle=0, radius=3cm];
    
    % Vẽ nửa đường tròn nhỏ với đường kính 8 cm (phía trên)
    \draw (N) arc[start angle=180, end angle=0, radius=2cm];
    
    % Vẽ đoạn thẳng đáy LM
    \draw (L) -- (M);
        \draw (2,0) arc[start angle=0, end angle=-180, radius=1cm];
    % Tô phần gạch sọc chỉ thuộc nửa đường tròn lớn mà không thuộc nửa đường tròn nhỏ
    \begin{scope}
        % Giới hạn vùng tô chỉ thuộc nửa đường tròn lớn mà không thuộc nửa đường tròn nhỏ
        \clip (L) arc[start angle=180, end angle=0, radius=3cm] -- (M) -- (L) -- cycle;
        \fill[pattern=north east lines] (L) arc[start angle=180, end angle=0, radius=3cm];
        
        % Loại bỏ vùng thuộc nửa đường tròn nhỏ ra khỏi vùng tô bằng một vùng trắng
        \fill[white] (N) arc[start angle=180, end angle=0, radius=2cm] -- (N) -- cycle;
    \end{scope}
    
    % Vẽ nửa đường tròn đối diện phía dưới đoạn LN
    % Tâm của nửa đường tròn là trung điểm của đoạn LN, với bán kính 2 cm

    
    % Đánh dấu các điểm
    \filldraw (L) circle (1pt) node[below left] {$L$};        % Điểm L
    \filldraw (N) circle (1pt) node[below] {$N$};             % Điểm N
    \filldraw (M) circle (1pt) node[below right] {$M$};       % Điểm M
    
    % Đánh dấu độ dài
    \node[below] at (1, 0) {$10$ cm};   % Độ dài LN
    \node[below] at (4, 0) {$20$ cm} ;% Độ dài NM
        \begin{scope}
        % Giới hạn vùng tô chỉ thuộc nửa đường tròn lớn mà không thuộc nửa đường tròn nhỏ
        \clip (2,0) arc[start angle=0, end angle=-180, radius=1cm] -- (N) -- (L) -- cycle;
        \fill[pattern=north east lines] (2,0) arc[start angle=0, end angle=-180, radius=1cm];
        
    \end{scope}
    \draw[fill=black] (4,0) circle (1pt);
\end{tikzpicture}
}
\end{minipage}

\end{tcolorbox}
\vspace{4pt}


\begin{tcolorbox}[
    breakable,
    colback=white,
    colframe=blue!60!black,
    boxrule=0.8pt,
    arc=3pt,
    title={\textbf{Câu \#368} \quad \textbf{MÃ CÂU: \texttt{[L09\_C5\_B6\_02]}} \quad | \quad \textbf{Nhóm: \texttt{HINH\_271}} \quad | \quad \textbf{MapID: \texttt{9H5H4-5}} \hfill \textcolor{yellow!90!white}{[Hình gốc .pyc]}}
]
\textbf{Dạng toán:} Thông hiểu tiếp tuyến của đường tròn \hfill \textit{(Trắc nghiệm -- Thông hiểu)}
\vspace{4pt}\hrule\vspace{6pt}
\begin{minipage}[t]{0.58\textwidth}
\vspace{0pt}
Cho hình vẽ sau. ${FQ}$, ${FK}$ lần lượt là tiếp tuyến của đường tròn ${(O)}$ tại ${Q}$ và tại ${K}$. Khẳng định nào sau đây sai?

\vspace{4pt}\noindent
\textbf{A.} ${OF}$ là đường trung trực của đoạn thẳng ${QK}$ \quad \textbf{B.} $\Delta {QFO}= \Delta {KFO} $ \quad \textbf{C.} ${OF}$ là tia phân giác của $\widehat{QOK}$ \quad \textbf{\textcolor{blue!75!black}{D. (Đ)}} $\Delta{QKF}$ là tam giác đều
\end{minipage}\hfill
\begin{minipage}[t]{0.39\textwidth}
\vspace{0pt}
\centering

\resizebox{\linewidth}{!}{%
\begin{tikzpicture}
    \coordinate (O) at (0, 0);
    \coordinate (M) at (5, 0);
    \coordinate (B) at (0.8, 1.83);
    \coordinate (A) at (0.8, -1.83);
    
    % Vẽ tứ giác
    \draw[thick] (O) -- (B) -- (M) -- (O);
    \draw[thick] (O) -- (A) -- (M) -- (O);
    
    % Vẽ đường tròn tâm O, bán kính OA
    \draw[thick] (O) circle[radius=1.97];  % Bán kính OA được tính là khoảng cách từ O đến A
    
    % Ghi nhãn các điểm
    \node[below left] at (O) {$O$};
    \node[above right] at (M) {$F$};
    \node[above left] at (B) {$K$};
    \node[below left] at (A) {$Q$};
    
    % Vẽ ký hiệu góc vuông
    \def\khvuong[size=#1](#2,#3,#4){%
            \draw ($(#3)!#1!(#2)$) --
            ($($(#3)!#1!(#2)$)+($(#3)!#1!(#4)$)-(#3)$)--($(#3)!#1!(#4)$);
    }
    \khvuong[size=6pt](O,A,M);
    \khvuong[size=6pt](O,B,M);
    
    % Đánh dấu các điểm
    \fill [black] (O) circle (1.5pt);
    \fill [black] (M) circle (1.5pt);
    \fill [black] (B) circle (1.5pt);
    \fill [black] (A) circle (1.5pt);
 \end{tikzpicture}
}
\end{minipage}

\end{tcolorbox}
\vspace{4pt}


\begin{tcolorbox}[
    breakable,
    colback=white,
    colframe=blue!60!black,
    boxrule=0.8pt,
    arc=3pt,
    title={\textbf{Câu \#369} \quad \textbf{MÃ CÂU: \texttt{[L09\_C5\_B6\_03]}} \quad | \quad \textbf{Nhóm: \texttt{HINH\_272}} \quad | \quad \textbf{MapID: \texttt{9H5H4-5}} \hfill \textcolor{yellow!90!white}{[Hình gốc .pyc]}}
]
\textbf{Dạng toán:} Thông hiểu tiếp tuyến của đường tròn kiểu 2 \hfill \textit{(Trắc nghiệm -- Thông hiểu)}
\vspace{4pt}\hrule\vspace{6pt}
\begin{minipage}[t]{0.58\textwidth}
\vspace{0pt}
Cho hình vẽ sau. Khẳng định nào sau đây đúng?

\vspace{4pt}\noindent
\textbf{A.} $\Delta{FPO} $ cân tại ${F}$ \quad \textbf{B.} $\Delta{FPK} =\Delta { FOK}$ \quad \textbf{C.} ${OK} ={PO}$ \quad \textbf{\textcolor{blue!75!black}{D. (Đ)}} $\Delta{FKO} =\Delta { FQO}$
\end{minipage}\hfill
\begin{minipage}[t]{0.39\textwidth}
\vspace{0pt}
\centering

\resizebox{\linewidth}{!}{%
\begin{tikzpicture}
    \coordinate (C) at (0, 0);
    \coordinate (G) at (5, 0);
    \coordinate (E) at (0.8, 1.83);
    \coordinate (F) at (0.8, -1.83);
     \coordinate (T) at (-1.18, 2.7);
    % Vẽ tứ giác
    \draw[thick] (G) -- (E) -- (T) -- (C)--(F)--(E)--(C);
    \draw[thick] (E) -- (F) ;
      \draw[thick] (G) -- (F) ;
       \draw[thick] (G) -- (C) ;
    % Vẽ đường tròn tâm O, bán kính OA
    \draw[thick] (C) circle[radius=1.97];  % Bán kính OA được tính là khoảng cách từ O đến A
    
    % Ghi nhãn các điểm
    \node[below right] at (G) {$F$};
    \node[above right] at (E) {$Q$};
    \node[below left] at (F) {$K$};
    \node[below left] at (C) {$O$};
    \node[below left] at (T) {$P$};
    % Vẽ ký hiệu góc vuông
    \def\khvuong[size=#1](#2,#3,#4){%
            \draw ($(#3)!#1!(#2)$) --
            ($($(#3)!#1!(#2)$)+($(#3)!#1!(#4)$)-(#3)$)--($(#3)!#1!(#4)$);
    }
    \khvuong[size=6pt](C,E,G);
    \khvuong[size=6pt](C,F,G);
    
    % Đánh dấu các điểm
    \fill [black] (G) circle (1.5pt);
    \fill [black] (E) circle (1.5pt);
    \fill [black] (F) circle (1.5pt);
    \fill [black] (C) circle (1.5pt);
     \fill [black] (T) circle (1.5pt);
 \end{tikzpicture}
}
\end{minipage}

\end{tcolorbox}
\vspace{4pt}


\begin{tcolorbox}[
    breakable,
    colback=white,
    colframe=blue!60!black,
    boxrule=0.8pt,
    arc=3pt,
    title={\textbf{Câu \#370} \quad \textbf{MÃ CÂU: \texttt{[L09\_C5\_B6\_04]}} \quad | \quad \textbf{Nhóm: \texttt{HINH\_273}} \quad | \quad \textbf{MapID: \texttt{9H5N4-1}} \hfill \textcolor{yellow!90!white}{[Hình gốc .pyc]}}
]
\textbf{Dạng toán:} Nhận biết tiếp tuyển của đường tròn bài toán liên quan đến tam giác vuông \hfill \textit{(Trắc nghiệm -- Nhận biết)}
\vspace{4pt}\hrule\vspace{6pt}
\begin{minipage}[t]{0.58\textwidth}
\vspace{0pt}
Cho $\Delta{FME}$ vuông tại ${E}$. Khi đó ${FE}$ là tiếp tuyến của đường tròn nào sau đây?

\vspace{4pt}\noindent
\textbf{A.} Đường tròn đường kính ${FM}$ \quad \textbf{B.} Đường tròn $(E,ME)$ \quad \textbf{C.} Đường tròn $(F,ME)$ \quad \textbf{\textcolor{blue!75!black}{D. (Đ)}} Đường tròn $(M,ME)$
\end{minipage}\hfill
\begin{minipage}[t]{0.39\textwidth}
\vspace{0pt}
\centering

\resizebox{\linewidth}{!}{%
\begin{tikzpicture}
        % Tọa độ cho tam giác vuông
        \coordinate (P) at (0, 0);
        \coordinate (M) at (4, 0);
        \coordinate (N) at (0, 3);

        % Vẽ các cạnh của tam giác
        \draw[thick] (M) -- (N) -- (P) -- cycle;

        % Ghi nhãn các điểm
        \node[below left] at (P) {$E$};
        \node[below right] at (M) {$F$};
        \node[above left] at (N) {$M$};

        % Vẽ dấu vuông tại điểm A
        \draw (0,0) rectangle +(0.3,0.3);
    \end{tikzpicture}
}
\end{minipage}

\end{tcolorbox}
\vspace{4pt}


\begin{tcolorbox}[
    breakable,
    colback=white,
    colframe=blue!60!black,
    boxrule=0.8pt,
    arc=3pt,
    title={\textbf{Câu \#371} \quad \textbf{MÃ CÂU: \texttt{[L09\_C5\_B6\_05]}} \quad | \quad \textbf{Nhóm: \texttt{HINH\_274}} \quad | \quad \textbf{MapID: \texttt{9H5N4-5}} \hfill \textcolor{yellow!90!white}{[Hình gốc .pyc]}}
]
\textbf{Dạng toán:} Tam giác ABC vuông tại A đường cao AH nhận biết tiếp tuyến trong bài \hfill \textit{(Trắc nghiệm -- Nhận biết)}
\vspace{4pt}\hrule\vspace{6pt}
\begin{minipage}[t]{0.58\textwidth}
\vspace{0pt}
Cho $\Delta{QKO}$ vuông tại ${Q}$, đường cao ${QH}$. Chọn khẳng định đúng?

\vspace{4pt}\noindent
\textbf{A.} ${QK}$ là tiếp tuyến của đường tròn $(Q,QO)$ \quad \textbf{B.} ${KO}$ là tiếp tuyến của đường tròn tâm ${H}$ \quad \textbf{C.} ${QK}$ là tiếp tuyến của đường tròn đường kính ${KO}$ \quad \textbf{\textcolor{blue!75!black}{D. (Đ)}} ${KO}$ là tiếp tuyến của đường tròn $(Q,QH)$
\end{minipage}\hfill
\begin{minipage}[t]{0.39\textwidth}
\vspace{0pt}
\centering

\resizebox{\linewidth}{!}{%
\begin{tikzpicture}[scale=0.6]  % Thu nhỏ toàn bộ hình vẽ
            % Tạo các điểm A, B, C, H
            \coordinate [label=above:{$Q$}] (A) at (0, 4);
            \coordinate [label=below left:{$K$}] (B) at (-3, 0);
            \coordinate [label=below right:{$O$}] (C) at (4, 0);
            \coordinate [label=below:{$H$}] (H) at (0, 0);
            
            % Vẽ các cạnh tam giác
            \draw (B) -- (A) -- (C) -- cycle;
            
            % Vẽ đường cao từ A xuống H
            \draw  (A) -- (H);
            
            % Vẽ góc vuông tại H
            \draw (H) rectangle ++(0.3, 0.3);   

            %Vẽ góc vuông tại A
            \pic[draw,angle radius=10]{right angle=B--A--C};
        \end{tikzpicture}
}
\end{minipage}

\end{tcolorbox}
\vspace{4pt}


\begin{tcolorbox}[
    breakable,
    colback=white,
    colframe=blue!60!black,
    boxrule=0.8pt,
    arc=3pt,
    title={\textbf{Câu \#372} \quad \textbf{MÃ CÂU: \texttt{[L09\_C5\_B7\_01]}} \quad | \quad \textbf{Nhóm: \texttt{HINH\_275}} \quad | \quad \textbf{MapID: \texttt{9H5H4-1}} \hfill \textcolor{yellow!90!white}{[Hình gốc .pyc]}}
]
\textbf{Dạng toán:} Tổng hợp 1 \hfill \textit{(Tự luận -- Thông hiểu)}
\vspace{4pt}\hrule\vspace{6pt}
\begin{minipage}[t]{0.58\textwidth}
\vspace{0pt}
Cho tam giác ${PQR}$ có $PQ<PR$ và đường cao ${PM}$.

a) Trong các điểm ${Q}$, ${M}$ và ${R}$, điểm nào nằm trong, điểm nào nằm trên và điểm nào nằm ngoài đường tròn $(P;PQ)$ ? Vì sao?

b) Xác định vị trí của điểm ${T}$ trên đoạn $PR$ trong mỗi trường hợp sau:

- Đường tròn $({P})$ và đường tròn $(R;RT)$ tiếp xúc với nhau;

 - Đường tròn $({P})$ và đường tròn $(R;RT)$ cắt nhau;

- Đường tròn $({P})$ và đường tròn $(R;RT)$ không giao nhau.\\
\end{minipage}\hfill
\begin{minipage}[t]{0.39\textwidth}
\vspace{0pt}
\centering

\resizebox{\linewidth}{!}{%
\begin{tikzpicture}[scale = .7]
     \path 
    (0,0) coordinate (B)
    (4,0) coordinate (C)
    (1.5,3) coordinate (A)
    ($(B)!(A)!(C)$) coordinate (H)
    ;
    \draw (A)--(B)--(C)--cycle (A)--(H);
    \pic[draw, angle radius=2mm]{right angle=C--H--A};
        \filldraw[black] (A) circle (2pt) node[above ] {$P$};
        \filldraw[black] (B) circle (2pt) node[above left] {$Q$};
        \filldraw[black] (C) circle (2pt) node[above right] {$R$};
        \filldraw[black] (H) circle (2pt) node[above right]{$M$};
    \end{tikzpicture}
}
\end{minipage}

\end{tcolorbox}
\vspace{4pt}


\begin{tcolorbox}[
    breakable,
    colback=white,
    colframe=blue!60!black,
    boxrule=0.8pt,
    arc=3pt,
    title={\textbf{Câu \#373} \quad \textbf{MÃ CÂU: \texttt{[L09\_C5\_B7\_02]}} \quad | \quad \textbf{Nhóm: \texttt{HINH\_276}} \quad | \quad \textbf{MapID: \texttt{9H5V4-1}} \hfill \textcolor{yellow!90!white}{[Hình gốc .pyc]}}
]
\textbf{Dạng toán:} Tổng hợp 2 \hfill \textit{(Tự luận -- Vận dụng)}
\vspace{4pt}\hrule\vspace{6pt}
\begin{minipage}[t]{0.58\textwidth}
\vspace{0pt}
Cho hình thang cân ${A B C D }$ $(A B  \parallel C D )$.

a) Chứng minh rằng đường trung trực ${y}$ của ${A B }$ cũng là đường trung trực của ${C D }$ (từ đó suy ra hai điểm ${A }$ và ${B }$ đối xứng với nhau, ${C }$ và ${D }$ đối xứng với nhau qua ${y}$).

b) Giải thích tại sao nếu một đường tròn đi qua 3 điểm ${A }$, ${B }$ và ${C }$ thì nó cũng đi qua điểm ${D }$.\\
\end{minipage}\hfill
\begin{minipage}[t]{0.39\textwidth}
\vspace{0pt}
\centering

\resizebox{\linewidth}{!}{%
\begin{tikzpicture}[scale=.7]
    \def\r{3}
    \path 
    (0:0) coordinate (O)
    (110:\r) coordinate (A)
    (70:\r) coordinate (B)
    (-40:\r) coordinate (C)
    (-140:\r) coordinate (D)
    (intersection of D--A and C--B) coordinate (S)
    ;
    \draw (A)--(B)--(C)--(D)--cycle (A)--(S)--(B) (S)--(O) node[midway, left]{$y$} (A)--(O)--(B) (C)--(O)--(D);
        \filldraw[black] (A) circle (2pt) node[above left] {$A $};
    \filldraw[black] (B) circle (2pt) node[above right] {$B $};
    \filldraw[black] (C) circle (2pt) node[above right] {$C $};
   \filldraw[black] (D) circle (2pt) node[above left]{$D $};
    \filldraw[black] (S) circle (2pt) node[above] {$M $};
     \filldraw[black] (O) circle (2pt) node[below] {$O $};
\end{tikzpicture}
}
\end{minipage}

\end{tcolorbox}
\vspace{4pt}


\begin{tcolorbox}[
    breakable,
    colback=white,
    colframe=blue!60!black,
    boxrule=0.8pt,
    arc=3pt,
    title={\textbf{Câu \#374} \quad \textbf{MÃ CÂU: \texttt{[L09\_C5\_B7\_03]}} \quad | \quad \textbf{Nhóm: \texttt{HINH\_277}} \quad | \quad \textbf{MapID: \texttt{9H5V4-1}} \hfill \textcolor{yellow!90!white}{[Hình gốc .pyc]}}
]
\textbf{Dạng toán:} Tổng hợp 3 \hfill \textit{(Tự luận -- Vận dụng)}
\vspace{4pt}\hrule\vspace{6pt}
\begin{minipage}[t]{0.58\textwidth}
\vspace{0pt}
Giả sử ${HF}$ là một dây song song với đường kính ${QK}$ của đường tròn (${O}$) sao cho $QKHF$ là một tứ giác lồi. Gọi ${E}$ là trung điểm của đoạn ${HF}$.

a) Chứng minh rằng $A$ đối xứng với ${K}$ và ${H}$ đối xứng với ${F}$ qua đường thẳng ${OE}$.

b) Chứng minh rằng tứ giác ${QKHF}$ là một hình thang cân.

c) Biết rằng $QK=22$ cm và $\widehat{HOF}=116^{\circ}$. Tính độ dài cung (nhỏ) ${QF}$ và cung (lớn) ${QKH}$.

d) Với giả thiết ở câu c, tính diện tích hình quạt tròn ứng với cung nhỏ ${KF}$.\\
\end{minipage}\hfill
\begin{minipage}[t]{0.39\textwidth}
\vspace{0pt}
\centering

\resizebox{\linewidth}{!}{%
\begin{tikzpicture}[scale=.7]
                \def\r{3}
                \path 
                (0:0) coordinate (O)
                (180:\r) coordinate (A)
                (0:\r) coordinate (B)
                (50:\r) coordinate (C)
                (130:\r) coordinate (D)
                (90:\r) coordinate (F)
                (-90:\r) coordinate (G)
                ($(C)!.5!(D)$) coordinate (E)
                ;
                \draw (O) circle (\r);
                \draw (D)--(A)--(B)--(C) (D)--(E) node[midway]{$/$} (C)--(E) node[midway]{$/$} (A)--(C) (B)--(D) (F)--(G) (C)--(O)--(D);
               
                   \filldraw[black] (A) circle (2pt) node[above left] {$Q$};
    \filldraw[black] (B) circle (2pt) node[above right] {$K$};
    \filldraw[black] (C) circle (2pt) node[above right] {$H$};
   \filldraw[black] (D) circle (2pt) node[above left]{$F$};
    \filldraw[black] (E) circle (2pt) node[below left] {$E$};
    \filldraw[black] (O) circle (2pt) node[below right] {$O$};
            \end{tikzpicture}
}
\end{minipage}

\end{tcolorbox}
\vspace{4pt}


\begin{tcolorbox}[
    breakable,
    colback=white,
    colframe=blue!60!black,
    boxrule=0.8pt,
    arc=3pt,
    title={\textbf{Câu \#375} \quad \textbf{MÃ CÂU: \texttt{[L09\_C5\_B7\_04]}} \quad | \quad \textbf{Nhóm: \texttt{HINH\_278}} \quad | \quad \textbf{MapID: \texttt{9H5V4-1}} \hfill \textcolor{yellow!90!white}{[Hình gốc .pyc]}}
]
\textbf{Dạng toán:} Tổng hợp 4 \hfill \textit{(Tự luận -- Vận dụng)}
\vspace{4pt}\hrule\vspace{6pt}
\begin{minipage}[t]{0.58\textwidth}
\vspace{0pt}
Cho tam giác vuông $IJQ$ $\left(\widehat{I}=90^{\circ} \right)$ có $\widehat{Q}=25^{\circ}$ và $IJ=5$ cm. Đường phân giác của góc ${J}$ cắt ${IQ}$ tại ${K}$.

a) Chứng minh rằng đường tròn $(K;KI)$ tiếp xúc với cạnh ${JQ}$.

b) Tính độ dài cung nằm trong góc ${JKQ}$ của đường tròn $(K;KI)$ và diện tích hình quạt tròn tương ứng với cung ấy.

c) Tính diện tích hình vành khuyên tạo bởi hai đường tròn $(K;KI)$ và $(K;KQ)$. 

 (Làm tròn các kết quả đến chữ số thập phân thứ hai, lấy $\pi = 3,14$)\\
\end{minipage}\hfill
\begin{minipage}[t]{0.39\textwidth}
\vspace{0pt}
\centering

\resizebox{\linewidth}{!}{%
\begin{tikzpicture}[scale=.8]
                \def\chanphangiac(#1,#2,#3)(#4){
                \path 
                ($(#2)!1mm!(#1)$) coordinate (#2#1)
                ($(#2)!1mm!(#3)$) coordinate (#2#3)
                ($(#2#1)!.5!(#2#3)$) coordinate (#2t)
                (intersection of #2--#2t and #1--#3) coordinate (#4);
                }
                \path 
                (0,0) coordinate (A)
                (0,3) coordinate (B)
                (5,0) coordinate (F)
                ($(B)!1mm!60:(A)$) coordinate (Bt)
                (intersection of B--Bt and A--F) coordinate (C)
                ;
                \chanphangiac(A,B,C)(D)
                \draw (A)--(B)--(C)--cycle (B)--(D);
                \draw (D) let \p1=($(D)-(A)$) in circle ({veclen(\x1,\y1)});
                \draw (D) let \p2=($(D)-(C)$) in circle ({veclen(\x2,\y2)});
                \path ($(B)!(D)!(C)$) coordinate (E);
                 \draw 
    (C) pic[draw,angle radius = 15] {angle = B--C--A} node[shift={(-15:-25pt)}]{$ 25^\circ $}   
    ;
        \filldraw[black] (A) circle (2pt) node[below left] {$I$};
    \filldraw[black] (B) circle (2pt) node[above left] {$J$};
    \filldraw[black] (C) circle (2pt) node[below right] {$Q$};
   \filldraw[black] (D) circle (2pt) node[below]{$K$};
    \filldraw[black] (E) circle (2pt) node[above right] {$E$};  
      % Vẽ góc vuông tại A
        \draw (A) rectangle +(0.5, 0.5);       
 \end{tikzpicture}
}
\end{minipage}

\end{tcolorbox}
\vspace{4pt}


\begin{tcolorbox}[
    breakable,
    colback=white,
    colframe=blue!60!black,
    boxrule=0.8pt,
    arc=3pt,
    title={\textbf{Câu \#376} \quad \textbf{MÃ CÂU: \texttt{[L09\_C5\_B7\_05]}} \quad | \quad \textbf{Nhóm: \texttt{HINH\_279}} \quad | \quad \textbf{MapID: \texttt{9H5V4-1}} \hfill \textcolor{yellow!90!white}{[Hình gốc .pyc]}}
]
\textbf{Dạng toán:} Tổng hợp 5 \hfill \textit{(Tự luận -- Vận dụng)}
\vspace{4pt}\hrule\vspace{6pt}
\begin{minipage}[t]{0.58\textwidth}
\vspace{0pt}
Từ điểm ${I}$ nằm ngoài đường tròn $(O)$, kẻ hai tiếp tuyến ${IQ}$ và ${IK}$ đến đường tròn (${Q}$ và ${K}$ là hai tiếp điểm).

a) Chứng minh rằng $IO \perp QK$.

b) Gọi ${H}$ là điểm đối xứng với ${Q}$ qua ${O}$. Chứng minh rằng $KH \parallel IO$.

c) Tính độ dài các cạnh của tam giác ${IQK}$, biết $OQ=25$ cm và $OI=65$ cm.\\
\end{minipage}\hfill
\begin{minipage}[t]{0.39\textwidth}
\vspace{0pt}
\centering

\resizebox{\linewidth}{!}{%
\begin{tikzpicture}[scale=.8]
    \def\r{3}
    \path 
    (0,0) coordinate (O)
    (-7,0) coordinate (P)
    ($(O)!.5!(P)$) coordinate (I);
    \path[name path=c1] (O) circle (\r);
    \path[name path=c2] (I) let\p1=($(I)-(O)$) in circle ({veclen(\x1,\y1)});
    \path[name intersections={of = c1 and c2, by={B,A}}]
    ($(A)!2!(O)$) coordinate (C)
    (intersection of A--B and O--P) coordinate (H);
    \draw (O) circle (\r) (A)--(C)--(B)--(P)--cycle (A)--(B) (O)--(P);
    \pic[draw, angle radius=2mm]{right angle = O--H--B};
    \filldraw[black] (A) circle (2pt) node[below] {$Q$};
    \filldraw[black] (B) circle (2pt) node[above left] {$K$};
    \filldraw[black] (C) circle (2pt) node[above] {$H$};
    \filldraw[black] (O) circle (2pt) node[above right]{$O$};
    \filldraw[black] (H) circle (2pt) node[below left] {$E$};
    \filldraw[black] (P) circle (2pt) node[above] {$I$};
    \end{tikzpicture}
}
\end{minipage}

\end{tcolorbox}
\vspace{4pt}


\begin{tcolorbox}[
    breakable,
    colback=white,
    colframe=blue!60!black,
    boxrule=0.8pt,
    arc=3pt,
    title={\textbf{Câu \#377} \quad \textbf{MÃ CÂU: \texttt{[L09\_C5\_B7\_06]}} \quad | \quad \textbf{Nhóm: \texttt{HINH\_280}} \quad | \quad \textbf{MapID: \texttt{9H5V4-1}} \hfill \textcolor{yellow!90!white}{[Hình gốc .pyc]}}
]
\textbf{Dạng toán:} Tổng hợp 6 \hfill \textit{(Tự luận -- Vận dụng)}
\vspace{4pt}\hrule\vspace{6pt}
\begin{minipage}[t]{0.58\textwidth}
\vspace{0pt}
Cho tam giác ${EFK}$ vuông tại ${E}$, đường cao ${EN}$. Từ ${F}$ và từ ${K}$ kẻ hai đường thẳng tiếp xúc với đường tròn $(E;EN)$ lần lượt tại ${H}$ và ${M}$. Chứng minh rằng:

 a) Hai điểm ${H}$ và ${M}$ đối xứng với nhau qua ${E}$;

b) ${HM}$ tiếp xúc với đường tròn đường kính ${FK}$.\\
\end{minipage}\hfill
\begin{minipage}[t]{0.39\textwidth}
\vspace{0pt}
\centering

\resizebox{\linewidth}{!}{%
\begin{tikzpicture}[scale=.8]
    \def\r{3}
    \path 
    (0:0) coordinate (O)
    (105:\r) coordinate (A)
    (180:\r) coordinate (B)
    (0:\r) coordinate (C)
    ($(B)!(A)!(C)$) coordinate (H);
    \path[name path=c1] (A) let\p1=($(A)-(H)$) in circle ({veclen(\x1,\y1)});
    \path[name path=c2] (B) let\p2=($(B)-(H)$) in circle ({veclen(\x2,\y2)});
    \path[name intersections={of= c1 and c2,by={D}}]
    ($(D)!2!(A)$) coordinate (E);
    \pgfresetboundingbox
    \draw (A) let\p1=($(A)-(H)$) in circle ({veclen(\x1,\y1)});
    \draw (A)--(B)--(C)--cycle (A)--(H) (B)--(D)--(E)--(C) (A)--(O);
    \pic[draw,angle radius=2mm]{right angle = C--H--A};
    \pic[draw,angle radius=2mm]{right angle = B--A--C} ;

   

       \filldraw[black] (A) circle (2pt) node[above right] {$E$};
    \filldraw[black] (B) circle (2pt) node[below left] {$F$};
    \filldraw[black] (C) circle (2pt) node[above right] {$K$};
   \filldraw[black] (D) circle (2pt) node[above left]{$H$};
    \filldraw[black] (H) circle (2pt) node[below left] {$N$};
    \filldraw[black] (O) circle (2pt) node[below] {$O$};
     \filldraw[black] (E) circle (2pt) node[above] {$M$};
\end{tikzpicture}
}
\end{minipage}

\end{tcolorbox}
\vspace{4pt}


\begin{tcolorbox}[
    breakable,
    colback=white,
    colframe=blue!60!black,
    boxrule=0.8pt,
    arc=3pt,
    title={\textbf{Câu \#378} \quad \textbf{MÃ CÂU: \texttt{[L09\_C5\_B7\_07]}} \quad | \quad \textbf{Nhóm: \texttt{HINH\_281}} \quad | \quad \textbf{MapID: \texttt{9H5C4-1}} \hfill \textcolor{yellow!90!white}{[Hình gốc .pyc]}}
]
\textbf{Dạng toán:} Tổng hợp 7 \hfill \textit{(Tự luận -- Vận dụng cao)}
\vspace{4pt}\hrule\vspace{6pt}
\begin{minipage}[t]{0.58\textwidth}
\vspace{0pt}
Cho đường tròn $({O})$, đường thẳng ${d}$ tiếp xúc với $({O})$ tại ${M}$, đường thẳng ${m}$ tiếp xúc với $({O})$ tại ${N}$ sao cho $d \parallel m$. Gọi ${P}$ là một điểm tuỳ ý thuộc $({O})$, khác ${M}$ và ${N}$. Tiếp tuyến ${t}$ của $({O})$ tại ${P}$ cắt ${d}$ và ${m}$ lần lượt tại ${I}$ và ${J}$.

a) Chứng minh ${MN}$ là một đường kính của $({O})$.

b) Gọi ${Q}$, ${E}$ và ${F}$ lần lượt là các điểm đối xứng với ${P}$, ${I}$ và ${J}$ qua tâm ${O}$. Chứng minh rằng $Q \in (O)$, $E \in m$ và $F \in d$.

c) Chứng minh rằng ${EF}$ tiếp xúc với $({O})$ tại ${Q}$.

d) Chứng minh tứ giác ${IJEF}$ là một hình thoi.\\
\end{minipage}\hfill
\begin{minipage}[t]{0.39\textwidth}
\vspace{0pt}
\centering

\resizebox{\linewidth}{!}{%
\begin{tikzpicture}[scale=.8]
                \def\r{3}
                \path 
                (0:0) coordinate (O)
                (180:\r) coordinate (A)
                (0:\r) coordinate (B)
                (-70:\r) coordinate (C)
                ($(A)!5cm!90:(O)$) coordinate (At) node[below right]{$d$}
                ($(At)!2!(A)$) coordinate (Az)
                ($(B)!5cm!90:(O)$) coordinate (Bt) 
                ($(Bt)!2!(B)$) coordinate (Bz) node[below left]{$m$}
                ($(C)!1mm!90:(O)$) coordinate (Ct)
                (intersection of C--Ct and Bz--Bt) coordinate (N) (intersection of C--Ct and Az--At) coordinate (M)
                ($(C)!2!(O)$) coordinate (D)
                ($(M)!2!(O)$) coordinate (P)
                ($(N)!2!(O)$) coordinate (Q);
                \pgfresetboundingbox
                \draw (O) circle (\r);
                \draw (Az)--(At) (Bz)--(Bt) (M)--(N) (A)--(B) (C)--(D) (M)--(P) (N)--(Q)--(P) (M)--(B) (A)--(P);
                   \filldraw[black] (A) circle (2pt) node[above left] {$M$};
    \filldraw[black] (B) circle (2pt) node[above right] {$N$};
    \filldraw[black] (C) circle (2pt) node[below] {$P$};
   \filldraw[black] (D) circle (2pt) node[above right]{$Q$};
    \filldraw[black] (O) circle (2pt) node[below left] {$O$};
    \filldraw[black] (P) circle (2pt) node[above right] {$E$};
     \filldraw[black] (Q) circle (2pt) node[above left] {$F$};
      \filldraw[black] (M) circle (2pt) node[above left] {$I$};
     \filldraw[black] (N) circle (2pt) node[above right] {$J$};
            \end{tikzpicture}
}
\end{minipage}

\end{tcolorbox}
\vspace{4pt}


\begin{tcolorbox}[
    breakable,
    colback=white,
    colframe=blue!60!black,
    boxrule=0.8pt,
    arc=3pt,
    title={\textbf{Câu \#379} \quad \textbf{MÃ CÂU: \texttt{[L09\_C5\_B7\_08]}} \quad | \quad \textbf{Nhóm: \texttt{HINH\_282}} \quad | \quad \textbf{MapID: \texttt{9H5C4-1}} \hfill \textcolor{yellow!90!white}{[Hình gốc .pyc]}}
]
\textbf{Dạng toán:} Tổng hợp 8 \hfill \textit{(Tự luận -- Vận dụng cao)}
\vspace{4pt}\hrule\vspace{6pt}
\begin{minipage}[t]{0.58\textwidth}
\vspace{0pt}
Cho hai đường tròn $(J;R)$ và $(I; R')$ tiếp xúc ngoài với nhau tại ${A}$, hai điểm $B \in (J)$ và $C \in (I)$ sao cho ${B}$ và ${C}$ nằm cùng phía đối với đường thẳng ${JI}$ và $JB\parallel IC$.

a) Chứng minh $\widehat{BAC}$ là góc vuông.

b) Cho biết $R=28$ cm, $R'=23$ cm và ${BC}$ cắt ${JI}$ tại ${D}$. Tính độ dài đoạn ${JD}$.\\
\end{minipage}\hfill
\begin{minipage}[t]{0.39\textwidth}
\vspace{0pt}
\centering

\resizebox{\linewidth}{!}{%
\begin{tikzpicture}[scale=.8]
    \def\r{3}
    \def\R{1.5}
    \path 
    (0,0) coordinate (O)
    (4.5,0) coordinate (O')
    (3,0) coordinate (A)
    (-1,3) coordinate (Bt)
    (3.5,3) coordinate (Ct)
    ($(A)!2!(O)$) coordinate (E)
    ;
    \path[name path=c1] (O) circle (\r);
    \path[name path = c2] (O') circle (\R);
    \path[name path = bt] (O)--(Bt);
    \path[name path = ct] (O')--(Ct);
    \path[name intersections={of= c1 and bt, by={B}}];
    \path[name intersections={of= c2 and ct, by={C}}]
    (intersection of O--O' and B--C) coordinate (D);
    \draw (O) circle (\r) (O') circle (\R);
    \draw (B)--(O)--(D)--cycle (O)--(E) (B)--(A)--(C)--(O');
        \node at (2.2, 0.2) {$1$};  
                \node at (4.2, 0.2) {$1$};  
                \node at (3.4, 0.2) {$2$};  
                \node at (0.2, 0.2) {$1$};  
                \node at (-0.5, 0.2) {$2$};  
                \node at (-0.4, 2.2) {$1$};  
 

    \filldraw[black] (A) circle (2pt) node[below left] {$A$};
    \filldraw[black] (B) circle (2pt) node[above left] {$B$};
    \filldraw[black] (C) circle (2pt) node[above left] {$C$};
   \filldraw[black] (D) circle (2pt) node[above right]{$D$};
    \filldraw[black] (O) circle (2pt) node[below ] {$J$};
    \filldraw[black] (O') circle (2pt) node[below] {$I$};
   
            \end{tikzpicture}
}
\end{minipage}

\end{tcolorbox}
\vspace{4pt}


\begin{tcolorbox}[
    breakable,
    colback=white,
    colframe=blue!60!black,
    boxrule=0.8pt,
    arc=3pt,
    title={\textbf{Câu \#380} \quad \textbf{MÃ CÂU: \texttt{[L09\_C5\_B7\_09]}} \quad | \quad \textbf{Nhóm: \texttt{HINH\_283}} \quad | \quad \textbf{MapID: \texttt{9H5C4-1}} \hfill \textcolor{yellow!90!white}{[Hình gốc .pyc]}}
]
\textbf{Dạng toán:} Tổng hợp 9 \hfill \textit{(Tự luận -- Vận dụng cao)}
\vspace{4pt}\hrule\vspace{6pt}
\begin{minipage}[t]{0.58\textwidth}
\vspace{0pt}
Cho đường tròn tâm ${O}$, đường kính ${MN}$. Một đường tròn $({N})$ cắt $({O})$ tại ${I}$ và ${J}$.

a) Chứng minh rằng ${MI}$ và ${MJ}$ là hai tiếp tuyến của $({N})$.

b) Đường thẳng qua ${N}$ và vuông góc với ${NI}$ cắt ${MJ}$ tại ${Q}$. Chứng minh hai điểm ${M}$ và ${N}$ đối xứng với nhau qua ${OQ}$.

c) Đường thẳng qua ${M}$ và vuông góc với ${MI}$ cắt ${NJ}$ tại ${K}$. Chứng minh ba điểm ${O}$, ${Q}$ và ${K}$ thẳng hàng.\\
\end{minipage}\hfill
\begin{minipage}[t]{0.39\textwidth}
\vspace{0pt}
\centering

\resizebox{\linewidth}{!}{%
\begin{tikzpicture}[scale=.8]
      
                \def\r{3}
                \def\R{2}
                \path 
                (0:0) coordinate (O)
                (180:\r) coordinate (M)
                (0:\r) coordinate (N)
                ($(O)!1.7!(N)$) coordinate (Ot)
                ;
                \path[name path=c1] (O) circle (\r);
                \path[name path = c2] (N) circle (\R);
                \path[name intersections={of= c1 and c2, by={A,B}}]
                ($(N)!1mm!90:(A)$) coordinate (Nt)
                (intersection of N--Nt and M--B) coordinate (C)
                ($(M)!1mm!90:(A)$) coordinate (Mt)
                (intersection of M--Mt and N--B) coordinate (D);
                \draw (O) circle (\r) (N) circle (\R);
                \draw (M)--(Ot) (M)--(A)--(N) (M)--(B)--(N)--(C) (M)--(D)--(B) (O)--(D);

    \pic[draw, angle radius=2mm]{right angle = D--M--A};
    \pic[draw, angle radius=2mm]{right angle = A--N--C};
    \filldraw[black] (A) circle (2pt) node[above right] {$I$};
    \filldraw[black] (B) circle (2pt) node[below right] {$J$};
    \filldraw[black] (C) circle (2pt) node[below left] {$Q$};
   \filldraw[black] (D) circle (2pt) node[below]{$K$};
    \filldraw[black] (O) circle (2pt) node[above ] {$O$};
        \filldraw[black] (M) circle (2pt) node[below  left] {$M$};
    \filldraw[black] (N) circle (2pt) node[below right] {$N$};
          
            \node at (2, -0.2) {$1$};  
                \node at (-2, 0.2) {$2$};  
                \node at (-2, -0.2) {$1$};  
            \end{tikzpicture}
}
\end{minipage}

\end{tcolorbox}
\vspace{4pt}


\subsection{Chương 6 - Hàm số y=$ax^2$. Phương trình bậc hai 1 ẩn}

\begin{tcolorbox}[
    breakable,
    colback=white,
    colframe=blue!60!black,
    boxrule=0.8pt,
    arc=3pt,
    title={\textbf{Câu \#381} \quad \textbf{MÃ CÂU: \texttt{[L09\_C6\_B1\_11]}} \quad | \quad \textbf{Nhóm: \texttt{HINH\_284}} \quad | \quad \textbf{MapID: \texttt{9D6H1-1}} \hfill \textcolor{yellow!90!white}{[Hình gốc .pyc]}}
]
\textbf{Dạng toán:} Cho đồ thị. Tìm công thức hàm số \hfill \textit{(Trắc nghiệm -- Thông hiểu)}
\vspace{4pt}\hrule\vspace{6pt}
\begin{minipage}[t]{0.58\textwidth}
\vspace{0pt}
Biết đường cong sau là đồ thị của hàm số $y=ax^{2}$. Tìm hệ số ${a}$

\vspace{4pt}\noindent
\textbf{A.} ${- \dfrac{3}{2}}$ \quad \textbf{B.} ${- \dfrac{5}{2}}$ \quad \textbf{\textcolor{blue!75!black}{C. (Đ)}} ${- \dfrac{1}{2}}$ \quad \textbf{D.} ${\dfrac{1}{2}}$
\end{minipage}\hfill
\begin{minipage}[t]{0.39\textwidth}
\vspace{0pt}
\centering

\resizebox{\linewidth}{!}{%
\begin{tikzpicture}[scale=1, font=\footnotesize, line join=round, line cap=round,>=stealth]
    \def\xmin{-6.5};\def\ymin{-7};\def\xmax{6.5};\def\ymax{1};
    \coordinate (O) at (0,0);
    \draw[->] (\xmin,0)--(\xmax,0) node[below]{$x$};
    \draw[->] (0,\ymin)--(0,\ymax) node[left]{$y$};
    \fill (O) node[below left]{$O$} circle(1pt);
    \clip ({\xmin-0.1},{\ymin-0.5}) rectangle ({\xmax+0.1},{\ymax+0.1});
    \foreach \x in {-6,-5,-4,-3,-2,-1,1,2,3,4,5,6}{
        \fill (\x,0) node[below]{$\x$} circle(1pt);
    }
    \foreach \y in {-6,-5,-4,-3,-2,-1,1}{
        \fill (0,\y) node[left]{$\y$} circle(1pt);
    }
    \draw[thick,samples=100] plot[domain=-5.3:5.3](\x,{-0.5*(\x)^2});

    \draw[dashed] (-2,0)--(-2,-2)--(0,-2);
    \end{tikzpicture}
}
\end{minipage}

\end{tcolorbox}
\vspace{4pt}


\begin{tcolorbox}[
    breakable,
    colback=white,
    colframe=blue!60!black,
    boxrule=0.8pt,
    arc=3pt,
    title={\textbf{Câu \#382} \quad \textbf{MÃ CÂU: \texttt{[L09\_C6\_B1\_31]}} \quad | \quad \textbf{Nhóm: \texttt{HINH\_285}} \quad | \quad \textbf{MapID: \texttt{9D6V1-1}} \hfill \textcolor{yellow!90!white}{[Hình gốc .pyc]}}
]
\textbf{Dạng toán:} TF-M3. Cho đồ thị y=$ax^2$ Tìm hệ số điểm thuộc dthi giá trị hàm số \hfill \textit{(Đúng-Sai -- Vận dụng)}
\vspace{4pt}\hrule\vspace{6pt}
\begin{minipage}[t]{0.58\textwidth}
\vspace{0pt}
Biết đường cong sau là đồ thị của hàm số $y=ax^{2}$.

\vspace{4pt}\noindent
\textbf{a)} Đồ thị hàm số không đi qua điểm có toạ độ $(1; -5)$\\[2pt]
\textbf{\textcolor{blue!75!black}{b) (Đ)}} $a=-5$\\[2pt]
\textbf{\textcolor{blue!75!black}{c) (Đ)}} Tung độ của điểm thuộc đồ thị hàm số có hoành độ ${- \dfrac{1}{5}}$ là  ${- \dfrac{1}{5}}$\\[2pt]
\textbf{\textcolor{blue!75!black}{d) (Đ)}} Các điểm thuộc đồ thị có tung độ $-80$ là $ \left( -4;-80 \right) $; $ \left( 4;-80 \right) $
\end{minipage}\hfill
\begin{minipage}[t]{0.39\textwidth}
\vspace{0pt}
\centering

\resizebox{\linewidth}{!}{%
\begin{tikzpicture}[scale=1, font=\footnotesize, line join=round, line cap=round,>=stealth]
    \def\xmin{-6.5};\def\ymin{-7};\def\xmax{6.5};\def\ymax{1};
    \coordinate (O) at (0,0);
    \draw[->] (\xmin,0)--(\xmax,0) node[below]{$x$};
    \draw[->] (0,\ymin)--(0,\ymax) node[left]{$y$};
    \fill (O) node[below left]{$O$} circle(1pt);
    \clip ({\xmin-0.1},{\ymin-0.5}) rectangle ({\xmax+0.1},{\ymax+0.1});
    \foreach \x in {-6,-5,-4,-3,-2,-1,1,2,3,4,5,6}{
        \fill (\x,0) node[below]{$\x$} circle(1pt);
        }
    \foreach \y in {-6,-5,-4,-3,-2,-1,1}{
        \fill (0,\y) node[left]{$\y$} circle(1pt);
    }
    \draw[thick,samples=100] plot[domain=-6.3:6.3](\x,{-5*(\x)^2}) ;

    \draw[dashed] (1,0)--(1,-5)--(0,-5);
    \end{tikzpicture}
}
\end{minipage}

\end{tcolorbox}
\vspace{4pt}


\begin{tcolorbox}[
    breakable,
    colback=white,
    colframe=blue!60!black,
    boxrule=0.8pt,
    arc=3pt,
    title={\textbf{Câu \#383} \quad \textbf{MÃ CÂU: \texttt{[L09\_C6\_B1\_24]}} \quad | \quad \textbf{Nhóm: \texttt{HINH\_286}} \quad | \quad \textbf{MapID: \texttt{9D6V1-1}} \hfill \textcolor{yellow!90!white}{[Hình gốc .pyc]}}
]
\textbf{Dạng toán:} Cho đồ thị giải các bài toán liên quan \hfill \textit{(Tự luận -- Vận dụng)}
\vspace{4pt}\hrule\vspace{6pt}
\begin{minipage}[t]{0.58\textwidth}
\vspace{0pt}
Biết đường cong sau là đồ thị của hàm số $y=ax^{2}$. 

 a) Tìm hệ số ${a}$.

 b) Tìm tung độ của điểm thuộc đồ thị hàm số có hoành độ $-1$ .

c) Tìm các điểm thuộc parabol có tung độ ${\dfrac{1}{180}}$ .
\end{minipage}\hfill
\begin{minipage}[t]{0.39\textwidth}
\vspace{0pt}
\centering

\resizebox{\linewidth}{!}{%
\begin{tikzpicture}[scale=1, font=\footnotesize, line join=round, line cap=round,>=stealth]
    \def\xmin{-6.5};\def\ymin{-1};\def\xmax{6.5};\def\ymax{7};
    \coordinate (O) at (0,0);
    \draw[->] (\xmin,0)--(\xmax,0) node[below]{$x$};
    \draw[->] (0,\ymin)--(0,\ymax) node[left]{$y$};
    \fill (O) node[below left]{$O$} circle(1pt);
    \clip ({\xmin-0.1},{\ymin-0.5}) rectangle ({\xmax+0.1},{\ymax+0.1});
    \foreach \x in {-6,-5,-4,-3,-2,-1,1,2,3,4,5,6}{
        \fill (\x,0) node[below]{$\x$} circle(1pt);
    }
    \foreach \y in {1,2,3,4,5,6}{
        \fill (0,\y) node[left]{$\y$} circle(1pt);
    }
    \draw[thick,samples=100] plot[domain=-5.3:5.3](\x,{0.2*(\x)^2});

    \draw[dashed] (5,0)--(5,5)--(0,5);
    \end{tikzpicture}
}
\end{minipage}

\end{tcolorbox}
\vspace{4pt}


\begin{tcolorbox}[
    breakable,
    colback=white,
    colframe=blue!60!black,
    boxrule=0.8pt,
    arc=3pt,
    title={\textbf{Câu \#384} \quad \textbf{MÃ CÂU: \texttt{[L09\_C6\_B1\_25]}} \quad | \quad \textbf{Nhóm: \texttt{HINH\_287}} \quad | \quad \textbf{MapID: \texttt{9D6C1-7}} \hfill \textcolor{yellow!90!white}{[Hình gốc .pyc]}}
]
\textbf{Dạng toán:} Toán thực tế liên quan đến đồ thị y=$ax^2$ với a<0 \hfill \textit{(Tự luận -- Vận dụng cao)}
\vspace{4pt}\hrule\vspace{6pt}
\begin{minipage}[t]{0.58\textwidth}
\vspace{0pt}
Bạn Linh lắp ghép một mô hình cổng vòm được thiết kế dạng parabol như Hình vẽ mô phỏng. Biết chiều rộng của mô hình là $AB=48$ cm và chiều cao của mô hình là $OI=1152$ cm.

Bạn Linh có một mô hình ô tô tải có chiều rộng $18$ cm và chiều cao $28$ cm. Hỏi mô hình ô tô tải có đi qua được mô hình cổng vòm trên không?
\end{minipage}\hfill
\begin{minipage}[t]{0.39\textwidth}
\vspace{0pt}
\centering

\resizebox{\linewidth}{!}{%
\begin{tikzpicture}
    
    % Vẽ parabol
    \draw[thick, blue] 
    plot[domain=-2.83:2.83, samples=100] (\x, {-0.5*\x*\x+4}); % Parabol y = -0.5x^2 + 4
    
    % Đánh dấu và ghi nhãn đỉnh parabol
    \fill[red] (0, 4) circle(2pt); % Đánh dấu đỉnh
    \node[above left] at (0, 4.2) {\( O \)}; % Nhãn đỉnh là I
        \fill[red] (0, 0) circle(2pt); % Đánh dấu đỉnh
    \node[above left ] at (0, 0.2) {\( I \)}; % Nhãn đỉnh là I
    % Đánh dấu và ghi nhãn hai chân cổng
    \fill[green] (-2.83, 0) circle(2pt); % Chân trái (x = -sqrt(8))
    \node[below left] at (-2.83, 0) {\( A \)}; % Nhãn chân trái là A
    
    \fill[green] (2.83, 0) circle(2pt); % Chân phải (x = sqrt(8))
    \node[below right] at (2.83, 0) {\( B \)}; % Nhãn chân phải là B
    
    % Nối hai chân cổng bằng nét đứt
    \draw[dashed, gray] (-2.83, 0) -- (2.83, 0); % 
        \draw[dashed, gray] (0, 0) -- (0, 4); 
        
\end{tikzpicture}
}
\end{minipage}

\end{tcolorbox}
\vspace{4pt}


\begin{tcolorbox}[
    breakable,
    colback=white,
    colframe=blue!60!black,
    boxrule=0.8pt,
    arc=3pt,
    title={\textbf{Câu \#385} \quad \textbf{MÃ CÂU: \texttt{[L09\_C6\_B1\_26]}} \quad | \quad \textbf{Nhóm: \texttt{HINH\_288}} \quad | \quad \textbf{MapID: \texttt{9D6C1-7}} \hfill \textcolor{yellow!90!white}{[Hình gốc .pyc]}}
]
\textbf{Dạng toán:} Toán thực tế liên quan đến đồ thị y=$ax^2$ với a>0 \hfill \textit{(Tự luận -- Vận dụng cao)}
\vspace{4pt}\hrule\vspace{6pt}
\begin{minipage}[t]{0.58\textwidth}
\vspace{0pt}
Một mô hình mô phỏng cây cầu treo có trụ tháp đôi cao $BD=AC=\dfrac{49}{5}$ m và cách nhau $AB=14$ m. Các dây cáp có dạng đồ thị của hàm số $y=ax^{2}$ như hình. Một dây nối hai điểm ${M}$ và ${N}$ trên dây cáp như hình. Biết dây nối cách mặt của cây cầu là ${\dfrac{5}{4}}$. Tính độ dài dây nối hai điểm ${M}$ và ${N}$.
\end{minipage}\hfill
\begin{minipage}[t]{0.39\textwidth}
\vspace{0pt}
\centering

\resizebox{\linewidth}{!}{%
\begin{tikzpicture}
            
            % Vẽ parabol
            \draw[thick, blue] 
            plot[domain=-4:4, samples=100] (\x, {0.12*\x*\x}); % Parabol y = -0.5x^2 + 4
            % Đánh dấu và ghi nhãn hai chân cổng
            \fill[green] (-4, 1.92) circle(2pt); % Chân trái (x = -sqrt(8))
            \node[below left] at (-4, 1.92) {\( A \)}; % Nhãn chân trái là A
            
            \fill[green] (4, 1.92) circle(2pt); % Chân phải (x = sqrt(8))
            \node[below right] at (4, 1.92) {\( B \)}; % Nhãn chân phải là B
            
            % Nối hai chân cổng bằng nét đứt
            \draw[dashed, gray] (-4, 1.92) -- (4, 1.92); % 
            \draw[thick](-4, 0) -- (4, 0); % 
                \draw[dashed, gray](-4, 1.92) -- (4, 1.92); % 
                \draw[thick](-4, 1.92) -- (-4, 0); %
                \draw[thick](4, 1.92) -- (4, 0); %
                    \fill[green] (-4, 0) circle(2pt); % Chân trái (x = -sqrt(8))
                \node[below left] at (-4, 0) {\( C \)}; % Nhãn chân trái là A
                
                \fill[green] (4, 0) circle(2pt); % Chân phải (x = sqrt(8))
                \node[below right] at (4, 0) {\( D \)}; % Nhãn chân phải là B
                    \draw[dashed, gray] (-2, 0.48) -- (2, 0.48); % 
                            \fill[green] (-2, 0.48) circle(2pt); % Chân trái (x = -sqrt(8))
                    \node[below left] at (-2, 0.48) {\( M \)}; % Nhãn chân trái là A
                    
                    \fill[green] (2, 0.48) circle(2pt); % Chân phải (x = sqrt(8))
                    \node[below right] at (2, 0.48) {\(N \)}; % Nhãn chân phải là B
        \end{tikzpicture}
}
\end{minipage}

\end{tcolorbox}
\vspace{4pt}


\begin{tcolorbox}[
    breakable,
    colback=white,
    colframe=blue!60!black,
    boxrule=0.8pt,
    arc=3pt,
    title={\textbf{Câu \#386} \quad \textbf{MÃ CÂU: \texttt{[L09\_C6\_B4\_35]}} \quad | \quad \textbf{Nhóm: \texttt{HINH\_289}} \quad | \quad \textbf{MapID: \texttt{9D6V4-7}} \hfill \textcolor{yellow!90!white}{[Hình gốc .pyc]}}
]
\textbf{Dạng toán:} Bài toán thực tế về diện tích đẩt xây nhà sân vườn \hfill \textit{(Trắc nghiệm -- Vận dụng)}
\vspace{4pt}\hrule\vspace{6pt}
\begin{minipage}[t]{0.58\textwidth}
\vspace{0pt}
Cô Dung có một mảnh đất hình chữ nhật với chiều dài $32$ m và chiều rộng $24$ m. Cô dự định xây nhà trên mảnh đất đó và dành một phần diện tích đất để làm sân vườn như Hình dưới đây. Biết diện tích đất làm nhà là ${270}m^{2}$. Hỏi ${x}$ bằng bao nhiêu mét?

\vspace{4pt}\noindent
\textbf{A.} $13$ \quad \textbf{B.} $11$ \quad \textbf{C.} $10$ \quad \textbf{\textcolor{blue!75!black}{D. (Đ)}} $9$
\end{minipage}\hfill
\begin{minipage}[t]{0.39\textwidth}
\vspace{0pt}
\centering

\resizebox{\linewidth}{!}{%
\begin{tikzpicture}
    % Vẽ hình sân vườn ngoài cùng
    \draw[thick] (0, 0) rectangle (10, 7.5);
    \node at (5, 6.5) {Sân vuon};
    
    % Vẽ hình nhà bên trong, góc dưới bên trái của nhà trùng với góc dưới bên trái của sân vườn
    \draw[thick] (0, 0) rectangle (7, 5.5);
    \node at (3.5, 2.25) {Nhà};
    
    % Ghi chú các chiều
    \draw[<->] (0, 8) -- (10, 8);
    \node at (5, 8.3) {$32$m};
    
    \draw[<->] (10.5, 0) -- (10.5, 7.5);
    \node at (11, 3.75) {$24$ m};
    
    \draw[<->] (0, 5.5) -- (0, 7.5);
    \node at (-0.4, 6.25) {x};
    
    \draw[<->] (7, -0.5) -- (10, -0.5);
    \node at (8.5, -0.8) {x+5};
    
\end{tikzpicture}
}
\end{minipage}

\end{tcolorbox}
\vspace{4pt}


\begin{tcolorbox}[
    breakable,
    colback=white,
    colframe=blue!60!black,
    boxrule=0.8pt,
    arc=3pt,
    title={\textbf{Câu \#387} \quad \textbf{MÃ CÂU: \texttt{[L09\_C6\_B4\_60]}} \quad | \quad \textbf{Nhóm: \texttt{HINH\_290}} \quad | \quad \textbf{MapID: \texttt{9D6V4-7}} \hfill \textcolor{yellow!90!white}{[Hình gốc .pyc]}}
]
\textbf{Dạng toán:} SA-M3. Bài toán thực tế về diện tích đẩt xây nhà, sân vườn \hfill \textit{(Trả lời ngắn -- Vận dụng)}
\vspace{4pt}\hrule\vspace{6pt}
\begin{minipage}[t]{0.58\textwidth}
\vspace{0pt}
Cô Dung có một mảnh đất hình chữ nhật với chiều dài $30$ m và chiều rộng $22$ m. Cô dự định xây nhà trên mảnh đất đó và dành một phần diện tích đất để làm sân vườn như Hình dưới đây. Biết diện tích đất làm nhà là ${260}m^{2}$. Hỏi ${x}$ bằng bao nhiêu mét?


\textbf{Đáp án điền số:} 	exttt{9}
\end{minipage}\hfill
\begin{minipage}[t]{0.39\textwidth}
\vspace{0pt}
\centering

\resizebox{\linewidth}{!}{%
\begin{tikzpicture}
    % Vẽ hình sân vườn ngoài cùng
    \draw[thick] (0, 0) rectangle (10, 7.5);
    \node at (5, 6.5) {Sân vuon};
    
    % Vẽ hình nhà bên trong, góc dưới bên trái của nhà trùng với góc dưới bên trái của sân vườn
    \draw[thick] (0, 0) rectangle (7, 5.5);
    \node at (3.5, 2.25) {Nhà};
    
    % Ghi chú các chiều
    \draw[<->] (0, 8) -- (10, 8);
    \node at (5, 8.3) {$30$m};
    
    \draw[<->] (10.5, 0) -- (10.5, 7.5);
    \node at (11, 3.75) {$22$ m};
    
    \draw[<->] (0, 5.5) -- (0, 7.5);
    \node at (-0.4, 6.25) {x};
    
    \draw[<->] (7, -0.5) -- (10, -0.5);
    \node at (8.5, -0.8) {x+1};
    
\end{tikzpicture}
}
\end{minipage}

\end{tcolorbox}
\vspace{4pt}


\begin{tcolorbox}[
    breakable,
    colback=white,
    colframe=blue!60!black,
    boxrule=0.8pt,
    arc=3pt,
    title={\textbf{Câu \#388} \quad \textbf{MÃ CÂU: \texttt{[L09\_C6\_B4\_21]}} \quad | \quad \textbf{Nhóm: \texttt{HINH\_291}} \quad | \quad \textbf{MapID: \texttt{9D6V4-7}} \hfill \textcolor{yellow!90!white}{[Hình gốc .pyc]}}
]
\textbf{Dạng toán:} Bài toán thực tế về diện tích đẩt xây nhà sân vườn \hfill \textit{(Tự luận -- Vận dụng)}
\vspace{4pt}\hrule\vspace{6pt}
\begin{minipage}[t]{0.58\textwidth}
\vspace{0pt}
Cô Dung có một mảnh đất hình chữ nhật với chiều dài $25$ m và chiều rộng $20$ m. Cô dự định xây nhà trên mảnh đất đó và dành một phần diện tích đất để làm sân vườn như Hình dưới đây. Biết diện tích đất làm nhà là ${221}m^{2}$. Hỏi ${x}$ bằng bao nhiêu mét?
\end{minipage}\hfill
\begin{minipage}[t]{0.39\textwidth}
\vspace{0pt}
\centering

\resizebox{\linewidth}{!}{%
\begin{tikzpicture}
    % Vẽ hình sân vườn ngoài cùng
    \draw[thick] (0, 0) rectangle (10, 7.5);
    \node at (5, 6.5) {Sân vuon};
    
    % Vẽ hình nhà bên trong, góc dưới bên trái của nhà trùng với góc dưới bên trái của sân vườn
    \draw[thick] (0, 0) rectangle (7, 5.5);
    \node at (3.5, 2.25) {Nhà};
    
    % Ghi chú các chiều
    \draw[<->] (0, 8) -- (10, 8);
    \node at (5, 8.3) {$25$m};
    
    \draw[<->] (10.5, 0) -- (10.5, 7.5);
    \node at (11, 3.75) {$20$ m};
    
    \draw[<->] (0, 5.5) -- (0, 7.5);
    \node at (-0.4, 6.25) {x};
    
    \draw[<->] (7, -0.5) -- (10, -0.5);
    \node at (8.5, -0.8) {x+1};
    
\end{tikzpicture}
}
\end{minipage}

\end{tcolorbox}
\vspace{4pt}


\subsection{Chương 7 - Tần số và tần số tương đối}

\begin{tcolorbox}[
    breakable,
    colback=white,
    colframe=blue!60!black,
    boxrule=0.8pt,
    arc=3pt,
    title={\textbf{Câu \#389} \quad \textbf{MÃ CÂU: \texttt{[L09\_C7\_B1\_12]}} \quad | \quad \textbf{Nhóm: \texttt{HINH\_292}} \quad | \quad \textbf{MapID: \texttt{9X7N1-1}} \hfill \textcolor{yellow!90!white}{[Hình gốc .pyc]}}
]
\textbf{Dạng toán:} CHo b.đồ cột t.kê số từ mới học mỗi ngày. Tìm tần số \hfill \textit{(Trắc nghiệm -- Nhận biết)}
\vspace{4pt}\hrule\vspace{6pt}
\begin{minipage}[t]{0.58\textwidth}
\vspace{0pt}
Vào đợt nghỉ hè vừa rồi, mỗi ngày bạn Lam đều học thêm một số từ vựng tiếng Anh mới. Số lượng từ vựng mới bạn Lam học mỗi ngày được biểu diễn ở biểu đồ cột như hình dưới.

Tìm tần số của giá trị $8$.

\vspace{4pt}\noindent
\textbf{A.} $7$ \quad \textbf{B.} $3$ \quad \textbf{\textcolor{blue!75!black}{C. (Đ)}} $4$ \quad \textbf{D.} $5$
\end{minipage}\hfill
\begin{minipage}[t]{0.39\textwidth}
\vspace{0pt}
\centering

\resizebox{\linewidth}{!}{%
\begin{tikzpicture}[y=0.3cm,x=0.8cm]
        %\draw[color=gray!50,line width=1.5pt]  (-1,-2.4)--(14.6,-2.4)--(14.6,6.5)--(-1,6.5)--cycle ;
        \foreach \i  in {2,4,...,10} 
        {
            \draw (0,\i) node[left]{$\i$};
            \draw (-0.1,\i)--(0.1,\i) ;
            \draw[gray!40] (0.5,\i)--(11.5,\i);
        }
        
        \foreach \x/\y/\z in {2/5/10,4/6/7,6/8/4,8/14/10,10/15/6} 
        {
            \draw ({\x},0) node[below]{$\y$};
            \draw ({\x},{\z}) node[above]{$\z$};
        }
        \draw [->] (0,0)--(11.5,0) node[below right]{Số từ mới};
        \draw [->] (0,0)--(0,11) node[above]{\text{Số ngày}};
        \draw [pattern=north east lines]  (1.6,0)--(1.6,10)--(2.4,10)--(2.4,0)--cycle   
        (3.6,0)--(3.6,7)--(4.4,7)--(4.4,0)--cycle  
        (5.6,0)--(5.6,4)--(6.4,4)--(6.4,0)--cycle   
         (7.6,0)--(7.6,10)--(8.4,10)--(8.4,0)--cycle   
         (9.6,0)--(9.6,6)--(10.4,6)--(10.4,0)--cycle   ;
        %[pattern=north east lines]
    \end{tikzpicture}
}
\end{minipage}

\end{tcolorbox}
\vspace{4pt}


\begin{tcolorbox}[
    breakable,
    colback=white,
    colframe=blue!60!black,
    boxrule=0.8pt,
    arc=3pt,
    title={\textbf{Câu \#390} \quad \textbf{MÃ CÂU: \texttt{[L09\_C7\_B1\_13]}} \quad | \quad \textbf{Nhóm: \texttt{HINH\_293}} \quad | \quad \textbf{MapID: \texttt{9X7N1-1}} \hfill \textcolor{yellow!90!white}{[Hình gốc .pyc]}}
]
\textbf{Dạng toán:} CHo b.đồ cột t.kê điểm thi. Tìm tần số \hfill \textit{(Trắc nghiệm -- Nhận biết)}
\vspace{4pt}\hrule\vspace{6pt}
\begin{minipage}[t]{0.58\textwidth}
\vspace{0pt}
Nhà trường thống kê điểm số của học sinh khối $7$ sau kì thi khảo sát chất lượng được biểu diễn ở biểu đồ cột như hình dưới.

Tìm tần số của giá trị $9$.

\vspace{4pt}\noindent
\textbf{\textcolor{blue!75!black}{A. (Đ)}} $13$ \quad \textbf{B.} $15$ \quad \textbf{C.} $11$ \quad \textbf{D.} $14$
\end{minipage}\hfill
\begin{minipage}[t]{0.39\textwidth}
\vspace{0pt}
\centering

\resizebox{\linewidth}{!}{%
\begin{tikzpicture}[y=0.1cm,x=0.7cm]
        %\draw[color=gray!50,line width=1.5pt]  (-1,-2.4)--(14.6,-2.4)--(14.6,6.5)--(-1,6.5)--cycle ;
        \foreach \i  in {10,20,...,50} 
        {
            \draw (0,\i) node[left]{$\i$};
            \draw (-0.1,\i)--(0.1,\i) ;
           
        }
        
        \foreach \x/\y/\z in {2/4/48,4/5/50,6/6/47,8/7/6,10/8/11,12/9/13,14/10/27} 
        {
            \draw ({\x},0) node[below]{$\y$};
            \draw ({\x},{\z}) node[above]{$\z$};
        }
        \draw [->] (0,0)--(15.5,0) node[below right]{Điểm};
        \draw [->] (0,0)--(0,54) node[above]{\text{Số học sinh}};
        \draw [pattern=north east lines]  (1.6,0)--(1.6,48)--(2.4,48)--(2.4,0)--cycle   
        (3.6,0)--(3.6,50)--(4.4,50)--(4.4,0)--cycle  
        (5.6,0)--(5.6,47)--(6.4,47)--(6.4,0)--cycle   
         (7.6,0)--(7.6,6)--(8.4,6)--(8.4,0)--cycle   
         (9.6,0)--(9.6,11)--(10.4,11)--(10.4,0)--cycle 
          (11.6,0)--(11.6,13)--(12.4,13)--(12.4,0)--cycle 
           (13.6,0)--(13.6,27)--(14.4,27)--(14.4,0)--cycle 
           ;
        %[pattern=north east lines]
    \end{tikzpicture}
}
\end{minipage}

\end{tcolorbox}
\vspace{4pt}


\begin{tcolorbox}[
    breakable,
    colback=white,
    colframe=blue!60!black,
    boxrule=0.8pt,
    arc=3pt,
    title={\textbf{Câu \#391} \quad \textbf{MÃ CÂU: \texttt{[L09\_C7\_B1\_18]}} \quad | \quad \textbf{Nhóm: \texttt{HINH\_294}} \quad | \quad \textbf{MapID: \texttt{9X7N1-1}} \hfill \textcolor{yellow!90!white}{[Hình gốc .pyc]}}
]
\textbf{Dạng toán:} CHo b.đồ cột t.kê số con. Tìm tổng số GD dc điều tra \hfill \textit{(Trắc nghiệm -- Nhận biết)}
\vspace{4pt}\hrule\vspace{6pt}
\begin{minipage}[t]{0.58\textwidth}
\vspace{0pt}
Số con của các hộ gia đình tại thôn X được thống kê bằng biểu đồ cột như hình dưới. Hỏi thôn X có tất cả bao nhiêu hộ gia đình tham gia điều tra.

\vspace{4pt}\noindent
\textbf{A.} $215$ \quad \textbf{B.} $211$ \quad \textbf{C.} $216$ \quad \textbf{\textcolor{blue!75!black}{D. (Đ)}} $213$
\end{minipage}\hfill
\begin{minipage}[t]{0.39\textwidth}
\vspace{0pt}
\centering

\resizebox{\linewidth}{!}{%
\begin{tikzpicture}[y=0.1cm,x=0.7cm]
        %\draw[color=gray!50,line width=1.5pt]  (-1,-2.4)--(14.6,-2.4)--(14.6,6.5)--(-1,6.5)--cycle ;
        \foreach \i  in {10,20,...,70} 
        {
            \draw (0,\i) node[left]{$\i$};
            \draw (-0.1,\i)--(0.1,\i) ;
           
        }
        
        \foreach \x/\y/\z in {2/0/4,4/1/64,6/2/65,8/3/6,10/4/23,12/5/12,14/6/39} 
        {
            \draw ({\x},0) node[below]{$\y$};
            \draw ({\x},{\z}) node[above]{$\z$};
        }
        \draw [->] (0,0)--(15.5,0) node[below right]{Số con};
        \draw [->] (0,0)--(0,74) node[above]{\text{Số hộ gia đình}};
        \draw [pattern=north east lines]  (1.6,0)--(1.6,4)--(2.4,4)--(2.4,0)--cycle   
        (3.6,0)--(3.6,64)--(4.4,64)--(4.4,0)--cycle  
        (5.6,0)--(5.6,65)--(6.4,65)--(6.4,0)--cycle   
         (7.6,0)--(7.6,6)--(8.4,6)--(8.4,0)--cycle   
         (9.6,0)--(9.6,23)--(10.4,23)--(10.4,0)--cycle 
          (11.6,0)--(11.6,12)--(12.4,12)--(12.4,0)--cycle 
           (13.6,0)--(13.6,39)--(14.4,39)--(14.4,0)--cycle 
           ;
        %[pattern=north east lines]
    \end{tikzpicture}
}
\end{minipage}

\end{tcolorbox}
\vspace{4pt}


\begin{tcolorbox}[
    breakable,
    colback=white,
    colframe=blue!60!black,
    boxrule=0.8pt,
    arc=3pt,
    title={\textbf{Câu \#392} \quad \textbf{MÃ CÂU: \texttt{[L09\_C7\_B1\_19]}} \quad | \quad \textbf{Nhóm: \texttt{HINH\_295}} \quad | \quad \textbf{MapID: \texttt{9X7H1-1}} \hfill \textcolor{yellow!90!white}{[Hình gốc .pyc]}}
]
\textbf{Dạng toán:} CHo b.đồ cột t.kê số con. Hầu hết các GD sinh bn con? \hfill \textit{(Trắc nghiệm -- Thông hiểu)}
\vspace{4pt}\hrule\vspace{6pt}
\begin{minipage}[t]{0.58\textwidth}
\vspace{0pt}
Số con của các hộ gia đình tại thôn X được thống kê bằng biểu đồ cột như hình dưới. Hầu hết các gia đình được điều tra có bao nhiêu con.

\vspace{4pt}\noindent
\textbf{\textcolor{blue!75!black}{A. (Đ)}} $2$ \quad \textbf{B.} $3$ \quad \textbf{C.} $4$ \quad \textbf{D.} $5$
\end{minipage}\hfill
\begin{minipage}[t]{0.39\textwidth}
\vspace{0pt}
\centering

\resizebox{\linewidth}{!}{%
\begin{tikzpicture}[y=0.1cm,x=0.7cm]
        %\draw[color=gray!50,line width=1.5pt]  (-1,-2.4)--(14.6,-2.4)--(14.6,6.5)--(-1,6.5)--cycle ;
        \foreach \i  in {10,20,...,40} 
        {
            \draw (0,\i) node[left]{$\i$};
            \draw (-0.1,\i)--(0.1,\i) ;
           
        }
        
        \foreach \x/\y/\z in {2/0/12,4/1/30,6/2/35,8/3/33,10/4/28,12/5/24,14/6/9} 
        {
            \draw ({\x},0) node[below]{$\y$};
            \draw ({\x},{\z}) node[above]{$\z$};
        }
        \draw [->] (0,0)--(15.5,0) node[below right]{Số con};
        \draw [->] (0,0)--(0,44) node[above]{\text{Số hộ gia đình}};
        \draw [pattern=north east lines]  (1.6,0)--(1.6,12)--(2.4,12)--(2.4,0)--cycle   
        (3.6,0)--(3.6,30)--(4.4,30)--(4.4,0)--cycle  
        (5.6,0)--(5.6,35)--(6.4,35)--(6.4,0)--cycle   
         (7.6,0)--(7.6,33)--(8.4,33)--(8.4,0)--cycle   
         (9.6,0)--(9.6,28)--(10.4,28)--(10.4,0)--cycle 
          (11.6,0)--(11.6,24)--(12.4,24)--(12.4,0)--cycle 
           (13.6,0)--(13.6,9)--(14.4,9)--(14.4,0)--cycle 
           ;
        %[pattern=north east lines]
    \end{tikzpicture}
}
\end{minipage}

\end{tcolorbox}
\vspace{4pt}


\begin{tcolorbox}[
    breakable,
    colback=white,
    colframe=blue!60!black,
    boxrule=0.8pt,
    arc=3pt,
    title={\textbf{Câu \#393} \quad \textbf{MÃ CÂU: \texttt{[L09\_C7\_B1\_06]}} \quad | \quad \textbf{Nhóm: \texttt{HINH\_296}} \quad | \quad \textbf{MapID: \texttt{9X7H1-1}} \hfill \textcolor{yellow!90!white}{[Hình gốc .pyc]}}
]
\textbf{Dạng toán:} Btoan về b.đồ cột số lượng từ vựng tiếng anh học mỗi ngày. \hfill \textit{(Tự luận -- Thông hiểu)}
\vspace{4pt}\hrule\vspace{6pt}
\begin{minipage}[t]{0.58\textwidth}
\vspace{0pt}
Vào đợt nghỉ hè vừa rồi, mỗi ngày bạn Hiếu đều học thêm một số từ vựng tiếng Anh mới. Số lượng từ vựng mới bạn Bình học mỗi ngày được biểu diễn ở biểu đồ cột như hình dưới.

a) Số lượng từ vựng mới mà bạn Hiếu học mỗi ngày nhận những giá trị nào? Tìm tần số của mỗi giá trị đó?

b) Bạn Hiếu đã học từ vựng tiếng Anh mới trong bao nhiêu ngày?

c) Có bao nhiêu ngày bạn Hiếu học nhiều hơn $12$ từ vựng mới?.
\end{minipage}\hfill
\begin{minipage}[t]{0.39\textwidth}
\vspace{0pt}
\centering

\resizebox{\linewidth}{!}{%
\begin{tikzpicture}[y=0.3cm,x=0.8cm]
        %\draw[color=gray!50,line width=1.5pt]  (-1,-2.4)--(14.6,-2.4)--(14.6,6.5)--(-1,6.5)--cycle ;
        \foreach \i  in {2,4,...,10} 
        {
            \draw (0,\i) node[left]{$\i$};
            \draw (-0.1,\i)--(0.1,\i) ;
            \draw[gray!40] (0.5,\i)--(11.5,\i);
        }
        
        \foreach \x/\y/\z in {2/7/8,4/9/9,6/12/10,8/15/3,10/19/7} 
        {
            \draw ({\x},0) node[below]{$\y$};
            \draw ({\x},{\z}) node[above]{$\z$};
        }
        \draw [->] (0,0)--(11.5,0) node[below right]{Số từ mới};
        \draw [->] (0,0)--(0,11) node[above]{\text{Số ngày}};
        \draw [pattern=north east lines]  (1.6,0)--(1.6,8)--(2.4,8)--(2.4,0)--cycle   
        (3.6,0)--(3.6,9)--(4.4,9)--(4.4,0)--cycle  
        (5.6,0)--(5.6,10)--(6.4,10)--(6.4,0)--cycle   
         (7.6,0)--(7.6,3)--(8.4,3)--(8.4,0)--cycle   
         (9.6,0)--(9.6,7)--(10.4,7)--(10.4,0)--cycle   ;
        %[pattern=north east lines]
    \end{tikzpicture}
}
\end{minipage}

\end{tcolorbox}
\vspace{4pt}


\begin{tcolorbox}[
    breakable,
    colback=white,
    colframe=blue!60!black,
    boxrule=0.8pt,
    arc=3pt,
    title={\textbf{Câu \#394} \quad \textbf{MÃ CÂU: \texttt{[L09\_C7\_B1\_23]}} \quad | \quad \textbf{Nhóm: \texttt{HINH\_297}} \quad | \quad \textbf{MapID: \texttt{9X7H1-1}} \hfill \textcolor{yellow!90!white}{[Hình gốc .pyc]}}
]
\textbf{Dạng toán:} Btoan về biểu đồ cột số năm công tác của các bác sĩ. \hfill \textit{(Tự luận -- Thông hiểu)}
\vspace{4pt}\hrule\vspace{6pt}
\begin{minipage}[t]{0.58\textwidth}
\vspace{0pt}
Biểu đồ bên dưới thống kê thời gian công tác (theo năm) của các y tá ở một phòng khám.

a) Các y tá của phòng khám có thời gian công tác nhận những giá trị nào? Tìm tần số của mỗi giá trị đó.

b) Phòng khám có tổng số bao nhiêu y tá?

c) Có bao nhiêu y tá đã công tác ở phòng khám nhiều hơn $17$ năm?
\end{minipage}\hfill
\begin{minipage}[t]{0.39\textwidth}
\vspace{0pt}
\centering

\resizebox{\linewidth}{!}{%
\begin{tikzpicture}[y=0.3cm,x=0.8cm]
        %\draw[color=gray!50,line width=1.5pt]  (-1,-2.4)--(14.6,-2.4)--(14.6,6.5)--(-1,6.5)--cycle ;
        \foreach \i  in {2,4,...,10} 
        {
            \draw (0,\i) node[left]{$\i$};
            \draw (-0.1,\i)--(0.1,\i) ;
            \draw[gray!40] (0.5,\i)--(11.5,\i);
        }
        
        \foreach \x/\y/\z in {2/6/5,4/11/3,6/14/7,8/16/7,10/17/10} 
        {
            \draw ({\x},0) node[below]{$\y$};
            \draw ({\x},{\z}) node[above]{$\z$};
        }
        \draw [->] (0,0)--(11.5,0) node[below right]{Số năm công tác};
        \draw [->] (0,0)--(0,11) node[above]{\text{Số y tá}};
        \draw [pattern=north east lines]  (1.6,0)--(1.6,5)--(2.4,5)--(2.4,0)--cycle   
        (3.6,0)--(3.6,3)--(4.4,3)--(4.4,0)--cycle  
        (5.6,0)--(5.6,7)--(6.4,7)--(6.4,0)--cycle   
         (7.6,0)--(7.6,7)--(8.4,7)--(8.4,0)--cycle   
         (9.6,0)--(9.6,10)--(10.4,10)--(10.4,0)--cycle   ;
        %[pattern=north east lines]
    \end{tikzpicture}
}
\end{minipage}

\end{tcolorbox}
\vspace{4pt}


\begin{tcolorbox}[
    breakable,
    colback=white,
    colframe=blue!60!black,
    boxrule=0.8pt,
    arc=3pt,
    title={\textbf{Câu \#395} \quad \textbf{MÃ CÂU: \texttt{[L09\_C7\_B1\_24]}} \quad | \quad \textbf{Nhóm: \texttt{HINH\_298}} \quad | \quad \textbf{MapID: \texttt{9X7H1-1}} \hfill \textcolor{yellow!90!white}{[Hình gốc .pyc]}}
]
\textbf{Dạng toán:} Btoan về biểu đồ cột số sản phẩm bán ra. \hfill \textit{(Tự luận -- Thông hiểu)}
\vspace{4pt}\hrule\vspace{6pt}
\begin{minipage}[t]{0.58\textwidth}
\vspace{0pt}
Biểu đồ bên dưới thống kê số lượng xi măng (đơn vị: Tấn) của một cửa hàng bán được trong $6$ tháng đầu năm.

a) Liệt kê số lượng xi măng bán được trong các tháng.

b) Tổng số xi măng bán được trong $6$ tháng đầu năm là bao nhiêu

c) Một người đưa ra nhận định: Số xi măng bán được trong tháng $2$ nhiều hơn ${27}\%$ số xi măng bán được trong $5$ tháng còn lại. Hỏi nhận định của người đó là đúng hay sai?
\end{minipage}\hfill
\begin{minipage}[t]{0.39\textwidth}
\vspace{0pt}
\centering

\resizebox{\linewidth}{!}{%
\begin{tikzpicture}[y=0.3cm,x=0.8cm]
        %\draw[color=gray!50,line width=1.5pt]  (-1,-2.4)--(14.6,-2.4)--(14.6,6.5)--(-1,6.5)--cycle ;
        \foreach \i  in {2,4,...,12} 
        {
            \draw (0,\i) node[left]{$\i$};
            \draw (-0.1,\i)--(0.1,\i) ;
            \draw[gray!40] (0.5,\i)--(13.5,\i);
        }
        
        \foreach \x/\y/\z in {2/1/2,4/2/6,6/3/8,8/4/12,10/5/9,12/6/12} 
        {
            \draw ({\x},0) node[below]{$\y$};
            \draw ({\x},{\z}) node[above]{$\z$};
        }
        \draw [->] (0,0)--(13.5,0) node[below right]{Tháng};
        \draw [->] (0,0)--(0,13) node[above]{\text{Sản lượng}};
        \draw [pattern=north east lines]  (1.6,0)--(1.6,2)--(2.4,2)--(2.4,0)--cycle   
        (3.6,0)--(3.6,6)--(4.4,6)--(4.4,0)--cycle  
        (5.6,0)--(5.6,8)--(6.4,8)--(6.4,0)--cycle   
         (7.6,0)--(7.6,12)--(8.4,12)--(8.4,0)--cycle   
         (9.6,0)--(9.6,9)--(10.4,9)--(10.4,0)--cycle   
         (11.6,0)--(11.6,12)--(12.4,12)--(12.4,0)--cycle   ;
        %[pattern=north east lines]
    \end{tikzpicture}
}
\end{minipage}

\end{tcolorbox}
\vspace{4pt}


\begin{tcolorbox}[
    breakable,
    colback=white,
    colframe=blue!60!black,
    boxrule=0.8pt,
    arc=3pt,
    title={\textbf{Câu \#396} \quad \textbf{MÃ CÂU: \texttt{[L09\_C7\_B2\_03]}} \quad | \quad \textbf{Nhóm: \texttt{HINH\_299}} \quad | \quad \textbf{MapID: \texttt{9X7H2-1}} \hfill \textcolor{yellow!90!white}{[Hình gốc .pyc]}}
]
\textbf{Dạng toán:} Cho b.đồ cột t.kê số từ mới học mỗi ngày. Tìm tần số t.đối \hfill \textit{(Trắc nghiệm -- Thông hiểu)}
\vspace{4pt}\hrule\vspace{6pt}
\begin{minipage}[t]{0.58\textwidth}
\vspace{0pt}
Vào đợt nghỉ hè vừa rồi, mỗi ngày bạn Minh đều học thêm một số từ vựng tiếng Anh mới. Số lượng từ vựng mới bạn Bình học mỗi ngày được biểu diễn ở biểu đồ cột như hình dưới.

Tìm tần số tương đối của giá trị $19$.

\vspace{4pt}\noindent
\textbf{\textcolor{blue!75!black}{A. (Đ)}} ${15,4\%}$ \quad \textbf{B.} ${14,9\%}$ \quad \textbf{C.} ${16,4\%}$ \quad \textbf{D.} ${15,9\%}$
\end{minipage}\hfill
\begin{minipage}[t]{0.39\textwidth}
\vspace{0pt}
\centering

\resizebox{\linewidth}{!}{%
\begin{tikzpicture}[y=0.3cm,x=0.8cm]
        %\draw[color=gray!50,line width=1.5pt]  (-1,-2.4)--(14.6,-2.4)--(14.6,6.5)--(-1,6.5)--cycle ;
        \foreach \i  in {2,4,...,16} 
        {
            \draw (0,\i) node[left]{$\i$};
            \draw (-0.1,\i)--(0.1,\i) ;
            \draw[gray!40] (0.5,\i)--(11.5,\i);
        }
        
        \foreach \x/\y/\z in {2/5/15,4/8/13,6/11/13,8/18/3,10/19/8} 
        {
            \draw ({\x},0) node[below]{$\y$};
            \draw ({\x},{\z}) node[above]{$\z$};
        }
        \draw [->] (0,0)--(11.5,0) node[below right]{Số từ mới};
        \draw [->] (0,0)--(0,17) node[above]{\text{Số ngày}};
        \draw [pattern=north east lines]  (1.6,0)--(1.6,15)--(2.4,15)--(2.4,0)--cycle   
        (3.6,0)--(3.6,13)--(4.4,13)--(4.4,0)--cycle  
        (5.6,0)--(5.6,13)--(6.4,13)--(6.4,0)--cycle   
         (7.6,0)--(7.6,3)--(8.4,3)--(8.4,0)--cycle   
         (9.6,0)--(9.6,8)--(10.4,8)--(10.4,0)--cycle   ;
        %[pattern=north east lines]
    \end{tikzpicture}
}
\end{minipage}

\end{tcolorbox}
\vspace{4pt}


\begin{tcolorbox}[
    breakable,
    colback=white,
    colframe=blue!60!black,
    boxrule=0.8pt,
    arc=3pt,
    title={\textbf{Câu \#397} \quad \textbf{MÃ CÂU: \texttt{[L09\_C7\_B2\_04]}} \quad | \quad \textbf{Nhóm: \texttt{HINH\_300}} \quad | \quad \textbf{MapID: \texttt{9X7H2-1}} \hfill \textcolor{yellow!90!white}{[Hình gốc .pyc]}}
]
\textbf{Dạng toán:} Cho b.đồ cột t.kê điểm thi. Tìm tần số t.đối \hfill \textit{(Trắc nghiệm -- Thông hiểu)}
\vspace{4pt}\hrule\vspace{6pt}
\begin{minipage}[t]{0.58\textwidth}
\vspace{0pt}
Nhà trường thống kê điểm số của học sinh khối $12$ sau kì thi khảo sát chất lượng được biểu diễn ở biểu đồ cột như hình dưới. Tìm tỉ lệ học sinh được $9$ điểm.

\vspace{4pt}\noindent
\textbf{A.} ${17,2\%}$ \quad \textbf{B.} ${15,7\%}$ \quad \textbf{\textcolor{blue!75!black}{C. (Đ)}} ${16,2\%}$ \quad \textbf{D.} ${16,7\%}$
\end{minipage}\hfill
\begin{minipage}[t]{0.39\textwidth}
\vspace{0pt}
\centering

\resizebox{\linewidth}{!}{%
\begin{tikzpicture}[y=0.1cm,x=0.7cm]
        %\draw[color=gray!50,line width=1.5pt]  (-1,-2.4)--(14.6,-2.4)--(14.6,6.5)--(-1,6.5)--cycle ;
        \foreach \i  in {10,20,...,60} 
        {
            \draw (0,\i) node[left]{$\i$};
            \draw (-0.1,\i)--(0.1,\i) ;
           
        }
        
        \foreach \x/\y/\z in {2/4/33,4/5/54,6/6/60,8/7/54,10/8/51,12/9/55,14/10/33} 
        {
            \draw ({\x},0) node[below]{$\y$};
            \draw ({\x},{\z}) node[above]{$\z$};
        }
        \draw [->] (0,0)--(15.5,0) node[below right]{Điểm};
        \draw [->] (0,0)--(0,64) node[above]{\text{Số học sinh}};
        \draw [pattern=north east lines]  (1.6,0)--(1.6,33)--(2.4,33)--(2.4,0)--cycle   
        (3.6,0)--(3.6,54)--(4.4,54)--(4.4,0)--cycle  
        (5.6,0)--(5.6,60)--(6.4,60)--(6.4,0)--cycle   
         (7.6,0)--(7.6,54)--(8.4,54)--(8.4,0)--cycle   
         (9.6,0)--(9.6,51)--(10.4,51)--(10.4,0)--cycle 
          (11.6,0)--(11.6,55)--(12.4,55)--(12.4,0)--cycle 
           (13.6,0)--(13.6,33)--(14.4,33)--(14.4,0)--cycle 
           ;
        %[pattern=north east lines]
    \end{tikzpicture}
}
\end{minipage}

\end{tcolorbox}
\vspace{4pt}


\begin{tcolorbox}[
    breakable,
    colback=white,
    colframe=blue!60!black,
    boxrule=0.8pt,
    arc=3pt,
    title={\textbf{Câu \#398} \quad \textbf{MÃ CÂU: \texttt{[L09\_C7\_B2\_10]}} \quad | \quad \textbf{Nhóm: \texttt{HINH\_301}} \quad | \quad \textbf{MapID: \texttt{9X7H2-1}} \hfill \textcolor{yellow!90!white}{[Hình gốc .pyc]}}
]
\textbf{Dạng toán:} Cho b.đồ tròn biểu diễn tstd của các ngôn ngữ lập trình. Tìm tần số của một ngồn ngữ \hfill \textit{(Trắc nghiệm -- Thông hiểu)}
\vspace{4pt}\hrule\vspace{6pt}
\begin{minipage}[t]{0.58\textwidth}
\vspace{0pt}
Biểu đồ hình quạt tròn dưới đây biểu diễn tần số tương đối của các ngôn ngữ lập trình được sử dụng khi viết $200$ phần mềm của một công ty công nghệ.

 Số phần mềm được viết bằng ngôn ngữ  Go là bao nhiêu?

\vspace{4pt}\noindent
\textbf{\textcolor{blue!75!black}{A. (Đ)}} $128$ \quad \textbf{B.} $126$ \quad \textbf{C.} $132$ \quad \textbf{D.} $129$
\end{minipage}\hfill
\begin{minipage}[t]{0.39\textwidth}
\vspace{0pt}
\centering

\resizebox{\linewidth}{!}{%
\begin{tikzpicture}
    \def\R{2.6}
    \draw[double](0,0)circle(\R);
   \def\X{ 64 /north east lines/ Go/1,20/vertical lines/CSS/2,6/dots/Python/3,10/checkerboard/C-Sharp/4}
    \pgfmathsetmacro{\a}{90}
    \foreach \x/\mau/\yeuto/\stt [count=\i] in \X{
        \pgfmathsetmacro{\g}{\x/100*360};
        \pgfmathsetmacro{\b}{\a-\g};
        
        \fill[pattern=\mau,opacity=1,draw=black](0,0)--({\b+\g}:\R)arc({\b+\g}:\b:\R)--cycle;
        \path({\b+\g}:0.7*\R)arc({\b+\g}:\b:0.5*\R)node[pos=0.7,fill=white]{\textbf{\x $\%$}};
        
        \global \let \a = \b;
        \fill[pattern=\mau,draw=black] (3.5, 2.5-\stt) rectangle +(.575,.575);
        \path(6.3,2.75-\stt)node[yscale=1.25,xscale=1.25,text width=1.5*\R cm,align=left,scale=0.75]{\yeuto};
    };
    \path(3.5,3.5)node[yscale=1.5,xscale=1.5,text width=4*\R cm,align=center,scale=0.65]{\textbf{Tần số tương đối của các ngôn ngữ lập trình}};
\end{tikzpicture}
}
\end{minipage}

\end{tcolorbox}
\vspace{4pt}


\begin{tcolorbox}[
    breakable,
    colback=white,
    colframe=blue!60!black,
    boxrule=0.8pt,
    arc=3pt,
    title={\textbf{Câu \#399} \quad \textbf{MÃ CÂU: \texttt{[L09\_C7\_B2\_11]}} \quad | \quad \textbf{Nhóm: \texttt{HINH\_302}} \quad | \quad \textbf{MapID: \texttt{9X7H2-1}} \hfill \textcolor{yellow!90!white}{[Hình gốc .pyc]}}
]
\textbf{Dạng toán:} Cho b.đồ tròn b.diễn tstd của tỉ lệ bình chọn cho dh cầu thủ xs. Tìm số bình chọn của 1 cầu thủ \hfill \textit{(Trắc nghiệm -- Thông hiểu)}
\vspace{4pt}\hrule\vspace{6pt}
\begin{minipage}[t]{0.58\textwidth}
\vspace{0pt}
Biểu đồ hình quạt tròn dưới đây biểu diễn tần số tương đối của tỉ lệ học sinh bình chọn cho danh hiệu cầu thủ xuất sắc nhất trong giải bóng đá của trường. Biết rằng có $1000$ học sinh tham gia bình chọn.

 Số phiếu bầu cho Thái là bao nhiêu?

\vspace{4pt}\noindent
\textbf{A.} $631$ \quad \textbf{\textcolor{blue!75!black}{B. (Đ)}} $630$ \quad \textbf{C.} $632$ \quad \textbf{D.} $629$
\end{minipage}\hfill
\begin{minipage}[t]{0.39\textwidth}
\vspace{0pt}
\centering

\resizebox{\linewidth}{!}{%
\begin{tikzpicture}
    \def\R{2.6}
    \draw[double](0,0)circle(\R);
    \def\X{63/north east lines/Thái/1,17/vertical lines/Hoà/2,6/dots/Trường/3,14/checkerboard/Công/4}
    \pgfmathsetmacro{\a}{90}
    \foreach \x/\mau/\yeuto/\stt [count=\i] in \X{
        \pgfmathsetmacro{\g}{\x/100*360};
        \pgfmathsetmacro{\b}{\a-\g};
        
        \fill[pattern=\mau,opacity=1,draw=black](0,0)--({\b+\g}:\R)arc({\b+\g}:\b:\R)--cycle;
        \path({\b+\g}:0.7*\R)arc({\b+\g}:\b:0.5*\R)node[pos=0.7,fill=white]{\textbf{\x $\%$}};
        
        \global \let \a = \b;
        \fill[pattern=\mau,draw=black] (3.5, 2.5-\stt) rectangle +(.575,.575);
        \path(6.3,2.75-\stt)node[yscale=1.25,xscale=1.25,text width=1.5*\R cm,align=left,scale=0.75]{\yeuto};
    };
    \path(3.5,3.5)node[yscale=1.5,xscale=1.5,text width=4*\R cm,align=center,scale=0.65]{\textbf{Tỉ lệ học sinh bình chọn}};
\end{tikzpicture}
}
\end{minipage}

\end{tcolorbox}
\vspace{4pt}


\begin{tcolorbox}[
    breakable,
    colback=white,
    colframe=blue!60!black,
    boxrule=0.8pt,
    arc=3pt,
    title={\textbf{Câu \#400} \quad \textbf{MÃ CÂU: \texttt{[L09\_C7\_B2\_15]}} \quad | \quad \textbf{Nhóm: \texttt{HINH\_303}} \quad | \quad \textbf{MapID: \texttt{9X7V2-1}} \hfill \textcolor{yellow!90!white}{[Hình gốc .pyc]}}
]
\textbf{Dạng toán:} TF-M3. Cho b.đồ tròn b.diễn tstd của tỉ lệ bình chọn cho dh cầu thủ xs. Tìm số phiếu bình chọn của các cầu thủ \hfill \textit{(Đúng-Sai -- Vận dụng)}
\vspace{4pt}\hrule\vspace{6pt}
\begin{minipage}[t]{0.58\textwidth}
\vspace{0pt}
Biểu đồ hình quạt tròn dưới đây biểu diễn tần số tương đối của tỉ lệ học sinh bình chọn cho danh hiệu cầu thủ xuất sắc nhất trong giải bóng đá của trường. Biết rằng có $1000$ học sinh tham gia bình chọn.

\vspace{4pt}\noindent
\textbf{\textcolor{blue!75!black}{a) (Đ)}} Số phiếu bình chọn cho Tân là $660$\\[2pt]
\textbf{\textcolor{blue!75!black}{b) (Đ)}} Số phiếu bình chọn cho Lam là $130$\\[2pt]
\textbf{c)} Số phiếu bình chọn cho Khôi là $95$\\[2pt]
\textbf{d)} Số phiếu bình chọn cho Hoà là $120$
\end{minipage}\hfill
\begin{minipage}[t]{0.39\textwidth}
\vspace{0pt}
\centering

\resizebox{\linewidth}{!}{%
\begin{tikzpicture}
    \def\R{2.6}
    \draw[double](0,0)circle(\R);
    \def\X{66/north east lines/Tân/1,13/vertical lines/Lam/2,10/dots/Khôi/3,11/checkerboard/Hoà/4}
    \pgfmathsetmacro{\a}{90}
    \foreach \x/\mau/\yeuto/\stt [count=\i] in \X{
        \pgfmathsetmacro{\g}{\x/100*360};
        \pgfmathsetmacro{\b}{\a-\g};
        
        \fill[pattern=\mau,opacity=1,draw=black](0,0)--({\b+\g}:\R)arc({\b+\g}:\b:\R)--cycle;
        \path({\b+\g}:0.7*\R)arc({\b+\g}:\b:0.5*\R)node[pos=0.7,fill=white]{\textbf{\x $\%$}};
        
        \global \let \a = \b;
        \fill[pattern=\mau,draw=black] (3.5, 2.5-\stt) rectangle +(.575,.575);
        \path(6.3,2.75-\stt)node[yscale=1.25,xscale=1.25,text width=1.5*\R cm,align=left,scale=0.75]{\yeuto};
    };
    \path(3.5,3.5)node[yscale=1.5,xscale=1.5,text width=4*\R cm,align=center,scale=0.65]{\textbf{Tỉ lệ học sinh bình chọn}};
\end{tikzpicture}
}
\end{minipage}

\end{tcolorbox}
\vspace{4pt}


\begin{tcolorbox}[
    breakable,
    colback=white,
    colframe=blue!60!black,
    boxrule=0.8pt,
    arc=3pt,
    title={\textbf{Câu \#401} \quad \textbf{MÃ CÂU: \texttt{[L09\_C7\_B2\_16]}} \quad | \quad \textbf{Nhóm: \texttt{HINH\_304}} \quad | \quad \textbf{MapID: \texttt{9X7V2-1}} \hfill \textcolor{yellow!90!white}{[Hình gốc .pyc]}}
]
\textbf{Dạng toán:} TF-M3. Cho biểu đồ tròn biểu diễn tstd của các ngôn ngữ lập trình. Tính số phần mềm viết bằng các ngồn ngữ \hfill \textit{(Đúng-Sai -- Vận dụng)}
\vspace{4pt}\hrule\vspace{6pt}
\begin{minipage}[t]{0.58\textwidth}
\vspace{0pt}
Biểu đồ hình quạt tròn dưới đây biểu diễn tần số tương đối của các ngôn ngữ lập trình được sử dụng khi viết $200$ phần mềm của một công ty công nghệ.

\vspace{4pt}\noindent
\textbf{a)} Số phần mềm viết bằng ngôn ngữ Python là $74$\\[2pt]
\textbf{b)} Số phần mềm viết bằng ngôn ngữ TypeScript là $51$\\[2pt]
\textbf{c)} Số phần mềm viết bằng ngôn ngữ Ruby là $45$\\[2pt]
\textbf{d)} Số phần mềm viết bằng ngôn ngữ Swift là $40$
\end{minipage}\hfill
\begin{minipage}[t]{0.39\textwidth}
\vspace{0pt}
\centering

\resizebox{\linewidth}{!}{%
\begin{tikzpicture}
    \def\R{2.6}
    \draw[double](0,0)circle(\R);
   \def\X{ 32 /north east lines/Python/1,28/vertical lines/TypeScript/2,25/dots/Ruby/3,15/checkerboard/Swift/4}
    \pgfmathsetmacro{\a}{90}
    \foreach \x/\mau/\yeuto/\stt [count=\i] in \X{
        \pgfmathsetmacro{\g}{\x/100*360};
        \pgfmathsetmacro{\b}{\a-\g};
        
        \fill[pattern=\mau,opacity=1,draw=black](0,0)--({\b+\g}:\R)arc({\b+\g}:\b:\R)--cycle;
        \path({\b+\g}:0.7*\R)arc({\b+\g}:\b:0.5*\R)node[pos=0.7,fill=white]{\textbf{\x $\%$}};
        
        \global \let \a = \b;
        \fill[pattern=\mau,draw=black] (3.5, 2.5-\stt) rectangle +(.575,.575);
        \path(6.3,2.75-\stt)node[yscale=1.25,xscale=1.25,text width=1.5*\R cm,align=left,scale=0.75]{\yeuto};
    };
    \path(3.5,3.5)node[yscale=1.5,xscale=1.5,text width=4*\R cm,align=center,scale=0.65]{\textbf{Tần số tương đối của các ngôn ngữ lập trình}};
\end{tikzpicture}
}
\end{minipage}

\end{tcolorbox}
\vspace{4pt}


\begin{tcolorbox}[
    breakable,
    colback=white,
    colframe=blue!60!black,
    boxrule=0.8pt,
    arc=3pt,
    title={\textbf{Câu \#402} \quad \textbf{MÃ CÂU: \texttt{[L09\_C7\_B2\_01]}} \quad | \quad \textbf{Nhóm: \texttt{HINH\_305}} \quad | \quad \textbf{MapID: \texttt{9X7V2-1}} \hfill \textcolor{yellow!90!white}{[Hình gốc .pyc]}}
]
\textbf{Dạng toán:} Cho b.đồ tròn b.diễn tstd của các ngôn ngữ lập trình. Tính tần số. \hfill \textit{(Tự luận -- Vận dụng)}
\vspace{4pt}\hrule\vspace{6pt}
\begin{minipage}[t]{0.58\textwidth}
\vspace{0pt}
Biểu đồ hình quạt tròn dưới đây biểu diễn tần số tương đối của các ngôn ngữ lập trình được sử dụng khi viết $400$ phần mềm của một công ty công nghệ. Biết rằng, mỗi phần mềm được viết bằng đúng một ngôn ngữ lập trình.

a) Ngôn ngữ lập trình nào được sử dụng phổ biến nhất khi viết $400$ phần mềm đó?

b) Hãy lập bảng tần số biểu diễn số liệu đã cho bởi biểu đồ trên.
\end{minipage}\hfill
\begin{minipage}[t]{0.39\textwidth}
\vspace{0pt}
\centering

\resizebox{\linewidth}{!}{%
\begin{tikzpicture}
    \def\R{2.6}
    \draw[double](0,0)circle(\R);
   \def\X{47/north east lines/C-Sharp/1,28/vertical lines/Python/2,10/dots/Ruby/3,15/checkerboard/C++/4}
    \pgfmathsetmacro{\a}{90}
    \foreach \x/\mau/\yeuto/\stt [count=\i] in \X{
        \pgfmathsetmacro{\g}{\x/100*360};
        \pgfmathsetmacro{\b}{\a-\g};
        
        \fill[pattern=\mau,opacity=1,draw=black](0,0)--({\b+\g}:\R)arc({\b+\g}:\b:\R)--cycle;
        \path({\b+\g}:0.7*\R)arc({\b+\g}:\b:0.5*\R)node[pos=0.7,fill=white]{\textbf{\x $\%$}};
        
        \global \let \a = \b;
        \fill[pattern=\mau,draw=black] (3.5, 2.5-\stt) rectangle +(.575,.575);
        \path(6.3,2.75-\stt)node[yscale=1.25,xscale=1.25,text width=1.5*\R cm,align=left,scale=0.75]{\yeuto};
    };
    \path(3.5,3.5)node[yscale=1.5,xscale=1.5,text width=4*\R cm,align=center,scale=0.65]{\textbf{Tần số tương đối của các ngôn ngữ lập trình}};
\end{tikzpicture}
}
\end{minipage}

\end{tcolorbox}
\vspace{4pt}


\begin{tcolorbox}[
    breakable,
    colback=white,
    colframe=blue!60!black,
    boxrule=0.8pt,
    arc=3pt,
    title={\textbf{Câu \#403} \quad \textbf{MÃ CÂU: \texttt{[L09\_C7\_B2\_02]}} \quad | \quad \textbf{Nhóm: \texttt{HINH\_306}} \quad | \quad \textbf{MapID: \texttt{9X7H2-1}} \hfill \textcolor{yellow!90!white}{[Hình gốc .pyc]}}
]
\textbf{Dạng toán:} Cho b.đồ tròn b.diễn tstd của tỉ lệ b.chọn cho cầu thủ xs. Tính tần số \hfill \textit{(Tự luận -- Thông hiểu)}
\vspace{4pt}\hrule\vspace{6pt}
\begin{minipage}[t]{0.58\textwidth}
\vspace{0pt}
Biểu đồ hình quạt tròn dưới đây biểu diễn tần số tương đối của tỉ lệ học sinh bình chọn cho danh hiệu cầu thủ xuất sắc nhất trong giải bóng đá của trường. Biết rằng có $2000$ học sinh tham gia bình chọn.

a) Cầu thủ nào được bình chọn nhiều nhất?

b) Hãy lập bảng tần số biểu diễn số liệu đã cho bởi biểu đồ trên.
\end{minipage}\hfill
\begin{minipage}[t]{0.39\textwidth}
\vspace{0pt}
\centering

\resizebox{\linewidth}{!}{%
\begin{tikzpicture}
    \def\R{2.6}
    \draw[double](0,0)circle(\R);
    \def\X{49/north east lines/Tân/1,19/vertical lines/Thành/2,16/dots/Khôi/3,16/checkerboard/Hoàng/4}
    \pgfmathsetmacro{\a}{90}
    \foreach \x/\mau/\yeuto/\stt [count=\i] in \X{
        \pgfmathsetmacro{\g}{\x/100*360};
        \pgfmathsetmacro{\b}{\a-\g};
        
        \fill[pattern=\mau,opacity=1,draw=black](0,0)--({\b+\g}:\R)arc({\b+\g}:\b:\R)--cycle;
        \path({\b+\g}:0.7*\R)arc({\b+\g}:\b:0.5*\R)node[pos=0.7,fill=white]{\textbf{\x $\%$}};
        
        \global \let \a = \b;
        \fill[pattern=\mau,draw=black] (3.5, 2.5-\stt) rectangle +(.575,.575);
        \path(6.3,2.75-\stt)node[yscale=1.25,xscale=1.25,text width=1.5*\R cm,align=left,scale=0.75]{\yeuto};
    };
    \path(3.5,3.5)node[yscale=1.5,xscale=1.5,text width=4*\R cm,align=center,scale=0.65]{\textbf{Tỉ lệ học sinh bình chọn}};
\end{tikzpicture}
}
\end{minipage}

\end{tcolorbox}
\vspace{4pt}


\subsection{Chương 8 - Xác suất}

\begin{tcolorbox}[
    breakable,
    colback=white,
    colframe=blue!60!black,
    boxrule=0.8pt,
    arc=3pt,
    title={\textbf{Câu \#404} \quad \textbf{MÃ CÂU: \texttt{[L09\_C8\_B2\_17]}} \quad | \quad \textbf{Nhóm: \texttt{HINH\_307}} \quad | \quad \textbf{MapID: \texttt{9X8V2-1}} \hfill \textcolor{yellow!90!white}{[Hình gốc .pyc]}}
]
\textbf{Dạng toán:} XS các btoán quay số từ 2 vòng quay và nhân lại với nhau(8 kiểu) \hfill \textit{(Trắc nghiệm -- Vận dụng)}
\vspace{4pt}\hrule\vspace{6pt}
\begin{minipage}[t]{0.58\textwidth}
\vspace{0pt}
Tấm bìa cứng ${A}$ hình tròn được chia thành ba hình quạt có diện tích bằng nhau, đánh số ${4; 1; 5}$ và tấm bìa cứng ${B}$ hình tròn được chia thành $4$ hình quạt có diện tích bằng nhau đánh số ${5; 1; 2; 6}$. Trục quay của ${A}$ và ${B}$ được gắn mũi tên ở tâm. Bạn Xuân quay tấm bìa ${A}$, bạn Mai quay tấm bìa ${B}$. Xác xuất để tích hai số ở hình quạt mà hai mũi tên chỉ vào là số lẻ là:

\vspace{4pt}\noindent
\textbf{A.} ${\dfrac{4}{9}}$ \quad \textbf{B.} ${\dfrac{4}{13}}$ \quad \textbf{\textcolor{blue!75!black}{C. (Đ)}} ${\dfrac{1}{3}}$ \quad \textbf{D.} ${\dfrac{4}{15}} $
\end{minipage}\hfill
\begin{minipage}[t]{0.39\textwidth}
\vspace{0pt}
\centering

\resizebox{\linewidth}{!}{%
\begin{tikzpicture}
    % Vẽ hình tròn thứ nhất (nhỏ lại)
    \draw[thick] (0,0) circle(2cm);
    
    % Vẽ các đường chia hình tròn thứ nhất thành 3 phần bằng nhau
    \foreach \angle in {0,120,240} {
        \draw[thick] (0,0) -- (\angle:2cm);
    }
    
    % Đánh số các phần của hình tròn thứ nhất
    \node at (60:1cm) {$4$};
    \node at (180:1cm) {$1$};
    \node at (-60:1cm){$5$};
    
    % Vẽ hình tròn thứ hai (nhỏ lại, có khoảng cách)
    \draw[thick] (5,0) circle(2cm);
    
    % Vẽ các đường chia hình tròn thứ hai thành 4 phần bằng nhau
    \foreach \angle in {0,90,180,270} {
        \draw[thick] (5,0) -- ++(\angle:2cm);
    }
    
    % Đánh số các phần của hình tròn thứ hai
    \node at (6,1) {$5$};
    \node at (6,-1) {$1$};
    \node at (4,-1) {$2$};
    \node at (4,1) {$6$};
\end{tikzpicture}
}
\end{minipage}

\vspace{3pt}\hrule\vspace{3pt}\footnotesize\textcolor{gray!80!black}{\textbf{Các mã dùng chung hình (HINH\_307):} \texttt{L09\_C8\_B2\_18, L09\_C8\_B2\_19}}
\end{tcolorbox}
\vspace{4pt}


\begin{tcolorbox}[
    breakable,
    colback=white,
    colframe=blue!60!black,
    boxrule=0.8pt,
    arc=3pt,
    title={\textbf{Câu \#405} \quad \textbf{MÃ CÂU: \texttt{[L09\_C8\_B2\_18]}} \quad | \quad \textbf{Nhóm: \texttt{HINH\_307}} \quad | \quad \textbf{MapID: \texttt{9X8V2-1}} \hfill \textcolor{yellow!90!white}{[Hình gốc .pyc]}}
]
\textbf{Dạng toán:} XS các btoán quay số từ 2 vòng quay và cộng lại với nhau(7 kiểu) \hfill \textit{(Trắc nghiệm -- Vận dụng)}
\vspace{4pt}\hrule\vspace{6pt}
\begin{minipage}[t]{0.58\textwidth}
\vspace{0pt}
Tấm bìa cứng ${A}$ hình tròn được chia thành ba hình quạt có diện tích bằng nhau, đánh số ${2; 8; 6}$ và tấm bìa cứng ${B}$ hình tròn được chia thành $4$ hình quạt có diện tích bằng nhau đánh số ${6; 1; 2; 8}$. Trục quay của ${A}$ và ${B}$ được gắn mũi tên ở tâm. Bạn h quay tấm bìa ${A}$, bạn u quay tấm bìa ${B}$. Xác xuất để tổng hai số ở hình quạt mà hai mũi tên chỉ vào là số lớn hơn $5$ và nhỏ hơn $15$ là:

\vspace{4pt}\noindent
\textbf{A.} ${\dfrac{9}{17}} $ \quad \textbf{B.} ${\dfrac{9}{8}}$ \quad \textbf{C.} ${\dfrac{9}{11}}$ \quad \textbf{\textcolor{blue!75!black}{D. (Đ)}} ${\dfrac{3}{4}}$
\end{minipage}\hfill
\begin{minipage}[t]{0.39\textwidth}
\vspace{0pt}
\centering

\resizebox{\linewidth}{!}{%
\begin{tikzpicture}
    % Vẽ hình tròn thứ nhất (nhỏ lại)
    \draw[thick] (0,0) circle(2cm);
    
    % Vẽ các đường chia hình tròn thứ nhất thành 3 phần bằng nhau
    \foreach \angle in {0,120,240} {
        \draw[thick] (0,0) -- (\angle:2cm);
    }
    
    % Đánh số các phần của hình tròn thứ nhất
    \node at (60:1cm) {$2$};
    \node at (180:1cm) {$8$};
    \node at (-60:1cm){$6$};
    
    % Vẽ hình tròn thứ hai (nhỏ lại, có khoảng cách)
    \draw[thick] (5,0) circle(2cm);
    
    % Vẽ các đường chia hình tròn thứ hai thành 4 phần bằng nhau
    \foreach \angle in {0,90,180,270} {
        \draw[thick] (5,0) -- ++(\angle:2cm);
    }
    
    % Đánh số các phần của hình tròn thứ hai
    \node at (6,1) {$6$};
    \node at (6,-1) {$1$};
    \node at (4,-1) {$2$};
    \node at (4,1) {$8$};
\end{tikzpicture}
}
\end{minipage}

\vspace{3pt}\hrule\vspace{3pt}\footnotesize\textcolor{gray!80!black}{\textbf{Các mã dùng chung hình (HINH\_307):} \texttt{L09\_C8\_B2\_17, L09\_C8\_B2\_19}}
\end{tcolorbox}
\vspace{4pt}


\begin{tcolorbox}[
    breakable,
    colback=white,
    colframe=blue!60!black,
    boxrule=0.8pt,
    arc=3pt,
    title={\textbf{Câu \#406} \quad \textbf{MÃ CÂU: \texttt{[L09\_C8\_B2\_19]}} \quad | \quad \textbf{Nhóm: \texttt{HINH\_307}} \quad | \quad \textbf{MapID: \texttt{9X8V2-1}} \hfill \textcolor{yellow!90!white}{[Hình gốc .pyc]}}
]
\textbf{Dạng toán:} XS các btoán quay số từ 2 vòng quay và ghép lại với nhau(7 kiểu) \hfill \textit{(Trắc nghiệm -- Vận dụng)}
\vspace{4pt}\hrule\vspace{6pt}
\begin{minipage}[t]{0.58\textwidth}
\vspace{0pt}
Tấm bìa cứng ${A}$ hình tròn được chia thành ba hình quạt có diện tích bằng nhau, đánh số ${4; 7; 3}$ và tấm bìa cứng ${B}$ hình tròn được chia thành $4$ hình quạt có diện tích bằng nhau đánh số ${4; 2; 5; 6}$. Trục quay của ${A}$ và ${B}$ được gắn mũi tên ở tâm. Bạn Hà quay tấm bìa ${A}$, bạn Duy quay tấm bìa ${B}$ và ghép hai số quay được của 2 bạn thành số có hai chữ số với chữ số bạn Hà quay được là chữ số hàng chục. Xác suất để kết quả là số chính phương là:

\vspace{4pt}\noindent
\textbf{\textcolor{blue!75!black}{A. (Đ)}} ${\dfrac{1}{12}}$ \quad \textbf{B.} ${\dfrac{1}{15}} $ \quad \textbf{C.} ${\dfrac{1}{16}}$ \quad \textbf{D.} ${\dfrac{1}{17}} $
\end{minipage}\hfill
\begin{minipage}[t]{0.39\textwidth}
\vspace{0pt}
\centering

\resizebox{\linewidth}{!}{%
\begin{tikzpicture}
    % Vẽ hình tròn thứ nhất (nhỏ lại)
    \draw[thick] (0,0) circle(2cm);
    
    % Vẽ các đường chia hình tròn thứ nhất thành 3 phần bằng nhau
    \foreach \angle in {0,120,240} {
        \draw[thick] (0,0) -- (\angle:2cm);
    }
    
    % Đánh số các phần của hình tròn thứ nhất
    \node at (60:1cm) {$4$};
    \node at (180:1cm) {$7$};
    \node at (-60:1cm){$3$};
    
    % Vẽ hình tròn thứ hai (nhỏ lại, có khoảng cách)
    \draw[thick] (5,0) circle(2cm);
    
    % Vẽ các đường chia hình tròn thứ hai thành 4 phần bằng nhau
    \foreach \angle in {0,90,180,270} {
        \draw[thick] (5,0) -- ++(\angle:2cm);
    }
    
    % Đánh số các phần của hình tròn thứ hai
    \node at (6,1) {$4$};
    \node at (6,-1) {$2$};
    \node at (4,-1) {$5$};
    \node at (4,1) {$6$};
\end{tikzpicture}
}
\end{minipage}

\vspace{3pt}\hrule\vspace{3pt}\footnotesize\textcolor{gray!80!black}{\textbf{Các mã dùng chung hình (HINH\_307):} \texttt{L09\_C8\_B2\_17, L09\_C8\_B2\_18}}
\end{tcolorbox}
\vspace{4pt}


\subsection{Chương 9 - Đường tròn nội tiếp-ngoại tiếp}

\begin{tcolorbox}[
    breakable,
    colback=white,
    colframe=blue!60!black,
    boxrule=0.8pt,
    arc=3pt,
    title={\textbf{Câu \#407} \quad \textbf{MÃ CÂU: \texttt{[L09\_C9\_B1\_09]}} \quad | \quad \textbf{Nhóm: \texttt{HINH\_308}} \quad | \quad \textbf{MapID: \texttt{9H9H1-1}} \hfill \textcolor{yellow!90!white}{[Hình gốc .pyc]}}
]
\textbf{Dạng toán:} Tính góc \hfill \textit{(Trắc nghiệm -- Thông hiểu)}
\vspace{4pt}\hrule\vspace{6pt}
\begin{minipage}[t]{0.58\textwidth}
\vspace{0pt}
Trên đường tròn tâm ${I}$ đường kính ${CF}$ lấy điểm ${G}$ sao cho $\widehat{CIG}=78^{\circ}$. Hãy tính các góc của $\triangle {CFG}$.

\vspace{4pt}\noindent
\textbf{\textcolor{blue!75!black}{A. (Đ)}} $\widehat{CGF}=90^{\circ}$;  $\widehat{CFG}=39^{\circ}$; $ \widehat{GCF}=51^{\circ}$ \quad \textbf{B.} $\widehat{CGF}=39^{\circ}$;  $\widehat{CFG}=51^{\circ}$; $ \widehat{GCF}=51^{\circ}$ \quad \textbf{C.} $\widehat{CGF}=51^{\circ}$;  $\widehat{CFG}=39^{\circ}$; $ \widehat{GCF}=51^{\circ}$ \quad \textbf{D.} $\widehat{CGF}=51^{\circ}$;  $\widehat{CFG}=39^{\circ}$; $ \widehat{GCF}=90^{\circ}$
\end{minipage}\hfill
\begin{minipage}[t]{0.39\textwidth}
\vspace{0pt}
\centering

\resizebox{\linewidth}{!}{%
\begin{tikzpicture}[scale=1, font=\footnotesize, line join=round, line cap=round, >=stealth]
                \def\r{2}   
                \path (0,0) coordinate (O)
                ($(O) + (180:\r)$) coordinate (A)
                ($(O) + (0:\r)$) coordinate (B)
                ($(O) + (110:\r)$) coordinate (C);
                \pgfresetboundingbox
                \draw (O) circle (\r)
                (A)--(B)--(C)--(A)--(O)--(B) (O)--(C);
                \draw ($(O) + (150:0.5)$) node{$78^\circ$};
                   \filldraw[black] (A) circle (1.5pt) node[above left ] {$C$};
    \filldraw[black] (B) circle (1.5pt) node[above right] {$F$};
    \filldraw[black] (C) circle (1.5pt) node[above ] {$G$};
   \filldraw[black] (O) circle (1.5pt) node[below]{$I$};
            \end{tikzpicture}
}
\end{minipage}

\vspace{3pt}\hrule\vspace{3pt}\footnotesize\textcolor{gray!80!black}{\textbf{Các mã dùng chung hình (HINH\_308):} \texttt{L09\_C9\_B1\_13}}
\end{tcolorbox}
\vspace{4pt}


\begin{tcolorbox}[
    breakable,
    colback=white,
    colframe=blue!60!black,
    boxrule=0.8pt,
    arc=3pt,
    title={\textbf{Câu \#408} \quad \textbf{MÃ CÂU: \texttt{[L09\_C9\_B1\_10]}} \quad | \quad \textbf{Nhóm: \texttt{HINH\_309}} \quad | \quad \textbf{MapID: \texttt{9H9V1-1}} \hfill \textcolor{yellow!90!white}{[Hình gốc .pyc]}}
]
\textbf{Dạng toán:} Tính số đo các góc \hfill \textit{(Trắc nghiệm -- Vận dụng)}
\vspace{4pt}\hrule\vspace{6pt}
\begin{minipage}[t]{0.58\textwidth}
\vspace{0pt}
Cho ${AE}$ và ${FG}$ là hai đường kính của đường tròn $({O})$. Biết rằng $\widehat{AOF} = 76^{\circ}$, tính số đo của $\widehat{AEF}$, $\widehat{AGF}$ và $\widehat{AEG}$.

\vspace{4pt}\noindent
\textbf{A.} $\widehat{AEF}=104^{\circ}$, $\widehat{AGF}=104^{\circ}$ và $\widehat{AEG}= 104^{\circ}$ \quad \textbf{\textcolor{blue!75!black}{B. (Đ)}} $\widehat{AEF}=38^{\circ}$, $\widehat{AGF}=38^{\circ}$ và $\widehat{AEG}= 52^{\circ}$ \quad \textbf{C.} $\widehat{AEF}=76^{\circ}$, $\widehat{AGF}=38^{\circ}$ và $\widehat{AEG}= 104^{\circ}$ \quad \textbf{D.} $\widehat{AEF}=76^{\circ}$, $\widehat{AGF}=76^{\circ}$ và $\widehat{AEG}= 52^{\circ}$
\end{minipage}\hfill
\begin{minipage}[t]{0.39\textwidth}
\vspace{0pt}
\centering

\resizebox{\linewidth}{!}{%
\begin{tikzpicture}[>=stealth,line join=round,line cap=round,font=\footnotesize,scale=1]
                \path 
                (0,0) coordinate (O)
                ($(O) + (-15:2)$) coordinate (A)
                ($(A)!2!(O)$) coordinate (B)
                ($(O) + (60:2)$) coordinate (C)
                ($(C)!2!(O)$) coordinate (D)
                ;
                \draw (O) circle (2cm);
                \draw (B)--(A)--(D)--(B)--(C)--(D);
                 \filldraw[black] (A) circle (1.5pt) node[above right] {$A$};
    \filldraw[black] (B) circle (1.5pt) node[above left] {$E$};
    \filldraw[black] (C) circle (1.5pt) node[above] {$F$};
   \filldraw[black] (D) circle (1.5pt) node[below left]{$G$};
    \filldraw[black] (O) circle (1.5pt) node[below] {$O$};
            \end{tikzpicture}
}
\end{minipage}

\vspace{3pt}\hrule\vspace{3pt}\footnotesize\textcolor{gray!80!black}{\textbf{Các mã dùng chung hình (HINH\_309):} \texttt{L09\_C9\_B1\_14}}
\end{tcolorbox}
\vspace{4pt}


\begin{tcolorbox}[
    breakable,
    colback=white,
    colframe=blue!60!black,
    boxrule=0.8pt,
    arc=3pt,
    title={\textbf{Câu \#409} \quad \textbf{MÃ CÂU: \texttt{[L09\_C9\_B1\_11]}} \quad | \quad \textbf{Nhóm: \texttt{HINH\_310}} \quad | \quad \textbf{MapID: \texttt{9H9V1-1}} \hfill \textcolor{yellow!90!white}{[Hình gốc .pyc]}}
]
\textbf{Dạng toán:} Tính số đo các góc 2 \hfill \textit{(Trắc nghiệm -- Vận dụng)}
\vspace{4pt}\hrule\vspace{6pt}
\begin{minipage}[t]{0.58\textwidth}
\vspace{0pt}
Cho đường tròn $({O})$ và hai dây cung ${AB}$, ${CF}$ cắt nhau tại ${H}$. Tính số đo các góc của tam giác ${AHC}$, biết rằng $\widehat{HBF} = 50^{\circ}$, $\widehat{HFB} = 57^{\circ}$.

\vspace{4pt}\noindent
\textbf{A.} $\widehat{ACH} =57^{\circ}$, $\widehat{CAH}=114^{\circ}$ và $\widehat{AHC}= 73^{\circ}$ \quad \textbf{\textcolor{blue!75!black}{B. (Đ)}} $\widehat{ACH}=50^{\circ}$, $\widehat{CAH}=57^{\circ}$ và $\widehat{AHC}= 73^{\circ}$ \quad \textbf{C.} $\widehat{ACH}=80^{\circ}$, $\widehat{CAH}=66^{\circ}$ và $\widehat{AHC}= 73^{\circ}$ \quad \textbf{D.} $\widehat{ACH} =50^{\circ}$, $\widehat{CAH}=114^{\circ}$ và $\widehat{AHC}= 73^{\circ}$
\end{minipage}\hfill
\begin{minipage}[t]{0.39\textwidth}
\vspace{0pt}
\centering

\resizebox{\linewidth}{!}{%
\begin{tikzpicture}[>=stealth,line join=round,line cap=round,font=\footnotesize,scale=1]
                \path 
                (0,0) coordinate (O)
                ($(O) + (32:2)$) coordinate (A)
                ($(O) + (205:2)$) coordinate (B)
                ($(O) + (-10:2)$) coordinate (C)
                ($(O) + (145:2)$) coordinate (D)
                (intersection of A--B and C--D) coordinate (X)
                ;
                \draw (O) circle (2cm);
                \draw (B)--(A)--(C)--(D)--(B);
                    \filldraw[black] (A) circle (1pt) node[above right] {$A$};
    \filldraw[black] (B) circle (1pt) node[above left] {$B$};
    \filldraw[black] (C) circle (1pt) node[above right] {$C$};
   \filldraw[black] (D) circle (1pt) node[above left]{$F$};
    \filldraw[black] (O) circle (1pt) node[below left] {$O$};
     \filldraw[black] (X) circle (1pt) node[above] {$H$};
                \fill[black] (D) +(-50:5mm) node{$57^\circ$};
                \fill[black] (B) +(50:5.5mm) node{$50^\circ$};
                
            \end{tikzpicture}
}
\end{minipage}

\vspace{3pt}\hrule\vspace{3pt}\footnotesize\textcolor{gray!80!black}{\textbf{Các mã dùng chung hình (HINH\_310):} \texttt{L09\_C9\_B1\_15, L09\_C9\_B1\_05}}
\end{tcolorbox}
\vspace{4pt}


\begin{tcolorbox}[
    breakable,
    colback=white,
    colframe=blue!60!black,
    boxrule=0.8pt,
    arc=3pt,
    title={\textbf{Câu \#410} \quad \textbf{MÃ CÂU: \texttt{[L09\_C9\_B1\_12]}} \quad | \quad \textbf{Nhóm: \texttt{HINH\_311}} \quad | \quad \textbf{MapID: \texttt{9H9V1-1}} \hfill \textcolor{yellow!90!white}{[Hình gốc .pyc]}}
]
\textbf{Dạng toán:} Tính số đo các góc 3 \hfill \textit{(Trắc nghiệm -- Vận dụng)}
\vspace{4pt}\hrule\vspace{6pt}
\begin{minipage}[t]{0.58\textwidth}
\vspace{0pt}
Cho hai điểm ${J}$, ${M}$ nằm trên đường tròn $({O})$ và cho điểm ${I}$ nằm trên cung lớn $\stackrel\frown{JM}$. Biết rằng $\widehat{OJI} = 26^{\circ}$, $\widehat{OMI} = 33^{\circ}$. Tính số đo các góc của tam giác ${IJM}$.

\vspace{4pt}\noindent
\textbf{\textcolor{blue!75!black}{A. (Đ)}} $\widehat{IJM}=57^{\circ}$, $\widehat{IMJ}=64^{\circ}$, $\widehat{JIM}=59^{\circ}$ \quad \textbf{B.} $\widehat{IJM}=57^{\circ}$, $\widehat{IMJ}=26^{\circ}$, $\widehat{JIM}=59^{\circ}$ \quad \textbf{C.} $\widehat{IJM}=57^{\circ}$, $\widehat{IMJ}=57^{\circ}$, $\widehat{JIM}=59^{\circ}$ \quad \textbf{D.} $\widehat{IJM}=33^{\circ}$, $\widehat{IMJ}=26^{\circ}$, $\widehat{JIM}=59^{\circ}$
\end{minipage}\hfill
\begin{minipage}[t]{0.39\textwidth}
\vspace{0pt}
\centering

\resizebox{\linewidth}{!}{%
\begin{tikzpicture}[>=stealth,line join=round,line cap=round,font=\footnotesize,scale=1]
                \path 
                (0,0) coordinate (O)
                ($(O) + (100:2)$) coordinate (A)
                ($(O) + (-15:2)$) coordinate (B)
                ($(O) + (-165:2)$) coordinate (C)
                ;
                \draw (O) circle (2cm);
                \draw (C)--(B)--(A)--(C)--(A)--(O)--(B) (O)--(C);
                     \node[above left] at (A) {$I$};
      \node[above right] at (B) {$J$};
       \node[above left] at (C) {$M$};
        \node[above left] at (O) {$O$};
                \fill[black] (B) +(150:7.3mm) node[scale=0.8]{$26^\circ$};
                \fill[black] (C) +(30:5.5mm) node[scale=0.8]{$33^\circ$};
            \end{tikzpicture}
}
\end{minipage}

\vspace{3pt}\hrule\vspace{3pt}\footnotesize\textcolor{gray!80!black}{\textbf{Các mã dùng chung hình (HINH\_311):} \texttt{L09\_C9\_B1\_16, L09\_C9\_B1\_06}}
\end{tcolorbox}
\vspace{4pt}


\begin{tcolorbox}[
    breakable,
    colback=white,
    colframe=blue!60!black,
    boxrule=0.8pt,
    arc=3pt,
    title={\textbf{Câu \#411} \quad \textbf{MÃ CÂU: \texttt{[L09\_C9\_B1\_13]}} \quad | \quad \textbf{Nhóm: \texttt{HINH\_308}} \quad | \quad \textbf{MapID: \texttt{9H9H1-1}} \hfill \textcolor{yellow!90!white}{[Hình gốc .pyc]}}
]
\textbf{Dạng toán:} TF-M2. Tính số đo góc1 \hfill \textit{(Đúng-Sai -- Thông hiểu)}
\vspace{4pt}\hrule\vspace{6pt}
\begin{minipage}[t]{0.58\textwidth}
\vspace{0pt}
Trên đường tròn tâm ${O}$ đường kính ${EG}$ lấy điểm ${H}$ sao cho $\widehat{EOH}=80^{\circ}$ . Xét tính đúng-sai của các khẳng định sau.

\vspace{4pt}\noindent
\textbf{\textcolor{blue!75!black}{a) (Đ)}} $\widehat{EHG}=90^{\circ}$\\[2pt]
\textbf{\textcolor{blue!75!black}{b) (Đ)}} $\widehat{EGH}= 40^{\circ}$\\[2pt]
\textbf{\textcolor{blue!75!black}{c) (Đ)}} $ \widehat{HEG}=50^{\circ}$\\[2pt]
\textbf{d)} $ \widehat{EHO}=40^{\circ}$
\end{minipage}\hfill
\begin{minipage}[t]{0.39\textwidth}
\vspace{0pt}
\centering

\resizebox{\linewidth}{!}{%
\begin{tikzpicture}[scale=1, font=\footnotesize, line join=round, line cap=round, >=stealth]
                \def\r{2}   
                \path (0,0) coordinate (O)
                ($(O) + (180:\r)$) coordinate (A)
                ($(O) + (0:\r)$) coordinate (B)
                ($(O) + (110:\r)$) coordinate (C);
                \pgfresetboundingbox
                \draw (O) circle (\r)
                (A)--(B)--(C)--(A)--(O)--(B) (O)--(C);
                \draw ($(O) + (150:0.5)$) node{$80^\circ$};
                   \filldraw[black] (A) circle (1.5pt) node[above left ] {$E$};
    \filldraw[black] (B) circle (1.5pt) node[above right] {$G$};
    \filldraw[black] (C) circle (1.5pt) node[above ] {$H$};
   \filldraw[black] (O) circle (1.5pt) node[below]{$O$};
            \end{tikzpicture}
}
\end{minipage}

\vspace{3pt}\hrule\vspace{3pt}\footnotesize\textcolor{gray!80!black}{\textbf{Các mã dùng chung hình (HINH\_308):} \texttt{L09\_C9\_B1\_09}}
\end{tcolorbox}
\vspace{4pt}


\begin{tcolorbox}[
    breakable,
    colback=white,
    colframe=blue!60!black,
    boxrule=0.8pt,
    arc=3pt,
    title={\textbf{Câu \#412} \quad \textbf{MÃ CÂU: \texttt{[L09\_C9\_B1\_14]}} \quad | \quad \textbf{Nhóm: \texttt{HINH\_309}} \quad | \quad \textbf{MapID: \texttt{9H9V1-1}} \hfill \textcolor{yellow!90!white}{[Hình gốc .pyc]}}
]
\textbf{Dạng toán:} TF-M3. Tính số đo các góc2 \hfill \textit{(Đúng-Sai -- Vận dụng)}
\vspace{4pt}\hrule\vspace{6pt}
\begin{minipage}[t]{0.58\textwidth}
\vspace{0pt}
Cho ${AD}$ và ${GH}$ là hai đường kính của đường tròn $({I})$. Biết rằng $\widehat{AIG} = 76^{\circ}$. Xét tính đúng-sai của các khẳng định sau.

\vspace{4pt}\noindent
\textbf{a)} $\widehat{ADG} = 76^{\circ} $\\[2pt]
\textbf{b)} $ \widehat{ADH}= 142^{\circ} $\\[2pt]
\textbf{\textcolor{blue!75!black}{c) (Đ)}} $\widehat{AIH}=  104^{\circ}$\\[2pt]
\textbf{d)} $\widehat{DGH}=  142^{\circ}$
\end{minipage}\hfill
\begin{minipage}[t]{0.39\textwidth}
\vspace{0pt}
\centering

\resizebox{\linewidth}{!}{%
\begin{tikzpicture}[>=stealth,line join=round,line cap=round,font=\footnotesize,scale=1]
                \path 
                (0,0) coordinate (O)
                ($(O) + (-15:2)$) coordinate (A)
                ($(A)!2!(O)$) coordinate (B)
                ($(O) + (60:2)$) coordinate (C)
                ($(C)!2!(O)$) coordinate (D)
                ;
                \draw (O) circle (2cm);
                \draw (B)--(A)--(D)--(B)--(C)--(D);
                 \filldraw[black] (A) circle (1.5pt) node[above right] {$A$};
    \filldraw[black] (B) circle (1.5pt) node[above left] {$D$};
    \filldraw[black] (C) circle (1.5pt) node[above] {$G$};
   \filldraw[black] (D) circle (1.5pt) node[below left]{$H$};
    \filldraw[black] (O) circle (1.5pt) node[below] {$I$};
            \end{tikzpicture}
}
\end{minipage}

\vspace{3pt}\hrule\vspace{3pt}\footnotesize\textcolor{gray!80!black}{\textbf{Các mã dùng chung hình (HINH\_309):} \texttt{L09\_C9\_B1\_10}}
\end{tcolorbox}
\vspace{4pt}


\begin{tcolorbox}[
    breakable,
    colback=white,
    colframe=blue!60!black,
    boxrule=0.8pt,
    arc=3pt,
    title={\textbf{Câu \#413} \quad \textbf{MÃ CÂU: \texttt{[L09\_C9\_B1\_15]}} \quad | \quad \textbf{Nhóm: \texttt{HINH\_310}} \quad | \quad \textbf{MapID: \texttt{9H9V1-1}} \hfill \textcolor{yellow!90!white}{[Hình gốc .pyc]}}
]
\textbf{Dạng toán:} TF-M3. Tính số đo các góc 3 \hfill \textit{(Đúng-Sai -- Vận dụng)}
\vspace{4pt}\hrule\vspace{6pt}
\begin{minipage}[t]{0.58\textwidth}
\vspace{0pt}
Cho đường tròn $({I})$ và hai dây cung ${BC}$, ${DE}$ cắt nhau tại ${H}$, biết rằng $\widehat{HCE} = 50^{\circ}$, $\widehat{HEC} = 60^{\circ}$. Xét tính đúng-sai của các khẳng định sau.

\vspace{4pt}\noindent
\textbf{\textcolor{blue!75!black}{a) (Đ)}} $\widehat{BHD} = 70^{\circ} $\\[2pt]
\textbf{\textcolor{blue!75!black}{b) (Đ)}} $\widehat{BDH} =  50^{\circ}$\\[2pt]
\textbf{\textcolor{blue!75!black}{c) (Đ)}} $\widehat{DBH} = 60^{\circ}$\\[2pt]
\textbf{\textcolor{blue!75!black}{d) (Đ)}} $\widehat{CHD}= 110^{\circ} $
\end{minipage}\hfill
\begin{minipage}[t]{0.39\textwidth}
\vspace{0pt}
\centering

\resizebox{\linewidth}{!}{%
\begin{tikzpicture}[>=stealth,line join=round,line cap=round,font=\footnotesize,scale=1]
                \path 
                (0,0) coordinate (O)
                ($(O) + (32:2)$) coordinate (A)
                ($(O) + (205:2)$) coordinate (B)
                ($(O) + (-10:2)$) coordinate (C)
                ($(O) + (145:2)$) coordinate (D)
                (intersection of A--B and C--D) coordinate (X)
                ;
                \draw (O) circle (2cm);
                \draw (B)--(A)--(C)--(D)--(B);
                    \filldraw[black] (A) circle (1pt) node[above right] {$B$};
    \filldraw[black] (B) circle (1pt) node[above left] {$C$};
    \filldraw[black] (C) circle (1pt) node[above right] {$D$};
   \filldraw[black] (D) circle (1pt) node[above left]{$E$};
    \filldraw[black] (O) circle (1pt) node[below left] {$I$};
     \filldraw[black] (X) circle (1pt) node[above] {$H$};
                \fill[black] (D) +(-50:5mm) node{$60^\circ$};
                \fill[black] (B) +(50:5.5mm) node{$50^\circ$};
                
            \end{tikzpicture}
}
\end{minipage}

\vspace{3pt}\hrule\vspace{3pt}\footnotesize\textcolor{gray!80!black}{\textbf{Các mã dùng chung hình (HINH\_310):} \texttt{L09\_C9\_B1\_11, L09\_C9\_B1\_05}}
\end{tcolorbox}
\vspace{4pt}


\begin{tcolorbox}[
    breakable,
    colback=white,
    colframe=blue!60!black,
    boxrule=0.8pt,
    arc=3pt,
    title={\textbf{Câu \#414} \quad \textbf{MÃ CÂU: \texttt{[L09\_C9\_B1\_16]}} \quad | \quad \textbf{Nhóm: \texttt{HINH\_311}} \quad | \quad \textbf{MapID: \texttt{9H9V1-1}} \hfill \textcolor{yellow!90!white}{[Hình gốc .pyc]}}
]
\textbf{Dạng toán:} TF-M3. Tính số đo các góc 4 \hfill \textit{(Đúng-Sai -- Vận dụng)}
\vspace{4pt}\hrule\vspace{6pt}
\begin{minipage}[t]{0.58\textwidth}
\vspace{0pt}
Cho hai điểm ${E}$, ${H}$ nằm trên đường tròn $({O'})$ và cho điểm ${B}$ nằm trên cung lớn $\stackrel\frown{EH}$. Biết rằng $\widehat{O'EB} = 32^{\circ}$, $\widehat{O'HB} = 38^{\circ}$. Xét tính đúng sai của các khẳng định sau.

\vspace{4pt}\noindent
\textbf{a)} $\widehat{BEH} = 42^{\circ}$\\[2pt]
\textbf{b)} $\widehat{BHE}= 68^{\circ}$\\[2pt]
\textbf{\textcolor{blue!75!black}{c) (Đ)}} $\widehat{EBH}=70^{\circ}$\\[2pt]
\textbf{d)} $\widehat{O'EH}=40^{\circ} $
\end{minipage}\hfill
\begin{minipage}[t]{0.39\textwidth}
\vspace{0pt}
\centering

\resizebox{\linewidth}{!}{%
\begin{tikzpicture}[>=stealth,line join=round,line cap=round,font=\footnotesize,scale=1]
                \path 
                (0,0) coordinate (O)
                ($(O) + (100:2)$) coordinate (A)
                ($(O) + (-15:2)$) coordinate (B)
                ($(O) + (-165:2)$) coordinate (C)
                ;
                \draw (O) circle (2cm);
                \draw (C)--(B)--(A)--(C)--(A)--(O)--(B) (O)--(C);
                     \node[above left] at (A) {$B$};
      \node[above right] at (B) {$E$};
       \node[above left] at (C) {$H$};
        \node[above left] at (O) {$O'$};
                \fill[black] (B) +(150:7.3mm) node[scale=0.8]{$32^\circ$};
                \fill[black] (C) +(30:5.5mm) node[scale=0.8]{$38^\circ$};
            \end{tikzpicture}
}
\end{minipage}

\vspace{3pt}\hrule\vspace{3pt}\footnotesize\textcolor{gray!80!black}{\textbf{Các mã dùng chung hình (HINH\_311):} \texttt{L09\_C9\_B1\_12, L09\_C9\_B1\_06}}
\end{tcolorbox}
\vspace{4pt}


\begin{tcolorbox}[
    breakable,
    colback=white,
    colframe=blue!60!black,
    boxrule=0.8pt,
    arc=3pt,
    title={\textbf{Câu \#415} \quad \textbf{MÃ CÂU: \texttt{[L09\_C9\_B1\_01]}} \quad | \quad \textbf{Nhóm: \texttt{HINH\_312}} \quad | \quad \textbf{MapID: \texttt{9H9N1-1}} \hfill \textcolor{yellow!90!white}{[Hình gốc .pyc]}}
]
\textbf{Dạng toán:} Nhận biết góc nội tiếp \hfill \textit{(Tự luận -- Nhận biết)}
\vspace{4pt}\hrule\vspace{6pt}
\begin{minipage}[t]{0.58\textwidth}
\vspace{0pt}
Cho tam giác ${ADE}$ có ba đỉnh nằm trên đường tròn $({O})$. Hãy chỉ ra các góc nội tiếp của đường tròn và nó chắn cung nào?
\end{minipage}\hfill
\begin{minipage}[t]{0.39\textwidth}
\vspace{0pt}
\centering

\resizebox{\linewidth}{!}{%
\begin{tikzpicture}[scale=1, font=\footnotesize, line join=round, line cap=round, >=stealth]
            \def\r{2}   
            \path (0,0) coordinate (O)
            ($(O) + (110:\r)$) coordinate (A)
            ($(O) + (-160:\r)$) coordinate (B)
            ($(O) + (-30:\r)$) coordinate (C);
            \pgfresetboundingbox
            \draw (O) circle (\r)
            (A)--(B)--(C)--(A)--(O)--(B) (O)--(C);
              \filldraw[black] (A) circle (2pt) node[above ] {$A$};
    \filldraw[black] (B) circle (2pt) node[above left] {$D$};
    \filldraw[black] (C) circle (2pt) node[above right] {$E$};
   \filldraw[black] (O) circle (2pt) node[below]{$O$};
        \end{tikzpicture}
}
\end{minipage}

\end{tcolorbox}
\vspace{4pt}


\begin{tcolorbox}[
    breakable,
    colback=white,
    colframe=blue!60!black,
    boxrule=0.8pt,
    arc=3pt,
    title={\textbf{Câu \#416} \quad \textbf{MÃ CÂU: \texttt{[L09\_C9\_B1\_02]}} \quad | \quad \textbf{Nhóm: \texttt{HINH\_313}} \quad | \quad \textbf{MapID: \texttt{9H9H1-1}} \hfill \textcolor{yellow!90!white}{[Hình gốc .pyc]}}
]
\textbf{Dạng toán:} Tính góc \hfill \textit{(Tự luận -- Thông hiểu)}
\vspace{4pt}\hrule\vspace{6pt}
\begin{minipage}[t]{0.58\textwidth}
\vspace{0pt}
Trên đường tròn tâm ${I}$ đường kính ${AD}$ lấy điểm ${G}$ sao cho $\widehat{AIG}=64^{\circ}$. Hãy tính các góc của $\triangle {ADG}$.
\end{minipage}\hfill
\begin{minipage}[t]{0.39\textwidth}
\vspace{0pt}
\centering

\resizebox{\linewidth}{!}{%
\begin{tikzpicture}[scale=1, font=\footnotesize, line join=round, line cap=round, >=stealth]
                \def\r{2}   
                \path (0,0) coordinate (O)
                ($(O) + (180:\r)$) coordinate (A)
                ($(O) + (0:\r)$) coordinate (B)
                ($(O) + (110:\r)$) coordinate (C);
                \pgfresetboundingbox
                \draw (O) circle (\r)
                (A)--(B)--(C)--(A)--(O)--(B) (O)--(C);
                \draw ($(O) + (150:0.5)$) node{$64^\circ$};
                   \filldraw[black] (A) circle (2pt) node[above left ] {$A$};
    \filldraw[black] (B) circle (2pt) node[above right] {$D$};
    \filldraw[black] (C) circle (2pt) node[above ] {$G$};
   \filldraw[black] (O) circle (2pt) node[below]{$I$};
            \end{tikzpicture}
}
\end{minipage}

\end{tcolorbox}
\vspace{4pt}


\begin{tcolorbox}[
    breakable,
    colback=white,
    colframe=blue!60!black,
    boxrule=0.8pt,
    arc=3pt,
    title={\textbf{Câu \#417} \quad \textbf{MÃ CÂU: \texttt{[L09\_C9\_B1\_03]}} \quad | \quad \textbf{Nhóm: \texttt{HINH\_314}} \quad | \quad \textbf{MapID: \texttt{9H9N1-1}} \hfill \textcolor{yellow!90!white}{[Hình gốc .pyc]}}
]
\textbf{Dạng toán:} Nhận biết góc nội tiếp qua các hình \hfill \textit{(Tự luận -- Nhận biết)}
\vspace{4pt}\hrule\vspace{6pt}
\begin{minipage}[t]{0.58\textwidth}
\vspace{0pt}
Hình nào dưới đây vẽ một góc nội tiếp của đường tròn?.
\end{minipage}\hfill
\begin{minipage}[t]{0.39\textwidth}
\vspace{0pt}
\centering

\resizebox{\linewidth}{!}{%
\begin{tikzpicture}[>=stealth,line join=round,line cap=round,font=\footnotesize,scale=1]
        %\clip (-2.4,-2.3) rectangle (2.4,4.6);
        \path 
        (0,0) coordinate (O)
        ($(O) + (-30:1.7)$) coordinate (A)
        ($(O) + (-115:1.7)$) coordinate (B)
        ;
        \draw (O) circle (1.7cm);
        \draw (A)--(O)--(B);
        \foreach \x in {O,A,B}\fill[black] (\x) circle (1pt);
        \fill (0,-2.1) node[below]{a)};
        % hinh b
        \path 
        (4.5,0) coordinate (O1)
        ($(O1) + (110:1.7)$) coordinate (A1)
        ($(A1) + (0,0.6)$) coordinate (a)
        ($(O1) + (-30:1.7)$) coordinate (B1)
        ($(O1) + (-115:1.7)$) coordinate (C1)
        ;
        \draw (O1) circle (1.7cm);
        \draw (B1)--(a)--(C1);
        \foreach \x in {O1,a,B1,C1}\fill[black] (\x) circle (1pt);
        \fill (5,-2.1) node[below]{b)};
        % hinh c
        \path 
        (9,0) coordinate (O2)
        ($(O2) + (110:1.7)$) coordinate (A2)
        ($(O2) + (-30:1.7)$) coordinate (B2)
        ($(O2) + (-115:1.7)$) coordinate (C2)
        ;
        \draw (O2) circle (1.7cm);
        \draw (B2)--(A2)--(C2);
        \foreach \x in {O2,A2,B2,C2}\fill[black] (\x) circle (1pt);
        \fill (9,-2.1) node[below]{c)};
        % hinh d
        \path 
        (13.5,0) coordinate (O3)
        ($(O3) + (90:1.7)$) coordinate (A3)
        ($(A3) + (0,-0.6)$) coordinate (a3)
        ($(O3) + (-90:1.7)$) coordinate (B3)
        ($(O3) + (-130:1.7)$) coordinate (C3)
        ;
        \draw (O3) circle (1.7cm);
        \draw (B3)--(a3)--(C3);
        \foreach \x in {O3,a3,B3,C3}\fill[black] (\x) circle (1pt);
        \fill (13.5,-2.1) node[below]{d)};
    \end{tikzpicture}
}
\end{minipage}

\end{tcolorbox}
\vspace{4pt}


\begin{tcolorbox}[
    breakable,
    colback=white,
    colframe=blue!60!black,
    boxrule=0.8pt,
    arc=3pt,
    title={\textbf{Câu \#418} \quad \textbf{MÃ CÂU: \texttt{[L09\_C9\_B1\_04]}} \quad | \quad \textbf{Nhóm: \texttt{HINH\_315}} \quad | \quad \textbf{MapID: \texttt{9H9H1-1}} \hfill \textcolor{yellow!90!white}{[Hình gốc .pyc]}}
]
\textbf{Dạng toán:} Tính số đo các góc 1 \hfill \textit{(Tự luận -- Thông hiểu)}
\vspace{4pt}\hrule\vspace{6pt}
\begin{minipage}[t]{0.58\textwidth}
\vspace{0pt}
Cho ${BC}$ và ${GH}$ là hai đường kính của đường tròn $({O})$. Biết rằng $\widehat{BOG} = 64^{\circ}$, tính số đo của $\widehat{BCG}$, $\widehat{BHG}$ và $\widehat{BCH}$.
\end{minipage}\hfill
\begin{minipage}[t]{0.39\textwidth}
\vspace{0pt}
\centering

\resizebox{\linewidth}{!}{%
\begin{tikzpicture}[>=stealth,line join=round,line cap=round,font=\footnotesize,scale=1]
                \path 
                (0,0) coordinate (O)
                ($(O) + (-15:2)$) coordinate (A)
                ($(A)!2!(O)$) coordinate (B)
                ($(O) + (60:2)$) coordinate (C)
                ($(C)!2!(O)$) coordinate (D)
                ;
                \draw (O) circle (2cm);
                \draw (B)--(A)--(D)--(B)--(C)--(D);
                 \filldraw[black] (A) circle (2pt) node[above right] {$B$};
    \filldraw[black] (B) circle (2pt) node[above left] {$C$};
    \filldraw[black] (C) circle (2pt) node[above] {$G$};
   \filldraw[black] (D) circle (2pt) node[above left]{$H$};
    \filldraw[black] (O) circle (2pt) node[below left] {$O$};
            \end{tikzpicture}
}
\end{minipage}

\end{tcolorbox}
\vspace{4pt}


\begin{tcolorbox}[
    breakable,
    colback=white,
    colframe=blue!60!black,
    boxrule=0.8pt,
    arc=3pt,
    title={\textbf{Câu \#419} \quad \textbf{MÃ CÂU: \texttt{[L09\_C9\_B1\_05]}} \quad | \quad \textbf{Nhóm: \texttt{HINH\_310}} \quad | \quad \textbf{MapID: \texttt{9H9H1-1}} \hfill \textcolor{yellow!90!white}{[Hình gốc .pyc]}}
]
\textbf{Dạng toán:} Tính số đo các góc 2 \hfill \textit{(Tự luận -- Thông hiểu)}
\vspace{4pt}\hrule\vspace{6pt}
\begin{minipage}[t]{0.58\textwidth}
\vspace{0pt}
Cho đường tròn $({I})$ và hai dây cung ${BC}$, ${DF}$ cắt nhau tại ${H}$. Tính số đo các góc của tam giác ${BHD}$, biết rằng $\widehat{HCF} = 57^{\circ}$, $\widehat{HFC} = 65^{\circ}$.
\end{minipage}\hfill
\begin{minipage}[t]{0.39\textwidth}
\vspace{0pt}
\centering

\resizebox{\linewidth}{!}{%
\begin{tikzpicture}[>=stealth,line join=round,line cap=round,font=\footnotesize,scale=1]
                \path 
                (0,0) coordinate (O)
                ($(O) + (32:2)$) coordinate (A)
                ($(O) + (205:2)$) coordinate (B)
                ($(O) + (-10:2)$) coordinate (C)
                ($(O) + (145:2)$) coordinate (D)
                (intersection of A--B and C--D) coordinate (X)
                ;
                \draw (O) circle (2cm);
                \draw (B)--(A)--(C)--(D)--(B);
                    \filldraw[black] (A) circle (1pt) node[above right] {$B$};
    \filldraw[black] (B) circle (1pt) node[above left] {$C$};
    \filldraw[black] (C) circle (1pt) node[above right] {$D$};
   \filldraw[black] (D) circle (1pt) node[above left]{$F$};
    \filldraw[black] (O) circle (1pt) node[below left] {$I$};
     \filldraw[black] (X) circle (1pt) node[above] {$H$};
                \fill[black] (D) +(-50:5mm) node{$65^\circ$};
                \fill[black] (B) +(50:5.5mm) node{$57^\circ$};
                
            \end{tikzpicture}
}
\end{minipage}

\vspace{3pt}\hrule\vspace{3pt}\footnotesize\textcolor{gray!80!black}{\textbf{Các mã dùng chung hình (HINH\_310):} \texttt{L09\_C9\_B1\_11, L09\_C9\_B1\_15}}
\end{tcolorbox}
\vspace{4pt}


\begin{tcolorbox}[
    breakable,
    colback=white,
    colframe=blue!60!black,
    boxrule=0.8pt,
    arc=3pt,
    title={\textbf{Câu \#420} \quad \textbf{MÃ CÂU: \texttt{[L09\_C9\_B1\_06]}} \quad | \quad \textbf{Nhóm: \texttt{HINH\_311}} \quad | \quad \textbf{MapID: \texttt{9H9H1-1}} \hfill \textcolor{yellow!90!white}{[Hình gốc .pyc]}}
]
\textbf{Dạng toán:} Tính số đo các góc 3 \hfill \textit{(Tự luận -- Thông hiểu)}
\vspace{4pt}\hrule\vspace{6pt}
\begin{minipage}[t]{0.58\textwidth}
\vspace{0pt}
Cho hai điểm ${F}$, ${H}$ nằm trên đường tròn $({O})$ và cho điểm ${E}$ nằm trên cung lớn $\stackrel\frown{FH}$. Biết rằng $\widehat{OFE} = 30^{\circ}$, $\widehat{OHE} = 35^{\circ}$. Tính số đo các góc của tam giác ${EFH}$.
\end{minipage}\hfill
\begin{minipage}[t]{0.39\textwidth}
\vspace{0pt}
\centering

\resizebox{\linewidth}{!}{%
\begin{tikzpicture}[>=stealth,line join=round,line cap=round,font=\footnotesize,scale=1]
                \path 
                (0,0) coordinate (O)
                ($(O) + (100:2)$) coordinate (A)
                ($(O) + (-15:2)$) coordinate (B)
                ($(O) + (-165:2)$) coordinate (C)
                ;
                \draw (O) circle (2cm);
                \draw (C)--(B)--(A)--(C)--(A)--(O)--(B) (O)--(C);
                     \node[above left] at (A) {$E$};
      \node[above right] at (B) {$F$};
       \node[above left] at (C) {$H$};
        \node[above left] at (O) {$O$};
                \fill[black] (B) +(150:7.3mm) node[scale=0.8]{$30^\circ$};
                \fill[black] (C) +(30:5.5mm) node[scale=0.8]{$35^\circ$};
            \end{tikzpicture}
}
\end{minipage}

\vspace{3pt}\hrule\vspace{3pt}\footnotesize\textcolor{gray!80!black}{\textbf{Các mã dùng chung hình (HINH\_311):} \texttt{L09\_C9\_B1\_12, L09\_C9\_B1\_16}}
\end{tcolorbox}
\vspace{4pt}


\begin{tcolorbox}[
    breakable,
    colback=white,
    colframe=blue!60!black,
    boxrule=0.8pt,
    arc=3pt,
    title={\textbf{Câu \#421} \quad \textbf{MÃ CÂU: \texttt{[L09\_C9\_B1\_07]}} \quad | \quad \textbf{Nhóm: \texttt{HINH\_316}} \quad | \quad \textbf{MapID: \texttt{9H9H1-1}} \hfill \textcolor{yellow!90!white}{[Hình gốc .pyc]}}
]
\textbf{Dạng toán:} Chứng minh 2 đoạn thẳng bằng nhau \hfill \textit{(Tự luận -- Thông hiểu)}
\vspace{4pt}\hrule\vspace{6pt}
\begin{minipage}[t]{0.58\textwidth}
\vspace{0pt}
Cho đường tròn $({O})$ có đường kính ${BC}$ và điểm ${D}$ thuộc đường tròn sao cho ${D}$ khác ${B}$, ${C}$. Lấy ${G}$ là điểm đối xứng của ${B}$ qua ${D}$. Chứng minh rằng $CG = CB$.
\end{minipage}\hfill
\begin{minipage}[t]{0.39\textwidth}
\vspace{0pt}
\centering

\resizebox{\linewidth}{!}{%
\begin{tikzpicture}[>=stealth,line join=round,line cap=round,font=\footnotesize,scale=1]
    \path 
    (0,0) coordinate (O)
    ($(O) + (-120:2)$) coordinate (A)
    ($(O) + (60:2)$) coordinate (B)
    ($(O) + (-30:2)$) coordinate (C)
    ($(A)!2!(C)$) coordinate (E)
    ;
    \draw (O) circle (2cm);
    \draw (A)--(E)--(B)--(A) (B)--(C);
        \filldraw[black] (A) circle (2pt) node[below left] {$B$};
    \filldraw[black] (B) circle (1pt) node[above left] {$C$};
    \filldraw[black] (C) circle (1pt) node[above right] {$D$};
   \filldraw[black] (E) circle (1pt) node[above right]{$G$};
    \filldraw[black] (O) circle (1pt) node[above] {$O$};
    
    \path 
    (C) edge node[midway, sloped, rotate=60, anchor=center] {$ - $} (A)
    (C) edge node[midway, sloped, rotate=60, anchor=center] {$ - $} (E);
    \draw pic [draw, angle radius=2mm] {right angle=A--C--B};
\end{tikzpicture}
}
\end{minipage}

\end{tcolorbox}
\vspace{4pt}


\begin{tcolorbox}[
    breakable,
    colback=white,
    colframe=blue!60!black,
    boxrule=0.8pt,
    arc=3pt,
    title={\textbf{Câu \#422} \quad \textbf{MÃ CÂU: \texttt{[L09\_C9\_B1\_08]}} \quad | \quad \textbf{Nhóm: \texttt{HINH\_317}} \quad | \quad \textbf{MapID: \texttt{9H9H1-1}} \hfill \textcolor{yellow!90!white}{[Hình gốc .pyc]}}
]
\textbf{Dạng toán:} Bài tổng hợp \hfill \textit{(Tự luận -- Thông hiểu)}
\vspace{4pt}\hrule\vspace{6pt}
\begin{minipage}[t]{0.58\textwidth}
\vspace{0pt}
Cho tam giác nhọn ${CDG}$ cân tại ${C}$. Đường tròn đường kính ${CD}$ cắt các cạnh ${CG}$, ${DG}$ của tam giác ${CDG}$ tại ${J}$ và ${K}$ (${J}$ khác ${C}$, ${K}$ khác ${D}$).

a) Chứng minh rằng tam giác ${GJK}$ cân tại ${K}$.

b) Cho ${DJ}$ cắt ${CK}$ tại ${N}$. Chứng minh rằng ${GN}$ vuông góc với ${CD}$.
\end{minipage}\hfill
\begin{minipage}[t]{0.39\textwidth}
\vspace{0pt}
\centering

\resizebox{\linewidth}{!}{%
\begin{tikzpicture}[>=stealth,line join=round,line cap=round,font=\footnotesize,scale=1]
                \path 
                (2,2.8) coordinate (A)
                (4,0) coordinate (B)
                (0,0) coordinate (C)
                ($(A)!0.5!(B)$) coordinate (O)
                ($(A)!(B)!(C)$) coordinate (X)
                ($(B)!(A)!(C)$) coordinate (Y)
                (intersection of B--X and A--Y) coordinate (K)
                ;
                \draw (O) circle (1.72cm);
                \draw (A)--(B)--(C)--(A)--(Y) (B)--(X) (C)--(K);
                  \filldraw[black] (A) circle (2pt) node[above] {$C$};
    \filldraw[black] (B) circle (1pt) node[below right] {$D$};
    \filldraw[black] (C) circle (1pt) node[below left] {$G$};
   \filldraw[black] (X) circle (1pt) node[above left]{$J$};
    \filldraw[black] (Y) circle (1pt) node[below] {$K$};
                \fill (O) circle(1pt);
                \foreach \a/\b/\c in {A/Y/B, B/X/A} {\draw pic [draw, angle radius=2mm] {right angle=\a--\b--\c};}
            \end{tikzpicture}
}
\end{minipage}

\end{tcolorbox}
\vspace{4pt}


\begin{tcolorbox}[
    breakable,
    colback=white,
    colframe=blue!60!black,
    boxrule=0.8pt,
    arc=3pt,
    title={\textbf{Câu \#423} \quad \textbf{MÃ CÂU: \texttt{[L09\_C9\_B2\_04]}} \quad | \quad \textbf{Nhóm: \texttt{HINH\_318}} \quad | \quad \textbf{MapID: \texttt{9H9V2-1}} \hfill \textcolor{yellow!90!white}{[Hình gốc .pyc]}}
]
\textbf{Dạng toán:} Tính bán kính đường tròn nội tiếp tam giác vuông khi biết 2 cạnh \hfill \textit{(Tự luận -- Vận dụng)}
\vspace{4pt}\hrule\vspace{6pt}
\begin{minipage}[t]{0.58\textwidth}
\vspace{0pt}
Cho tam giác ${ABC}$ vuông tại ${A}$ có $AB={28}$ cm, $AC={45}$ cm. Bán kính đường tròn nội tiếp tam giác là
\end{minipage}\hfill
\begin{minipage}[t]{0.39\textwidth}
\vspace{0pt}
\centering

\resizebox{\linewidth}{!}{%
\begin{tikzpicture}[line cap=round,line join=round,scale=1,font=\footnotesize]
            \def\r{2.5}
            \path (0,0) coordinate (O)
            (-150:\r) coordinate (A)
            (150:\r) coordinate (B)
            (-30:\r) coordinate (C)
            ($ (A)!1cm!(B) $) coordinate (Nt)
            ($ (A)!1cm!(C) $) coordinate (Mt)
            ($ (Mt)!0.5!(Nt) $) coordinate (Xt)
            ($ (B)!1cm!(C) $) coordinate (Kt)
            ($ (B)!1cm!(A) $) coordinate (Lt)
            ($ (Kt)!0.5!(Lt) $) coordinate (Yt)
            (intersection of A--Xt and B--Yt) coordinate (I)
            ($ (A)!(I)!(B) $) coordinate (F)
            ($ (B)!(I)!(C) $) coordinate (D)
            ($ (C)!(I)!(A) $) coordinate (E);
            \draw (A)--(B)--(C)--cycle (E)--(I)--(D) (I)--(F) (B)--(I)--(C) (I)--(A);
            \draw let \p1=($  (I) -  (D) $) in  (I)  circle ({veclen(\x1,\y1)});
            \foreach \x/\g in {A/-135,B/140,C/-20,I/30}{\fill (\x) circle (1pt)node[shift={(\g:.3)}]{$\x$};}
            \pic[draw,angle radius=8pt]{right angle=C--A--B};
            \path (I)--(D) node[midway,left=0pt]{$r$};
            \path (A)--(B)node[midway,left]{$28cm$};
            \path (A)--(C)node[midway,below]{$45cm$};
    \end{tikzpicture}
}
\end{minipage}

\end{tcolorbox}
\vspace{4pt}


\begin{tcolorbox}[
    breakable,
    colback=white,
    colframe=blue!60!black,
    boxrule=0.8pt,
    arc=3pt,
    title={\textbf{Câu \#424} \quad \textbf{MÃ CÂU: \texttt{[L09\_C9\_B2\_05]}} \quad | \quad \textbf{Nhóm: \texttt{HINH\_319}} \quad | \quad \textbf{MapID: \texttt{9H9V2-1}} \hfill \textcolor{yellow!90!white}{[Hình gốc .pyc]}}
]
\textbf{Dạng toán:} Chứng minh các góc bằng nhau. \hfill \textit{(Tự luận -- Vận dụng)}
\vspace{4pt}\hrule\vspace{6pt}
\begin{minipage}[t]{0.58\textwidth}
\vspace{0pt}
Cho tam giác ${BCE}$ có trực tâm ${G}$ và nội tiếp đường tròn $({O})$. Chứng minh rằng:

a) $\widehat{OCE}=90^{\circ}-\widehat{CBE}$;

b) $\widehat{CBG}=\widehat{OBE}$.
\end{minipage}\hfill
\begin{minipage}[t]{0.39\textwidth}
\vspace{0pt}
\centering

\resizebox{\linewidth}{!}{%
\begin{tikzpicture}[line cap=round,line join=round,scale=1,font=\footnotesize]
\def\r{2}
\path (0,0) coordinate (O)
(120:\r) coordinate (A)
(-150:\r) coordinate (B)
(-30:\r) coordinate (C)
($(A)!(B)!(C)$) coordinate (E)
($(B)!(A)!(C)$) coordinate (D)
(intersection of A--D and B--E) coordinate (H)
;
\draw (A)--(B)--(C)--cycle (O) circle (\r) (A)--(D) (B)--(O)--(C);
    \filldraw[black] (A) circle (1pt) node[above right] {$B$};
    \filldraw[black] (B) circle (1pt) node[below left] {$C$};
    \filldraw[black] (C) circle (1pt) node[below right] {$E$};
   \filldraw[black] (D) circle (1pt) node[below]{$F$};
    \filldraw[black] (H) circle (1pt) node[below left] {$G$};
      \filldraw[black] (O) circle (1pt) node[above] {$O$};
\end{tikzpicture}
}
\end{minipage}

\end{tcolorbox}
\vspace{4pt}


\begin{tcolorbox}[
    breakable,
    colback=white,
    colframe=blue!60!black,
    boxrule=0.8pt,
    arc=3pt,
    title={\textbf{Câu \#425} \quad \textbf{MÃ CÂU: \texttt{[L09\_C9\_B2\_06]}} \quad | \quad \textbf{Nhóm: \texttt{HINH\_320}} \quad | \quad \textbf{MapID: \texttt{9H9V2-1}} \hfill \textcolor{yellow!90!white}{[Hình gốc .pyc]}}
]
\textbf{Dạng toán:} Chứng minh các góc bằng nhau 2 \hfill \textit{(Tự luận -- Vận dụng)}
\vspace{4pt}\hrule\vspace{6pt}
\begin{minipage}[t]{0.58\textwidth}
\vspace{0pt}
Cho tam giác ${AGH}$ ngoại tiếp đường tròn $({I})$. 

 Chứng minh rằng $\widehat{GIH}=90^{\circ}+\dfrac{\widehat{GAH} }{2};\widehat{HIA}=90^{\circ}+\dfrac{\widehat{HGA}}{2};\widehat{AIG}=90^{\circ}+\dfrac{\widehat{AHG}}{2}$
\end{minipage}\hfill
\begin{minipage}[t]{0.39\textwidth}
\vspace{0pt}
\centering

\resizebox{\linewidth}{!}{%
\begin{tikzpicture}[line cap=round,line join=round,scale=1,font=\footnotesize]  
            \def\r{2.5}
            \path (0,0) coordinate (O)
            (120:\r) coordinate (A)
            (-160:\r) coordinate (B)
            (-20:\r) coordinate (C)
            ($ (A)!1cm!(B) $) coordinate (Nt)
            ($ (A)!1cm!(C) $) coordinate (Mt)
            ($ (Mt)!0.5!(Nt) $) coordinate (Xt)
            ($ (B)!1cm!(C) $) coordinate (Kt)
            ($ (B)!1cm!(A) $) coordinate (Lt)
            ($ (Kt)!0.5!(Lt) $) coordinate (Yt)
            (intersection of A--Xt and B--Yt) coordinate (I)
            ($ (A)!(I)!(B) $) coordinate (F)
            ($ (B)!(I)!(C) $) coordinate (D)
            ($ (C)!(I)!(A) $) coordinate (E);
            \draw (A)--(B)--(C)--cycle (B)--(I)--(C) (A)--(I);
            \draw let \p1=($  (I) -  (D) $) in  (I)  circle ({veclen(\x1,\y1)});
    \filldraw[black] (A) circle (1pt) node[above right] {$A$};
    \filldraw[black] (B) circle (1pt) node[below left] {$G$};
    \filldraw[black] (C) circle (1pt) node[below right] {$H$};
   \filldraw[black] (I) circle (1pt) node[below]{$I$};
   
\end{tikzpicture}
}
\end{minipage}

\end{tcolorbox}
\vspace{4pt}


\begin{tcolorbox}[
    breakable,
    colback=white,
    colframe=blue!60!black,
    boxrule=0.8pt,
    arc=3pt,
    title={\textbf{Câu \#426} \quad \textbf{MÃ CÂU: \texttt{[L09\_C9\_B2\_07]}} \quad | \quad \textbf{Nhóm: \texttt{HINH\_321}} \quad | \quad \textbf{MapID: \texttt{9H9H2-1}} \hfill \textcolor{yellow!90!white}{[Hình gốc .pyc]}}
]
\textbf{Dạng toán:} Chứng minh tam giác đồng dạng \hfill \textit{(Tự luận -- Thông hiểu)}
\vspace{4pt}\hrule\vspace{6pt}
\begin{minipage}[t]{0.58\textwidth}
\vspace{0pt}
Cho tam giác nhọn ${ACE}$ nội tiếp đường tròn $({I})$. Cho điểm ${H}$ trên cạnh ${CE}$ của tam giác ${ACE}$ và điểm ${F}$ trên cung nhỏ ${CE}$ của $({I})$ sao cho $\widehat{CAF}=\widehat{HAE}$. Chứng minh rằng $\triangle AHC \backsim \triangle AEF$.
\end{minipage}\hfill
\begin{minipage}[t]{0.39\textwidth}
\vspace{0pt}
\centering

\resizebox{\linewidth}{!}{%
\begin{tikzpicture}[line cap=round,line join=round,scale=1,font=\footnotesize]
\def\r{2.5}
\path (0,0) coordinate (O)
(120:\r) coordinate (A)
(-150:\r) coordinate (B)
(-30:\r) coordinate (C)
(-70:\r) coordinate (AM)
(intersection of A--AM and B--C) coordinate (M)
(-110:\r) coordinate (D)
;
\draw (A)--(B)--(C)--cycle (O) circle (\r) (M)--(A)--(D)--(C);
\pic[draw,angle radius=14pt]{angle=B--A--D};
\pic[draw,angle radius=14pt]{angle=M--A--C};
    \filldraw[black] (A) circle (1pt) node[above right] {$A$};
    \filldraw[black] (B) circle (1pt) node[below left] {$C$};
    \filldraw[black] (C) circle (1pt) node[above right] {$E$};
   \filldraw[black] (D) circle (1pt) node[below left]{$F$};
    \filldraw[black] (M) circle (1pt) node[below left] {$H$};
    \filldraw[black] (O) circle (1pt) node[above] {$I$};
    \end{tikzpicture}
}
\end{minipage}

\end{tcolorbox}
\vspace{4pt}


\begin{tcolorbox}[
    breakable,
    colback=white,
    colframe=blue!60!black,
    boxrule=0.8pt,
    arc=3pt,
    title={\textbf{Câu \#427} \quad \textbf{MÃ CÂU: \texttt{[L09\_C9\_B2\_08]}} \quad | \quad \textbf{Nhóm: \texttt{HINH\_322}} \quad | \quad \textbf{MapID: \texttt{9H9V2-1}} \hfill \textcolor{yellow!90!white}{[Hình gốc .pyc]}}
]
\textbf{Dạng toán:} Tính số đo góc \hfill \textit{(Tự luận -- Vận dụng)}
\vspace{4pt}\hrule\vspace{6pt}
\begin{minipage}[t]{0.58\textwidth}
\vspace{0pt}
Cho tam giác ${AEF}$ vuông tại ${E}$ nội tiếp đường tròn $({I})$ và đường kính ${EH}$. Tính số đo của góc ${EAF}$, biết rằng $\widehat{EAF}=17 \widehat{FEH}$.
\end{minipage}\hfill
\begin{minipage}[t]{0.39\textwidth}
\vspace{0pt}
\centering

\resizebox{\linewidth}{!}{%
\begin{tikzpicture}[line cap=round,line join=round,scale=1,font=\footnotesize]

\def\r{2.5}
\path (0,0) coordinate (O)
(120:\r) coordinate (A)
(180:\r) coordinate (B)
(-60:\r) coordinate (C)
(0:\r) coordinate (D)
;
\draw (C)--(B)--(A)--(D)--(B) (A)--(C) (O) circle (\r);

    \filldraw[black] (A) circle (1pt) node[above right] {$A$};
    \filldraw[black] (B) circle (1pt) node[below left] {$E$};
    \filldraw[black] (C) circle (1pt) node[below right] {$F$};
   \filldraw[black] (D) circle (1pt) node[below right]{$H$};
    \filldraw[black] (O) circle (1pt) node[above] {$I$};
    \end{tikzpicture}
}
\end{minipage}

\end{tcolorbox}
\vspace{4pt}


\begin{tcolorbox}[
    breakable,
    colback=white,
    colframe=blue!60!black,
    boxrule=0.8pt,
    arc=3pt,
    title={\textbf{Câu \#428} \quad \textbf{MÃ CÂU: \texttt{[L09\_C9\_B2\_09]}} \quad | \quad \textbf{Nhóm: \texttt{HINH\_323}} \quad | \quad \textbf{MapID: \texttt{9H9V2-1}} \hfill \textcolor{yellow!90!white}{[Hình gốc .pyc]}}
]
\textbf{Dạng toán:} Chứng minh một điểm là tâm đường tròn ngoại tiếp tam giác \hfill \textit{(Tự luận -- Vận dụng)}
\vspace{4pt}\hrule\vspace{6pt}
\begin{minipage}[t]{0.58\textwidth}
\vspace{0pt}
Cho tam giác ${DEG}$ nội tiếp đường tròn $({I})$ và ngoại tiếp đường tròn $({I'})$. Tia ${DI'}$ cắt $({I})$ tại ${H}$ (khác ${D}$). Chứng minh rằng ${H}$ là tâm của đường tròn ngoại tiếp tam giác ${EI'G}$.
\end{minipage}\hfill
\begin{minipage}[t]{0.39\textwidth}
\vspace{0pt}
\centering

\resizebox{\linewidth}{!}{%
\begin{tikzpicture}[line cap=round,line join=round,scale=1,font=\footnotesize]

    \def\r{2.5}
    \path (0,0) coordinate (O)
    (140:\r) coordinate (A)
    (-160:\r) coordinate (B)
    (-20:\r) coordinate (C)
    ($ (A)!1cm!(B) $) coordinate (Nt)
    ($ (A)!1cm!(C) $) coordinate (Mt)
    ($ (Mt)!0.5!(Nt) $) coordinate (Xt)
    ($ (B)!1cm!(C) $) coordinate (Kt)
    ($ (B)!1cm!(A) $) coordinate (Lt)
    ($ (Kt)!0.5!(Lt) $) coordinate (Yt)
    (intersection of A--Xt and B--Yt) coordinate (I)
    ($ (A)!(I)!(B) $) coordinate (F)
    ($ (B)!(I)!(C) $) coordinate (D)
    ($ (C)!(I)!(A) $) coordinate (E)
    ($(A)!(O)!(I)$) coordinate (Ot)
    ($(A)!2!(Ot)$) coordinate (X)
    ;
    \draw (A)--(B)--(C)--cycle (B)--(I)--(C) (A)--(X) (B)--(X)--(C) (O) circle (\r);
    \draw let \p1=($  (I) -  (D) $) in  (I)  circle ({veclen(\x1,\y1)});
\filldraw[black] (I) circle (1pt) node[above] {$I'$};
    \filldraw[black] (A) circle (1pt) node[above] {$D$};
    \filldraw[black] (B) circle (1pt) node[below left] {$E$};
    \filldraw[black] (C) circle (1pt) node[below right] {$G$};
   \filldraw[black] (X) circle (1pt) node[below right]{$H$};
    \filldraw[black] (O) circle (1pt) node[above] {$I$};
    \end{tikzpicture}
}
\end{minipage}

\end{tcolorbox}
\vspace{4pt}


\begin{tcolorbox}[
    breakable,
    colback=white,
    colframe=blue!60!black,
    boxrule=0.8pt,
    arc=3pt,
    title={\textbf{Câu \#429} \quad \textbf{MÃ CÂU: \texttt{[L09\_C9\_B2\_10]}} \quad | \quad \textbf{Nhóm: \texttt{HINH\_324}} \quad | \quad \textbf{MapID: \texttt{9H9C2-1}} \hfill \textcolor{yellow!90!white}{[Hình gốc .pyc]}}
]
\textbf{Dạng toán:} Chứng minh tỉ số giữa R R' r r' cảu 2 tam giác đồng dạng \hfill \textit{(Tự luận -- Vận dụng cao)}
\vspace{4pt}\hrule\vspace{6pt}
\begin{minipage}[t]{0.58\textwidth}
\vspace{0pt}
Cho $\triangle ABD \backsim \triangle PRU$ với tỉ số đồng dạng $k>0$. Gọi ${(I, R)}$ và $\left( O' , R'\right)$ lần lượt là đường tròn ngoại tiếp các tam giác ${ABD}$ và ${PRU}$. Gọi $(G, r)$ và $\left(Y, r'\right)$ lần lượt là đường tròn nội tiếp các tam giác ${ABD}$ và ${PRU}$. Chứng minh rằng $\dfrac{R}{R'}=\dfrac{r}{r'}=k$.
\end{minipage}\hfill
\begin{minipage}[t]{0.39\textwidth}
\vspace{0pt}
\centering

\resizebox{\linewidth}{!}{%
\begin{tikzpicture}[line cap=round,line join=round,scale=1,font=\footnotesize]

    \def\r{2.5}
    \path (0,0) coordinate (O)
    (130:\r) coordinate (A)
    (-160:\r) coordinate (B)
    (-20:\r) coordinate (C)
    ($ (A)!1cm!(B) $) coordinate (Nt)
    ($ (A)!1cm!(C) $) coordinate (Mt)
    ($ (Mt)!0.5!(Nt) $) coordinate (Xt)
    ($ (B)!1cm!(C) $) coordinate (Kt)
    ($ (B)!1cm!(A) $) coordinate (Lt)
    ($ (Kt)!0.5!(Lt) $) coordinate (Yt)
    (intersection of A--Xt and B--Yt) coordinate (I)
    ($ (A)!(I)!(B) $) coordinate (F)
    ($ (B)!(I)!(C) $) coordinate (D)
    ($ (C)!(I)!(A) $) coordinate (E)
    ($(A)!(O)!(I)$) coordinate (Ot)
    ;
    \draw (A)--(B)--(C)--cycle (B)--(I)--(C) (I)--(D) (B)--(O)--(C) (O) circle (\r);
    \draw let \p1=($  (I) -  (D) $) in  (I)  circle ({veclen(\x1,\y1)});
    \filldraw[black] (I) circle (1pt) node[above] {$G$};
    \filldraw[black] (A) circle (1pt) node[above] {$A$};
    \filldraw[black] (B) circle (1pt) node[below left] {$B$};
    \filldraw[black] (C) circle (1pt) node[below right] {$D$};
   \filldraw[black] (D) circle (1pt) node[below right]{$F$};
    \filldraw[black] (O) circle (1pt) node[above] {$I$};
\end{tikzpicture}
}
\end{minipage}

\end{tcolorbox}
\vspace{4pt}


\begin{tcolorbox}[
    breakable,
    colback=white,
    colframe=blue!60!black,
    boxrule=0.8pt,
    arc=3pt,
    title={\textbf{Câu \#430} \quad \textbf{MÃ CÂU: \texttt{[L09\_C9\_B3\_01]}} \quad | \quad \textbf{Nhóm: \texttt{HINH\_325}} \quad | \quad \textbf{MapID: \texttt{9H9N3-1}} \hfill \textcolor{yellow!90!white}{[Hình gốc .pyc]}}
]
\textbf{Dạng toán:} Chứng minh các góc bằng nhau \hfill \textit{(Tự luận -- Nhận biết)}
\vspace{4pt}\hrule\vspace{6pt}
\begin{minipage}[t]{0.58\textwidth}
\vspace{0pt}
Cho tam giác nhọn ${BGH}$ có các đường cao ${BI}$, ${GJ}$, ${HL}$ cắt nhau tại ${M}$. Chứng minh rằng:

a) $\widehat{BJL}=\widehat{BGH};\widehat{BLJ}=\widehat{BHG}$;

b) $\widehat{GMH}=180^{\circ}-\widehat{GBH};\widehat{HMB}=180^{\circ}-\widehat{HGB};\widehat{BMG}=180^{\circ}-\widehat{BHG}$.
\end{minipage}\hfill
\begin{minipage}[t]{0.39\textwidth}
\vspace{0pt}
\centering

\resizebox{\linewidth}{!}{%
\begin{tikzpicture}[scale=1, font=\footnotesize, line join=round, line cap=round, >=stealth]
    \def \R{2}
    \path 
    (0:0) coordinate (O)
    +(180:\R) coordinate (B)
    +(0:\R) coordinate (C)
    +(80:\R) coordinate (E)
    +(145:\R) coordinate (F)
    (intersection of F--B and C--E) coordinate (A)
    (intersection of E--B and C--F) coordinate (H)
    (intersection of H--A and C--B) coordinate (D)
    ($(A)!.5!(H)$) coordinate (I)
    
    ;
    \draw (O) circle (\R)
    (B)--(A)--(C)--cycle
    (F)--(E)--(D)--cycle
    (A)--(D)
    let \p1=($(I)-(A)$) in (I) circle ({veclen(\x1,\y1)})
    (B)--(H)--(C)
    \foreach \x/\y/\z in {B/F/C,A/E/B}{
        pic[draw, thin, angle radius = 6pt]{right angle = \x--\y--\z}
    }
    (F)--(H)--(E)
    ;
      \filldraw[black] (E) circle (1pt) node[above right] {$J$};
    \filldraw[black] (A) circle (1pt) node[above] {$B$};
    \filldraw[black] (B) circle (1pt) node[below left] {$G$};
    \filldraw[black] (C) circle (1pt) node[below right] {$H$};
   \filldraw[black] (D) circle (1pt) node[below right]{$I$};
    \filldraw[black] (F) circle (1pt) node[below left] {$L$};
     \filldraw[black] (H) circle (1pt) node[above left] {$M$};
    \fill (O) circle (1pt);
\end{tikzpicture}
}
\end{minipage}

\end{tcolorbox}
\vspace{4pt}


\begin{tcolorbox}[
    breakable,
    colback=white,
    colframe=blue!60!black,
    boxrule=0.8pt,
    arc=3pt,
    title={\textbf{Câu \#431} \quad \textbf{MÃ CÂU: \texttt{[L09\_C9\_B3\_02]}} \quad | \quad \textbf{Nhóm: \texttt{HINH\_326}} \quad | \quad \textbf{MapID: \texttt{9H9H3-1}} \hfill \textcolor{yellow!90!white}{[Hình gốc .pyc]}}
]
\textbf{Dạng toán:} Sử dụng góc nội tiếp để CM \hfill \textit{(Tự luận -- Thông hiểu)}
\vspace{4pt}\hrule\vspace{6pt}
\begin{minipage}[t]{0.58\textwidth}
\vspace{0pt}
Cho tứ giác ${ACDE}$ nội tiếp đường tròn $({O})$ sao cho hai tia ${AC}$ và ${ED}$ cắt nhau tại điểm ${H}$, hai đường chéo ${AD}$ và ${CE}$ cắt nhau tại điểm ${F}$. Kí hiệu $\stackrel\frown{AE}$ là cung ${AE}$ không chứa ${C}$ và $\stackrel\frown{CD}$ là cung ${CD}$ không chứa ${A}$. Chứng minh rằng:

a) $\widehat{CHD}=\dfrac{1}{2} \left( sd\stackrel\frown{AE}-sd\stackrel\frown{CD} \right)$;

b) $\widehat{CFD}=\dfrac{1}{2} \left( sd\stackrel\frown{AE}+sd\stackrel\frown{CD} \right)$.
\end{minipage}\hfill
\begin{minipage}[t]{0.39\textwidth}
\vspace{0pt}
\centering

\resizebox{\linewidth}{!}{%
\begin{tikzpicture}[scale=1, font=\footnotesize, line join=round, line cap=round, >=stealth]
    \def \R{2}
    \path 
    (0:0) coordinate (O)
    +(130:\R) coordinate (A)
    +(20:\R) coordinate (B)
    +(-20:\R) coordinate (C)
    +(-60:\R) coordinate (D)
    (intersection of A--C and B--D) coordinate (H)
    (intersection of A--B and C--D) coordinate (K)
    ;
    \draw (O) circle (\R)
    (A)--(B)--(C)--(D)--cycle
    (A)--(C)
    (B)--(D)
    (B)--(K)--(C)
    ;
    
          \filldraw[black] (O) circle (1pt) node[above right] {$O$};
    \filldraw[black] (A) circle (1pt) node[above] {$A$};
    \filldraw[black] (B) circle (1pt) node[above right] {$C$};
    \filldraw[black] (C) circle (1pt) node[below right] {$D$};
   \filldraw[black] (D) circle (1pt) node[below right]{$E$};
    \filldraw[black] (K) circle (1pt) node[below right] {$H$};
     \filldraw[black] (H) circle (1pt) node[above left] {$F$};
 
\end{tikzpicture}
}
\end{minipage}

\end{tcolorbox}
\vspace{4pt}


\begin{tcolorbox}[
    breakable,
    colback=white,
    colframe=blue!60!black,
    boxrule=0.8pt,
    arc=3pt,
    title={\textbf{Câu \#432} \quad \textbf{MÃ CÂU: \texttt{[L09\_C9\_B3\_03]}} \quad | \quad \textbf{Nhóm: \texttt{HINH\_327}} \quad | \quad \textbf{MapID: \texttt{9H9V3-1}} \hfill \textcolor{yellow!90!white}{[Hình gốc .pyc]}}
]
\textbf{Dạng toán:} CM tam giác đồng dạng \hfill \textit{(Tự luận -- Vận dụng)}
\vspace{4pt}\hrule\vspace{6pt}
\begin{minipage}[t]{0.58\textwidth}
\vspace{0pt}
Cho tứ giác ${BCDE}$ nội tiếp đường tròn $({O})$ sao cho hai tia ${BC}$ và ${ED}$ cắt nhau tại điểm ${I}$. Chứng minh rằng:

a) $\triangle IBE\backsim \triangle IDC;\triangle IBD\backsim \triangle IEC$;

b) $\dfrac{ID}{IE}=\dfrac{BD}{BE}\cdot\dfrac{CD}{CE}$.
\end{minipage}\hfill
\begin{minipage}[t]{0.39\textwidth}
\vspace{0pt}
\centering

\resizebox{\linewidth}{!}{%
\begin{tikzpicture}[scale=1, font=\footnotesize, line join=round, line cap=round, >=stealth]
        \def \R{2}
        \path 
        (0:0) coordinate (O)
        +(130:\R) coordinate (A)
        +(20:\R) coordinate (B)
        +(-20:\R) coordinate (C)
        +(-60:\R) coordinate (D)
        (intersection of A--C and B--D) coordinate (H)
        (intersection of A--B and C--D) coordinate (I)
        ;
        \draw (O) circle (\R)
        (A)--(B)--(C)--(D)--cycle
        (A)--(C)
        (B)--(D)
        (B)--(I)--(C)
        ;
          \filldraw[black] (O) circle (1pt) node[above right] {$O$};
    \filldraw[black] (A) circle (1pt) node[above] {$B$};
    \filldraw[black] (B) circle (1pt) node[above right] {$C$};
    \filldraw[black] (C) circle (1pt) node[below right] {$D$};
   \filldraw[black] (D) circle (1pt) node[below right]{$E$};
    \filldraw[black] (I) circle (1pt) node[below right] {$I$};
    
    \end{tikzpicture}
}
\end{minipage}

\end{tcolorbox}
\vspace{4pt}


\begin{tcolorbox}[
    breakable,
    colback=white,
    colframe=blue!60!black,
    boxrule=0.8pt,
    arc=3pt,
    title={\textbf{Câu \#433} \quad \textbf{MÃ CÂU: \texttt{[L09\_C9\_B3\_04]}} \quad | \quad \textbf{Nhóm: \texttt{HINH\_328}} \quad | \quad \textbf{MapID: \texttt{9H9V3-1}} \hfill \textcolor{yellow!90!white}{[Hình gốc .pyc]}}
]
\textbf{Dạng toán:} CM tam giác đồng dạng 2 \hfill \textit{(Tự luận -- Vận dụng)}
\vspace{4pt}\hrule\vspace{6pt}
\begin{minipage}[t]{0.58\textwidth}
\vspace{0pt}
Cho tứ giác ${ABCE}$ nội tiếp đường tròn $({O})$, hai đường chéo ${AC}$ và ${BE}$ cắt nhau tại điểm ${H}$. Chứng minh rằng:

a) $\triangle HAE\backsim\triangle HBC$; $\triangle HAB\backsim\triangle HEC$;

b) $\dfrac{HA}{HC}=\dfrac{BA}{BC}\cdot\dfrac{EA}{EC}$.
\end{minipage}\hfill
\begin{minipage}[t]{0.39\textwidth}
\vspace{0pt}
\centering

\resizebox{\linewidth}{!}{%
\begin{tikzpicture}[scale=1, font=\footnotesize, line join=round, line cap=round, >=stealth]
        \def \R{2}
        \path 
        (0:0) coordinate (O)
        +(130:\R) coordinate (A)
        +(20:\R) coordinate (B)
        +(-20:\R) coordinate (C)
        +(-60:\R) coordinate (D)
        (intersection of A--C and B--D) coordinate (J)
        (intersection of A--B and C--D) coordinate (I)
        ;
        \draw (O) circle (\R)
        (A)--(B)--(C)--(D)--cycle
        (A)--(C)
        (B)--(D)
        ;
      
          \filldraw[black] (O) circle (1pt) node[above right] {$O$};
    \filldraw[black] (A) circle (1pt) node[above] {$A$};
    \filldraw[black] (B) circle (1pt) node[above right] {$B$};
    \filldraw[black] (C) circle (1pt) node[below right] {$C$};
   \filldraw[black] (D) circle (1pt) node[below right]{$E$};
    \filldraw[black] (J) circle (1pt) node[below left] {$H$};
    
    \end{tikzpicture}
}
\end{minipage}

\end{tcolorbox}
\vspace{4pt}


\begin{tcolorbox}[
    breakable,
    colback=white,
    colframe=blue!60!black,
    boxrule=0.8pt,
    arc=3pt,
    title={\textbf{Câu \#434} \quad \textbf{MÃ CÂU: \texttt{[L09\_C9\_B3\_05]}} \quad | \quad \textbf{Nhóm: \texttt{HINH\_329}} \quad | \quad \textbf{MapID: \texttt{9H9V3-1}} \hfill \textcolor{yellow!90!white}{[Hình gốc .pyc]}}
]
\textbf{Dạng toán:} CM tứ giác nội tiếp \hfill \textit{(Tự luận -- Vận dụng)}
\vspace{4pt}\hrule\vspace{6pt}
\begin{minipage}[t]{0.58\textwidth}
\vspace{0pt}
Cho hình vuông ${BHKL}$ có cạnh bằng $38$ cm. Gọi ${Q}$, ${S}$, ${T}$, ${M}$ lần lượt là trung điểm của ${BH}$, ${HK}$, ${KL}$, ${LB}$. Chứng minh rằng tứ giác ${QSTM}$ nội tiếp một đường tròn và tìm bán kính, chu vi của đường tròn đó.
\end{minipage}\hfill
\begin{minipage}[t]{0.39\textwidth}
\vspace{0pt}
\centering

\resizebox{\linewidth}{!}{%
\begin{tikzpicture}[scale=1, font=\footnotesize, line join=round, line cap=round, >=stealth]
      
    \tikzset{
        markx/.pic={\draw[thin] (45:.1)--(-135:.1) (-45:.1)--(135:.1);},
        markl/.pic={\draw[thin] (90:.1)--(-90:.1);},
        markll/.pic={
            \draw[thin, shift={(180:.02)}] (90:.1)--(-90:.1);
            \draw[thin, shift={(0:.02)}] (90:.1)--(-90:.1);
        },
        marklll/.pic={
            \draw[thin, shift={(180:.03)}] (90:.1)--(-90:.1);
            \draw[thin] (90:.1)--(-90:.1);
            \draw[thin, shift={(0:.03)}] (90:.1)--(-90:.1);
        }
    }
    \def\canh{4}
    \path 
    (0:0) coordinate (A)
    ++(0:\canh) coordinate (B)
    ++(-90:\canh) coordinate (C)
    ($(A)+(C)-(B)$) coordinate (D)
    ($(A)!.5!(B)$) coordinate (M)
    ($(C)!.5!(B)$) coordinate (N)
    ($(C)!.5!(D)$) coordinate (P)
    ($(A)!.5!(D)$) coordinate (Q)
    (A)--(M) pic[pos = .5, sloped]{markl}
    (B)--(M) pic[pos = .5, sloped]{markl}
    (C)--(N) pic[pos = .5, sloped]{markl}
    (C)--(P) pic[pos = .5, sloped]{markl}
    (D)--(P) pic[pos = .5, sloped]{markl}
    (Q)--(D) pic[pos = .5, sloped]{markl}
    (A)--(Q) pic[pos = .5, sloped]{markl}
    ;
    \draw (A)--(B)--(C)--(D)--cycle
    (M)--(N)--(P)--(Q)--cycle
    \foreach \x/\y/\z in {D/A/B,A/B/C,C/D/A,B/C/D}{
        pic[draw, thin, angle radius = 6pt]{right angle = \x--\y--\z}
    } ;
    
          \filldraw[black] (Q) circle (1pt) node[above left] {$M$};
    \filldraw[black] (A) circle (1pt) node[above left] {$B$};
    \filldraw[black] (B) circle (1pt) node[above right] {$H$};
    \filldraw[black] (C) circle (1pt) node[below right] {$K$};
   \filldraw[black] (D) circle (1pt) node[below left]{$L$};
    \filldraw[black] (M) circle (1pt) node[above] {$Q$};
       \filldraw[black] (N) circle (1pt) node[below right]{$S$};
    \filldraw[black] (P) circle (1pt) node[below] {$T$};
    \end{tikzpicture}
}
\end{minipage}

\end{tcolorbox}
\vspace{4pt}


\begin{tcolorbox}[
    breakable,
    colback=white,
    colframe=blue!60!black,
    boxrule=0.8pt,
    arc=3pt,
    title={\textbf{Câu \#435} \quad \textbf{MÃ CÂU: \texttt{[L09\_C9\_B3\_06]}} \quad | \quad \textbf{Nhóm: \texttt{HINH\_330}} \quad | \quad \textbf{MapID: \texttt{9H9V3-1}} \hfill \textcolor{yellow!90!white}{[Hình gốc .pyc]}}
]
\textbf{Dạng toán:} Tính diện tích hình viên phân \hfill \textit{(Tự luận -- Vận dụng)}
\vspace{4pt}\hrule\vspace{6pt}
\begin{minipage}[t]{0.58\textwidth}
\vspace{0pt}
Cho hình vuông ${ADEG}$ có cạnh bằng $18$ cm và nội tiếp đường tròn $({I})$ như Hình vẽ. Tính tổng diện tích của bốn hình viên phân được giới hạn bởi các cạnh hình vuông (phần tô đậm trong hình).
\end{minipage}\hfill
\begin{minipage}[t]{0.39\textwidth}
\vspace{0pt}
\centering

\resizebox{\linewidth}{!}{%
\begin{tikzpicture}[scale=1, font=\footnotesize, line join=round, line cap=round, >=stealth]
    \def \R{2}
    \path 
    (0:0) coordinate (O)
    +(135:\R) coordinate (A)
    +(45:\R) coordinate (B)
    +(-45:\R) coordinate (C)
    +(-135:\R) coordinate (D)
    ;
    \fill[gray,opacity=.5] (O) circle (\R);
    \fill[white] (A)--(B)--(C)--(D)--cycle;
    \draw (O) circle (\R)
    (A)--(B)--(C)--(D)--cycle
    (A)--(C)
    (B)--(D)
    
    ;
    \path
    (A)--(B) node[midway,below] {$18$ cm}
    ;
    \filldraw[black] (A) circle (1pt) node[above left] {$A$};
    \filldraw[black] (B) circle (1pt) node[above right] {$D$};
    \filldraw[black] (C) circle (1pt) node[below right] {$E$};
   \filldraw[black] (D) circle (1pt) node[below left]{$G$};
   
    \end{tikzpicture}
}
\end{minipage}

\end{tcolorbox}
\vspace{4pt}


\begin{tcolorbox}[
    breakable,
    colback=white,
    colframe=blue!60!black,
    boxrule=0.8pt,
    arc=3pt,
    title={\textbf{Câu \#436} \quad \textbf{MÃ CÂU: \texttt{[L09\_C9\_B3\_07]}} \quad | \quad \textbf{Nhóm: \texttt{HINH\_331}} \quad | \quad \textbf{MapID: \texttt{9H9V3-1}} \hfill \textcolor{yellow!90!white}{[Hình gốc .pyc]}}
]
\textbf{Dạng toán:} CM tứ giác nội tiếp \hfill \textit{(Tự luận -- Vận dụng)}
\vspace{4pt}\hrule\vspace{6pt}
\begin{minipage}[t]{0.58\textwidth}
\vspace{0pt}
Cho hình thoi ${BCEH}$ có $BE=12$ cm, $CH=4$ cm. Gọi ${Q}$, ${S}$, ${T}$, ${X}$ lần lượt là trung điểm của ${BC}$, ${CE}$, ${EH}$, ${HB}$. Chứng minh rằng tứ giác ${QSTX}$ nội tiếp một đường tròn và tìm bán kính của đường tròn đó.
\end{minipage}\hfill
\begin{minipage}[t]{0.39\textwidth}
\vspace{0pt}
\centering

\resizebox{\linewidth}{!}{%
\begin{tikzpicture}[scale=1, font=\footnotesize, line join=round, line cap=round, >=stealth]
   
    \tikzset{
        markx/.pic={\draw[thin] (45:.1)--(-135:.1) (-45:.1)--(135:.1);},
        markl/.pic={\draw[thin] (90:.1)--(-90:.1);},
        markll/.pic={
            \draw[thin, shift={(180:.02)}] (90:.1)--(-90:.1);
            \draw[thin, shift={(0:.02)}] (90:.1)--(-90:.1);
        },
        marklll/.pic={
            \draw[thin, shift={(180:.03)}] (90:.1)--(-90:.1);
            \draw[thin] (90:.1)--(-90:.1);
            \draw[thin, shift={(0:.03)}] (90:.1)--(-90:.1);
        }
    }
    \def\canh{4}
    \path 
    (0:0) coordinate (A)
    ++(30:\canh) coordinate (B)
    ++(-30:\canh) coordinate (C)
    ($(A)+(C)-(B)$) coordinate (D)
    
    
    
    ($(A)!.5!(B)$) coordinate (M)
    ($(C)!.5!(B)$) coordinate (N)
    ($(C)!.5!(D)$) coordinate (P)
    ($(A)!.5!(D)$) coordinate (Q)
    ($(A)!.5!(C)$) coordinate (O)
    
    (A)--(M) pic[pos = .5, sloped]{markl}
    (B)--(M) pic[pos = .5, sloped]{markl}
    (C)--(N) pic[pos = .5, sloped]{markl}
    (C)--(P) pic[pos = .5, sloped]{markl}
    (D)--(P) pic[pos = .5, sloped]{markl}
    (Q)--(D) pic[pos = .5, sloped]{markl}
    (A)--(Q) pic[pos = .5, sloped]{markl}
    
    ;
    \draw (A)--(B)--(C)--(D)--cycle
    (M)--(N)--(P)--(Q)--cycle
    (A)--(C)
    (B)--(D)
    (M)--(P)
    (Q)--(N)
    
    ;
\filldraw[black] (O) circle (1pt) node[above left] {$O$};
    \filldraw[black] (A) circle (1pt) node[above left] {$B$};
    \filldraw[black] (B) circle (1pt) node[above right] {$C$};
    \filldraw[black] (C) circle (1pt) node[below right] {$E$};
   \filldraw[black] (D) circle (1pt) node[below left]{$H$};
       \filldraw[black] (M) circle (1pt) node[above left] {$Q$};
    \filldraw[black] (N) circle (1pt) node[above right] {$S$};
    \filldraw[black] (P) circle (1pt) node[below right] {$T$};
   \filldraw[black] (Q) circle (1pt) node[below left]{$X$};
    \end{tikzpicture}
}
\end{minipage}

\end{tcolorbox}
\vspace{4pt}


\begin{tcolorbox}[
    breakable,
    colback=white,
    colframe=blue!60!black,
    boxrule=0.8pt,
    arc=3pt,
    title={\textbf{Câu \#437} \quad \textbf{MÃ CÂU: \texttt{[L09\_C9\_B3\_08]}} \quad | \quad \textbf{Nhóm: \texttt{HINH\_332}} \quad | \quad \textbf{MapID: \texttt{9H9V3-1}} \hfill \textcolor{yellow!90!white}{[Hình gốc .pyc]}}
]
\textbf{Dạng toán:} Bài toán về Hình vuông, HCN nội tiếp \hfill \textit{(Tự luận -- Vận dụng)}
\vspace{4pt}\hrule\vspace{6pt}
\begin{minipage}[t]{0.58\textwidth}
\vspace{0pt}
Cho hình chữ nhật ${ACDE}$ nội tiếp đường tròn $({I})$ với $AC=30$ cm; $AE=72$ cm. Vẽ một hình vuông nội tiếp $({I})$. Tính diện tích của hình vuông đó.
\end{minipage}\hfill
\begin{minipage}[t]{0.39\textwidth}
\vspace{0pt}
\centering

\resizebox{\linewidth}{!}{%
\begin{tikzpicture}[scale=1, font=\footnotesize, line join=round, line cap=round, >=stealth]
    \def \R{2}
    \path 
    (0:0) coordinate (O)
    +(-150:\R) coordinate (A)
    +(-30:\R) coordinate (D)
    +(30:\R) coordinate (C)
    +(150:\R) coordinate (B)
    +(0:\R) coordinate (1)
    +(90:\R) coordinate (2)
    +(180:\R) coordinate (3)
    +(270:\R) coordinate (4)
    (B)--(A) node[midway,right] {$30$ cm}
    (D)--(A) node[midway,above] {$72$ cm}
    ;
    \draw (O) circle (\R)
    (A)--(B)--(C)--(D)--cycle
    (1)--(2)--(3)--(4)--cycle
    (B)--(D)
    
    ;

\filldraw[black] (O) circle (1pt) node[above left] {$I$};
    \filldraw[black] (A) circle (1pt) node[below left] {$A$};
    \filldraw[black] (B) circle (1pt) node[above left] {$C$};
    \filldraw[black] (C) circle (1pt) node[above right] {$D$};
   \filldraw[black] (D) circle (1pt) node[below right]{$E$};
   
    \end{tikzpicture}
}
\end{minipage}

\end{tcolorbox}
\vspace{4pt}


\begin{tcolorbox}[
    breakable,
    colback=white,
    colframe=blue!60!black,
    boxrule=0.8pt,
    arc=3pt,
    title={\textbf{Câu \#438} \quad \textbf{MÃ CÂU: \texttt{[L09\_C9\_B3\_09]}} \quad | \quad \textbf{Nhóm: \texttt{HINH\_333}} \quad | \quad \textbf{MapID: \texttt{9H9V3-1}} \hfill \textcolor{yellow!90!white}{[Hình gốc .pyc]}}
]
\textbf{Dạng toán:} Bài toán về Hình vuông, HCN nội tiếp \hfill \textit{(Tự luận -- Vận dụng)}
\vspace{4pt}\hrule\vspace{6pt}
\begin{minipage}[t]{0.58\textwidth}
\vspace{0pt}
Cho tam giác nhọn ${DEG}$ có các đường cao ${EQ}$, ${GS}$. Một đường tròn $({I})$ đi qua hai điểm ${Q}$, ${S}$ và cắt các tia đối của hai tia ${ES}$, ${GQ}$ lần lượt tại ${V}$ và ${W}$. Chứng minh rằng ${VW}$ song song với ${EG}$.
\end{minipage}\hfill
\begin{minipage}[t]{0.39\textwidth}
\vspace{0pt}
\centering

\resizebox{\linewidth}{!}{%
\begin{tikzpicture}[scale=1, font=\footnotesize, line join=round, line cap=round, >=stealth]
        \def \R{2}
        \path 
        (0:0) coordinate (O)
        +(-10:\R) coordinate (Y)
        +(-170:\R) coordinate (X)
        +(90:\R) coordinate (E)
        +(120:\R) coordinate (F)
        (intersection of X--F and Y--E) coordinate (A)
        ($(E)!2mm!90:(A)$) coordinate (1)
        ($(F)!2mm!-90:(A)$) coordinate (2)
        (intersection of X--A and 1--E) coordinate (B)
        (intersection of Y--A and 2--F) coordinate (C)
        ;
        \draw (O) circle (\R)
        (X)--(A)--(Y)--cycle
        (B)--(C)
        (B)--(E)--(F)--(C)
        \foreach \x/\y/\z in {B/F/C,B/E/C}{
            pic[draw, thin, angle radius = 6pt]{right angle = \x--\y--\z}
        }
        ;
    \filldraw[black] (A) circle (1pt) node[above] {$D$};
    \filldraw[black] (B) circle (1pt) node[above left] {$E$};
    \filldraw[black] (C) circle (1pt) node[above right] {$G$};
   \filldraw[black] (E) circle (1pt) node[above right]{$Q$};
     \filldraw[black] (F) circle (1pt) node[above left] {$S$};
    \filldraw[black] (X) circle (1pt) node[above left] {$V$};
    \filldraw[black] (Y) circle (1pt) node[above right] {$W$};

    \end{tikzpicture}
}
\end{minipage}

\end{tcolorbox}
\vspace{4pt}


\begin{tcolorbox}[
    breakable,
    colback=white,
    colframe=blue!60!black,
    boxrule=0.8pt,
    arc=3pt,
    title={\textbf{Câu \#439} \quad \textbf{MÃ CÂU: \texttt{[L09\_C9\_B4\_01]}} \quad | \quad \textbf{Nhóm: \texttt{HINH\_334}} \quad | \quad \textbf{MapID: \texttt{9H9V4-1}} \hfill \textcolor{yellow!90!white}{[Hình gốc .pyc]}}
]
\textbf{Dạng toán:} Tính chu vi và diện tích của lục giác đều biết nó cùng nội tiếp với tam giác đều cạnh a \hfill \textit{(Trắc nghiệm -- Vận dụng)}
\vspace{4pt}\hrule\vspace{6pt}
\begin{minipage}[t]{0.58\textwidth}
\vspace{0pt}
Cho một lục giác đều và một tam giác đều cùng nội tiếp một đường tròn. Biết rằng tam giác đều có cạnh bằng $3$ cm. Tính chu vi và diện tích của hình lục giác đều đã cho.

\vspace{4pt}\noindent
\textbf{A.} Chu vi là $4 \sqrt{3} \left( \mathrm{cm}\right)$, diện tích là $9 \sqrt{3}\left( \mathrm{cm^{2}}\right) $ \quad \textbf{B.} Chu vi là $2 \sqrt{3} \left( \mathrm{cm}\right)$, diện tích là $\dfrac{9 \sqrt{3}}{2}\left( \mathrm{cm^{2}}\right) $ \quad \textbf{C.} Chu vi là $\sqrt{3} \left( \mathrm{cm}\right)$, diện tích là $\dfrac{9 \sqrt{3}}{4}\left( \mathrm{cm^{2}}\right) $ \quad \textbf{\textcolor{blue!75!black}{D. (Đ)}} Chu vi là $6 \sqrt{3} \left( \mathrm{cm}\right)$, diện tích là $\dfrac{27 \sqrt{3}}{2}\left( \mathrm{cm^{2}}\right) $
\end{minipage}\hfill
\begin{minipage}[t]{0.39\textwidth}
\vspace{0pt}
\centering

\resizebox{\linewidth}{!}{%
\begin{tikzpicture}[scale=1, font=\footnotesize,>=stealth]
                %Gán số liệu.
                \def\canhAB{4};
                \pgfmathsetmacro\r{(\canhAB*sqrt(3))/3};
                %Gán tọa độ.
                \coordinate (A) at (90:\r);
                \coordinate (B) at (210:\r);
                \coordinate (C) at (-30:\r);
                \coordinate (O) at ($1/3*(A)+1/3*(B)+1/3*(C)$);
                \draw  let \p1=($(A)-(O)$) in (O) circle({veclen(\x1,\y1)});
                %Vẽ tam giác ABC.
                \draw (A)--(B)--(C)--cycle;
                \draw[dashed, blue] (O)--(A);
                %Gán nhãn.
                %\foreach \x/\y in {A/90,B/210,C/-30, O/90}{\fill (\x) circle(1pt) ($(\x)+(\y:0.3cm)$) node{$\x$};}
                % Vẽ lục giác đều
                \foreach \tendinh/\goc in {A/0,B/60,C/120,D/180,E/240,F/300}{
                    \coordinate (\tendinh) at (\goc:\r);
                    \fill (\tendinh) circle(1pt) ($(\tendinh)+(\goc:0.3cm)$);
                }
                %Vẽ đa giác đều.
                \draw (A)--(B)--(C)--(D)--(E)--(F)--cycle;
                \draw[dashed, red] (O)--(A)(O)--(B)(O)--(C)(O)--(D)(O)--(E)(O)--(F);
            \end{tikzpicture}
}
\end{minipage}

\vspace{3pt}\hrule\vspace{3pt}\footnotesize\textcolor{gray!80!black}{\textbf{Các mã dùng chung hình (HINH\_334):} \texttt{L09\_C9\_B4\_02}}
\end{tcolorbox}
\vspace{4pt}


\begin{tcolorbox}[
    breakable,
    colback=white,
    colframe=blue!60!black,
    boxrule=0.8pt,
    arc=3pt,
    title={\textbf{Câu \#440} \quad \textbf{MÃ CÂU: \texttt{[L09\_C9\_B4\_02]}} \quad | \quad \textbf{Nhóm: \texttt{HINH\_334}} \quad | \quad \textbf{MapID: \texttt{9H9V4-1}} \hfill \textcolor{yellow!90!white}{[Hình gốc .pyc]}}
]
\textbf{Dạng toán:} Tính chu vi và diện tích của lục giác đều biết nó cùng nội tiếp với tam giác đều cạnh a \hfill \textit{(Tự luận -- Vận dụng)}
\vspace{4pt}\hrule\vspace{6pt}
\begin{minipage}[t]{0.58\textwidth}
\vspace{0pt}
Cho một lục giác đều và một tam giác đều cùng nội tiếp một đường tròn. Biết rằng tam giác đều có cạnh bằng $20$ cm. Tính chu vi và diện tích của hình lục giác đều đã cho.
\end{minipage}\hfill
\begin{minipage}[t]{0.39\textwidth}
\vspace{0pt}
\centering

\resizebox{\linewidth}{!}{%
\begin{tikzpicture}[scale=1, font=\footnotesize,>=stealth]
                %Gán số liệu.
                \def\canhAB{4};
                \pgfmathsetmacro\r{(\canhAB*sqrt(3))/3};
                %Gán tọa độ.
                \coordinate (A) at (90:\r);
                \coordinate (B) at (210:\r);
                \coordinate (C) at (-30:\r);
                \coordinate (O) at ($1/3*(A)+1/3*(B)+1/3*(C)$);
                \draw  let \p1=($(A)-(O)$) in (O) circle({veclen(\x1,\y1)});
                %Vẽ tam giác ABC.
                \draw (A)--(B)--(C)--cycle;
                \draw[dashed, blue] (O)--(A);
                %Gán nhãn.
                %\foreach \x/\y in {A/90,B/210,C/-30, O/90}{\fill (\x) circle(1pt) ($(\x)+(\y:0.3cm)$) node{$\x$};}
                % Vẽ lục giác đều
                \foreach \tendinh/\goc in {A/0,B/60,C/120,D/180,E/240,F/300}{
                    \coordinate (\tendinh) at (\goc:\r);
                    \fill (\tendinh) circle(1pt) ($(\tendinh)+(\goc:0.3cm)$);
                }
                %Vẽ đa giác đều.
                \draw (A)--(B)--(C)--(D)--(E)--(F)--cycle;
                \draw[dashed, red] (O)--(A)(O)--(B)(O)--(C)(O)--(D)(O)--(E)(O)--(F);
            \end{tikzpicture}
}
\end{minipage}

\vspace{3pt}\hrule\vspace{3pt}\footnotesize\textcolor{gray!80!black}{\textbf{Các mã dùng chung hình (HINH\_334):} \texttt{L09\_C9\_B4\_01}}
\end{tcolorbox}
\vspace{4pt}


\begin{tcolorbox}[
    breakable,
    colback=white,
    colframe=blue!60!black,
    boxrule=0.8pt,
    arc=3pt,
    title={\textbf{Câu \#441} \quad \textbf{MÃ CÂU: \texttt{[L09\_C9\_B5\_01]}} \quad | \quad \textbf{Nhóm: \texttt{HINH\_335}} \quad | \quad \textbf{MapID: \texttt{9H9V3-1}} \hfill \textcolor{yellow!90!white}{[Hình gốc .pyc]}}
]
\textbf{Dạng toán:} Tổng hợp 1 \hfill \textit{(Tự luận -- Vận dụng)}
\vspace{4pt}\hrule\vspace{6pt}
\begin{minipage}[t]{0.58\textwidth}
\vspace{0pt}
Cho nửa đường tròn ${(O)}$ đường kính ${CG}$. Điểm ${U}$ nằm trên nửa đường tròn ${(U \neq C, G)}$. Tiếp tuyến tại ${U}$ cắt tiếp tuyến tại ${C}$ và ${G}$ của đường tròn ${(O)}$ lần lượt tại ${J}$ và ${K}$. 

a) Chứng minh rằng: tứ giác ${CJUO}$ nội tiếp.

b) Gọi ${W}$ là giao điểm ${JK}$ và ${CG}$. Chứng minh: ${WC \cdot WO = WJ \cdot WU}$.

c) Gọi ${L}$ là giao điểm của ${CU}$ và ${GK}$; ${P}$ là giao điểm của ${CJ}$ và ${GU}$. Chứng minh: Ba điểm ${L}$, ${P}$, ${W}$ thẳng hàng.
\end{minipage}\hfill
\begin{minipage}[t]{0.39\textwidth}
\vspace{0pt}
\centering

\resizebox{\linewidth}{!}{%
\begin{tikzpicture}[line cap=round,line join=round,font=\footnotesize]
        \def\r{3}
        \path 
        (0,0) coordinate(A)
        ++(0:2*\r) coordinate (B)
        ($(A)!0.5!(B)$) coordinate (O)
        (O)++(120:\r) coordinate (M)
        ($(M)!1!90:(O)$) coordinate (Mx)
        (B)++(90:1) coordinate (Bx)
        (intersection of B--Bx and M--Mx) coordinate (D)
        (intersection of A--B and M--Mx) coordinate (P)
        (A)++(90:1) coordinate (Ax)
        (intersection of A--Ax and B--M) coordinate (F)
        (intersection of M--P and A--F) coordinate (C)
        (intersection of B--D and A--M) coordinate (E)
        ;
        \draw (O) circle(\r) (A)--(B)(O)--(M) (M)--(A) (M)--(B)(A)--(P)--(D)--(B)(A)--(F)--(M)--(E)--(D);
          \foreach \x/\y/\z in {F/A/B,O/M/D,D/B/A}{\draw pic[draw,angle radius=3mm]{right angle=\x--\y--\z}; }
       
    \filldraw[black] (A) circle (1pt) node[below left] {$C$};
    \filldraw[black] (B) circle (1pt) node[above right] {$G$};
    \filldraw[black] (C) circle (1pt) node[above left] {$J$};
   \filldraw[black] (D) circle (1pt) node[above right]{$K$};
    \filldraw[black] (E) circle (1pt) node[above] {$L$};
    \filldraw[black] (F) circle (1pt) node[above] {$P$};
    \filldraw[black] (M) circle (1pt) node[above] {$U$};
   \filldraw[black] (P) circle (1pt) node[left]{$W$};
    \filldraw[black] (O) circle (1pt) node[below ] {$O$};
    \end{tikzpicture}
}
\end{minipage}

\end{tcolorbox}
\vspace{4pt}


\begin{tcolorbox}[
    breakable,
    colback=white,
    colframe=blue!60!black,
    boxrule=0.8pt,
    arc=3pt,
    title={\textbf{Câu \#442} \quad \textbf{MÃ CÂU: \texttt{[L09\_C9\_B5\_02]}} \quad | \quad \textbf{Nhóm: \texttt{HINH\_336}} \quad | \quad \textbf{MapID: \texttt{9H9V3-1}} \hfill \textcolor{yellow!90!white}{[Hình gốc .pyc]}}
]
\textbf{Dạng toán:} Tổng hợp 2 \hfill \textit{(Tự luận -- Vận dụng)}
\vspace{4pt}\hrule\vspace{6pt}
\begin{minipage}[t]{0.58\textwidth}
\vspace{0pt}
Cho đường tròn ${(O)}$ đường kính ${BE}$. Điểm ${A}$ thuộc đường tròn thỏa mãn ${AB < AE}$ (${A}$ khác ${B}$, ${A}$ khác ${E}$). Lấy điểm ${F}$ trên cạnh ${AE}$ sao cho ${AF = AB}$. Vẽ hình vuông ${BAFG}$. Tia ${AG}$ cắt $({O})$ tại ${H}$.

a) Tam giác ${HBE}$ là tam giác gì? Tại sao?

b) Chứng minh: $\widehat{HFE} = \widehat{HEF}$.

c) Vẽ tia ${B}x$ là tiếp tuyến của $({O})$ tại ${B}$, ${B}x$ cắt ${EH}$ tại ${L}$. Chứng minh ba điểm ${F}$, ${G}$, ${L}$ thẳng hàng.
\end{minipage}\hfill
\begin{minipage}[t]{0.39\textwidth}
\vspace{0pt}
\centering

\resizebox{\linewidth}{!}{%
\begin{tikzpicture}[line cap=round,line join=round,font=\footnotesize] 
  \def\r{3}
    \path
    (0,0) coordinate (B)
    ++(0:2*\r) coordinate (C)
    ($(B)!0.5!(C)$) coordinate (O)
    (O)++(120:\r) coordinate (A)
 ($(B)!1!-90:(A)$) coordinate (Bx) 
 (B)++(90:2cm) coordinate (By) ;
 \path[name path = dtron1] let \p1=($(A) - (B)$) in (A) circle ({veclen(\x1,\y1)}); 
 \path[name path=dthang] (A)--(C); 
 \path[name path=dtron] let \p1=($ (A) - (B) $) in (A) circle ({veclen(\x1,\y1)}); 
 \path[name intersections={of=dtron and dthang, by={D}}];
  \path ($(D)!1!90:(A)$) coordinate (Dx) (intersection of B--Bx and D--Dx) coordinate (E) ($(E)!(O)!(A)$) coordinate (E1) ($(A)!2!(E1)$) coordinate (F) (intersection of C--F and E--D) coordinate (I) ; 
 \draw (A)--(B)--(C)--cycle (B)--(E)--(D)(B)--(I)--(E)(B)--(F)--(I)--cycle (D)--(F)--(C)(A)--(F)(B)--(By); 
 \draw let \p2=($(O) - (A)$) in (O) circle ({veclen(\x2,\y2)}); 

    \filldraw[black] (A) circle (1pt) node[above] {$A$};
    \filldraw[black] (B) circle (1pt) node[left] {$B$};
    \filldraw[black] (C) circle (1pt) node[above right] {$E$};
   \filldraw[black] (D) circle (1pt) node[above right]{$F$};
    \filldraw[black] (E) circle (1pt) node[below left] {$G$};
    \filldraw[black] (F) circle (1pt) node[below] {$H$};
     \filldraw[black] (I) circle (1pt) node[below] {$L$};
    \filldraw[black] (O) circle (1pt) node[below ] {$O$};
  \end{tikzpicture}
}
\end{minipage}

\end{tcolorbox}
\vspace{4pt}


\begin{tcolorbox}[
    breakable,
    colback=white,
    colframe=blue!60!black,
    boxrule=0.8pt,
    arc=3pt,
    title={\textbf{Câu \#443} \quad \textbf{MÃ CÂU: \texttt{[L09\_C9\_B5\_03]}} \quad | \quad \textbf{Nhóm: \texttt{HINH\_337}} \quad | \quad \textbf{MapID: \texttt{9H9V3-1}} \hfill \textcolor{yellow!90!white}{[Hình gốc .pyc]}}
]
\textbf{Dạng toán:} Tổng hợp 3 \hfill \textit{(Tự luận -- Vận dụng)}
\vspace{4pt}\hrule\vspace{6pt}
\begin{minipage}[t]{0.58\textwidth}
\vspace{0pt}
Cho đường tròn tâm ${O}$ , bán kính R  và một đường thẳng d không cắt đường tròn tâm ${O}$. Dựng đường thẳng ${OV}$ vuông góc với đường thẳng d tại điểm ${V}$. Trên đường thẳng d lấy điểm ${N}$ (khác điểm ${V}$), qua ${N}$ vẽ hai tiếp tuyến ${NT}$ và ${NS}$ với đường tròn $({O})$ , ( ${T}$ và ${S}$ là các tiếp điểm) sao cho ${T}$ và ${V}$ nằm về hai phía  của đường thẳng ${ON}$.

a) Chứng minh tứ giác ${NTOV}$ nội tiếp được trong đường tròn.

b) Đường thẳng ${TS}$ cắt đường thẳng ${OV}$ tại điểm ${D}$. Chứng minh rằng ${DT \cdot DS=DV\cdot DO}$.

c) Khi ${ON=2R; OV=R\sqrt{3} }$. Tính diện tích $\Delta {NTD}$ theo R.
\end{minipage}\hfill
\begin{minipage}[t]{0.39\textwidth}
\vspace{0pt}
\centering

\resizebox{\linewidth}{!}{%
\begin{tikzpicture}[line cap=round,line join=round,font=\footnotesize] 
        \def\r{3}
        \path 
        (0,0) coordinate (O)
        (O)++(150:2.5*\r) coordinate (K)
        (O)++(30:\r) coordinate (O1)
        (K)++(20:3*\r) coordinate (Mx)
        ;
        \begin{scope}[overlay]
            \path ($(K)!0.5!(O)$) coordinate (K1);
            \path[name path=dtronK1] let\p2=($(K1)-(K)$) in (K1) circle({veclen(\x2,\y2)});
            \path[name path=dtronT] let \p1=($(O)-(O1)$) in (O) circle ({veclen(\x1,\y1)});
            \path[name intersections={of=dtronK1 and dtronT, by={B,A}}];
            \draw (A)--(K)--(B);
            \draw (O) circle(\r) (K)--(O) (O)--(A) (O)--(B)--(A);
        \end{scope}
        \path 
        ($(K)!(O)!(Mx)$) coordinate (H)
        (intersection of O--H and A--B) coordinate (I)
        (intersection of O--K and A--B) coordinate (M)
        ;
        \draw (O)--(H)(K)--(Mx)node[above]{$d$} (A)--(H)(K)--(I);
        \foreach \x/\y/\z in {K/B/O,K/A/O,O/H/K}{\draw pic[draw,angle radius=3mm]{right angle=\x--\y--\z};}
  
    \filldraw[black] (I) circle (1pt) node[above right] {$D$};
    \filldraw[black] (M) circle (1pt) node[below right] {$E$};
    \filldraw[black] (K) circle (1pt) node[above] {$N$};
   \filldraw[black] (B) circle (1pt) node[above right]{$S$};
    \filldraw[black] (A) circle (1pt) node[below left] {$T$};
     \filldraw[black] (H) circle (1pt) node[above] {$V$};
    \filldraw[black] (O) circle (1pt) node[below ] {$O$};
  \end{tikzpicture}
}
\end{minipage}

\end{tcolorbox}
\vspace{4pt}


\begin{tcolorbox}[
    breakable,
    colback=white,
    colframe=blue!60!black,
    boxrule=0.8pt,
    arc=3pt,
    title={\textbf{Câu \#444} \quad \textbf{MÃ CÂU: \texttt{[L09\_C9\_B5\_04]}} \quad | \quad \textbf{Nhóm: \texttt{HINH\_338}} \quad | \quad \textbf{MapID: \texttt{9H9V3-1}} \hfill \textcolor{yellow!90!white}{[Hình gốc .pyc]}}
]
\textbf{Dạng toán:} Tổng hợp 4 \hfill \textit{(Tự luận -- Vận dụng)}
\vspace{4pt}\hrule\vspace{6pt}
\begin{minipage}[t]{0.58\textwidth}
\vspace{0pt}
Cho đường tròn ${(O; R)}$ và dây ${FR < 2R}$. Trên cung lớn ${FR}$ lấy điểm ${J}$ sao cho ${JF < JR}$. Các đường cao ${JW}$ và ${FT}$ của tam giác ${JFR}$ cắt nhau tại ${B}$.

a) Chứng minh tứ giác ${JFWT}$ nội tiếp đường tròn và xác định tâm của đường tròn đó.

b) Chứng minh: ${RW \cdot RF} = {RT \cdot RJ}$.

c) Đường tròn ngoại tiếp tam giác ${RWT}$ cắt ${(O; R)}$ tại điểm ${Q}$ ($Q \neq R$). Vẽ đường kính ${RE}$ của ${(O; R)}$ và gọi ${E}$ là trung điểm của ${JF}$. Chứng minh 3 điểm ${E}$, ${E}$, ${Q}$ thẳng hàng.
\end{minipage}\hfill
\begin{minipage}[t]{0.39\textwidth}
\vspace{0pt}
\centering

\resizebox{\linewidth}{!}{%
\begin{tikzpicture}[line cap=round,line join=round,font=\footnotesize] 
    \def\canh{4}
    \path 
    (0,0) coordinate (O)
    ++(-30:\canh) coordinate (C)
    ++(0:1) coordinate (C1)
    ($(C)!(O)!(C1)$) coordinate (C2)
    ($(C)!2!(C2)$) coordinate (B)
    (O)++(120:\canh) coordinate (A)
    ($(B)!(A)!(C)$) coordinate (D)
    ($(A)!(B)!(C)$) coordinate (F)
    (intersection of A--D and B--F) coordinate (I)
    ($(C)!2!(O)$) coordinate (K)
    ;
    \path ($(D)!(C)!(F)$) coordinate (CH)
    ($(C)!(D)!(F)$) coordinate (DH)
    (intersection of C--CH and D--DH) coordinate (Ht)
    (barycentric cs:C=1,D=1,F=1) coordinate (Gt)
    (barycentric cs:Gt=3,Ht=-1) coordinate (J);
    \draw[name path = dtronJ] let \p1=($(J) - (C)$) in (J) circle ({veclen(\x1,\y1)});
    \path[name path=dtron1] let \p1=($ (O) - (B) $) in (O) circle ({veclen(\x1,\y1)});
    \path[name path=dtron2] let \p1=($ (J) - (C) $) in (J) circle ({veclen(\x1,\y1)});
    \path[name intersections={of=dtron1 and dtron2, by={H, H1}}];
    \path 
    ($(C)!2!(J)$) coordinate (I)
    (intersection of H--I and B--A) coordinate (E)
    ;
    \draw[name path=dtron1] let \p1=($(O)-(B)$) in (O) circle ({veclen(\x1,\y1)})
    (B)--(K)--(A)--(C)--cycle (A)--(B)--(D)--(F)--cycle (K)--(C) (K)--(H) (K)--(D)--(A) (E)--(D) (E)--(F)--(J)--(D) (C)--(I)
    ;
    
    \filldraw[black] (C) circle (1pt) node[above right] {$R$};
    \filldraw[black] (E) circle (1pt) node[below right] {$E$};
    \filldraw[black] (K) circle (1pt) node[above] {$E$};
   \filldraw[black] (B) circle (1pt) node[above left]{$F$};
    \filldraw[black] (A) circle (1pt) node[above left] {$J$};
     \filldraw[black] (H) circle (1pt) node[below] {$Q$};
    \filldraw[black] (O) circle (1pt) node[below ] {$O$};
  \filldraw[black] (I) circle (1pt) node[below left] {$B$};
     \filldraw[black] (F) circle (1pt) node[above] {$T$};
    \filldraw[black] (D) circle (1pt) node[below left] {$W$};
  \filldraw[black] (J) circle (1pt) node[below left] {$U$};
  
  \end{tikzpicture}
}
\end{minipage}

\end{tcolorbox}
\vspace{4pt}


\begin{tcolorbox}[
    breakable,
    colback=white,
    colframe=blue!60!black,
    boxrule=0.8pt,
    arc=3pt,
    title={\textbf{Câu \#445} \quad \textbf{MÃ CÂU: \texttt{[L09\_C9\_B5\_05]}} \quad | \quad \textbf{Nhóm: \texttt{HINH\_339}} \quad | \quad \textbf{MapID: \texttt{9H9V3-1}} \hfill \textcolor{yellow!90!white}{[Hình gốc .pyc]}}
]
\textbf{Dạng toán:} Tổng hợp 5 \hfill \textit{(Tự luận -- Vận dụng)}
\vspace{4pt}\hrule\vspace{6pt}
\begin{minipage}[t]{0.58\textwidth}
\vspace{0pt}
Cho $\triangle {ACE}$ có 3 góc nhọn và đường cao ${CG}$. Gọi ${Y, W}$ lần lượt là chân đường vuông góc kẻ từ điểm ${G}$ đến ${AC, CE}$.

a) Chứng minh tứ giác ${CYGW}$ nội tiếp;

b) Chứng minh: ${CY \cdot CA = CW \cdot CE}$;

c) Gọi ${M}$ là chân đường vuông góc kẻ từ điểm ${E}$ đến đường thẳng ${AC}$, ${I}$ là trung điểm của đoạn thẳng ${GM}$. Chứng minh rằng ${Y, I, W}$ thẳng hàng.
\end{minipage}\hfill
\begin{minipage}[t]{0.39\textwidth}
\vspace{0pt}
\centering

\resizebox{\linewidth}{!}{%
\begin{tikzpicture}[declare function={a = 8; goc = 65; b = 0.65*a;},font=\footnotesize,line cap=round,line join=round]
    \path
    (goc:b) coordinate (A)
    (0,0) coordinate (B)
    (a,0) coordinate (C)
    ($(A)!(B)!(C)$) coordinate (E)
    ($(A)!(E)!(B)$) coordinate (H)
    ($(B)!(E)!(C)$) coordinate (K)
    ($(A)!(C)!(B)$) coordinate (F)
    ($(E)!.5!(F)$) coordinate (I)
    (intersection of H--I and C--F) coordinate (M)
    (intersection of E--K and C--F) coordinate (N)
    ;
    \path ($(H)!(E)!(K)$) coordinate (EH)
    ($(E)!(H)!(K)$) coordinate (HH)
    (intersection of E--EH and H--HH) coordinate (Ht)
    (barycentric cs:E=1,H=1,K=1) coordinate (Gt)
    (barycentric cs:Gt=3,Ht=-1) coordinate (J);
    \draw[name path = dtron2] let \p1=($(J) - (E)$) in (J) circle ({veclen(\x1,\y1)});
    \draw (A)--(B)--(C)--cycle (B)--(E)--(H) (C)--(F)(H)--(K)(E)--(F) (E)--(M) (E)--(N)--(K);
   \filldraw[black] (C) circle (1pt) node[below right] {$E$};
    \filldraw[black] (E) circle (1pt) node[above right] {$G$};
    \filldraw[black] (K) circle (1pt) node[below] {$W$};
   \filldraw[black] (B) circle (1pt) node[below left]{$C$};
    \filldraw[black] (A) circle (1pt) node[above left] {$A$};
     \filldraw[black] (H) circle (1pt) node[below left] {$Y$};
  \filldraw[black] (N) circle (1pt) node[below left] {$S$};
     \filldraw[black] (F) circle (1pt) node[left] {$M$};
  \filldraw[black] (M) circle (1pt) node[below left] {$U$};
    \foreach \x/\y/\z in {E/H/F,B/F/C,E/M/F,E/K/C}{\draw pic[draw,angle radius=2mm]{right angle=\x--\y--\z};}
\end{tikzpicture}
}
\end{minipage}

\end{tcolorbox}
\vspace{4pt}


\begin{tcolorbox}[
    breakable,
    colback=white,
    colframe=blue!60!black,
    boxrule=0.8pt,
    arc=3pt,
    title={\textbf{Câu \#446} \quad \textbf{MÃ CÂU: \texttt{[L09\_C9\_B5\_06]}} \quad | \quad \textbf{Nhóm: \texttt{HINH\_340}} \quad | \quad \textbf{MapID: \texttt{9H9V3-1}} \hfill \textcolor{yellow!90!white}{[Hình gốc .pyc]}}
]
\textbf{Dạng toán:} Tổng hợp 6 \hfill \textit{(Tự luận -- Vận dụng)}
\vspace{4pt}\hrule\vspace{6pt}
\begin{minipage}[t]{0.58\textwidth}
\vspace{0pt}
Cho tam giác ${ABC}$ nhọn (${AB<AC}$) nội tiếp đường tròn $(O)$. Hai đường cao ${BD}$ và ${CH}$ của tam giác ${ABC}$ cắt nhau tại điểm ${V}$. Gọi ${U}$ là trung điểm ${BC}$. 

a) Chứng minh $\triangle {ADH}$ đồng dạng $\triangle {ABC}$.

b) Chứng minh đường thẳng ${OA}$ vuông góc với đường thẳng ${DH}$.

c) Đường phân giác góc $\widehat {HVB}$ cắt ${AB}$ và ${AC}$ lần lượt tại ${P}$ và ${M}$. Gọi ${W}$ là trung điểm của ${PM}$, ${X}$ là trung điểm của ${AV}$. Chứng minh tứ giác ${AHVW}$ nội tiếp và ba điểm ${W, X, U}$ thẳng hàng.
\end{minipage}\hfill
\begin{minipage}[t]{0.39\textwidth}
\vspace{0pt}
\centering

\resizebox{\linewidth}{!}{%
\begin{tikzpicture}[declare function={a = 8; goc = 65; b = 0.65*a;},font=\footnotesize,line cap=round,line join=round]
    \path
    (goc:b) coordinate (A)
    (0,0) coordinate (B)
    (a,0) coordinate (C)
    ($(A)!(B)!(C)$) coordinate (E)
    ($(A)!(C)!(B)$) coordinate (F)
    ($(C)!(A)!(B)$) coordinate (D)
    (intersection of B--E and C--F) coordinate (H)
    ($(B)!.5!(C)$) coordinate (K)
    ($(A)!.5!(H)$) coordinate (J)
    ($(H)!1cm!(F)$) coordinate (HF)
    ($(H)!1cm!(B)$) coordinate (HB)
    ($(HF)!0.5!(HB)$) coordinate (Ht)
    (intersection of H--Ht and F--B) coordinate (M)
    (intersection of H--M and A--C) coordinate (N)
    (intersection of M--N and J--K) coordinate (I)
    ;
    \path ($(B)!(A)!(C)$) coordinate (AH)
    ($(A)!(B)!(C)$) coordinate (BH)
    (intersection of A--AH and B--BH) coordinate (Ht)
    (barycentric cs:A=1,B=1,C=1) coordinate (Gt)
    (barycentric cs:Gt=3,Ht=-1) coordinate (O);
    \draw (A)--(B)--(C)--cycle 
    let \p1=($(O) - (A)$) in (O) circle ({veclen(\x1,\y1)})
    (A)--(E)--(F)--cycle (A)--(D) (C)--(F) (B)--(E) (M)--(N) (A)--(O)
    ;
    \draw[dashed] (J)--(K);
      \filldraw[black] (J) circle (1pt) node[below left] {$X$};
     \filldraw[black] (C) circle (1pt) node[below right] {$C$};
    \filldraw[black] (E) circle (1pt) node[above right] {$D$};
    \filldraw[black] (K) circle (1pt) node[below] {$U$};
   \filldraw[black] (B) circle (1pt) node[below left]{$B$};
    \filldraw[black] (A) circle (1pt) node[above left] {$A$};
     \filldraw[black] (H) circle (1pt) node[below left] {$V$};
  \filldraw[black] (N) circle (1pt) node[above right] {$M$};
     \filldraw[black] (F) circle (1pt) node[left] {$H$};
  \filldraw[black] (M) circle (1pt) node[below left] {$P$};
  \filldraw[black] (I) circle (1pt) node[above right] {$W$};
  \filldraw[black] (O) circle (1pt) node[below right] {$O$};
    \foreach \x/\y/\z in {B/E/C,A/F/C,A/D/C}{\draw pic[draw,angle radius=2mm]{right angle=\x--\y--\z};}
\end{tikzpicture}
}
\end{minipage}

\end{tcolorbox}
\vspace{4pt}


\begin{tcolorbox}[
    breakable,
    colback=white,
    colframe=blue!60!black,
    boxrule=0.8pt,
    arc=3pt,
    title={\textbf{Câu \#447} \quad \textbf{MÃ CÂU: \texttt{[L09\_C9\_B5\_07]}} \quad | \quad \textbf{Nhóm: \texttt{HINH\_341}} \quad | \quad \textbf{MapID: \texttt{9H9V3-1}} \hfill \textcolor{yellow!90!white}{[Hình gốc .pyc]}}
]
\textbf{Dạng toán:} Tổng hợp 7 \hfill \textit{(Tự luận -- Vận dụng)}
\vspace{4pt}\hrule\vspace{6pt}
\begin{minipage}[t]{0.58\textwidth}
\vspace{0pt}
Cho tam giác ${CHI}$ nhọn $({CH < CI})$ có đường cao ${CK}$ và đường phân giác trong ${CY}$ (${K, Y}$ thuộc cạnh ${HI}$). Kẻ ${YS \perp CH}$ tại ${S}$, ${YP \perp CI}$ tại ${P}$.

a) Chứng minh bốn điểm ${Y, S, K, P}$ cùng nằm trên một đường tròn.

b) Chứng minh: $\widehat{HKS} = \widehat{YKP}$.

c) Qua ${Y}$ kẻ đường thẳng vuông góc với ${HI}$ cắt ${SP}$ tại ${W}$, ${CW}$ cắt ${HI}$ tại ${V}$. Chứng minh ${V}$ là trung điểm của ${HI}$.
\end{minipage}\hfill
\begin{minipage}[t]{0.39\textwidth}
\vspace{0pt}
\centering

\resizebox{\linewidth}{!}{%
\begin{tikzpicture}[declare function={a = 8; goc = 65; b = 0.65*a;},font=\footnotesize,line cap=round,line join=round]
    \path
    (goc:b) coordinate (A)
    (0,0) coordinate (B)
    (a,0) coordinate (C)
    ($(B)!(A)!(C)$) coordinate (D)
    ($(A)!1cm!(B)$) coordinate (AB)
    ($(A)!1cm!(C)$) coordinate (AC)
    ($(AB)!0.5!(AC)$) coordinate (At)
    (intersection of A--At and B--C) coordinate (O)
    ($(A)!(O)!(B)$) coordinate (M)
    ($(A)!(O)!(C)$) coordinate (N)
    ($(M)!(O)!(D)$) coordinate (OH)
    ($(O)!(M)!(D)$) coordinate (MH)
    (intersection of O--OH and M--MH) coordinate (Ht)
    (barycentric cs:O=1,M=1,D=1) coordinate (Gt)
    (barycentric cs:Gt=3,Ht=-1) coordinate (J)
    (O)++(90:1cm) coordinate (I1)
    (intersection of O--I1 and M--N) coordinate (I)
    (intersection of A--I and B--C) coordinate (K)
    ($(I)+(B)-(O)$) coordinate (P1)
    (intersection of A--B and P1--I) coordinate (P)
    (intersection of A--C and P1--I) coordinate (Q)
    ;
    \pgfresetboundingbox;
    \draw (A)--(B)--(C)--cycle (A)--(D) (A)--(O) (O)--(M) (O)--(N) (P)--(Q) (M)--(N) (A)--(D) (I)--(O)(P)--(O) (M)--(O)(M)--(D)(A)--(K);
    \draw[name path = dtron1] let \p1=($(J) - (O)$) in (J) circle ({veclen(\x1,\y1)});
  \filldraw[black] (D) circle (1pt) node[below left] {$K$};
     \filldraw[black] (C) circle (1pt) node[below right] {$I$};
    \filldraw[black] (P) circle (1pt) node[above left] {$L$};
    \filldraw[black] (K) circle (1pt) node[below] {$V$};
   \filldraw[black] (B) circle (1pt) node[below left]{$H$};
    \filldraw[black] (A) circle (1pt) node[above left] {$C$};
  \filldraw[black] (N) circle (1pt) node[above right] {$P$};
     \filldraw[black] (Q) circle (1pt) node[right] {$O$};
  \filldraw[black] (M) circle (1pt) node[below left] {$S$};
  \filldraw[black] (I) circle (1pt) node[above right] {$W$};
  \filldraw[black] (O) circle (1pt) node[below right] {$Y$};
    \foreach \x/\y/\z in {A/D/C,I/O/C,O/M/A,O/N/A}{\draw pic[draw,angle radius=2mm]{right angle=\x--\y--\z};}

\end{tikzpicture}
}
\end{minipage}

\end{tcolorbox}
\vspace{4pt}


\begin{tcolorbox}[
    breakable,
    colback=white,
    colframe=blue!60!black,
    boxrule=0.8pt,
    arc=3pt,
    title={\textbf{Câu \#448} \quad \textbf{MÃ CÂU: \texttt{[L09\_C9\_B5\_08]}} \quad | \quad \textbf{Nhóm: \texttt{HINH\_342}} \quad | \quad \textbf{MapID: \texttt{9H9V3-1}} \hfill \textcolor{yellow!90!white}{[Hình gốc .pyc]}}
]
\textbf{Dạng toán:} Tổng hợp 8 \hfill \textit{(Tự luận -- Vận dụng)}
\vspace{4pt}\hrule\vspace{6pt}
\begin{minipage}[t]{0.58\textwidth}
\vspace{0pt}
Cho $({O;R})$ đường kính ${AC}$. Bán kính ${DO}$ vuông góc với ${AC}$, ${L}$ là điểm bất kì trên cung nhỏ ${AD}$ (${L}$ khác ${A}$ và ${D}$), ${LC}$ cắt ${AD}$ tại ${I}$. Gọi ${T}$ là hình chiếu của ${I}$ trên ${AC}$.

 Chứng minh bốn điểm ${D, C, I, T}$ cùng thuộc một đường tròn.

 Chứng minh ${DA}$ là phân giác ${LDT}$.

 Kẻ ${A}x$ là tiếp tuyến của nửa đường tròn tại ${A}$. Lấy ${P \in Ax}$ sao cho $\dfrac{AP \cdot LC}{LA} = R$. Chứng minh ${PC}$ đi qua trung điểm của đoạn thẳng ${IT}$.
\end{minipage}\hfill
\begin{minipage}[t]{0.39\textwidth}
\vspace{0pt}
\centering

\resizebox{\linewidth}{!}{%
\begin{tikzpicture}[line cap=round,line join=round,font=\footnotesize,>=stealth]
    \def\r{3}
    \path 
    (0,0) coordinate (A)
    ++(0:\r) coordinate (O)
    ++(0:\r) coordinate (B)
    (O)++(90:\r) coordinate (C)
    (O)++(130:\r) coordinate (M)
    (A)++(90:1.5*\r) coordinate (x)node[left]{$x$}
    (intersection of A--x and B--M) coordinate (Q)
    (intersection of C--A and B--Q) coordinate (H)
    ($(A)!(H)!(B)$) coordinate (K)
    ($(H)!.5!(K)$) coordinate (F)
    (intersection of B--F and A--Q) coordinate (P)
    ($(H)!.5!(B)$) coordinate (I)
    ;
    \draw[name path = dtron1] let \p1=($(O) - (A)$) in (O) circle ({veclen(\x1,\y1)})
    (C)--(O) (C)--(A) (A)--(x) (C)--(M) (A)--(B)--(P) (B)--(Q) (H)--(K) (C)--(K)
    ;
       \filldraw[black] (A) circle (1pt) node[above left] {$A$};
    \filldraw[black] (B) circle (1pt) node[above right] {$C$};
    \filldraw[black] (C) circle (1pt) node[above] {$D$};
   \filldraw[black] (P) circle (1pt) node[above left]{$P$};
    \filldraw[black] (F) circle (1pt) node[above right] {$K$};
     \filldraw[black] (Q) circle (1pt) node[above left] {$R$};
 \filldraw[black] (M) circle (1pt) node[above] {$L$};
    \filldraw[black] (I) circle (1pt) node[above] {$U$};
    \filldraw[black] (H) circle (1pt) node[above left] {$I$};
   \filldraw[black] (K) circle (1pt) node[below]{$T$};
    \filldraw[black] (O) circle (1pt) node[below ] {$O$};
    \draw pic[draw, angle radius=2mm]{right angle=H--K--B};
\end{tikzpicture}
}
\end{minipage}

\end{tcolorbox}
\vspace{4pt}


\begin{tcolorbox}[
    breakable,
    colback=white,
    colframe=blue!60!black,
    boxrule=0.8pt,
    arc=3pt,
    title={\textbf{Câu \#449} \quad \textbf{MÃ CÂU: \texttt{[L09\_C9\_B5\_09]}} \quad | \quad \textbf{Nhóm: \texttt{HINH\_343}} \quad | \quad \textbf{MapID: \texttt{9H9V3-1}} \hfill \textcolor{yellow!90!white}{[Hình gốc .pyc]}}
]
\textbf{Dạng toán:} Tổng hợp 9 \hfill \textit{(Tự luận -- Vận dụng)}
\vspace{4pt}\hrule\vspace{6pt}
\begin{minipage}[t]{0.58\textwidth}
\vspace{0pt}
Từ điểm ${A}$ nằm ngoài $({O})$ vẽ hai tiếp tuyến ${AF, AG}$ với đường tròn (${F, G}$ là các tiếp điểm). Kẻ đường kính ${GH}$ của $({O})$.

a) Chứng minh ${FH \parallel AO}$.

b) ${AH}$ cắt $({O})$ tại ${I \ (A,I,H}$ theo thứ tự). Chứng minh rằng ${AF^{2} = AI \cdot AH}$.

c) Vẽ ${FP \perp HG}$ tại ${P}$. Gọi ${K}$ là trung điểm của ${FP}$. Chứng minh ba điểm ${A, K, H}$ thẳng hàng.
\end{minipage}\hfill
\begin{minipage}[t]{0.39\textwidth}
\vspace{0pt}
\centering

\resizebox{\linewidth}{!}{%
\begin{tikzpicture}[line cap=round,line join=round,font=\footnotesize,>=stealth]
    \def\r{3}
    \path 
    (0,0) coordinate (O)
    ++(0:\r) coordinate (M1)
    (O)++(0:3*\r) coordinate (A)
    ;
    \begin{scope}[overlay]
        \draw[name path = dtronO] let \p1=($ (O) - (M1) $) in (O) circle ({veclen(\x1,\y1)});
        \path ($(O)!0.5!(A)$) coordinate (OA);
        \path[name path=dtronT] let \p1=($(OA) - (O)$) in (OA) circle ({veclen(\x1,\y1)});
        \path[name intersections={of=dtronO and dtronT, by={B,C}}];
        \draw (B)--(A)--(C);
    \end{scope}
    \path 
    ($(C)!2!(O)$) coordinate (D)
    ($(C)!(B)!(D)$) coordinate (H)
    (intersection of A--D and B--H) coordinate (I)
    ($(D)!(O)!(I)$) coordinate (H1)
    ($(D)!2!(H1)$) coordinate (E)
    (intersection of O--A and B--C) coordinate (M)
    ;
    \draw (C)--(D)--(B)--cycle (B)--(H) (A)--(D) (B)--(O)--(A);
   
         \filldraw[black] (A) circle (1pt) node[above right] {$A$};
    \filldraw[black] (B) circle (1pt) node[above right] {$F$};
    \filldraw[black] (C) circle (1pt) node[below] {$G$};
   \filldraw[black] (D) circle (1pt) node[above left]{$H$};
    \filldraw[black] (E) circle (1pt) node[above right] {$I$};
     
 \filldraw[black] (M) circle (1pt) node[above right] {$R$};
    \filldraw[black] (I) circle (1pt) node[above] {$K$};
    \filldraw[black] (H) circle (1pt) node[above left] {$P$};

    \filldraw[black] (O) circle (1pt) node[below left ] {$O$};
\end{tikzpicture}
}
\end{minipage}

\end{tcolorbox}
\vspace{4pt}


\begin{tcolorbox}[
    breakable,
    colback=white,
    colframe=blue!60!black,
    boxrule=0.8pt,
    arc=3pt,
    title={\textbf{Câu \#450} \quad \textbf{MÃ CÂU: \texttt{[L09\_C9\_B5\_10]}} \quad | \quad \textbf{Nhóm: \texttt{HINH\_344}} \quad | \quad \textbf{MapID: \texttt{9H9V3-1}} \hfill \textcolor{yellow!90!white}{[Hình gốc .pyc]}}
]
\textbf{Dạng toán:} Tổng hợp 10 \hfill \textit{(Tự luận -- Vận dụng)}
\vspace{4pt}\hrule\vspace{6pt}
\begin{minipage}[t]{0.58\textwidth}
\vspace{0pt}
Cho tam giác ${BFI}$ nhọn với ${BF > BI}$. Các đường cao ${FR, IV}$ cắt nhau tại ${P}$.

a) Chứng minh tứ giác ${BRPV}$ nội tiếp.

b) Gọi ${K}$ là giao điểm của ${BP}$ và ${FI}$. Chứng minh ${KB}$ phân giác của $\widehat {RKV}$.

c) Đường thẳng qua ${K}$ và song song với ${RV}$ cắt ${BF, BI}$ lần lượt tại ${L, M}$. Chứng minh ${K}$ là trung điểm ${LM}$.
\end{minipage}\hfill
\begin{minipage}[t]{0.39\textwidth}
\vspace{0pt}
\centering

\resizebox{\linewidth}{!}{%
\begin{tikzpicture}[declare function={a = 6; goc = 65; b = 0.65*a;},line cap=round,line join=round,font=\footnotesize,>=stealth]
    \path
    (goc:b) coordinate (A)
    (0,0) coordinate (C)
    (a,0) coordinate (B)
    ($(A)!(B)!(C)$) coordinate (M)
    ($(A)!(C)!(B)$) coordinate (N)
    (intersection of B--M and C--N) coordinate (H)
    (intersection of A--H and B--C) coordinate (D)
    ($(N)+(D)-(M)$) coordinate (D1)
    (intersection of N--C and D--D1) coordinate (J)
    (intersection of D--D1 and A--B) coordinate (I)
    ($(B)!.5!(C)$) coordinate (O1)
    ;
    \draw[name path = dtronO] let \p1=($ (O1) - (B) $) in (O1) circle ({veclen(\x1,\y1)});
    \draw (A)--(B)--(C)--cycle (B)--(M) (C)--(N) (I)--(J) (J)--(N) (M)--(D)--(N)--cycle (A)--(D);
         \filldraw[black] (A) circle (1pt) node[above right] {$B$};
    \filldraw[black] (B) circle (1pt) node[above right] {$F$};
    \filldraw[black] (C) circle (1pt) node[above left] {$I$};
   \filldraw[black] (D) circle (1pt) node[below]{$K$};
    \filldraw[black] (N) circle (1pt) node[above right] {$V$};
     
 \filldraw[black] (M) circle (1pt) node[above left] {$R$};
    \filldraw[black] (I) circle (1pt) node[above] {$L$};
    \filldraw[black] (H) circle (1pt) node[below right] {$P$};

    \filldraw[black] (J) circle (1pt) node[below left ] {$M$};

    \foreach \x/\y/\z in {A/D/C,C/N/B,C/M/B}{\draw pic[draw,angle radius=2mm]{right angle=\x--\y--\z};}
\end{tikzpicture}
}
\end{minipage}

\end{tcolorbox}
\vspace{4pt}


\begin{tcolorbox}[
    breakable,
    colback=white,
    colframe=blue!60!black,
    boxrule=0.8pt,
    arc=3pt,
    title={\textbf{Câu \#451} \quad \textbf{MÃ CÂU: \texttt{[L09\_C9\_B5\_11]}} \quad | \quad \textbf{Nhóm: \texttt{HINH\_345}} \quad | \quad \textbf{MapID: \texttt{9H9V3-1}} \hfill \textcolor{yellow!90!white}{[Hình gốc .pyc]}}
]
\textbf{Dạng toán:} Tổng hợp 11 \hfill \textit{(Tự luận -- Vận dụng)}
\vspace{4pt}\hrule\vspace{6pt}
\begin{minipage}[t]{0.58\textwidth}
\vspace{0pt}
Cho tam giác ${ABD}$ nhọn. Đường tròn $({O})$ đường kính ${BD}$ cắt ${AB}$, ${AD}$ lần lượt tại ${J}$ và ${G}$; ${BG}$ cắt ${DJ}$ tại ${N}$, ${AN}$ cắt ${BD}$ tại ${Y}$. Từ ${A}$ kẻ tiếp tuyến ${AM}$, ${AW}$ của đường tròn $(O)$ (${M, W}$ là tiếp điểm).

a) Chứng minh tứ giác ${AJNG}$ nội tiếp. 

b) Chứng minh ${AB.BJ = BY.BD}$, từ đó suy ra ${AB.BJ + AD.DG} = BD^{2}$. 

c) Chứng minh ba điểm ${M}$, ${N}$, ${W}$ thẳng hàng.
\end{minipage}\hfill
\begin{minipage}[t]{0.39\textwidth}
\vspace{0pt}
\centering

\resizebox{\linewidth}{!}{%
\begin{tikzpicture}[declare function={a = 5; goc = 65; b = 0.8*a;}]
    \path
    (goc:b) coordinate (A)
    (0,0) coordinate (C)
    (a,0) coordinate (B)
    ($(B)!(A)!(C)$) coordinate (I)
    ($(A)!(B)!(C)$) coordinate (D)
    ($(A)!(C)!(B)$) coordinate (E)
    ($(B)!.5!(C)$) coordinate (O)
    (intersection of D--B and C--E) coordinate (H)
    ;
    \draw (A)--(B)--(C)--cycle ;
    \begin{scope}[overlay]
        \draw[name path = dtronO] let \p1=($ (O) - (B) $) in (O) circle ({veclen(\x1,\y1)});
        \path ($(O)!0.5!(A)$) coordinate (OA);
        \path[name path=dtronT] let \p1=($(OA) - (O)$) in (OA) circle ({veclen(\x1,\y1)});
        \path[name intersections={of=dtronO and dtronT, by={M,N}}];
        \draw (M)--(A)--(N);
    \end{scope}
    \draw (O)--(M) (O)--(N) (A)--(I) (B)--(D) (C)--(E) (M)--(N) (I)--(M) (I)--(N);
           \filldraw[black] (A) circle (1pt) node[above right] {$A$};
    \filldraw[black] (B) circle (1pt) node[below right] {$B$};
    \filldraw[black] (C) circle (1pt) node[below left] {$D$};
   \filldraw[black] (D) circle (1pt) node[above left]{$G$};
    \filldraw[black] (N) circle (1pt) node[above left] {$W$};
 \filldraw[black] (M) circle (1pt) node[above right] {$M$};
    \filldraw[black] (I) circle (1pt) node[below] {$Y$};
   \filldraw[black] (O) circle (1pt) node[below] {$O$};
    \filldraw[black] (E) circle (1pt) node[above right ] {$J$};
 \filldraw[black] (H) circle (1pt) node[below right] {$N$};

    \foreach \x/\y/\z in {A/N/O,A/M/O,B/D/A,A/E/C,A/I/B}{\draw pic[draw,angle radius=2mm]{right angle=\x--\y--\z};}
\end{tikzpicture}
}
\end{minipage}

\end{tcolorbox}
\vspace{4pt}


\begin{tcolorbox}[
    breakable,
    colback=white,
    colframe=blue!60!black,
    boxrule=0.8pt,
    arc=3pt,
    title={\textbf{Câu \#452} \quad \textbf{MÃ CÂU: \texttt{[L09\_C9\_B5\_12]}} \quad | \quad \textbf{Nhóm: \texttt{HINH\_346}} \quad | \quad \textbf{MapID: \texttt{9H9V3-1}} \hfill \textcolor{yellow!90!white}{[Hình gốc .pyc]}}
]
\textbf{Dạng toán:} Tổng hợp 12 \hfill \textit{(Tự luận -- Vận dụng)}
\vspace{4pt}\hrule\vspace{6pt}
\begin{minipage}[t]{0.58\textwidth}
\vspace{0pt}
Cho đường tròn $({O})$, dây ${IJ}$ cố định. Gọi ${E}$ là điểm chính giữa cung nhỏ ${IJ}$, kẻ đường kính ${DE}$ cắt ${IJ}$ tại ${S}$. Lấy điểm ${M}$ bất kỳ trên cung lớn ${IJ}$, ${ME}$ cắt ${IJ}$ tại ${K}$. Đường thẳng ${DM}$ cắt đường thẳng ${IJ}$ tại ${T}$.

a) Chứng minh: Tứ giác ${TMSE}$ nội tiếp.

b) Chứng minh: ${DM \cdot DT = DS \cdot DE}$.

c) Gọi ${W}$ là giao điểm của đường thẳng ${DK}$ và ${ET}$. Kẻ ${WL \perp DE}$ cắt ${DE}$ tại ${L}$, cắt đường tròn $(O)$ tại ${Q}$. Chứng minh ${Q, S, M}$ thẳng hàng.
\end{minipage}\hfill
\begin{minipage}[t]{0.39\textwidth}
\vspace{0pt}
\centering

\resizebox{\linewidth}{!}{%
\begin{tikzpicture}[line cap=round,line join=round,font=\footnotesize,>=stealth,,declare function={r=3;}]
    \path 
    (0,0) coordinate (O)
    ++(-30:r) coordinate (D)
    (O)++(-150:r) coordinate (C)
    (O)++(-90:r) coordinate (B)
    (O)++(-50:r) coordinate (N)
    ($(O)!(N)!(B)$) coordinate (M)
    ($(N)!2!(M)$) coordinate (K)
    (O)++(90:r) coordinate (A)
    (intersection of A--B and C--D) coordinate (I)
    ($(I)!(O)!(N)$) coordinate (H1)
    ($(N)!2!(H1)$) coordinate (H)
    (intersection of A--H and K--B) coordinate (P)
    ($(P)!(A)!(B)$) coordinate (K)
    (intersection of A--K and B--H) coordinate (E)
    ;
    \draw (C)--(D) (O) circle(r) (C)--(D) (K)--(N) (K)--(A)--(P)--(B) (A)--(B) (H)--(N) (H)--(B)(P)--(C)(I)--(K);
         \filldraw[black] (A) circle (1pt) node[above right] {$D$};
    \filldraw[black] (B) circle (1pt) node[below right] {$E$};
    \filldraw[black] (C) circle (1pt) node[below left] {$I$};
   \filldraw[black] (D) circle (1pt) node[right]{$J$};
    \filldraw[black] (N) circle (1pt) node[below right] {$Q$};
 \filldraw[black] (M) circle (1pt) node[above right] {$L$};
    \filldraw[black] (I) circle (1pt) node[above right] {$S$};
   \filldraw[black] (O) circle (1pt) node[right] {$O$};
    \filldraw[black] (E) circle (1pt) node[above right ] {$K$};
 \filldraw[black] (H) circle (1pt) node[left] {$M$};
 \filldraw[black] (P) circle (1pt) node[below left] {$T$};
 \filldraw[black] (K) circle (1pt) node[below left] {$W$};
    \foreach \x/\y/\z in {A/I/P,A/M/N}{\draw pic[draw,angle radius=2mm]{right angle=\x--\y--\z};}
\end{tikzpicture}
}
\end{minipage}

\end{tcolorbox}
\vspace{4pt}


\begin{tcolorbox}[
    breakable,
    colback=white,
    colframe=blue!60!black,
    boxrule=0.8pt,
    arc=3pt,
    title={\textbf{Câu \#453} \quad \textbf{MÃ CÂU: \texttt{[L09\_C9\_B5\_13]}} \quad | \quad \textbf{Nhóm: \texttt{HINH\_347}} \quad | \quad \textbf{MapID: \texttt{9H9V3-1}} \hfill \textcolor{yellow!90!white}{[Hình gốc .pyc]}}
]
\textbf{Dạng toán:} Tổng hợp 13 \hfill \textit{(Tự luận -- Vận dụng)}
\vspace{4pt}\hrule\vspace{6pt}
\begin{minipage}[t]{0.58\textwidth}
\vspace{0pt}
Cho đường tròn tâm $(O)$ và dây ${FG}$ cố định không đi qua ${O}$. Trên cung lớn ${FG}$ lấy điểm ${E}$ sao cho ${EF < EG}$. Kẻ đường kính ${EV}$, ${I}$ là hình chiếu của ${G}$ trên ${EV}$, ${R}$ là trung điểm của ${FG}$.

a) Chứng minh bốn điểm ${G, I, R, O}$ cùng thuộc một đường tròn.

b) Kẻ ${EH \perp FG}$ tại ${H}$. Chứng minh ${EH.EV = EF.EG}$ và $\triangle {RHI}$ cân.

c) Gọi ${X}$ là hình chiếu của ${F}$ trên ${EV}$. Chứng minh khi ${E}$ di chuyển trên cung lớn ${FG}$ thì tâm đường tròn ngoại tiếp $\triangle {HIX}$ là một điểm cố định.
\end{minipage}\hfill
\begin{minipage}[t]{0.39\textwidth}
\vspace{0pt}
\centering

\resizebox{\linewidth}{!}{%
\begin{tikzpicture}[declare function={a = 5; goc = 65; b = 0.65*a;},font=\footnotesize,line cap=round,line join=round,>=stealth]
    \path
    (goc:b) coordinate (A)
    (0,0) coordinate (B)
    (a,0) coordinate (C)
    ($(B)!(A)!(C)$) coordinate (D)
    ($(A)!.5!(C)$) coordinate (Q)
    ($(B)!(A)!(C)$) coordinate (AH)
    ($(A)!(B)!(C)$) coordinate (BH)
    (intersection of A--AH and B--BH) coordinate (Ht)
    (barycentric cs:A=1,B=1,C=1) coordinate (Gt)
    (barycentric cs:Gt=3,Ht=-1) coordinate (O)
    ($(A)!2!(O)$) coordinate (K)
    ($(A)!(B)!(O)$) coordinate (F)
    ($(A)!(C)!(O)$) coordinate (E)
    ($(B)!.5!(C)$) coordinate (M)
    ($(B)!.5!(O)$) coordinate (P)
    ($(C)!.5!(O)$) coordinate (I)
    ;
    \draw[name path = dtron1] let \p1=($(O) - (A)$) in (O) circle ({veclen(\x1,\y1)});
    \draw (A)--(B)--(C)--cycle (B)--(F) (B)--(O)--(C) (B)--(F) (C)--(E) (A)--(K)(A)--(D) (B)--(K)(F)--(M)--(I)--(E)--(M) (D)--(E);
    
   
        \filldraw[black] (A) circle (1pt) node[above right] {$E$};
    \filldraw[black] (B) circle (1pt) node[below right] {$F$};
    \filldraw[black] (C) circle (1pt) node[below right] {$G$};
   \filldraw[black] (D) circle (1pt) node[below]{$H$};
    \filldraw[black] (Q) circle (1pt) node[above right] {$U$};
 \filldraw[black] (M) circle (1pt) node[below left] {$R$};
    \filldraw[black] (I) circle (1pt) node[above right] {$N$};
   \filldraw[black] (O) circle (1pt) node[right] {$O$};
    \filldraw[black] (E) circle (1pt) node[below left ] {$I$};
 \filldraw[black] (F) circle (1pt) node[above right] {$X$};
\filldraw[black] (P) circle (1pt) node[above left] {$T$};
\filldraw[black] (K) circle (1pt) node[below left] {$V$};

    \foreach \x/\y/\z in {A/D/C,B/F/K,C/E/F}{\draw pic[draw,angle radius=2mm]{right angle=\x--\y--\z};}
\end{tikzpicture}
}
\end{minipage}

\end{tcolorbox}
\vspace{4pt}


\begin{tcolorbox}[
    breakable,
    colback=white,
    colframe=blue!60!black,
    boxrule=0.8pt,
    arc=3pt,
    title={\textbf{Câu \#454} \quad \textbf{MÃ CÂU: \texttt{[L09\_C9\_B5\_14]}} \quad | \quad \textbf{Nhóm: \texttt{HINH\_348}} \quad | \quad \textbf{MapID: \texttt{9H9V3-1}} \hfill \textcolor{yellow!90!white}{[Hình gốc .pyc]}}
]
\textbf{Dạng toán:} Tổng hợp 14 \hfill \textit{(Tự luận -- Vận dụng)}
\vspace{4pt}\hrule\vspace{6pt}
\begin{minipage}[t]{0.58\textwidth}
\vspace{0pt}
Tứ giác ${FJKL}$ nội tiếp đường tròn tâm ${O}$ đường kính ${FL}$. Hai đường chéo ${FK, JL}$ cắt nhau tại ${N}$. Từ ${N}$ kẻ ${NP}$ vuông góc với ${FL}$ (${P \in FL}$). Đường thẳng ${KP}$ cắt đường tròn tại điểm thứ hai là ${S}$. Giao điểm của ${JL}$ và ${KP}$ là ${V}$. Chứng minh:

a) ${KNPL}$ là tứ giác nội tiếp.

b) ${PF}$ là tia phân giác của $\widehat {JPS}$.

c) ${JN \cdot LV = NV \cdot JL}$.
\end{minipage}\hfill
\begin{minipage}[t]{0.39\textwidth}
\vspace{0pt}
\centering

\resizebox{\linewidth}{!}{%
\begin{tikzpicture}[declare function={r = 3; goc = 65; b = 0.65*a;},font=\footnotesize,line cap=round,line join=round,>=stealth]
    \path 
    (0,0) coordinate (A)
    ++(0:r) coordinate (O)
    ++(0:r) coordinate (D)
    (O)++(70:r) coordinate (C)
    (O)++(130:r) coordinate (B)
    (intersection of A--C and B--D) coordinate (E)
    ($(A)!(E)!(D)$) coordinate (F)
    ($(F)!(O)!(C)$) coordinate (M1)
    ($(C)!2!(M1)$) coordinate (M)
    (intersection of B--D and C--F) coordinate (N)
    ($(E)!.5!(D)$) coordinate (I)
    ;
    \draw (O) circle(r);
    \draw (A)--(B)--(C)--(D)--cycle (A)--(C)--(M) (F)--(B)--(D)(A)--(D)(E)--(F)--(I)--(C);
    \foreach \x/\y/\z in {E/F/D,A/C/D,A/B/D}{\draw pic[draw,angle radius=2mm]{right angle=\x--\y--\z};}
        \filldraw[black] (A) circle (1pt) node[above left] {$F$};
    \filldraw[black] (B) circle (1pt) node[above] {$J$};
    \filldraw[black] (C) circle (1pt) node[above right] {$K$};
   \filldraw[black] (D) circle (1pt) node[below left]{$L$};
 \filldraw[black] (M) circle (1pt) node[below left] {$S$};
    \filldraw[black] (I) circle (1pt) node[above right] {$W$};
   \filldraw[black] (O) circle (1pt) node[below] {$O$};
    \filldraw[black] (E) circle (1pt) node[above] {$N$};
 \filldraw[black] (F) circle (1pt) node[below] {$P$};
\filldraw[black] (N) circle (1pt) node[left] {$V$};
\end{tikzpicture}
}
\end{minipage}

\end{tcolorbox}
\vspace{4pt}


\begin{tcolorbox}[
    breakable,
    colback=white,
    colframe=blue!60!black,
    boxrule=0.8pt,
    arc=3pt,
    title={\textbf{Câu \#455} \quad \textbf{MÃ CÂU: \texttt{[L09\_C9\_B5\_15]}} \quad | \quad \textbf{Nhóm: \texttt{HINH\_349}} \quad | \quad \textbf{MapID: \texttt{9H9V3-1}} \hfill \textcolor{yellow!90!white}{[Hình gốc .pyc]}}
]
\textbf{Dạng toán:} Tổng hợp 15 \hfill \textit{(Tự luận -- Vận dụng)}
\vspace{4pt}\hrule\vspace{6pt}
\begin{minipage}[t]{0.58\textwidth}
\vspace{0pt}
Trên nửa đường tròn ${(O, R)}$ đường kính ${GJ}$, lấy điểm ${D}$ sao cho ${GD = R}$.

a) Chứng minh tam giác ${DGJ}$ vuông tại ${D}$ và tính số đo các góc ${G, J}$ của tam giác vuông ${DGJ}$.

b) Qua ${G}$ kẻ tiếp tuyến với nửa đường tròn $(O)$, nó cắt tia ${JD}$ tại ${K}$. Qua ${K}$ kẻ tiếp tuyến ${KM}$ với nửa đường tròn ${(O)}$ (${M}$ là tiếp điểm). Gọi ${Z}$ là giao điểm của ${OK}$ và ${GM}$. Chứng minh rằng ${OK \perp GM}$ và ${KZ \cdot KO = KD \cdot KJ}$.

c) Kẻ ${MY}$ vuông góc với ${GJ}$ tại ${Y}$. ${MY}$ cắt ${JK}$ tại ${W}$. Chứng minh ${ZW}$ song song với ${GJ}$.
\end{minipage}\hfill
\begin{minipage}[t]{0.39\textwidth}
\vspace{0pt}
\centering

\resizebox{\linewidth}{!}{%
\begin{tikzpicture}[line cap=round,line join=round,font=\footnotesize,>=stealth,declare function={r = 3;}]
    \path 
    (0,0) coordinate (B)
    ++(0:2*r) coordinate (C)
    ($(B)!.5!(C)$) coordinate (O)
    (O)++(120:r) coordinate (A)
    (B)++(90:r) coordinate (B1)
    (intersection of C--A and B--B1) coordinate (D)
    ($(B)!2!(D)$) coordinate (F)
    ($(F)!(O)!(C)$) coordinate (O1)
    ($(C)!2!(O1)$) coordinate (E)
    ($(B)!(E)!(C)$) coordinate (H)
    (intersection of E--B and O--D) coordinate (I)
    (intersection of E--H and C--D) coordinate (G)
    ;
    \draw (A)--(B)--(C)--cycle (B)--(F)--(C) (D)--(E)--(H) (E)--(B) (E)--(O) (I)--(G) (A)--(O)--(D)
    ;
    \draw[name path = dtronO] let \p1=($ (O) - (B) $) in (O) circle ({veclen(\x1,\y1)});
      \filldraw[black] (A) circle (1pt) node[above left] {$D$};
    \filldraw[black] (B) circle (1pt) node[above left] {$G$};
    \filldraw[black] (C) circle (1pt) node[above right] {$J$};
    \filldraw[black] (D) circle (1pt) node[below left]{$K$};
  \filldraw[black] (F) circle (1pt) node[below left] {$S$};
    \filldraw[black] (I) circle (1pt) node[below] {$Z$};
   \filldraw[black] (O) circle (1pt) node[below] {$O$};
     \filldraw[black] (E) circle (1pt) node[above] {$M$};
 \filldraw[black] (G) circle (1pt) node[above right] {$W$};
 \filldraw[black] (H) circle (1pt) node[below] {$Y$};
\end{tikzpicture}
}
\end{minipage}

\end{tcolorbox}
\vspace{4pt}


\begin{tcolorbox}[
    breakable,
    colback=white,
    colframe=blue!60!black,
    boxrule=0.8pt,
    arc=3pt,
    title={\textbf{Câu \#456} \quad \textbf{MÃ CÂU: \texttt{[L09\_C9\_B5\_16]}} \quad | \quad \textbf{Nhóm: \texttt{HINH\_350}} \quad | \quad \textbf{MapID: \texttt{9H9V3-1}} \hfill \textcolor{yellow!90!white}{[Hình gốc .pyc]}}
]
\textbf{Dạng toán:} Tổng hợp 16 \hfill \textit{(Tự luận -- Vận dụng)}
\vspace{4pt}\hrule\vspace{6pt}
\begin{minipage}[t]{0.58\textwidth}
\vspace{0pt}
Cho đường tròn $(O;R)$, đường kính ${BC}$ vuông góc với dây ${DF}$ tại điểm ${I}$ (${I}$ nằm giữa ${B}$ và ${O}$). Lấy điểm ${G}$ bất kỳ trên cung nhỏ ${CD}$ (${G}$ khác ${C}$ và ${D}$). ${BG}$ cắt ${DF}$ tại ${Z}$.

a) Chứng minh bốn điểm ${Z, G, C, I}$ cùng thuộc một đường tròn.

b) Chứng minh ${BZ \cdot BG = BI \cdot BC}$.

c) Gọi ${T}$ là giao điểm của tia ${CG}$ và tia ${FD}$, ${W}$ là giao điểm của ${BT}$ và ${CZ}$. Chứng minh ${IZ}$ là phân giác của $\widehat {GIW}$. Chứng minh ${OW}$ là tiếp tuyến của đường tròn ngoại tiếp tam giác ${TWG}$.
\end{minipage}\hfill
\begin{minipage}[t]{0.39\textwidth}
\vspace{0pt}
\centering

\resizebox{\linewidth}{!}{%
\begin{tikzpicture}[line cap=round,line join=round,font=\footnotesize,>=stealth,declare function={r = 3;}]
    \path 
    (0,0) coordinate (A)
    ++(0:2*r) coordinate (B)
    ($(A)!0.5!(B)$) coordinate (O)
    (O)++(110:r) coordinate (D)
    (O)++(-110:r) coordinate (C)
    (O)++(-50:r) coordinate (E)
    (intersection of B--E and C--D) coordinate (P)
    ($(A)!(O)!(P)$) coordinate (Q1)
    ($(A)!2!(Q1)$) coordinate (Q)
    (intersection of C--D and A--B) coordinate (I)
    (intersection of B--Q and A--E) coordinate (K)
    ;
    \draw[name path = dtronO] let \p1=($ (O) - (B) $) in (O) circle ({veclen(\x1,\y1)});
    \draw (A)--(P)--(B)--cycle (P)--(D)(Q)--(I) (Q)--(O)(A)--(E)(B)--(Q)(I)--(E);
   
         \filldraw[black] (A) circle (1pt) node[above left] {$B$};
    \filldraw[black] (B) circle (1pt) node[above right] {$C$};
    \filldraw[black] (C) circle (1pt) node[above right] {$D$};
    \filldraw[black] (D) circle (1pt) node[above]{$F$};
  \filldraw[black] (Q) circle (1pt) node[below left] {$W$};
    \filldraw[black] (I) circle (1pt) node[above right] {$I$};
   \filldraw[black] (O) circle (1pt) node[above] {$O$};
     \filldraw[black] (E) circle (1pt) node[below right] {$G$};
 \filldraw[black] (P) circle (1pt) node[below] {$T$};
 \filldraw[black] (K) circle (1pt) node[below left] {$Z$};
\end{tikzpicture}
}
\end{minipage}

\end{tcolorbox}
\vspace{4pt}


\end{document}
